\documentclass[11pt,a4paper]{article}
\usepackage[T2A]{fontenc}
\usepackage[utf8]{inputenc}
\usepackage[english,russian]{babel}
\usepackage{amsmath,amssymb,amsthm,mathtools}
\usepackage[left=1.9cm,right=1.9cm,top=1.8cm,bottom=1.9cm]{geometry}
\usepackage{enumitem}
\usepackage[most]{tcolorbox}
\usepackage[unicode,hidelinks]{hyperref}
\usepackage{titlesec}
\setlist{nosep,leftmargin=1.5em}
\setlength{\parindent}{0pt}
\setlength{\parskip}{3.5pt}
\allowdisplaybreaks
\titleformat{\section}{\Large\bfseries\color{blue!40!black}}{}{0em}{}
\titleformat{\subsection}{\large\bfseries\color{blue!30!black}}{\thesubsection.}{0.5em}{}
\titlespacing*{\section}{0pt}{14pt}{6pt}

% ---------- ticket headings ----------
% \ticket{N}{Название билета}   — обычный билет
% \ticketx{41*}{Название}       — вопрос со звёздочкой (доп. материал)
\newcommand{\ticket}[2]{\clearpage\setcounter{subsection}{0}\renewcommand{\thesubsection}{#1.\arabic{subsection}}%
  \section[Билет #1. #2]{Билет #1.\quad #2}}
\newcommand{\ticketx}[2]{\clearpage\setcounter{subsection}{0}\renewcommand{\thesubsection}{#1.\arabic{subsection}}%
  \section[Вопрос #1 #2]{Вопрос #1\quad #2}}
\newcommand{\prelim}[2]{\clearpage\setcounter{subsection}{0}\renewcommand{\thesubsection}{#1.\arabic{subsection}}%
  \section[#2]{#2}}

% ---------- boxes ----------
% \begin{defs} ... \end{defs}              — рамка «Определения»
% \begin{thmbox}{Теорема 13.1 (принцип симметрии)} ... \end{thmbox}  — формулировка
% \begin{detail}{Заголовок} ... \end{detail}  — восполненная техническая деталь
% \begin{remarkbox}{Заголовок} ... \end{remarkbox} — замечания/примеры
\newtcolorbox{defs}[1][Определения]{breakable,colback=blue!3,colframe=blue!45!black,boxrule=0.5pt,arc=1mm,
  left=4pt,right=4pt,top=3pt,bottom=3pt,title={#1},fonttitle=\bfseries\small}
\newtcolorbox{thmbox}[1]{breakable,colback=orange!4,colframe=orange!60!black,boxrule=0.5pt,arc=1mm,
  left=4pt,right=4pt,top=3pt,bottom=3pt,title={#1},fonttitle=\bfseries\small}
\newtcolorbox{detail}[1]{breakable,colback=green!3,colframe=green!35!black,boxrule=0.4pt,arc=1mm,
  left=4pt,right=4pt,top=3pt,bottom=3pt,title={\textit{Деталь:} #1},fonttitle=\bfseries\small}
\newtcolorbox{remarkbox}[1]{breakable,colback=gray!4,colframe=gray!55!black,boxrule=0.4pt,arc=1mm,
  left=4pt,right=4pt,top=3pt,bottom=3pt,title={#1},fonttitle=\bfseries\small}

% ---------- macros ----------
\newcommand{\CC}{\mathbb C}
\newcommand{\RR}{\mathbb R}
\newcommand{\NN}{\mathbb N}
\newcommand{\ZZ}{\mathbb Z}
\newcommand{\QQ}{\mathbb Q}
\newcommand{\DD}{\mathbb D}          % единичный круг
\newcommand{\TT}{\mathbb T}          % единичная окружность
\newcommand{\Ch}{\widehat{\mathbb C}} % расширенная плоскость
\newcommand{\Sc}{\mathcal S}         % класс S
\newcommand{\eps}{\varepsilon}
\newcommand{\dist}{\operatorname{dist}}
\newcommand{\diam}{\operatorname{diam}}
\newcommand{\Ind}{\operatorname{Ind}}
\newcommand{\Res}{\operatorname*{Res}}
\newcommand{\Arg}{\operatorname{Arg}}
\newcommand{\Ln}{\operatorname{Ln}}
\newcommand{\Arcsin}{\operatorname{Arcsin}}
\newcommand{\Arctg}{\operatorname{Arctg}}
\newcommand{\Th}{\operatorname{th}}
\newcommand{\CAut}{\operatorname{CAut}}
\newcommand{\Zh}{\mathop{\text{Ж}}\nolimits}
\newcommand{\clos}[1]{\overline{#1}}
\newcommand{\ord}{\operatorname{ord}}
\renewcommand{\Re}{\operatorname{Re}}
\renewcommand{\Im}{\operatorname{Im}}
\renewcommand{\le}{\leqslant}
\renewcommand{\ge}{\geqslant}
\renewcommand{\leq}{\leqslant}
\renewcommand{\geq}{\geqslant}
\title{\textbf{Теория функций комплексного переменного}\\[4pt]
\large 3 курс, весенний семестр 2025/2026, лектор К.\,Ю.~Федоровский\\[2pt]
\large Подробные ответы на билеты к экзамену}
\date{}
\begin{document}
\maketitle
\thispagestyle{empty}
\begin{center}
\begin{minipage}{0.88\textwidth}\small
Текст составлен по слайдам лекций (главы 12--17) и программе экзамена. Каждый билет устроен одинаково:
\begin{itemize}
\item рамка \textbf{«Определения»} --- всё, что нужно уметь точно сформулировать;
\item \textbf{формулировки} теорем --- в оранжевых рамках, с \emph{нумерацией лекций} (Предложение 13.4, Теорема 16.3 и т.\,п.), чтобы удобно сверяться со слайдами;
\item \textbf{полные доказательства}; места, которые на лекциях оставлены как «упражнение», «легко видеть», «проверка деталей оставляется читателю», расписаны и, как правило, выделены зелёными рамками \textbf{«Деталь»};
\item серые рамки --- замечания, примеры, найденные опечатки слайдов и краткий \textbf{план ответа} на экзамене.
\end{itemize}
Раздел 0 содержит необходимые сведения из первого семестра и главу 12 (лемма Шварца, формула Йенсена, лемма Бореля--Каратеодори и др.), на которые опираются билеты. Вопросы 41*--43* в билеты не входят, но могут быть заданы как дополнительные; они помещены рядом с соответствующими темами.
\end{minipage}
\end{center}
\tableofcontents
% ===== part 1a =====
% ======================================================================
% ЧАСТЬ 1a. Предварительные сведения и глава 12
% ======================================================================
\prelim{0}{Предварительные сведения}

Этот раздел содержит: (1) обозначения и соглашения; (2) краткую сводку фактов первого семестра, на которые
опирается весенний курс (с точными формулировками и короткими доказательствами или идеями доказательств);
(3) полное изложение главы~12 лекций (слайды 258--272): лемма Шварца и её следствия, произведения Бляшке,
формула Йенсена, теорема Каратеодори об аппроксимации, лемма Бореля--Каратеодори и малая теорема Пикара
для целых функций класса $\mathcal E$. Нумерация утверждений главы~12 совпадает с нумерацией на слайдах;
факты первого семестра занумерованы как <<Факт 0.$k$>>.

\subsection{Обозначения и соглашения}

\begin{defs}[Основные обозначения]
\begin{itemize}
\item $B(a,r)=\{z:|z-a|<r\}$ --- открытый круг, $\clos{B}(a,r)$ --- замкнутый круг, $C(a,r)=\{z:|z-a|=r\}$ ---
окружность; $\DD=B(0,1)$, $\TT=C(0,1)$; $\CC_+=\{\Im z>0\}$, $\CC_-=\{\Im z<0\}$.
\item $\Ch=\CC\cup\{\infty\}$ --- расширенная комплексная плоскость (сфера Римана) со сферической метрикой;
окрестности $\infty$ --- множества $\{|z|>R\}\cup\{\infty\}$. Для $E\subset\Ch$ через $\partial^\sharp E$
обозначается граница $E$ в $\Ch$ (например, $\partial^\sharp\CC=\{\infty\}$, $\partial^\sharp\Ch=\varnothing$).
\item $U(a)$ --- (какая-то) окрестность точки $a$, $V(a)=U(a)\setminus\{a\}$ --- проколотая окрестность.
\item \emph{Область} --- непустое открытое связное множество. $H(D)$ --- класс функций, голоморфных в области
$D\subset\Ch$; $H(a)$ --- функции, голоморфные в некоторой окрестности точки $a$; $H(E)$ для произвольного
$E$ --- функции, голоморфные в некоторой (своей) окрестности $E$. Функции из $H(\CC)$ называются \emph{целыми}.
\item \emph{Голоморфность в $\infty$}: $f$ определена и конечна в $U(\infty)$ и $g(w)=f(1/w)$, $g(0)=f(\infty)$,
голоморфна в точке $0$. Отображение $f$ со значениями в $\Ch$ \emph{голоморфно (мероморфно) в точке $a$},
если в подходящих картах ($z\mapsto z$ или $z\mapsto 1/z$ в прообразе и образе) оно голоморфно.
\item \emph{Конформность.} $f$ конформна в точке $a\in\CC$, если $f$ $\CC$-дифференцируема в $a$ и $f'(a)\ne0$;
в точке $\infty$ и в точках, где $f=\infty$, --- в смысле карт $1/z$. $\mathrm{Conf}(D)$ --- функции,
локально конформные и взаимно однозначные в $D$; $\mathrm{Conf}(D,G)$ --- функции $f\in\mathrm{Conf}(D)$ с $f(D)=G$
(конформные изоморфизмы $D$ на $G$); $\CAut(D)=\mathrm{Conf}(D,D)$.
\item $\Lambda$ --- множество дробно-линейных отображений (ДЛО) $z\mapsto\frac{az+b}{cz+d}$, $ad-bc\ne0$.
Для $a\in\DD$ положим
\[
\varphi_a(z)=\frac{z-a}{1-\bar a z}\qquad\text{(элементарный множитель Бляшке).}
\]
\item $f_n\to f$ \emph{локально равномерно} в $D$, если $f_n\to f$ равномерно на каждом компакте $K\subset D$;
$\|f\|_K=\sup_K|f|$.
\item $\mathfrak z(f,D)$ --- число нулей $f$ в $D$ с учётом кратности; $\mathfrak p(f,D)$ --- число полюсов с учётом
кратности.
\item \emph{Жорданова кривая} --- образ окружности при непрерывном инъективном отображении в $\CC$
(т.е. путь $\gamma:[0,1]\to\CC$, $\gamma(s)=\gamma(t)\Leftrightarrow s=t$ или $\{s,t\}=\{0,1\}$).
$\operatorname{Int}\Gamma$ --- ограниченная компонента $\CC\setminus\Gamma$; \emph{жорданова область} ---
область вида $\operatorname{Int}\Gamma$.
\item \emph{Допустимая область (Д.О.)} --- область $D=D_0\setminus(\clos{D_1}\cup\dots\cup\clos{D_M})$, где
$D_0,\dots,D_M$ --- жордановы области, $\clos{D_m}\subset D_0$, $\clos{D_m}\cap\clos{D_k}=\varnothing$ ($m\ne k$);
$\partial^+D=\partial^+D_0-\sum_{m\ge1}\partial^+D_m$ (положительная ориентация: область остаётся слева).
\end{itemize}
\end{defs}

\subsection{Сводка фактов первого семестра}

\begin{thmbox}{Факт 0.1 (лемма Гурса, интегральная теорема Коши, теорема Мореры)}
\textbf{(а) Лемма Гурса.} Если $f\in H(G)$, $G\subset\CC$ открыто, $\Delta\subset G$ --- замкнутый треугольник,
то $\int_{\partial\Delta}f\,dz=0$.\\
\textbf{(б) Теорема Коши для выпуклой области.} Если $G$ --- выпуклое открытое множество, $f\in H(G)$, то $f$
имеет в $G$ первообразную $F(z)=\int_{[z_0,z]}f(\zeta)\,d\zeta$, и $\int_\gamma f\,dz=0$ для любого замкнутого
кусочно-гладкого пути $\gamma$ в $G$.\\
\textbf{(в) Теорема Коши для Д.О.} Если $D$ --- Д.О. со спрямляемой границей, $f\in H(D)\cap C(\clos D)$, то
$\int_{\partial^+D}f\,dz=0$. Более общо (гомологическая форма): если $\gamma$ --- замкнутый путь в области $D$ и
$\Ind(\gamma,a)=0$ для всех $a\notin D$, то $\int_\gamma f\,dz=0$ для всех $f\in H(D)$.\\
\textbf{(г) Теорема Мореры.} Если $f\in C(G)$ и $\int_{\partial\Delta}f\,dz=0$ для любого замкнутого треугольника
$\Delta\subset G$, то $f\in H(G)$.
\end{thmbox}
\textit{Идея.} (а) деление треугольника на 4 подобных и оценка через $\CC$-дифференцируемость в предельной точке;
(б) из (а): $F'=f$, интеграл от производной по замкнутому пути равен нулю; (г) в каждом круге $B\subset G$
функция $F(z)=\int_{[a,z]}f$ корректно задаёт первообразную ($F'=f$), а производная голоморфной функции
голоморфна (Факт 0.2).

\begin{thmbox}{Факт 0.2 (интегральная формула Коши, ряд Тейлора, неравенства Коши, теорема Лиувилля)}
Пусть $f\in H(G)$, $\clos B(a,R)\subset G$. Тогда при $|z-a|<R$
\[
f(z)=\frac1{2\pi i}\int_{C(a,R)}\frac{f(\zeta)\,d\zeta}{\zeta-z},\qquad
f^{(n)}(z)=\frac{n!}{2\pi i}\int_{C(a,R)}\frac{f(\zeta)\,d\zeta}{(\zeta-z)^{n+1}},
\]
$f(z)=\sum_{n\ge0}c_n(z-a)^n$ в $B(a,R)$ (и в наибольшем круге с центром $a$, лежащем в $G$), $c_n=f^{(n)}(a)/n!$, и
справедливы \emph{неравенства Коши}
\[
|c_n|\le \frac{M(r)}{r^n},\qquad M(r)=\max_{|z-a|=r}|f(z)|,\quad 0<r<R .
\]
\textbf{Теорема Лиувилля.} Ограниченная целая функция постоянна. Более общо, если $f$ целая и
$|f(z)|\le C|z|^p$ при $|z|\ge r_0$, то $f$ --- многочлен степени $\le p$.
\end{thmbox}
\textit{Доказательство последнего.} По неравенствам Коши с $a=0$: $|c_n|\le Cr^{p}/r^{n}\to0$ при $r\to\infty$, если
$n>p$. $\square$

\begin{thmbox}{Факт 0.3 (теорема Вейерштрасса, Теорема 7.5 лекций)}
Если $f_n\in H(D)$ и $f_n\to f$ локально равномерно в области $D$, то $f\in H(D)$ и $f_n^{(m)}\to f^{(m)}$ локально
равномерно в $D$ при каждом $m\in\NN$.
\end{thmbox}
\textit{Доказательство.} $f$ непрерывна; для треугольника $\Delta\subset D$:
$\int_{\partial\Delta}f=\lim\int_{\partial\Delta}f_n=0$, по теореме Мореры $f\in H(D)$. Если $\clos B(a,r+d)\subset D$
и $z_0\in\clos B(a,r)$, то по формуле Коши
$|f_n'(z_0)-f'(z_0)|\le\frac1{2\pi}\cdot\frac{2\pi(r+d)}{d^2}\|f_n-f\|_{C(a,r+d)}\to0$ равномерно по $z_0$;
любой компакт покрывается конечным числом таких кругов; далее индукция по $m$. $\square$

\begin{thmbox}{Факт 0.4 (нули и теорема единственности)}
Пусть $f\in H(D)$, $D$ --- область, $f\not\equiv0$. Тогда нули $f$ изолированы: если $f(a)=0$, то
$f(z)=(z-a)^n g(z)$ в $U(a)$, где $n\in\NN$ (порядок нуля), $g\in H(U(a))$, $g(a)\ne0$; множество нулей не имеет
предельных точек в $D$. \textbf{Теорема единственности:} если $f,g\in H(D)$ совпадают на множестве, имеющем
предельную точку в $D$, то $f\equiv g$ в $D$.
\end{thmbox}
\textit{Идея.} Множество точек, где все производные $f$ обращаются в нуль, открыто (ряд Тейлора) и замкнуто, значит,
пусто или совпадает с $D$; если же в $a$ не все производные нулевые, то первый ненулевой коэффициент Тейлора
даёт порядок нуля. $\square$

\begin{thmbox}{Факт 0.5 (теорема о среднем и принцип максимума модуля)}
(а) Если $f\in H(\clos B(a,r))$, то $f(a)=\frac1{2\pi}\int_0^{2\pi}f(a+re^{i\vartheta})\,d\vartheta$; то же верно для
$\Re f$ и $\Im f$, а также для любой вещественной функции вида $\Re F$, $F\in H(\clos B(a,r))$ (например, для
$\ln|F|$, если $F$ не имеет нулей в $\clos B(a,r)$: $\ln|F|=\Re\ln F$).\\
(б) \textbf{Принцип максимума модуля:} если $f\in H(D)$ и $|f|$ имеет локальный максимум в точке $D$, то
$f\equiv\mathrm{const}$. Следовательно, если $D$ ограничена и $f\in H(D)\cap C(\clos D)$, то
$\max_{\clos D}|f|=\max_{\partial D}|f|$.\\
(в) Если $f\in H(D)\cap C(\clos D)$, $D$ ограничена, то $\max_{\clos D}\Re f=\max_{\partial D}\Re f$, и если
$\Re f$ достигает максимума во внутренней точке, то $f\equiv\mathrm{const}$.
\end{thmbox}
\textit{Доказательство (б).} Пусть $|f(z)|\le|f(a)|$ в $B(a,\delta)\subset D$. По (а) при $0<r<\delta$
\[
|f(a)|\le\frac1{2\pi}\int_0^{2\pi}|f(a+re^{i\vartheta})|\,d\vartheta\le|f(a)|,
\]
поэтому непрерывная функция $|f(a)|-|f(a+re^{i\vartheta})|\ge0$ имеет нулевой интеграл, т.е. $|f|\equiv|f(a)|$ на
$C(a,r)$ при всех $r<\delta$, и $|f|$ постоянен в $B(a,\delta)$. Тогда $f(B(a,\delta))$ лежит на окружности и не
может быть открытым множеством, поэтому по принципу сохранения области (Факт 0.11 ниже, доказательство которого
не использует (б)) $f$ постоянна в $B(a,\delta)$, а по теореме единственности --- в $D$. Второе утверждение:
$|f|$ достигает максимума на компакте $\clos D$; если точка максимума внутренняя, то $f$ --- константа. $\square$

\textit{Доказательство (в).} $|e^{f}|=e^{\Re f}$; применяем (б) к $e^f$; если $e^f\equiv\mathrm{const}$, то
$f'e^f\equiv0$, т.е. $f'\equiv0$. $\square$

\begin{thmbox}{Факт 0.6 (изолированные особые точки)}
Пусть $f\in H(V(a))$, $a\in\Ch$. Точка $a$ --- \emph{устранимая}, если существует конечный $\lim_{z\to a}f(z)$
(эквивалентно: $f$ ограничена в $V(a)$ --- теорема Римана); \emph{полюс}, если $\lim_{z\to a} f(z)=\infty$
(тогда $f=g/(z-a)^p$, $g(a)\ne0$, $p$ --- порядок полюса; при $a=\infty$ --- $f(z)=z^pg(z)$, $g(\infty)\ne0$);
\emph{существенно особая}, если предела нет. \textbf{Теорема Сохоцкого--Казорати--Вейерштрасса:} если $a$ ---
существенно особая точка, то $f(W)$ плотно в $\CC$ для любой проколотой окрестности $W\subset V(a)$.
\end{thmbox}
\textit{Идея (теорема Римана об устранимой особенности, $a\in\CC$).} Если $|f|\le M$ в $V(a)$, то функция
$h(z)=(z-a)^2f(z)$, $h(a)=0$, $\CC$-дифференцируема и в точке $a$ ($h'(a)=\lim(z-a)f(z)=0$), т.е. $h\in H(U(a))$,
$h(a)=h'(a)=0$, и $f=h/(z-a)^2$ продолжается голоморфно в $a$. \textit{Идея (Сохоцкий).} Если
$B(w_0,\eps)\cap f(W)=\varnothing$, то $1/(f-w_0)$ ограничена в $W$, по теореме Римана имеет в $a$ устранимую
особенность, и тогда $a$ --- устранимая точка или полюс для $f$. $\square$

\begin{thmbox}{Факт 0.7 (индекс пути, теорема Жордана)}
(а) Для замкнутого пути $\gamma$ и $a\notin\gamma^*$ определён индекс $\Ind(\gamma,a)\in\ZZ$; для спрямляемого
$\gamma$: $\Ind(\gamma,a)=\frac1{2\pi i}\int_\gamma\frac{dz}{z-a}$, для произвольного --- через вписанную ломаную
(Леммы 7.1--7.3 лекций); $\Ind(\gamma,a)=\frac1{2\pi}\Delta_\gamma\Arg(z-a)$. Индекс постоянен на компонентах
связности $\CC\setminus\gamma^*$ и равен нулю на неограниченной компоненте.\\
(б) \textbf{Теорема Жордана.} Если $\Gamma$ --- жорданова кривая в $\CC$, то $\CC\setminus\Gamma$ состоит ровно из
двух компонент: ограниченной $\operatorname{Int}\Gamma$ и неограниченной $\operatorname{Ext}\Gamma$, причём
$\Gamma=\partial\operatorname{Int}\Gamma=\partial\operatorname{Ext}\Gamma$. \textbf{Дополнение (об индексе):}
существует $\eps\in\{1,-1\}$ такое, что $\Ind(\Gamma,b)=\eps$ при всех $b\in\operatorname{Int}\Gamma$ (и
$\Ind(\Gamma,b)=0$ при $b\in\operatorname{Ext}\Gamma$); при $\eps=1$ ориентация $\Gamma$ называется положительной.
\end{thmbox}

\begin{thmbox}{Факт 0.8 (односвязность)}
Для области $D\subset\CC$ равносильны: (a) $\Ch\setminus D$ связно; (b) $\Ind(\gamma,a)=0$ для всех замкнутых путей
$\gamma$ в $D$ и всех $a\notin D$; (c) всякий замкнутый путь в $D$ гомотопен в $D$ постоянному. Такие области
называются \emph{односвязными}; для областей в $\Ch$ односвязность определяется условием (c) (или (a)).
Жорданова область односвязна ($\Ch\setminus\operatorname{Int}\Gamma=\Gamma\cup\operatorname{Ext}\Gamma\cup\{\infty\}$
связно). Свойство (c), очевидно, сохраняется при гомеоморфизмах.
\end{thmbox}
\textit{Комментарий.} (a)$\Rightarrow$(b): если $a\notin D$, то связное множество $\Ch\setminus D$ содержит $a$ и
$\infty$ и не пересекает $\gamma^*$, поэтому $a$ лежит в неограниченной компоненте $\CC\setminus\gamma^*$.
(b)$\Rightarrow$(c) удобно получать из теоремы Римана (билеты 10--11): область со свойством (b) конформно
эквивалентна $\DD$ или $\CC$. (c)$\Rightarrow$(b): индекс инвариантен относительно гомотопий.

\begin{thmbox}{Факт 0.9 (ветви логарифма и корня; Лемма 10.1 и Лемма 12.3 лекций)}
Пусть $D$ --- односвязная область в $\CC$, $f\in H(D)$, $f(z)\ne0$ в $D$. Тогда существует $g\in H(D)$ с
$f=e^g$ (голоморфная ветвь $\ln f$), и для любого $m\in\NN$, $m\ge2$, существует $h\in H(D)$ с $f=h^m$ (например,
$h=e^{g/m}$). Ветвь логарифма определяется однозначно с точностью до слагаемого $2\pi ik$, $k\in\ZZ$, ветвь корня ---
с точностью до множителя $e^{2\pi ik/m}$.
\end{thmbox}
\textit{Доказательство.} $f'/f\in H(D)$; в силу (b) Факта 0.8 и гомологической теоремы Коши интегралы $f'/f$ по
замкнутым путям равны нулю, поэтому $f'/f$ имеет первообразную $g_0$ в $D$. Тогда
$(fe^{-g_0})'=e^{-g_0}(f'-fg_0')=0$, т.е. $f=Ce^{g_0}$, $C\ne0$; положим $g=g_0+c$, $e^c=C$. Если $e^{g_1}=e^{g_2}$,
то $(g_1-g_2)/(2\pi i)$ --- непрерывная целочисленная функция на связном $D$, т.е. константа. $\square$

\begin{thmbox}{Факт 0.10 (принцип аргумента, теорема Руше; Теоремы 10.1, 10.2, Предложение 10.2)}
(а) Пусть $D$ --- Д.О., $f$ мероморфна в $D$ (с конечным числом полюсов), $f\in C(\clos D\setminus P(f,D))$ и
$f\ne0$ на $\partial D$. Тогда
\[
\mathfrak z(f,D)-\mathfrak p(f,D)=\frac1{2\pi}\Delta_{\partial^+D}\Arg f .
\]
Если $D$ --- жорданова область, $\partial^+D$ задаётся путём $\gamma$, то правая часть равна
$\Ind(f\circ\gamma,0)$.\\
(б) \textbf{Теорема Руше.} Пусть $D$ --- Д.О., $f,g\in H(D)\cap C(\clos D)$, $|g|<|f|$ на $\partial D$. Тогда
$\mathfrak z(f+g,D)=\mathfrak z(f,D)$.
\end{thmbox}
\textit{Идея.} Для $f$, голоморфной в окрестности $\clos D$ и Д.О. со спрямляемой границей,
$\frac1{2\pi i}\int_{\partial^+D}\frac{f'}{f}dz$ равен сумме логарифмических вычетов, т.е. $\mathfrak z-\mathfrak p$, и
равен $\frac1{2\pi}\Delta\Arg f$ (Лемма 10.2). Общий случай --- аппроксимацией $D$ изнутри. В (б):
$\Delta\Arg(f+g)=\Delta\Arg f+\Delta\Arg(1+g/f)$, а $1+g/f$ на $\partial D$ лежит в правой полуплоскости, так что
$\Delta\Arg(1+g/f)=0$ (Лемма 10.3). $\square$

\begin{thmbox}{Факт 0.11 (принцип сохранения области, Теорема 10.3)}
Если $D\subset\CC$ --- область, $f\in H(D)$, $f\not\equiv\mathrm{const}$, то $f(D)$ --- область. Следствие: если
$f\in\mathrm{Conf}(D)$ (или просто $f\in H(D)$ инъективна), то $f$ --- гомеоморфизм $D$ на область $f(D)$.
\end{thmbox}
\textit{Доказательство.} Пусть $w_0=f(z_0)$. Нуль $z_0$ функции $f-w_0$ изолирован: найдётся $\delta>0$ с
$\clos B(z_0,\delta)\subset D$, $\eps:=\min_{C(z_0,\delta)}|f-w_0|>0$. Для $|w-w_0|<\eps$ по теореме Руше
($f-w=(f-w_0)+(w_0-w)$) $\mathfrak z(f-w,B(z_0,\delta))=\mathfrak z(f-w_0,B(z_0,\delta))\ge1$, т.е.
$B(w_0,\eps)\subset f(D)$. Связность $f(D)$ --- из непрерывности. Для следствия: открытость $f$ означает
непрерывность $f^{-1}$. $\square$

\begin{thmbox}{Факт 0.12 (теорема об обратной функции; Предложение 2.6 и его следствия)}
(а) Пусть $f:U(a)\to U(b)$ --- гомеоморфизм, $b=f(a)$, $g=f^{-1}$. Если $f$ $\CC$-дифференцируема в $a$ и
$f'(a)\ne0$, то $g$ $\CC$-дифференцируема в $b$ и $g'(b)=1/f'(a)$.\\
(б) Если $f\in H(a)$ и $f'(a)\ne0$, то $f$ \emph{локально обратима}: существуют окрестности $U\ni a$, $\tilde U\ni f(a)$
и $g\in H(\tilde U)$ с $g(f(z))=z$, $z\in U$.\\
(в) Если $f\in H(D)$ инъективна, то $f'\ne0$ в $D$, $G=f(D)$ --- область и $f^{-1}\in H(G)$, т.е.
$f\in\mathrm{Conf}(D,G)$.
\end{thmbox}
\textit{Доказательство.} (а) Если $b+\Delta w=f(a+\Delta z)$, то $\Delta z\to0$ при $\Delta w\to0$ (непрерывность $g$),
и $\frac{g(b+\Delta w)-g(b)}{\Delta w}=\frac{\Delta z}{f(a+\Delta z)-f(a)}\to\frac1{f'(a)}$ (здесь $\Delta z\ne0$
при $\Delta w\ne0$ в силу инъективности).

(б) Выберем $\rho>0$ так, что $\clos B(a,\rho)$ лежит в области голоморфности $f$ и
$|f'(z)-f'(a)|<\frac12|f'(a)|$ в $B(a,\rho)$. Для $z_1\ne z_2$ из $B(a,\rho)$, интегрируя по отрезку $I=[z_1,z_2]$,
\[
|f(z_2)-f(z_1)-f'(a)(z_2-z_1)|=\Bigl|\int_I(f'(w)-f'(a))\,dw\Bigr|<\tfrac12|f'(a)|\,|z_2-z_1|,
\]
откуда $f(z_1)\ne f(z_2)$. Итак, $f$ инъективна в $U=B(a,\rho)$; по Факту 0.11 $\tilde U=f(U)$ --- область и
$f:U\to\tilde U$ --- гомеоморфизм, а по (а) обратная функция $\CC$-дифференцируема в каждой точке $\tilde U$
(во всех точках $U$ имеем $f'\ne0$ по выбору $\rho$).

(в) Пусть $f'(z_0)=0$. Так как $f$ инъективна, $f\not\equiv\mathrm{const}$, и
$f(z)-f(z_0)=(z-z_0)^ng(z)$, $n\ge2$, $g(z_0)\ne0$. Нули $f'$ и $f-f(z_0)$ изолированы, поэтому найдётся
$\delta>0$ такое, что в $\clos B(z_0,\delta)\setminus\{z_0\}$ нет ни нулей $f'$, ни нулей $f-f(z_0)$. Пусть
$\eps=\min_{C(z_0,\delta)}|f-f(z_0)|>0$ и $0<|w-f(z_0)|<\eps$. По теореме Руше $f-w$ имеет в $B(z_0,\delta)$
ровно $n\ge2$ нулей с учётом кратности; $z_0$ среди них нет, а в остальных точках $f'\ne0$, т.е. все нули
простые и их не меньше двух --- противоречие с инъективностью. Значит, $f'\ne0$ в $D$, $G=f(D)$ --- область
(Факт 0.11), $f^{-1}$ непрерывна и по (а) голоморфна в $G$. $\square$

\begin{thmbox}{Факт 0.13 (дробно-линейные отображения)}
(а) Всякое $W\in\Lambda$ --- биекция $\Ch\to\Ch$, конформная во всех точках $\Ch$ (с соглашением
$W(-d/c)=\infty$, $W(\infty)=a/c$); $\Lambda$ --- группа относительно композиции; $W$ определяется матрицей
$\begin{pmatrix}a&b\\c&d\end{pmatrix}$ с точностью до ненулевого множителя.\\
(б) Для любых трёх различных точек $z_1,z_2,z_3\in\Ch$ и трёх различных $w_1,w_2,w_3\in\Ch$ существует ровно одно
$W\in\Lambda$ с $W(z_k)=w_k$.\\
(в) \textbf{Круговое свойство:} $W$ переводит обобщённые окружности (окружности и прямые, дополненные $\infty$) в
обобщённые окружности; для любых двух обобщённых окружностей существует $W\in\Lambda$, переводящее одну в другую.\\
(г) \textbf{Сохранение симметрии (Лемма 3.3):} точки, симметричные относительно обобщённой окружности $\Gamma$,
переходят в точки, симметричные относительно $W(\Gamma)$. Симметрия относительно $\Ch\supset\RR\cup\{\infty\}$ ---
это $z\mapsto\bar z$; относительно $C(z_0,R)$ --- $z\mapsto z_0+R^2/(\bar z-\bar z_0)$.\\
(д) Для $a\in\DD$: $\varphi_a\in\Lambda$, $\varphi_a(a)=0$, $\varphi_a(0)=-a$, $\varphi_a^{-1}=\varphi_{-a}$,
$\varphi_a'(z)=\frac{1-|a|^2}{(1-\bar az)^2}$ и
\begin{equation}\label{t0:phi}
1-|\varphi_a(z)|^2=\frac{(1-|a|^2)(1-|z|^2)}{|1-\bar az|^2}\qquad(z\ne1/\bar a),
\end{equation}
так что $\varphi_a(\DD)=\DD$, $\varphi_a(\TT)=\TT$, $\varphi_a\in\CAut(\DD)$.
\end{thmbox}
\textit{Проверка (д).} $|1-\bar az|^2-|z-a|^2=1-2\Re(\bar az)+|a|^2|z|^2-|z|^2+2\Re(\bar a z)-|a|^2
=(1-|a|^2)(1-|z|^2)$, откуда \eqref{t0:phi}. Далее
$\varphi_{-a}(\varphi_a(z))=\frac{(z-a)+a(1-\bar az)}{(1-\bar az)+\bar a(z-a)}=\frac{(1-|a|^2)z}{1-|a|^2}=z$.
Из \eqref{t0:phi}: $|\varphi_a(z)|<1\Leftrightarrow|z|<1$ и $|\varphi_a(z)|=1\Leftrightarrow|z|=1$; вместе с
биективностью $\varphi_a$ на $\Ch$ это даёт $\varphi_a(\DD)=\DD$. $\square$

\begin{thmbox}{Факт 0.14 (функция Жуковского)}
$\Zh(z)=\frac12\bigl(z+\frac1z\bigr)$; $\Zh(z_1)=\Zh(z_2)$, $z_1\ne z_2$ $\Leftrightarrow$ $z_1z_2=1$;
$\Zh'(z)=\frac12(1-z^{-2})$ обращается в нуль лишь в $\pm1$. Поэтому $\Zh$ конформно отображает
$\{|z|>1\}\cup\{\infty\}$ на $\Ch\setminus[-1,1]$ (и $\DD$ --- на ту же область), а $\CC_+$ --- на
$\CC\setminus\bigl((-\infty,-1]\cup[1,+\infty)\bigr)$; окружность $\TT$ переходит в отрезок $[-1,1]$, проходимый
дважды: $\Zh(e^{it})=\cos t$.
\end{thmbox}
\textit{Проверка первого утверждения.} $\Zh(z_1)-\Zh(z_2)=\frac12(z_1-z_2)\bigl(1-\frac1{z_1z_2}\bigr)$. Значение
$w=\Zh(z)$ принимается в точках $z$ и $1/z$ --- корнях уравнения $z^2-2wz+1=0$; при $w\notin[-1,1]$ ровно один
из корней лежит в $\{|z|>1\}$ (их произведение равно $1$, а $|z|=1$ влечёт $w=\Re z\in[-1,1]$). $\square$

\begin{thmbox}{Факт 0.15 (прямой принцип соответствия границ, простейшая форма)}
Пусть $D$, $G\subset\CC$ --- ограниченные области, $f\in\mathrm{Conf}(D,G)$ и $f$ продолжается до непрерывного
отображения $\clos D\to\CC$. Тогда $f(\clos D)=\clos G$ и $f(\partial D)=\partial G$. Если, кроме того, $D$ и $G$
--- жордановы области, а продолжение инъективно на $\partial D$, то $f$ --- гомеоморфизм $\clos D$ на $\clos G$,
и при положительном обходе $\partial D$ кривая $f(\partial D)$ обходится положительно.
\end{thmbox}
\textit{Доказательство.} $f(\clos D)$ компактно и содержит $G$, значит, содержит $\clos G$; с другой стороны,
$f(\clos D)\subset\clos{f(D)}=\clos G$ по непрерывности. Пусть $\zeta\in\partial D$ и $f(\zeta)\in G$; возьмём
$z_n\in D$, $z_n\to\zeta$. Тогда $f(z_n)\to f(\zeta)$ и по непрерывности $f^{-1}$ на $G$ (Факт 0.11)
$z_n=f^{-1}(f(z_n))\to f^{-1}(f(\zeta))\in D$, т.е. $\zeta\in D$ --- противоречие. Итак, $f(\partial D)\subset\partial G$,
а так как $f(D)=G$ и $f(\clos D)=\clos G$, то $f(\partial D)\supset\clos G\setminus G=\partial G$. При инъективности на
$\partial D$ продолжение --- непрерывная биекция компакта $\clos D$ на $\clos G$, т.е. гомеоморфизм. Ориентация:
для $w\in G$ по принципу аргумента $\Ind(f\circ\gamma,w)=\mathfrak z(f-w,D)=1>0$. $\square$
\begin{remarkbox}{Замечание}
Тот факт, что \emph{всякое} конформное отображение жордановых областей продолжается до гомеоморфизма замыканий,
--- это теорема Каратеодори (билеты 18--20). Обратное утверждение (из взаимной однозначности на границе следует
однолистность внутри) --- \emph{обратный} принцип соответствия границ, он доказан в билете 2.
\end{remarkbox}

% ======================================================================
% Глава 12
% ======================================================================
\subsection{Глава 12: лемма Шварца и смежные вопросы (слайды 258--272)}

\begin{defs}[Определения главы 12]
\begin{itemize}
\item $\varphi_a(z)=\dfrac{z-a}{1-\bar az}$, $a\in\DD$ --- \emph{элементарный множитель Бляшке} (Факт 0.13(д)):
$\varphi_a\in\CAut(\DD)$, $\varphi_a(a)=0$, $\varphi_a^{-1}=\varphi_{-a}$.
\item \emph{Конечное произведение Бляшке} степени $n\in\ZZ_+$ --- функция
$B(z)=e^{i\alpha}\prod_{k=1}^n\varphi_{a_k}(z)$, $\alpha\in\RR$, $a_k\in\DD$ (точки $a_k$ могут повторяться;
при $n=0$ это унимодулярная константа $e^{i\alpha}$). $B$ голоморфна в окрестности $\clos\DD$ (полюсы
$1/\bar a_k$ лежат вне $\clos\DD$), $|B|=1$ на $\TT$, $|B|<1$ в $\DD$ при $n\ge1$, нули $B$ в $\DD$ --- ровно
$a_1,\dots,a_n$ (с кратностями). \textit{Пример:} $z^n$, $\varphi_a$, $z\varphi_a(z)$.
\item \emph{Псевдогиперболическое расстояние} в $\DD$: $\rho(z,w)=\Bigl|\dfrac{z-w}{1-\bar wz}\Bigr|=|\varphi_w(z)|$,
$z,w\in\DD$; $0\le\rho<1$, $\rho(z,w)=\rho(w,z)$, $\rho(z,w)=0\Leftrightarrow z=w$, $\rho(z,0)=|z|$.
\emph{Псевдогиперболический круг} $K(a,r)=\{z\in\DD:\rho(z,a)<r\}$, $0<r<1$.
\item Для $f\in H(\clos B(0,R))$ и $0\le r\le R$: $M(r)=M_f(r)=\max_{|z|=r}|f(z)|$,
$A(r)=A_f(r)=\max_{|z|=r}\Re f(z)$ (при $r=0$: $M(0)=|f(0)|$, $A(0)=\Re f(0)$). По принципу максимума
(Факт 0.5) $M$ и $A$ не убывают и $M(r)=\max_{|z|\le r}|f|$, $A(r)=\max_{|z|\le r}\Re f$.
\item \emph{Итерированный логарифм}: $\log_{(1)}t=\log t$, $\log_{(n+1)}t=\log\log_{(n)}t$ (определён при
достаточно больших $t$; возрастает и $\to+\infty$ при $t\to+\infty$).
\item $\mathcal E_{(n)}$ --- класс целых функций $f$ с
$\displaystyle\limsup_{r\to\infty}\frac{\log_{(n)}M_f(r)}{\log r}<\infty$;
$\mathcal E=\bigcup_{n\ge1}\mathcal E_{(n)}$. Для непостоянной целой $f$ имеем $M_f(r)\uparrow+\infty$ (теорема
Лиувилля), поэтому $\log_{(n)}M_f(r)$ определён при больших $r$; постоянные функции считаем входящими во все
$\mathcal E_{(n)}$. \textit{Примеры:} многочлены $\in\mathcal E_{(1)}$; $e^z\in\mathcal E_{(2)}$
($\log\log M=\log r$); $\exp(e^z)\in\mathcal E_{(3)}$.
\end{itemize}
\end{defs}

\subsubsection*{Лемма Шварца}

\begin{thmbox}{Предложение 12.1 (лемма Шварца)}
Пусть $f\in H(\DD)$, $|f(z)|\le1$ при $|z|<1$ и $f(0)=0$. Тогда $|f(z)|\le|z|$ при всех $z\in\DD$ и
$|f'(0)|\le1$. Если $|f(z_0)|=|z_0|$ хотя бы в одной точке $z_0\in\DD\setminus\{0\}$ или $|f'(0)|=1$, то
$f(z)=e^{i\alpha}z$, $\alpha\in\RR$.
\end{thmbox}
\begin{proof}
Положим $g(z)=f(z)/z$ при $z\in\DD\setminus\{0\}$, $g(0)=f'(0)$. Так как $f(0)=0$, то $f(z)=\sum_{n\ge1}c_nz^n$
в $\DD$ и $g(z)=\sum_{n\ge1}c_nz^{n-1}$, т.е. $g\in H(\DD)$ (особенность в нуле устранима).

Фиксируем $z\in\DD$ и $r\in(|z|,1)$. По принципу максимума модуля в круге $\clos B(0,r)$ (Факт 0.5(б))
\[
|g(z)|\le\max_{|\zeta|=r}|g(\zeta)|=\max_{|\zeta|=r}\frac{|f(\zeta)|}{r}\le\frac1r .
\]
Устремляя $r\to1-$, получаем $|g(z)|\le1$ при всех $z\in\DD$, т.е. $|f(z)|\le|z|$ и $|f'(0)|=|g(0)|\le1$.

Если $|f(z_0)|=|z_0|$, $z_0\ne0$, или $|f'(0)|=1$, то $|g|$ достигает в точке $z_0$ (соответственно $0$)
своего наибольшего значения $1$ на $\DD$; по принципу максимума $g\equiv c$, $|c|=1$, т.е. $c=e^{i\alpha}$ и
$f(z)=e^{i\alpha}z$.
\end{proof}

\begin{remarkbox}{Замечание (уточнение формулировки)}
На слайде сказано <<если равенство выполнено хотя бы в одной точке $z_0\in\DD$>>. Для неравенства
$|f(z)|\le|z|$ точку $z_0=0$ надо исключить: там равенство $0=0$ выполнено всегда.
\end{remarkbox}

\begin{detail}{Упражнение: вывести лемму Шварца из неравенств Коши}
Пусть $f=\sum_{n\ge1}c_nz^n$. \emph{Оценка $|f'(0)|\le1$:} по неравенству Коши
$|c_1|\le M(r)/r\le1/r$ при всех $r<1$, откуда $|f'(0)|=|c_1|\le1$.

\emph{Оценка $|f(z)|\le|z|$.} Функция $g=f/z$ (см. выше) голоморфна в $\DD$ и $|g|\le1/r$ на $C(0,r)$.
Фиксируем $z_0$, $|z_0|<r<1$, и $k\in\NN$. Применим к $g^k$ интегральную формулу Коши (оценку того же типа, что
и неравенства Коши):
\[
|g(z_0)|^k=\Bigl|\frac1{2\pi i}\int_{C(0,r)}\frac{g(\zeta)^k\,d\zeta}{\zeta-z_0}\Bigr|
\le\frac1{2\pi}\cdot\frac{2\pi r}{r-|z_0|}\cdot\frac1{r^k}.
\]
Извлекая корень степени $k$ и устремляя $k\to\infty$, получаем $|g(z_0)|\le1/r$, затем $r\to1-$:
$|g(z_0)|\le1$.

\emph{Случай равенства <<на уровне коэффициентов>>.} Для $h=\sum d_n(z-z_0)^n\in H(\clos B(z_0,s))$ имеет место
равенство Парсеваля $\frac1{2\pi}\int_0^{2\pi}|h(z_0+se^{i\vartheta})|^2d\vartheta=\sum_n|d_n|^2s^{2n}$
(перемножаем равномерно сходящиеся ряды для $h$ и $\bar h$ и пользуемся ортогональностью
$\int_0^{2\pi}e^{i(n-m)\vartheta}d\vartheta=2\pi\delta_{nm}$). Если $|g(z_0)|=1$ ($z_0\in\DD$), то для
$g(z)=\sum d_n(z-z_0)^n$ в $B(z_0,1-|z_0|)$ получаем $1+\sum_{n\ge1}|d_n|^2s^{2n}\le1$, т.е. $d_n=0$ при $n\ge1$,
$g$ постоянна вблизи $z_0$, а по теореме единственности --- в $\DD$. Это даёт и случай $|f'(0)|=|g(0)|=1$.
\end{detail}

\begin{detail}{Задача: вывести теорему Лиувилля из леммы Шварца}
Пусть $f$ целая и $|f|\le M$. Для $R>0$ функция $F(\zeta)=\dfrac{f(R\zeta)-f(0)}{2M}$ голоморфна в $\DD$,
$|F|\le1$, $F(0)=0$ (если $M=0$, доказывать нечего). По лемме Шварца $|F(\zeta)|\le|\zeta|$, т.е.
$|f(z)-f(0)|\le 2M|z|/R$ при $|z|<R$. Фиксируя $z$ и устремляя $R\to\infty$, получаем $f(z)=f(0)$.
\end{detail}

\begin{thmbox}{Замечание 12.1 (модификации леммы Шварца)}
1) Пусть $f\in H(B(a,R))$, $|f(z)|\le M$ при $z\in B(a,R)$ и $f(a)=0$. Тогда $|f(z)|\le M|z-a|/R$ при всех
$z\in B(a,R)$.\\
2) Пусть в условиях леммы Шварца $f$ имеет в $0$ нуль порядка $m\ge1$. Тогда $|f(z)|\le|z|^m$ при $z\in\DD$, а
если равенство выполнено хотя бы в одной точке $z_0\in\DD\setminus\{0\}$, то $f=e^{i\alpha}z^m$, $\alpha\in\RR$.
Кроме того, $|f^{(m)}(0)|\le m!$, с равенством только для $f=e^{i\alpha}z^m$.
\end{thmbox}
\begin{proof}
1) При $M=0$ утверждение очевидно. Иначе $F(\zeta)=f(a+R\zeta)/M$ удовлетворяет условиям леммы Шварца, и
$|f(z)|=M|F((z-a)/R)|\le M|z-a|/R$.

2) $f(z)=z^mg(z)$, $g\in H(\DD)$ (если $f\equiv0$, всё очевидно; <<нуль порядка $\ge m$>> тоже подходит).
Как в доказательстве Предложения 12.1: при $|z|<r<1$ имеем $|g(z)|\le\max_{|\zeta|=r}|f(\zeta)|/r^m\le r^{-m}$,
и при $r\to1-$ получаем $|g|\le1$ в $\DD$, т.е. $|f(z)|\le|z|^m$. Равенство в точке $z_0\ne0$ (или равенство
$|g(0)|=|f^{(m)}(0)|/m!=1$) означает, что $|g|$ достигает максимума $1$ внутри $\DD$, откуда $g\equiv e^{i\alpha}$.
\end{proof}

\subsubsection*{Конформные автоморфизмы круга}

\begin{thmbox}{Предложение 12.2 (общий вид конформных автоморфизмов $\DD$)}
Пусть $f$ --- конформное отображение круга $\DD$ на себя. Тогда
\[
f(z)=e^{i\alpha}\,\frac{z-a}{1-\bar az}=e^{i\alpha}\varphi_a(z),\qquad a\in\DD,\ \alpha\in\RR .
\]
Обратно, все такие функции принадлежат $\CAut(\DD)$. Представление единственно: $a=f^{-1}(0)$, $e^{i\alpha}$
определяется по $f$ однозначно.
\end{thmbox}
\begin{proof}
По условию $f\in H(\DD)$, $f$ --- биекция $\DD$ на $\DD$, и $f^{-1}\in H(\DD)$ (Факт 0.12(в)). Пусть
$a=f^{-1}(0)\in\DD$. Положим $h=f\circ\varphi_{-a}$; так как $\varphi_{-a}\in\CAut(\DD)$ и $\varphi_{-a}(0)=a$, то $h$
--- голоморфная биекция $\DD$ на $\DD$, $h(0)=0$, $|h|<1$. По лемме Шварца $|h(z)|\le|z|$. Функция $h^{-1}$
также голоморфна в $\DD$, $|h^{-1}|<1$, $h^{-1}(0)=0$, и по лемме Шварца $|h^{-1}(w)|\le|w|$; подставляя
$w=h(z)$, получаем $|z|\le|h(z)|$. Итак, $|h(z)|=|z|$ для всех $z$, и по случаю равенства в лемме Шварца
$h(z)=e^{i\alpha}z$. Тогда $f=h\circ\varphi_{-a}^{-1}=h\circ\varphi_a=e^{i\alpha}\varphi_a$.

Обратно, $\varphi_a\in\CAut(\DD)$ (Факт 0.13(д)), поворот $z\mapsto e^{i\alpha}z$ --- тоже, а $\CAut(\DD)$ --- группа.
Единственность: из $e^{i\alpha}\varphi_a=e^{i\beta}\varphi_b$ подстановкой $z=a$ получаем $\varphi_b(a)=0$, т.е.
$b=a$, а затем $e^{i\alpha}=e^{i\beta}$.
\end{proof}
\begin{remarkbox}{Следствия}
(1) $\CAut(\DD)$ --- трёхпараметрическое семейство (параметры: $a\in\DD$ --- 2 вещественных, $\alpha$ --- 1).
(2) Автоморфизмы с $f(0)=0$ --- в точности повороты $e^{i\alpha}z$. (3) Автоморфизм, сохраняющий точку $a$, имеет
вид $\varphi_a^{-1}(e^{i\alpha}\varphi_a(z))$. (4) Все конформные автоморфизмы круга дробно-линейны.
\end{remarkbox}

\subsubsection*{Лемма Шварца--Пика}

\begin{thmbox}{Предложение 12.3 (лемма Шварца--Пика)}
Пусть $f\in H(\DD)$ и $|f(z)|<1$ при $|z|<1$. Тогда
\begin{equation}\label{t0:pick}
\Bigl|\frac{f(z)-f(w)}{1-\overline{f(w)}f(z)}\Bigr|\le\Bigl|\frac{z-w}{1-\bar wz}\Bigr|,\qquad z\ne w,\ z,w\in\DD,
\end{equation}
т.е. $\rho(f(z),f(w))\le\rho(z,w)$, и
\begin{equation}\label{t0:pickd}
|f'(z)|\le\frac{1-|f(z)|^2}{1-|z|^2},\qquad z\in\DD .
\end{equation}
Равенство в \eqref{t0:pick} хотя бы для одной пары $z\ne w$ или в \eqref{t0:pickd} хотя бы в одной точке
имеет место тогда и только тогда, когда $f\in\CAut(\DD)$ (т.е. $f$ --- ДЛО круга $\DD$ на себя); в этом случае
равенства выполняются тождественно.
\end{thmbox}
\begin{proof}
Фиксируем $w\in\DD$, положим $b=f(w)\in\DD$ и рассмотрим
$\varphi(z)=\varphi_{-w}(z)=\dfrac{z+w}{1+\bar wz}$, $\psi(u)=\varphi_b(u)=\dfrac{u-b}{1-\bar bu}$,
$g=\psi\circ f\circ\varphi$. Тогда $g\in H(\DD)$, $|g|<1$ (так как $|f|<1$ и $\psi(\DD)=\DD$), и
$g(0)=\psi(f(w))=0$. По лемме Шварца
\begin{equation}\label{t0:pickg}
|g(\zeta)|\le|\zeta|\quad(\zeta\in\DD),\qquad|g'(0)|\le1 .
\end{equation}
\emph{Первое неравенство.} Подставим $\zeta=\varphi^{-1}(z)=\varphi_w(z)$: $g(\varphi_w(z))=\psi(f(z))=\varphi_b(f(z))$,
и \eqref{t0:pickg} превращается в $|\varphi_{f(w)}(f(z))|\le|\varphi_w(z)|$, а это и есть \eqref{t0:pick}.

\emph{Второе неравенство.} По правилу дифференцирования композиции
$g'(0)=\psi'(f(w))\,f'(w)\,\varphi'(0)$. Из Факта 0.13(д):
$\varphi_c'(z)=\dfrac{1-|c|^2}{(1-\bar cz)^2}$, поэтому
\[
\varphi'(0)=\varphi_{-w}'(0)=1-|w|^2,\qquad
\psi'(b)=\frac{1-|b|^2}{(1-|b|^2)^2}=\frac1{1-|f(w)|^2}.
\]
Значит, $|g'(0)|=\dfrac{(1-|w|^2)|f'(w)|}{1-|f(w)|^2}\le1$, что есть \eqref{t0:pickd} в точке $w$.

\emph{Случаи равенства.} Если в \eqref{t0:pick} равенство при некоторых $z\ne w$, то $|g(\zeta_0)|=|\zeta_0|$ при
$\zeta_0=\varphi_w(z)\ne0$; если в \eqref{t0:pickd} равенство в точке $w$, то $|g'(0)|=1$. В обоих случаях по
лемме Шварца $g(\zeta)=e^{i\alpha}\zeta$, и $f=\psi^{-1}\circ g\circ\varphi^{-1}=\varphi_{-b}\circ(e^{i\alpha}\,\cdot)\circ\varphi_w$
--- композиция автоморфизмов круга, т.е. $f\in\CAut(\DD)$.

Обратно, пусть $f=e^{i\alpha}\varphi_a$. Равенство в \eqref{t0:pick} при всех $z,w$ --- это инвариантность $\rho$
(см. деталь ниже). Для \eqref{t0:pickd}: $|f'(z)|=\dfrac{1-|a|^2}{|1-\bar az|^2}$, а по \eqref{t0:phi}
$1-|f(z)|^2=\dfrac{(1-|a|^2)(1-|z|^2)}{|1-\bar az|^2}$, так что $\dfrac{1-|f(z)|^2}{1-|z|^2}=|f'(z)|$.
\end{proof}

\begin{detail}{Два тождества для $\varphi_a$}
Для $a\in\DD$ и $z,w\in\DD$:
\begin{gather}
\varphi_a(z)-\varphi_a(w)=\frac{(1-|a|^2)(z-w)}{(1-\bar az)(1-\bar aw)},\label{t0:id1}\\
1-\overline{\varphi_a(w)}\,\varphi_a(z)=\frac{(1-|a|^2)(1-\bar wz)}{(1-a\bar w)(1-\bar az)}.\label{t0:id2}
\end{gather}
\emph{Проверка \eqref{t0:id1}:} числитель разности
$(z-a)(1-\bar aw)-(w-a)(1-\bar az)=z-\bar azw-a+|a|^2w-w+\bar azw+a-|a|^2z=(1-|a|^2)(z-w)$.\\
\emph{Проверка \eqref{t0:id2}:} приводя к общему знаменателю $(1-a\bar w)(1-\bar az)$, получаем числитель
\[
(1-a\bar w)(1-\bar az)-(\bar w-\bar a)(z-a)=1-\bar az-a\bar w+|a|^2\bar wz-\bar wz+a\bar w+\bar az-|a|^2
=(1-|a|^2)(1-\bar wz).
\]
\end{detail}

\begin{detail}{Упражнение: инвариантность $\rho$ относительно $\CAut(\DD)$}
Пусть $\varphi\in\CAut(\DD)$; по Предложению 12.2 $\varphi=e^{i\alpha}\varphi_a$. Поворот не меняет $\rho$:
$\rho(e^{i\alpha}z,e^{i\alpha}w)=|e^{i\alpha}(z-w)|/|1-e^{-i\alpha}\bar we^{i\alpha}z|=\rho(z,w)$. Для $\varphi_a$
делим \eqref{t0:id1} на \eqref{t0:id2}:
\[
\frac{\varphi_a(z)-\varphi_a(w)}{1-\overline{\varphi_a(w)}\varphi_a(z)}=\frac{z-w}{1-\bar wz}\cdot
\frac{1-a\bar w}{1-\bar aw},
\]
а последняя дробь имеет модуль $1$ (числитель и знаменатель комплексно сопряжены). Значит,
$\rho(\varphi(z),\varphi(w))=\rho(z,w)$.\\
\emph{Другой способ:} по Шварцу--Пику $\rho(\varphi(z),\varphi(w))\le\rho(z,w)$, а применяя то же к
$\varphi^{-1}\in\CAut(\DD)$, получаем обратное неравенство.
\end{detail}

\begin{detail}{Задача: неравенство треугольника для $\rho$}
Докажем более сильное неравенство: для $z,w,u\in\DD$
\begin{equation}\label{t0:tri}
\rho(z,w)\le\frac{\rho(z,u)+\rho(u,w)}{1+\rho(z,u)\rho(u,w)}\le\rho(z,u)+\rho(u,w).
\end{equation}
По инвариантности (применяем $\varphi_u$) можно считать $u=0$; обозначим $s=|z|=\rho(z,0)$, $t=|w|=\rho(0,w)$.
Из \eqref{t0:phi} (с $a=w$) и неравенства $|1-\bar wz|\le1+st$:
\[
1-\rho(z,w)^2=\frac{(1-s^2)(1-t^2)}{|1-\bar wz|^2}\ge\frac{(1-s^2)(1-t^2)}{(1+st)^2}
=1-\Bigl(\frac{s+t}{1+st}\Bigr)^2,
\]
поскольку $(1+st)^2-(1-s^2)(1-t^2)=2st+s^2+t^2=(s+t)^2$. Это первое неравенство \eqref{t0:tri}; второе
очевидно, так как $1+st\ge1$. Итак, $\rho$ --- метрика на $\DD$.
\end{detail}

\begin{detail}{Упражнение: $K(a,r)$ --- евклидов круг}
Пусть $a\in\DD$, $0<r<1$. Тогда $K(a,r)=B(b,R)$, где
\[
b=\frac{(1-r^2)\,a}{1-r^2|a|^2},\qquad R=\frac{r(1-|a|^2)}{1-r^2|a|^2}.
\]
\emph{Вывод.} $z\in K(a,r)\Leftrightarrow|z-a|^2<r^2|1-\bar az|^2$ (для $z\in\DD$; заметим, что
$K(a,r)=\varphi_{-a}(B(0,r))$, и найденный ниже круг лежит в $\DD$, так как $|b|+R<1$, см. ниже). Раскрываем:
\[
|z|^2-2\Re(\bar az)+|a|^2<r^2\bigl(1-2\Re(\bar az)+|a|^2|z|^2\bigr)
\iff(1-r^2|a|^2)|z|^2-2(1-r^2)\Re(\bar az)+|a|^2-r^2<0 .
\]
Деля на $1-r^2|a|^2>0$: $|z|^2-2\Re(\bar bz)+\dfrac{|a|^2-r^2}{1-r^2|a|^2}<0$, т.е.
$|z-b|^2<|b|^2-\dfrac{|a|^2-r^2}{1-r^2|a|^2}$. Числитель правой части после приведения к знаменателю
$(1-r^2|a|^2)^2$:
\[
|a|^2(1-r^2)^2-(|a|^2-r^2)(1-r^2|a|^2)=r^2-2r^2|a|^2+r^2|a|^4=r^2(1-|a|^2)^2,
\]
что и даёт $R$. Далее, $|b|+R=\dfrac{|a|(1-r^2)+r(1-|a|^2)}{1-r^2|a|^2}=\dfrac{(|a|+r)(1-r|a|)}{(1-r|a|)(1+r|a|)}
=\dfrac{|a|+r}{1+r|a|}<1$. Так как неравенство $|z-a|<r|1-\bar az|$ при $|z|\ge1$ невозможно (по \eqref{t0:phi}
при $|z|\ge1$ имеем $|\varphi_a(z)|\ge1$), никакой потери точек нет. Аналогично $\{\rho(z,a)\le r\}=\clos B(b,R)$.
\end{detail}

\begin{detail}{Задача: $|f(z)|\le\dfrac{|f(0)|+|z|}{1+|f(0)||z|}$ для $f\in H(\DD)$, $|f|<1$}
\emph{Способ 1 (указание слайда).} По Шварцу--Пику $\rho(f(z),f(0))\le\rho(z,0)=|z|$ (при $f(z)=f(0)$ это
тривиально), т.е. $f(z)\in\{\rho(\cdot,f(0))\le|z|\}=\clos B(b,R)$ с $a=f(0)$, $r=|z|$ (при $z\ne0$). Поэтому
$|f(z)|\le|b|+R=\dfrac{|f(0)|+|z|}{1+|f(0)||z|}$ по вычислению предыдущего упражнения.\\
\emph{Способ 2.} По \eqref{t0:tri} с $u=f(0)$: $|f(z)|=\rho(f(z),0)\le\dfrac{s+t}{1+st}$, где
$s=\rho(f(z),f(0))\le|z|$, $t=|f(0)|$; функция $s\mapsto\frac{s+t}{1+st}$ возрастает на $[0,1)$
(производная $\frac{1-t^2}{(1+st)^2}>0$), так что можно заменить $s$ на $|z|$.
\end{detail}

\subsubsection*{Конечные произведения Бляшке}

\begin{thmbox}{Лемма 12.1 (характеризация конечных произведений Бляшке)}
Пусть $F\in H(\DD)\cap C(\clos\DD)$, $F$ имеет в $\DD$ конечное число нулей и $|F(z)|=1$ при $z\in\TT$. Тогда
\[
F(z)=e^{i\alpha}\prod_{k=1}^n\frac{z-a_k}{1-\bar a_kz},\qquad\alpha\in\RR,\ n\in\ZZ_+,\ a_k\in\DD .
\]
\end{thmbox}
\begin{proof}
Пусть $a_1,\dots,a_n$ --- все нули $F$ в $\DD$ с учётом кратности ($n=0$, если нулей нет), и
$B=\prod_{k=1}^n\varphi_{a_k}$. Функция $B$ голоморфна в окрестности $\clos\DD$, $|B|=1$ на $\TT$, и её нули в
$\clos\DD$ --- в точности $a_k$ с теми же кратностями. Положим $G=F/B$ на $\clos\DD\setminus\{a_1,\dots,a_n\}$.
\begin{itemize}
\item В каждой точке $a_k$ кратности $m$ имеем $F=(z-a_k)^mF_1$, $B=(z-a_k)^mB_1$, $B_1(a_k)\ne0$, поэтому
$G=F_1/B_1$ продолжается голоморфно в $a_k$, и $G(a_k)\ne0$ (так как $F_1(a_k)\ne0$). Итак, $G\in H(\DD)$ и $G$ не
имеет нулей в $\DD$.
\item На $\TT$ функция $B$ непрерывна и не обращается в нуль, поэтому $G$ непрерывна на $\clos\DD$, и
$|G|=|F|/|B|=1$ на $\TT$; в частности, $G\ne0$ на $\clos\DD$.
\end{itemize}
По принципу максимума (Факт 0.5(б)) $|G|\le\max_\TT|G|=1$ в $\DD$. Функция $1/G$ также лежит в
$H(\DD)\cap C(\clos\DD)$ и $|1/G|=1$ на $\TT$, поэтому $|1/G|\le1$, т.е. $|G|\ge1$. Итак, $|G|\equiv1$ в $\DD$;
$|G|$ достигает максимума во внутренней точке, и $G\equiv e^{i\alpha}$. Значит, $F=e^{i\alpha}B$.
\end{proof}
\begin{remarkbox}{Замечания}
(1) На слайде написано $n\in\NN$; нужно допускать и $n=0$ (тогда $F\equiv e^{i\alpha}$; пример: $F\equiv1$).
(2) Условие конечности числа нулей на самом деле следует из остальных: так как $|F|=1$ на $\TT$ и $F$ равномерно
непрерывна на $\clos\DD$, то $|F|>1/2$ в кольце $1-\delta<|z|\le1$, и нули лежат в компакте
$\clos B(0,1-\delta)$, где их конечное число (теорема единственности; $F\not\equiv0$).
(3) Обратно, всякое конечное произведение Бляшке удовлетворяет условиям леммы, так что лемма даёт их
характеризацию: \emph{это в точности функции из $H(\DD)\cap C(\clos\DD)$ с $|F|=1$ на $\TT$}.
\end{remarkbox}

\subsubsection*{Формула Йенсена}

\begin{thmbox}{Лемма 12.2 (формула Йенсена)}
Пусть $f\in H(B(0,R))$, $f(0)\ne0$, и $\{a_n\}$ --- все нули $f$ в $B(0,R)$, занумерованные в порядке неубывания
модулей с учётом кратности. Тогда при любом $r\in(0,R)$
\begin{equation}\label{t0:jensen}
\frac1{2\pi}\int_0^{2\pi}\ln|f(re^{i\vartheta})|\,d\vartheta=\ln|f(0)|+\sum_{n:\,|a_n|<r}\ln\frac r{|a_n|}.
\end{equation}
\end{thmbox}
\begin{remarkbox}{Замечание}
Нулей может быть бесконечно много (на слайде $N\in\NN$; правильнее $N\in\NN\cup\{0,\infty\}$), но в каждом
круге $\clos B(0,r)$, $r<R$, их конечное число (теорема единственности), поэтому сумма в \eqref{t0:jensen}
конечна. Слагаемые с $|a_n|=r$ равны нулю, так что сумму можно брать и по $|a_n|\le r$.
\end{remarkbox}
\begin{proof}
\textbf{Шаг 1: $r\ne|a_n|$ для всех $n$.} Положим
\[
B(z)=\prod_{n:\,|a_n|<r}\frac{r(z-a_n)}{r^2-\bar a_nz}=\prod_{n:\,|a_n|<r}\varphi_{a_n/r}\Bigl(\frac zr\Bigr).
\]
Каждый множитель --- ДЛО, переводящее $C(0,r)$ в $\TT$ (так как $\varphi_c(\TT)=\TT$ при $c=a_n/r\in\DD$), поэтому
$|B(re^{i\vartheta})|=1$. Полюсы множителей --- точки $r^2/\bar a_n$ с $|r^2/\bar a_n|=r^2/|a_n|>r$ (при $a_n=0$
полюса нет --- но $a_n\ne0$, так как $f(0)\ne0$). Выберем $r'>r$ так, что $r'<R$, $r'<r^2/|a_n|$ для всех
$|a_n|<r$ и $r'<|a_m|$ для всех $|a_m|>r$ (возможно, так как в $\clos B(0,(r+R)/2)$ лишь конечное число нулей,
и ни один из них не лежит на $C(0,r)$). Тогда $B\in H(B(0,r'))$, а $F=f/B$ после устранения особенностей в
точках $a_n$ (нули $B$ и $f$ там одного порядка, как в доказательстве Леммы 12.1) голоморфна в $B(0,r')$ и не
имеет там нулей: нули $f$ в $B(0,r')$ --- это ровно $a_n$ с $|a_n|<r$.

Круг $B(0,r')$ односвязен, поэтому (Факт 0.9) существует $L\in H(B(0,r'))$, $e^L=F$, и
$\ln|F|=\Re L$. По теореме о среднем (Факт 0.5(а)) для $\Re L$ на $C(0,r)$:
\[
\ln|F(0)|=\frac1{2\pi}\int_0^{2\pi}\ln|F(re^{i\vartheta})|\,d\vartheta=\frac1{2\pi}\int_0^{2\pi}\ln|f(re^{i\vartheta})|\,d\vartheta,
\]
так как $|F|=|f|/|B|=|f|$ на $C(0,r)$. Наконец, $B(0)=\prod_{|a_n|<r}\frac{-ra_n}{r^2}=\prod(-a_n/r)$, поэтому
$\ln|F(0)|=\ln|f(0)|-\sum_{|a_n|<r}\ln\frac{|a_n|}r=\ln|f(0)|+\sum_{|a_n|<r}\ln\frac r{|a_n|}$. Формула
\eqref{t0:jensen} доказана для таких $r$.

\textbf{Шаг 2: вспомогательный интеграл.} Для $c\in\CC$, $r>0$:
\begin{equation}\label{t0:logint}
\frac1{2\pi}\int_0^{2\pi}\ln|re^{i\vartheta}-c|\,d\vartheta=\ln\max(r,|c|).
\end{equation}
\begin{itemize}
\item $|c|>r$: $\ln|re^{i\vartheta}-c|=\ln|c|+\ln|1-re^{i\vartheta}/c|$, а $z\mapsto1-z/c$ не имеет нулей в
$B(0,|c|)\supset\clos B(0,r)$; по теореме о среднем для $\ln|1-z/c|$ (Факт 0.5(а)) среднее второго слагаемого равно
$\ln|1-0|=0$.
\item $|c|<r$: $|re^{i\vartheta}-c|=r|1-ce^{-i\vartheta}/r|=r|1-\bar ce^{i\vartheta}/r|$ (модуль сопряжённого
числа), так что $\ln|re^{i\vartheta}-c|=\ln r+\ln|1-\bar ce^{i\vartheta}/r|$. Функция
$z\mapsto1-\bar cz/r$ не имеет нулей в $\clos\DD$ (так как $|c|/r<1$), и среднее $\ln|1-\bar ce^{i\vartheta}/r|$ по
теореме о среднем равно $0$.
\item $|c|=r$, $c=re^{i\theta_0}$: $|re^{i\vartheta}-c|=r|1-e^{i(\vartheta-\theta_0)}|=2r|\sin\frac{\vartheta-\theta_0}2|$.
В силу $2\pi$-периодичности достаточно проверить $\int_0^{2\pi}\ln(2|\sin\frac\vartheta2|)\,d\vartheta=0$, т.е.
$I:=\int_0^{\pi}\ln\sin u\,du=-\pi\ln2$ (замена $\vartheta=2u$). Интеграл сходится ($\ln\sin u\sim\ln u$ при
$u\to0+$, аналогично у $\pi$). Пусть $J=\int_0^{\pi/2}\ln\sin u\,du=\int_0^{\pi/2}\ln\cos u\,du$; тогда $I=2J$ и
\[
2J=\int_0^{\pi/2}\ln\frac{\sin2u}2\,du=-\frac\pi2\ln2+\frac12\int_0^{\pi}\ln\sin v\,dv=-\frac\pi2\ln2+J,
\]
откуда $J=-\frac\pi2\ln2$, $I=-\pi\ln2$, и $\int_0^{2\pi}\ln(2|\sin\frac\vartheta2|)d\vartheta=2\pi\ln2+2I=0$.
\end{itemize}

\textbf{Шаг 3: $r=|a_k|$ для некоторых $k$} (то, что на слайде оставлено как упражнение). Пусть $b_1,\dots,b_m$ ---
нули $f$ на $C(0,r)$ (с учётом кратности). Тогда $f(z)=h(z)\prod_{j=1}^m(z-b_j)$, где $h\in H(B(0,R))$ не имеет
нулей на $C(0,r)$, $h(0)=f(0)/\prod_j(-b_j)\ne0$, и нули $h$ в $B(0,r)$ --- те же $a_n$, $|a_n|<r$. По шагу 1
\[
\frac1{2\pi}\int_0^{2\pi}\ln|h(re^{i\vartheta})|\,d\vartheta=\ln|h(0)|+\sum_{|a_n|<r}\ln\frac r{|a_n|}
=\ln|f(0)|-m\ln r+\sum_{|a_n|<r}\ln\frac r{|a_n|}.
\]
Функция $\ln|f(re^{i\vartheta})|=\ln|h(re^{i\vartheta})|+\sum_j\ln|re^{i\vartheta}-b_j|$ имеет лишь конечное число
логарифмических (интегрируемых) особенностей, и по \eqref{t0:logint} каждый интеграл
$\frac1{2\pi}\int_0^{2\pi}\ln|re^{i\vartheta}-b_j|d\vartheta=\ln r$. Складывая, получаем \eqref{t0:jensen}.
\end{proof}
\begin{remarkbox}{Следствия}
(1) Правая часть \eqref{t0:jensen} равна $\ln|f(0)|+\int_0^r\frac{n(t)}t\,dt$, где $n(t)$ --- число нулей в
$\clos B(0,t)$; в частности, среднее $\ln|f|$ по окружности --- непрерывная неубывающая функция $r$.
(2) Оценка числа нулей: если $|f|\le M$ в $B(0,R)$, то число нулей в $\clos B(0,r)$ не больше
$\ln\frac M{|f(0)|}\big/\ln\frac Rr$ (подставить в \eqref{t0:jensen} радиус $R'\in(r,R)$ и оставить только нули
с $|a_n|\le r$, каждый из которых даёт вклад $\ge\ln(R'/r)$; затем $R'\to R$).
\end{remarkbox}

\subsubsection*{Теорема Каратеодори об аппроксимации}

\begin{thmbox}{Предложение 12.4 (теорема Каратеодори об аппроксимации)}
Пусть $f\in H(\DD)$ и $|f(z)|\le1$ при $|z|<1$. Тогда найдётся последовательность $\{B_k\}$ конечных
произведений Бляшке такая, что $B_k(z)\to f(z)$ при $k\to\infty$ при всех $z\in\DD$; более того, сходимость
локально равномерна в $\DD$ и $|f(z)-B_k(z)|\le2|z|^{k+1}$.
\end{thmbox}
Обозначим $\mathfrak B=\{f\in H(\DD):|f(z)|\le1\ \text{при}\ |z|<1\}$ и $c_k(g)$ --- $k$-й тейлоровский
коэффициент функции $g$ в нуле. Ключевой шаг --- интерполяционная лемма.

\begin{thmbox}{Лемма (интерполяция коэффициентов)}
Для любой $f\in\mathfrak B$ и любого $n\in\ZZ_+$ существует конечное произведение Бляшке $B_n$ степени $\le n+1$
такое, что $c_k(B_n)=c_k(f)$ при $k=0,1,\dots,n$.
\end{thmbox}
\begin{proof}
Индукция по $n$ (утверждение доказывается сразу для всех $f\in\mathfrak B$). Положим $c=f(0)=c_0(f)$.
Если $|c|=1$, то $|f|$ достигает максимума в точке $0\in\DD$, и по принципу максимума $f\equiv c$; тогда при
любом $n$ годится $B_n\equiv c$ (степень $0$). Далее считаем $|c|<1$.

\emph{База $n=0$:} $B_0(z)=\dfrac{z+c}{1+\bar cz}=\varphi_{-c}(z)$, $B_0(0)=c$.

\emph{Переход $n-1\to n$} ($n\ge1$). Рассмотрим
\[
g(z)=\frac{f(z)-c}{z\bigl(1-\bar cf(z)\bigr)}\quad(z\ne0),\qquad g(0)=\frac{f'(0)}{1-|c|^2}.
\]
\emph{Проверка $g\in\mathfrak B$ (упражнение слайда).} Знаменатель $1-\bar cf(z)\ne0$, так как $|\bar cf(z)|\le|c|<1$.
Функция $u=\varphi_c\circ f$ голоморфна в $\DD$, $u(0)=0$ и $|u|\le1$ (так как $|\varphi_c(w)|\le1$ при $|w|\le1$ по
\eqref{t0:phi}). По лемме Шварца $|u(z)|\le|z|$, а $g(z)=u(z)/z$; значит, $g$ имеет в нуле устранимую
особенность, $g\in H(\DD)$ и $|g|\le1$, т.е. $g\in\mathfrak B$. (Если вдобавок $f\in C(\clos\DD)$, то и
$g\in C(\clos\DD)$, так что класс $\mathfrak B_C=\mathfrak B\cap C(\clos\DD)$ слайда также сохраняется.)

Обращая формулу для $g$: $zg(1-\bar cf)=f-c\iff f(1+\bar czg)=zg+c$, т.е.
\begin{equation}\label{t0:fg}
f(z)=\frac{zg(z)+c}{1+\bar czg(z)}=\varphi_{-c}\bigl(zg(z)\bigr)
\end{equation}
(знаменатель $\ne0$, так как $|\bar czg|<1$). По предположению индукции найдётся конечное произведение Бляшке
$\tilde B=\tilde B_{n-1}$ степени $m\le n$ с $c_k(\tilde B)=c_k(g)$ при $k=0,\dots,n-1$. Положим
\[
B(z)=\frac{z\tilde B(z)+c}{1+\bar cz\tilde B(z)}=\varphi_{-c}\bigl(z\tilde B(z)\bigr).
\]
\emph{$B$ --- конечное произведение Бляшке степени $m+1$ (упражнение слайда).} Функция $P=z\tilde B$ ---
произведение Бляшке степени $m+1$: голоморфна в окрестности $\clos\DD$, $|P|=1$ на $\TT$, $|P|<1$ в $\DD$. Значит,
$B=\varphi_{-c}\circ P$ голоморфна в окрестности $\clos\DD$ (знаменатель $1+\bar cP\ne0$ на $\clos\DD$), и
$|B|=1$ на $\TT$, так как $\varphi_{-c}(\TT)=\TT$. Нули $B$ в $\DD$ --- это нули функции $P+c$; по теореме Руше
($|c|<1=|P|$ на $\TT$) их столько же, сколько нулей у $P$, т.е. $m+1$ (конечное число). По Лемме 12.1 $B$ ---
конечное произведение Бляшке, и его степень (число нулей) равна $m+1\le n+1$.

\emph{Совпадение коэффициентов.} По \eqref{t0:fg} и тождеству \eqref{t0:id1} (с $a=-c$):
\[
f(z)-B(z)=\varphi_{-c}(zg)-\varphi_{-c}(z\tilde B)=\frac{(1-|c|^2)\,z\bigl(g(z)-\tilde B(z)\bigr)}
{\bigl(1+\bar czg(z)\bigr)\bigl(1+\bar cz\tilde B(z)\bigr)} .
\]
Знаменатель голоморфен и не обращается в нуль в $\DD$, а $g-\tilde B$ имеет в нуле нуль порядка $\ge n$
(совпадают коэффициенты с номерами $0,\dots,n-1$). Значит, $f-B$ имеет в нуле нуль порядка $\ge n+1$, т.е.
$c_k(B)=c_k(f)$ при $k=0,\dots,n$. Полагаем $B_n:=B$.
\end{proof}
\begin{remarkbox}{Замечание (неточность слайда)}
На слайде сказано, что $f-B$ имеет в нуле <<ноль порядка $n$>>; на самом деле порядок не меньше $n+1$ (иначе
совпадали бы лишь коэффициенты $c_0,\dots,c_{n-1}$, и индукция не замкнулась бы). Кроме того, предварительная
редукция к $f\in C(\clos\DD)$ (через $f(rz)$) не нужна: приведённое рассуждение использует только $|f|\le1$ в $\DD$.
\end{remarkbox}
\begin{proof}[Доказательство Предложения 12.4]
Пусть $B_n$ из леммы. Функция $h=(f-B_n)/2$ голоморфна в $\DD$, $|h|\le(|f|+|B_n|)/2\le1$, и имеет в нуле нуль
порядка $\ge n+1$ (или $h\equiv0$). По Замечанию 12.1(2) $|h(z)|\le|z|^{n+1}$, т.е.
$|f(z)-B_n(z)|\le2|z|^{n+1}$. Следовательно, $\sup_{|z|\le r}|f-B_n|\le2r^{n+1}\to0$ при каждом $r<1$: $B_n\to f$
локально равномерно в $\DD$ (в частности, поточечно). Этот последний шаг на слайде опущен.
\end{proof}

\subsubsection*{Лемма Бореля--Каратеодори}

\begin{thmbox}{Предложение 12.5 (лемма Бореля--Каратеодори)}
Пусть $f\in H(\clos B(0,R))$, $M(r)=\max_{|z|=r}|f(z)|$, $A(r)=\max_{|z|=r}\Re f(z)$. Тогда при $0<r<R$
\[
M(r)\le\frac{R+r}{R-r}\,|f(0)|+\frac{2r}{R-r}\,A(R).
\]
\end{thmbox}
\begin{proof}
\textbf{Постоянная $f\equiv c$.} Нужно $(R-r)|c|\le(R+r)|c|+2r\Re c$, т.е. $0\le2r(|c|+\Re c)$ --- верно.
Далее $f\not\equiv\mathrm{const}$.

\textbf{Случай $f(0)=0$.} \emph{Проверка $A(R)>0$ (упражнение слайда).} Функция $\Re f$ --- вещественная часть
голоморфной функции, по теореме о среднем (Факт 0.5(а)) $\frac1{2\pi}\int_0^{2\pi}\Re f(Re^{i\vartheta})d\vartheta
=\Re f(0)=0$. Если бы $A(R)\le0$, то непрерывная функция $\Re f\le0$ на $C(0,R)$ имела бы нулевое среднее, т.е.
$\Re f\equiv0$ на $C(0,R)$; по Факту 0.5(в), применённому к $f$ и к $-f$, $\max_{\clos B(0,R)}|\Re f|=0$, т.е.
$\Re f\equiv0$ в круге, и $f$ постоянна (непостоянная голоморфная функция открыта, Факт 0.11, а образ лежит на
мнимой оси) --- противоречие. Итак, $A(R)>0=A(0)$.

\emph{Знаменатель (упражнение слайда).} По Факту 0.5(в) имеем $\Re f(z)\le A(R)$ при всех $|z|\le R$.
Поэтому $\Re(2A(R)-f(z))\ge A(R)>0$, т.е. знаменатель $2A(R)-f(z)$ не обращается в нуль
на $\clos B(0,R)$. Значит,
\[
g(z)=\frac{f(z)}{2A(R)-f(z)}\in H(B(0,R))\cap C(\clos B(0,R)),\qquad g(0)=0 .
\]
\emph{Оценка $|g|\le1$.} При $f=u+iv$: $|g|^2=\dfrac{u^2+v^2}{(2A(R)-u)^2+v^2}\le1\iff u^2\le(2A(R)-u)^2
\iff4A(R)\bigl(A(R)-u\bigr)\ge0$, что верно, так как $A(R)>0$ и $u\le A(R)$.

По Замечанию 12.1(1) (с $M=1$, $a=0$) $|g(z)|\le|z|/R$; при $|z|=r$ имеем $|g(z)|\le r/R<1$. Из
$g(2A(R)-f)=f$ находим $f=\dfrac{2A(R)g}{1+g}$ ($1+g\ne0$, так как $|g|<1$), и при $|z|=r$
\[
|f(z)|\le\frac{2A(R)|g(z)|}{1-|g(z)|}\le\frac{2A(R)\,r/R}{1-r/R}=\frac{2r}{R-r}A(R)
\]
(функция $t\mapsto t/(1-t)$ возрастает на $[0,1)$). Это и есть утверждение при $f(0)=0$.

\textbf{Случай $f(0)\ne0$.} Применим доказанное к $h=f-f(0)$ ($h(0)=0$; если $h\equiv0$, то $f$ постоянна).
Так как $A_h(R)=A(R)-\Re f(0)\le A(R)+|f(0)|$, то при $|z|=r$
\[
|f(z)|\le|f(0)|+|h(z)|\le|f(0)|+\frac{2r}{R-r}\bigl(A(R)+|f(0)|\bigr)=\frac{R+r}{R-r}|f(0)|+\frac{2r}{R-r}A(R).
\qedhere
\]
\end{proof}
\begin{remarkbox}{Смысл леммы}
Модуль голоморфной функции оценивается через \emph{одностороннюю} оценку её вещественной части (сверху!) на
большей окружности. Это ключ к доказательству малой теоремы Пикара ниже: если $f=e^g$, то $A_g=\log M_f$
известна, а $M_g$ --- нет.
\end{remarkbox}

\begin{detail}{Упражнение слайда: может ли $|f(\frac{1+i}2)|=14$?}
Условие: $f\in H(\clos\DD)$, $f(0)=0$, $|\Re f(z)|\le e^x$ при $z=x+iy\in\DD$ (по непрерывности --- и на
$\clos\DD$). \textbf{Ответ: нет.} Точка $z_0=\frac{1+i}2$ имеет модуль $r=1/\sqrt2$. На $\TT$ имеем
$\Re f(z)\le e^x\le e$, т.е. $A(1)\le e$. По Предложению 12.5 с $R=1$, $f(0)=0$:
\[
|f(z_0)|\le M(r)\le\frac{2r}{1-r}A(1)\le\frac{2/\sqrt2}{1-1/\sqrt2}\,e=\frac{2}{\sqrt2-1}\,e=2(\sqrt2+1)e\approx13{,}13<14 .
\]
(Использование меньших $R\in(r,1)$ с $A(R)\le e^R$ не улучшает оценку: функция $R\mapsto\ln\frac{2re^R}{R-r}$
имеет производную $1-\frac1{R-r}<0$ при $R-r<1$, так что минимум при $R=1$.)
\end{detail}

\subsubsection*{Малая теорема Пикара для класса $\mathcal E$}

\begin{thmbox}{Теорема 12.1 (малая теорема Пикара)}
Пусть $f$ --- целая функция, не являющаяся постоянной. Тогда $f$ принимает в $\CC$ все значения, за
исключением, быть может, одного.
\end{thmbox}
Пример $e^z\ne0$ показывает, что одно исключение возможно. В лекциях теорема доказывается для
$f\in\mathcal E=\bigcup_n\mathcal E_{(n)}$ (класс целых функций обобщённого порядка не выше итерированной
экспоненты), с помощью леммы Бореля--Каратеодори.

\begin{thmbox}{Лемма 12.3}
Пусть $D$ --- односвязная область в $\CC$, $f\in H(D)$ и $f(z)\ne0$ при $z\in D$. Тогда существует $g\in H(D)$ с
$f=\exp(g)$, а также для любого $m\in\NN$, $m\ge2$, существует $h\in H(D)$ с $f=h^m$.
\end{thmbox}
Это Факт 0.9 (там же доказательство): $g$ --- первообразная $f'/f$, подправленная константой; $h=e^{g/m}$.

\begin{detail}{Вспомогательные свойства итерированного логарифма и классов $\mathcal E_{(n)}$}
\textbf{(i)} \emph{Для $C\ge1$ и $n\ge1$: $\log_{(n)}(Ct)\le\log_{(n)}t+\log C$ при всех достаточно больших $t$.}
При $n=1$ это равенство. Если неравенство верно для $n$ и $t$ так велико, что $\log_{(n)}t\ge1$, то, пользуясь
$\log(x+a)=\log x+\log(1+a/x)\le\log x+a/x\le\log x+a$ при $x\ge1$, $a\ge0$:
$\log_{(n+1)}(Ct)\le\log\bigl(\log_{(n)}t+\log C\bigr)\le\log_{(n+1)}t+\log C$.

\textbf{(ii)} \emph{$\mathcal E_{(1)}$ --- многочлены.} Если $\limsup\frac{\log M_f(r)}{\log r}<p<\infty$, то
$M_f(r)\le r^p$ при $r\ge r_0$, и по неравенствам Коши $|c_n|\le r^{p-n}\to0$ при $n>p$: $f$ --- многочлен.
Обратно, для многочлена степени $d$: $M_f(r)\le Cr^d$ при $r\ge1$, и $\limsup\le d$.

\textbf{(iii)} \emph{Если $f\in\mathcal E_{(n)}$, $a\in\CC$, то $f-a\in\mathcal E_{(n)}$.} Для непостоянной $f$
имеем $M_f(r)\to\infty$, поэтому $M_{f-a}(r)\le M_f(r)+|a|\le2M_f(r)$ при больших $r$, и по (i)
$\log_{(n)}M_{f-a}(r)\le\log_{(n)}M_f(r)+\log2$; деление на $\log r\to\infty$ не меняет конечности $\limsup$.

\textbf{(iv)} \emph{Если $f=e^g$, $g$ целая, то $M_f(r)=e^{A_g(r)}$}, так как $|f|=e^{\Re g}$ и $t\mapsto e^t$
возрастает. Отсюда $\log_{(n+1)}M_f(r)=\log_{(n)}A_g(r)$ при больших $r$.

\textbf{(v)} \emph{Упражнение слайда: если $g$ целая и непостоянная, то $A_g(t)\to\infty$.} $A_g$ не убывает
(Факт 0.5(в)); если $A_g\le C$, то $\Re g\le C$ всюду, $|e^g|\le e^C$, по теореме Лиувилля $e^g\equiv\mathrm{const}$,
откуда $g'e^g\equiv0$, $g'\equiv0$ --- противоречие.
\end{detail}

\begin{proof}[Доказательство Теоремы 12.1 для $f\in\mathcal E$]
Докажем индукцией по $n$: \emph{всякая непостоянная $f\in\mathcal E_{(n)}$ принимает все значения из $\CC$,
кроме, быть может, одного.}

\emph{База $n=1$.} По (ii) $f$ --- непостоянный многочлен, и по основной теореме алгебры уравнение $f(z)=a$
имеет корни при любом $a\in\CC$.

\emph{Переход $N\to N+1$.} Пусть $f\in\mathcal E_{(N+1)}$ непостоянна и $f(z)\ne a$ при всех $z$. По (iii)
$f-a\in\mathcal E_{(N+1)}$, поэтому, заменив $f$ на $f-a$, считаем $a=0$. Так как $\CC$ односвязна и $f\ne0$, по
Лемме 12.3 $f=e^g$, $g$ целая и непостоянная. По (iv)
\[
\limsup_{r\to\infty}\frac{\log_{(N)}A_g(r)}{\log r}=\limsup_{r\to\infty}\frac{\log_{(N+1)}M_f(r)}{\log r}=:L<\infty .
\]
Лемма Бореля--Каратеодори для $g$ с $R=2r$ ($\frac{R+r}{R-r}=3$, $\frac{2r}{R-r}=2$):
$M_g(r)\le2A_g(2r)+3|g(0)|$. По (v) $A_g(2r)\ge3|g(0)|$ при $r\ge r_0$, и тогда $M_g(r)\le3A_g(2r)$. По монотонности
$\log_{(N)}$ и (i) при больших $r$
\[
\frac{\log_{(N)}M_g(r)}{\log r}\le\frac{\log_{(N)}A_g(2r)+\log3}{\log r}
=\frac{\log_{(N)}A_g(2r)}{\log(2r)}\cdot\frac{\log(2r)}{\log r}+\frac{\log3}{\log r}.
\]
Так как $\frac{\log(2r)}{\log r}\to1$, $\frac{\log 3}{\log r}\to0$, а $\limsup_{r}\frac{\log_{(N)}A_g(2r)}{\log(2r)}=L$
(это тот же верхний предел по $s=2r\to\infty$), получаем $\limsup\frac{\log_{(N)}M_g(r)}{\log r}\le L<\infty$,
т.е. $g\in\mathcal E_{(N)}$.

По предположению индукции $g$ принимает все значения, кроме, быть может, одного. Пусть $w\in\CC$. Среди
попарно различных чисел $w+2\pi ik$, $k\in\ZZ$, $g$ пропускает не более одного, поэтому найдутся бесконечно
много $k$ и точки $z_k$ (попарно различные, так как различны $g(z_k)$) с $g(z_k)=w+2\pi ik$; тогда
$f(z_k)=e^{w+2\pi ik}=e^w$. Так как $e^w$ пробегает все $\CC\setminus\{0\}$, $f$ принимает каждое ненулевое
значение (даже бесконечно много раз) и пропускает только $0$. Возвращаясь к исходной функции ($f-a$ вместо $f$),
получаем: $f$ пропускает только значение $a$.
\end{proof}
\begin{remarkbox}{Замечания}
(1) Попутно доказано больше: непостоянная $f\in\mathcal E$, пропускающая значение $a$, принимает любое
другое значение бесконечно много раз. (2) Класс $\mathcal E_{(2)}$ --- целые функции конечного порядка;
$\mathcal E$ очень широк, но не содержит всех целых функций (например, функций, растущих быстрее любой
конечной башни экспонент), поэтому общая малая теорема Пикара требует иных методов (модулярная функция или
теорема Блоха; см. вопросы со звёздочкой).
\end{remarkbox}

\begin{remarkbox}{Коротко: что нужно уметь по главе 12}
\begin{enumerate}
\item Лемма Шварца: $g=f/z$, принцип максимума на $C(0,r)$, $r\to1$; случай равенства.
\item Модификации: круг $B(a,R)$ и оценка $M|z-a|/R$; нуль порядка $m$ даёт $|z|^m$.
\item $\CAut(\DD)=\{e^{i\alpha}\varphi_a\}$: Шварц для $h=f\circ\varphi_{-a}$ и для $h^{-1}$.
\item Шварц--Пик: $g=\varphi_{f(w)}\circ f\circ\varphi_{-w}$; формулы $\varphi_c'$; $\rho$ --- инвариантная метрика,
$K(a,r)$ --- евклидов круг.
\item Лемма 12.1: деление на произведение Бляшке и принцип максимума для $G$ и $1/G$.
\item Йенсен: $F=f/B$, $\ln|F|=\Re\ln F$, теорема о среднем; нули на окружности --- через
$\frac1{2\pi}\int\ln|re^{i\vartheta}-c|=\ln\max(r,|c|)$.
\item Каратеодори: $f=\varphi_{-c}(zg)$, индукция по числу совпадающих коэффициентов, $|f-B_n|\le2|z|^{n+1}$.
\item Борель--Каратеодори: $g=f/(2A(R)-f)$, Шварц, обращение.
\item Пикар для $\mathcal E$: $f=e^g$, $A_g=\log M_f$, Борель--Каратеодори ($R=2r$) понижает класс на единицу.
\end{enumerate}
\end{remarkbox}

% ===== part 1b =====
% ===================== Часть 1b: билеты 1--4 =====================

\ticket{1}{Однолистные функции. Критерий локальной однолистности}

\begin{defs}
\textbf{Окрестности.} $B(a,r)=\{z:|z-a|<r\}$ --- открытый круг, $\clos{B}(a,r)$ --- замкнутый круг,
$C(a,r)=\{|z-a|=r\}$ --- окружность; $U(a)$ --- какая-нибудь окрестность точки $a$.

\textbf{Голоморфность в точке.} $f\in H(a)$, если $f$ голоморфна в некотором круге $B(a,r)$, $r>0$.
$H(D)$ --- множество функций, голоморфных в открытом множестве $D$.

\textbf{Конформность в точке (Опр.~2.7 лекций).} Пусть $f$ дифференцируема в точке $a$ в вещественном смысле
($f\in D_{\RR}(a)$). Говорят, что $f$ \emph{конформна в точке $a$} ($f\in\mathrm{Conf}(a)$), если дифференциал
$df(a)$ (как линейное отображение $\Delta z\mapsto df(a)(\Delta z)$) --- композиция гомотетии с положительным
коэффициентом и поворота. \emph{Критерий (Предл.~2.7):} $f\in\mathrm{Conf}(a)\iff f$ комплексно дифференцируема
в $a$ и $f'(a)\ne0$; тогда $df(a)\Delta z=f'(a)\Delta z$ --- растяжение в $|f'(a)|$ раз и поворот на $\arg f'(a)$.

\textbf{Локально конформное и конформное отображения (Опр.~2.9).} $f\in\mathrm{Conf}_0(D)$ (локально конформна в
области $D$), если $f\in\mathrm{Conf}(z)$ для всех $z\in D$; $f\in\mathrm{Conf}(D)$ (конформна в $D$), если $f$
локально конформна и взаимно однозначна в $D$. Для области $D\subset\CC$:
$f\in\mathrm{Conf}_0(D)\iff f\in H(D)$ и $f'\ne0$ в $D$; $f\in \mathrm{Conf}(D)\iff f\in H(D)$, $f'\ne 0$ в $D$
и $f$ инъективна. Запись $f\in\mathrm{Conf}(D,G)$ означает: $f\in \mathrm{Conf}(D)$ и $f(D)=G$
(\emph{конформное отображение $D$ на $G$}).

\textbf{Однолистность (Опр.~13.1).} Пусть $a\in\CC$, $f\in H(a)$. Функция $f$ \emph{однолистна в точке $a$}
($f\in\mathcal U(a)$), если существует $\delta>0$ такое, что $f$ взаимно однозначна (инъективна) в $B(a,\delta)$.
Пусть $D\subset\CC$ --- область, $f\in H(D)$. Функция $f$ \emph{локально однолистна в $D$} ($f\in\mathcal U_\ell(D)$),
если она однолистна в каждой точке $D$; $f$ \emph{однолистна в $D$} ($f\in\mathcal U(D)$), если $f$ инъективна в $D$.

\textbf{Локальная обратимость.} $f\in H(a)$ \emph{локально обратима в точке $a$}, если существуют окрестности
$U=U(a)$, $\widetilde U=U(f(a))$ и $g\in H(\widetilde U)$ такие, что $g(f(z))=z$ для всех $z\in U$.

\textbf{Порядок нуля.} Если $f\in H(a)$, $f\not\equiv f(a)$ вблизи $a$, то существуют $n\in\NN$ и $g\in H(a)$,
$g(a)\ne0$, с $f(z)-f(a)=(z-a)^ng(z)$; $n$ --- кратность значения $f(a)$ в точке $a$ (порядок нуля $f-f(a)$).
Ясно, что $n=1\iff f'(a)\neq0$.

\textbf{Примеры.} $e^z\in\mathcal U_\ell(\CC)$ ($(e^z)'=e^z\ne0$), но $e^z\notin\mathcal U(\CC)$, так как $e^{z+2\pi i}=e^z$;
$e^z\in\mathcal U(\{|\Im z|<\pi\})$. Функция $z^2$ однолистна в каждой точке $a\ne0$ и не однолистна в точке $0$
($(\pm\eps)^2$ совпадают); $z^2\in\mathcal U(\{\Re z>0\})$. Дробно-линейная функция однолистна в $\CC$ без полюса.
\end{defs}

Очевидные соотношения: $\mathcal U(D)\subset\mathcal U_\ell(D)$ (если $f$ инъективна во всей $D$, то она инъективна
и в любом круге $B(a,\delta)\subset D$), но обратное включение неверно (пример $e^z$).

\subsection{Критерий однолистности в точке}

\begin{thmbox}{Предложение 13.1 (критерии однолистности)}
\begin{enumerate}
\item Пусть $a\in\CC$, $f\in H(a)$. Тогда $f\in\mathcal U(a)\iff f'(a)\ne0$.
\item Пусть $D\subset\CC$ --- область, $f\in H(D)$. Тогда $f\in \mathcal U(D)\iff f$ конформна в $D$ (т.е.
$f\in\mathrm{Conf}(D)$). Более того, если $f\in\mathcal U(D)$, то $G=f(D)$ --- область и $f\in \mathrm{Conf}(D,G)$, причём
обратная функция $f^{-1}$ голоморфна в $G$ и $(f^{-1})'(w)=1/f'(f^{-1}(w))$.
\end{enumerate}
В частности, $f\in\mathcal U_\ell(D)\iff f'(z)\neq0$ при всех $z\in D\iff f\in\mathrm{Conf}_0(D)$.
\end{thmbox}

\begin{proof}
\textbf{1), достаточность ($f'(a)\ne0\Rightarrow f\in\mathcal U(a)$).}
Пусть $f\in H(B(a,r))$, $r>0$, и $f'(a)\ne0$. Так как $f'\in H(B(a,r))$ (производная голоморфной функции голоморфна)
и, в частности, непрерывна в $a$, найдётся $\rho\in(0,r)$ такое, что
\begin{equation}\label{t1:near}
|f'(z)-f'(a)|<\tfrac12|f'(a)|\qquad\text{при } z\in B(a,\rho).
\end{equation}
Возьмём $z_1,z_2\in B(a,\rho)$, $z_1\ne z_2$, и обозначим $I=[z_1,z_2]$ --- отрезок; он лежит в $B(a,\rho)$ по
выпуклости круга. Так как $f$ --- первообразная для $f'$ в круге, по формуле Ньютона--Лейбница
\[
f(z_2)-f(z_1)=\int_I f'(w)\,dw=\int_I\bigl(f'(w)-f'(a)\bigr)\,dw+f'(a)(z_2-z_1)
\]
(мы воспользовались тем, что $\int_I dw=z_2-z_1$). Отсюда по стандартной оценке интеграла и \eqref{t1:near}
\[
\bigl|f(z_2)-f(z_1)-f'(a)(z_2-z_1)\bigr|\le\int_I|f'(w)-f'(a)|\,|dw|<\tfrac12|f'(a)|\,|z_2-z_1|
\]
(строгое неравенство: подынтегральная функция непрерывна и строго меньше $\frac12|f'(a)|$ на компакте $I$, а
длина $I$ равна $|z_2-z_1|>0$). По неравенству треугольника
\[
|f(z_2)-f(z_1)|\ge|f'(a)|\,|z_2-z_1|-\bigl|f(z_2)-f(z_1)-f'(a)(z_2-z_1)\bigr|>\tfrac12|f'(a)|\,|z_2-z_1|>0 .
\]
Значит, $f(z_1)\ne f(z_2)$, т.е. $f$ инъективна в $B(a,\rho)$ и $f\in\mathcal U(a)$.
Заметим попутно, что из \eqref{t1:near} следует $f'(z)\ne0$ при $z\in B(a,\rho)$, и мы получили оценку
\begin{equation}\label{t1:bilip}
|f(z_2)-f(z_1)|\ge\tfrac12|f'(a)|\,|z_2-z_1|,\qquad z_1,z_2\in B(a,\rho).
\end{equation}

\textbf{1), необходимость ($f'(a)=0\Rightarrow f\notin\mathcal U(a)$).}
Пусть $f\in H(B(a,r))$ и $f'(a)=0$. Если $f\equiv f(a)$ в $B(a,r)$, то $f$, очевидно, не инъективна ни в каком круге
$B(a,\delta)$. Пусть $f\not\equiv f(a)$. Тогда $a$ --- изолированный нуль функции $f-f(a)$ конечного порядка $n$, и
$n\ge2$, поскольку $f'(a)=0$:
\[
f(z)-f(a)=(z-a)^ng(z),\qquad g\in H(B(a,r)),\ g(a)\ne0,\ n\ge2.
\]
По непрерывности $g$ найдётся $\delta\in(0,r)$ такое, что $g(z)\ne0$ при $z\in B(a,\delta)$. Круг $B(a,\delta)$ односвязен,
поэтому по Лемме~12.3 (существование голоморфного корня из функции без нулей в односвязной области) найдётся
$g_1\in H(B(a,\delta))$ с $g=g_1^{\,n}$. Ясно, что $g_1(a)\ne0$. Положим $h_1(z)=(z-a)g_1(z)$; тогда
\[
f(z)=f(a)+\bigl(h_1(z)\bigr)^n,\qquad z\in B(a,\delta),\qquad h_1(a)=0,\quad h_1'(a)=g_1(a)\ne0 .
\]
По уже доказанной достаточности найдётся $\rho\in(0,\delta)$ такое, что $h_1$ инъективна в $B(a,\rho)$.

Пусть теперь $\delta'>0$ произвольно; покажем, что $f$ не инъективна в $B(a,\delta')$. Положим
$\rho'=\min(\rho,\delta')$. Функция $h_1$ непостоянна в $B(a,\rho')$, поэтому по принципу сохранения области
(теорема об открытости) $h_1(B(a,\rho'))$ --- открытое множество, содержащее $h_1(a)=0$; значит,
$h_1(B(a,\rho'))\supset B(0,t)$ при некотором $t>0$. Возьмём $w_1=t/2$ и $w_2=w_1e^{2\pi i/n}$. Так как $n\ge2$,
то $w_2\ne w_1$, но $w_1^n=w_2^n$. Найдутся $z_1,z_2\in B(a,\rho')$ с $h_1(z_j)=w_j$; поскольку $w_1\ne w_2$, то
$z_1\ne z_2$. При этом
\[
f(z_1)=f(a)+w_1^n=f(a)+w_2^n=f(z_2).
\]
Итак, в любом круге $B(a,\delta')$ функция $f$ принимает некоторое значение в двух различных точках, т.е.
$f\notin\mathcal U(a)$.

\textbf{2).} Пусть $f\in\mathcal U(D)$. Для каждого $a\in D$ возьмём круг $B(a,\delta)\subset D$; $f$ инъективна в нём,
т.е. $f\in\mathcal U(a)$, и по п.~1) $f'(a)\ne0$. Итак, $f\in H(D)$, $f'\ne0$ в $D$ и $f$ инъективна, т.е.
$f\in\mathrm{Conf}(D)$ по Предл.~2.9. Обратно, если $f\in\mathrm{Conf}(D)$, то $f$ инъективна по определению, т.е.
$f\in\mathcal U(D)$.

Пусть $f\in\mathcal U(D)$. Функция $f$ непостоянна, поэтому по принципу сохранения области $G=f(D)$ --- область, и
$f$ переводит открытые подмножества $D$ в открытые множества (применяем ту же теорему к ограничениям $f$ на круги).
Значит, $f^{-1}:G\to D$ непрерывна (прообраз открытого множества $V\subset D$ при $f^{-1}$ равен $f(V)$ и открыт),
т.е. $f$ --- гомеоморфизм $D$ на $G$. Голоморфность $f^{-1}$: пусть $w_0\in G$, $z_0=f^{-1}(w_0)$, $w\to w_0$,
$w\ne w_0$, $z=f^{-1}(w)$. Тогда $z\ne z_0$ (инъективность), $z\to z_0$ (непрерывность $f^{-1}$), и
\[
\frac{f^{-1}(w)-f^{-1}(w_0)}{w-w_0}=\frac{z-z_0}{f(z)-f(z_0)}\longrightarrow\frac1{f'(z_0)},
\]
так как $f'(z_0)\neq0$. Следовательно, $f^{-1}\in H(G)$, $(f^{-1})'=1/f'\circ f^{-1}\ne0$, и $f\in\mathrm{Conf}(D,G)$.

Последнее утверждение (о $\mathcal U_\ell$) --- это п.~1), применённый к каждой точке $a\in D$, плюс Предл.~2.9.
\end{proof}

\begin{remarkbox}{Сравнение с вещественным случаем}
Для отображений $\RR\to\RR$ условие $f'(a)\ne0$ лишь \emph{достаточно} для локальной инъективности
($x\mapsto x^3$ инъективна, хотя $f'(0)=0$). Для голоморфных функций оно ещё и \emph{необходимо}: если $f'(a)=0$, то
вблизи $a$ функция «наматывает» окрестность $n\ge2$ раз. Заметим также, что якобиан отображения $f=u+iv$ равен
$u_xv_y-u_yv_x=u_x^2+v_x^2=|f'|^2$ (условия Коши--Римана), так что достаточность в п.~1) можно было бы получить
и из вещественной теоремы об обратной функции; приведённое доказательство элементарно и даёт явную оценку
\eqref{t1:bilip}.
\end{remarkbox}

\subsection{Локальная структура голоморфной функции в критической точке}

Доказательство необходимости даёт больше: точное описание поведения $f$ около точки, где $f'(a)=0$.

\begin{thmbox}{Следствие (локальная $n$-листность)}
Пусть $f\in H(a)$, $f\not\equiv\mathrm{const}$, и $a$ --- точка кратности $n\ge1$ значения $f(a)$
(т.е. $f'(a)=\ldots=f^{(n-1)}(a)=0\ne f^{(n)}(a)$). Тогда существуют окрестность $V$ точки $a$ и $t>0$ такие, что
каждое значение $w\in B(f(a),t^n)\setminus\{f(a)\}$ принимается в $V$ ровно в $n$ различных точках, в каждой из
них $f'\ne0$, а значение $f(a)$ принимается в $V$ только в точке $a$.
\end{thmbox}
\begin{proof}
При $n=1$ это следует из достаточности в Предл.~13.1 (с $V=f^{-1}(B(f(a),t))\cap B(a,\rho)$, где $B(f(a),t)\subset f(B(a,\rho))$).
Пусть $n\ge2$; в обозначениях доказательства $f=f(a)+h_1^n$ на $B(a,\rho)$, где $h_1$ инъективна в $B(a,\rho)$ и
$h_1'\neq0$ там (см. \eqref{t1:near} для $h_1$). Выберем $t>0$ с $B(0,t)\subset h_1(B(a,\rho))$ и положим
$V=h_1^{-1}(B(0,t))\cap B(a,\rho)$ --- открытая окрестность $a$, и $h_1:V\to B(0,t)$ --- биекция. Пусть
$0<|w-f(a)|<t^n$. Уравнение $s^n=w-f(a)$ имеет ровно $n$ различных корней $s_1,\ldots,s_n$, все они по модулю
равны $|w-f(a)|^{1/n}<t$, т.е. лежат в $B(0,t)$. Поэтому решения уравнения $f(z)=w$ в $V$ --- это ровно $n$ различных
точек $z_k=h_1^{-1}(s_k)$, и $f'(z_k)=n\,s_k^{n-1}h_1'(z_k)\ne0$. Наконец, $f(z)=f(a)\iff h_1(z)=0\iff z=a$ в $V$.
\end{proof}

\subsection{Локальная обратимость}

\begin{thmbox}{Утверждение (критерий локальной обратимости)}
Функция $f\in H(a)$ локально обратима в точке $a$ тогда и только тогда, когда $f'(a)\ne0$.
\end{thmbox}
\begin{proof}
Если $g(f(z))=z$ при $z\in U$, то $f$ инъективна в $U$ (из $f(z_1)=f(z_2)$ следует $z_1=g(f(z_1))=g(f(z_2))=z_2$),
т.е. $f\in\mathcal U(a)$, и по Предл.~13.1 $f'(a)\ne0$.

Обратно, пусть $f'(a)\ne0$. По доказательству Предл.~13.1 $f$ инъективна в $U=B(a,\rho)$ и $f'\neq0$ в $U$. По п.~2)
Предл.~13.1 (применённому к области $U$) $f(U)$ --- область, содержащая $f(a)$, и $f^{-1}\in H(f(U))$. Положим
$\widetilde U=f(U)$, $g=f^{-1}$; тогда $g(f(z))=z$ при $z\in U$.
\end{proof}

\begin{remarkbox}{Теорема об обратной функции в голоморфном варианте}
Итак, если $f'(a)\ne0$, то $f$ --- биголоморфное отображение некоторой окрестности $U$ точки $a$ на окрестность
$\widetilde U$ точки $f(a)$, и $(f^{-1})'(f(a))=1/f'(a)$. Кроме того, $f^{-1}$ даётся формулой
$f^{-1}(w)=\frac1{2\pi i}\int_{C(a,r)}\frac{zf'(z)}{f(z)-w}\,dz$ для $w$, близких к $f(a)$ (вычет функции
$zf'(z)/(f(z)-w)$ в единственном простом нуле $z_w$ функции $f-w$ внутри $C(a,r)$ равен $z_w$), откуда ещё раз видна
голоморфность $f^{-1}$.
\end{remarkbox}

\begin{remarkbox}{Примеры применения критерия}
\begin{enumerate}
\item $f(z)=z^n$, $n\ge2$: $f'(z)=nz^{n-1}$, поэтому $f\in\mathcal U_\ell(\CC\setminus\{0\})$ и $f\notin\mathcal U(0)$;
вблизи нуля $f$ ровно $n$-листна (Следствие выше).
\item $f(z)=\cos z$: $f'(z)=-\sin z=0\iff z=\pi k$; значит, $\cos z$ однолистна в каждой точке, кроме $\pi k$,
$k\in\ZZ$; вблизи $z=0$ имеем $\cos z=1-2\sin^2(z/2)=1+h_1(z)^2$, $h_1(z)=i\sqrt2\sin(z/2)$, --- двулистность.
\item Функция Жуковского $\Zh(z)=\frac12(z+1/z)$: $\Zh'(z)=\frac12(1-1/z^2)=0\iff z=\pm1$. Поэтому $\Zh$ не однолистна
ни в какой области, содержащей $\pm1$; она однолистна в $\DD\setminus\{0\}$ и в $\{|z|>1\}$, т.к.
$\Zh(z_1)=\Zh(z_2)\iff z_1=z_2$ или $z_1z_2=1$.
\item Однолистность в точке не означает однолистности в области: $e^z$ (см. выше). Даже если $f\in \mathcal U_\ell(D)$
и $f$ инъективна на $\partial D$, без дополнительных условий нельзя гарантировать $f\in\mathcal U(D)$; достаточные
условия глобальной однолистности --- предмет билетов 2 и 3.
\end{enumerate}
\end{remarkbox}

\begin{remarkbox}{Коротко: план ответа на экзамене}
\begin{enumerate}
\item Определения: $\mathcal U(a)$, $\mathcal U_\ell(D)$, $\mathcal U(D)$; конформность в точке и в области; примеры
$e^z$, $z^2$.
\item Предл.~13.1(1), $\Leftarrow$: выбираем $\rho$ с $|f'-f'(a)|<\frac12|f'(a)|$ в $B(a,\rho)$, пишем
$f(z_2)-f(z_1)=\int_{[z_1,z_2]}(f'-f'(a))+f'(a)(z_2-z_1)$, оцениваем.
\item $\Rightarrow$: при $f'(a)=0$ имеем $f-f(a)=(z-a)^ng$, $n\ge2$; извлекаем корень $g=g_1^n$ (Лемма 12.3), $h_1=(z-a)g_1$
однолистна, $f=f(a)+h_1^n$; по открытости $h_1$ берём $w_1\neq w_2$ с $w_1^n=w_2^n$.
\item Следствие: около точки кратности $n$ функция ровно $n$-листна.
\item Предл.~13.1(2): $\mathcal U(D)=\mathrm{Conf}(D)$; $f(D)$ --- область, $f^{-1}$ голоморфна, $(f^{-1})'=1/f'$.
\item Критерий локальной обратимости: $\Leftrightarrow f'(a)\ne0$.
\end{enumerate}
\end{remarkbox}

% =====================================================================
\ticket{2}{Обратный принцип соответствия границ и~теорема\texorpdfstring{\\}{ }Носиро--Варшавского (однолистность голоморфной в~выпуклой области функции с~условием \texorpdfstring{$\Re f'>0$}{Re f'>0})}

\begin{defs}
\textbf{Жорданова кривая.} Замкнутая жорданова кривая --- образ $\Gamma=\gamma([0,1])$ непрерывного пути
$\gamma:[0,1]\to\CC$ с $\gamma(0)=\gamma(1)$, инъективного на $[0,1)$ (т.е. гомеоморфный образ окружности).

\textbf{Теорема Жордана.} Если $\Gamma\subset\CC$ --- замкнутая жорданова кривая, то $\Ch\setminus\Gamma$ состоит
ровно из двух компонент связности: ограниченной $\mathrm{Int}(\Gamma)$ и неограниченной $\mathrm{Ext}(\Gamma)\ni\infty$;
при этом $\partial\,\mathrm{Int}(\Gamma)=\partial\,\mathrm{Ext}(\Gamma)=\Gamma$.
\emph{Дополнение (об индексе):} если $\gamma$ --- параметризация $\Gamma$, то $\Ind(\gamma,b)=0$ при
$b\in\mathrm{Ext}(\Gamma)$ и $\Ind(\gamma,b)=\pm1$ (один и тот же знак для всех $b\in\mathrm{Int}(\Gamma)$) при
$b\in\mathrm{Int}(\Gamma)$; знак $+$ соответствует положительной ориентации.

\textbf{Жорданова область} --- область $D=\mathrm{Int}(\Gamma)$; её граница $\partial D=\Gamma$, $\partial^+D$ --- граница,
ориентированная положительно (так, что $\Ind(\partial^+D,b)=1$ при $b\in D$).

\textbf{Индекс и приращение аргумента.} Для замкнутого пути $\sigma:[0,1]\to\CC\setminus\{b\}$:
$\Ind(\sigma,b)=\frac1{2\pi}\Delta_\sigma\Arg(w-b)$, где $\Delta_\sigma\Arg(w-b)=\alpha(1)-\alpha(0)$ для любой
непрерывной ветви $\alpha(t)$ аргумента $\sigma(t)-b$. Для функции $f$, непрерывной и не обращающейся в нуль на
$\partial D$, $\Delta_{\partial^+D}\Arg f$ --- приращение аргумента $f$ вдоль $\partial^+D$, т.е.
$\Delta_{f\circ\gamma}\Arg w$, где $\gamma$ --- параметризация $\partial^+D$.

\textbf{Число нулей.} $z(f,D)$ --- число нулей $f$ в $D$ с учётом кратности.

\textbf{Выпуклая область}: вместе с любыми двумя точками $a,b$ содержит отрезок $[a,b]=\{a+t(b-a):t\in[0,1]\}$.
\end{defs}

\subsection{Принцип аргумента (напоминание)}

\begin{thmbox}{Теорема 10.2 (усиленный принцип аргумента; частный случай)}
Пусть $D\subset\CC$ --- жорданова (более общо: допустимая) область, $f\in H(D)\cap C(\clos D)$ и $f\ne0$ на
$\partial D$. Тогда $f$ имеет в $D$ конечное число нулей и
\[
z(f,D)=\frac1{2\pi}\Delta_{\partial^+D}\Arg f .
\]
\end{thmbox}
Для гладкой (кусочно гладкой) границы и $f$, голоморфной в окрестности $\clos D$, это --- интегральная формула
$z(f,D)=\frac1{2\pi i}\int_{\partial^+D}\frac{f'}{f}dz$ (теорема о логарифмических вычетах + Лемма 10.2); общий
случай $f\in C(\clos D)$ в лекциях выводится из усиленной интегральной теоремы Коши (Теорема 7.4). Мы используем
Теорему 10.2 как известный факт первого семестра.

\subsection{Обратный принцип соответствия границ}

\begin{thmbox}{Предложение 10.3 (обратный принцип соответствия границ)}
Пусть $D\subset\CC$ --- жорданова область, $\Gamma=\partial D$, и пусть $f\in H(D)\cap C(\clos D)$ инъективна на $\Gamma$.
Тогда $\Upsilon=f(\Gamma)$ --- замкнутая жорданова кривая, $f\in\mathcal U(D)$ и $f$ отображает $D$ конформно на
жорданову область $G=\mathrm{Int}(\Upsilon)$: $f\in\mathrm{Conf}(D,G)$. Более того, $f$ --- гомеоморфизм $\clos D$
на $\clos G$.
\end{thmbox}

\begin{proof}
\textbf{Шаг 1: $\Upsilon$ --- жорданова кривая.} Пусть $\gamma:[0,1]\to\CC$ --- параметризация $\Gamma$ (непрерывна,
$\gamma(0)=\gamma(1)$, инъективна на $[0,1)$). Тогда $\sigma=f\circ\gamma$ непрерывна, $\sigma(0)=\sigma(1)$ и
$\sigma$ инъективна на $[0,1)$ как композиция инъективных отображений ($f$ инъективна на $\Gamma$). Значит,
$\Upsilon=\sigma([0,1])=f(\Gamma)$ --- замкнутая жорданова кривая, и по теореме Жордана определены
$G=\mathrm{Int}(\Upsilon)$ и $E=\mathrm{Ext}(\Upsilon)$, $\CC=G\sqcup\Upsilon\sqcup(E\setminus\{\infty\})$.

Отметим, что $f$ непостоянна в $D$: иначе по непрерывности $f$ была бы постоянна и на $\Gamma=\partial D$, что
противоречит инъективности $f$ на $\Gamma$ (на $\Gamma$ бесконечно много точек).

\textbf{Шаг 2: подсчёт прообразов точек $b\notin\Upsilon$.} Пусть $b\in\CC\setminus\Upsilon$. Функция $f-b$ лежит
в $H(D)\cap C(\clos D)$ и не обращается в нуль на $\Gamma$ (так как $f(\Gamma)=\Upsilon\not\ni b$). По Теореме 10.2
\[
z(f-b,D)=\frac1{2\pi}\Delta_{\partial^+D}\Arg(f-b)=\frac1{2\pi}\Delta_{\sigma}\Arg(w-b)=\Ind(\sigma,b),
\]
где $\sigma=f\circ\gamma$ --- параметризация $\Upsilon$ (здесь $\gamma$ ориентирован положительно относительно $D$).
По дополнению к теореме Жордана:
\begin{itemize}
\item если $b\in E$, то $\Ind(\sigma,b)=0$, т.е. $z(f-b,D)=0$: значение $b$ в $D$ не принимается;
\item если $b\in G$, то $\Ind(\sigma,b)=\eps$, где $\eps\in\{1,-1\}$ не зависит от $b\in G$. Поскольку
$z(f-b,D)\ge0$, то $\eps=1$. Значит, $z(f-b,D)=1$: значение $b$ принимается в $D$ ровно один раз, причём
с кратностью $1$.
\end{itemize}

\begin{detail}{значения из $\Upsilon$ не принимаются в $D$ (на слайде этот момент пропущен)}
Пусть $f(z_0)=b\in\Upsilon$ для некоторого $z_0\in D$. По принципу сохранения области $f(D)$ открыто, т.е.
$B(b,\eps)\subset f(D)$ при некотором $\eps>0$. Так как $b\in\Upsilon=\partial E$, круг $B(b,\eps)$ пересекает $E$, т.е.
некоторая точка $b'\in E$ принимается функцией $f$ в $D$, что противоречит доказанному ($z(f-b',D)=0$).
\end{detail}

\textbf{Шаг 3: вывод.} Из шага 2 и детали: $f(D)\subset G$, и каждое $b\in G$ имеет ровно один прообраз в $D$.
Значит, $f:D\to G$ --- биекция, $f\in\mathcal U(D)$, и по Предл.~13.1 $f\in\mathrm{Conf}(D,G)$ (в частности,
$f'\ne0$ в $D$, что видно и из того, что кратность каждого прообраза равна $1$).

\textbf{Шаг 4: гомеоморфизм замыканий.} Отображение $f:\clos D\to\CC$ непрерывно и инъективно: на $D$ и на $\Gamma$
по отдельности инъективно, а $f(D)=G$ и $f(\Gamma)=\Upsilon$ не пересекаются. Его образ равен $G\cup\Upsilon=\clos G$.
Непрерывная биекция компакта $\clos D$ на хаусдорфово пространство $\clos G$ является гомеоморфизмом.
\end{proof}

\begin{remarkbox}{Замечание}
Важно, что условие инъективности наложено \emph{только на границе}: однолистность внутри --- следствие.
Это --- первое из достаточных условий однолистности.
\end{remarkbox}

\begin{remarkbox}{Пример применения}
Пусть $f(z)=z+\frac{z^2}{2}$ в $\DD$. На $\TT$: если $f(z_1)=f(z_2)$, $z_1\ne z_2$, то
$(z_1-z_2)\bigl(1+\frac{z_1+z_2}2\bigr)=0$, т.е. $z_1+z_2=-2$, что для $|z_1|=|z_2|=1$ возможно лишь при
$z_1=z_2=-1$. Значит, $f$ инъективна на $\TT$, и по Предл.~10.3 $f$ однолистна в $\DD$ и отображает $\DD$ на
внутренность кривой $\{e^{it}+\frac12e^{2it}\}$ (кардиоидоподобной кривой с точкой возврата $f(-1)=-\frac12$; заметим,
$f'(-1)=0$, но $-1\notin\DD$). Это же следует и из Предл.~13.2: $\Re f'(z)=\Re(1+z)>0$ в $\DD$.
\end{remarkbox}

\subsection{Теорема Носиро--Варшавского}

\begin{thmbox}{Предложение 13.2 (условие однолистности Носиро--Варшавского)}
Пусть $D\subset\CC$ --- выпуклая область, а $f\in H(D)$ такова, что $\Re f'(z)>0$ при всех $z\in D$.
Тогда $f\in\mathcal U(D)$.
\end{thmbox}

\begin{proof}
Пусть $a,b\in D$, $a\ne b$. Так как $D$ выпукла, $[a,b]\subset D$. Функция $f$ --- первообразная $f'$ в $D$, и
путь $t\mapsto a+(b-a)t$, $t\in[0,1]$, гладкий, поэтому
\[
f(b)-f(a)=\int_{[a,b]}f'(\zeta)\,d\zeta=(b-a)\int_0^1f'\bigl(a+(b-a)t\bigr)\,dt .
\]
Разделим на $b-a\neq 0$ и возьмём вещественную часть (она перестановочна с интегралом от непрерывной функции):
\[
\Re\frac{f(b)-f(a)}{b-a}=\int_0^1\Re f'\bigl(a+(b-a)t\bigr)\,dt>0,
\]
так как подынтегральная функция непрерывна и строго положительна на $[0,1]$. Следовательно, $f(b)\ne f(a)$.
\end{proof}

\begin{remarkbox}{Замечания к Предл.~13.2}
\begin{enumerate}
\item Условие можно ослабить до $\Re\bigl(e^{i\theta}f'(z)\bigr)>0$ в $D$ при некотором $\theta\in\RR$ (применить
предложение к $e^{i\theta}f$). Можно также требовать лишь $\Re f'\ge0$ в $D$ и $f'\not\equiv0$: если $\Re f'$
обращается в нуль в точке $D$, то гармоническая функция $\Re f'\ge0$ достигает там минимума и по принципу
минимума $\Re f'\equiv0$, т.е. $f'\equiv ic$, $c\in\RR\setminus\{0\}$, и $f(z)=icz+d$ однолистна; иначе $\Re f'>0$.
\item Геометрически: $\Re\frac{f(b)-f(a)}{b-a}>0$ означает, что угол между векторами $b-a$ и $f(b)-f(a)$ меньше
$\pi/2$: $f$ «не поворачивает хорды более чем на прямой угол».
\item \textbf{Выпуклость существенна.} Пусть $f=2\Zh$, т.е. $f(z)=z+\frac1z$, $f'(z)=1-\frac1{z^2}$, так что
$\Re f'(z)>0\iff\Re\frac1{z^2}<1$; при $z=re^{i\varphi}$ это условие $\cos2\varphi<r^2$. Рассмотрим путь $L$: дуга
окружности $|z|=2$ от $2i$ через $-2$ до $-2i$ (там $\cos2\varphi\le1<4$), затем отрезок мнимой оси от $-2i$ до
$-\frac i2$ (там $z=-is$, $\frac1{z^2}=-\frac1{s^2}<1$). Компакт $L$ лежит в открытом множестве
$\{z\ne0:\Re(1/z^2)<1\}$, поэтому при малом $\delta>0$ его $\delta$-окрестность $D=\bigcup_{z\in L}B(z,\delta)$ ---
(невыпуклая) область, в которой $\Re f'>0$. Но $f(2i)=2i-\frac i2=\frac{3i}2=-\frac i2+2i=f(-\frac i2)$, так что
$f\notin\mathcal U(D)$.
\end{enumerate}
\end{remarkbox}

\subsection{Жордановость образа при непрерывности вплоть до границы}

\begin{thmbox}{Замечание к Предл.~13.2 (лекции)}
Если $f\in H(\DD)\cap C(\clos\DD)$ и $\Re f'(z)>0$ при всех $z\in\DD$, то $f$ инъективна на $\clos\DD$, $f(\TT)$ ---
замкнутая жорданова кривая и $f(\DD)=\mathrm{Int}\,f(\TT)$ --- жорданова область.
\end{thmbox}
\begin{proof}
На $\DD$ инъективность --- Предл.~13.2. Пусть $a,b\in\TT$, $a\ne b$. Тогда интервал $(a,b)=\{a+t(b-a):0<t<1\}$
лежит в $\DD$ (круг строго выпукл: $|a+t(b-a)|<1$ при $0<t<1$, поскольку $|(1-t)a+tb|\le1$ с равенством только при
$a=b$). Сам отрезок $[a,b]$ в $\DD$ не лежит, и интеграл от $f'$ по нему, вообще говоря, несобственный; поэтому
аккуратно: при $0<s<\frac12$ точки $a_s=a+s(b-a)$, $b_s=a+(1-s)(b-a)$ лежат в $\DD$, и как в доказательстве
Предл.~13.2
\[
\Re\frac{f(b_s)-f(a_s)}{b-a}=\int_s^{1-s}\Re f'\bigl(a+(b-a)t\bigr)\,dt\ \ge\ \int_{1/4}^{3/4}\Re f'\bigl(a+(b-a)t\bigr)\,dt=:c>0
\]
при $s<\frac14$ (подынтегральная функция положительна, а $[\frac14,\frac34]\subset[s,1-s]$). Устремляя $s\to0+$ и
пользуясь непрерывностью $f$ в $\clos\DD$, получаем $\Re\frac{f(b)-f(a)}{b-a}\ge c>0$, т.е. $f(a)\ne f(b)$.
Итак, $f$ инъективна на $\TT$, и по Предл.~10.3 $f(\TT)$ --- жорданова кривая, а $f$ отображает $\DD$ конформно на
$\mathrm{Int}\,f(\TT)$. Наконец, $f(\DD)\cap f(\TT)=\varnothing$ (Предл.~10.3), так что $f$ инъективна и на $\clos\DD$.
\end{proof}

\begin{remarkbox}{Замечание}
То же рассуждение проходит для любой ограниченной строго выпуклой области $D$ (интервал между двумя граничными
точками лежит в $D$). Для произвольной выпуклой области отрезок между граничными точками может лежать на
$\partial D$ (например, сторона квадрата), и там рассуждение не работает.
\end{remarkbox}

\begin{remarkbox}{Коротко: план ответа на экзамене}
\begin{enumerate}
\item Определения: жорданова кривая и область, теорема Жордана и индекс жордановой кривой ($\pm1$ внутри, $0$ снаружи).
\item Принцип аргумента для $f\in H(D)\cap C(\clos D)$: $z(f,D)=\frac1{2\pi}\Delta_{\partial^+ D}\Arg f$.
\item Предл.~10.3: $\Upsilon=f(\Gamma)$ жорданова; для $b\notin\Upsilon$ $z(f-b,D)=\Ind(f\circ\gamma,b)$, равное $0$ вне
и $1$ внутри (знак фиксируется неотрицательностью числа нулей).
\item Значения на $\Upsilon$ не принимаются (открытость $f(D)$ + $\Upsilon=\partial\,\mathrm{Ext}$). Итог: биекция $D\to G$.
\item Предл.~13.2: $f(b)-f(a)=(b-a)\int_0^1 f'(a+(b-a)t)dt$, $\Re\frac{f(b)-f(a)}{b-a}>0$.
\item Выпуклость существенна (пример с $z+1/z$); при $f\in C(\clos\DD)$ образ круга --- жорданова область.
\end{enumerate}
\end{remarkbox}

% =====================================================================
\ticket{3}{Достаточное условие однолистности голоморфной непостоянной в~области функции с~условием \texorpdfstring{$f(z)\to\Gamma$}{f(z)->Gamma} при \texorpdfstring{$z\to\partial D$}{z->dD}}

\begin{defs}
\textbf{Расширенная плоскость.} $\Ch=\CC\cup\{\infty\}$ (сфера Римана) --- компактное метрическое пространство;
окрестности $\infty$: $\{|z|>R\}\cup\{\infty\}$. Для $\Omega\subset\Ch$ замыкание $\clos\Omega$ и граница
$\partial\Omega$ берутся в $\Ch$. Функция $f$ \emph{голоморфна в $\infty$}, если $f(\infty)\in\CC$ и $w\mapsto f(1/w)$
голоморфна в окрестности $w=0$ (с $f(1/0):=f(\infty)$). $H(\Omega)$ для $\Omega\subset\Ch$ --- функции, голоморфные
в каждой точке $\Omega$ (в частности, конечнозначные).

\textbf{Условие «$f(z)\to\Gamma$ при $z\to\partial\Omega$».} Если $z_n\in\Omega$ и все предельные точки
последовательности $\{z_n\}$ (в $\Ch$) лежат на $\partial\Omega$, то все предельные точки последовательности
$\{f(z_n)\}$ (в $\Ch$) лежат на $\Gamma$. Для ограниченной кривой $\Gamma$ это равносильно:
$\dist(f(z),\Gamma)\to0$, когда $z$ стремится к $\partial\Omega$ (в сферической метрике).

\textbf{Жорданова кривая, теорема Жордана} --- см. билет 2: $\Ch\setminus\Gamma=D\sqcup D_\infty$, где
$D=D(\Gamma)=\mathrm{Int}(\Gamma)$ --- ограниченная компонента, $\infty\in D_\infty$, и
$\partial D=\partial D_\infty=\Gamma$.

\textbf{Кратность значения.} Значение $w$ принимается функцией $f$ в точке $z_0\in\Omega\cap\CC$ с кратностью $k$,
если $z_0$ --- нуль порядка $k$ функции $f-w$. «$f$ принимает значение $w$ ровно один раз с учётом кратности»
означает: существует единственная точка $z_w\in\Omega$ с $f(z_w)=w$, и в ней $f'(z_w)\ne0$.

\textbf{Односвязность.} Область $\Omega\subset\Ch$ односвязна, если любой замкнутый путь в $\Omega$ гомотопен в $\Omega$
постоянному. Для областей в $\Ch$ это равносильно связности $\Ch\setminus\Omega$ (факт первого семестра).
Односвязность сохраняется при гомеоморфизмах.
\end{defs}

\subsection{Обратный принцип соответствия границ}

Этот принцип --- исторически первое достаточное условие однолистности и прообраз Предл.~13.3; по программе его
нужно уметь формулировать и доказывать и в этом билете.

\begin{thmbox}{Предложение 10.3 (обратный принцип соответствия границ)}
Пусть $D\subset\CC$ --- жорданова область, $\Gamma=\partial D$, и пусть $f\in H(D)\cap C(\clos D)$ взаимно однозначна
на $\Gamma$. Тогда $f\in\mathcal U(D)$ и $f$ отображает $D$ конформно на область $\mathrm{Int}(f(\Gamma))$.
\end{thmbox}
\begin{proof}
Пусть $\gamma:[0,1]\to\Gamma$ --- положительно ориентированная параметризация $\Gamma$. Тогда $\sigma=f\circ\gamma$
непрерывна, замкнута и инъективна на $[0,1)$, т.е. $\Upsilon=f(\Gamma)$ --- замкнутая жорданова кривая; пусть
$G=\mathrm{Int}(\Upsilon)$, $E=\mathrm{Ext}(\Upsilon)$. Для $b\notin\Upsilon$ функция $f-b\in H(D)\cap C(\clos D)$ не
обращается в нуль на $\Gamma$, и по усиленному принципу аргумента (Теорема 10.2)
\[
z(f-b,D)=\tfrac1{2\pi}\Delta_{\partial^+D}\Arg(f-b)=\tfrac1{2\pi}\Delta_\sigma\Arg(w-b)=\Ind(\sigma,b).
\]
По дополнению к теореме Жордана $\Ind(\sigma,b)=0$ при $b\in E$ и $\Ind(\sigma,b)=\eps\in\{\pm1\}$ (одно и то же для всех
$b\in G$) при $b\in G$; так как число нулей неотрицательно, $\eps=1$. Итак, точки $E$ не принимаются в $D$, а каждая
точка $G$ принимается ровно один раз и с кратностью $1$. Точки $b\in\Upsilon$ тоже не принимаются: $f$ непостоянна
(она инъективна на $\Gamma$), поэтому $f(D)$ открыто, и если бы $b=f(z_0)\in\Upsilon$, $z_0\in D$, то $f(D)$ содержало бы
круг $B(b,\eps)$, а он пересекает $E$ (так как $b\in\partial E$) --- противоречие. Значит, $f$ --- биекция $D$ на $G$, и
по Предл.~13.1 $f\in\mathrm{Conf}(D,G)$. (Подробнее --- билет 2.)
\end{proof}

\begin{remarkbox}{Связь с Предл.~13.3}
В условиях Предл.~10.3 выполнено $f(z)\to\Upsilon$ при $z\to\partial D$: если предельные точки $\{z_n\}\subset D$ лежат
на $\Gamma$, то для любой подпоследовательности $z_{n_k}\to\zeta\in\Gamma$ по непрерывности $f$ в $\clos D$ имеем
$f(z_{n_k})\to f(\zeta)\in\Upsilon$. Поэтому первое утверждение Предл.~10.3 ($f(D)=\mathrm{Int}\,\Upsilon$) --- частный
случай Предл.~13.3. Предл.~13.3 гораздо гибче: не требуется ни непрерывности $f$ вплоть до границы, ни жордановости
(и даже односвязности) области $\Omega$; односвязность $\Omega$ получается как \emph{следствие}.
\end{remarkbox}

\subsection{Формулировка и вспомогательные факты}

\begin{thmbox}{Предложение 13.3 (условие однолистности)}
Пусть $\Omega\subset\Ch$ --- область, $\Omega\ne\Ch$, $\Gamma$ --- замкнутая жорданова кривая в $\CC$ и $D=D(\Gamma)$.
Пусть функция $f\in H(\Omega)$ \emph{непостоянна} и $f(z)\to\Gamma$ при $z\to\partial\Omega$. Тогда $f(\Omega)=D$.
Более того, если $f$ принимает в $\Omega$ какое-то значение из $D$ ровно один раз (с учётом кратности), то $f$
инъективна (т.е. $f\in\mathrm{Conf}(\Omega,D)$), а $\Omega$ односвязна.
\end{thmbox}

\begin{remarkbox}{Замечание (уточнение формулировки слайда)}
На слайде не оговорены условия «$f$ непостоянна» и «$\Omega\ne\Ch$». Они необходимы: постоянная $f\equiv c\in\Gamma$
удовлетворяет условию $f(z)\to\Gamma$, но $f(\Omega)=\{c\}\ne D$; при $\Omega=\Ch$ граница пуста, условие выполнено
бессодержательно, а $f\in H(\Ch)$ постоянна по теореме Лиувилля. В программе (билет 3) условие непостоянства
указано явно. Кроме того, на слайде 279 в фразе «так как $\partial\Omega\subset\Gamma$» имеется в виду
$\partial f(\Omega)\subset\Gamma$.
\end{remarkbox}

\begin{thmbox}{Лемма А (топологическая)}
Пусть $X$ --- топологическое пространство, $A\subset X$, $E\subset X$ связно, $E\cap A\ne\varnothing$ и $E\setminus A\ne\varnothing$.
Тогда $E\cap\partial A\ne\varnothing$.
\end{thmbox}
\begin{proof}
Пусть $E\cap\partial A=\varnothing$. Пространство $X$ есть дизъюнктное объединение $\mathrm{int}A$, $\partial A$ и
$\mathrm{ext}A=X\setminus\clos A$. Тогда $E=(E\cap\mathrm{int}A)\sqcup(E\cap\mathrm{ext}A)$ --- объединение двух
относительно открытых в $E$ множеств. Точка из $E\cap A$ лежит в $\mathrm{int}A$ (в $\partial A$ она не лежит), а точка
из $E\setminus A$ лежит в $\mathrm{ext}A$; значит, оба множества непусты, что противоречит связности $E$.
\end{proof}

\begin{detail}{сведение к случаю $\Omega\subset\CC$}
Если $\infty\in\Omega$, выберем $p\in\Ch\setminus\Omega$ (существует, так как $\Omega\ne\Ch$); $p\ne\infty$. ДЛО
$T(z)=1/(z-p)$ --- гомеоморфизм $\Ch$ на себя, $T(p)=\infty$, поэтому $\Omega'=T(\Omega)\subset\CC$ --- область,
$T(\partial\Omega)=\partial\Omega'$, и $T$ сохраняет сходимость последовательностей и их предельные точки. Функция
$f_1=f\circ T^{-1}$ голоморфна в $\Omega'$ (в точке $T(\infty)=0$ --- по определению голоморфности $f$ в $\infty$, так как
$T^{-1}(w)=p+1/w$), непостоянна, $f_1(\Omega')=f(\Omega)$, $f_1(z)\to\Gamma$ при $z\to\partial\Omega'$; число прообразов
каждого значения с учётом кратности у $f$ и $f_1$ одинаково (кратность в $\infty$ определяется через локальную
координату, а $T$ --- локально конформна). Наконец, инъективность $f$ и односвязность $\Omega$ равносильны
инъективности $f_1$ и односвязности $\Omega'$. Поэтому \textbf{далее считаем $\Omega\subset\CC$.}
\end{detail}

\subsection{Доказательство: $f(\Omega)=D$}

\begin{proof}[Доказательство первого утверждения]
По принципу сохранения области $f(\Omega)$ --- область в $\CC$; в частности, $f(\Omega)$ открыто, и
$\partial f(\Omega)\cap f(\Omega)=\varnothing$.

\textbf{Шаг 1: $\partial f(\Omega)\subset\Gamma$} (граница в $\Ch$). Пусть $w\in\partial f(\Omega)$ (случай $w=\infty$ не
исключается). Найдутся $z_n\in\Omega$ с $f(z_n)\to w$. По компактности $\Ch$, переходя к подпоследовательности,
считаем $z_n\to\zeta\in\clos\Omega$. Если $\zeta\in\Omega$, то $f(z_n)\to f(\zeta)$ по непрерывности, т.е.
$w=f(\zeta)\in f(\Omega)$, что невозможно, так как $w\in\partial f(\Omega)$ и $f(\Omega)$ открыто. Значит, $\zeta\in\partial\Omega$;
у последовательности $\{z_n\}$ единственная предельная точка $\zeta\in\partial\Omega$, и по условию все предельные точки
$\{f(z_n)\}$, т.е. $w$, лежат на $\Gamma$. (В частности, $\infty\notin\partial f(\Omega)$, и $f$ ограничена.)

\textbf{Шаг 2: $f(\Omega)\subset D$.} Пусть $f(\Omega)\not\subset D$. Тогда $f(\Omega)$ пересекает $\Gamma$ или $D_\infty$. Если
$f(\Omega)\ni b\in\Gamma$, то, так как $f(\Omega)$ открыто, а $b\in\Gamma=\partial D_\infty$, множество $f(\Omega)$ пересекает
$D_\infty$. Итак, в любом случае $D_\infty\cap f(\Omega)\ne\varnothing$. С другой стороны, $\infty\in D_\infty\setminus f(\Omega)$
($f$ конечнозначна). Множество $D_\infty$ связно, и по Лемме А $D_\infty\cap\partial f(\Omega)\ne\varnothing$. Но
$\partial f(\Omega)\subset\Gamma$ и $\Gamma\cap D_\infty=\varnothing$ --- противоречие.

\textbf{Шаг 3: $f(\Omega)=D$.} Если $D\setminus f(\Omega)\ne\varnothing$, то связное множество $D$ пересекает $f(\Omega)$
(по шагу 2 $\varnothing\neq f(\Omega)\subset D$) и его дополнение, и по Лемме А $D\cap\partial f(\Omega)\neq\varnothing$, что
противоречит $\partial f(\Omega)\subset\Gamma$, $\Gamma\cap D=\varnothing$.
\end{proof}

\subsection{Доказательство: однолистность и односвязность}

\begin{proof}[Доказательство второго утверждения]
Положим
\[
\widetilde\Omega=\bigl\{z\in\Omega:\ f'(z)\ne0\ \text{и}\ f(z)\ne f(z')\ \text{для всех}\ z'\in\Omega,\ z'\ne z\bigr\}
\]
--- множество точек, в которых $f$ «однократна». По условию существует $a\in D$, принимаемое ровно один раз с учётом
кратности: $f(z_a)=a$ для единственной $z_a\in\Omega$, и $f'(z_a)\ne0$. Значит, $z_a\in\widetilde\Omega\ne\varnothing$.
Докажем, что $\widetilde\Omega=\Omega$; тогда $f$ инъективна и $f'\ne0$ в $\Omega$.

Пусть $\Omega\setminus\widetilde\Omega\neq\varnothing$. Так как $\Omega$ связна, по Лемме А (с $X=\CC$, $E=\Omega$,
$A=\widetilde\Omega$) найдётся $z_0\in\Omega\cap\partial\widetilde\Omega$. Каждая окрестность $z_0$ содержит точки
$\widetilde\Omega$ (так как $z_0\in\partial\widetilde\Omega$); мы получим противоречие, построив окрестность $z_0$ без точек
$\widetilde\Omega$. Поскольку $z_0\in\partial\widetilde\Omega$ и $\Omega$ открыто, существуют $z_n\in\Omega\setminus\widetilde\Omega$,
$z_n\to z_0$ (либо $z_0\notin\widetilde\Omega$ и можно взять $z_n=z_0$, либо $z_0\in\widetilde\Omega$, и тогда любая окрестность
$z_0$ содержит точки $\CC\setminus\widetilde\Omega$, которые при малой окрестности лежат в $\Omega$).
Точка $z_n\notin\widetilde\Omega$ означает: либо (i) $f'(z_n)=0$, либо (ii) существует $w_n\in\Omega$, $w_n\ne z_n$,
$f(w_n)=f(z_n)$. Возможны случаи.

\textbf{Случай I: (i) выполнено для бесконечного числа $n$.} Тогда по непрерывности $f'$ получаем $f'(z_0)=0$ ---
переходим к шагу «$f'(z_0)=0$» ниже.

\textbf{Случай II: (ii) выполнено для всех $n$, начиная с некоторого.} По компактности $\Ch$ можно считать $w_n\to w_0\in\clos\Omega$.
Так как $f(w_n)=f(z_n)\to f(z_0)\in f(\Omega)=D$, точка $w_0$ не может лежать на $\partial\Omega$: иначе по условию
$f(z)\to\Gamma$ (применённому к последовательности $\{w_n\}$ с единственной предельной точкой $w_0\in\partial\Omega$)
предел $f(z_0)$ лежал бы на $\Gamma$, а $\Gamma\cap D=\varnothing$. Итак, $w_0\in\Omega$ и $f(w_0)=f(z_0)$.

\emph{Подслучай II.1: $w_0\ne z_0$.} Возьмём непересекающиеся круги $U_0=B(z_0,r)$, $V_0=B(w_0,r)$, лежащие в $\Omega$.
Функция $f$ непостоянна на каждом из них (иначе по теореме единственности $f$ постоянна в $\Omega$), поэтому по
принципу сохранения области $f(U_0)$ и $f(V_0)$ --- открытые множества, содержащие точку $c=f(z_0)=f(w_0)$. Значит,
найдётся $\eps>0$ с $W=B(c,\eps)\subset f(U_0)\cap f(V_0)$. Окрестность $N=U_0\cap f^{-1}(W)$ точки $z_0$ открыта
(непрерывность $f$). Любая точка $\zeta\in N$ не лежит в $\widetilde\Omega$: значение $f(\zeta)\in W\subset f(V_0)$
принимается и в некоторой точке $\zeta'\in V_0$, а $\zeta'\ne\zeta$, так как $U_0\cap V_0=\varnothing$. Итак,
$N\cap\widetilde\Omega=\varnothing$ --- противоречие с $z_0\in\partial\widetilde\Omega$.

\emph{Подслучай II.2: $w_0=z_0$.} Покажем, что $f'(z_0)=0$. Пусть $\eps>0$. Возьмём $\rho>0$ с $B(z_0,\rho)\subset\Omega$
и $|f'(\zeta)-f'(z_0)|<\eps$ при $\zeta\in B(z_0,\rho)$ (непрерывность $f'$). Так как $z_N\to z_0$ и $w_N\to z_0$,
найдётся $N$ с $z_N,w_N\in B(z_0,\rho)$; тогда $[z_N,w_N]\subset B(z_0,\rho)$ и
\[
0=f(w_N)-f(z_N)=\int_{[z_N,w_N]}\bigl(f'(\zeta)-f'(z_0)\bigr)\,d\zeta+f'(z_0)(w_N-z_N).
\]
Отсюда $|f'(z_0)|\,|w_N-z_N|\le\eps|w_N-z_N|$, и, поскольку $w_N\ne z_N$, $|f'(z_0)|\le\eps$. В силу произвольности
$\eps$ имеем $f'(z_0)=0$.

\textbf{Шаг «$f'(z_0)=0$» (завершает случаи I и II.2).} Так как $f$ непостоянна, $f(z)-f(z_0)=(z-z_0)^kg(z)$ в некотором
круге $B(z_0,\rho)\subset\Omega$, где $g\in H(B(z_0,\rho))$, $g(z_0)\neq0$, и $k\ge2$ (так как $f'(z_0)=0$). Выберем
$\delta\in(0,\rho)$ так, что $g\ne0$ на $\clos B(z_0,\delta)$, и положим
\[
\eps=\delta^k\min_{|z-z_0|=\delta}|g(z)|>0 .
\]
Пусть $w\in B(f(z_0),\eps)$. На окружности $|z-z_0|=\delta$ имеем $|f(z)-f(z_0)|=\delta^k|g(z)|\ge\eps>|f(z_0)-w|$.
По теореме Руше функция $f(z)-w=\bigl(f(z)-f(z_0)\bigr)+\bigl(f(z_0)-w\bigr)$ имеет в $B(z_0,\delta)$ столько же нулей
(с учётом кратности), сколько $f(z)-f(z_0)=(z-z_0)^kg(z)$, т.е. $k\ge2$. Теперь пусть
$\zeta\in N:=B(z_0,\delta)\cap f^{-1}(B(f(z_0),\eps))$ --- это окрестность $z_0$. Для $w=f(\zeta)$ функция $f-w$
имеет в $B(z_0,\delta)$ не менее двух нулей с учётом кратности, один из которых $\zeta$. Значит, либо $\zeta$ --- кратный
нуль, т.е. $f'(\zeta)=0$, либо есть ещё нуль $\zeta'\ne\zeta$, т.е. $f(\zeta')=f(\zeta)$. В обоих случаях
$\zeta\notin\widetilde\Omega$. Итак, $N\cap\widetilde\Omega=\varnothing$ --- противоречие с $z_0\in\partial\widetilde\Omega$.

Все случаи приводят к противоречию, значит, $\widetilde\Omega=\Omega$: $f$ инъективна в $\Omega$ и $f'\ne0$, т.е. (по
Предл.~13.1) $f\in\mathrm{Conf}(\Omega,D)$, и $f:\Omega\to D$ --- гомеоморфизм.

\textbf{Односвязность $\Omega$.} Область $D=\mathrm{Int}(\Gamma)$ односвязна: $\Ch\setminus D=\Gamma\cup D_\infty=\clos{D_\infty}$
(так как $\partial D_\infty=\Gamma$) --- замыкание связного множества, а значит, связно. Пусть $\gamma$ --- замкнутый путь
в $\Omega$; $f\circ\gamma$ --- замкнутый путь в $D$, гомотопный в $D$ постоянному посредством гомотопии $H(s,t)$; тогда
$f^{-1}\circ H$ --- гомотопия $\gamma$ к постоянному пути в $\Omega$ (непрерывна, так как $f^{-1}$ непрерывна).
Следовательно, $\Omega$ односвязна.
\end{proof}

\begin{remarkbox}{Другой взгляд: «степень» отображения}
Доказательство фактически показывает, что $f:\Omega\to D$ --- \emph{собственное} отображение (прообраз компакта
$K\subset D$ компактен в $\Omega$: иначе нашлась бы последовательность $z_n\in f^{-1}(K)$ с предельной точкой на
$\partial\Omega$, и $f(z_n)$ имели бы предельные точки на $\Gamma$, а не в $K$). Для собственного непостоянного
голоморфного отображения число прообразов $n(w)$ точки $w\in D$ с учётом кратности конечно и локально постоянно
(по теореме Руше), следовательно, постоянно в связной $D$. Условие «некоторое значение принимается ровно один раз»
означает $n\equiv1$, т.е. инъективность. В доказательстве выше этот факт получен «вручную» через множество
$\widetilde\Omega$.
\end{remarkbox}

\subsection{Примеры}

\begin{remarkbox}{Примеры}
\begin{enumerate}
\item \textbf{Преобразование Кэли.} $\Omega=\{\Im z>0\}$, $f(z)=\frac{z-i}{z+i}$, $\Gamma=\TT$, $D=\DD$. Функция $f$ ---
ДЛФ, непрерывная на $\clos\Omega\subset\Ch$ (полюс $-i\notin\clos\Omega$), причём $|f(x)|=\frac{|x-i|}{|x+i|}=1$ при $x\in\RR$
и $f(\infty)=1$, т.е. $|f|=1$ на $\partial\Omega=\RR\cup\{\infty\}$. Поэтому $f(z)\to\TT$ при $z\to\partial\Omega$. Значение $0$
принимается только в точке $z=i$, и $f'(i)=\frac{2i}{(2i)^2}=\frac1{2i}\ne0$. По Предл.~13.3 $f$ конформно отображает
верхнюю полуплоскость на $\DD$. Здесь $\infty\in\partial\Omega$ --- случай, не охватываемый Предл.~10.3.
\item \textbf{Конечные произведения Бляшке.} $B(z)=e^{i\theta}\prod_{j=1}^n\frac{z-a_j}{1-\clos{a_j}z}$, $a_j\in\DD$.
$B\in C(\clos\DD)$, $|B|=1$ на $\TT$, поэтому $B(z)\to\TT$ при $|z|\to1$, и $B(\DD)=\DD$. Значение $0$ принимается
$n$ раз с учётом кратности; при $n=1$ Предл.~13.3 даёт однолистность ($B\in\CAut(\DD)$).
\item \textbf{Условие «ровно один раз» существенно.} $f(z)=z^2$ в $\DD$: $f(z)\to\TT$, $f(\DD)=\DD$, но каждое значение
$w\ne0$ принимается дважды, а $0$ --- один раз, но с кратностью $2$; $f$ не однолистна.
\item \textbf{Неодносвязная $\Omega$.} $\Omega=\{\frac12<|z|<2\}$, $f=\Zh$, $\Zh(z)=\frac12(z+\frac1z)$. Обе граничные
окружности $|z|=2$ и $|z|=\frac12$ переходят в эллипс $\Gamma=\{\frac{u^2}{(5/4)^2}+\frac{v^2}{(3/4)^2}=1\}$ (при
$z=2e^{it}$: $\Zh=\frac54\cos t+\frac{3i}4\sin t$; при $z=\frac12e^{it}$: $\Zh=\frac54\cos t-\frac{3i}4\sin t$), и по
непрерывности $f$ в $\clos\Omega$ выполнено $f(z)\to\Gamma$. По Предл.~13.3 $f(\Omega)=\mathrm{Int}\,\Gamma$. Однолистности
нет: $\Zh(z)=\Zh(1/z)$, и $z,1/z\in\Omega$ одновременно; каждое значение принимается дважды (с учётом кратности) ---
что согласуется с неодносвязностью кольца.
\end{enumerate}
\end{remarkbox}

\begin{remarkbox}{Коротко: план ответа на экзамене}
\begin{enumerate}
\item Обратный принцип соответствия границ: формулировка и доказательство через принцип аргумента и индекс жордановой кривой.
\item Формулировка Предл.~13.3; смысл условия $f(z)\to\Gamma$; необходимость непостоянства $f$.
\item $\partial f(\Omega)\subset\Gamma$ (компактность $\Ch$ + открытость $f(\Omega)$).
\item $f(\Omega)\subset D$: иначе $D_\infty$ пересекает $f(\Omega)$ и его дополнение ($\infty$), значит и $\partial f(\Omega)$.
Затем $f(\Omega)=D$ аналогично.
\item $\widetilde\Omega$ --- точки однократности; $\widetilde\Omega\ne\varnothing$; граничная точка $z_0\in\Omega\cap\partial\widetilde\Omega$.
\item Пары $z_n\neq w_n$ с $f(z_n)=f(w_n)$; $w_0\in\Omega$ (иначе $f(z_0)\in\Gamma$); если $w_0\ne z_0$ --- открытость;
если $w_0=z_0$ --- $f'(z_0)=0$, затем Руше: $f$ не менее чем двулистна около $z_0$.
\item Итог: $f$ --- гомеоморфизм $\Omega\to D$, значит $\Omega$ односвязна. Примеры: Кэли, Бляшке, $z^2$, $\Zh$ в кольце.
\end{enumerate}
\end{remarkbox}

% =====================================================================
\ticket{4}{Теорема Ландау о круге однолистности}

\begin{defs}
\textbf{Элементарный множитель Бляшке.} Для $w\in\DD$
\[
\varphi_w(z)=\frac{z-w}{1-\clos wz}.
\]
Это ДЛФ с полюсом $1/\clos w\notin\clos\DD$ (при $w=0$: $\varphi_0(z)=z$), $\varphi_w(w)=0$, $\varphi_w(0)=-w$,
$\varphi_w'(z)=\frac{1-|w|^2}{(1-\clos wz)^2}$. Ключевое тождество (проверяется раскрытием скобок):
\begin{equation}\label{t4:id}
1-|\varphi_w(z)|^2=\frac{(1-|w|^2)(1-|z|^2)}{|1-\clos wz|^2}.
\end{equation}
\begin{detail}{проверка \eqref{t4:id}}
$|1-\clos wz|^2-|z-w|^2=(1-2\Re(\clos wz)+|w|^2|z|^2)-(|z|^2-2\Re(\clos wz)+|w|^2)=1-|z|^2-|w|^2+|w|^2|z|^2
=(1-|w|^2)(1-|z|^2)$; делим на $|1-\clos wz|^2$.
\end{detail}
Следовательно, $|\varphi_w|<1$ в $\DD$, $|\varphi_w|=1$ на $\TT$, $|\varphi_w|>1$ при $|z|>1$; $\varphi_w$ --- биекция $\DD$ на
$\DD$ с обратной $\varphi_w^{-1}=\varphi_{-w}$, т.е. $\varphi_w\in\CAut(\DD)$.

\textbf{$\CAut(\DD)$} --- группа конформных автоморфизмов круга: $\CAut(\DD)=\{e^{i\theta}\varphi_w:\ \theta\in\RR,\ w\in\DD\}$.

\textbf{Неевклидовы круги.} Для $a\in\DD$ и $0<r<1$
\[
K(a,r)=\{z\in\DD:\ |\varphi_a(z)|<r\}=\varphi_a^{-1}\bigl(B(0,r)\bigr)
\]
--- прообраз круга $B(0,r)$ при автоморфизме $\varphi_a$ (круг в псевдогиперболической метрике
$\delta(z,a)=|\varphi_a(z)|$ с центром $a$ и радиусом $r$).

\textbf{Однолистность в круге:} $f\in\mathcal U(B(0,R))$ --- $f$ инъективна в $B(0,R)$ (билет 1).
\end{defs}

\subsection{Лемма Шварца и лемма Шварца--Пика (напоминание)}

\begin{thmbox}{Лемма Шварца (Предл.~12.1)}
Пусть $f\in H(\DD)$, $|f|\le1$ в $\DD$ и $f(0)=0$. Тогда $|f(z)|\le|z|$ в $\DD$ и $|f'(0)|\le1$; если же
$|f(z_0)|=|z_0|$ при некотором $z_0\in\DD\setminus\{0\}$ или $|f'(0)|=1$, то $f(z)=e^{i\theta}z$.
\end{thmbox}
\begin{proof}[Схема]
$F(z)=f(z)/z$ ($F(0)=f'(0)$) голоморфна в $\DD$; по принципу максимума при $|z|\le r<1$ имеем
$|F(z)|\le\max_{|\zeta|=r}|f(\zeta)|/r\le1/r$; устремляя $r\to1$, получаем $|F|\le1$. Если $|F|$ достигает значения $1$
внутри $\DD$, то $F\equiv e^{i\theta}$ по принципу максимума.
\end{proof}

\begin{thmbox}{Лемма Шварца--Пика (Предл.~12.3)}
Пусть $g\in H(\DD)$, $g(\DD)\subset\DD$. Тогда для всех $z,w\in\DD$
\[
\left|\frac{g(z)-g(w)}{1-\clos{g(w)}g(z)}\right|\le\left|\frac{z-w}{1-\clos wz}\right|,
\quad\text{т.е.}\quad |\varphi_{g(w)}(g(z))|\le|\varphi_w(z)|.
\]
\end{thmbox}
\begin{proof}
Функция $h=\varphi_{g(w)}\circ g\circ\varphi_{-w}$ голоморфна в $\DD$, $h(\DD)\subset\DD$ и
$h(0)=\varphi_{g(w)}(g(w))=0$. По лемме Шварца $|h(\zeta)|\le|\zeta|$; подставляем $\zeta=\varphi_w(z)$ (тогда
$\varphi_{-w}(\zeta)=z$).
\end{proof}

\begin{detail}{если $f\in H(\DD)$, $|f|\le1$ и $f(0)=0$, то $f(\DD)\subset\DD$}
Если $|f(z_0)|=1$ в точке $z_0\in\DD$, то $|f|$ достигает максимума внутри и $f$ постоянна (принцип максимума), т.е.
$|f|\equiv1$, что противоречит $f(0)=0$. Поэтому к таким $f$ (и их композициям с автоморфизмами) применима лемма
Шварца--Пика.
\end{detail}

\subsection{Формулировка теоремы}

Пусть $f\in H(\DD)$, $f(0)=0$, $|f|\le1$ в $\DD$. По лемме Шварца $|f'(0)|\le1$, а если $f'(0)\ne0$, то $f$
однолистна в некотором круге $B(0,r)$ (билет 1). Естественный вопрос: \emph{как связаны $r$ и $|f'(0)|$?}

\begin{thmbox}{Предложение 13.4 (теорема Ландау о круге однолистности)}
Пусть $f\in H(\DD)$, $f(0)=0$ и $|f(z)|\le1$ при $z\in\DD$, и пусть $|f'(0)|\ge\alpha$ при некотором $\alpha\in(0,1)$.
Тогда $f\in\mathcal U\bigl(B(0,R_\alpha)\bigr)$, где
\[
R_\alpha=\frac{1-\sqrt{1-\alpha^2}}{\alpha}=\frac{\alpha}{1+\sqrt{1-\alpha^2}},
\]
причём это значение $R_\alpha$ максимально возможное: для любой области $G$ с условием
$B(0,R_\alpha)\subsetneq G\subset\DD$ найдётся функция указанного вида, не однолистная в $G$.
\end{thmbox}

\begin{detail}{свойства числа $R_\alpha$}
Равенство двух выражений: $\frac{1-\sqrt{1-\alpha^2}}{\alpha}\cdot\frac{1+\sqrt{1-\alpha^2}}{1+\sqrt{1-\alpha^2}}
=\frac{1-(1-\alpha^2)}{\alpha(1+\sqrt{1-\alpha^2})}=\frac{\alpha}{1+\sqrt{1-\alpha^2}}$. Отсюда $0<R_\alpha<\alpha<1$.

$R_\alpha$ --- меньший корень уравнения $\alpha R^2-2R+\alpha=0$, т.е. уравнения
\begin{equation}\label{t4:Ralpha}
\frac{2R}{1+R^2}=\alpha
\end{equation}
(корни $R=\frac{1\pm\sqrt{1-\alpha^2}}{\alpha}$, их произведение равно $1$, так что ровно один корень лежит в $(0,1)$).
Функция $\psi(R)=\frac{2R}{1+R^2}$ строго возрастает на $[0,1]$ ($\psi'(R)=\frac{2(1-R^2)}{(1+R^2)^2}>0$ при $R<1$),
$\psi(0)=0$, $\psi(1)=1$. Поэтому $R_\alpha=\psi^{-1}(\alpha)$ возрастает по $\alpha$, $R_\alpha\to1$ при $\alpha\to1$ и
$R_\alpha\sim\alpha/2$ при $\alpha\to0$.
\end{detail}

Доказательство (по существу --- следствие леммы Шварца--Пика) основано на двух наблюдениях.

\subsection{Наблюдение 1: оценка Ландау}

\begin{thmbox}{Оценка Ландау}
Пусть $f\in H(\DD)$, $f(0)=0$, $|f|\le1$ в $\DD$. Если $a,b\in\DD$, $a\ne b$, и $f(a)=f(b)=:c$, то $|c|\le|a|\,|b|$.
\end{thmbox}
\begin{proof}
Если $f$ постоянна, то $f\equiv0$ и $c=0$ --- доказывать нечего. Пусть $f$ непостоянна; тогда $f(\DD)\subset\DD$ (деталь
выше), в частности $|c|<1$. Положим
\[
g(z)=\varphi_c(f(z))=\frac{f(z)-c}{1-\clos cf(z)} .
\]
Тогда $g\in H(\DD)$ (знаменатель не обращается в нуль, так как $|\clos cf(z)|<1$), $g(\DD)\subset\DD$ (композиция $f$ и
автоморфизма $\varphi_c$ круга), и $g(a)=g(b)=0$, $g(0)=\varphi_c(0)=-c$.

\textbf{Шаг 1.} По лемме Шварца--Пика для $g$ с $w=a$ (учитывая $g(a)=0$, $\varphi_0=\mathrm{id}$):
\[
|g(z)|=\left|\frac{g(z)-g(a)}{1-\clos{g(a)}g(z)}\right|\le\left|\frac{z-a}{1-\clos az}\right|=|\varphi_a(z)|,\qquad z\in\DD.
\]
Функция $g_1=g/\varphi_a$ голоморфна в $\DD\setminus\{a\}$, а в точке $a$ (простом нуле $\varphi_a$, нуле $g$) имеет
устранимую особенность; после доопределения $g_1\in H(\DD)$, и $|g_1|\le1$ в $\DD\setminus\{a\}$, а по непрерывности ---
во всём $\DD$. Кроме того, $g_1(b)=g(b)/\varphi_a(b)=0$ (здесь $\varphi_a(b)\ne0$, так как $b\ne a$). Значит, $g_1$
непостоянна или $\equiv0$; в любом случае $|g_1|\le1$ и, если $g_1\not\equiv0$, то $g_1(\DD)\subset\DD$ (принцип максимума).

\textbf{Шаг 2.} Если $g_1\equiv0$, то $g\equiv0$, т.е. $f\equiv c$ --- исключено. Применим лемму Шварца--Пика к $g_1$ с $w=b$:
$|g_1(z)|\le|\varphi_b(z)|$. Как и выше, $h=g_1/\varphi_b\in H(\DD)$ и $|h|\le1$ в $\DD$. Итак,
\[
g(z)=\varphi_a(z)\varphi_b(z)h(z),\qquad |h|\le1 .
\]

\textbf{Шаг 3.} При $z=0$: $-c=g(0)=\varphi_a(0)\varphi_b(0)h(0)=(-a)(-b)h(0)=ab\,h(0)$, откуда $|c|=|a|\,|b|\,|h(0)|\le|a|\,|b|$.
\end{proof}

\begin{remarkbox}{Замечание к слайду 282}
На слайде функция $h=g/(\varphi_a\varphi_b)$ вводится сразу, а неравенство $|h|\le1$ обосновывается ссылкой на лемму
Шварца--Пика в виде $|g|\le|\varphi_a|$ «и аналогично с заменой $a$ на $b$». Из двух неравенств
$|g|\le|\varphi_a|$ и $|g|\le|\varphi_b|$ неравенство $|g|\le|\varphi_a\varphi_b|$ напрямую \emph{не} следует;
корректно действовать в два шага, как выше: сначала разделить на $\varphi_a$, а затем применить Шварца--Пика к $g_1=g/\varphi_a$.
(Альтернатива: $|h(z)|\le\max_{|\zeta|=r}\frac{|g(\zeta)|}{|\varphi_a(\zeta)\varphi_b(\zeta)|}\le\frac1{\min_{|\zeta|=r}|\varphi_a(\zeta)\varphi_b(\zeta)|}\to1$
при $r\to1$ по принципу максимума, так как $|\varphi_a\varphi_b|\to1$ равномерно при $|\zeta|\to1$.)
\end{remarkbox}

\subsection{Наблюдение 2: неевклидовы круги и образ $F=f(z)/z$}

\begin{thmbox}{Лемма (неевклидов круг --- евклидов круг)}
Пусть $a\in\DD$, $0<r<1$. Тогда $K(a,r)=B(\tilde a,\tilde r)$, где
\[
\tilde a=w(a,r):=\frac{1-r^2}{1-r^2|a|^2}\,a,\qquad \tilde r=\rho(a,r):=r\,\frac{1-|a|^2}{1-r^2|a|^2}.
\]
Кроме того, $\inf\{|z|: z\in B(\tilde a,\tilde r)\}=|\tilde a|-\tilde r=\dfrac{|a|-r}{1-r|a|}$, если $|a|\ge r$.
\end{thmbox}
\begin{proof}
\textbf{Вычисление круга} (упражнение со слайда). Пусть $z\in\CC$, $z\ne1/\clos a$. Неравенство $|\varphi_a(z)|<r$
равносильно $|z-a|^2<r^2|1-\clos az|^2$, т.е.
\[
|z|^2-2\Re(\clos az)+|a|^2<r^2\bigl(1-2\Re(\clos az)+|a|^2|z|^2\bigr)
\iff (1-r^2|a|^2)|z|^2-2(1-r^2)\Re(\clos az)+|a|^2-r^2<0 .
\]
Поделим на $1-r^2|a|^2>0$ и заметим, что $(1-r^2)\clos a/(1-r^2|a|^2)=\clos{\tilde a}$:
\[
|z|^2-2\Re(\clos{\tilde a}z)+\frac{|a|^2-r^2}{1-r^2|a|^2}<0\iff
|z-\tilde a|^2<|\tilde a|^2-\frac{|a|^2-r^2}{1-r^2|a|^2}.
\]
Вычислим правую часть:
\[
|\tilde a|^2-\frac{|a|^2-r^2}{1-r^2|a|^2}=\frac{(1-r^2)^2|a|^2-(|a|^2-r^2)(1-r^2|a|^2)}{(1-r^2|a|^2)^2}.
\]
Числитель:
$(1-2r^2+r^4)|a|^2-\bigl(|a|^2-r^2|a|^4-r^2+r^4|a|^2\bigr)=r^2-2r^2|a|^2+r^2|a|^4=r^2(1-|a|^2)^2$.
Значит, правая часть равна $\tilde r^{\,2}$ и $\{z\in\CC\setminus\{1/\clos a\}:|\varphi_a(z)|<r\}=B(\tilde a,\tilde r)$
(точка $1/\clos a$ в круг не попадает: при $z\to1/\clos a$ $|\varphi_a(z)|\to\infty$; при $a=0$ всё тривиально:
$\tilde a=0$, $\tilde r=r$). По \eqref{t4:id} при $|z|\ge1$ имеем $|\varphi_a(z)|\ge1>r$, поэтому этот круг лежит в
$\DD$ и совпадает с $K(a,r)$.

\textbf{Минимум модуля.} Для $z\in B(\tilde a,\tilde r)$: $|z|\ge|\tilde a|-|z-\tilde a|>|\tilde a|-\tilde r$, и точная нижняя
грань достигается при приближении к точке $\tilde a(1-\tilde r/|\tilde a|)$ (если $|\tilde a|\ge\tilde r$). Обозначим
$t=|a|$. Тогда
\[
|\tilde a|-\tilde r=\frac{(1-r^2)t-r(1-t^2)}{1-r^2t^2}=\frac{(t-r)+rt(t-r)}{(1-rt)(1+rt)}=\frac{(t-r)(1+rt)}{(1-rt)(1+rt)}=\frac{t-r}{1-rt},
\]
так как $(1-r^2)t-r(1-t^2)=t-r^2t-r+rt^2=(t-r)+rt(t-r)$. Условие $|\tilde a|\ge\tilde r$ равносильно $t\ge r$.
\end{proof}

\begin{thmbox}{Лемма (образ $F$)}
Пусть $f\in H(\DD)$, $|f|\le1$ в $\DD$, $f(0)=0$, и $f(z)\not\equiv e^{i\vartheta}z$ ни при каком $\vartheta\in\RR$.
Положим $F(z)=f(z)/z$ (и $F(0)=f'(0)$). Тогда $F\in H(\DD)$, $F(\DD)\subset\DD$, $|f'(0)|<1$, и при всех $r\in(0,1)$
\[
F\bigl(B(0,r)\bigr)\subset K\bigl(f'(0),r\bigr).
\]
\end{thmbox}
\begin{proof}
$F$ голоморфна в $\DD$ (устранимая особенность в нуле, $F(0)=f'(0)$), и по лемме Шварца $|F|\le1$. Если $|F(z_0)|=1$
при некотором $z_0\in\DD$, то по принципу максимума $F\equiv e^{i\vartheta}$, т.е. $f(z)=e^{i\vartheta}z$ --- исключено.
Значит, $F(\DD)\subset\DD$, в частности $|f'(0)|=|F(0)|<1$. По лемме Шварца--Пика для $F$ с $w=0$:
\[
\bigl|\varphi_{f'(0)}(F(z))\bigr|=\left|\frac{F(z)-f'(0)}{1-\clos{f'(0)}F(z)}\right|\le|\varphi_0(z)|=|z| .
\]
Если $|z|<r$, то $|\varphi_{f'(0)}(F(z))|<r$, т.е. $F(z)\in K(f'(0),r)$.
\end{proof}

\subsection{Доказательство теоремы Ландау}

\begin{proof}[Доказательство первого утверждения (однолистность в $B(0,R_\alpha)$)]
Если $f(z)=e^{i\vartheta}z$, то $f$ однолистна во всём $\DD$. Пусть $f\not\equiv e^{i\vartheta}z$; положим $R=R_\alpha$,
$t=|f'(0)|$; тогда $\alpha\le t<1$ (лемма об образе $F$). Предположим противное: $f\notin\mathcal U(B(0,R))$, т.е.
существуют $z_1,z_2\in B(0,R)$, $z_1\ne z_2$, $f(z_1)=f(z_2)$. Из двух различных точек хотя бы одна ненулевая;
пусть (переобозначив) $z_1\ne0$.

\textbf{(а) Верхняя оценка $|F(z_1)|$.} По оценке Ландау $|f(z_1)|\le|z_1|\,|z_2|<R|z_1|$ (так как $|z_2|<R$ и $|z_1|>0$).
Поэтому $|F(z_1)|=|f(z_1)|/|z_1|<R$.

\textbf{(б) Нижняя оценка $|F(z_1)|$.} По лемме об образе $F$ (с $r=R$) и лемме о неевклидовом круге
\[
F(z_1)\in F\bigl(B(0,R)\bigr)\subset K(f'(0),R)=B(w_0,\rho_0),\qquad w_0=w(f'(0),R),\ \rho_0=\rho(f'(0),R).
\]
Так как $t\ge\alpha>R$ (см. деталь о $R_\alpha$: $R_\alpha<\alpha$), то $|w_0|>\rho_0$, и для каждой точки
$\zeta\in B(w_0,\rho_0)$
\[
|\zeta|>|w_0|-\rho_0=\frac{t-R}{1-Rt}.
\]
\textbf{(в) Противоречие.} Из (а) и (б): $\dfrac{t-R}{1-Rt}<|F(z_1)|<R$. Так как $1-Rt>0$, это равносильно
$t-R<R-R^2t$, т.е. $t(1+R^2)<2R$, т.е.
\[
|f'(0)|=t<\frac{2R}{1+R^2}=\alpha
\]
в силу \eqref{t4:Ralpha}. Это противоречит условию $|f'(0)|\ge\alpha$. Значит, $f$ однолистна в $B(0,R_\alpha)$.
\end{proof}

\begin{remarkbox}{Замечание (опечатка на слайде 283)}
В выражении $|w_0|-\rho_0=\frac{|f'(0)|-R}{1-R|f'(0)|}$ знаменатель содержит $|f'(0)|$; в первой строке слайда 284
модуль в знаменателе пропущен ($1-Rf'(0)$). Также на слайде 284 вводится $A=1/\alpha$, но далее не используется.
\end{remarkbox}

\subsection{Точность константы: экстремальная функция}

\begin{proof}[Доказательство второго утверждения (максимальность $R_\alpha$)]
Пусть $\vartheta\in\RR$. Рассмотрим
\[
f(z)=z\,\frac{z-\alpha e^{i\vartheta}}{\alpha z-e^{i\vartheta}} .
\]
\textbf{Проверка условий теоремы.} Так как $\alpha z-e^{i\vartheta}=-e^{i\vartheta}\bigl(1-\alpha e^{-i\vartheta}z\bigr)$ и
$\clos{\alpha e^{i\vartheta}}=\alpha e^{-i\vartheta}$, имеем
\[
f(z)=-e^{-i\vartheta}\,z\,\varphi_{\alpha e^{i\vartheta}}(z)
\]
--- произведение Бляшке из двух множителей (с нулями $0$ и $\alpha e^{i\vartheta}$). Полюс $e^{i\vartheta}/\alpha$ лежит вне
$\clos\DD$, так что $f\in H(\DD)$; $|f(z)|=|z|\,|\varphi_{\alpha e^{i\vartheta}}(z)|<1$ в $\DD$ (по \eqref{t4:id}); $f(0)=0$; и
\[
f'(0)=\lim_{z\to0}\frac{f(z)}z=\frac{0-\alpha e^{i\vartheta}}{0-e^{i\vartheta}}=\alpha,
\qquad |f'(0)|=\alpha .
\]
Итак, $f$ удовлетворяет всем условиям Предл.~13.4.

\textbf{Вычисление $f'$.} Сделаем замену $z=e^{i\vartheta}u$. Тогда $z-\alpha e^{i\vartheta}=e^{i\vartheta}(u-\alpha)$,
$\alpha z-e^{i\vartheta}=e^{i\vartheta}(\alpha u-1)$, и
\[
f(e^{i\vartheta}u)=e^{i\vartheta}p(u),\qquad p(u)=\frac{u(u-\alpha)}{\alpha u-1}=\frac{u^2-\alpha u}{\alpha u-1}.
\]
По правилу дифференцирования частного
\[
p'(u)=\frac{(2u-\alpha)(\alpha u-1)-\alpha(u^2-\alpha u)}{(\alpha u-1)^2}
=\frac{2\alpha u^2-2u-\alpha^2u+\alpha-\alpha u^2+\alpha^2u}{(\alpha u-1)^2}=\frac{\alpha u^2-2u+\alpha}{(\alpha u-1)^2}.
\]
Дифференцируя тождество $f(e^{i\vartheta}u)=e^{i\vartheta}p(u)$ по $u$, получаем $e^{i\vartheta}f'(e^{i\vartheta}u)=e^{i\vartheta}p'(u)$, т.е.
$f'(z)=p'(e^{-i\vartheta}z)$. Числитель $\alpha u^2-2u+\alpha$ обращается в нуль при
$u=\frac{1\pm\sqrt{1-\alpha^2}}{\alpha}$; корень со знаком минус есть $R_\alpha\in(0,1)$, а знаменатель при $u=R_\alpha$
не равен нулю ($\alpha R_\alpha<1$). Следовательно,
\[
f'(R_\alpha e^{i\vartheta})=p'(R_\alpha)=0 .
\]
Контрольная проверка через \eqref{t4:Ralpha}: $p'(u)=0\iff \alpha(1+u^2)=2u\iff \frac{2u}{1+u^2}=\alpha$.

\textbf{Неоднолистность в $G$.} Пусть $G$ --- область, $B(0,R_\alpha)\subsetneq G\subset\DD$. Покажем, что $G$ содержит точку
окружности $|z|=R_\alpha$. Иначе $G=\bigl(G\cap B(0,R_\alpha)\bigr)\sqcup\bigl(G\cap\{|z|>R_\alpha\}\bigr)$ --- разбиение на два
открытых множества, первое из которых равно $B(0,R_\alpha)\ne\varnothing$, а второе непусто, так как $G\ne B(0,R_\alpha)$; это
противоречит связности $G$. Итак, $R_\alpha e^{i\vartheta}\in G$ при некотором $\vartheta$. Для функции $f$ с этим $\vartheta$
имеем $f'(R_\alpha e^{i\vartheta})=0$, поэтому по Предл.~13.1 $f$ не однолистна даже в точке $R_\alpha e^{i\vartheta}$, а тем
более не однолистна в $G$.
\end{proof}

\begin{remarkbox}{Комментарии}
\begin{enumerate}
\item Экстремальная функция однолистна в $B(0,R_\alpha)$ (по уже доказанному), но не в большем круге: при
$R>R_\alpha$ в круге $B(0,R)$ лежит критическая точка $R_\alpha e^{i\vartheta}$. По следствию из билета 1 вблизи неё $f$
двулистна.
\item Иная запись: при $\vartheta=0$ $f(z)=z\frac{z-\alpha}{\alpha z-1}=\frac{z(\alpha-z)}{1-\alpha z}$. Пример: $\alpha=\frac45$:
$\sqrt{1-\alpha^2}=\frac35$, $R_\alpha=\frac{1-3/5}{4/5}=\frac12$, и действительно $\frac{2\cdot\frac12}{1+\frac14}=\frac45$;
$p'(u)=\frac{\frac45u^2-2u+\frac45}{(\frac45u-1)^2}$ обращается в нуль при $u=\frac12$ и $u=2$.
\item Теорема Ландау --- «количественная» версия критерия однолистности: радиус однолистности оценивается снизу только
через $|f'(0)|$ при нормировке $|f|\le1$, $f(0)=0$; без ограничения $|f|\le1$ такой оценки нет
(например, $f_n(z)=\frac1n\bigl(e^{nz}-1\bigr)$: $f_n'(0)=1$, но $f_n$ не однолистна в $B(0,R)$ при $n>\pi/R$, так как
$f_n(\pm i\pi/n)$ совпадают).
\end{enumerate}
\end{remarkbox}

\begin{remarkbox}{Коротко: план ответа на экзамене}
\begin{enumerate}
\item Постановка: $f\in H(\DD)$, $f(0)=0$, $|f|\le1$, $|f'(0)|\ge\alpha$ $\Rightarrow$ однолистность в $B(0,R_\alpha)$,
$R_\alpha=\frac{1-\sqrt{1-\alpha^2}}\alpha$ --- корень $\frac{2R}{1+R^2}=\alpha$.
\item Оценка Ландау: $f(a)=f(b)=c$ $\Rightarrow$ $|c|\le|a||b|$; через $g=\varphi_c\circ f=\varphi_a\varphi_bh$, $|h|\le1$, $g(0)=-c$.
\item Неевклидов круг $K(a,r)=B(\tilde a,\tilde r)$ с явными формулами; $\min|z|=\frac{|a|-r}{1-r|a|}$.
\item $F=f/z$: $F(B(0,r))\subset K(f'(0),r)$ (Шварц--Пик).
\item От противного: $f(z_1)=f(z_2)$ в $B(0,R)$ $\Rightarrow$ $|F(z_1)|<R$, но $|F(z_1)|>\frac{t-R}{1-Rt}$ $\Rightarrow$ $t<\frac{2R}{1+R^2}=\alpha$.
\item Точность: $f(z)=z\frac{z-\alpha e^{i\vartheta}}{\alpha z-e^{i\vartheta}}$, $f'(0)=\alpha$, $f'(R_\alpha e^{i\vartheta})=0$;
любая область $G\supsetneq B(0,R_\alpha)$ содержит точку окружности $|z|=R_\alpha$.
\end{enumerate}
\end{remarkbox}

% ===== part 1c =====
% ======================================================================
% ЧАСТЬ 1c. Билеты 5, 6, 7
% ======================================================================

\ticket{5}{Принцип симметрии Римана--Шварца. Пример отыскания конформного отображения с~использованием принципа симметрии}

Материал: слайды 285--287 (Теорема 13.1), а также сведения первого семестра о дробно-линейных
отображениях и симметрии (Леммы 3.2--3.4 лекций). Пример в п.~5.5 на слайдах отсутствует --- он требуется
программой и разобран полностью.

\begin{defs}
\begin{itemize}
\item $\Ch=\CC\cup\{\infty\}$ --- расширенная плоскость; окрестности $\infty$ --- множества
$\{|z|>R\}\cup\{\infty\}$. Отображение $f$ со значениями в $\Ch$, заданное в окрестности точки $a\in\Ch$,
\emph{голоморфно в $a$}, если оно голоморфно в подходящих картах: в прообразе берём координату $z-a$
(или $1/z$, если $a=\infty$), в образе --- координату $w$ (если $f(a)\ne\infty$) или $1/w$ (если $f(a)=\infty$).
Например, $f(z)=1/(z-a)$ голоморфна (как отображение в $\Ch$) в каждой точке $\Ch$.
\item \emph{Конформный изоморфизм.} Для областей $D,W\subset\Ch$ через $\mathrm{Conf}(D,W)$ обозначается
множество биекций $f\colon D\to W$, конформных (в смысле карт) в каждой точке $D$. Для $D,W\subset\CC$:
$f\in\mathrm{Conf}(D,W)$ $\iff$ $f\in H(D)$, $f$ взаимно однозначна и $f(D)=W$ (условие $f'\ne0$ тогда
выполнено автоматически по Предложению 13.1, см. билет 1).
\item \emph{Голоморфное продолжение.} Пусть $D\subset\Ch$ --- область, $E\subset D$, $h\colon E\to\CC$.
Если $f\in H(D)$ и $f|_E=h$, то $f$ --- \emph{голоморфное продолжение} $h$ с $E$ на $D$. Если $E$ имеет
предельную точку в $D$, то такое продолжение единственно (теорема единственности).
\item \emph{Обобщённая окружность} в $\Ch$ --- окружность $C(z_0,R)$ или прямая в $\CC$, дополненная точкой $\infty$.
Обозначения: $\widehat\RR=\RR\cup\{\infty\}$, $i\widehat\RR=i\RR\cup\{\infty\}$. Дополнение $\Ch\setminus\Gamma$
обобщённой окружности $\Gamma$ состоит ровно из двух областей (открытый круг и внешность замкнутого круга
в $\Ch$, либо две открытые полуплоскости). Говорят, что множество $D$ \emph{лежит по одну сторону от $\Gamma$},
если $D$ содержится в одной из этих двух компонент. $\CC_\pm=\{\pm\Im z>0\}$.
\item \emph{Дуга} обобщённой окружности $\Gamma$ --- непустое открытое (в $\Gamma$) связное подмножество
$\gamma\subset\Gamma$. Это либо сама $\Gamma$, либо $\Gamma$ без точки, либо <<обычная>> открытая дуга.
\item \emph{Симметрия} (Определение 3.2 лекций). Точки $z_1,z_2$ симметричны относительно окружности
$C(z_0,R)$, если они лежат на одном луче с началом $z_0$ и $|z_1-z_0|\,|z_2-z_0|=R^2$; кроме того,
$z_0$ и $\infty$ симметричны. Эквивалентно:
\[
z_2=s_\Gamma(z_1):=z_0+\frac{R^2}{\overline{z_1}-\overline{z_0}} .
\]
Относительно прямой --- обычная зеркальная симметрия, $\infty$ симметрична сама себе. Примеры:
$s_{\widehat\RR}(z)=\bar z$, \ $s_{i\widehat\RR}(z)=-\bar z$, \ $s_{\TT}(z)=1/\bar z$.
Отображение $s_\Gamma\colon\Ch\to\Ch$ называется \emph{симметрией относительно $\Gamma$}.
\item Для множества $E\subset\Ch$ положим $E^*=s_\Gamma(E)$ --- множество, \emph{симметричное} $E$ относительно $\Gamma$.
\item \emph{Лемма 3.3} (сохранение симметрии): если $\varphi$ --- ДЛО (дробно-линейное отображение), а $z_1,z_2$
симметричны относительно $\Gamma$, то $\varphi(z_1),\varphi(z_2)$ симметричны относительно $\varphi(\Gamma)$, т.е.
$\varphi\circ s_\Gamma=s_{\varphi(\Gamma)}\circ\varphi$. \emph{Лемма 3.2}: ДЛО переводит обобщённые окружности
в обобщённые окружности. \emph{Лемма 3.4}: для любых двух троек различных точек $\Ch$ существует ровно одно ДЛО,
переводящее первую тройку во вторую.
\item \emph{Лемма Гурса}: если $F\in H(G)$, $G\subset\CC$ открыто, $\Delta\subset G$ --- замкнутый треугольник,
то $\int_{\partial\Delta}F\,dz=0$. \emph{Теорема Мореры}: если $F\in C(G)$ и $\int_{\partial\Delta}F\,dz=0$
для любого замкнутого треугольника $\Delta\subset G$, то $F\in H(G)$. (Доказательство: в каждом круге
$B\subset G$ функция $\Phi(z)=\int_{[z_B,z]}F$, $z_B$ --- центр $B$, по условию на треугольниках удовлетворяет
$\Phi(z+h)-\Phi(z)=\int_{[z,z+h]}F$, откуда $\Phi'=F$; значит, $\Phi\in H(B)$ и $F=\Phi'\in H(B)$.)
\end{itemize}
\end{defs}

\subsection{Свойства симметрии относительно обобщённой окружности}

\begin{thmbox}{Лемма 5.1 (свойства отображения $s_\Gamma$)}
Пусть $\Gamma$ --- обобщённая окружность. Тогда
\begin{enumerate}
\item[(а)] если $\varphi$ --- ДЛО и $\varphi(\Gamma)=\widehat\RR$, то $s_\Gamma=\varphi^{-1}\circ s_{\widehat\RR}\circ\varphi$,
где $s_{\widehat\RR}(z)=\bar z$;
\item[(б)] $s_\Gamma$ --- гомеоморфизм $\Ch$ на себя, $s_\Gamma\circ s_\Gamma=\mathrm{id}$, $s_\Gamma(z)=z$ при $z\in\Gamma$;
\item[(в)] $s_\Gamma$ переставляет две компоненты множества $\Ch\setminus\Gamma$;
\item[(г)] для любого ДЛО $\psi$ выполнено $\psi\circ s_\Gamma=s_{\psi(\Gamma)}\circ\psi$.
\end{enumerate}
\end{thmbox}
\begin{proof}
(г) --- это Лемма 3.3. (а): применяем (г) к $\psi=\varphi$: $\varphi\circ s_\Gamma=s_{\widehat\RR}\circ\varphi$,
откуда $s_\Gamma=\varphi^{-1}\circ s_{\widehat\RR}\circ\varphi$. ДЛО $\varphi$ с $\varphi(\Gamma)=\widehat\RR$ существует по Лемме 3.2.
(б): отображение $z\mapsto\bar z$ --- гомеоморфизм $\Ch$ (с $\infty\mapsto\infty$), инволюция, тождественно на $\widehat\RR$;
по (а) $s_\Gamma$ --- композиция гомеоморфизмов, $s_\Gamma\circ s_\Gamma=\varphi^{-1}\circ\bar{\ }\circ\bar{\ }\circ\varphi=\mathrm{id}$,
и при $z\in\Gamma$: $\varphi(z)\in\widehat\RR$, $\overline{\varphi(z)}=\varphi(z)$, т.е. $s_\Gamma(z)=z$.
(в): $\varphi$ --- гомеоморфизм $\Ch$, переводящий $\Gamma$ в $\widehat\RR$, поэтому компоненты $\Ch\setminus\Gamma$ --- это
$\varphi^{-1}(\CC_+)$ и $\varphi^{-1}(\CC_-)$; сопряжение переставляет $\CC_+$ и $\CC_-$, значит, по (а) $s_\Gamma$
переставляет $\varphi^{-1}(\CC_+)$ и $\varphi^{-1}(\CC_-)$.
\end{proof}

\begin{thmbox}{Лемма 5.2 (нормализация)}
Пусть $\gamma\ne\Gamma$ --- дуга обобщённой окружности $\Gamma$, а область $D\subset\Ch$ лежит по одну сторону от $\Gamma$.
Тогда существует ДЛО $\varphi$ такое, что $\varphi(\Gamma)=\widehat\RR$, $\varphi(\gamma)\subset\RR$ (т.е. $\infty\notin\varphi(\gamma)$)
и $\varphi(D)\subset\CC_+$.
\end{thmbox}
\begin{proof}
Так как $\gamma\ne\Gamma$, найдётся $p\in\Gamma\setminus\gamma$. Возьмём на $\Gamma$ ещё две точки $q,r$, отличные от $p$ и друг от друга,
и по Лемме 3.4 --- ДЛО $\varphi_0$ с $\varphi_0(p)=\infty$, $\varphi_0(q)=0$, $\varphi_0(r)=1$. По Лемме 3.2 $\varphi_0(\Gamma)$ ---
обобщённая окружность, проходящая через $\infty,0,1$, т.е. $\widehat\RR$. Так как $p\notin\gamma$ и $\varphi_0$ инъективно,
$\infty\notin\varphi_0(\gamma)$, т.е. $\varphi_0(\gamma)\subset\RR$. Область $D$ лежит в одной компоненте $\Ch\setminus\Gamma$,
поэтому связное множество $\varphi_0(D)$ лежит в одной компоненте $\Ch\setminus\widehat\RR$, т.е. в $\CC_+$ или в $\CC_-$.
В первом случае берём $\varphi=\varphi_0$, во втором $\varphi=-\varphi_0$ (отображение $w\mapsto-w$ сохраняет $\widehat\RR$ и $\RR$
и переводит $\CC_-$ в $\CC_+$).
\end{proof}

\subsection{Формулировка теоремы}

\begin{thmbox}{Теорема 13.1 (принцип симметрии Римана--Шварца)}
Пусть $D$ и $W$ --- области в $\Ch$. Предположим, что $\partial D$ и $\partial W$ содержат дуги $\gamma$ и $\sigma$ обобщённых
окружностей $\Gamma$ и $\Sigma$ соответственно ($\gamma,\sigma$ непусты, открыты и связны в $\Gamma$ и $\Sigma$), причём $D$ и $W$
расположены по одну сторону от $\Gamma$ и $\Sigma$ соответственно. Пусть множества
\[
\widetilde D=D\cup\gamma\cup D^*,\qquad \widetilde W=W\cup\sigma\cup W^*,
\]
где $D^*$ и $W^*$ --- области, симметричные $D$ и $W$ относительно $\Gamma$ и $\Sigma$, являются областями.
Наконец, пусть $f\in\mathrm{Conf}(D,W)$, причём $f\in C(D\cup\gamma)$ и $f$ гомеоморфно отображает $\gamma$ на $\sigma$.
Тогда существует $\widetilde f\in\mathrm{Conf}(\widetilde D,\widetilde W)$ такая, что $\widetilde f|_D=f$. Более того, точки,
симметричные относительно $\Gamma$, переходят при отображении $\widetilde f$ в точки, симметричные относительно $\Sigma$:
\[
\widetilde f\bigl(s_\Gamma(z)\bigr)=s_\Sigma\bigl(\widetilde f(z)\bigr),\qquad z\in\widetilde D
\]
(в частности, $\widetilde f(D^*)=W^*$). Указанная $\widetilde f$ единственна.
\end{thmbox}

\begin{remarkbox}{Замечания к формулировке}
\begin{itemize}
\item Множества $D$, $\gamma$, $D^*$ попарно не пересекаются: $D$ лежит в одной компоненте $\Ch\setminus\Gamma$, $D^*$ --- в другой
(Лемма 5.1(в)), $\gamma\subset\Gamma$. Аналогично для $W,\sigma,W^*$. Поэтому формула
<<$\widetilde f=f$ на $D\cup\gamma$, $\widetilde f=s_\Sigma\circ f\circ s_\Gamma$ на $D^*$>> корректно задаёт функцию;
именно она и окажется искомым продолжением.
\item <<$f\in C(D\cup\gamma)$>> понимается для отображений в $\Ch$ (со сферической метрикой);
значения $f$ на $\gamma$ --- это предельные значения $f$ из $D$.
\item Требование, чтобы $\widetilde D$ было областью, существенно: оно означает, что каждая точка $\gamma$ имеет окрестность,
лежащую в $D\cup\gamma\cup D^*$, т.е. к точкам $\gamma$ область $D$ примыкает <<с одной стороны целиком>>.
\end{itemize}
\end{remarkbox}

\subsection{Доказательство теоремы 13.1}

\begin{proof}
Сначала разберём случай $\gamma\ne\Gamma$; тогда и $\sigma\ne\Sigma$ (если бы $\sigma=\Sigma$, то $\sigma$ была бы гомеоморфна
окружности, а $\gamma$ --- нет: $\gamma$ гомеоморфна интервалу, так как $\Gamma$ минус точка гомеоморфна $\RR$, а связное открытое
подмножество $\RR$ --- интервал; но окружность не гомеоморфна интервалу --- после удаления точки окружность остаётся связной,
а интервал --- нет). Случай $\gamma=\Gamma$ разобран в шаге 5.

\textbf{Шаг 1 (сведение к случаю $\Gamma=\Sigma=\widehat\RR$).} По Лемме 5.2 существуют ДЛО $\varphi$ и $\psi$ с условиями
\[
\varphi(\Gamma)=\widehat\RR,\ \ \varphi(\gamma)\subset\RR,\ \ \varphi(D)\subset\CC_+;\qquad
\psi(\Sigma)=\widehat\RR,\ \ \psi(\sigma)\subset\RR,\ \ \psi(W)\subset\CC_+ .
\]
Положим $D_1=\varphi(D)$, $\gamma_1=\varphi(\gamma)$, $W_1=\psi(W)$, $\sigma_1=\psi(\sigma)$ и
\[
h=\psi\circ f\circ\varphi^{-1}\in\mathrm{Conf}(D_1,W_1).
\]
Проверим, что для $h$ выполнены все условия теоремы с $\Gamma=\Sigma=\widehat\RR$:
\begin{itemize}
\item $D_1,W_1\subset\CC_+$ --- области в $\CC$ (лежат по одну сторону от $\widehat\RR$), $\gamma_1,\sigma_1$ --- дуги $\widehat\RR$,
лежащие в $\RR$, т.е. открытые интервалы $\RR$ (связные открытые подмножества $\RR$); $\gamma_1\subset\partial D_1$,
$\sigma_1\subset\partial W_1$, так как ДЛО --- гомеоморфизмы $\Ch$.
\item По Лемме 5.1(г) $\varphi(D^*)=\varphi(s_\Gamma(D))=\overline{\varphi(D)}=D_1^*$, где теперь $^*$ --- сопряжение.
Поэтому $\varphi(\widetilde D)=D_1\cup\gamma_1\cup D_1^*=:\widetilde D_1$ --- область (образ области при гомеоморфизме);
аналогично $\psi(\widetilde W)=W_1\cup\sigma_1\cup W_1^*=:\widetilde W_1$ --- область.
\item $h\in C(D_1\cup\gamma_1)$ и $h$ гомеоморфно отображает $\gamma_1$ на $\sigma_1$ (композиция гомеоморфизмов).
\item Значения $h$ конечны: $h(D_1)=W_1\subset\CC_+$, $h(\gamma_1)=\sigma_1\subset\RR$. Итак, $h\in H(D_1)$, $h$ инъективна,
$h(D_1)=W_1$.
\end{itemize}
Если мы построим $\widetilde h\in\mathrm{Conf}(\widetilde D_1,\widetilde W_1)$ с $\widetilde h|_{D_1}=h$ и
$\widetilde h(\bar z)=\overline{\widetilde h(z)}$, то функция $\widetilde f=\psi^{-1}\circ\widetilde h\circ\varphi$ лежит в
$\mathrm{Conf}(\widetilde D,\widetilde W)$ (композиция конформных биекций), $\widetilde f|_D=\psi^{-1}\circ h\circ\varphi=f$, и
по Лемме 5.1(г), применённой к $\varphi$ (переводит $\Gamma$ в $\widehat\RR$) и к $\psi^{-1}$ (переводит $\widehat\RR$ в $\Sigma$):
\[
\widetilde f\circ s_\Gamma=\psi^{-1}\circ\widetilde h\circ\varphi\circ s_\Gamma
=\psi^{-1}\circ\widetilde h\circ s_{\widehat\RR}\circ\varphi
=\psi^{-1}\circ s_{\widehat\RR}\circ\widetilde h\circ\varphi
=s_\Sigma\circ\psi^{-1}\circ\widetilde h\circ\varphi=s_\Sigma\circ\widetilde f .
\]
Итак, достаточно доказать теорему в \emph{нормализованной ситуации}: $D\subset\CC_+$, $W\subset\CC_+$, $\gamma,\sigma\subset\RR$ ---
интервалы, $f\in H(D)$ --- биекция $D$ на $W$, $f\in C(D\cup\gamma)$, $f$ --- гомеоморфизм $\gamma$ на $\sigma$
(в частности, $f(x)\in\RR$ при $x\in\gamma$). Ниже до конца шага 4 считаем, что мы в этой ситуации; симметрия --- это $z\mapsto\bar z$.
Отметим, что
\begin{equation}\label{t5:halves}
\widetilde D\cap\CC_+=D,\qquad \widetilde D\cap\RR=\gamma,\qquad \widetilde D\cap\CC_-=D^*,
\end{equation}
так как $D\subset\CC_+$, $\gamma\subset\RR$, $D^*\subset\CC_-$.

\textbf{Шаг 2 (определение продолжения, голоморфность вне $\RR$, непрерывность).} Положим
\[
\widetilde f(z)=\begin{cases} f(z), & z\in D\cup\gamma,\\[2pt] \overline{f(\bar z)}, & z\in D^*.\end{cases}
\]
\emph{$\widetilde f\in H(D^*)$.} Пусть $z_0\in D^*$, тогда $\bar z_0\in D$; выберем $r>0$ с $B(\bar z_0,r)\subset D$ и разложим
$f(w)=\sum_{k\ge0}c_k(w-\bar z_0)^k$ в $B(\bar z_0,r)$. Тогда $B(z_0,r)\subset D^*$, и для $z\in B(z_0,r)$ точка $\bar z\in B(\bar z_0,r)$,
\[
\widetilde f(z)=\overline{\sum_{k\ge0}c_k(\bar z-\bar z_0)^k}=\sum_{k\ge0}\overline{c_k}\,(z-z_0)^k ,
\]
т.е. $\widetilde f$ представима сходящимся степенным рядом в $B(z_0,r)$. Значит, $\widetilde f\in H(D^*)$; конечно, $\widetilde f\in H(D)$.

\emph{$\widetilde f\in C(\widetilde D)$.} В точках открытых множеств $D$ и $D^*$ это ясно. Пусть $x\in\gamma$. Так как $\widetilde D$ открыто,
найдётся $r>0$ с $B(x,r)\subset\widetilde D$; по \eqref{t5:halves} $B(x,r)\cap\clos{\CC_+}\subset D\cup\gamma$ и
$B(x,r)\cap\clos{\CC_-}\subset D^*\cup\gamma$. На $B(x,r)\cap\clos{\CC_+}$ функция $\widetilde f=f$ непрерывна в $x$ (по условию
$f\in C(D\cup\gamma)$). На $B(x,r)\cap\clos{\CC_-}$ имеем $\widetilde f(z)=\overline{f(\bar z)}$ (при $z\in\gamma$ это верно, так как
$f(z)\in\RR$), и это непрерывно в $x$ как композиция непрерывных отображений, причём значение в $x$ равно
$\overline{f(x)}=f(x)$. Функция, непрерывная в точке $x$ на каждом из двух замкнутых (относительно $B(x,r)$) множеств,
покрывающих $B(x,r)$, непрерывна в $x$. Итак, $\widetilde f\in C(\widetilde D)$.

\textbf{Шаг 3 (голоморфность в $\widetilde D$: теорема Мореры).} По теореме Мореры достаточно проверить, что
$\int_{\partial\Delta}\widetilde f\,dz=0$ для любого невырожденного замкнутого треугольника $\Delta\subset\widetilde D$
(для вырожденного <<треугольника>> --- отрезка --- интеграл равен нулю тривиально: каждый участок проходится дважды в противоположных
направлениях).

\emph{Случай 1: $\Delta\cap\RR=\varnothing$.} Выпуклое множество $\Delta$ лежит тогда в $\CC_+$ или в $\CC_-$, т.е. по
\eqref{t5:halves} в $D$ или в $D^*$, где $\widetilde f$ голоморфна; интеграл равен нулю по лемме Гурса.

\emph{Случай 2: $\Delta\subset\clos{\CC_+}$ и $\Delta\cap\RR\ne\varnothing$} (аналогично $\Delta\subset\clos{\CC_-}$).
Пусть $c$ --- центр тяжести $\Delta$; это внутренняя точка невырожденного треугольника, поэтому $\Im c>0$.
Для $t\in(0,1)$ положим $\Delta_t=c+t(\Delta-c)$ (гомотетия с центром $c$). Тогда:
\begin{itemize}
\item $\Delta_t\subset\Delta$ (выпуклость $\Delta$: $c+t(z-c)=(1-t)c+tz$);
\item $\Delta_t\subset\CC_+$: при $z\in\Delta$ имеем $\Im\bigl((1-t)c+tz\bigr)=(1-t)\Im c+t\Im z\ge(1-t)\Im c>0$.
\end{itemize}
Значит, $\Delta_t\subset\widetilde D\cap\CC_+=D$, и по лемме Гурса $\int_{\partial\Delta_t}\widetilde f\,dz=0$. Замена
$\zeta=c+t(z-c)$ переводит $\partial\Delta$ в $\partial\Delta_t$ с сохранением ориентации, $d\zeta=t\,dz$:
\[
0=\int_{\partial\Delta_t}\widetilde f(\zeta)\,d\zeta=t\int_{\partial\Delta}\widetilde f\bigl(c+t(z-c)\bigr)\,dz .
\]
Функция $\widetilde f$ равномерно непрерывна на компакте $\Delta$, а $|c+t(z-c)-z|=(1-t)|z-c|\le(1-t)\diam\Delta$.
Поэтому $\widetilde f(c+t(z-c))\to\widetilde f(z)$ равномерно по $z\in\partial\Delta$ при $t\to1-$, и, переходя к пределу,
получаем $\int_{\partial\Delta}\widetilde f\,dz=0$.

\emph{Случай 3: $\Delta$ содержит точки и из $\CC_+$, и из $\CC_-$.} Прямая $\RR$ разрезает $\Delta$ на два выпуклых многоугольника
$P_+=\Delta\cap\clos{\CC_+}$ и $P_-=\Delta\cap\clos{\CC_-}$ (треугольники или четырёхугольники) с общей стороной $I=\Delta\cap\RR$.
При положительной ориентации границ отрезок $I$ входит в $\partial P_+$ и $\partial P_-$ с противоположными направлениями, а
остальные части $\partial P_\pm$ в сумме дают $\partial\Delta$; поэтому
\[
\int_{\partial\Delta}\widetilde f\,dz=\int_{\partial P_+}\widetilde f\,dz+\int_{\partial P_-}\widetilde f\,dz .
\]
Каждый из $P_\pm$ --- выпуклый многоугольник с непустой внутренностью, лежащий в $\clos{\CC_\pm}\cap\widetilde D$, и рассуждение
случая 2 дословно проходит для $P_\pm$ (гомотетия к внутренней точке $c_\pm$ с $\pm\Im c_\pm>0$; интеграл по границе
выпуклого многоугольника $P\subset D$ от голоморфной функции равен нулю, так как $P$ разбивается диагоналями из одной вершины на
треугольники, лежащие в $P$, и к каждому применима лемма Гурса). Итак, оба интеграла равны нулю.

Во всех случаях $\int_{\partial\Delta}\widetilde f\,dz=0$, и по теореме Мореры $\widetilde f\in H(\widetilde D)$.

\textbf{Шаг 4 (однолистность, образ, симметрия).} \emph{Свойство симметрии:} $\widetilde f(\bar z)=\overline{\widetilde f(z)}$ при
всех $z\in\widetilde D$. Действительно, при $z\in D^*$ это определение; при $z\in D$ имеем $\bar z\in D^*$ и
$\widetilde f(\bar z)=\overline{f(z)}$; при $z\in\gamma$: $\bar z=z$ и $\widetilde f(z)=f(z)\in\RR$.

\emph{Образ.} $\widetilde f(D)=W$, $\widetilde f(\gamma)=\sigma$, $\widetilde f(D^*)=\{\overline{f(w)}:w\in D\}=\{\bar\omega:\omega\in W\}=W^*$. Итак, $\widetilde f(\widetilde D)=W\cup\sigma\cup W^*=\widetilde W$.

\emph{Инъективность.} Пусть $\widetilde f(z_1)=\widetilde f(z_2)$. Так как $W\subset\CC_+$, $\sigma\subset\RR$, $W^*\subset\CC_-$
попарно не пересекаются, точки $z_1,z_2$ лежат в одном и том же из множеств $D$, $\gamma$, $D^*$. Если $z_1,z_2\in D$ --- то
$z_1=z_2$ по однолистности $f$; если $z_1,z_2\in\gamma$ --- по инъективности $f|_\gamma$ (гомеоморфизм); если $z_1,z_2\in D^*$,
то $\overline{f(\bar z_1)}=\overline{f(\bar z_2)}$, т.е. $f(\bar z_1)=f(\bar z_2)$, $\bar z_1=\bar z_2$, $z_1=z_2$.

Итак, $\widetilde f\in H(\widetilde D)$ --- биекция $\widetilde D$ на $\widetilde W$; по Предложению 13.1 (однолистная голоморфная
функция имеет ненулевую производную) $\widetilde f\in\mathrm{Conf}(\widetilde D,\widetilde W)$. Вместе с шагом 1 это доказывает
теорему при $\gamma\ne\Gamma$.

\textbf{Шаг 5 (случай $\gamma=\Gamma$ --- упражнение со слайда 286).} Тогда, как отмечено выше, и $\sigma=\Sigma$
(обратная импликация доказывается так же). Фиксируем точку $p\in\Gamma$ и положим $\gamma_1=\Gamma\setminus\{p\}$,
$\sigma_1=\Sigma\setminus\{f(p)\}=f(\gamma_1)$. Тогда $\gamma_1$ --- дуга $\Gamma$, $\gamma_1\ne\Gamma$ (связность $\Gamma$ без точки:
это образ $\RR$ при гомеоморфизме $\varphi^{-1}$, где $\varphi(p)=\infty$), аналогично $\sigma_1$. Множество
$D\cup\gamma_1\cup D^*=\widetilde D\setminus\{p\}$ --- область (открытое множество без точки остаётся открытым, а связная открытая
область на плоскости или сфере после удаления точки остаётся связной), аналогично $\widetilde W\setminus\{f(p)\}$. Все остальные
условия теоремы для пары $(\gamma_1,\sigma_1)$ выполнены, и по уже доказанному существует
$\widetilde f_1\in\mathrm{Conf}\bigl(\widetilde D\setminus\{p\},\widetilde W\setminus\{f(p)\}\bigr)$ с $\widetilde f_1|_D=f$ и
$\widetilde f_1\circ s_\Gamma=s_\Sigma\circ\widetilde f_1$. Доопределим $\widetilde f(p)=f(p)$, $\widetilde f=\widetilde f_1$ вне $p$.

\begin{detail}{$\widetilde f$ непрерывна и голоморфна в точке $p$}
Непрерывность: при $z\to p$, $z\in D\cup\gamma_1$, имеем $\widetilde f(z)=f(z)\to f(p)$ (так как $f\in C(D\cup\Gamma)$). При $z\to p$,
$z\in D^*$, имеем $s_\Gamma(z)\in D$, $s_\Gamma(z)\to s_\Gamma(p)=p$ и
$\widetilde f(z)=s_\Sigma\bigl(f(s_\Gamma(z))\bigr)\to s_\Sigma(f(p))=f(p)$, так как $f(p)\in\Sigma$ неподвижна для $s_\Sigma$.
Итак, $\widetilde f$ непрерывна в $p$ (как отображение в $\Ch$).

Голоморфность: перейдём к картам. Пусть $\zeta$ --- локальная координата в $p$ ($\zeta=z-p$ или $1/z$, если $p=\infty$), $\omega$ ---
локальная координата в $f(p)$ ($\omega=w-f(p)$ или $1/w$, если $f(p)=\infty$). В этих координатах $\widetilde f$ записывается функцией
$g(\zeta)$, определённой и голоморфной в проколотой окрестности $0$ (там $\widetilde f=\widetilde f_1\ne f(p)$ по инъективности,
так что $g$ конечна), непрерывной в $0$ и $g(0)=0$. Значит, $g$ ограничена около $0$, особая точка устранима
(теорема Римана об устранимой особенности), и по непрерывности значение продолжения в $0$ совпадает с $g(0)$. Итак, $\widetilde f$
голоморфна в $p$.
\end{detail}

Далее, $\widetilde f$ --- биекция $\widetilde D\to\widetilde W$ (биекция $\widetilde f_1$ дополнена соответствием $p\mapsto f(p)$),
голоморфна и инъективна, значит, конформна в каждой точке (Предложение 13.1 в картах). Соотношение
$\widetilde f\circ s_\Gamma=s_\Sigma\circ\widetilde f$ в точке $p$ верно, так как $s_\Gamma(p)=p$, $s_\Sigma(f(p))=f(p)$.

\textbf{Шаг 6 (единственность).} Пусть $\widetilde f,\widehat f\in\mathrm{Conf}(\widetilde D,\widetilde W)$ и $\widetilde f=\widehat f=f$ на $D$.
Множество $\{z\in\widetilde D:\widetilde f(z)=\widehat f(z)\}$ содержит непустое открытое множество $D$. По теореме единственности
(для голоморфных отображений в $\Ch$ её можно применить в картах: множество точек, в окрестности которых $\widetilde f\equiv\widehat f$,
открыто, непусто и, по теореме единственности в картах, замкнуто в $\widetilde D$) и связности $\widetilde D$ получаем
$\widetilde f\equiv\widehat f$.
\end{proof}

\begin{remarkbox}{Замечание (о доказательстве на слайде 287)}
На слайде треугольник $\Delta$, пересекающий $\RR$, разрезается прямыми $\Im z=\pm\delta$ на части $X_\delta$, $Y_\delta$ (лежащие в
$D$ и $D^*$, где интеграл равен нулю по теореме Коши) и <<полосу>> $K_\delta=\Delta\cap\{|\Im z|\le\delta\}$, интеграл по границе
которой стремится к нулю при $\delta\to0$ в силу равномерной непрерывности $\widetilde f$ (горизонтальные стороны $K_\delta$
сливаются в отрезок $\Delta\cap\RR$, проходимый в противоположных направлениях, а наклонные участки имеют длину $O(\delta)$).
Приведённый выше вариант с гомотетией --- то же рассуждение, но без анализа геометрии полосы. Случай $\Delta\cap\RR\ne\varnothing$
при $\Delta\not\subset D\cup D^*$ (на слайде написано <<$\Delta\cap\gamma\ne\varnothing$>>) --- это одно и то же, так как
$\Delta\cap\RR\subset\widetilde D\cap\RR=\gamma$.
\end{remarkbox}

\subsection{Принцип симметрии для неоднолистных функций}

Шаги 2--3 доказательства не использовали однолистность. Это даёт часто используемое следствие.

\begin{thmbox}{Следствие 5.3 (принцип симметрии для функций)}
Пусть $D\subset\CC_+$ --- область, $\gamma\subset\partial D\cap\RR$ --- интервал, и $\widetilde D=D\cup\gamma\cup D^*$ --- область.
Если $f\in H(D)\cap C(D\cup\gamma)$ и $f(\gamma)\subset\RR$, то функция $\widetilde f$, равная $f$ на $D\cup\gamma$ и
$\overline{f(\bar z)}$ на $D^*$, голоморфна в $\widetilde D$. В частности, если $f\in H(D)$ непрерывна вплоть до $\gamma$ и
$f(\gamma)$ лежит на обобщённой окружности $\Sigma$, то (после применения ДЛО, переводящего $\Sigma$ в $\widehat\RR$) $f$
голоморфно продолжается через $\gamma$ по формуле $\widetilde f=s_\Sigma\circ f\circ s_{\widehat\RR}$ (если значения конечны).
\end{thmbox}

\subsection{Пример: конформное отображение внешности креста}

\textbf{Задача.} Пусть $a,b>0$. Найти конформное отображение области
\[
E=\Ch\setminus\bigl([-a,a]\cup[-ib,ib]\bigr)\qquad\text{(внешность <<креста>>)}
\]
на внешность единичного круга $\{\zeta\in\Ch:|\zeta|>1\}$ с условием $\infty\mapsto\infty$.

Положим $c=\sqrt{a^2+b^2}$. \textbf{Ответ:}
\begin{equation}\label{t5:answer}
F(z)=\frac{\sqrt{z^2-a^2}+\sqrt{z^2+b^2}}{\sqrt{a^2+b^2}},
\end{equation}
где $\sqrt{z^2-a^2}$ --- ветвь в $\Ch\setminus[-a,a]$, а $\sqrt{z^2+b^2}$ --- ветвь в $\Ch\setminus[-ib,ib]$, обе ведущие себя
как $z+O(1/z)$ при $z\to\infty$. Ниже --- вывод с помощью принципа симметрии и проверка.

\textbf{Идея.} Область $E$ симметрична относительно мнимой оси. Её правая половина
\[
E_+=E\cap\{\Re z>0\}=\{\Re z>0\}\setminus(0,a]
\]
--- <<полуплоскость с разрезом>> --- отображается на полуплоскость элементарными функциями. Граница $E_+$ на мнимой оси
содержит дугу
\[
\gamma=\{iy:\ |y|>b\}\cup\{\infty\}\subset\Gamma=i\widehat\RR ,
\]
а $E=E_+\cup\gamma\cup E_+^*$, где $E_+^*=\{-\bar z:z\in E_+\}=\{\Re z<0\}\setminus[-a,0)$. Если отображение $E_+$ на
полуплоскость переводит $\gamma$ в дугу прямой, то принцип симметрии сразу даёт отображение всей $E$.

\textbf{Шаг 1: отображение $E_+$.} Рассмотрим функцию $f(z)=\sqrt{z^2-a^2}$ с \emph{главной} ветвью
корня $\sqrt{u}=|u|^{1/2}e^{i\arg u/2}$, $\arg u\in(-\pi,\pi)$, определённой в $\CC\setminus(-\infty,0]$.

\begin{detail}{$f\in\mathrm{Conf}\bigl(E_+,\{\Re w>0\}\bigr)$}
Отображение $z\mapsto z^2$ конформно переводит $\{\Re z>0\}$ на $\CC\setminus(-\infty,0]$ (обратное --- главная ветвь $\sqrt{\ }$).
При этом $(0,a]\mapsto(0,a^2]$, поэтому $E_+\mapsto\CC\setminus(-\infty,a^2]$ (биективно). Сдвиг $u\mapsto u-a^2$ переводит
$\CC\setminus(-\infty,a^2]$ на $\CC\setminus(-\infty,0]$, а главная ветвь корня конформно отображает $\CC\setminus(-\infty,0]$
на $\{\Re w>0\}$. Композиция трёх конформных биекций --- конформная биекция.
\end{detail}

\begin{detail}{$f$ непрерывно продолжается на $E_+\cup\gamma$ и гомеоморфно отображает $\gamma$ на $\sigma$}
Главная ветвь корня непрерывно продолжается на $\clos{\CC_+}\setminus\{0\}$ (с $\arg u\in[0,\pi]$) и на $\clos{\CC_-}\setminus\{0\}$
(с $\arg u\in[-\pi,0]$). Пусть $z=x+iy$, $x\ge0$, $y>b$. Тогда $\Im(z^2-a^2)=2xy\ge0$ и $z^2-a^2\ne0$, поэтому
$f$ непрерывна на $\{x\ge0,\ y>b\}$ (с первым продолжением корня), и при $x=0$: $z^2-a^2=-(y^2+a^2)$, $\arg=\pi$,
\[
f(iy)=i\sqrt{y^2+a^2}\qquad(y>b).
\]
Аналогично при $y<-b$ (здесь $\Im(z^2-a^2)=2xy\le0$): $f(iy)=-i\sqrt{y^2+a^2}$. Наконец, $|f(z)|^2=|z^2-a^2|\to\infty$ при
$z\to\infty$, так что полагаем $f(\infty)=\infty$. Итак, $f(iy)=i\,\mathrm{sgn}(y)\sqrt{y^2+a^2}$ на $\gamma$. Функция
$y\mapsto\mathrm{sgn}(y)\sqrt{y^2+a^2}$ строго возрастает на $(-\infty,-b)$ и на $(b,+\infty)$ и биективно отображает
$\{|y|>b\}$ на $\{|t|>c\}$, $c=\sqrt{a^2+b^2}$, причём $|t|\to\infty$ $\iff$ $|y|\to\infty$. Значит, $f$ гомеоморфно отображает
$\gamma$ на дугу
\[
\sigma=\{it:\ |t|>c\}\cup\{\infty\}\subset\Sigma=i\widehat\RR .
\]
\end{detail}

\textbf{Шаг 2: применение принципа симметрии.} Проверим условия Теоремы 13.1 для $D=E_+$, $W=\{\Re w>0\}$,
$\Gamma=\Sigma=i\widehat\RR$ (симметрия $s(z)=-\bar z$):
\begin{itemize}
\item $\gamma\subset\partial E_+$, $\sigma\subset\partial W$ --- дуги (открыты и связны в $i\widehat\RR$: окрестность $\infty$ на $i\widehat\RR$
связывает два луча); $E_+$ и $W$ лежат в правой полуплоскости, т.е. по одну сторону от $i\widehat\RR$;
\item $\widetilde D=E_+\cup\gamma\cup E_+^*=E$ --- область: действительно, $E\cap i\widehat\RR=\gamma$, $E\cap\{\Re z>0\}=E_+$,
$E\cap\{\Re z<0\}=E_+^*$;
\item $\widetilde W=W\cup\sigma\cup W^*=\Ch\setminus[-ic,ic]$ --- область;
\item $f\in\mathrm{Conf}(E_+,W)$, $f\in C(E_+\cup\gamma)$, $f\colon\gamma\to\sigma$ --- гомеоморфизм (шаг 1).
\end{itemize}
По Теореме 13.1 существует $\widetilde f\in\mathrm{Conf}\bigl(E,\ \Ch\setminus[-ic,ic]\bigr)$, $\widetilde f|_{E_+}=f$,
$\widetilde f(-\bar z)=-\overline{\widetilde f(z)}$. Явно, при $\Re z<0$:
\[
\widetilde f(z)=-\overline{f(-\bar z)}=-\overline{\sqrt{\bar z^{\,2}-a^2}}=-\sqrt{z^2-a^2}
\]
(главная ветвь корня перестановочна с сопряжением: $\sqrt{\bar u}=\overline{\sqrt u}$ при $u\notin(-\infty,0]$).

\begin{detail}{$\widetilde f$ --- это ветвь $\sqrt{z^2-a^2}$ в $\Ch\setminus[-a,a]$ с асимптотикой $z$ на $\infty$}
Положим $R_a(z)=z\sqrt{1-a^2/z^2}$ (главная ветвь корня). Так как $1-a^2/z^2\in(-\infty,0]$ $\iff$ $a^2/z^2\in[1,+\infty)$
$\iff$ $z^2\in(0,a^2]$ $\iff$ $z\in[-a,a]\setminus\{0\}$, функция $R_a$ голоморфна в $\CC\setminus[-a,a]$, имеет простой полюс
в $\infty$, $R_a^2=z^2-a^2$ и $R_a(z)=z-\frac{a^2}{2z}+O(z^{-3})$ при $z\to\infty$. С другой стороны, $\widetilde f\in H(E)$
(в конечных точках) и $\widetilde f^{\,2}=z^2-a^2$ на открытом множестве $E_+\cup E_+^*$, а по непрерывности --- всюду в $E$.
Отношение $\widetilde f/R_a$ непрерывно в связной области $E\setminus\{\infty\}$ (знаменатель там не обращается в нуль) и
принимает лишь значения $\pm1$, значит, постоянно; при $z=x>a$ обе функции положительны. Итак, $\widetilde f=R_a$ на $E$.
\end{detail}

\textbf{Шаг 3: от внешности отрезка к внешности круга.} Отображение $\omega=-i\,\widetilde f(z)/c$ переводит $E$ конформно на
$\Ch\setminus[-1,1]$ (поворот на $-\pi/2$ и растяжение переводят $[-ic,ic]$ в $[-1,1]$). Функция Жуковского
$\Zh(\zeta)=\frac12\bigl(\zeta+\frac1\zeta\bigr)$ конформно отображает $\{|\zeta|>1\}$ на $\Ch\setminus[-1,1]$ (с $\infty\mapsto\infty$),
обратное отображение
\[
\zeta=\Zh^{-1}(\omega)=\omega+s(\omega),\qquad s(\omega)^2=\omega^2-1,\qquad s(\omega)\sim\omega\ (\omega\to\infty).
\]
Подставим $\omega=-i\widetilde f/c$: $\omega^2-1=-\dfrac{z^2-a^2}{c^2}-1=-\dfrac{z^2+b^2}{c^2}$ (так как $c^2-a^2=b^2$).
Положим $g(z)=i c\, s(\omega(z))$; тогда $g\in H(E)$, $g^2=-c^2(\omega^2-1)=z^2+b^2$, и $g(z)\sim ic\,\omega=\widetilde f(z)\sim z$ при
$z\to\infty$. Как и выше, $g=R_b'$ на $E$, где $R_b'(z)=z\sqrt{1+b^2/z^2}$ (главная ветвь) --- голоморфная в $\CC\setminus[-ib,ib]$ ветвь
$\sqrt{z^2+b^2}$ с асимптотикой $z$ (здесь $1+b^2/z^2\in(-\infty,0]$ $\iff$ $z^2\in[-b^2,0)$ $\iff$ $z\in[-ib,ib]\setminus\{0\}$;
знак фиксируется асимптотикой: $g/R_b'\equiv\pm1$ и $g/R_b'\to1$ при $z\to\infty$). Получаем
\[
\zeta=\omega+s(\omega)=-\frac ic\bigl(\widetilde f(z)+g(z)\bigr),
\]
и это конформное отображение $E$ на $\{|\zeta|>1\}$. Умножение на $i$ сохраняет внешность круга, поэтому
$F=i\zeta=(\widetilde f+g)/c$, т.е. \eqref{t5:answer}, --- искомое отображение, $F(\infty)=\infty$, $F(z)=\frac{2}{c}z+O(1/z)$.

\begin{detail}{Проверка ответа}
Обозначим $u=\sqrt{z^2-a^2}$, $v=\sqrt{z^2+b^2}$ (указанные ветви). Тогда $v^2-u^2=a^2+b^2=c^2$, т.е. $(v+u)(v-u)=c^2$, откуда
\[
\frac1F=\frac{v-u}{c},\qquad F+\frac1F=\frac{2v}{c},\qquad F-\frac1F=\frac{2u}{c}.
\]
\begin{itemize}
\item \emph{Однолистность} (независимая проверка): если $F(z_1)=F(z_2)$, то $u(z_1)=u(z_2)$, а $u=\widetilde f$ однолистна в $E$.
\item \emph{Граница креста переходит в $\TT$.} На берегах отрезка $[-a,a]$ ($z=x$, $|x|\le a$) предельные значения:
$u=\pm i\sqrt{a^2-x^2}$, $v=\pm\sqrt{x^2+b^2}$ (вещественно), поэтому $|F|^2=\frac{(a^2-x^2)+(x^2+b^2)}{c^2}=1$.
На берегах отрезка $[-ib,ib]$ ($z=iy$, $|y|\le b$): $u=\pm i\sqrt{y^2+a^2}$, $v=\pm\sqrt{b^2-y^2}$, и
$|F|^2=\frac{(y^2+a^2)+(b^2-y^2)}{c^2}=1$.
\item \emph{Контрольные значения:} $F(a)=\frac{0+c}{c}=1$, $F(ib)=\frac{ic+0}{c}=i$ (так как $u(ib)=i\sqrt{b^2+a^2}$ по шагу 1);
$F(x)>1$ при $x>a$ (обе ветви положительны и $v>c$); из соотношения $F-1/F=2u/c$ видно также, что
при $w=F(z)$: $z^2=a^2+\frac{c^2}4\bigl(w-w^{-1}\bigr)^2$ --- явная формула для обратного отображения.
\item \emph{Частные случаи.} При $b=0$ (крест вырождается в отрезок $[-a,a]$, $c=a$): $F(z)=\frac{z+\sqrt{z^2-a^2}}{a}$ ---
обратная функция Жуковского для отрезка $[-a,a]$. При $a=0$: $F(z)=\frac{z+\sqrt{z^2+b^2}}{b}$ --- аналогично для $[-ib,ib]$.
\end{itemize}
\end{detail}

\begin{remarkbox}{Чем здесь помог принцип симметрии}
Формула для $\sqrt{z^2-a^2}$ в $E$ <<угадывается>> легко, но доказать \emph{однолистность} в $E$ и найти \emph{образ} всей
области $E$ напрямую неудобно. Принцип симметрии сводит всё к половине области $E_+$, где отображение --- композиция
элементарных конформных отображений, и к проверке поведения на дуге $\gamma$ оси симметрии: однолистность в $E$ и
равенство $\widetilde f(E)=\Ch\setminus[-ic,ic]$ получаются автоматически. Общий рецепт: если область симметрична
относительно обобщённой окружности $\Gamma$, отображаем <<половину>> на область, симметричную относительно некоторой
$\Sigma$, так, чтобы дуга $\gamma=\widetilde D\cap\Gamma$ перешла в дугу $\Sigma$; далее --- Теорема 13.1. Попутно получаем
\emph{конформный радиус} внешности креста в $\infty$ (ёмкость креста): $\lim_{z\to\infty}z/F(z)=c/2=\frac12\sqrt{a^2+b^2}$.
\end{remarkbox}

\begin{remarkbox}{Коротко: план ответа на экзамене}
\begin{enumerate}
\item Определения: голоморфное продолжение, обобщённая окружность, дуга, симметрия $s_\Gamma$ (формулы), $D^*$.
\item Формулировка Теоремы 13.1 (все условия: дуги, одна сторона, $\widetilde D,\widetilde W$ --- области, $f\in C(D\cup\gamma)$,
$f|_\gamma$ --- гомеоморфизм на $\sigma$).
\item Сведение ДЛО к $\Gamma=\Sigma=\widehat\RR$, $D,W\subset\CC_+$; ДЛО сохраняют симметрию (Лемма 3.3).
\item Определение $\widetilde f(z)=\overline{f(\bar z)}$ на $D^*$; голоморфность в $D^*$ через ряды; непрерывность на $\gamma$
(значения на $\gamma$ вещественны).
\item Теорема Мореры: треугольники, не пересекающие $\RR$ (Гурса), касающиеся $\RR$ (гомотетия/полоска и равномерная
непрерывность), пересекающие $\RR$ (разрезание).
\item Однолистность по случаям $D,\gamma,D^*$; образ $W\cup\sigma\cup W^*$; симметрия; единственность.
\item Пример: внешность креста; $\sqrt{z^2-a^2}$ на правой половине, $\gamma\mapsto\sigma$, симметрия $\Rightarrow$ $E\to\Ch\setminus[-ic,ic]$,
далее обратная функция Жуковского; ответ \eqref{t5:answer} и проверка $|F|=1$ на кресте.
\end{enumerate}
\end{remarkbox}

% ======================================================================
\ticket{6}{Теорема Гурвица о нулях и теорема о локально равномерно сходящейся последовательности однолистных функций}

Материал: слайды 288--289 (Предложения 13.5, 13.6, 13.7 и замечание к теореме Гурвица). Предложение 13.7 на слайдах
оставлено как задача --- ниже оно доказано полностью.

\begin{defs}
\begin{itemize}
\item $D\subset\CC$ --- \emph{область} (открытое связное множество); $H(D)$ --- голоморфные в $D$ функции;
$B(a,r)=\{|z-a|<r\}$, $\clos B(a,r)$ --- замкнутый круг, $C(a,r)=\{|z-a|=r\}$; $\|f\|_K=\max_K|f|$.
\item \emph{Локально равномерная сходимость}: $f_n\to f$ локально равномерно в $D$, если $\|f_n-f\|_K\to0$ для каждого
компакта $K\subset D$ (эквивалентно: каждая точка $D$ имеет окрестность, в которой $f_n\to f$ равномерно).
\item $U(a)$ --- функции, однолистные в точке $a$ (взаимно однозначные в некотором круге $B(a,\delta)$);
$U(D)$ --- функции из $H(D)$, взаимно однозначные в $D$ (\emph{однолистные}).
\item $\mathfrak z(f,B)$ --- число нулей $f$ в множестве $B$ с учётом кратности.
\item \emph{Теорема Вейерштрасса}: если $f_n\in H(D)$ и $f_n\to f$ локально равномерно в $D$, то $f\in H(D)$ и
$f_n^{(k)}\to f^{(k)}$ локально равномерно в $D$ при каждом $k\in\NN$.
\item \emph{Теорема единственности}: если $f\in H(D)$, $f\not\equiv0$, то нули $f$ изолированы (не имеют предельных точек в $D$);
каждый нуль $a$ имеет конечную кратность $m$: $f(z)=(z-a)^mg(z)$, $g\in H(D)$, $g(a)\ne0$.
\item \emph{Теорема Руше}: если $f,g$ голоморфны в окрестности замкнутого круга $\clos B$ и $|g(z)|<|f(z)|$ на $\partial B$, то
$\mathfrak z(f+g,B)=\mathfrak z(f,B)$. (Следует из принципа аргумента: $\mathfrak z(f,B)=\frac1{2\pi}\Delta_{\partial B}\arg f$, и
$\Delta_{\partial B}\arg(f+g)=\Delta_{\partial B}\arg f+\Delta_{\partial B}\arg(1+g/f)=\Delta_{\partial B}\arg f$, так как
$|g/f|<1$ на $\partial B$ и кривая $1+g/f$ не обходит $0$.)
\item \emph{Предложение 13.1} (билет 1): $f\in U(a)\iff f'(a)\ne0$; при этом если $|f'(z)-f'(a)|<\frac12|f'(a)|$ в круге $B(a,\rho)$,
то $f$ инъективна в $B(a,\rho)$.
\end{itemize}
\end{defs}

\subsection{Теорема Гурвица}

\begin{thmbox}{Предложение 13.5 (теорема Гурвица)}
Пусть $D\subset\CC$ --- область, $a\in D$, $f_n\in H(D)$, $n\in\NN$, и $f_n\to f$ локально равномерно в $D$. Если $f$ не
постоянна в $D$ и $f(a)=0$, то для любого $r>0$ найдётся $N\in\NN$ такое, что при всех $n>N$ выполнено
$\mathfrak z\bigl(f_n,B(a,r)\bigr)>0$ (т.е. $f_n$ имеет нуль в $B(a,r)$).
\end{thmbox}
\begin{proof}
По теореме Вейерштрасса $f\in H(D)$. Так как $f$ не постоянна, а $D$ связна, $f\not\equiv0$; по теореме единственности нуль $a$
изолирован: существует $\delta_0>0$ с $\clos B(a,\delta_0)\subset D$ и $f(z)\ne0$ при $0<|z-a|\le\delta_0$. Выберем
\[
0<\delta<\min\{r,\ \delta_0\}\quad(\text{тогда } \delta<\dist(a,\partial D)).
\]
Окружность $C(a,\delta)$ --- компакт, на нём $f$ непрерывна и не обращается в нуль, поэтому
$\rho=\min_{z\in C(a,\delta)}|f(z)|>0$. Так как $f_n\to f$ равномерно на компакте $\clos B(a,\delta)\subset D$, найдётся $N$ такое,
что $|f_n(z)-f(z)|<\rho$ при $z\in\clos B(a,\delta)$, $n>N$.

Применим теорему Руше в круге $B(a,\delta)$ к паре $f$ и $g:=f_n-f$ (обе голоморфны в окрестности $\clos B(a,\delta)$): на
$C(a,\delta)$ имеем $|g(z)|<\rho\le|f(z)|$. Поэтому
\[
\mathfrak z\bigl(f_n,B(a,\delta)\bigr)=\mathfrak z\bigl(f+g,B(a,\delta)\bigr)=\mathfrak z\bigl(f,B(a,\delta)\bigr)\ge1
\]
(в $B(a,\delta)$ лежит нуль $a$ функции $f$). Итак, при $n>N$ функция $f_n$ имеет нуль в $B(a,\delta)\subset B(a,r)$.
\end{proof}

\begin{thmbox}{Замечание к Предложению 13.5 (усиленная форма)}
Пусть $D$, $f_n$, $f$ такие же, как выше, и $f\not\equiv0$ в $D$. Пусть $B=B(a,R)$ --- круг с $\clos B\subset D$ и $f(z)\ne0$ при
$z\in C(a,R)$. Тогда существует $N\in\NN$ такое, что $\mathfrak z(f_n,B)=\mathfrak z(f,B)$ при всех $n>N$.
\end{thmbox}
\begin{proof}
Тот же аргумент: $\rho=\min_{C(a,R)}|f|>0$, при $n>N$ имеем $|f_n-f|<\rho$ на $C(a,R)$, и по теореме Руше
$\mathfrak z(f_n,B)=\mathfrak z(f+(f_n-f),B)=\mathfrak z(f,B)$. (Число $\mathfrak z(f,B)$ конечно: нули $f\not\equiv0$ не накапливаются
на компакте $\clos B$.)
\end{proof}

\begin{remarkbox}{Следствия и примеры}
\begin{itemize}
\item \textbf{Следствие (нули предела).} Если $f_n\in H(D)$ не имеют нулей в $D$ и $f_n\to f$ локально равномерно, то либо
$f\equiv0$, либо $f$ не имеет нулей в $D$. \emph{Доказательство}: если $f\not\equiv0$ и $f(a)=0$, то $f$ не постоянна, и по
Предложению 13.5 у $f_n$ есть нули при больших $n$ --- противоречие.
\item \textbf{Следствие (пропущенные значения).} Если все $f_n$ не принимают значения $w_0$, то $f$ либо $\equiv w_0$, либо не
принимает $w_0$ (применить предыдущее к $f_n-w_0$).
\item \textbf{Необходимость непостоянства.} $f_n(z)=e^z/n$ не имеют нулей, $f_n\to0$ локально равномерно в $\CC$, а $f\equiv0$:
утверждение теоремы для $f\equiv0$ ложно.
\item \textbf{Обратное неверно.} Нули допредельных функций могут <<уйти на границу>>: в $D=\DD$ функции $f_n(z)=z-(1-\frac1n)$
имеют нули в $\DD$, а предел $f(z)=z-1$ нулей в $\DD$ не имеет. Это не противоречит теореме: она утверждает, что нули
\emph{предела} наследуются допредельными функциями, а не наоборот.
\item \textbf{Кратности.} Если $a$ --- нуль $f$ кратности $m$, то из усиленной формы следует, что при больших $n$ у $f_n$ в малом
круге $B(a,\delta)$ ровно $m$ нулей с учётом кратности (например, $f_n(z)=z^2-\frac1n\to z^2$: два простых нуля
$\pm1/\sqrt n$ сливаются в двукратный нуль).
\end{itemize}
\end{remarkbox}

\subsection{Предел последовательности однолистных функций}

\begin{thmbox}{Предложение 13.6}
Пусть $D$ --- область в $\CC$, $f_n\in U(D)$, $n\in\NN$. Если $f_n\to f$ локально равномерно в $D$, то либо $f\in U(D)$, либо $f$
постоянна.
\end{thmbox}
\begin{proof}
По теореме Вейерштрасса $f\in H(D)$. Предположим, что $f$ не постоянна, но и не однолистна в $D$, т.е. найдутся $z_1,z_2\in D$,
$z_1\ne z_2$, с $f(z_1)=f(z_2)$. Положим
\[
g_n(z):=f_n(z)-f_n(z_1),\qquad g(z):=f(z)-f(z_1).
\]
Тогда $g_n\in H(D)$ и $g_n\to g$ локально равномерно в $D$: на компакте $K$ имеем
$\|g_n-g\|_K\le\|f_n-f\|_K+|f_n(z_1)-f(z_1)|\le2\|f_n-f\|_{K\cup\{z_1\}}\to0$. Функция $g$ не постоянна (так как $f$ не постоянна)
и $g(z_2)=f(z_2)-f(z_1)=0$. Пусть $r=\frac12|z_1-z_2|>0$. По теореме Гурвица (применённой к $g_n\to g$ в точке $z_2$) при всех
достаточно больших $n$ найдутся точки $z_n\in B(z_2,r)$, в которых $g_n(z_n)=0$; из доказательства теоремы Гурвица видно,
что $z_n\in B(z_2,\delta)$ с $\clos B(z_2,\delta)\subset D$, т.е. $z_n\in D$. Так как $|z_n-z_2|<r=\frac12|z_1-z_2|$, имеем
$z_n\ne z_1$. Но $g_n(z_n)=0$ означает $f_n(z_n)=f_n(z_1)$ при $z_n\ne z_1$, что противоречит однолистности $f_n$ в $D$.
\end{proof}

\begin{remarkbox}{Замечания к Предложению 13.6}
\begin{itemize}
\item Постоянный предел действительно возможен: $f_n(z)=z/n\in U(\CC)$, $f_n\to0$ локально равномерно.
\item \textbf{Типичное применение} (теорема Римана, компактность класса $\Sc$, см. билеты 9--11): если $f_n\in U(D)$,
$f_n\to f$ локально равномерно и $f_n'(a)=1$ (или $|f_n'(a)|\ge c>0$), то по теореме Вейерштрасса $f'(a)=\lim f_n'(a)\ne0$, значит,
$f$ не постоянна и $f\in U(D)$.
\item Аналогично (с тем же доказательством через следствие о пропущенных значениях): предел однолистных функций, не принимающих
значение $w_0$, либо постоянен, либо однолистен и не принимает $w_0$.
\end{itemize}
\end{remarkbox}

\subsection{Однолистность допредельных функций на компактно вложенных кругах}

\begin{thmbox}{Предложение 13.7}
Пусть $D\subset\CC$ --- область, $f\in U(D)$, а $f_n\in H(D)$, $n\in\NN$, и $f_n\to f$ локально равномерно в $D$ при $n\to\infty$.
Тогда для любого круга $B$ с $\clos B\subset D$ найдётся $N\in\NN$ такое, что $f_n\in U(B)$ при всех $n>N$.
\end{thmbox}

\begin{proof}[Доказательство (задача со слайда 289)]
Пусть $B$ --- круг, $\clos B\subset D$. Предположим противное: существует бесконечно много номеров $n$, для которых $f_n$ не
инъективна в $B$. Переходя к подпоследовательности (она по-прежнему сходится к $f$ локально равномерно), можно считать, что
для \emph{каждого} $n$ найдутся точки $z_n,w_n\in B$, $z_n\ne w_n$, с $f_n(z_n)=f_n(w_n)$. Так как $\clos B$ --- компакт, переходя
ещё раз к подпоследовательности, считаем $z_n\to z^*\in\clos B$, $w_n\to w^*\in\clos B$.

\begin{detail}{Вспомогательный факт: $f_n(\zeta_n)\to f(\zeta^*)$, если $\zeta_n\to\zeta^*\in\clos B$, $\zeta_n\in\clos B$}
$|f_n(\zeta_n)-f(\zeta^*)|\le|f_n(\zeta_n)-f(\zeta_n)|+|f(\zeta_n)-f(\zeta^*)|\le\|f_n-f\|_{\clos B}+|f(\zeta_n)-f(\zeta^*)|\to0$
(равномерная сходимость на компакте $\clos B\subset D$ и непрерывность $f$).
\end{detail}

\emph{Случай 1: $z^*\ne w^*$.} По вспомогательному факту $f(z^*)=\lim f_n(z_n)=\lim f_n(w_n)=f(w^*)$, причём $z^*,w^*\in\clos B\subset D$,
$z^*\ne w^*$. Это противоречит однолистности $f$ в $D$.

\emph{Случай 2: $z^*=w^*=:p\in\clos B\subset D$.} Так как $f\in U(D)$, по Предложению 13.1 $f'(p)\ne0$. Обозначим $\eta=|f'(p)|>0$.
По непрерывности $f'$ найдётся $\rho>0$ такое, что $\clos B(p,\rho)\subset D$ и $|f'(z)-f'(p)|<\eta/4$ при $z\in B(p,\rho)$.
По теореме Вейерштрасса $f_n'\to f'$ равномерно на компакте $\clos B(p,\rho)$, поэтому при $n>N_1$ выполнено
$|f_n'(z)-f'(z)|<\eta/4$ на $\clos B(p,\rho)$. Тогда при $n>N_1$, $z\in B(p,\rho)$:
\[
|f_n'(z)-f'(p)|\le|f_n'(z)-f'(z)|+|f'(z)-f'(p)|<\frac\eta2=\frac12|f'(p)| .
\]
Отсюда $f_n$ инъективна в $B(p,\rho)$: для $\zeta_1\ne\zeta_2$ из $B(p,\rho)$ (круг выпукл, $I=[\zeta_1,\zeta_2]\subset B(p,\rho)$)
\[
\bigl|f_n(\zeta_2)-f_n(\zeta_1)-f'(p)(\zeta_2-\zeta_1)\bigr|=\Bigl|\int_I\bigl(f_n'(\zeta)-f'(p)\bigr)\,d\zeta\Bigr|
<\frac{|f'(p)|}{2}\,|\zeta_2-\zeta_1| ,
\]
поэтому $|f_n(\zeta_2)-f_n(\zeta_1)|>\frac12|f'(p)|\,|\zeta_2-\zeta_1|>0$. Но $z_n,w_n\to p$, так что при больших $n$ (больших $N_1$)
$z_n,w_n\in B(p,\rho)$, $z_n\ne w_n$ и $f_n(z_n)=f_n(w_n)$ --- противоречие.

В обоих случаях получено противоречие; значит, $f_n\in U(B)$ при всех достаточно больших $n$.
\end{proof}

\begin{remarkbox}{Замечания к Предложению 13.7}
\begin{itemize}
\item \textbf{Условие $\clos B\subset D$ существенно}, в частности, утверждение неверно для самой области $D$: $f_n(z)=z+z^2/n\to z$
локально равномерно в $D=\CC$, $f(z)=z\in U(\CC)$, но ни одна $f_n$ не однолистна в $\CC$ ($f_n(0)=f_n(-n)=0$). Точки
неоднолистности <<уходят на бесконечность>> --- т.е. на границу области; именно этот эффект исключает компактность $\clos B$.
Пример в круге: $D=\DD$, $f_n(z)=z+\frac{n+1}{2n}z^2\to f(z)=z+\frac12z^2$ локально равномерно (даже равномерно в $\DD$).
Здесь $f\in U(\DD)$: $\Re f'(z)=\Re(1+z)>0$ в выпуклой области $\DD$ (теорема Носиро--Варшавского, билет 2). Но
$f_n'(z)=1+\frac{n+1}{n}z$ обращается в нуль в точке $-\frac{n}{n+1}\in\DD$, поэтому по Предложению 13.1 $f_n$ не однолистна
даже локально, т.е. $f_n\notin U(\DD)$ ни при каком $n$. При этом, в согласии с Предложением 13.7, для круга $B=B(0,r)$, $r<1$,
имеем $f_n\in U(B)$ при $\frac n{n+1}>r$ (по той же теореме Носиро--Варшавского: $\Re f_n'(z)\ge1-\frac{n+1}nr>0$ при
$|z|<r$ и $n$ достаточно большом).
\item Предложения 13.6 и 13.7 вместе: <<однолистность устойчива при локально равномерной сходимости>> --- предел однолистных
однолистен (или постоянен), и, обратно, вблизи однолистной функции голоморфные функции однолистны на компактно вложенных кругах
(и, по тому же доказательству, на любом компакте $K\subset D$: в рассуждении нужна была лишь компактность $\clos B$).
\end{itemize}
\end{remarkbox}

\begin{remarkbox}{Коротко: план ответа на экзамене}
\begin{enumerate}
\item Определения: локально равномерная сходимость, $U(D)$, $\mathfrak z(f,B)$; теоремы Вейерштрасса и Руше.
\item Теорема Гурвица: выбор $\delta$ (изолированность нуля), $\rho=\min_{C(a,\delta)}|f|>0$, равномерная сходимость на
окружности, Руше $\Rightarrow$ $\mathfrak z(f_n,B(a,\delta))=\mathfrak z(f,B(a,\delta))\ge1$.
\item Усиленная форма: совпадение числа нулей; следствие о нулях/пропущенных значениях предела; пример $e^z/n$.
\item Предложение 13.6: $g_n=f_n-f_n(z_1)$, $g=f-f(z_1)$, нуль $g$ в $z_2$, Гурвиц в круге $B(z_2,\frac12|z_1-z_2|)$ даёт
$f_n(z_n)=f_n(z_1)$, $z_n\ne z_1$.
\item Пример $z/n$ (предел постоянен); применение с нормировкой $f_n'(a)=1$.
\item Предложение 13.7: от противного, пары $z_n\ne w_n$ с $f_n(z_n)=f_n(w_n)$, предельные точки; случай $z^*\ne w^*$ --- противоречие
с однолистностью $f$; случай $z^*=w^*=p$ --- $f'(p)\ne0$, $f_n'\to f'$, оценка $|f_n'-f'(p)|<\frac12|f'(p)|$ даёт инъективность $f_n$
около $p$.
\item Контрпример $z+z^2/n$ в $\CC$: условие $\clos B\subset D$ нужно.
\end{enumerate}
\end{remarkbox}

% ======================================================================
\ticket{7}{Вычисление групп конформных автоморфизмов областей $\DD$, $\CC$ и $\Ch$. Отсутствие конформной эквивалентности этих областей}

Материал: слайды 290--292, а также Предложения 12.1, 12.2 (лемма Шварца и автоморфизмы круга, слайды 259--261).
Упражнения слайдов (групповая структура, изоморфизм групп автоморфизмов конформно эквивалентных областей, продолжение
автоморфизма $\CC$ до автоморфизма $\Ch$) доказаны ниже; рассуждение о <<размерностях>> дополнено строгим доказательством.

\begin{defs}
\begin{itemize}
\item $\Ch=\CC\cup\{\infty\}$ --- компактное топологическое пространство (сфера Римана); $\DD=\{|z|<1\}$.
Для $E\subset\Ch$ через $\partial^\sharp E$ обозначается граница $E$ в $\Ch$: $\partial^\sharp\Ch=\varnothing$,
$\partial^\sharp\CC=\{\infty\}$, $\partial^\sharp\DD=\TT$.
\item \emph{Голоморфность в $\infty$ и в точках, где значение равно $\infty$} понимается в картах ($z\mapsto1/z$ в прообразе,
$w\mapsto1/w$ в образе). Например, $f$ голоморфна в $\infty$ и $f(\infty)=\infty$, если $f(z)\to\infty$ при $z\to\infty$ и функция
$1/f(1/\zeta)$ (доопределённая нулём в $0$) голоморфна в окрестности $\zeta=0$.
\item $\mathrm{Conf}(D_1,D_2)$ ($D_1,D_2\subset\Ch$ --- области) --- \emph{конформные изоморфизмы}: биекции $D_1\to D_2$, конформные
(в картах) в каждой точке. Области $D_1,D_2$ \emph{конформно эквивалентны}, если $\mathrm{Conf}(D_1,D_2)\ne\varnothing$.
\item $\CAut(D)=\mathrm{Conf}(D,D)$ --- \emph{группа конформных автоморфизмов} области $D$ (операция --- композиция).
\item $\Lambda$ --- множество дробно-линейных отображений (ДЛО) $z\mapsto\frac{az+b}{cz+d}$, $a,b,c,d\in\CC$, $ad-bc\ne0$
(со стандартными соглашениями: $-d/c\mapsto\infty$, $\infty\mapsto a/c$). Матрицы $M$ и $\lambda M$ ($\lambda\ne0$) задают одно
и то же ДЛО, и композиции ДЛО соответствует произведение матриц; $\Lambda\cong GL(2,\CC)/\{\lambda I\}=PSL(2,\CC)$.
\item $\mathrm{Aff}=\{z\mapsto az+b:\ a\in\CC\setminus\{0\},\ b\in\CC\}$ --- аффинные отображения.
\item $\varphi_a(z)=\dfrac{z-a}{1-\bar az}$, $a\in\DD$ --- элементарный множитель Бляшке.
\item \emph{Изолированные особые точки}: устранимая (функция ограничена около точки), полюс ($f\to\infty$), существенная
(иначе). \emph{Теорема Сохоцкого--Казорати--Вейерштрасса}: в любой проколотой окрестности существенной особой точки образ
функции плотен в $\CC$. \emph{Теорема Лиувилля}: ограниченная целая функция постоянна. \emph{Принцип сохранения области}:
непостоянная голоморфная функция переводит открытые множества в открытые.
\item \emph{Изоморфизм групп} $\Phi\colon G_1\to G_2$ --- биекция с $\Phi(gh)=\Phi(g)\Phi(h)$. \emph{Инволюция} --- элемент $g\ne e$
с $g^2=e$.
\end{itemize}
\end{defs}

\subsection{Общие факты о группах автоморфизмов}

\begin{thmbox}{Предложение 7.1 (упражнения со слайда 290)}
\begin{enumerate}
\item[(а)] Для любой области $D\subset\Ch$ множество $\CAut(D)$ --- группа относительно композиции.
\item[(б)] Если $F\in\mathrm{Conf}(D_1,D_2)$, то отображение $\Phi_F\colon\CAut(D_1)\to\CAut(D_2)$, $\Phi_F(g)=F\circ g\circ F^{-1}$, ---
изоморфизм групп. В частности, конформно эквивалентные области имеют изоморфные группы автоморфизмов.
\end{enumerate}
\end{thmbox}
\begin{proof}
\emph{Предварительно:} если $F\in\mathrm{Conf}(D_1,D_2)$, то $F^{-1}\in\mathrm{Conf}(D_2,D_1)$. Действительно, $F^{-1}$ --- биекция;
в окрестности любой точки $w_0=F(z_0)$ (в картах) $F$ голоморфна с $F'(z_0)\ne0$, и по теореме об обратной функции
существует голоморфная локальная обратная $G$ в окрестности $w_0$ с $G(w_0)=z_0$, $G'(w_0)=1/F'(z_0)\ne0$; так как $F$ ---
биекция, $G$ совпадает с $F^{-1}$ в окрестности $w_0$ (точнее: $F$ инъективна, а $F\circ G=\mathrm{id}$ около $w_0$, поэтому
$G=F^{-1}$ там). Значит, $F^{-1}$ конформна в каждой точке.
Кроме того, композиция конформных изоморфизмов $D_1\to D_2\to D_3$ --- конформный изоморфизм (биекция; голоморфность
композиции и правило цепочки $(F_2\circ F_1)'=F_2'(F_1)F_1'\ne0$ в картах).

(а) Композиция элементов $\CAut(D)$ лежит в $\CAut(D)$; $\mathrm{id}\in\CAut(D)$; обратный элемент лежит в $\CAut(D)$ (см. выше);
ассоциативность композиции очевидна.

(б) $\Phi_F(g)=F\circ g\circ F^{-1}$ --- композиция конформных изоморфизмов $D_2\to D_1\to D_1\to D_2$, т.е. $\Phi_F(g)\in\CAut(D_2)$.
Гомоморфизм: $\Phi_F(g\circ h)=F g h F^{-1}=(F g F^{-1})(F h F^{-1})=\Phi_F(g)\Phi_F(h)$. Биективность: обратное отображение ---
$\Phi_{F^{-1}}\colon\tilde g\mapsto F^{-1}\circ\tilde g\circ F$, так как $\Phi_{F^{-1}}\circ\Phi_F=\mathrm{id}$ и $\Phi_F\circ\Phi_{F^{-1}}=\mathrm{id}$.
\end{proof}

\subsection{Группа $\CAut(\DD)$}

\begin{thmbox}{Предложение 12.1 (лемма Шварца)}
Пусть $f\in H(\DD)$, $|f(z)|\le1$ при $|z|<1$ и $f(0)=0$. Тогда $|f(z)|\le|z|$ при всех $z\in\DD$ и $|f'(0)|\le1$. Если хотя бы одно
из соответствующих равенств выполнено хотя бы в одной точке $z_0\in\DD\setminus\{0\}$ (соответственно $|f'(0)|=1$), то
$f(z)=e^{i\alpha}z$, $\alpha\in\RR$.
\end{thmbox}
\begin{proof}
Функция $g(z)=f(z)/z$, $g(0)=f'(0)$, голоморфна в $\DD$ (устранимая особенность). При $|z|<r<1$ по принципу максимума
$|g(z)|\le\max_{|\zeta|=r}|f(\zeta)|/r\le1/r$; при $r\to1-$ получаем $|g(z)|\le1$, т.е. $|f(z)|\le|z|$ и $|f'(0)|=|g(0)|\le1$. Если
$|g(z_0)|=1$ в некоторой точке $z_0\in\DD$ (это и есть случай равенства), то $|g|$ достигает максимума внутри $\DD$, и по принципу
максимума $g\equiv e^{i\alpha}$.
\end{proof}

\begin{detail}{$\varphi_a\in\CAut(\DD)$ при $a\in\DD$, и $\varphi_a^{-1}=\varphi_{-a}$}
$\varphi_a\in\Lambda$ (определитель матрицы $\begin{pmatrix}1&-a\\-\bar a&1\end{pmatrix}$ равен $1-|a|^2\ne0$), полюс $1/\bar a\notin\clos\DD$.
Тождество
\[
1-|\varphi_a(z)|^2=\frac{|1-\bar az|^2-|z-a|^2}{|1-\bar az|^2}=\frac{(1-|a|^2)(1-|z|^2)}{|1-\bar az|^2}
\]
(раскрываем: $|1-\bar az|^2=1-2\Re(\bar az)+|a|^2|z|^2$, $|z-a|^2=|z|^2-2\Re(\bar az)+|a|^2$, разность равна
$1+|a|^2|z|^2-|z|^2-|a|^2=(1-|a|^2)(1-|z|^2)$) показывает, что $\varphi_a(\DD)\subset\DD$. Прямая проверка:
$\varphi_{-a}(\varphi_a(z))=\dfrac{\frac{z-a}{1-\bar az}+a}{1+\bar a\frac{z-a}{1-\bar az}}=\dfrac{z-a+a-|a|^2z}{1-\bar az+\bar az-|a|^2}=z$, и
аналогично $\varphi_a\circ\varphi_{-a}=\mathrm{id}$. Значит, $\varphi_a$ --- биекция $\DD$ на $\DD$, голоморфная, т.е. $\varphi_a\in\CAut(\DD)$.
\end{detail}

\begin{thmbox}{Предложение 12.2 (общий вид конформных автоморфизмов $\DD$)}
\[
\CAut(\DD)=\Bigl\{z\mapsto e^{i\alpha}\frac{z-a}{1-\bar az}:\ a\in\DD,\ \alpha\in\RR\Bigr\}\subset\Lambda .
\]
Представление $f=e^{i\alpha}\varphi_a$ элемента $f\in\CAut(\DD)$ единственно (с точностью до $\alpha\bmod2\pi$): $a=f^{-1}(0)$,
$e^{i\alpha}=f'(a)(1-|a|^2)$.
\end{thmbox}
\begin{proof}
Включение <<$\supset$>>: $\varphi_a\in\CAut(\DD)$ (см. деталь выше), поворот $z\mapsto e^{i\alpha}z$ лежит в $\CAut(\DD)$, композиция --- тоже.

Включение <<$\subset$>>: пусть $f\in\CAut(\DD)$; так как $f$ --- биекция $\DD$ на себя, существует (единственная) $a\in\DD$ с $f(a)=0$.
Положим $h=f\circ\varphi_{-a}$. Тогда $h\in\CAut(\DD)$, $h(0)=f(a)=0$. По лемме Шварца $|h(z)|\le|z|$. Обратная функция
$h^{-1}\in\CAut(\DD)$ (Предложение 7.1) также удовлетворяет условиям леммы Шварца, поэтому $|h^{-1}(w)|\le|w|$; при $w=h(z)$:
$|z|\le|h(z)|$. Итак, $|h(z)|=|z|$ при всех $z\in\DD$, и по случаю равенства в лемме Шварца $h(z)=e^{i\alpha}z$. Следовательно,
$f=h\circ\varphi_{-a}^{-1}=h\circ\varphi_a$, т.е. $f(z)=e^{i\alpha}\varphi_a(z)$.

Единственность: из $f=e^{i\alpha}\varphi_a$ следует $f(a)=0$, т.е. $a=f^{-1}(0)$ определено однозначно; далее,
$\varphi_a'(z)=\dfrac{1-|a|^2}{(1-\bar az)^2}$, $\varphi_a'(a)=\dfrac1{1-|a|^2}$, так что $e^{i\alpha}=f'(a)(1-|a|^2)$.
\end{proof}

\textbf{<<Размерность>>.} Отображение $(a,e^{i\alpha})\mapsto e^{i\alpha}\varphi_a$ --- биекция $\DD\times\TT\to\CAut(\DD)$. Пространство
параметров $\DD\times\TT$ --- трёхмерное вещественное многообразие (два вещественных параметра $\Re a,\Im a$ и угол $\alpha$), поэтому
говорят, что $\dim_\RR\CAut(\DD)=3$. (Как группа Ли, $\CAut(\DD)\cong PSU(1,1)\cong PSL(2,\RR)$.)

\subsection{Группа $\CAut(\CC)$}

\begin{thmbox}{Предложение 7.2 ($\CAut(\CC)=\mathrm{Aff}$)}
\[
\CAut(\CC)=\{z\mapsto az+b:\ a,b\in\CC,\ a\ne0\}\subset\Lambda,\qquad \dim_\RR\CAut(\CC)=4 .
\]
\end{thmbox}

На слайдах вычисление сводится к $\CAut(\Ch)$: каждый $f\in\CAut(\CC)$ продолжается до $f^*\in\CAut(\Ch)$ условием $f^*(\infty)=\infty$
(<<проверка необходимых деталей оставляется в качестве упражнения>>). Проведём её.

\begin{thmbox}{Лемма 7.3 (упражнение со слайда 292)}
Пусть $f\in\CAut(\CC)$. Тогда $f(z)\to\infty$ при $z\to\infty$, и отображение $f^*$, $f^*|_\CC=f$, $f^*(\infty)=\infty$, лежит в $\CAut(\Ch)$.
\end{thmbox}
\begin{proof}
$f$ --- целая функция, и $\infty$ --- её изолированная особая точка. Выясним её тип.

\emph{Не устранимая.} Иначе $f$ ограничена в $\{|z|>R\}$; на компакте $\clos B(0,R)$ она ограничена как непрерывная; значит, $f$ ---
ограниченная целая функция и по теореме Лиувилля постоянна, что противоречит инъективности.

\emph{Не существенная.} Иначе по теореме Сохоцкого--Казорати--Вейерштрасса множество $f(\{|z|>1\})$ плотно в $\CC$. По принципу
сохранения области $f(\DD)$ --- непустое открытое множество, поэтому оно пересекается с плотным множеством $f(\{|z|>1\})$:
$f(z_1)=f(z_2)$ при $|z_1|<1<|z_2|$, $z_1\ne z_2$ --- противоречие с инъективностью.

Значит, $\infty$ --- полюс $f$, т.е. $f(z)\to\infty$ при $z\to\infty$. Положим $f^*(\infty)=\infty$. Тогда $f^*\colon\Ch\to\Ch$ ---
биекция (так как $f\colon\CC\to\CC$ --- биекция). Голоморфность $f^*$ в $\infty$: функция $q(\zeta)=1/f(1/\zeta)$ определена и
голоморфна в проколотой окрестности $0$ (там $f(1/\zeta)\ne0$ при малых $\zeta$, так как $f\to\infty$), и $q(\zeta)\to0$ при
$\zeta\to0$; по теореме об устранимой особенности $q$ голоморфна в $0$, $q(0)=0$. Наконец, $q$ инъективна около $0$ (композиция
инъективных отображений), поэтому $q'(0)\ne0$ по Предложению 13.1, т.е. $f^*$ конформна в $\infty$. Итак, $f^*\in\CAut(\Ch)$.
\end{proof}

\begin{proof}[Доказательство Предложения 7.2]
Аффинное отображение $z\mapsto az+b$, $a\ne0$, --- голоморфная биекция $\CC\to\CC$ (обратное $w\mapsto(w-b)/a$), т.е. лежит в
$\CAut(\CC)$. Обратно, если $f\in\CAut(\CC)$, то по Лемме 7.3 $f^*\in\CAut(\Ch)$, и по Предложению 7.4 ниже
$f^*(z)=\frac{\alpha z+\beta}{\gamma z+\delta}$. Условие $f^*(\infty)=\infty$ означает $\gamma=0$ (иначе $f^*(\infty)=\alpha/\gamma\ne\infty$);
тогда $\delta\ne0$ и $f(z)=w_1z+w_0$ с $w_1=\alpha/\delta\ne0$, $w_0=\beta/\delta$. (Предложение 7.4 доказывается независимо от
Предложения 7.2, так что круга в рассуждении нет.)

Размерность: $(a,b)\mapsto(z\mapsto az+b)$ --- биекция $(\CC\setminus\{0\})\times\CC\to\CAut(\CC)$, пространство параметров ---
открытое подмножество $\CC^2\cong\RR^4$, так что $\dim_\RR\CAut(\CC)=4$.
\end{proof}

\begin{remarkbox}{Замечание (опечатка на слайде 292)}
На слайде написано $f^*(z)=w_1z+w_1$; правильно $f^*(z)=w_1z+w_0$ (свободный член произволен и не связан с коэффициентом при $z$).
\end{remarkbox}

\begin{remarkbox}{Другое доказательство Предложения 7.2 (без $\Ch$)}
Из доказательства Леммы 7.3: $\infty$ --- полюс целой функции $f$, значит, $f$ --- многочлен (главная часть ряда Лорана в $\infty$
конечна, а остаток --- ограниченная целая функция, т.е. константа). Если $\deg f=n\ge2$, то $f'$ --- многочлен степени $n-1\ge1$ и
имеет нуль $z_0$; по Предложению 13.1 $f$ не однолистна в $z_0$ --- противоречие. Значит, $\deg f=1$.
\end{remarkbox}

\subsection{Группа $\CAut(\Ch)$}

\begin{thmbox}{Предложение 7.4 ($\CAut(\Ch)=\Lambda$)}
\[
\CAut(\Ch)=\Bigl\{z\mapsto\frac{az+b}{cz+d}:\ a,b,c,d\in\CC,\ ad-bc\ne0\Bigr\}=\Lambda,\qquad \dim_\RR\CAut(\Ch)=6 .
\]
\end{thmbox}
\begin{proof}
\emph{$\Lambda\subset\CAut(\Ch)$} (<<обратное включение очевидно>>): ДЛО $T(z)=\frac{az+b}{cz+d}$ --- биекция $\Ch\to\Ch$ (обратное ---
$\frac{dw-b}{-cw+a}\in\Lambda$), и оно конформно в каждой точке $\Ch$: в конечных точках $z\ne-d/c$ производная
$T'(z)=\frac{ad-bc}{(cz+d)^2}\ne0$; в остальных точках (полюс $-d/c$ и $\infty$) --- в картах: $1/T(z)=\frac{cz+d}{az+b}$ имеет в $-d/c$
производную $\frac{c(az+b)-a(cz+d)}{(az+b)^2}=\frac{-(ad-bc)}{(az+b)^2}\ne0$, а $T(1/\zeta)=\frac{a+b\zeta}{c+d\zeta}$ --- снова ДЛО от
$\zeta$, конформное в $\zeta=0$ по уже разобранному.

\emph{$\CAut(\Ch)\subset\Lambda$.} Пусть $f\in\CAut(\Ch)$, $a:=f^{-1}(\infty)$ (единственная точка, так как $f$ --- биекция).

\textbf{Случай $a\ne\infty$.} Тогда $f$ конечна на $\Ch\setminus\{a\}$ и голоморфна там (в частности, голоморфна в $\infty$, т.е.
$f(z)\to f(\infty)\in\CC$ при $z\to\infty$), а в $a$ имеет полюс ($f(z)\to\infty$ при $z\to a$ по непрерывности). Проверим, что это полюс
первого порядка. Функция $g=1/f$, доопределённая по непрерывности $g(a)=0$, голоморфна в окрестности $a$ (это и есть голоморфность
$f$ в $a$ в смысле карт). Если бы нуль $g$ в $a$ имел порядок $\ge2$, то $g'(a)=0$, и по Предложению 13.1 $g$ не была бы однолистна ни
в какой окрестности $a$; но $g=(w\mapsto1/w)\circ f$ --- композиция инъективных отображений, т.е. инъективна --- противоречие.
Итак, $g(z)=(z-a)g_1(z)$, $g_1$ голоморфна около $a$, $g_1(a)\ne0$, и
\[
f(z)=\frac1{(z-a)g_1(z)}=\frac{w_0}{z-a}+\bigl(\text{голоморфная около $a$ функция}\bigr),\qquad w_0=\frac1{g_1(a)}\ne0 .
\]
Рассмотрим $F(z)=f(z)-\frac{w_0}{z-a}$. Она голоморфна в $\CC\setminus\{a\}$, имеет устранимую особенность в $a$ (выше), т.е. $F$ ---
целая; при $z\to\infty$: $F(z)\to f(\infty)-0$, так что $F$ ограничена в окрестности $\infty$, а значит, и на всём $\CC$. По теореме
Лиувилля $F\equiv w_1$. Итак,
\[
f(z)=\frac{w_0}{z-a}+w_1=\frac{w_1z+(w_0-aw_1)}{z-a},\qquad \det\begin{pmatrix}w_1&w_0-aw_1\\1&-a\end{pmatrix}=-w_0\ne0,
\]
т.е. $f\in\Lambda$.

\textbf{Случай $a=\infty$.} Тогда $f\colon\CC\to\CC$ --- биекция, голоморфная, т.е. $f|_\CC\in\CAut(\CC)$, и $f$ голоморфна в $\infty$ с
$f(\infty)=\infty$: функция $q(\zeta)=1/f(1/\zeta)$ голоморфна около $0$, $q(0)=0$, и, как выше (инъективность и Предложение 13.1),
$q'(0)\ne0$: $q(\zeta)=\zeta q_1(\zeta)$, $q_1(0)\ne0$. Тогда $v=1/q_1$ голоморфна около $0$ и
$f(z)=\frac{1}{q(1/z)}=z\,v(1/z)=v(0)z+v'(0)+O(1/z)$ при $z\to\infty$. Функция $F(z)=f(z)-v(0)z$ --- целая и ограниченная около $\infty$
(стремится к $v'(0)$), значит, ограниченная на $\CC$ и по Лиувиллю постоянна: $f(z)=w_1z+w_0$, $w_1=v(0)\ne0$. Итак, $f\in\Lambda$.

\emph{Размерность.} Отображение $\begin{pmatrix}a&b\\c&d\end{pmatrix}\mapsto\frac{az+b}{cz+d}$ переводит $GL(2,\CC)$ (открытое
подмножество $\CC^4$) на $\Lambda$, и два набора задают одно ДЛО тогда и только тогда, когда они пропорциональны.
\begin{detail}{ДЛО определяет матрицу с точностью до множителя}
Пусть матрицы $M,M'$ задают одно и то же ДЛО. Так как композиции ДЛО соответствует произведение матриц, матрица $M'^{-1}M$
задаёт тождественное ДЛО. Если $\frac{\alpha z+\beta}{\gamma z+\delta}\equiv z$, то $\gamma z^2+(\delta-\alpha)z-\beta\equiv0$, т.е.
$\gamma=\beta=0$, $\alpha=\delta$. Значит, $M'^{-1}M=\alpha I$, $M=\alpha M'$.
\end{detail}
Таким образом, $\Lambda$ параметризуется классами пропорциональности ненулевых матриц с $\det\ne0$ --- открытым подмножеством
комплексного проективного пространства $\CC P^3$ комплексной размерности $4-1=3$; $\dim_\RR\Lambda=6$. (Эквивалентно: нормируя
$ad-bc=1$, получаем $\Lambda\cong SL(2,\CC)/\{\pm I\}$, а $SL(2,\CC)$ задаётся в $\CC^4$ одним комплексным уравнением.)
\end{proof}

\subsection{Попарная неэквивалентность $\DD$, $\CC$, $\Ch$}

\begin{thmbox}{Теорема 7.5}
Области $\DD$, $\CC$ и $\Ch$ попарно конформно не эквивалентны.
\end{thmbox}

\begin{proof}[Первое доказательство (прямое)]
\emph{$\Ch$ не эквивалентна ни $\CC$, ни $\DD$.} Конформный изоморфизм --- в частности, непрерывная сюръекция. Непрерывный образ
компакта компактен, а $\Ch$ компактна (гомеоморфна сфере). Множества $\CC$ и $\DD$ не компактны: последовательность $z_n=n$
(соответственно $z_n=1-\frac1n$) не имеет предельной точки в $\CC$ (соответственно в $\DD$). Значит, не существует даже непрерывного
отображения $\Ch$ \emph{на} $\CC$ или $\DD$.

\emph{$\CC$ не эквивалентна $\DD$.} Если $F\in\mathrm{Conf}(\CC,\DD)$, то $F$ --- целая функция с $|F|<1$, по теореме Лиувилля
постоянна, что противоречит биективности. (Обратное отображение $\DD\to\CC$ не нужно рассматривать отдельно: эквивалентность
симметрична, так как $F^{-1}\in\mathrm{Conf}(\DD,\CC)$ влечёт $F\in\mathrm{Conf}(\CC,\DD)$.)
\end{proof}

\begin{proof}[Второе доказательство (через группы автоморфизмов)]
По Предложению 7.1(б) конформно эквивалентные области имеют изоморфные группы автоморфизмов. Покажем, что группы
$G_\DD=\CAut(\DD)$, $G_\CC=\mathrm{Aff}$, $G_{\Ch}=\Lambda$ попарно не изоморфны. Используем свойства, сохраняющиеся при
изоморфизмах групп (изоморфизм переводит инволюции в инволюции, коммутирующие элементы --- в коммутирующие):
\begin{itemize}
\item[(P1)] в группе есть две различные коммутирующие инволюции;
\item[(P2)] для любых инволюций $s_1,s_2,s_3,s_4$ элементы $s_1s_2$ и $s_3s_4$ коммутируют.
\end{itemize}

\emph{Лемма об инволюциях в $\Lambda$.} Пусть $T\in\Lambda$, $T\ne\mathrm{id}$, с матрицей $M$. Тогда $T^2=\mathrm{id}\iff\operatorname{tr}M=0$.
Действительно, по теореме Гамильтона--Кэли $M^2=(\operatorname{tr}M)M-(\det M)I$. Если $\operatorname{tr}M=0$, то $M^2=-(\det M)I$ ---
скаляр, т.е. $T^2=\mathrm{id}$. Обратно, если $M^2=\lambda I$, то $(\operatorname{tr}M)M=(\lambda+\det M)I$; при $\operatorname{tr}M\ne0$
матрица $M$ скалярна и $T=\mathrm{id}$, что исключено.

\emph{$\Lambda$ обладает (P1):} $z\mapsto-z$ и $z\mapsto1/z$ --- различные инволюции, и $-\frac1z=\frac1{-z}$, т.е. они коммутируют.

\emph{$\mathrm{Aff}$ не обладает (P1), но обладает (P2).} Инволюции в $\mathrm{Aff}$: $a(az+b)+b=z$ $\iff$ $a^2=1$, $(a+1)b=0$; при $a=1$
получаем $b=0$ (тождество), поэтому инволюции --- это в точности $s_b(z)=-z+b$, $b\in\CC$. Имеем
$s_b\circ s_{b'}(z)=-(-z+b')+b=z+(b-b')$, $s_{b'}\circ s_b(z)=z+(b'-b)$; они равны лишь при $b=b'$, так что (P1) не выполнено.
Произведение двух инволюций --- сдвиг $z\mapsto z+(b-b')$, а сдвиги коммутируют; значит, (P2) выполнено.

\emph{$\CAut(\DD)$ не обладает ни (P1), ни (P2).} Инволюции: для $f=e^{i\alpha}\varphi_a$ матрица
$\begin{pmatrix}e^{i\alpha}&-ae^{i\alpha}\\-\bar a&1\end{pmatrix}$ имеет след $e^{i\alpha}+1$, который равен нулю лишь при $e^{i\alpha}=-1$.
Значит, инволюции в $\CAut(\DD)$ --- это в точности
\[
\iota_a(z)=\frac{a-z}{1-\bar az},\qquad a\in\DD .
\]
\begin{detail}{У $\iota_a$ ровно одна неподвижная точка в $\DD$}
Уравнение $\iota_a(z)=z$ равносильно $\bar az^2-2z+a=0$. При $a=0$ единственный корень $z=0$. При $a\ne0$ корни
$z_\pm=\frac{1\pm\sqrt{1-|a|^2}}{\bar a}$ (вещественный корень $\sqrt{1-|a|^2}\in(0,1)$), и
$|z_-|=\frac{1-\sqrt{1-|a|^2}}{|a|}=\frac{|a|}{1+\sqrt{1-|a|^2}}<1$, тогда как $|z_+|=\frac{1+\sqrt{1-|a|^2}}{|a|}>1$. Итак, в $\DD$
ровно одна неподвижная точка.
\end{detail}
\emph{(P1) не выполнено:} пусть $g,h$ --- коммутирующие инволюции в $\CAut(\DD)$, $p\in\DD$ --- неподвижная точка $g$. Тогда
$g(h(p))=h(g(p))=h(p)$, т.е. $h(p)\in\DD$ --- неподвижная точка $g$, и по единственности $h(p)=p$. Сопряжём автоморфизмом $\varphi_p$
(Предложение 7.1(б) с $F=\varphi_p$): $g_0=\varphi_p\circ g\circ\varphi_p^{-1}$ и $h_0=\varphi_p\circ h\circ\varphi_p^{-1}$ --- инволюции в
$\CAut(\DD)$, оставляющие $0$ неподвижным. Но $\iota_a(0)=a$, так что инволюция, фиксирующая $0$, есть $\iota_0(z)=-z$. Итак,
$g_0=h_0$, откуда $g=h$.

\emph{(P2) не выполнено:} при $t\in(0,1)$ рассмотрим инволюции $\iota_0,\iota_t,\iota_0,\iota_{it}$ и
\[
A=\iota_0\circ\iota_t\colon z\mapsto\frac{z-t}{1-tz},\qquad B=\iota_0\circ\iota_{it}\colon z\mapsto\frac{z-it}{1+itz}.
\]
Тогда $B(0)=-it$, $A(0)=-t$ и
\[
A(B(0))=\frac{-it-t}{1+it^2},\qquad B(A(0))=\frac{-t-it}{1-it^2}.
\]
Числители равны и отличны от нуля, а знаменатели различны ($t\ne0$), поэтому $A\circ B\ne B\circ A$.

Итак: $\Lambda$ обладает (P1), а $\mathrm{Aff}$ и $\CAut(\DD)$ --- нет; $\mathrm{Aff}$ обладает (P2), а $\CAut(\DD)$ --- нет. Значит, три группы
попарно не изоморфны, и по Предложению 7.1(б) области $\DD$, $\CC$, $\Ch$ попарно конформно не эквивалентны.
\end{proof}

\begin{remarkbox}{Замечание (о рассуждении с размерностями на слайде 292)}
На слайде неизоморфность групп выводится из различия <<размерностей>> $3$, $4$, $6$. Для \emph{абстрактных} групп такое
рассуждение, строго говоря, недостаточно: абстрактный изоморфизм не обязан быть непрерывным, а размерность не является
алгебраическим инвариантом (например, аддитивные группы $\RR$ и $\RR^2$ изоморфны как абстрактные группы --- как векторные
пространства над $\QQ$ одной и той же размерности континуум, что использует базис Гамеля). Рассуждение становится строгим, если
заметить, что изоморфизм $\Phi_F(g)=F\circ g\circ F^{-1}$ из Предложения 7.1(б) является гомеоморфизмом относительно топологии
локально равномерной сходимости (композиция и обращение непрерывны в этой топологии на группах автоморфизмов), а
топологическая размерность многообразия --- топологический инвариант (теорема Брауэра об инвариантности области). Приведённые
выше первое и второе доказательства полностью элементарны и этой тонкости не требуют.
\end{remarkbox}

\begin{remarkbox}{Замечание (связь с границей и теоремой Римана)}
Все три области $\DD$, $\CC$, $\Ch$ односвязны, а $\partial^\sharp\Ch=\varnothing$, $\partial^\sharp\CC=\{\infty\}$ (одна точка),
$\partial^\sharp\DD=\TT$ (более двух точек). По теореме Римана (Теорема 13.2, билеты 10--11) любая односвязная область $D\subset\Ch$,
граница которой содержит более одной точки, конформно эквивалентна $\DD$; если $\partial^\sharp D$ --- одна точка, то $D$ дробно-линейно
эквивалентна $\CC$ (точку границы отправить в $\infty$); если $\partial^\sharp D=\varnothing$, то $D=\Ch$. Теорема 7.5 показывает, что
эти три класса действительно различны: каждая односвязная область в $\Ch$ конформно эквивалентна \emph{ровно одной} из областей
$\DD$, $\CC$, $\Ch$.
\end{remarkbox}

\begin{remarkbox}{Коротко: план ответа на экзамене}
\begin{enumerate}
\item Определения: $\mathrm{Conf}(D_1,D_2)$, $\CAut(D)$ --- группа; $F\in\mathrm{Conf}(D_1,D_2)\Rightarrow g\mapsto FgF^{-1}$ ---
изоморфизм $\CAut(D_1)\cong\CAut(D_2)$.
\item $\CAut(\DD)=\{e^{i\alpha}\varphi_a\}$: $\varphi_a\in\CAut(\DD)$ (тождество для $1-|\varphi_a|^2$); для $f$ с $f(a)=0$ лемма Шварца к
$h=f\circ\varphi_{-a}$ и $h^{-1}$ $\Rightarrow$ $|h(z)|=|z|$ $\Rightarrow$ поворот. Параметры $(a,\alpha)$, размерность 3.
\item $\CAut(\Ch)=\Lambda$: $a=f^{-1}(\infty)$; $1/f$ однолистна около $a$ $\Rightarrow$ полюс первого порядка; вычитаем главную часть,
Лиувилль $\Rightarrow$ $f=\frac{w_0}{z-a}+w_1$ (или $w_1z+w_0$ при $a=\infty$). Размерность 6.
\item $\CAut(\CC)$: $\infty$ --- полюс (не устранимая --- Лиувилль, не существенная --- Сохоцкий + открытость + инъективность) $\Rightarrow$
$f^*\in\CAut(\Ch)$, $f^*(\infty)=\infty$ $\Rightarrow$ $f=az+b$. Размерность 4.
\item Неэквивалентность: $\Ch$ компактна, $\CC,\DD$ --- нет; $\CC\to\DD$ --- Лиувилль. Через группы: инволюции и коммутирование
(или размерности, понимаемые для топологических групп).
\item $\partial^\sharp$: пустая, одна точка, более двух точек; связь с теоремой Римана.
\end{enumerate}
\end{remarkbox}

% ===== part 2a =====
% ===================================================================
%  Часть 2a: билеты 8--11
% ===================================================================

% -------------------------------------------------------------------
\ticket{8}{Локальная равномерная ограниченность и локальная равностепенная непрерывность семейств функций. Связь этих понятий для семейств голоморфных функций}
\label{t8:begin}

\begin{defs}
\textbf{Область} $D\subset\CC$ --- открытое связное множество. $C(D)$ --- множество функций, непрерывных в $D$;
$H(D)$ --- множество функций, голоморфных в $D$.

\textbf{Компакт} $K\subset D$ --- замкнутое ограниченное подмножество $\CC$, лежащее в $D$. Для $f\in C(D)$ и компакта
$K\subset D$ положим $\|f\|_K:=\max_{z\in K}|f(z)|$ (равномерная норма на $K$; максимум достигается по теореме Вейерштрасса).

\textbf{Расстояние до границы.} Для компакта $K\subset D$ положим
$\dist(K,\partial D):=\inf\{|z-w|:\ z\in K,\ w\in\CC\setminus D\}$ (при $D=\CC$ считаем $\dist(K,\partial D)=+\infty$).
Так как $K$ компактен, а $\CC\setminus D$ замкнуто и не пересекается с $K$, то $\dist(K,\partial D)>0$.
(Инфимум по $w\in\CC\setminus D$ совпадает с инфимумом по $w\in\partial D$: отрезок от точки $z\in K$ до $w\notin D$
пересекает $\partial D$ в точке, не более далёкой от $z$, чем $w$.)

\textbf{Замкнутая $\rho$-окрестность компакта:} $K_\rho:=\{z\in\CC:\ \dist(z,K)\le\rho\}$. Это компакт (замкнуто как прообраз
$[0,\rho]$ при непрерывной функции $z\mapsto\dist(z,K)$, ограничено), и при $0<\rho<\dist(K,\partial D)$ имеем $K_\rho\subset D$:
если $z\in K_\rho$, то $|z-w|\le\rho$ для некоторой $w\in K$ (инфимум по компакту достигается), и при $z\notin D$ было бы
$\dist(K,\partial D)\le\rho$. Кроме того, $\clos{B(w,\sigma)}\subset K_{\rho+\sigma}$ для любой $w\in K_\rho$.

\textbf{Локально равномерная сходимость:} $f_n\to f$ локально равномерно в $D$, если $\|f_n-f\|_K\to0$ для каждого компакта
$K\subset D$ (эквивалентно: у каждой точки есть круг, на котором сходимость равномерна).

\textbf{Определение 13.2.} Пусть $\mathcal F\subset C(D)$ --- семейство функций.
\begin{enumerate}
\item[1)] $\mathcal F$ \emph{локально равномерно ограничено} в $D$, если для любого компакта $K\subset D$ найдётся $M=M(K)>0$ такое,
что $\|f\|_K\le M$ для \emph{всех} $f\in\mathcal F$.
\item[2)] $\mathcal F$ \emph{локально равностепенно непрерывно} в $D$, если для любого компакта $K\subset D$ и любого $\eps>0$ найдётся
$\delta=\delta(K,\eps)>0$ такое, что для любых $z_1,z_2\in K$ с $|z_1-z_2|<\delta$ и для \emph{любой} $f\in\mathcal F$
выполнено $|f(z_1)-f(z_2)|<\eps$.
\end{enumerate}
Смысл: в 1) константа $M$ одна на всё семейство, в 2) число $\delta$ одно на всё семейство (для одной функции это
просто равномерная непрерывность на $K$). Слово <<локально>> означает, что требования предъявляются лишь на компактах
внутри $D$, а не на всей $D$.

\textbf{Примеры.}
(a) $\{z^n\}_{n\in\NN}$ в $\DD$: $\|z^n\|_K\le1$, семейство локально равномерно ограничено (и, как мы увидим, локально
равностепенно непрерывно); но <<на всём $\DD$>> равностепенной непрерывности нет: для точек $z_1=1-\frac1n$,
$z_2=1-\frac1{n^2}$ имеем $|z_1-z_2|<\frac1n$, а $z_1^n=(1-\frac1n)^n\to e^{-1}$, $z_2^n=(1-\frac1{n^2})^n\to1$, так что
$|z_2^n-z_1^n|\to1-e^{-1}>0$. Точки уходят к границе $\DD$ --- поэтому в определении важна локальность.
(b) $\{nz\}_{n\in\NN}$ в $\CC$ не является локально равномерно ограниченным (на компакте $\{|z|\le1\}$ имеем $\|nz\|=n$).
(c) $\{f_n\equiv n\}$ --- локально равностепенно непрерывно (все $f_n(z_1)-f_n(z_2)=0$), но не ограничено.
(d) $f_n(x+iy)=\sin(nx)$ --- непрерывные (не голоморфные) функции, $|f_n|\le1$, но семейство не равностепенно непрерывно
на компакте $K=[0,1]$: точки $0$ и $\frac{\pi}{2n}$ лежат в $K$ (при $n\ge2$) и сколь угодно близки, а
$|f_n(\frac{\pi}{2n})-f_n(0)|=1$.
\end{defs}

\begin{remarkbox}{Что утверждается в билете}
Для произвольных непрерывных функций свойства 1) и 2) независимы (примеры (c), (d)). Для \emph{голоморфных}
функций 1) влечёт 2) (Предложение 13.8). Обратно, 2) вместе с ограниченностью семейства в одной точке влечёт 1) --- это
верно уже для непрерывных функций (Дополнение ниже).
\end{remarkbox}

\subsection{Предложение 13.8}

\begin{thmbox}{Предложение 13.8}
Пусть $D$ --- область в $\CC$, $\mathcal F\subset H(D)$. Если $\mathcal F$ локально равномерно ограничено в $D$, то $\mathcal F$
локально равностепенно непрерывно в $D$. Более того, семейство производных $\mathcal F'=\{f':f\in\mathcal F\}$ также локально
равномерно ограничено в $D$.
\end{thmbox}

Напомним используемый факт (интегральная формула Коши для производной): если $f\in H(D)$ и $\clos{B(w,\rho)}\subset D$, то
\[
f'(w)=\frac1{2\pi i}\int_{C(w,\rho)}\frac{f(z)\,dz}{(z-w)^2},\qquad C(w,\rho)=\{|z-w|=\rho\},
\]
откуда (\emph{неравенство Коши}) $|f'(w)|\le \dfrac1{2\pi}\cdot\dfrac{\max_{C(w,\rho)}|f|}{\rho^2}\cdot 2\pi\rho=\dfrac{\max_{C(w,\rho)}|f|}{\rho}$.

\begin{proof}
Зафиксируем компакт $K\subset D$ и положим $d:=\min\{1,\dist(K,\partial D)\}>0$. Тогда $d/2<\dist(K,\partial D)$,
так что $K_{d/4}\subset K_{d/2}\subset D$ --- компакты.

\emph{Шаг 1: оценка производных.} По локальной равномерной ограниченности найдётся $M>0$ такое, что $\|f\|_{K_{d/2}}\le M$
для всех $f\in\mathcal F$. Пусть $w\in K_{d/4}$. Тогда $\clos{B(w,d/4)}\subset K_{d/2}$ (если $|z-w|\le d/4$ и
$\dist(w,K)\le d/4$, то $\dist(z,K)\le d/2$). По неравенству Коши с $\rho=d/4$:
\[
|f'(w)|=\Bigl|\frac1{2\pi i}\int_{C(w,d/4)}\frac{f(z)\,dz}{(z-w)^2}\Bigr|\le\frac1{2\pi}\cdot\frac{M}{(d/4)^2}\cdot2\pi\frac d4=\frac{4M}{d}=:M'
\]
для всех $w\in K_{d/4}$ и всех $f\in\mathcal F$. В частности (так как $K\subset K_{d/4}$ и $K$ произволен) семейство
$\mathcal F'$ локально равномерно ограничено.

\emph{Шаг 2: равностепенная непрерывность.} Пусть $\eps>0$; положим $\delta:=\min\{d/4,\ \eps/M'\}$. Пусть $z_1,z_2\in K$,
$|z_1-z_2|<\delta$. Тогда отрезок $[z_1,z_2]$ лежит в $\clos{B(z_1,d/4)}\subset K_{d/4}$, а на нём $|f'|\le M'$. По формуле
Ньютона--Лейбница (у голоморфной $f$ производная $f'$ непрерывна, $f$ --- её первообразная)
\[
|f(z_1)-f(z_2)|=\Bigl|\int_{[z_2,z_1]}f'(z)\,dz\Bigr|\le M'\,|z_1-z_2|<M'\delta\le\eps
\]
для всех $f\in\mathcal F$. Число $\delta$ зависит только от $K$ и $\eps$ (через $d$ и $M=M(K_{d/2})$), но не от $f$.
\end{proof}

\begin{remarkbox}{Замечание (к записи на слайде)}
На слайде $K_\rho$ определено как $\bigcup_{z\in K}B(z,\rho)$; если круги открытые, это множество открыто и не является
компактом. Нужно брать замкнутую окрестность $K_\rho=\{\dist(z,K)\le\rho\}$ (или объединение замкнутых кругов, что то же
самое), как сделано выше.
\end{remarkbox}

\begin{remarkbox}{Где использована голоморфность}
Только в шаге 1: оценка $|f'|$ через $\max|f|$ на чуть большем компакте --- это специфически голоморфное явление
(неравенство Коши). Для непрерывных функций ограниченность не даёт контроля <<наклона>> (пример $\sin(nx)$).
\end{remarkbox}

\subsection{Обратная связь (дополнение)}

\begin{thmbox}{Предложение (дополнение к билету)}
Пусть $D$ --- область, $\mathcal F\subset C(D)$ --- локально равностепенно непрерывное семейство, и пусть
$\sup_{f\in\mathcal F}|f(z_0)|<\infty$ для некоторой точки $z_0\in D$. Тогда $\mathcal F$ локально равномерно ограничено в $D$.
\end{thmbox}

\begin{proof}
\emph{(i) Поточечная ограниченность.} Пусть $U:=\{z\in D:\ \sup_{f\in\mathcal F}|f(z)|<\infty\}$; $z_0\in U$. Пусть $z\in D$
и $\rho>0$ такое, что $\clos{B(z,\rho)}\subset D$. Применим равностепенную непрерывность на компакте $\clos{B(z,\rho)}$
с $\eps=1$: найдётся $\delta>0$ такое, что $|f(z')-f(z)|<1$ при $|z'-z|<\delta$, $z'\in\clos{B(z,\rho)}$, для всех $f$.
Значит, для $z'\in B(z,\min\{\delta,\rho\})$ числа $\sup_f|f(z')|$ и $\sup_f|f(z)|$ конечны или бесконечны одновременно.
Поэтому и $U$, и $D\setminus U$ открыты. Так как $D$ связна и $U\ne\varnothing$, то $U=D$.

\emph{(ii) Равномерная ограниченность на компакте.} Пусть $K\subset D$ --- компакт, $\delta=\delta(K,1)$ из определения
равностепенной непрерывности. Покроем $K$ конечным числом кругов $B(z_j,\delta)$, $z_j\in K$, $j=1,\dots,J$. Для $z\in K$
выберем $j$ с $|z-z_j|<\delta$; тогда $|f(z)|\le|f(z_j)|+1\le\max_j\sup_{g\in\mathcal F}|g(z_j)|+1$ для всех $f\in\mathcal F$.
Правая часть конечна по (i) и не зависит от $f$ и $z$.
\end{proof}

Итак, для семейств голоморфных функций в области:
\[
\text{лок. равн. ограниченность}\ \Longleftrightarrow\ \text{лок. равност. непрерывность}+\text{ограниченность в одной точке}.
\]
Импликация <<$\Rightarrow$>> --- Предложение 13.8; <<$\Leftarrow$>> --- дополнение (голоморфность не нужна).

\begin{remarkbox}{Коротко: план ответа на экзамене}
\begin{enumerate}
\item Определения 13.2 (одна константа $M$ / одно $\delta$ на всё семейство, на каждом компакте).
\item Примеры: $f_n\equiv n$ (непрерывность без ограниченности), $\sin(nx)$ (ограниченность без непрерывности).
\item Предложение 13.8: ограниченность $\Rightarrow$ равностепенная непрерывность для голоморфных.
\item Доказательство: $d=\min\{1,\dist(K,\partial D)\}$, компакты $K_{d/4}\subset K_{d/2}\subset D$, $\|f\|_{K_{d/2}}\le M$.
\item Неравенство Коши: $|f'|\le 4M/d$ на $K_{d/4}$.
\item Ньютон--Лейбниц на отрезке $[z_1,z_2]\subset K_{d/4}$ при $|z_1-z_2|<\delta=\min\{d/4,\eps d/(4M)\}$.
\item Обратно: равностепенная непрерывность + ограниченность в точке $\Rightarrow$ локальная ограниченность (связность $D$).
\item Значение: вместе с теоремой Арцела--Асколи даёт теорему Монтеля (билет 9).
\end{enumerate}
\end{remarkbox}

% -------------------------------------------------------------------
\ticket{9}{Понятие компактности семейства голоморфных функций. Теорема Монтеля. Непрерывный функционал на компактном семействе функций}

\begin{defs}
Всюду $D\subset\CC$ --- область, $\mathcal F\subset H(D)$. Локально равномерная ограниченность и локальная равностепенная
непрерывность --- как в билете 8 (Определение 13.2): для каждого компакта $K\subset D$ одна константа $M$ с $\|f\|_K\le M$
$\forall f\in\mathcal F$, соответственно одно $\delta=\delta(K,\eps)$ с $|f(z_1)-f(z_2)|<\eps$ при $z_{1,2}\in K$,
$|z_1-z_2|<\delta$, $\forall f\in\mathcal F$.

\textbf{Определение 13.3.}
1) $\mathcal F$ \emph{компактно в $D$} (в другой терминологии --- \emph{нормально}), если из любой последовательности
$(f_n)\subset\mathcal F$ можно выделить подпоследовательность, сходящуюся локально равномерно в $D$. Предельная функция
голоморфна (теорема Вейерштрасса), но не обязана лежать в $\mathcal F$.

2) $\mathcal F$ \emph{компактно в себе}, если оно компактно и предел любой локально равномерно сходящейся
подпоследовательности из 1) лежит в $\mathcal F$.

Пример: $\{f\in H(\DD):|f|\le1\}$ компактно в себе (по теореме Монтеля компактно, а неравенство $|f|\le1$ сохраняется при
поточечном пределе); $\{cz:\ 0<c<1\}$ компактно, но не компактно в себе ($cz\to z$ при $c\to1$).

\textbf{Определение 13.4.} Функционал $\Phi:\mathcal F\to\CC$ \emph{непрерывен на $\mathcal F$}, если для любой
последовательности $f_n\in\mathcal F$, сходящейся локально равномерно в $D$ к $f\in\mathcal F$, выполнено $\Phi(f_n)\to\Phi(f)$.

\textbf{Теорема Вейерштрасса} (прошлый семестр). Если $f_n\in H(D)$ и $f_n\to f$ локально равномерно в $D$, то $f\in H(D)$ и
$f_n^{(m)}\to f^{(m)}$ локально равномерно в $D$ при каждом $m\in\NN$.

\textbf{Всюду плотное множество} $E\subset D$: каждый непустой открытый круг $B\subset D$ содержит точку $E$. Пример:
$A=\{x+iy\in D:\ x,y\in\QQ\}$ --- счётно и всюду плотно в $D$ (на слайде обозначено $\QQ\cap D$, имеются в виду
<<рациональные>> точки $\QQ+i\QQ$).
\end{defs}

\subsection{Теорема Арцела--Асколи (нужная версия) и Лемма 13.1}

\begin{thmbox}{Лемма 13.1}
Пусть $\mathcal F\subset C(D)$ локально равностепенно непрерывно в $D$, $E\subset D$ всюду плотно, и последовательность
$(f_n)\subset\mathcal F$ сходится в каждой точке $E$. Тогда $(f_n)$ сходится локально равномерно в $D$ (к некоторой
непрерывной функции).
\end{thmbox}

\begin{proof}
Зафиксируем компакт $K\subset D$ и $\eps>0$. Положим $\rho:=\min\{1,\tfrac12\dist(K,\partial D)\}$, так что
$K_\rho=\{\dist(z,K)\le\rho\}$ --- компакт в $D$ (билет 8). По локальной равностепенной непрерывности (на компакте $K_\rho$)
найдётся $\delta>0$: $|f_n(z_1)-f_n(z_2)|<\eps$ при $z_1,z_2\in K_\rho$, $|z_1-z_2|<\delta$, для всех $n$.

Положим $r:=\min\{\delta/2,\rho\}$. Круги $B(w,r)$, $w\in K$, покрывают $K$; выберем конечное подпокрытие
$B_1=B(w_1,r),\dots,B_J=B(w_J,r)$. Так как $r\le\rho$, то $B_j\subset K_\rho$; так как $\diam B_j\le 2r\le\delta$, то
\[
|f_n(z_1)-f_n(z_2)|<\eps\qquad\text{при }z_1,z_2\in B_j,\ j=1,\dots,J,\ n\in\NN .
\]
В каждом $B_j$ (открытый круг в $D$) выберем $a_j\in E\cap B_j$ --- это возможно по плотности $E$. Числовые
последовательности $(f_n(a_j))_n$, $j=1,\dots,J$, сходятся, значит фундаментальны; их конечное число, поэтому найдётся
$N$ такое, что $|f_n(a_j)-f_m(a_j)|<\eps$ при $n,m>N$ и всех $j$.

Пусть $z\in K$; выберем $j$ с $z\in B_j$. Тогда при $n,m>N$
\[
|f_n(z)-f_m(z)|\le|f_n(z)-f_n(a_j)|+|f_n(a_j)-f_m(a_j)|+|f_m(a_j)-f_m(z)|<3\eps .
\]
Итак, $(f_n)$ равномерно фундаментальна на $K$ и по критерию Коши сходится равномерно на $K$. Предел $f$ определён в каждой
точке $D$ (каждая точка лежит в некотором компакте), не зависит от выбора $K$ и непрерывен на $D$ как равномерный предел
непрерывных функций на каждом замкнутом круге в $D$.
\end{proof}

\begin{remarkbox}{Замечание (к слайду)}
На слайде равностепенная непрерывность применяется к точкам $z_1,z_2\in K$, но точки $a_j\in E\cap B_j$ могут не лежать в $K$.
Поэтому её нужно применять на чуть большем компакте $K_\rho\supset B_j$, как сделано выше (именно для этого и введено $\rho$).
\end{remarkbox}

\begin{thmbox}{Теорема (Арцела--Асколи, версия для областей)}
Пусть $D\subset\CC$ --- область, $\mathcal F\subset C(D)$ локально равностепенно непрерывно и поточечно ограничено
на некотором счётном всюду плотном множестве $A\subset D$: $\sup_{f\in\mathcal F}|f(a)|<\infty$ для всех $a\in A$. Тогда
из любой последовательности $(f_n)\subset\mathcal F$ можно выделить подпоследовательность, сходящуюся локально равномерно в $D$.
\end{thmbox}

\begin{proof}
По Лемме 13.1 достаточно выделить подпоследовательность, сходящуюся в каждой точке $A=\{a_1,a_2,\dots\}$. Это делается
\emph{диагональным процессом}. Удобно следить за индексами: подпоследовательность задаётся строго возрастающим отображением
$N:\NN\to\NN$, $k\mapsto N(k)$; для такого $N$ всегда $N(k)\ge k$.

\emph{Шаг $p=1$.} Числовая последовательность $(f_n(a_1))_n$ ограничена, поэтому у неё есть сходящаяся
подпоследовательность (теорема Больцано--Вейерштрасса), т.е.\ найдётся строго возрастающее $N_1$ такое, что
последовательность $(f_{N_1(k)}(a_1))_k$ сходится.

\emph{Шаг $p\ge2$.} Пусть построено строго возрастающее $N_{p-1}$. Последовательность $(f_{N_{p-1}(k)}(a_p))_k$ ограничена;
выберем строго возрастающее $\sigma_p$ так, что $(f_{N_{p-1}(\sigma_p(k))}(a_p))_k$ сходится, и положим $N_p:=N_{p-1}\circ\sigma_p$.

По построению $(f_{N_p(k)})_k$ --- подпоследовательность $(f_{N_{p-1}(k)})_k$, и $(f_{N_p(k)}(a_l))_k$ сходится при всех $l\le p$
(подпоследовательность сходящейся последовательности сходится).

\emph{Диагональ.} Положим $m_n:=N_n(n)$. Тогда
$m_{n+1}=N_n(\sigma_{n+1}(n+1))\ge N_n(n+1)>N_n(n)=m_n$, т.е. $(f_{m_n})$ --- подпоследовательность $(f_n)$. Зафиксируем $p$.
При $n\ge p$ имеем $N_n=N_p\circ\sigma_{p+1}\circ\dots\circ\sigma_n$, так что каждое $m_n$ ($n\ge p$) есть значение $N_p$:
$m_n=N_p(k_n)$; поскольку $m_n$ строго возрастает, а $N_p$ строго возрастает, то и $k_n$ строго возрастает. Значит,
$(f_{m_n}(a_p))_{n\ge p}$ --- подпоследовательность сходящейся последовательности $(f_{N_p(k)}(a_p))_k$ и сходится.
Итак, $(f_{m_n})$ сходится в каждой точке $A$, и по Лемме 13.1 --- локально равномерно в $D$.
\end{proof}

\subsection{Теорема Монтеля}

\begin{thmbox}{Теорема 13.3 (теорема Монтеля)}
Пусть $D\subset\CC$ --- область, $\mathcal F\subset H(D)$. Если $\mathcal F$ локально равномерно ограничено в $D$, то
$\mathcal F$ компактно в $D$.
\end{thmbox}

\begin{proof}
По Предложению 13.8 (билет 8) $\mathcal F$ локально равностепенно непрерывно. Оно поточечно ограничено (каждая точка ---
компакт). Пусть $A=(\QQ+i\QQ)\cap D$ --- счётное всюду плотное множество. По теореме Арцела--Асколи (доказанной выше
через диагональный процесс и Лемму 13.1) из любой последовательности $(f_n)\subset\mathcal F$ выделяется
подпоследовательность, сходящаяся локально равномерно в $D$; её предел голоморфен по теореме Вейерштрасса.
\end{proof}

\begin{thmbox}{Дополнение (обращение теоремы Монтеля)}
Если $\mathcal F\subset H(D)$ компактно в $D$, то $\mathcal F$ локально равномерно ограничено. Таким образом,
компактность семейства голоморфных функций равносильна его локальной равномерной ограниченности.
\end{thmbox}
\begin{proof}
Пусть нет: для некоторого компакта $K$ найдутся $f_n\in\mathcal F$ с $\|f_n\|_K>n$. Выделим подпоследовательность
$f_{n_k}\to f$ равномерно на $K$; $f$ непрерывна на $K$, значит $\|f\|_K<\infty$, и $\|f_{n_k}\|_K\le\|f\|_K+1$ при больших $k$,
что противоречит $\|f_{n_k}\|_K>n_k\to\infty$.
\end{proof}

\begin{remarkbox}{Пример: класс $\Sc$ компактен в себе}
$\Sc=\{f\in\mathcal U(\DD):\ f(0)=0,\ f'(0)=1\}$ (однолистные нормированные функции, билеты 13--14). По второй теореме
растяжения (билет 14) $|f(z)|\le |z|/(1-|z|)^2$, так что $\Sc$ локально равномерно ограничено и по Монтелю компактно. Если
$f_n\in\Sc$, $f_n\to f$ локально равномерно, то $f(0)=0$, $f'(0)=\lim f_n'(0)=1$ (Вейерштрасс), так что $f\not\equiv\mathrm{const}$,
и по теореме о пределе однолистных функций (Предложение 13.6, билет 6) $f$ однолистна. Значит, $f\in\Sc$.
\end{remarkbox}

\subsection{Непрерывные функционалы}

\begin{thmbox}{Пример непрерывного функционала}
При любых $a\in D$ и $m\in\NN\cup\{0\}$ функционал $\Phi(f)=f^{(m)}(a)$ непрерывен на $H(D)$.
\end{thmbox}
\begin{proof}
Пусть $f_n\to f$ локально равномерно, $\clos{B(a,\rho)}\subset D$. По формуле Коши для производных
\[
|f_n^{(m)}(a)-f^{(m)}(a)|=\Bigl|\frac{m!}{2\pi i}\int_{C(a,\rho)}\frac{f_n(z)-f(z)}{(z-a)^{m+1}}\,dz\Bigr|
\le\frac{m!}{\rho^m}\,\|f_n-f\|_{C(a,\rho)}\to0. \qedhere
\]
\end{proof}

Другие примеры: $f\mapsto|f'(a)|$, $f\mapsto\Re f''(a)$, $f\mapsto\int_{C(a,\rho)}|f|\,|dz|$ --- непрерывны (композиции
непрерывных функций с предыдущими функционалами или прямая оценка). Функционал $f\mapsto\sup_D|f|$ на семействе ограниченных
функций, вообще говоря, не непрерывен относительно локально равномерной сходимости.

\begin{thmbox}{Лемма 13.2}
Пусть $\mathcal F\subset H(D)$ --- непустое компактное в себе семейство, $\Phi$ --- непрерывный функционал на $\mathcal F$.
Тогда $|\Phi|$ ограничен на $\mathcal F$ и существует $f_0\in\mathcal F$ такая, что $|\Phi(f)|\le|\Phi(f_0)|$ для всех
$f\in\mathcal F$.
\end{thmbox}

\begin{proof}
Пусть $M:=\sup\{|\Phi(f)|:f\in\mathcal F\}\in[0,+\infty]$. По определению супремума найдётся последовательность $f_n\in\mathcal F$
с $|\Phi(f_n)|\to M$ (если $M=+\infty$, берём $|\Phi(f_n)|>n$; если $M<\infty$ --- $|\Phi(f_n)|>M-\frac1n$). Так как
$\mathcal F$ компактно в себе, найдутся подпоследовательность $f_{n_k}$ и $f_0\in\mathcal F$ с $f_{n_k}\to f_0$ локально
равномерно. По непрерывности $\Phi(f_{n_k})\to\Phi(f_0)$, значит $|\Phi(f_{n_k})|\to|\Phi(f_0)|$. С другой стороны,
$|\Phi(f_{n_k})|\to M$. Поэтому $M=|\Phi(f_0)|<\infty$, и супремум достигается на $f_0$.
\end{proof}

\begin{remarkbox}{Почему нужна компактность \emph{в себе}}
На семействе $\{cz:0<c<1\}\subset H(\DD)$ (компактном, но не в себе) непрерывный функционал $\Phi(f)=f'(0)$ имеет
$\sup|\Phi|=1$, но не достигает его. В теореме Римана (билеты 10--11) именно поэтому семейство приходится выбирать
так, чтобы предельная функция не могла выродиться в константу.
\end{remarkbox}

\begin{remarkbox}{Коротко: план ответа на экзамене}
\begin{enumerate}
\item Определение 13.3: компактность (нормальность) и компактность в себе; примеры.
\item Теорема Монтеля: лок. равн. ограниченность $\Rightarrow$ компактность.
\item Предложение 13.8 (билет 8): ограниченность $\Rightarrow$ равностепенная непрерывность.
\item Лемма 13.1: равностепенная непрерывность + сходимость на плотном множестве $\Rightarrow$ локально равномерная
сходимость (конечное покрытие кругами диаметра $<\delta$, $3\eps$-рассуждение).
\item Диагональный процесс на счётном плотном множестве $(\QQ+i\QQ)\cap D$.
\item Предел голоморфен по Вейерштрассу; обратная теорема.
\item Определение 13.4, пример $f\mapsto f^{(m)}(a)$ (формула Коши).
\item Лемма 13.2: непрерывный функционал на компактном в себе семействе достигает $\max|\Phi|$.
\end{enumerate}
\end{remarkbox}

% -------------------------------------------------------------------
\ticket{10}{Теорема Римана. Схема доказательства и подробное проведение рассуждения о существовании однолистной функции, отображающей данную область в единичный круг}

\begin{defs}
\textbf{Область} --- открытое связное множество в $\CC$. Область $D\subset\CC$ \emph{односвязна}, если любой замкнутый путь в $D$
гомотопен в $D$ постоянному пути (для плоских областей это равносильно связности $\Ch\setminus D$). Пример: круг, полуплоскость,
плоскость с разрезом по лучу односвязны; кольцо и проколотая плоскость $\CC\setminus\{p\}$ --- нет.

\textbf{Однолистные функции.} $\mathcal U(D)$ --- множество функций $f\in H(D)$, инъективных в $D$ (\emph{однолистных}).
Если $f\in\mathcal U(D)$, то $f'(z)\ne0$ при всех $z\in D$ (критерий локальной однолистности, билет 1), $f(D)$ --- область
(принцип сохранения области) и $f^{-1}\in H(f(D))$. Отображение $f\in\mathcal U(D)$ с $f(D)=G$ называется
\emph{конформным отображением} $D$ на $G$ (конформным изоморфизмом), $f\in\operatorname{Conf}(D,G)$.

\textbf{Автоморфизмы круга} (Предложение 12.2): $\CAut(\DD)=\bigl\{e^{i\alpha}\psi_b:\ b\in\DD,\ \alpha\in\RR\bigr\}$, где
\[
\psi_b(w)=\frac{w-b}{1-\bar b w},\qquad \psi_b(b)=0,\quad \psi_b(0)=-b,\quad \psi_b^{-1}=\psi_{-b}.
\]

\textbf{Ветви корня} (Лемма 12.3, следствие Леммы 10.1). Если $D$ односвязна, $F\in H(D)$ и $F(z)\ne0$ в $D$, то существует
$h\in H(D)$ с $h^2=F$ (голоморфная ветвь $\sqrt{F}$); вторая ветвь --- $-h$.

\textbf{Принцип сохранения области (открытости).} Непостоянная голоморфная в области функция отображает открытые множества
в открытые. \textbf{Принцип максимума модуля:} если $F\in H(D)$ непостоянна, то $|F|$ не достигает в $D$ максимума;
в частности, из $|F|\le1$ в $D$ следует $|F|<1$ в $D$.

\textbf{Сходимость.} $f_n\to f$ \emph{локально равномерно} в $D$, если равномерно на каждом компакте $K\subset D$.
\emph{Теорема Вейерштрасса}: тогда $f\in H(D)$ и $f_n'\to f'$ локально равномерно.
\emph{Предложение 13.6} (следствие теоремы Гурвица, билет 6): если $f_n\in\mathcal U(D)$, $f_n\to f$ локально равномерно в
области $D$, то $f$ либо однолистна, либо постоянна.

\textbf{Компактность} (билет 9). Семейство $\mathcal F\subset H(D)$ \emph{компактно в себе}, если из любой последовательности
$f_n\in\mathcal F$ можно выделить подпоследовательность, сходящуюся локально равномерно к функции \emph{из} $\mathcal F$.
\emph{Теорема Монтеля (13.3)}: локально равномерно ограниченное семейство компактно (подпоследовательность сходится,
предел, вообще говоря, вне $\mathcal F$). Функционал $\Phi$ на $\mathcal F$ \emph{непрерывен}, если $\Phi(f_n)\to\Phi(f)$ при
$f_n\to f$ локально равномерно ($f_n,f\in\mathcal F$). \emph{Лемма 13.2}: непрерывный функционал на компактном в себе семействе
ограничен по модулю, и $|\Phi|$ достигает на $\mathcal F$ своей верхней грани.
\end{defs}

\subsection{Формулировка и обсуждение условий}

\begin{thmbox}{Теорема 13.2 (теорема Римана)}
Пусть $D$ --- односвязная область в $\CC$ такая, что $\partial D$ состоит более чем из одной точки. Тогда существует
однолистная голоморфная в $D$ функция $f$, отображающая $D$ конформно на единичный круг: $f(D)=\DD$.

Более того, для любой точки $a\in D$ и любого $\vartheta\in[0,2\pi)$ однолистная функция $f$ в $D$ такая, что $f(D)=\DD$,
$f(a)=0$ и $\arg f'(a)=\vartheta$, единственна.
\end{thmbox}

\begin{remarkbox}{Об условиях теоремы}
\textbf{1. Условие на границу необходимо.} Если $D=\CC$ (тогда $\partial D=\varnothing$), то конформного отображения на $\DD$
нет: функция $f\in H(\CC)$ с $|f|<1$ постоянна по теореме Лиувилля. Также нет отображения $\Ch$ на $\DD$ (компакт не
переходит в некомпактное множество при непрерывном отображении на).

\textbf{2. Для односвязной $D\subset\CC$ условие <<$\partial D$ содержит более одной точки>> равносильно $D\ne\CC$.}
Действительно, если $\partial D=\varnothing$, то $\clos D=D$, т.е.\ $D$ открыто и замкнуто в связном $\CC$, и $D=\CC$.
Пусть $\partial D=\{p\}$. Тогда $p\notin D$ (точки открытого $D$ --- внутренние, а не граничные), т.е.\
$D\subset\Omega:=\CC\setminus\{p\}$; $D$ открыто в $\Omega$ и замкнуто в $\Omega$, так как
$\clos D\cap\Omega=(D\cup\{p\})\cap\Omega=D$. Поскольку $\Omega$ связна, $D=\Omega=\CC\setminus\{p\}$ --- но эта область
не односвязна: для окружности $\gamma=C(p,1)$ имеем $\int_\gamma\frac{dz}{z-p}=2\pi i\ne0$, тогда как для стягиваемого
пути интеграл от голоморфной функции равен нулю (теорема Коши). Итак, у односвязной $D\ne\CC$ граница содержит не менее
двух точек.

\textbf{3. (К слайду 293.)} В тексте после теоремы сказано <<с нетривиальной границей, содержащей более двух точек>>;
достаточно (как и в самой Теореме 13.2) более \emph{одной} точки. Фактически, как видно из п.~2, нужно лишь $D\ne\CC$.
\end{remarkbox}

\begin{remarkbox}{Следствие: конформный радиус (подробнее --- билет 11)}
Пусть $f$ --- отображение из теоремы с $f(a)=0$, $\vartheta=0$, т.е.\ $f'(a)>0$. Тогда $f_a:=f/f'(a)$ --- единственная функция
из $\mathcal U(D)$ с $f_a(a)=0$, $f_a'(a)=1$, отображающая $D$ на круг $B(0,R)$; здесь $R=R(D,a)=1/f'(a)$ --- \emph{конформный
радиус} области $D$ в точке $a$.
\end{remarkbox}

\subsection{Схема доказательства}

\begin{enumerate}
\item[0.] \emph{Сведение к $\vartheta=0$.} Если $f$ --- нужное отображение с $f'(a)>0$, то $e^{i\vartheta}f$ --- конформное
отображение $D$ на $\DD$ (поворот --- автоморфизм $\DD$), $e^{i\vartheta}f(a)=0$, $\arg\bigl(e^{i\vartheta}f'(a)\bigr)=\vartheta$.
Обратно, если $f$ отвечает $\vartheta$, то $e^{-i\vartheta}f$ отвечает $0$; поэтому и единственность достаточно доказать при $\vartheta=0$.
\item \emph{Класс-кандидат.} $\mathcal F=\{g\in\mathcal U(D):\ |g(z)|\le1\ \forall z\in D,\ g(a)=0\}$. Искомое отображение
(если оно есть) лежит в $\mathcal F$. Главное --- доказать, что $\mathcal F\ne\varnothing$ (построить функцию $g_0$),
\textbf{это и есть существование однолистной функции, отображающей $D$ \emph{в} круг}.
\item \emph{Компактный подкласс.} $\mathcal F_0=\{f\in\mathcal F:\ f'(a)\ge|g_0'(a)|\}$ (неравенство подразумевает, что
$f'(a)$ вещественно). Он непуст и компактен в себе (Монтель + Вейерштрасс + Гурвиц).
\item \emph{Экстремальная задача.} Функционал $\Phi(h)=h'(a)$ непрерывен на $\mathcal F_0$; по Лемме 13.2 существует
$f\in\mathcal F_0$ с $f'(a)=\max_{h\in\mathcal F_0}h'(a)$.
\item \emph{Сюръективность} (билет 11). Если бы $f(D)\ne\DD$, то, взяв $b\in\DD\setminus f(D)$, с помощью автоморфизмов
$\psi_b,\psi_c$ и ветви квадратного корня построили бы $\varphi_3\in\mathcal F_0$ с $\varphi_3'(a)>f'(a)$ ---
противоречие с экстремальностью.
\item \emph{Единственность} (билет 11): если $f_1,f_2$ --- два таких отображения, то $f_1\circ f_2^{-1}$ --- автоморфизм
$\DD$, сохраняющий $0$, с положительной производной в $0$; по лемме Шварца это тождество.
\end{enumerate}

Ниже подробно проводятся шаги 1--3.

\subsection{Шаг 1: существование однолистной функции, отображающей $D$ в единичный круг}

\begin{thmbox}{Утверждение 1 ($\mathcal F\ne\varnothing$)}
В условиях Теоремы 13.2 для любой точки $a\in D$ существует $g_0\in\mathcal U(D)$ такая, что $g_0(a)=0$ и
$|g_0(z)|\le1$ при всех $z\in D$.
\end{thmbox}

\begin{proof}
\emph{(а) Две точки вне $D$.} По условию в $\partial D$ есть две различные точки $w^*\ne w^{**}$. Так как $D$ открыто,
$\partial D\cap D=\varnothing$, поэтому $w^*,w^{**}\in\CC\setminus D$.

\emph{(б) Ветви корней.} Функции $z-w^*$ и $z-w^{**}$ голоморфны и не обращаются в нуль в $D$ (ибо $w^*,w^{**}\notin D$).
Так как $D$ односвязна, по Лемме 12.3 существуют $F^*,F^{**}\in H(D)$ с
\[
(F^*(z))^2=z-w^*,\qquad (F^{**}(z))^2=z-w^{**},\qquad z\in D,
\]
причём $F^*,F^{**}$ нигде в $D$ не обращаются в нуль (иначе обратилось бы в нуль и их квадраты). Положим
\[
g_1:=\frac{F^*}{F^{**}}\in H(D),\qquad g_2:=-g_1,\qquad\text{так что}\qquad g_1(z)^2=g_2(z)^2=L(z):=\frac{z-w^*}{z-w^{**}} .
\]

\emph{(в) Инъективность дробно-линейной $L$.} Если $L(z_1)=L(z_2)$, $z_1,z_2\ne w^{**}$, то
$(z_1-w^*)(z_2-w^{**})=(z_2-w^*)(z_1-w^{**})$. Раскрывая скобки и сокращая $z_1z_2$ и $w^*w^{**}$, получаем
\[
-z_1w^{**}-z_2w^*=-z_2w^{**}-z_1w^*\iff (z_1-z_2)(w^*-w^{**})=0\iff z_1=z_2,
\]
так как $w^*\ne w^{**}$. Кроме того, $L(z)=0\iff z=w^*$.

\emph{(г) $g_1,g_2\in\mathcal U(D)$.} Если $g_1(z_1)=g_1(z_2)$, то, возводя в квадрат, $L(z_1)=L(z_2)$, откуда $z_1=z_2$ по (в).
Для $g_2=-g_1$: $g_2(z_1)=g_2(z_2)\iff g_1(z_1)=g_1(z_2)$, и всё сводится к $g_1$.

\emph{(д) $g_1(D)\cap g_2(D)=\varnothing$.} Пусть $g_1(z_1)=g_2(z_2)=-g_1(z_2)$. Возводя в квадрат, $L(z_1)=L(z_2)$, так что
$z_1=z_2=:z$ по (в). Тогда $g_1(z)=-g_1(z)$, т.е.\ $g_1(z)=0$, и $L(z)=g_1(z)^2=0$, откуда $z=w^*$ --- но $w^*\notin D$.
Противоречие.

\emph{(е) Круг, не пересекающийся с $g_1(D)$.} Функция $g_2$ однолистна, значит, непостоянна; по принципу сохранения области
$g_2(D)$ --- непустое открытое множество. Возьмём $w_0\in g_2(D)$ и $r>0$ так, что $B(w_0,r)\subset g_2(D)$. По (д)
$B(w_0,r)\cap g_1(D)=\varnothing$, т.е.
\[
|g_1(z)-w_0|\ge r\qquad\text{при всех }z\in D .
\]

\emph{(ж) Функция $g_0$.} Положим
\[
g_0(z)=\frac r2\Bigl(\frac1{g_1(z)-w_0}-\frac1{g_1(a)-w_0}\Bigr),\qquad z\in D .
\]
Знаменатели не обращаются в нуль по (е), поэтому $g_0\in H(D)$. Далее:
\begin{itemize}
\item $g_0=T\circ g_1$, где $T(w)=\frac r2\bigl(\frac1{w-w_0}-\frac1{g_1(a)-w_0}\bigr)$ --- дробно-линейная функция, инъективная
на $\CC\setminus\{w_0\}$ (если $T(w)=T(w')$, то $\frac1{w-w_0}=\frac1{w'-w_0}$ и $w=w'$). Композиция инъективных отображений
инъективна, так что $g_0\in\mathcal U(D)$;
\item $g_0(a)=0$ очевидно;
\item по (е) $\bigl|\frac1{g_1(z)-w_0}\bigr|\le\frac1r$ при всех $z\in D$, в том числе при $z=a$; значит,
\[
|g_0(z)|\le\frac r2\Bigl(\frac1r+\frac1r\Bigr)=1,\qquad z\in D .
\]
\end{itemize}
Итак, $g_0\in\mathcal F$.
\end{proof}

\begin{detail}{производная $g_0$ в точке $a$ не равна нулю}
Так как $g_0$ однолистна, $g_0'(a)\ne0$ (критерий локальной однолистности). Можно и явно:
$g_0'(a)=-\frac r2\,\frac{g_1'(a)}{(g_1(a)-w_0)^2}$, а $g_1'(a)\ne0$, поскольку из $g_1^2=L$ имеем
$2g_1g_1'=L'=\frac{w^*-w^{**}}{(z-w^{**})^2}\ne0$. Это нужно в шаге 2: число $|g_0'(a)|>0$ служит нижней границей.
\end{detail}

\begin{remarkbox}{Смысл построения}
Однолистная функция $g_1$ (ветвь $\sqrt{L}$) обладает тем свойством, что её образ не пересекается со своим <<отражением>>
$-g_1(D)$; следовательно, у образа $g_1(D)$ есть <<дыра>> --- круг $B(w_0,r)$. Инверсия $w\mapsto\frac1{w-w_0}$ переводит
внешность этого круга в ограниченное множество (круг радиуса $1/r$), после чего остаётся сдвинуть $a$ в $0$ и подобрать
масштаб. Подчеркнём, что ограниченность не даётся даром: например, для $D=\CC\setminus(-\infty,0]$ тождественное отображение
однолистно, но неограниченно. Можно обойтись и одной точкой
$w^*\notin D$: ветвь $h=\sqrt{z-w^*}$ однолистна и $h(D)\cap(-h(D))=\varnothing$ по тем же соображениям.
\end{remarkbox}

\subsection{Шаг 2: компактный в себе подкласс $\mathcal F_0$}

Положим
\[
\mathcal F_0=\bigl\{f\in\mathcal F:\ f'(a)\ \text{вещественно и}\ f'(a)\ge|g_0'(a)|\bigr\}.
\]

\begin{thmbox}{Утверждение 2}
$\mathcal F_0\ne\varnothing$, и семейство $\mathcal F_0$ компактно в себе в $D$.
\end{thmbox}

\begin{proof}
\emph{Непустота.} Пусть $\alpha=-\arg g_0'(a)$ (определено, так как $g_0'(a)\ne0$). Функция $e^{i\alpha}g_0$ однолистна,
$|e^{i\alpha}g_0|=|g_0|\le1$, $e^{i\alpha}g_0(a)=0$ и $(e^{i\alpha}g_0)'(a)=e^{i\alpha}g_0'(a)=|g_0'(a)|$ --- вещественное
число, $\ge|g_0'(a)|$. Значит, $e^{i\alpha}g_0\in\mathcal F_0$.

\emph{Компактность.} Все $f\in\mathcal F_0$ удовлетворяют $|f|\le1$ во всей $D$, так что $\mathcal F_0$ локально равномерно
ограничено (с $M(K)=1$ для любого компакта $K$). По теореме Монтеля $\mathcal F_0$ компактно в $D$.

\emph{Замкнутость (компактность в себе).} Пусть $f_n\in\mathcal F_0$ и $f_n\to f_0$ локально равномерно в $D$ (любая
подпоследовательность из определения компактности такова). Проверим, что $f_0\in\mathcal F_0$:
\begin{itemize}
\item $f_0\in H(D)$ и $f_n'\to f_0'$ локально равномерно (теорема Вейерштрасса);
\item $f_0(a)=\lim f_n(a)=0$; $|f_0(z)|=\lim|f_n(z)|\le1$ при всех $z\in D$;
\item $f_0'(a)=\lim_{n\to\infty}f_n'(a)$ --- предел вещественных чисел, не меньших $|g_0'(a)|$; поэтому $f_0'(a)$ вещественно и
$f_0'(a)\ge|g_0'(a)|>0$;
\item в частности, $f_0'(a)\ne0$, т.е.\ $f_0$ \emph{не постоянна}. По Предложению 13.6 (Гурвиц) предел однолистных
функций $f_n$ либо постоянен, либо однолистен; значит, $f_0\in\mathcal U(D)$.
\end{itemize}
Итак, $f_0\in\mathcal F$ и $f_0'(a)\ge|g_0'(a)|$, т.е.\ $f_0\in\mathcal F_0$. Кроме того, так как $f_0$ непостоянна и
$|f_0|\le1$, по принципу максимума $|f_0(z)|<1$ в $D$.
\end{proof}

\begin{remarkbox}{Зачем нужен $\mathcal F_0$, а не $\mathcal F$}
Семейство $\mathcal F$ компактно, но \emph{не} компактно в себе: $g_0/n\in\mathcal F$ и $g_0/n\to0$ равномерно, а
$0\notin\mathcal U(D)$. Однолистность теряется при предельном переходе ровно тогда, когда предел --- константа
(Предложение 13.6). Условие $f'(a)\ge|g_0'(a)|>0$ запрещает вырождение в константу; оно же не мешает экстремальной задаче,
поскольку мы ищем \emph{максимум} $f'(a)$, а $g_0$ (повёрнутая) уже даёт значение $|g_0'(a)|$.
\end{remarkbox}

\subsection{Шаг 3: существование экстремальной функции}

\begin{thmbox}{Утверждение 3}
Существует $f\in\mathcal F_0$ такая, что $f'(a)=\max_{h\in\mathcal F_0}h'(a)$. При этом $f'(a)\ge|h'(a)|$ для \emph{всех}
$h\in\mathcal F$.
\end{thmbox}

\begin{proof}
Функционал $\Phi(h)=h'(a)$ непрерывен на $\mathcal F_0$: если $h_n\to h$ локально равномерно, то по теореме Вейерштрасса
(или по формуле Коши $h_n'(a)-h'(a)=\frac1{2\pi i}\int_{C(a,\rho)}\frac{h_n(z)-h(z)}{(z-a)^2}dz$, где $\clos{B(a,\rho)}\subset D$)
\[
|h_n'(a)-h'(a)|\le\frac1\rho\max_{|z-a|=\rho}|h_n(z)-h(z)|\to0 .
\]
По Лемме 13.2 (применимой, так как $\mathcal F_0$ непусто и компактно в себе) найдётся $f\in\mathcal F_0$ с
$|\Phi(h)|\le|\Phi(f)|$ при всех $h\in\mathcal F_0$. На $\mathcal F_0$ значения $\Phi$ неотрицательны, так что
$h'(a)\le f'(a)$, т.е.\ $f'(a)=\max_{\mathcal F_0}\Phi$.

Пусть теперь $h\in\mathcal F$ произвольна. Если $|h'(a)|<|g_0'(a)|$, то $|h'(a)|<|g_0'(a)|\le f'(a)$. Если же
$|h'(a)|\ge|g_0'(a)|$ (в частности, $h'(a)\ne0$), то $\tilde h=e^{-i\arg h'(a)}h\in\mathcal F_0$ и $|h'(a)|=\tilde h'(a)\le f'(a)$.
\end{proof}

\begin{remarkbox}{Конечность максимума явно}
Лемма 13.2 сама гарантирует конечность $\sup$. Её можно увидеть и непосредственно: при $\clos{B(a,\rho)}\subset D$ из
неравенства Коши $|h'(a)|\le\frac1\rho\max_{|z-a|=\rho}|h|\le\frac1\rho$ для всех $h\in\mathcal F$; значит,
$f'(a)\le1/\dist(a,\partial D)$.
\end{remarkbox}

\subsection{Окончание доказательства (кратко)}
Остаётся (билет 11) доказать, что экстремальная функция $f$ отображает $D$ \emph{на} $\DD$: иначе при $b\in\DD\setminus f(D)$
функция $\varphi_3=e^{i\beta}\psi_c\circ\sqrt{\psi_b\circ f}$ (где $c$ --- значение корня в точке $a$) лежит в $\mathcal F_0$
и имеет $\varphi_3'(a)=\frac{1+|b|}{2\sqrt{|b|}}f'(a)>f'(a)$. Единственность следует из леммы Шварца,
применённой к $f_1\circ f_2^{-1}$ и $f_2\circ f_1^{-1}$.

\begin{remarkbox}{Коротко: план ответа на экзамене}
\begin{enumerate}
\item Формулировка Теоремы 13.2; необходимость условия ($D=\CC$ --- Лиувилль); для односвязной $D$ условие $\iff D\ne\CC$.
\item Сведение к $\vartheta=0$ поворотом; схема: $\mathcal F\to\mathcal F_0\to$ экстремаль $\to$ сюръективность $\to$
единственность.
\item $w^*\ne w^{**}\notin D$; ветви $F^*,F^{**}$ корней (Лемма 12.3, односвязность!); $g_1=F^*/F^{**}$, $g_1^2=L$ --- ДЛФ.
\item $g_1,\ g_2=-g_1$ однолистны, $g_1(D)\cap g_2(D)=\varnothing$ (возведение в квадрат + $L(z)=0\iff z=w^*$).
\item Открытость $g_2(D)\Rightarrow B(w_0,r)\subset g_2(D)$, $|g_1-w_0|\ge r$; $g_0=\frac r2\bigl(\frac1{g_1-w_0}-\frac1{g_1(a)-w_0}\bigr)$,
$|g_0|\le1$, $g_0(a)=0$.
\item $\mathcal F_0=\{f\in\mathcal F: f'(a)\ge|g_0'(a)|\}$; непуст ($e^{i\alpha}g_0$); Монтель $\Rightarrow$ компактно.
\item Предел: $f_0'(a)\ge|g_0'(a)|>0\Rightarrow$ не константа $\Rightarrow$ однолистна (Гурвиц, Предл.~13.6) $\Rightarrow f_0\in\mathcal F_0$.
\item $\Phi(h)=h'(a)$ непрерывен; Лемма 13.2 $\Rightarrow$ экстремаль $f$; далее --- билет 11.
\end{enumerate}
\end{remarkbox}

% -------------------------------------------------------------------
\ticket{11}{Теорема Римана. Схема доказательства и подробное обоснование того факта, что экстремальная функция $f_0$ \allowbreak функционала $f\mapsto f'(a)$ отображает данную область на весь единичный круг}

\begin{defs}
\textbf{Обозначения} (как в билете 10). $D\subset\CC$ --- односвязная область, $\partial D$ содержит более одной точки, $a\in D$;
$\mathcal U(D)$ --- однолистные (инъективные голоморфные) функции в $D$;
\[
\mathcal F=\{g\in\mathcal U(D):\ |g|\le1\text{ в }D,\ g(a)=0\},\qquad
\mathcal F_0=\{f\in\mathcal F:\ f'(a)\in\RR,\ f'(a)\ge|g_0'(a)|\},
\]
где $g_0\in\mathcal F$ --- функция, построенная в билете 10 (так что $\mathcal F\ne\varnothing$, $g_0'(a)\ne0$).
\textbf{Экстремальная функция} --- такая $f\in\mathcal F_0$, что $f'(a)=\max_{h\in\mathcal F_0}h'(a)$ (существует по Лемме 13.2,
так как $\mathcal F_0$ компактно в себе, а $h\mapsto h'(a)$ непрерывен; билет 10, шаги 2--3).
В билете она обозначена $f_0$; ниже, как на слайдах, пишем $f$.

\textbf{Автоморфизмы круга.} Для $b\in\DD$ положим $\psi_b(w)=\dfrac{w-b}{1-\bar bw}$. Все автоморфизмы $\DD$ имеют вид
$e^{i\alpha}\psi_b$ (Предложение 12.2). Свойства $\psi_b$ доказаны ниже (Лемма А).

\textbf{Лемма Шварца} (Предложение 12.1). Если $h\in H(\DD)$, $|h|\le1$, $h(0)=0$, то $|h(w)|\le|w|$ в $\DD$ и $|h'(0)|\le1$;
если хотя бы в одной точке $w_0\ne0$ выполнено $|h(w_0)|=|w_0|$ или $|h'(0)|=1$, то $h(w)=e^{i\alpha}w$.

\textbf{Ветви корня} (Лемма 12.3). В односвязной области у голоморфной функции без нулей есть голоморфная ветвь квадратного корня.

\textbf{Конформный радиус.} Если $f$ --- конформное отображение $D$ на $\DD$ с $f(a)=0$, $f'(a)>0$, то
$R(D,a):=1/f'(a)$ --- конформный радиус $D$ в точке $a$ (корректность --- в п.~11.5).
\end{defs}

\subsection{Схема доказательства теоремы Римана}

\begin{thmbox}{Теорема 13.2 (теорема Римана)}
Пусть $D\subset\CC$ --- односвязная область, $\partial D$ содержит более одной точки. Тогда существует $f\in\mathcal U(D)$ с
$f(D)=\DD$; при любых $a\in D$, $\vartheta\in[0,2\pi)$ такая $f$ с $f(a)=0$, $\arg f'(a)=\vartheta$ единственна.
\end{thmbox}

Схема (подробно шаги 1--3 --- билет 10):
\begin{enumerate}
\item Достаточно случая $\vartheta=0$ (домножение на $e^{i\vartheta}$).
\item $\mathcal F\ne\varnothing$: ветви $\sqrt{z-w^*}$, $\sqrt{z-w^{**}}$, $g_1=\sqrt{(z-w^*)/(z-w^{**})}$, $g_1(D)\cap(-g_1(D))=\varnothing$,
круг $B(w_0,r)\subset-g_1(D)$, $g_0=\frac r2\bigl(\frac1{g_1-w_0}-\frac1{g_1(a)-w_0}\bigr)\in\mathcal F$.
\item $\mathcal F_0$ непусто и компактно в себе (Монтель; предел имеет $f'(a)\ge|g_0'(a)|>0$, значит не константа и, по
Гурвицу, однолистен).
\item Функционал $h\mapsto h'(a)$ непрерывен, по Лемме 13.2 существует экстремаль $f\in\mathcal F_0$.
\item \textbf{Экстремаль сюръективна: $f(D)=\DD$} --- основной предмет билета (п.~11.3).
\item Единственность --- лемма Шварца (п.~11.4).
\end{enumerate}

\subsection{Вспомогательные факты об автоморфизмах круга}

\begin{thmbox}{Лемма А (свойства $\psi_b$)}
Пусть $b\in\DD$. Тогда
\begin{enumerate}
\item[(1)] $\psi_b$ голоморфна в круге $B(0,1/|b|)\supset\clos\DD$ (при $b=0$ --- во всей $\CC$), и
\[
1-|\psi_b(w)|^2=\frac{(1-|b|^2)(1-|w|^2)}{|1-\bar bw|^2};\qquad\text{в частности, }\psi_b(\DD)\subset\DD;
\]
\item[(2)] $\psi_b$ --- биекция $\DD$ на $\DD$, $\psi_b^{-1}=\psi_{-b}$; $\psi_b(w)=0\iff w=b$; $\psi_b(0)=-b$;
\item[(3)] производная:
\[
\psi_b'(w)=\frac{1-|b|^2}{(1-\bar bw)^2},\qquad\text{так что}\qquad \psi_b'(0)=1-|b|^2,\qquad \psi_b'(b)=\frac1{1-|b|^2}.
\]
\end{enumerate}
\end{thmbox}

\begin{proof}
(1) Знаменатель $1-\bar bw=0$ лишь при $w=1/\bar b$, $|1/\bar b|>1$. Далее
\begin{align*}
|1-\bar bw|^2-|w-b|^2&=\bigl(1-2\Re(\bar bw)+|b|^2|w|^2\bigr)-\bigl(|w|^2-2\Re(\bar b w)+|b|^2\bigr)\\
&=1-|w|^2-|b|^2+|b|^2|w|^2=(1-|b|^2)(1-|w|^2);
\end{align*}
делим на $|1-\bar bw|^2$. При $|w|<1$ правая часть положительна, т.е.\ $|\psi_b(w)|<1$.

(2) Прямая подстановка: при $u=\psi_b(w)$
\[
\psi_{-b}(u)=\frac{u+b}{1+\bar bu}=\frac{(w-b)+b(1-\bar bw)}{(1-\bar bw)+\bar b(w-b)}=\frac{w(1-|b|^2)}{1-|b|^2}=w,
\]
и аналогично $\psi_b\circ\psi_{-b}=\mathrm{id}$ (замена $b\to-b$). Так как обе функции переводят $\DD$ в $\DD$ (по (1)),
$\psi_b:\DD\to\DD$ --- биекция. Равенства $\psi_b(b)=0$, $\psi_b(0)=-b$ очевидны.

(3) По формуле производной частного
\[
\psi_b'(w)=\frac{(1-\bar bw)+\bar b(w-b)}{(1-\bar bw)^2}=\frac{1-|b|^2}{(1-\bar bw)^2};
\]
при $w=0$ получаем $1-|b|^2$, при $w=b$ получаем $\frac{1-|b|^2}{(1-|b|^2)^2}=\frac1{1-|b|^2}$.
\end{proof}

\subsection{Экстремальная функция отображает $D$ на весь круг}

\begin{thmbox}{Утверждение 4 (основное утверждение билета)}
Пусть $f\in\mathcal F_0$ --- экстремальная функция: $f'(a)\ge h'(a)$ для всех $h\in\mathcal F_0$. Тогда $f(D)=\DD$.
\end{thmbox}

\begin{proof}
\emph{Шаг 0: $f(D)\subset\DD$.} Функция $f$ однолистна, значит, непостоянна, и $|f|\le1$ в $D$. По принципу максимума модуля
$|f(z)|<1$ при всех $z\in D$. Итак, $f(D)\subset\DD$, и остаётся исключить случай $f(D)\ne\DD$.

Предположим противное: существует $b\in\DD\setminus f(D)$. Так как $f(a)=0\in f(D)$, то $b\ne0$, т.е.\ $0<|b|<1$.

\emph{Шаг 1: функция $\varphi_0=\psi_b\circ f$.} Положим $\varphi_0(z)=\psi_b(f(z))=\dfrac{f(z)-b}{1-\bar bf(z)}$, $z\in D$.
\begin{itemize}
\item $\varphi_0\in H(D)$: $f(D)\subset\DD$, а $\psi_b$ голоморфна на $\DD$ (Лемма А(1));
\item $\varphi_0\in\mathcal U(D)$: композиция инъективных $f$ и $\psi_b$ (Лемма А(2));
\item $\varphi_0(D)\subset\DD$ (Лемма А(1));
\item $\varphi_0(z)\ne0$ при всех $z\in D$: $\psi_b(w)=0\iff w=b$, а $b\notin f(D)$;
\item $\varphi_0(a)=\psi_b(0)=-b$, и по правилу дифференцирования композиции и Лемме А(3)
\[
\varphi_0'(a)=\psi_b'(f(a))\,f'(a)=\psi_b'(0)\,f'(a)=(1-|b|^2)\,f'(a).
\]
\end{itemize}

\emph{Шаг 2: квадратный корень $\varphi_1=\sqrt{\varphi_0}$.} Так как $D$ односвязна, а $\varphi_0\in H(D)$ не имеет нулей,
по Лемме 12.3 существует $\varphi_1\in H(D)$ с $\varphi_1^2=\varphi_0$. Свойства:
\begin{itemize}
\item $\varphi_1\in\mathcal U(D)$: если $\varphi_1(z_1)=\varphi_1(z_2)$, то $\varphi_0(z_1)=\varphi_1(z_1)^2=\varphi_1(z_2)^2=\varphi_0(z_2)$,
откуда $z_1=z_2$ (однолистность $\varphi_0$);
\item $\varphi_1(D)\subset\DD$: $|\varphi_1(z)|=|\varphi_0(z)|^{1/2}<1$;
\item $c:=\varphi_1(a)$ --- одно из значений $\sqrt{-b}$ (на слайде $c=S(-b)$): $c^2=\varphi_0(a)=-b$, так что
$|c|=\sqrt{|b|}\in(0,1)$ и $|c|^2=|b|$;
\item дифференцируя тождество $\varphi_1^2=\varphi_0$, получаем $2\varphi_1\varphi_1'=\varphi_0'$; в точке $a$ (где $\varphi_1(a)=c\ne0$)
\[
\varphi_1'(a)=\frac{\varphi_0'(a)}{2\varphi_1(a)}=\frac{f'(a)\,(1-|b|^2)}{2c}.
\]
\end{itemize}

\emph{Шаг 3: возвращаем ноль в точку $a$: $\varphi_2=\psi_c\circ\varphi_1$.} Положим
\[
\varphi_2(z)=\psi_c(\varphi_1(z))=\frac{\varphi_1(z)-c}{1-\bar c\,\varphi_1(z)},\qquad z\in D .
\]
Так как $c\in\DD$ и $\varphi_1(D)\subset\DD$, по Лемме А $\varphi_2\in H(D)$, $\varphi_2$ однолистна (композиция инъективных),
$\varphi_2(D)\subset\DD$ и $\varphi_2(a)=\psi_c(c)=0$. Производная (Лемма А(3): $\psi_c'(c)=1/(1-|c|^2)$, и $|c|^2=|b|$):
\[
\varphi_2'(a)=\psi_c'(c)\,\varphi_1'(a)=\frac{\varphi_1'(a)}{1-|c|^2}=\frac{f'(a)\,(1-|b|^2)}{2c\,(1-|b|)}
=\frac{f'(a)\,(1+|b|)}{2c},
\]
где использовано $1-|b|^2=(1-|b|)(1+|b|)$. В частности, $\varphi_2'(a)\ne0$ и
\[
|\varphi_2'(a)|=\frac{1+|b|}{2|c|}\,f'(a)=\frac{1+|b|}{2\sqrt{|b|}}\,f'(a).
\]

\emph{Шаг 4: ключевое неравенство.} При $0<|b|<1$
\[
\frac{1+|b|}{2\sqrt{|b|}}>1\iff 1+|b|-2\sqrt{|b|}>0\iff\bigl(1-\sqrt{|b|}\bigr)^2>0,
\]
что верно, так как $\sqrt{|b|}\ne1$. Учитывая $f'(a)>0$ (ибо $f'(a)\ge|g_0'(a)|>0$), получаем $|\varphi_2'(a)|>f'(a)$.

\emph{Шаг 5: поворот и противоречие.} Пусть $\beta=-\arg\varphi_2'(a)$ и $\varphi_3=e^{i\beta}\varphi_2$. Тогда
\begin{itemize}
\item $\varphi_3\in\mathcal U(D)$, $|\varphi_3|=|\varphi_2|<1$ в $D$, $\varphi_3(a)=0$, т.е.\ $\varphi_3\in\mathcal F$;
\item $\varphi_3'(a)=e^{i\beta}\varphi_2'(a)=|\varphi_2'(a)|=\dfrac{(1+|b|)\,f'(a)}{2\sqrt{|b|}}$ --- вещественное число, и
$\varphi_3'(a)>f'(a)\ge|g_0'(a)|$, так что $\varphi_3\in\mathcal F_0$.
\end{itemize}
Но тогда $\varphi_3'(a)>f'(a)=\max_{h\in\mathcal F_0}h'(a)$ --- противоречие с экстремальностью $f$. Значит, $\DD\setminus f(D)=\varnothing$,
т.е.\ $f(D)=\DD$, и $f$ --- конформное отображение $D$ на $\DD$ с $f(a)=0$, $f'(a)>0$.
\end{proof}

\begin{detail}{сводка явных формул (как на слайде 302)}
\[
\varphi_0=\psi_b\circ f,\quad \varphi_1=\sqrt{\varphi_0},\quad c=\varphi_1(a)=\sqrt{-b},\quad \varphi_2=\psi_c\circ\varphi_1,\quad
\varphi_3=e^{i\beta}\varphi_2;
\]
\begin{gather*}
\varphi_1'(a)=\frac{f'(a)(1-|b|^2)}{2c},\qquad
\varphi_3'(a)=\frac{f'(a)(1+|b|)}{2\sqrt{|b|}},\\
\varphi_2'(a)=\frac{\varphi_1'(a)}{1-|\varphi_1(a)|^2}=\frac{f'(a)(1-|b|^2)}{2c(1-|b|)}=\frac{f'(a)(1+|b|)}{2c}.
\end{gather*}
Формулы слайда верны; в них $S(-b)$ обозначает выбранное значение корня $c=\varphi_1(a)$.
\end{detail}

\begin{remarkbox}{Почему производная растёт: взгляд через лемму Шварца}
Выразим $f$ через $\varphi_2$. Из $\varphi_2=\psi_c\circ\varphi_1$ следует $\varphi_1=\psi_{-c}\circ\varphi_2$, затем
$\varphi_0=\varphi_1^2$ и $f=\psi_{-b}\circ\varphi_0$. Итак, $f=G\circ\varphi_2$, где
\[
G(w)=\psi_{-b}\bigl((\psi_{-c}(w))^2\bigr),\qquad G:\DD\to\DD,\qquad G(0)=\psi_{-b}(c^2)=\psi_{-b}(-b)=0 .
\]
Функция $G$ голоморфна в $\DD$, но \emph{не} однолистна (возведение в квадрат склеивает $\pm u$), поэтому не является
поворотом, и по лемме Шварца $|G'(0)|<1$. По Лемме А(3) $\psi_{-c}'(0)=1-|c|^2$, $\psi_{-b}'(-b)=\frac1{1-|b|^2}$, откуда
\[
G'(0)=\psi_{-b}'(-b)\cdot2\psi_{-c}(0)\cdot\psi_{-c}'(0)=\frac{2c(1-|b|)}{1-|b|^2}=\frac{2c}{1+|b|},\qquad |G'(0)|=\frac{2\sqrt{|b|}}{1+|b|}<1,
\]
и $f'(a)=G'(0)\varphi_2'(a)$ даёт $|\varphi_2'(a)|=f'(a)/|G'(0)|>f'(a)$ --- в согласии с шагами 3--4. Мораль: если
$f(D)\ne\DD$, то $f$ <<пропускается>> через неоднолистное сжимающее отображение круга и потому не может максимизировать
$|f'(a)|$.
\end{remarkbox}

\begin{remarkbox}{Где использованы условия теоремы}
Односвязность $D$ --- дважды: при построении $g_0$ (корни $\sqrt{z-w^*}$, билет 10) и при извлечении корня $\sqrt{\varphi_0}$
(шаг 2). Связность --- в теореме Гурвица и принципе максимума. Наличие двух граничных точек --- при построении $g_0$.
Однолистность $\varphi_0$ существенна: из неё следует однолистность корня $\varphi_1$.
\end{remarkbox}

\subsection{Единственность}

\begin{thmbox}{Утверждение 5 (единственность)}
Пусть $f_1,f_2\in\mathcal U(D)$, $f_1(D)=f_2(D)=\DD$, $f_1(a)=f_2(a)=0$ и $\arg f_1'(a)=\arg f_2'(a)=\vartheta$. Тогда $f_1\equiv f_2$.
\end{thmbox}

\begin{proof}
Обратные отображения $f_j^{-1}:\DD\to D$ голоморфны (обратная к однолистной голоморфной функции голоморфна).
Положим $h=f_1\circ f_2^{-1}$. Тогда $h\in H(\DD)$, $|h|<1$, $h(0)=f_1(f_2^{-1}(0))=f_1(a)=0$. По лемме Шварца $|h(w)|\le|w|$
при $w\in\DD$; подставляя $w=f_2(z)$, получаем $|f_1(z)|\le|f_2(z)|$ при всех $z\in D$. Применяя то же рассуждение к
$\tilde h=f_2\circ f_1^{-1}$, получаем $|f_2(z)|\le|f_1(z)|$. Значит, $|f_1|=|f_2|$ в $D$.

Рассмотрим $f^*=f_1/f_2$. Так как $f_2$ однолистна и $f_2(a)=0$, то $f_2(z)\ne0$ при $z\ne a$, и $f^*\in H(D\setminus\{a\})$.
В точке $a$ обе функции имеют нуль первого порядка ($f_j'(a)\ne0$ по однолистности): $f_j(z)=(z-a)u_j(z)$, $u_j\in H(D)$,
$u_j(a)=f_j'(a)\ne0$, $u_2\ne0$ всюду в $D$. Поэтому $f^*=u_1/u_2\in H(D)$ (особенность в $a$ устранима). При $z\ne a$
$|f^*(z)|=1$, по непрерывности и в точке $a$. Голоморфная функция с постоянным модулем постоянна (иначе образ $f^*(D)$ был
бы открыт по принципу сохранения области, а он лежит на окружности $\TT$). Итак, $f^*\equiv e^{i\sigma}$, $f_1=e^{i\sigma}f_2$.

Отсюда $f_1'(a)=e^{i\sigma}f_2'(a)$. Так как $f_1'(a)=|f_1'(a)|e^{i\vartheta}$ и $f_2'(a)=|f_2'(a)|e^{i\vartheta}$, то
$e^{i\sigma}=|f_1'(a)|/|f_2'(a)|>0$, а $|e^{i\sigma}|=1$; значит, $e^{i\sigma}=1$ и $f_1\equiv f_2$.
\end{proof}

\begin{remarkbox}{Другой вариант концовки}
Из $|h(w)|\le|w|$ и $|h^{-1}(w)|\le|w|$ (здесь $h^{-1}=\tilde h$) следует $|h(w)|=|w|$ при всех $w$, и по случаю равенства в
лемме Шварца $h(w)=e^{i\sigma}w$, т.е.\ $f_1=e^{i\sigma}f_2$; далее --- как выше.
\end{remarkbox}

\subsection{Конформный радиус}

\begin{thmbox}{Следствие (нормированное отображение и конформный радиус)}
Для любой $a\in D$ существует единственная функция $f_a\in\mathcal U(D)$ с $f_a(a)=0$, $f_a'(a)=1$, отображающая $D$ на
некоторый круг $B(0,R)$. При этом $f_a=f/f'(a)$, где $f$ --- отображение из теоремы Римана с $\vartheta=0$, и
$R=R(D,a)=1/f'(a)$ --- \emph{конформный радиус} $D$ в точке $a$.
\end{thmbox}

\begin{proof}
Существование: $f/f'(a)$ однолистна, переводит $a$ в $0$, имеет производную $1$ в $a$ и отображает $D$ на
$\frac1{f'(a)}\DD=B(0,1/f'(a))$. Единственность: если $g\in\mathcal U(D)$, $g(D)=B(0,R)$, $g(a)=0$, $g'(a)=1$, то $g/R$ ---
конформное отображение $D$ на $\DD$ с $(g/R)(a)=0$, $(g/R)'(a)=1/R>0$; по единственности в теореме Римана $g/R=f$, откуда
$1/R=f'(a)$ и $g=f/f'(a)$.
\end{proof}

\begin{remarkbox}{Свойства и примеры конформного радиуса}
\textbf{(1) Экстремальное свойство.}
\[
\frac1{R(D,a)}=f'(a)=\max\bigl\{|h'(a)|:\ h\in H(D),\ h(D)\subset\DD,\ h(a)=0\bigr\}
\]
(здесь $h$ не обязательно однолистны). Действительно, $h\circ f^{-1}:\DD\to\DD$, $0\mapsto0$, и по лемме Шварца
$|(h\circ f^{-1})'(0)|=|h'(a)|/f'(a)\le1$; максимум достигается при $h=f$. Это объясняет выбор функционала $h\mapsto h'(a)$.

\textbf{(2) Монотонность.} Если $a\in D_1\subset D_2$ (обе области удовлетворяют условиям теоремы), то $R(D_1,a)\le R(D_2,a)$:
отображение $f_2|_{D_1}$ голоморфно, переводит $D_1$ в $\DD$ и $a$ в $0$, поэтому по (1) $f_2'(a)\le f_1'(a)$.

\textbf{(3) Примеры.} Для круга $D=B(0,\rho)$ и $a\in D$ возьмём
\[
f(z)=\psi_{a/\rho}(z/\rho)=\frac{\rho(z-a)}{\rho^2-\bar az},\qquad
f'(a)=\frac{\rho}{\rho^2-|a|^2}>0,\qquad\text{так что}\qquad R(B(0,\rho),a)=\frac{\rho^2-|a|^2}{\rho};
\]
в центре $R=\rho$. Для полуплоскости $\Pi=\{\Im z>0\}$ и $a\in\Pi$:
\[
f(z)=\frac{z-a}{z-\bar a},\qquad f'(z)=\frac{a-\bar a}{(z-\bar a)^2},\qquad f'(a)=\frac{1}{a-\bar a}=\frac1{2i\Im a},
\]
поэтому $R(\Pi,a)=1/|f'(a)|=2\Im a$ (при любом конформном отображении $D$ на $\DD$ с $f(a)=0$ величина $|f'(a)|$
одна и та же, так как по единственности такие отображения отличаются поворотом).

\textbf{(4) Оценки через расстояние до границы.} Пусть $d=\dist(a,\partial D)$. Так как $B(a,d)\subset D$, по (2), (3)
$R(D,a)\ge R(B(a,d),a)=d$. Обратно, пусть $R=R(D,a)$ и $g=f_a^{-1}:B(0,R)\to D$, $g(0)=a$, $g'(0)=1$. Функция
$F(w)=\frac1R\bigl(g(Rw)-a\bigr)$ однолистна в $\DD$, $F(0)=0$, $F'(0)=g'(0)=1$, т.е.\ $F\in\Sc$. По теореме Кёбе об $1/4$
(билет 13) $F(\DD)\supset B(0,\frac14)$, поэтому $D=g(B(0,R))=a+R\,F(\DD)\supset B(a,\frac R4)$, откуда $d\ge R/4$. Итак,
$d\le R(D,a)\le4d$.
\end{remarkbox}

\begin{remarkbox}{Коротко: план ответа на экзамене}
\begin{enumerate}
\item Формулировка теоремы Римана, сведение к $\vartheta=0$; схема: $\mathcal F\ne\varnothing$ ($g_0$), $\mathcal F_0$ компактно
в себе, экстремаль $f$ (Монтель + Лемма 13.2).
\item $f(D)\subset\DD$ (принцип максимума); пусть $b\in\DD\setminus f(D)$, $b\ne0$.
\item $\varphi_0=\psi_b\circ f$: однолистна, без нулей, $\varphi_0(a)=-b$, $\varphi_0'(a)=(1-|b|^2)f'(a)$.
\item $\varphi_1=\sqrt{\varphi_0}$ (односвязность!): однолистна, $|\varphi_1|<1$, $c=\varphi_1(a)$, $|c|=\sqrt{|b|}$,
$\varphi_1'(a)=\frac{(1-|b|^2)f'(a)}{2c}$.
\item $\varphi_2=\psi_c\circ\varphi_1$: $\varphi_2(a)=0$, $\varphi_2'(a)=\frac{(1+|b|)f'(a)}{2c}$.
\item $\frac{1+|b|}{2\sqrt{|b|}}>1\iff(1-\sqrt{|b|})^2>0$; $\varphi_3=e^{i\beta}\varphi_2\in\mathcal F_0$, $\varphi_3'(a)>f'(a)$ --- противоречие.
\item Смысл: $f=G\circ\varphi_2$, $G$ --- неоднолистное отображение $\DD\to\DD$, $G(0)=0$, по Шварцу $|G'(0)|<1$.
\item Единственность: Шварц для $f_1\circ f_2^{-1}$ и $f_2\circ f_1^{-1}$ $\Rightarrow|f_1|=|f_2|$, $f_1/f_2\equiv e^{i\sigma}=1$.
\item Конформный радиус $R(D,a)=1/f'(a)$; примеры (круг, полуплоскость).
\end{enumerate}
\end{remarkbox}

% ===== part 2b =====
% ===================================================================
%  Часть 2b: билеты 12--15 и вопрос 41*
% ===================================================================

% -------------------------------------------------------------------
\ticket{12}{Теорема площадей и её следствие для функций класса $H(\DD\setminus\{0\})$, взаимно однозначных в единичном круге $\DD$}
\label{t12:begin}

\begin{defs}
\textbf{Однолистные функции.} Для области $G\subset\CC$ через $\mathcal U(G)$ обозначается множество функций $f\in H(G)$,
\emph{однолистных} (т.е. инъективных) в $G$. Аналогично говорят об однолистности мероморфной функции (значение $\infty$
тоже должно приниматься не более одного раза).
\emph{Пример:} $z\mapsto z^2$ однолистна в правой полуплоскости, но не в $\DD$; $e^z$ однолистна в полосе $|\Im z|<\pi$.

\textbf{Известный факт (прошлый семестр).} Если $f\in\mathcal U(G)$, то $f'(z)\ne0$ при всех $z\in G$
(если $f'(z_0)=0$, то $f-f(z_0)$ имеет в $z_0$ нуль кратности $m\ge2$, и по принципу аргумента близкие к $f(z_0)$ значения
принимаются вблизи $z_0$ ровно $m$ раз). Кроме того, $f(G)$ --- область, а $f^{-1}:f(G)\to G$ голоморфна.

\textbf{Внешность круга.} $\Delta:=\{z\in\CC:\ |z|>1\}=\CC\setminus\clos{\DD}$. На слайдах пишется $\mathcal U(\CC\setminus\DD)$;
имеется в виду однолистность во \emph{внешности} единичного круга, т.е. в $\Delta$.

\textbf{Класс $\Sigma$ (условие $(\Sigma)$).} Говорят, что $g\in\mathcal U(\Delta)$ \emph{удовлетворяет условию $(\Sigma)$}, если
её лорановское разложение в $\Delta$ (т.е. в окрестности $\infty$) имеет вид
\[
  g(z)=z+b_0+\sum_{k=1}^{\infty}b_kz^{-k},\qquad |z|>1, \tag{$\Sigma$}
\]
т.е. $g$ имеет в точке $\infty$ простой полюс с <<вычетом>> (коэффициентом при $z$) равным $1$. Ряд сходится абсолютно
и равномерно на каждой окружности $|z|=r>1$ (лорановский ряд в кольце $1<|z|<\infty$).
\emph{Примеры:} $g(z)=z$; $g(z)=z+1/z$ (функция Жуковского, однолистна в $\Delta$, см. ниже);
$g(z)=z+b/z$ при $|b|\le1$.

\textbf{Множество $K_g$.} $K_g:=\CC\setminus g(\Delta)$ --- <<дополнение образа>>. Для $g(z)=z$ имеем $K_g=\clos{\DD}$,
для $g(z)=z+1/z$ --- $K_g=[-2,2]$.

\textbf{Кривые и площадь.} \emph{Жорданова кривая} --- непрерывный инъективный образ окружности. По \emph{теореме Жордана}
дополнение $\CC\setminus\Gamma$ замкнутой жордановой кривой $\Gamma$ состоит ровно из двух компонент: ограниченной
$D(\Gamma)$ (внутренность) и неограниченной (внешность), и $\Gamma$ --- общая граница обеих. $m_2$ --- плоская мера Лебега (площадь).
\emph{Формула Грина:} если $\Gamma$ --- кусочно-гладкая жорданова кривая, обходимая положительно (т.е. $\Ind_\Gamma(w)=1$ при $w\in D(\Gamma)$), то
\[
 m_2(D(\Gamma))=\iint_{D(\Gamma)}dx\,dy=\frac12\oint_\Gamma(x\,dy-y\,dx)=\frac1{2i}\oint_\Gamma\bar\zeta\,d\zeta,
\]
так как $\bar\zeta\,d\zeta=(x-iy)(dx+i\,dy)=(x\,dx+y\,dy)+i(x\,dy-y\,dx)$, а $\oint_\Gamma(x\,dx+y\,dy)=\oint_\Gamma d\bigl(\tfrac{x^2+y^2}2\bigr)=0$.

\textbf{Класс $H(\DD\setminus\{0\})$} --- функции, голоморфные в проколотом круге. В следствии рассматриваются $F\in H(\DD\setminus\{0\})$
с простым полюсом в $0$ вида $F(z)=\frac1z+\sum_{n\ge0}\alpha_nz^n$, взаимно однозначные в $\DD$ (как мероморфные функции,
$F(0)=\infty$). \emph{Пример:} $F(z)=\frac1z+z$; $F(z)=\frac1z$.
\end{defs}

\subsection{Свойства множества $K_g$}

\begin{thmbox}{Утверждение (со слайда 304)}
Пусть $g\in\mathcal U(\Delta)$ удовлетворяет $(\Sigma)$. Тогда $g$ отображает $\Delta$ на $\CC\setminus K_g$, где $K_g$ --- непустой
связный компакт, не разделяющий плоскость (т.е. $\CC\setminus K_g$ связно).
\end{thmbox}

\begin{proof}
Доопределим $g(\infty)=\infty$. Функция $G(u):=1/g(1/u)$ в окрестности $u=0$ равна
$u/(1+b_0u+b_1u^2+\dots)$, т.е. голоморфна в $0$ с $G(0)=0$, $G'(0)=1$; значит, $g$ --- голоморфное (в смысле сферы Римана)
отображение $\Delta\cup\{\infty\}\to\Ch$, инъективное (значение $\infty$ принимается только в $\infty$, так как $g$ конечна в $\Delta$).

\emph{Компактность.} Образ $g(\Delta\cup\{\infty\})$ открыт в $\Ch$ (теорема об открытости отображения) и содержит $\infty$; поэтому он
содержит $\{|w|>R\}$ при некотором $R$, и $K_g=\Ch\setminus g(\Delta\cup\{\infty\})$ замкнуто и лежит в $\clos{B(0,R)}$, т.е.
компактно.

\emph{Не разделяет плоскость:} $\CC\setminus K_g=g(\Delta)$ связно как непрерывный образ связного множества.

\emph{Связность и непустота.} Положим $\Gamma_r=g(\{|z|=r\})$, $r>1$; ниже (в доказательстве теоремы площадей) показано, что
$K_g=\bigcap_{r>1}\clos{G_r}$, где $G_r=D(\Gamma_r)$, причём $\clos{G_{r}}\subset\clos{G_{r'}}$ при $r<r'$.
Убывающее пересечение непустых связных компактов $\clos{G_r}$ непусто и связно (стандартный факт: если бы $K=\bigcap K_r=A\sqcup B$
с непересекающимися непустыми компактами $A,B$, то взяв непересекающиеся открытые $U\supset A$, $V\supset B$, получили бы,
что компакты $K_r\setminus(U\cup V)$ убывают и имеют пустое пересечение, значит, пусты при $r$ близких к $1$, и тогда связный $K_r$
лежит в $U\sqcup V$ и пересекает оба --- противоречие). Непустота: если бы $K_g=\varnothing$, то пересечение компактов
$\clos{G_r}$ было бы пусто, т.е. $\clos{G_r}=\varnothing$ при некотором $r$, что абсурдно.
\end{proof}

\subsection{Теорема площадей}

\begin{thmbox}{Предложение 13.9 (теорема площадей)}
Пусть $g\in\mathcal U(\Delta)$ удовлетворяет условию $(\Sigma)$. Тогда
\[
 \sum_{k=1}^{\infty}k|b_k|^2\le1,
\]
причём равенство имеет место тогда и только тогда, когда $m_2(K_g)=0$. Более точно,
\[
 m_2(K_g)=\pi\Bigl(1-\sum_{k=1}^\infty k|b_k|^2\Bigr).
\]
\end{thmbox}

\begin{proof}
\textbf{Шаг 1: кривые $\Gamma_r$ и области $G_r$.} Пусть $r>1$, $C_r=\{|z|=r\}$, $\Gamma_r:=g(C_r)$. Так как $g$ непрерывна и
инъективна на $C_r$, $\Gamma_r$ --- замкнутая жорданова кривая; она даже аналитическая и гладкая: $t\mapsto g(re^{it})$ имеет
производную $ire^{it}g'(re^{it})\ne0$ (так как $g'\ne0$ для однолистной $g$). Пусть $G_r=D(\Gamma_r)$ --- её внутренность,
$E_r$ --- внешность. Обозначим $A_r=\{1<|z|<r\}$, $\Delta_r=\{|z|>r\}$.

\begin{detail}{$g(\Delta_r)=E_r$ и $G_r=K_g\sqcup g(A_r)$}
Множество $g(\Delta_r)$ связно, не пересекается с $\Gamma_r$ (инъективность) и неограниченно ($g(z)\to\infty$ при $z\to\infty$), поэтому
$g(\Delta_r)$ лежит в одной компоненте $\CC\setminus\Gamma_r$, причём не в ограниченной, т.е. $g(\Delta_r)\subset E_r$.
Далее, $g(\Delta_r)$ открыто (открытость голоморфного отображения) и замкнуто в $E_r$: если $g(z_n)\to w\in E_r$, $|z_n|>r$, то
последовательность $z_n$ ограничена (иначе $|g(z_{n_j})|\to\infty$), поэтому по подпоследовательности $z_n\to z$, $|z|\ge r$,
и $w=g(z)$; если бы $|z|=r$, то $w\in\Gamma_r$, а $w\in E_r$. Значит, $|z|>r$, $w\in g(\Delta_r)$. Так как $E_r$ связно, а
$g(\Delta_r)\ne\varnothing$, то $g(\Delta_r)=E_r$.

Итак, $\CC=G_r\sqcup\Gamma_r\sqcup E_r$, где $\Gamma_r\sqcup E_r=g(\{|z|\ge r\})$. С другой стороны, $\CC=K_g\sqcup g(A_r)\sqcup g(\{|z|\ge r\})$
(по определению $K_g$ и инъективности $g$). Сравнивая, получаем $G_r=K_g\sqcup g(A_r)$. В частности, $K_g\subset G_r$ при каждом $r>1$ и
$\clos{G_r}=G_r\cup\Gamma_r=K_g\sqcup g(\{1<|z|\le r\})$, откуда $\bigcap_{r>1}\clos{G_r}=K_g$ и монотонность $\clos{G_r}$ по $r$.
\end{detail}

\begin{detail}{Ориентация: кривая $\Gamma_r$ (образ положительно ориентированной $C_r$) обходит $G_r$ положительно}
Возьмём точку $w_0\in G_r$ вида $w_0=g(z_0)$, $1<|z_0|<r$. Тогда
\[
 \Ind_{\Gamma_r}(w_0)=\frac1{2\pi i}\oint_{\Gamma_r}\frac{d\zeta}{\zeta-w_0}=\frac1{2\pi i}\oint_{|z|=r}\frac{g'(z)\,dz}{g(z)-w_0}.
\]
Функция $g'(z)/(g(z)-w_0)$ голоморфна в $\{|z|>r'\}$ при $|z_0|<r'<r$ (знаменатель обращается в нуль только в $z_0$), поэтому
интеграл по $|z|=r$ равен интегралу по $|z|=R$ для любого $R>r$ (теорема Коши для кольца). При $z\to\infty$
\[
 \frac{g'(z)}{g(z)-w_0}=\frac{1+O(z^{-2})}{z\,(1+(b_0-w_0)z^{-1}+O(z^{-2}))}=\frac1z+O(z^{-2}),
\]
так что интеграл по $|z|=R$ стремится к $2\pi i$ при $R\to\infty$, а он от $R$ не зависит. Итак, $\Ind_{\Gamma_r}(w_0)=1$; индекс
постоянен на связной $G_r$, значит, равен $1$ на всей $G_r$, т.е. обход положительный, и формула Грина применима со знаком $+$.
\end{detail}

\textbf{Шаг 2: вычисление площади $G_r$.} По формуле Грина, делая замену $\zeta=g(z)$, $z=re^{it}$, $dz=iz\,dt$:
\[
 m_2(G_r)=\frac1{2i}\oint_{\Gamma_r}\bar\zeta\,d\zeta=\frac1{2i}\oint_{|z|=r}\overline{g(z)}\,g'(z)\,dz
 =\frac12\int_0^{2\pi}\overline{g(re^{it})}\cdot\bigl(z g'(z)\bigr)\big|_{z=re^{it}}\,dt.
\]
Разложим оба множителя в ряды Фурье по $e^{ijt}$. Так как $g'(z)=1-\sum_{k\ge1}kb_kz^{-k-1}$, то
\[
 \overline{g(re^{it})}=re^{-it}+\bar b_0+\sum_{k=1}^\infty \bar b_k r^{-k}e^{ikt},\qquad
 zg'(z)\big|_{z=re^{it}}=re^{it}-\sum_{m=1}^\infty m b_m r^{-m}e^{-imt}.
\]
Оба ряда сходятся абсолютно: $\sum_k |b_k|r^{-k}<\infty$ и $\sum_m m|b_m|r^{-m}<\infty$, поскольку лорановские ряды для $g$ и $g'$
абсолютно сходятся при $|z|=r>1$ (можно взять $1<r_1<r$, тогда $|b_k|\le Mr_1^{k}$ по неравенствам Коши, и ряды
мажорируются геометрическими). Поэтому произведение рядов --- абсолютно сходящийся двойной ряд, его можно почленно интегрировать.
Используем $\int_0^{2\pi}e^{i(j+m)t}\,dt=2\pi$ при $j+m=0$ и $0$ иначе. Выживают лишь пары показателей с суммой $0$:
\begin{itemize}
\item $re^{-it}\cdot re^{it}$: даёт $r^2$;
\item $\bar b_k r^{-k}e^{ikt}\cdot\bigl(-kb_kr^{-k}e^{-ikt}\bigr)$, $k\ge1$: даёт $-k|b_k|^2r^{-2k}$;
\item член $\bar b_0$ (показатель $0$) в пару не входит: во втором ряду нет постоянного члена; прочие произведения имеют
ненулевую сумму показателей.
\end{itemize}
Итак,
\[
 m_2(G_r)=\frac12\cdot2\pi\Bigl(r^2-\sum_{k=1}^\infty k|b_k|^2r^{-2k}\Bigr)=\pi\Bigl(r^2-\sum_{k=1}^\infty \frac{k|b_k|^2}{r^{2k}}\Bigr).
 \tag{12.1}\label{t12:area}
\]

\textbf{Шаг 3: предельный переход $r\to1+$.} Так как $m_2(G_r)\ge0$, то при любом $N$
\[
 \sum_{k=1}^N\frac{k|b_k|^2}{r^{2k}}\le\sum_{k=1}^\infty\frac{k|b_k|^2}{r^{2k}}\le r^2 .
\]
Устремляя $r\to1+$ (конечная сумма непрерывна по $r$), получаем $\sum_{k=1}^N k|b_k|^2\le1$ при всех $N$, т.е.
$\sum_{k=1}^\infty k|b_k|^2\le1$.

Для формулы с $m_2(K_g)$: при $r\downarrow1$ члены $k|b_k|^2r^{-2k}$ монотонно возрастают к $k|b_k|^2$, поэтому по теореме
о монотонной сходимости для рядов (теорема Беппо Леви) $\sum_k k|b_k|^2r^{-2k}\to\sum_kk|b_k|^2$, и из \eqref{t12:area}
\[
 m_2(G_r)\xrightarrow[r\to1+]{}\pi\Bigl(1-\sum_{k=1}^\infty k|b_k|^2\Bigr).
\]
С другой стороны, по детали шага 1, $G_r=K_g\sqcup g(A_r)$, т.е. $m_2(G_r)=m_2(K_g)+m_2(g(A_r))$. Множества $g(A_r)$ убывают
при $r\downarrow1$, имеют конечную меру ($\le m_2(G_r)<\infty$) и пустое пересечение ($\bigcap_{r>1}A_r=\varnothing$). По непрерывности
меры сверху $m_2(g(A_r))\to0$. Следовательно, $m_2(G_r)\to m_2(K_g)$, и
\[
 m_2(K_g)=\pi\Bigl(1-\sum_{k=1}^\infty k|b_k|^2\Bigr).
\]
Отсюда равенство $\sum k|b_k|^2=1$ равносильно $m_2(K_g)=0$.
\end{proof}

\begin{remarkbox}{Замечание (к слайду 305)}
На слайде доказательство заканчивается фразой <<Так как $m_2(K)>0$, то отсюда вытекает требуемое утверждение>>. Строгого
неравенства здесь быть не должно: $m_2(K_g)$ может равняться нулю (например, $K_g=[-2,2]$ для $g(z)=z+1/z$). Неравенство
$\sum k|b_k|^2\le1$ следует из $m_2(G_r)\ge0$, а критерий равенства --- из точной формулы для $m_2(K_g)$, полученной выше.
\end{remarkbox}

\begin{remarkbox}{Смысл теоремы}
Площадь образа кольца $A_r$ равна $\iint_{A_r}|g'|^2\,dm_2=\pi\bigl[(r^2-1)+\sum k|b_k|^2(1-r^{-2k})\bigr]$, а теорема утверждает, что
<<дырка>> $K_g$ имеет площадь $\pi(1-\sum k|b_k|^2)\ge0$: функция класса $\Sigma$ не может <<растянуть>> внешность круга сильнее,
чем позволяет площадь $\pi$ единичного круга.
\end{remarkbox}

\subsection{Следствия теоремы площадей}

\begin{thmbox}{Следствие 1}
Пусть $g$ такая, как в теореме площадей. Тогда $|b_1|\le1$, причём равенство достигается только на функциях вида
$g(z)=z+b_0+b_1/z$, где $|b_1|=1$. Такая функция отображает $\Delta$ на внешность отрезка длины $4$.
\end{thmbox}

\begin{proof}
$|b_1|^2=1\cdot|b_1|^2\le\sum_k k|b_k|^2\le1$. Если $|b_1|=1$, то $\sum_{k\ge2}k|b_k|^2\le1-|b_1|^2=0$, т.е. $b_k=0$ при $k\ge2$,
и $g(z)=z+b_0+b_1/z$.
\end{proof}

\begin{detail}{Упражнение со слайда: $z+b_0+b_1/z$ при $|b_1|=1$ отображает $\Delta$ на внешность отрезка длины $4$}
Пусть $b_1=e^{i\alpha}$. Положим $z=e^{i\alpha/2}u$; тогда $|z|>1\iff|u|>1$ и
\[
 g(z)=e^{i\alpha/2}u+b_0+\frac{e^{i\alpha}}{e^{i\alpha/2}u}=b_0+e^{i\alpha/2}\Bigl(u+\frac1u\Bigr).
\]
Значит, достаточно доказать, что $J(u)=u+\frac1u$ (функция Жуковского) взаимно однозначно отображает $\{|u|>1\}$ на $\CC\setminus[-2,2]$.
Для $w\in\CC$ уравнение $u+1/u=w$ равносильно $u^2-wu+1=0$; у него два корня $u_1,u_2$ (с учётом кратности), $u_1u_2=1$.
Если $|u|=1$, $u=e^{it}$, то $J(u)=2\cos t\in[-2,2]$. Поэтому:
\begin{itemize}
\item если $w\notin[-2,2]$, то ни один корень не лежит на $\TT$, а так как $|u_1||u_2|=1$, ровно один корень лежит в $\{|u|>1\}$;
\item если $w=2\cos t\in[-2,2]$, то корни --- $e^{\pm it}$ (их произведение $1$, сумма $2\cos t$), оба на $\TT$, и в $\{|u|>1\}$ корней нет.
\end{itemize}
Итак, $J:\{|u|>1\}\to\CC\setminus[-2,2]$ --- биекция. Следовательно, $g(\Delta)=\CC\setminus I$, где
$I=b_0+e^{i\alpha/2}[-2,2]$ --- отрезок с концами $b_0\pm2e^{i\alpha/2}$, длины $4$. (Здесь $K_g=I$ и $m_2(K_g)=0$ --- в согласии с
критерием равенства.)
\end{detail}

\begin{thmbox}{Следствие 2}
Пусть $g$ такая, как в теореме площадей. Тогда $|b_n|\le n^{-1/2}$ при $n\ge1$. При $n>1$ эта оценка не точна.
\end{thmbox}

\begin{proof}
$n|b_n|^2\le\sum_k k|b_k|^2\le1$. Неточность: если бы $|b_n|=n^{-1/2}$, то, как в следствии 1, все остальные $b_k$ ($k\ge1$, $k\ne n$)
равнялись бы нулю, и $g(z)=z+b_0+b_nz^{-n}$. Но тогда $g'(z)=1-nb_nz^{-n-1}$ обращается в нуль при $z^{n+1}=nb_n$, т.е. при
$|z|=(n|b_n|)^{1/(n+1)}=n^{1/(2(n+1))}>1$ (так как $n>1$), т.е. в точке из $\Delta$, что невозможно для однолистной $g$.
Пример слайда: $z+n^{-1/2}z^{-n}$ --- производная имеет нули в $\Delta$, функция не однолистна.
\end{proof}

\begin{thmbox}{Следствие 3 (для функций класса $H(\DD\setminus\{0\})$)}
Пусть $F\in H(\DD\setminus\{0\})$ взаимно однозначна в $\DD$ (как мероморфная функция) и
\[
 F(z)=\frac1z+\sum_{n=0}^\infty\alpha_nz^n,\qquad 0<|z|<1.
\]
Тогда $\displaystyle\sum_{n=1}^\infty n|\alpha_n|^2\le1$. В частности, $|\alpha_1|\le1$.
\end{thmbox}

\begin{proof}
Положим $g(z):=F(1/z)$ при $|z|>1$. Отображение $z\mapsto1/z$ --- биекция $\Delta\to\DD\setminus\{0\}$, а $F$ инъективна на
$\DD\setminus\{0\}$, так что $g\in\mathcal U(\Delta)$, причём
\[
 g(z)=z+\alpha_0+\sum_{n=1}^\infty\alpha_nz^{-n},\qquad |z|>1,
\]
т.е. $g$ удовлетворяет $(\Sigma)$ с $b_0=\alpha_0$, $b_n=\alpha_n$. По теореме площадей $\sum n|\alpha_n|^2\le1$.
\end{proof}

\begin{remarkbox}{Как это используется дальше}
Следствие 3 применяется (билет 13) к $F=1/h$, где $h(z)=\sqrt{f(z^2)}$, $f\in\Sc$; из $|\alpha_1|\le1$ получается оценка Бибербаха $|a_2|\le2$.
Обратите внимание: $F$ однолистна в $\DD$ тогда и только тогда, когда $1/F$ однолистна в $\DD$; так что следствие 3 --- это утверждение
о функциях $h=1/F\in\mathcal U(\DD)$ с $h(0)=0$, $h'(0)=1$, не обращающихся в нуль вне нуля.
\end{remarkbox}

\begin{remarkbox}{Коротко: план ответа на экзамене}
\begin{enumerate}
\item Класс $(\Sigma)$: $g\in\mathcal U(\Delta)$, $g(z)=z+b_0+\sum b_kz^{-k}$; $K_g=\CC\setminus g(\Delta)$ --- связный компакт.
\item $\Gamma_r=g(\{|z|=r\})$ --- гладкая жорданова кривая, $G_r=D(\Gamma_r)\supset K_g$, обход положительный.
\item Формула Грина: $m_2(G_r)=\frac1{2i}\oint_{\Gamma_r}\bar\zeta\,d\zeta=\frac12\int_0^{2\pi}\overline{g}\,zg'\,dt$.
\item Перемножение рядов Фурье: $m_2(G_r)=\pi(r^2-\sum k|b_k|^2r^{-2k})\ge0$.
\item $r\to1$: $\sum k|b_k|^2\le1$; точнее $m_2(K_g)=\pi(1-\sum k|b_k|^2)$, отсюда критерий равенства.
\item Следствия: $|b_1|\le1$ (равенство --- отрезок длины 4, функция Жуковского), $|b_n|\le n^{-1/2}$.
\item Для $F(z)=1/z+\sum\alpha_nz^n$, однолистной в $\DD$: $g(z)=F(1/z)\in(\Sigma)$, значит $\sum n|\alpha_n|^2\le1$.
\end{enumerate}
\end{remarkbox}

% -------------------------------------------------------------------
\ticket{13}{Теорема Бибербаха (оценка $|a_2|\le2$ для функций класса $\Sc$) и теорема Кёбе (об $1/4$)}
\label{t13:begin}

\begin{defs}
\textbf{Однолистные функции} $\mathcal U(G)$ --- голоморфные инъективные функции в области $G$; для них $f'\ne0$ в $G$ и
$f^{-1}:f(G)\to G$ голоморфна (см. билет 12).

\textbf{Класс $\Sc$:} $\Sc=\{f\in\mathcal U(\DD):\ f(0)=0,\ f'(0)=1\}$, т.е.
\[
 f(z)=z+\sum_{k=2}^\infty a_kz^k,\qquad |z|<1 .
\]
\emph{Нормировка не ограничивает общности:} если $f\in\mathcal U(\DD)$, то $\frac{f(z)-f(0)}{f'(0)}\in\Sc$ ($f'(0)\ne0$).
\emph{Примеры:} $z$; $\frac{z}{1-z}$ (дробно-линейная, $\DD\to\{\Re w>-\tfrac12\}$); $z-\frac{z^2}{2}$ (производная $1-z\ne0$ в $\DD$,
и $f(z_1)-f(z_2)=(z_1-z_2)\bigl(1-\frac{z_1+z_2}2\bigr)\ne0$ при $z_1\ne z_2$).

\textbf{Функция Кёбе:}
\[
 k(z)=\frac{z}{(1-z)^2}=\sum_{n=1}^\infty nz^n=\frac14\Bigl(\frac{1+z}{1-z}\Bigr)^2-\frac14 .
\]
Её \emph{вращения}: $k_\vartheta(z)=e^{-i\vartheta}k(e^{i\vartheta}z)=\frac{z}{(1-e^{i\vartheta}z)^2}$, $\vartheta\in\RR$, тоже лежат в $\Sc$, и у них
$a_n=ne^{i(n-1)\vartheta}$, $|a_n|=n$.

\textbf{Однозначные ветви корня и логарифма.} Если $\varphi\in H(G)$, $G$ односвязна и $\varphi\ne0$ в $G$, то существует
$L\in H(G)$ с $e^L=\varphi$ (ветвь $\log\varphi$), и тогда $\psi=e^{L/2}$ --- голоморфная ветвь $\sqrt\varphi$; ветвь фиксируется
значением в одной точке.

\textbf{Нечётная функция:} $h(-z)=-h(z)$; её тейлоровский ряд содержит лишь нечётные степени.
\end{defs}

\subsection{Предварительная лемма: <<квадратно-корневое преобразование>>}

\begin{thmbox}{Лемма 13.3}
Пусть $f\in\Sc$. Тогда найдётся $h\in\Sc$ такая, что $(h(z))^2=f(z^2)$ при $z\in\DD$. Функция $h$ нечётна.
\end{thmbox}

\begin{proof}
Функция $\varphi(z):=f(z)/z$ (доопределённая $\varphi(0)=f'(0)=1$) голоморфна в $\DD$ (устранимая особенность). Она не имеет нулей:
$\varphi(0)=1$, а при $z\ne0$ имеем $f(z)\ne f(0)=0$ по инъективности. Так как $\DD$ односвязен, существует $\psi\in H(\DD)$ с
$\psi^2=\varphi$, $\psi(0)=1$ (см. определения); $\psi$ также не имеет нулей. Положим
\[
 h(z):=z\,\psi(z^2).
\]
Тогда $h^2(z)=z^2\psi(z^2)^2=z^2\varphi(z^2)=f(z^2)$, $h(0)=0$, $h'(0)=\psi(0)=1$, и $h(-z)=-z\psi(z^2)=-h(z)$.

\emph{Однолистность.} Пусть $h(z_1)=h(z_2)$. Тогда $f(z_1^2)=h(z_1)^2=h(z_2)^2=f(z_2^2)$, откуда $z_1^2=z_2^2$ (инъективность $f$),
т.е. $z_1=\pm z_2$. Если $z_1=-z_2$, то $h(z_1)=-h(z_2)$ (нечётность), а вместе с $h(z_1)=h(z_2)$ это даёт $h(z_1)=0$.
Но $h(z)=z\psi(z^2)=0$ только при $z=0$ ($\psi\ne0$). Значит, $z_1=0=z_2$. В любом случае $z_1=z_2$.
\end{proof}

\begin{detail}{Первые коэффициенты $h$}
Из $\varphi(\zeta)=1+a_2\zeta+a_3\zeta^2+\dots$ и $\psi^2=\varphi$, $\psi(0)=1$, получаем $\psi(\zeta)=1+c\zeta+O(\zeta^2)$ с
$(1+c\zeta)^2=1+2c\zeta+\dots$, т.е. $2c=a_2$. Следовательно,
\[
 h(z)=z\Bigl(1+\frac{a_2}{2}z^2+O(z^4)\Bigr)=z+\frac{a_2}2z^3+O(z^5).
\]
\end{detail}

\subsection{Теорема Бибербаха $|a_2|\le2$}

\begin{thmbox}{Лемма 13.4 (оценка Бибербаха)}
Пусть $f\in\Sc$, $f(z)=z+\sum_{k\ge2}a_kz^k$. Тогда $|a_2|\le2$. Равенство $|a_2|=2$ имеет место тогда и только тогда, когда
$f$ --- вращение функции Кёбе: $f(z)=e^{-i\vartheta}k(e^{i\vartheta}z)$, $\vartheta\in\RR$.
\end{thmbox}

\begin{proof}
Пусть $h$ --- функция из леммы 13.3. Положим $F(z):=1/h(z)$. Так как $h$ имеет нуль только в $0$ (простой), $F\in H(\DD\setminus\{0\})$
и имеет простой полюс в $0$; $F$ инъективна в $\DD$ вместе с $h$. По детали выше
\[
 F(z)=\frac1{z\bigl(1+\frac{a_2}2z^2+O(z^4)\bigr)}=\frac1z\Bigl(1-\frac{a_2}2z^2+O(z^4)\Bigr)=\frac1z-\frac{a_2}{2}z+O(z^3).
\]
Итак, $F(z)=\frac1z+\sum_{n\ge0}\alpha_nz^n$ с $\alpha_0=0$, $\alpha_1=-a_2/2$. По следствию 3 теоремы площадей (билет 12)
$\sum n|\alpha_n|^2\le1$, в частности $|\alpha_1|\le1$, т.е. $|a_2|\le2$.

\emph{Случай равенства.} Если $|a_2|=2$, то $|\alpha_1|=1$ и $\sum_{n\ge2}n|\alpha_n|^2\le0$, т.е. $\alpha_n=0$ при $n\ge2$; кроме того $\alpha_0=0$
($F$ нечётна вместе с $h$, так что все $\alpha_{2m}=0$). Значит, $F(z)=\frac1z-\frac{a_2}2z$. Положим $e^{i\vartheta}:=a_2/2$. Тогда
\[
 h(z)=\frac{z}{1-e^{i\vartheta}z^2},\qquad f(z^2)=h(z)^2=\frac{z^2}{(1-e^{i\vartheta}z^2)^2},
\]
т.е. $f(\zeta)=\frac{\zeta}{(1-e^{i\vartheta}\zeta)^2}=e^{-i\vartheta}k(e^{i\vartheta}\zeta)$ при всех $\zeta\in\DD$ (каждое $\zeta\in\DD$ есть $z^2$, $z\in\DD$).
Обратно, для вращений функции Кёбе $a_2=2e^{i\vartheta}$, $|a_2|=2$.
\end{proof}

\begin{detail}{Функция Кёбе лежит в $\Sc$ и $k(\DD)=\CC\setminus(-\infty,-\tfrac14]$}
Тождество $\frac14\bigl(\frac{1+z}{1-z}\bigr)^2-\frac14=\frac{(1+z)^2-(1-z)^2}{4(1-z)^2}=\frac{4z}{4(1-z)^2}=k(z)$. Отображение
$z\mapsto\frac{1+z}{1-z}$ --- дробно-линейная биекция $\DD$ на правую полуплоскость $\Pi=\{\Re u>0\}$ (точки $\TT$ переходят в мнимую ось,
$0\mapsto1$). Отображение $u\mapsto u^2$ --- биекция $\Pi$ на $\CC\setminus(-\infty,0]$ (обратная --- главная ветвь корня). Наконец,
$v\mapsto\frac14v-\frac14$ переводит $\CC\setminus(-\infty,0]$ в $\CC\setminus(-\infty,-\frac14]$. Композиция инъективна, так что $k\in\mathcal U(\DD)$;
$k(0)=0$, $k'(0)=1$ (из ряда $\sum nz^n$). Значит, $k\in\Sc$ и $k(\DD)=\CC\setminus(-\infty,-\tfrac14]$, в частности $-\frac14\notin k(\DD)$.
Для вращения $k_\vartheta(\DD)=e^{-i\vartheta}k(\DD)=\CC\setminus e^{-i\vartheta}(-\infty,-\frac14]$.
\end{detail}

\subsection{Теорема Кёбе об $1/4$}

\begin{thmbox}{Предложение 13.10 (теорема Кёбе об $1/4$)}
Пусть $f\in\Sc$. Тогда $B(0,1/4)\subset f(\DD)$. Константа $1/4$ точная: для функции Кёбе $-\frac14\notin k(\DD)$.
Более того, если $w\notin f(\DD)$ и $|w|=\frac14$, то $f$ --- вращение функции Кёбе.
\end{thmbox}

\begin{proof}
Пусть $w\in\CC\setminus f(\DD)$ (если таких нет, доказывать нечего); $w\ne0$, так как $0=f(0)$. Положим $q=1/w$ и
\[
 \varphi(z):=\frac{f(z)}{1-qf(z)}=\frac{w\,f(z)}{w-f(z)},\qquad z\in\DD.
\]
\begin{detail}{$\varphi\in\Sc$ и $\varphi(z)=z+(a_2+w^{-1})z^2+\dots$}
Знаменатель $w-f(z)\ne0$, так как $f$ не принимает значение $w$; значит, $\varphi\in H(\DD)$. Далее $\varphi=T\circ f$, где
$T(\zeta)=\frac{w\zeta}{w-\zeta}$ --- дробно-линейное отображение (определитель матрицы $\begin{pmatrix}w&0\\-1&w\end{pmatrix}$ равен $w^2\ne0$),
инъективное на $\CC\setminus\{w\}\supset f(\DD)$. Значит, $\varphi$ инъективна как композиция инъективных. $\varphi(0)=0$.
Разложение: при малых $z$ имеем $|qf(z)|<1$ и
\[
 \varphi(z)=f(z)\bigl(1+qf(z)+q^2f(z)^2+\dots\bigr)=(z+a_2z^2+\dots)(1+qz+O(z^2))=z+(a_2+q)z^2+O(z^3).
\]
Итак, $\varphi'(0)=1$, $\varphi''(0)=2(a_2+w^{-1})$, $\varphi\in\Sc$ и второй коэффициент $\varphi$ равен $a_2+\frac1w$.
\end{detail}
Применяя лемму 13.4 к $\varphi$ и к $f$, получаем $|a_2+\frac1w|\le2$ и $|a_2|\le2$, откуда
\[
 \Bigl|\frac1w\Bigr|\le\Bigl|a_2+\frac1w\Bigr|+|a_2|\le4,\qquad\text{т.е.}\quad |w|\ge\frac14 .
\]
Итак, каждая точка, не лежащая в $f(\DD)$, удовлетворяет $|w|\ge\frac14$, т.е. $B(0,\frac14)\subset f(\DD)$.

\emph{Случай $|w|=\frac14$.} Тогда в цепочке $4=|1/w|\le|a_2+1/w|+|a_2|\le2+2$ все неравенства --- равенства, в частности $|a_2|=2$, и
по лемме 13.4 $f$ --- вращение функции Кёбе.

\emph{Точность} следует из детали про $k(\DD)$: $-\frac14\notin k(\DD)$, так что никакой круг $B(0,\rho)$ с $\rho>\frac14$ не лежит в $k(\DD)$.
\end{proof}

\begin{remarkbox}{Задача со слайда 308: где достигаются равенства}
\begin{itemize}
\item \emph{Теорема площадей:} $\sum k|b_k|^2=1\iff m_2(K_g)=0$ (билет 12); например, $g(z)=z+b_0+e^{i\alpha}/z$, $K_g$ --- отрезок
(пример, связанный с функцией Кёбе, --- ниже).
\item \emph{Оценка Бибербаха:} $|a_2|=2$ ровно для $k_\vartheta(z)=e^{-i\vartheta}k(e^{i\vartheta}z)$ (лемма 13.4).
\item \emph{Теорема Кёбе:} радиус $\frac14$ достигается (граничная точка образа на расстоянии $\frac14$ от нуля) ровно для вращений функции Кёбе.
\end{itemize}
Для $f=k$: $h(z)=\frac{z}{1-z^2}$, $F(z)=\frac1z-z$, $g(z)=F(1/z)=z-\frac1z$; при $z=iu$: $g=i(u+\frac1u)$, т.е. $K_g=[-2i,2i]$ --- отрезок
длины 4, $m_2(K_g)=0$, $\sum n|\alpha_n|^2=1$.
\end{remarkbox}

\begin{remarkbox}{Задача со слайда 308: без однолистности теорема Кёбе неверна}
Пусть $f_n(z)=\frac{e^{nz}-1}{n}$. Тогда $f_n\in H(\DD)$, $f_n(0)=0$, $f_n'(0)=1$, но $f_n$ не принимает значение $-\frac1n$ (так как $e^{nz}\ne0$).
Поэтому никакой круг $B(0,r_0)$ с фиксированным $r_0>0$ не может содержаться во всех $f_n(\DD)$ (при $\frac1n<r_0$ точка $-\frac1n\in B(0,r_0)\setminus f_n(\DD)$).
Конечно, $f_n$ не однолистны в $\DD$ при $n\ge4$: $e^{nz}$ принимает равные значения в точках $z$ и $z+2\pi i/n$, а $|2\pi i/n|<2$.
\end{remarkbox}

\begin{remarkbox}{Гипотеза Бибербаха (1916) и её история}
Для всех $f\in\Sc$ и всех $n$: $|a_n|\le n$, причём $|a_n|=n$ (при каком-нибудь $n\ge2$) лишь для вращений функции Кёбе.
Лемма 13.4 --- случай $n=2$. Литтлвуд (1925) доказал $|a_n|<en$ (вопрос 41*). Полностью гипотеза доказана Л.~де Бранжем в 1984 г.
\end{remarkbox}

\begin{remarkbox}{Коротко: план ответа на экзамене}
\begin{enumerate}
\item Класс $\Sc$, функция Кёбе $k(z)=z/(1-z)^2=\sum nz^n$, $k(\DD)=\CC\setminus(-\infty,-\frac14]$.
\item Лемма 13.3: $f(z)=z\varphi(z)$, $\varphi\ne0$, $\psi=\sqrt\varphi$, $h(z)=z\psi(z^2)$ однолистна и нечётна, $h^2=f(z^2)$.
\item $h=z+\frac{a_2}2z^3+\dots$, $F=1/h=\frac1z-\frac{a_2}2z+\dots$; по следствию теоремы площадей $|a_2/2|\le1$.
\item Равенство: $F=\frac1z-e^{i\vartheta}z$, откуда $f=k_\vartheta$.
\item Кёбе: при $w\notin f(\DD)$ функция $\varphi=\frac{wf}{w-f}\in\Sc$, её $a_2$ равен $a_2+\frac1w$.
\item $|1/w|\le|a_2+1/w|+|a_2|\le4$, т.е. $|w|\ge\frac14$; точность --- функция Кёбе.
\item Без однолистности неверно: $(e^{nz}-1)/n$ не принимает $-1/n$.
\end{enumerate}
\end{remarkbox}

% -------------------------------------------------------------------
\ticket{14}{Первая и вторая теоремы растяжения для однолистных функций в круге}
\label{t14:begin}

\begin{defs}
\textbf{Класс $\Sc$} $=\{f\in\mathcal U(\DD):\ f(0)=0,\ f'(0)=1\}$, $f(z)=z+a_2z^2+\dots$ (билет 13). Для однолистной $f$ имеем $f'\ne0$.
Если $f\in\mathcal U(\DD)$, то $g=\frac{f-f(0)}{f'(0)}\in\Sc$ и $g'=f'/f'(0)$, $g(z)=\frac{f(z)-f(0)}{f'(0)}$.

\textbf{Оценка Бибербаха} (лемма 13.4, билет 13): $|a_2|\le2$ для $f\in\Sc$. \textbf{Теорема Кёбе} (предложение 13.10):
$B(0,\tfrac14)\subset f(\DD)$ для $f\in\Sc$.

\textbf{Автоморфизмы круга:} при $w\in\DD$ отображение $\varphi_w(z)=\frac{z+w}{1+\bar wz}$ --- голоморфная биекция $\DD\to\DD$
(обратная --- $\varphi_{-w}$), $\varphi_w(0)=w$.

\textbf{Коэффициент растяжения:} $|f'(z)|$ --- локальный коэффициент растяжения отображения $z\mapsto f(z)$:
$|f(z+h)-f(z)|=|f'(z)||h|+o(|h|)$. Теоремы растяжения (искажения) оценивают $|f'(z)|$ и $|f(z)|$ для $f\in\Sc$ через $|z|$.

\textbf{Преобразование Кёбе} функции $f\in\Sc$ с параметром $w\in\DD$:
\[
 f_w(z)=\frac{f\bigl(\frac{z+w}{1+\bar wz}\bigr)-f(w)}{(1-|w|^2)f'(w)},\qquad z\in\DD .
\]
Смысл: <<перенести>> точку $w$ в $0$ автоморфизмом круга и снова нормировать. $a_2(w)$ --- второй тейлоровский коэффициент $f_w$.

\textbf{Функция Кёбе} $k(z)=\frac{z}{(1-z)^2}$, $k'(z)=\frac{1+z}{(1-z)^3}$ (проверка: $\frac{(1-z)^2+2z(1-z)}{(1-z)^4}=\frac{1+z}{(1-z)^3}$).
\end{defs}

\subsection{Преобразование Кёбе и лемма 13.5}

\begin{thmbox}{Упражнения со слайда 309}
Пусть $f\in\Sc$, $w\in\DD$. Тогда (а) $f_w\in\Sc$; (б) второй коэффициент $f_w$ равен
\[
 a_2(w)=\frac12\Bigl((1-|w|^2)\frac{f''(w)}{f'(w)}-2\bar w\Bigr).
\]
\end{thmbox}

\begin{proof}
Обозначим $\varphi=\varphi_w$. Вычислим производные:
\[
 \varphi'(z)=\frac{(1+\bar wz)-\bar w(z+w)}{(1+\bar wz)^2}=\frac{1-|w|^2}{(1+\bar wz)^2},\qquad
 \varphi''(z)=\frac{-2\bar w(1-|w|^2)}{(1+\bar wz)^3},
\]
так что $\varphi(0)=w$, $\varphi'(0)=1-|w|^2$, $\varphi''(0)=-2\bar w(1-|w|^2)$.

(а) $f\circ\varphi$ голоморфна и инъективна в $\DD$ (композиция биекции $\varphi:\DD\to\DD$ и инъективной $f$). Знаменатель $(1-|w|^2)f'(w)$ --- ненулевая
константа ($|w|<1$, $f'(w)\ne0$). Значит, $f_w=\frac{f\circ\varphi-f(w)}{(1-|w|^2)f'(w)}$ голоморфна и инъективна (аффинная функция от инъективной), причём
\[
 f_w(0)=\frac{f(w)-f(w)}{\dots}=0,\qquad f_w'(0)=\frac{f'(\varphi(0))\varphi'(0)}{(1-|w|^2)f'(w)}=\frac{f'(w)(1-|w|^2)}{(1-|w|^2)f'(w)}=1 .
\]
Итак, $f_w\in\Sc$.

(б) По правилу дифференцирования композиции $(f\circ\varphi)''=f''(\varphi)\,(\varphi')^2+f'(\varphi)\,\varphi''$, поэтому
\[
 f_w''(0)=\frac{f''(w)(1-|w|^2)^2-2\bar w(1-|w|^2)f'(w)}{(1-|w|^2)f'(w)}=(1-|w|^2)\frac{f''(w)}{f'(w)}-2\bar w .
\]
Второй тейлоровский коэффициент равен $a_2(w)=\frac12f_w''(0)$, что и даёт формулу.
\end{proof}

\begin{thmbox}{Лемма 13.5}
Если $f\in\Sc$, то
\[
 \Bigl|(1-|z|^2)\frac{f''(z)}{f'(z)}-2\bar z\Bigr|\le4\qquad\text{при } |z|<1 .
\]
Для функции Кёбе $k$ при вещественных $z\in(-1,1)$ здесь достигается равенство.
\end{thmbox}

\begin{proof}
По упражнению (а) $f_z\in\Sc$, по оценке Бибербаха $|a_2(z)|\le2$, а по (б) $2a_2(z)$ --- это в точности выражение под модулем. Значит, оно
по модулю $\le4$.

\emph{Равенство для $k$.} $\log k'(z)=\log(1+z)-3\log(1-z)$ (ветви, равные $0$ в $0$), откуда
\[
 \frac{k''(z)}{k'(z)}=\frac1{1+z}+\frac3{1-z}=\frac{(1-z)+3(1+z)}{1-z^2}=\frac{4+2z}{1-z^2}.
\]
При $z=x\in(-1,1)$: $(1-x^2)\frac{4+2x}{1-x^2}-2x=4$. (Для вращения $k_\vartheta(z)=e^{-i\vartheta}k(e^{i\vartheta}z)$ равенство достигается на
диаметре $\{te^{-i\vartheta}:\ -1<t<1\}$.)
\end{proof}

\subsection{Первая теорема растяжения}

\begin{thmbox}{Теорема 13.4 (первая теорема растяжения)}
Пусть $f\in\mathcal U(\DD)$. Тогда для любого $z\in\DD$
\[
 |f'(0)|\frac{1-|z|}{(1+|z|)^3}\le|f'(z)|\le|f'(0)|\frac{1+|z|}{(1-|z|)^3}.
\]
Обе оценки точны: для $f=k$ верхняя достигается при $z=r\in[0,1)$, нижняя --- при $z=-r$.
\end{thmbox}

\begin{proof}
\textbf{Сведение к $\Sc$.} $g=\frac{f-f(0)}{f'(0)}\in\Sc$ и $|g'(z)|=|f'(z)|/|f'(0)|$; поэтому достаточно доказать для $f\in\Sc$ неравенства
\[
 \frac{1-r}{(1+r)^3}\le|f'(z)|\le\frac{1+r}{(1-r)^3},\qquad r=|z|.
\]

\textbf{Шаг 1: оценка $\Re\frac{zf''}{f'}$.} Пусть $z\in\DD$, $r=|z|$. Имеем тождество
\[
 \frac{z}{1-|z|^2}\Bigl((1-|z|^2)\frac{f''(z)}{f'(z)}-2\bar z\Bigr)=z\frac{f''(z)}{f'(z)}-\frac{2|z|^2}{1-|z|^2}
\]
(так как $z\bar z=|z|^2$). Беря модули и применяя лемму 13.5:
\[
 \Bigl|z\frac{f''(z)}{f'(z)}-\frac{2r^2}{1-r^2}\Bigr|=\frac{r}{1-r^2}\Bigl|(1-r^2)\frac{f''(z)}{f'(z)}-2\bar z\Bigr|\le\frac{4r}{1-r^2}.
 \tag{14.1}\label{t14:disc}
\]
Итак, число $z\frac{f''(z)}{f'(z)}$ лежит в замкнутом круге с центром $\frac{2r^2}{1-r^2}$ (вещественным) и радиусом $\frac{4r}{1-r^2}$; поскольку
$|\Re\zeta-c|\le|\zeta-c|$ при вещественном $c$,
\[
 \frac{2r^2-4r}{1-r^2}\le\Re\frac{zf''(z)}{f'(z)}\le\frac{2r^2+4r}{1-r^2}\qquad(|z|=r<1).
 \tag{14.2}\label{t14:re}
\]

\textbf{Шаг 2: переход к радиальной производной.}
\begin{detail}{Упражнение со слайда 311: $\Re(zh'(z))=r\,\Re\frac{\partial h}{\partial r}$ при $z=re^{i\vartheta}$}
Для $h\in H(\DD)$ и фиксированного $\vartheta$ функция $r\mapsto h(re^{i\vartheta})$ дифференцируема и по цепному правилу
$\frac{\partial}{\partial r}h(re^{i\vartheta})=h'(re^{i\vartheta})\cdot e^{i\vartheta}$. Умножая на $r$: $r\frac{\partial h}{\partial r}=re^{i\vartheta}h'(z)=zh'(z)$.
Это верно даже без $\Re$; взяв вещественную часть, получаем нужное.
\end{detail}
Так как $f'\ne0$ в односвязном $\DD$, существует голоморфная ветвь $L=\log f'$ с $L(0)=\log1=0$, и $\Re L=\log|f'|$, $L'=f''/f'$.
По упражнению с $h=L$:
\[
 \Re\frac{zf''(z)}{f'(z)}=\Re\bigl(zL'(z)\bigr)=r\frac{\partial}{\partial r}\Re L(re^{i\vartheta})=r\frac{\partial}{\partial r}\log|f'(re^{i\vartheta})|.
\]
Делим \eqref{t14:re} на $r>0$:
\[
 \frac{2r-4}{1-r^2}\le\frac{\partial}{\partial r}\log|f'(re^{i\vartheta})|\le\frac{2r+4}{1-r^2},\qquad 0<r<1 .
 \tag{14.3}\label{t14:dr}
\]

\textbf{Шаг 3: интегрирование.} Зафиксируем $\vartheta$ и положим $u(r)=\log|f'(re^{i\vartheta})|$; $u\in C^1[0,1)$, $u(0)=\log|f'(0)|=0$.
Разложим на простейшие дроби:
\[
 \frac{2t+4}{1-t^2}=\frac{3}{1-t}+\frac{1}{1+t},\qquad \frac{2t-4}{1-t^2}=-\frac{1}{1-t}-\frac{3}{1+t}
\]
(проверка: $3(1+t)+(1-t)=2t+4$; $-(1+t)-3(1-t)=2t-4$). Обе функции непрерывны на $[0,1)$, так что, интегрируя \eqref{t14:dr} по $t\in[0,R]$,
$0<R<1$, получаем ($u(R)=\int_0^Ru'(t)\,dt$):
\begin{align*}
 u(R)&\le\int_0^R\Bigl(\frac3{1-t}+\frac1{1+t}\Bigr)dt=-3\log(1-R)+\log(1+R)=\log\frac{1+R}{(1-R)^3},\\
 u(R)&\ge\int_0^R\Bigl(-\frac1{1-t}-\frac3{1+t}\Bigr)dt=\log(1-R)-3\log(1+R)=\log\frac{1-R}{(1+R)^3}.
\end{align*}
Применяя $\exp$ (монотонная функция), при $z=Re^{i\vartheta}$:
\[
 \frac{1-R}{(1+R)^3}\le|f'(z)|\le\frac{1+R}{(1-R)^3}.
\]
При $z=0$ неравенства очевидны ($|f'(0)|=1$).

\textbf{Точность.} Для $k$: $k'(r)=\frac{1+r}{(1-r)^3}$, $k'(-r)=\frac{1-r}{(1+r)^3}$ --- верхняя и нижняя оценки достигаются.
\end{proof}

\begin{remarkbox}{Почему для функции Кёбе всё обращается в равенства}
При $z=\pm r$ для $f=k$ в лемме 13.5 стоит равенство, а выражение $z\frac{k''}{k'}=\frac{4z+2z^2}{1-z^2}$ при $z=r$ равно
$\frac{2r^2+4r}{1-r^2}$ (правый край \eqref{t14:re}), при $z=-r$ --- $\frac{2r^2-4r}{1-r^2}$ (левый край). Значит, на лучах $[0,1)$ и $(-1,0]$
все неравенства доказательства превращаются в равенства.
\end{remarkbox}

\subsection{Вторая теорема растяжения (теорема роста)}

\begin{thmbox}{Теорема 13.5 (вторая теорема растяжения)}
Пусть $f\in\mathcal U(\DD)$. Тогда для любого $z\in\DD$
\[
 |f'(0)|\frac{|z|}{(1+|z|)^2}\le|f(z)-f(0)|\le|f'(0)|\frac{|z|}{(1-|z|)^2}.
\]
Обе оценки точны: для $f=k$ верхняя достигается при $z=r$, нижняя --- при $z=-r$.
\end{thmbox}

\begin{proof}
Как и выше, достаточно рассмотреть $f\in\Sc$ (для $g=\frac{f-f(0)}{f'(0)}$ имеем $|g(z)|=\frac{|f(z)-f(0)|}{|f'(0)|}$). Пусть $z=re^{i\vartheta}$, $0<r<1$
(при $z=0$ всё очевидно). Нам понадобятся первообразные:
\[
 \frac{d}{dt}\frac{t}{(1-t)^2}=\frac{(1-t)+2t}{(1-t)^3}=\frac{1+t}{(1-t)^3},\qquad
 \frac{d}{dt}\frac{t}{(1+t)^2}=\frac{(1+t)-2t}{(1+t)^3}=\frac{1-t}{(1+t)^3}.
\]

\textbf{Верхняя оценка.} По формуле Ньютона--Лейбница вдоль радиуса $t\mapsto te^{i\vartheta}$: $f(z)=f(z)-f(0)=\int_0^rf'(te^{i\vartheta})e^{i\vartheta}\,dt$.
По первой теореме растяжения
\[
 |f(z)|\le\int_0^r|f'(te^{i\vartheta})|\,dt\le\int_0^r\frac{1+t}{(1-t)^3}\,dt=\frac{r}{(1-r)^2}.
\]

\textbf{Нижняя оценка.} Отметим, что $\frac{r}{(1+r)^2}<\frac14$ при $r\in(0,1)$, так как $(1+r)^2-4r=(1-r)^2>0$.

\emph{Случай 1: $|f(z)|\ge\frac14$.} Тогда $|f(z)|\ge\frac14>\frac{r}{(1+r)^2}$.

\emph{Случай 2: $|f(z)|<\frac14$.} По теореме Кёбе $B(0,\frac14)\subset f(\DD)$, поэтому отрезок $[0,f(z)]\subset B(0,\frac14)\subset f(\DD)$. Функция
$f^{-1}:f(\DD)\to\DD$ голоморфна. Положим
\[
 \gamma(s):=f^{-1}\bigl(s\,f(z)\bigr),\qquad s\in[0,1].
\]
Это гладкая кривая в $\DD$ (прообраз отрезка) с $\gamma(0)=0$, $\gamma(1)=z$. Дифференцируя тождество $f(\gamma(s))=sf(z)$, получаем
$f'(\gamma(s))\gamma'(s)=f(z)$ --- т.е. аргумент $f'(\zeta)\,d\zeta$ вдоль $\gamma$ постоянен (равен $\arg f(z)$). Поэтому
\[
 |f(z)|=\int_0^1|f(z)|\,ds=\int_0^1|f'(\gamma(s))|\,|\gamma'(s)|\,ds=\int_\gamma|f'(\zeta)|\,|d\zeta| .
\]
Пусть $\rho(s)=|\gamma(s)|$. Так как $\gamma(s)=0$ лишь при $s=0$ (при $s>0$ имеем $f(\gamma(s))=sf(z)\ne0$, ибо $z\ne0$), функция $\rho$ непрерывна на
$[0,1]$, гладкая на $(0,1]$ и $\rho'(s)=\Re\bigl(\overline{\gamma(s)}\gamma'(s)\bigr)/|\gamma(s)|$, откуда $\rho'(s)\le|\gamma'(s)|$. Положим
$\phi(t)=\frac{1-t}{(1+t)^3}>0$, $\Phi(t)=\frac{t}{(1+t)^2}$, $\Phi'=\phi$. По первой теореме растяжения $|f'(\gamma(s))|\ge\phi(\rho(s))$, так что при $0<\delta<1$
\[
 \int_\delta^1|f'(\gamma(s))||\gamma'(s)|\,ds\ge\int_\delta^1\phi(\rho(s))\rho'(s)\,ds=\Phi(\rho(1))-\Phi(\rho(\delta))=\Phi(r)-\Phi(\rho(\delta)).
\]
Устремляя $\delta\to0$ ($\rho(\delta)\to0$, $\Phi(0)=0$), получаем
\[
 |f(z)|=\int_\gamma|f'(\zeta)||d\zeta|\ge\Phi(r)=\int_0^r\frac{1-t}{(1+t)^3}\,dt=\frac{r}{(1+r)^2}.
\]

\textbf{Точность.} $k(r)=\frac{r}{(1-r)^2}$, $|k(-r)|=\frac{r}{(1+r)^2}$.
\end{proof}

\begin{remarkbox}{Замечания}
\begin{itemize}
\item Смысл нижней оценки слайда: $|f(z)|=\int_\gamma|f'||d\zeta|$, а путь $\gamma$ от $0$ до $z$ <<не короче>> радиуса: при движении по $\gamma$ модуль $|\zeta|$
пробегает все значения от $0$ до $r$, и оценка $|f'(\zeta)|\ge\phi(|\zeta|)$ интегрируется по $d|\zeta|\le|d\zeta|$.
\item Наивная попытка оценить $|f(z)|$ снизу интегралом вдоль радиуса не проходит: $|\int f'\,d\zeta|$ может быть много меньше $\int|f'||d\zeta|$.
Именно поэтому берётся прообраз отрезка, вдоль которого аргумент $f'(\zeta)d\zeta$ постоянен; для этого нужна теорема Кёбе.
\item Следствие: при $|z|\to1$ нижняя оценка даёт $\liminf|f(z)|\ge\frac14$, т.е. снова теорему Кёбе (но она использовалась в доказательстве).
\item Из второй теоремы следует, что класс $\Sc$ локально равномерно ограничен ($|f(z)|\le\frac{r}{(1-r)^2}$ при $|z|\le r$) и, по теореме Монтеля,
компактен в себе (предельная функция однолистна по теореме Гурвица и $f'(0)=1$).
\end{itemize}
\end{remarkbox}

\begin{remarkbox}{Коротко: план ответа на экзамене}
\begin{enumerate}
\item Преобразование Кёбе $f_w$; проверка $f_w\in\Sc$; $a_2(w)=\frac12\bigl((1-|w|^2)\frac{f''(w)}{f'(w)}-2\bar w\bigr)$.
\item Бибербах $\Rightarrow$ лемма 13.5: $|(1-|z|^2)\frac{f''}{f'}-2\bar z|\le4$.
\item Умножаем на $\frac{z}{1-|z|^2}$: $|z\frac{f''}{f'}-\frac{2r^2}{1-r^2}|\le\frac{4r}{1-r^2}$, отсюда оценка $\Re\frac{zf''}{f'}$.
\item $\Re\frac{zf''}{f'}=r\partial_r\log|f'|$; интегрируем по $r$: $\frac{1-r}{(1+r)^3}\le|f'|\le\frac{1+r}{(1-r)^3}$.
\item Вторая теорема, верх: интеграл $|f'|$ по радиусу.
\item Низ: если $|f(z)|\ge\frac14$ --- ясно; иначе прообраз отрезка $[0,f(z)]$ (Кёбе), $|f(z)|=\int_\gamma|f'||d\zeta|\ge\int_0^r\frac{1-t}{(1+t)^3}dt$.
\item Общий случай $f\in\mathcal U(\DD)$: нормировка $\frac{f-f(0)}{f'(0)}$. Экстремальна функция Кёбе ($z=\pm r$).
\end{enumerate}
\end{remarkbox}

% -------------------------------------------------------------------
\ticket{15}{Радиус выпуклости для класса $\Sc$ однолистных функций в круге $\DD$ с нормировками $f(0)=0$ и $f'(0)=1$}
\label{t15:begin}

\begin{defs}
\textbf{Класс $\Sc$} $=\{f\in\mathcal U(\DD):\ f(0)=0,\ f'(0)=1\}$ (билет 13); для $f\in\Sc$ имеем $f'\ne0$ в $\DD$.
Обозначения: $B_r=\{|z|<r\}$, $C_r=\{|z|=r\}$.

\textbf{Выпуклое множество:} $E\subset\CC$ выпукло, если вместе с любыми двумя точками содержит соединяющий их отрезок.
\emph{Пример:} $k(B_r)$ при малых $r$ выпукло, а $k(\DD)=\CC\setminus(-\infty,-\frac14]$ --- нет (середина отрезка между $-1+i$ и $-1-i$
равна $-1\notin k(\DD)$).

\textbf{Звёздное (относительно $0$) множество:} вместе с каждой точкой $w$ содержит отрезок $[0,w]$.

\textbf{Радиус выпуклости класса $\Sc$:}
\[
 R_C(\Sc):=\sup\{r\in(0,1]:\ f(B_r)\text{ выпукла для всех } f\in\Sc\}.
\]
\textbf{Радиус звёздности:} $R_S(\Sc):=\sup\{r:\ f(B_r)\text{ звёздна относительно }0\ \text{для всех } f\in\Sc\}$ (доп. материал).

\textbf{Касательный угол кривой.} Для гладкой кривой $\Gamma(\vartheta)$ с $\Gamma'(\vartheta)\ne0$ \emph{угол касательной} $\tau(\vartheta)$ --- непрерывная ветвь
$\arg\Gamma'(\vartheta)$. Для $\Gamma(\vartheta)=f(re^{i\vartheta})$: $\Gamma'(\vartheta)=f'(z)\cdot iz$, $z=re^{i\vartheta}$, поэтому
$\tau(\vartheta)=\vartheta+\frac\pi2+\arg f'(re^{i\vartheta})$ --- касательная к $C_r$ в точке $z$ образует угол $\vartheta+\frac\pi2$ с осью $\RR$,
и отображение $f$ поворачивает её на $\arg f'(z)$.

\textbf{Функция Кёбе} $k(z)=\frac{z}{(1-z)^2}$, $\frac{k''(z)}{k'(z)}=\frac{4+2z}{1-z^2}$ (билет 14).
\end{defs}

\begin{thmbox}{Предложение 13.13}
$R_C(\Sc)=2-\sqrt3\approx0{,}268$. Точнее: при $r\le2-\sqrt3$ область $f(B_r)$ выпукла для каждой $f\in\Sc$, а при $r\in(2-\sqrt3,1)$ область
$k(B_r)$ не выпукла.
\end{thmbox}

\subsection{Аналитическое условие выпуклости}

\begin{detail}{Упражнение со слайда 320: $\frac{\partial}{\partial\vartheta}\bigl(\frac\pi2+\vartheta+\arg f'(z)\bigr)=1+\Re\frac{zf''(z)}{f'(z)}$, $z=re^{i\vartheta}$}
Пусть $L$ --- голоморфная ветвь $\log f'$ в $\DD$ (существует, так как $f'\ne0$, $\DD$ односвязен); тогда $\arg f'(z)=\Im L(z)$ --- непрерывная (и
$2\pi$-периодическая по $\vartheta$) ветвь аргумента. Так как $\frac{\partial z}{\partial\vartheta}=ire^{i\vartheta}=iz$, по цепному правилу
\[
 \frac{\partial}{\partial\vartheta}L(re^{i\vartheta})=L'(z)\cdot iz=iz\frac{f''(z)}{f'(z)},\qquad
 \frac{\partial}{\partial\vartheta}\Im L=\Im\Bigl(iz\frac{f''(z)}{f'(z)}\Bigr)=\Re\Bigl(z\frac{f''(z)}{f'(z)}\Bigr).
\]
Прибавляя $\frac{\partial}{\partial\vartheta}(\frac\pi2+\vartheta)=1$, получаем утверждение. Итак,
\[
 \tau'(\vartheta)=1+\Re\frac{zf''(z)}{f'(z)},\qquad \tau(\vartheta+2\pi)=\tau(\vartheta)+2\pi .
\]
\end{detail}

На слайде <<принцип выпуклости>> (область $f(B_r)$ выпукла $\iff$ касательная к $f(C_r)$ вращается в одну сторону) считается известным. Для
доказательства предложения нам нужна только одна его половина, и её мы докажем полностью; вторая половина (для функции Кёбе)
заменена ниже прямым рассуждением.

\begin{thmbox}{Лемма (достаточное условие выпуклости)}
Пусть $f\in\mathcal U(\DD)$, $0<r<1$, и $1+\Re\frac{zf''(z)}{f'(z)}>0$ при всех $z\in C_r$. Тогда область $f(B_r)$ выпукла.
\end{thmbox}

\begin{proof}
Пусть $\Gamma(\vartheta)=f(re^{i\vartheta})$, $\vartheta\in\RR$ ($2\pi$-периодическая). Это гладкая жорданова кривая ($f$ инъективна на $C_r$, $\Gamma'\ne0$).

\emph{(1) $f(B_r)$ --- внутренность $D(\Gamma)$.} Множество $f(B_r)$ открыто, связно, ограничено и не пересекается с $\Gamma$. Оно замкнуто в $\CC\setminus\Gamma$:
если $f(z_n)\to w\notin\Gamma$, $z_n\in B_r$, то по подпоследовательности $z_n\to z\in\clos{B_r}$, $w=f(z)$, и $z\notin C_r$ (иначе $w\in\Gamma$). Значит, $f(B_r)$ ---
компонента связности $\CC\setminus\Gamma$, причём ограниченная, т.е. $f(B_r)=D(\Gamma)$ (теорема Жордана).

\emph{(2) Кривая лежит по одну сторону от каждой своей касательной.} По упражнению $\tau'>0$, т.е. $\tau$ строго возрастает, и $\tau(\vartheta_0+2\pi)=\tau(\vartheta_0)+2\pi$.
Зафиксируем $\vartheta_0$ и положим
\[
 \Phi(\vartheta)=\Im\bigl[(\Gamma(\vartheta)-\Gamma(\vartheta_0))e^{-i\tau(\vartheta_0)}\bigr],\qquad
 \Phi'(\vartheta)=\Im\bigl[\Gamma'(\vartheta)e^{-i\tau(\vartheta_0)}\bigr]=|\Gamma'(\vartheta)|\sin\bigl(\tau(\vartheta)-\tau(\vartheta_0)\bigr).
\]
На $(\vartheta_0,\vartheta_0+2\pi)$ разность $\tau(\vartheta)-\tau(\vartheta_0)$ строго возрастает от $0$ до $2\pi$; пусть $\vartheta_1$ --- точка, где она равна $\pi$.
Тогда $\Phi'>0$ на $(\vartheta_0,\vartheta_1)$, $\Phi'<0$ на $(\vartheta_1,\vartheta_0+2\pi)$, а $\Phi(\vartheta_0)=\Phi(\vartheta_0+2\pi)=0$. Поэтому $\Phi>0$ на
$(\vartheta_0,\vartheta_0+2\pi)$. Итак, $\Gamma$ лежит в замкнутой полуплоскости
$H_{\vartheta_0}=\{w:\ \Im[(w-\Gamma(\vartheta_0))e^{-i\tau(\vartheta_0)}]\ge0\}$ (слева от касательной) и касается её границы лишь в точке $\Gamma(\vartheta_0)$.

\emph{(3) $D(\Gamma)$ выпукла.} Открытая полуплоскость $\CC\setminus H_{\vartheta_0}$ связна, неограниченна и не пересекает $\Gamma$, поэтому лежит во внешности $\Gamma$;
значит, $D(\Gamma)\subset H_{\vartheta_0}$ для всех $\vartheta_0$, т.е. $D(\Gamma)\subset K:=\bigcap_{\vartheta_0}H_{\vartheta_0}$, а так как $D(\Gamma)$ открыта ---
$D(\Gamma)\subset K^\circ$ (внутренность $K$). Множество $K$ выпукло (пересечение полуплоскостей), значит, $K^\circ$ выпукло и, следовательно, связно.
$K^\circ$ не пересекает $\Gamma$: точка $\Gamma(\vartheta_0)$ лежит на границе $H_{\vartheta_0}$, а значит, не во внутренности $K$. Поэтому связное $K^\circ$
лежит в одной компоненте $\CC\setminus\Gamma$; так как $K^\circ\supset D(\Gamma)\ne\varnothing$, это $D(\Gamma)$. Итак, $D(\Gamma)=K^\circ$ --- выпуклое множество.
\end{proof}

\subsection{Доказательство предложения 13.13}

\begin{proof}
\textbf{Шаг 1: $R_C(\Sc)\ge2-\sqrt3$.} Пусть $f\in\Sc$, $0<r<1$, $|z|=r$. Как показано в доказательстве первой теоремы растяжения (следствие леммы 13.5,
билет 14, неравенство (14.2)),
\[
 \Re\frac{zf''(z)}{f'(z)}\ge\frac{2r^2-4r}{1-r^2},\qquad\text{откуда}\qquad
 1+\Re\frac{zf''(z)}{f'(z)}\ge\frac{2r^2-4r+1-r^2}{1-r^2}=\frac{r^2-4r+1}{1-r^2}.
\]
Корни трёхчлена $r^2-4r+1$ --- числа $2\pm\sqrt3$; при $0<r<2-\sqrt3$ он положителен. По лемме область $f(B_r)$ выпукла при всех
$r<2-\sqrt3$ и всех $f\in\Sc$. При $r=2-\sqrt3$: $f(B_r)=\bigcup_{\rho<r}f(B_\rho)$ --- объединение возрастающей цепочки выпуклых множеств,
оно выпукло (любые две точки лежат в одном $f(B_\rho)$ вместе с отрезком). Значит, $R_C(\Sc)\ge2-\sqrt3$.

\textbf{Шаг 2: для функции Кёбе.} Вычисление (непосредственно, <<как легко проверить>>):
\[
 1+\frac{zk''(z)}{k'(z)}=1+\frac{z(4+2z)}{1-z^2}=\frac{1-z^2+4z+2z^2}{1-z^2}=\frac{z^2+4z+1}{1-z^2}.
\]
При $z=-r$ получаем $\frac{r^2-4r+1}{1-r^2}$, что отрицательно при $r\in(2-\sqrt3,1)$ (между корнями $2\pm\sqrt3$, а $2+\sqrt3>1$). Таким образом,
касательная к $k(C_r)$ в точке $k(-r)$ вращается <<в обратную сторону>>. Докажем прямо, что $k(B_r)$ не выпукла.

\begin{detail}{$k(B_r)$ не выпукла при $2-\sqrt3<r<1$}
(а) \emph{$k(B_\rho)\cap\RR=k((-\rho,\rho))$ для $\rho<1$.} Коэффициенты $k$ вещественны, так что $k(\bar z)=\overline{k(z)}$. Если $k(z)\in\RR$, то
$k(\bar z)=k(z)$, и по инъективности $z=\bar z\in\RR$. Функция $k$ на $(-1,1)$ строго возрастает ($k'(x)=\frac{1+x}{(1-x)^3}>0$), поэтому
$k(B_\rho)\cap\RR=(k(-\rho),k(\rho))$, в частности $\inf\bigl(k(B_\rho)\cap\RR\bigr)=k(-\rho)=-\frac{\rho}{(1+\rho)^2}$.

(б) \emph{Вторая производная вдоль окружности.} Пусть $u(\eps):=\Re k(-re^{i\eps})$. Так как $\overline{k(-re^{i\eps})}=k(-re^{-i\eps})$, функция $u$ чётна,
$u(0)=k(-r)$. Для $z(\eps)=-re^{i\eps}$ имеем $z'=iz$ и
\[
 \frac{d}{d\eps}k(z)=izk'(z),\qquad \frac{d^2}{d\eps^2}k(z)=i\cdot iz\,k'(z)+iz\cdot k''(z)\,iz=-zk'(z)\Bigl(1+\frac{zk''(z)}{k'(z)}\Bigr).
\]
При $\eps=0$, $z=-r$: $-zk'(z)=rk'(-r)=\frac{r(1-r)}{(1+r)^3}>0$, а скобка равна $\frac{r^2-4r+1}{1-r^2}<0$. Значит, $u''(0)<0$, $u'(0)=0$ (чётность), и
найдётся $\eps\in(0,\pi)$ с $u(\eps)<u(0)=k(-r)$.

(в) \emph{Противоречие с выпуклостью.} Функция $\rho\mapsto\Re k(-\rho e^{i\eps})+\frac{\rho}{(1+\rho)^2}$ непрерывна на $(0,1)$ и в точке $\rho=r$ отрицательна; поэтому
найдётся $\rho<r$ с $\Re k(-\rho e^{i\eps})<-\frac{r}{(1+r)^2}=k(-r)$. Точки $w=k(-\rho e^{i\eps})$ и $\bar w=k(-\rho e^{-i\eps})$ лежат в $k(B_r)$, а середина отрезка
$[w,\bar w]$ --- вещественное число $\Re w<k(-r)=\inf\bigl(k(B_r)\cap\RR\bigr)$, т.е. $\Re w\notin k(B_r)$. Значит, $k(B_r)$ не выпукла.
\end{detail}

Итак, при каждом $r>2-\sqrt3$ найдётся $f=k\in\Sc$ с невыпуклой $f(B_r)$; вместе с шагом 1 получаем $R_C(\Sc)=2-\sqrt3$.
\end{proof}

\begin{remarkbox}{Замечание (опечатка на слайде 321)}
На слайде написано $\frac{zk''(z)}{k'(z)}+1=\frac{z^2-4z+1}{1-z^2}$. Для $k(z)=z/(1-z)^2$ правильно $\frac{z^2+4z+1}{1-z^2}$ (см. вычисление выше); формула слайда
верна для функции $\frac{z}{(1+z)^2}=-k(-z)$. На вывод это не влияет: значение $\frac{r^2-4r+1}{1-r^2}$ получается при $z=-r$ для $k$ (при $z=r$ для
$\frac{z}{(1+z)^2}$).
\end{remarkbox}

\begin{remarkbox}{Радиус звёздности (доп. материал, предложение 13.14, без доказательства)}
$R_S(\Sc)=\Th(\pi/4)\approx0{,}66$. Звёздность $f(B_r)$ относительно $0$ равносильна тому, что $\arg f(re^{i\vartheta})$ не убывает по $\vartheta$, т.е.
$\frac{\partial}{\partial\vartheta}\arg f(re^{i\vartheta})=\Re\frac{zf'(z)}{f(z)}\ge0$ на $C_r$ (вычисление то же, что в упражнении выше, с $\log f(z)/z$ вместо $\log f'$);
эквивалентно $\bigl|\arg\frac{zf'(z)}{f(z)}\bigr|\le\frac\pi2$. Заметим, что $R_C(\Sc)<R_S(\Sc)$, что естественно: выпуклая область, содержащая $0$, звёздна относительно $0$.
\end{remarkbox}

\begin{remarkbox}{Коротко: план ответа на экзамене}
\begin{enumerate}
\item Определение $R_C(\Sc)$; ответ $2-\sqrt3$.
\item Касательная к $f(C_r)$ в $f(z)$ имеет угол $\tau=\vartheta+\frac\pi2+\arg f'(z)$; $\tau'(\vartheta)=1+\Re\frac{zf''}{f'}$.
\item Принцип выпуклости: если $\tau'>0$ (полный поворот $2\pi$), кривая лежит по одну сторону от каждой касательной, область --- пересечение полуплоскостей.
\item Из леммы 13.5 (билет 14): $\Re\frac{zf''}{f'}\ge\frac{2r^2-4r}{1-r^2}$, т.е. $1+\Re\frac{zf''}{f'}\ge\frac{r^2-4r+1}{1-r^2}>0$ при $r<2-\sqrt3$.
\item Функция Кёбе: $1+\frac{zk''}{k'}=\frac{z^2+4z+1}{1-z^2}$, при $z=-r$ отрицательно для $r>2-\sqrt3$.
\item Невыпуклость $k(B_r)$: симметрия относительно $\RR$ и $u''(0)<0$ для $u(\eps)=\Re k(-re^{i\eps})$.
\end{enumerate}
\end{remarkbox}

% -------------------------------------------------------------------
\ticketx{41*}{Оценка Литтлвуда в задаче Бибербаха}
\label{t41:begin}

\begin{defs}
\textbf{Класс $\Sc$} $=\{f\in\mathcal U(\DD):\ f(0)=0,\ f'(0)=1\}$, $f(z)=z+\sum_{n\ge2}a_nz^n$ (билет 13).

\textbf{Задача (гипотеза) Бибербаха:} $|a_n|\le n$ для всех $f\in\Sc$ и $n\ge2$ (доказана де Бранжем в 1984 г.); для функции Кёбе
$k(z)=\sum nz^n$ имеем $a_n=n$. Лемма 13.4: $|a_2|\le2$.

\textbf{Вторая теорема растяжения} (теорема 13.5, билет 14): для $f\in\Sc$ и $|z|=r<1$: $\frac{r}{(1+r)^2}\le|f(z)|\le\frac{r}{(1-r)^2}$.

\textbf{Лемма 13.3} (билет 13): для $f\in\Sc$ функция $h(z)=\sqrt{f(z^2)}$ (ветвь с $h'(0)=1$) лежит в $\Sc$ и нечётна.

\textbf{Формулы для коэффициентов.} Для $F(z)=\sum_{n\ge0}c_nz^n\in H(\DD)$ и $0<r<1$:
\[
 c_n=\frac1{2\pi i}\oint_{|z|=r}\frac{F(z)}{z^{n+1}}\,dz=\frac1{2\pi r^n}\int_0^{2\pi}F(re^{it})e^{-int}\,dt,
\qquad
 \frac1{2\pi}\int_0^{2\pi}|F(re^{it})|^2\,dt=\sum_{n\ge0}|c_n|^2r^{2n}
\]
(формула Коши и равенство Парсеваля: ряд $\sum c_nr^ne^{int}$ сходится равномерно по $t$, и $\int_0^{2\pi}e^{i(n-m)t}dt=2\pi\delta_{nm}$).

\textbf{Площадь образа.} Если $h\in\mathcal U(B_r)$, то $m_2(h(B_r))=\iint_{B_r}|h'(z)|^2\,dx\,dy$ (замена переменных: якобиан отображения
$h=u+iv$ равен $u_xv_y-u_yv_x=u_x^2+v_x^2=|h'|^2$ по условиям Коши--Римана).
\end{defs}

\begin{thmbox}{Предложение 13.11 (оценка Литтлвуда, 1925)}
Если $f\in\Sc$, то $|a_n|<e\,n$ при всех $n\ge2$ (и $|a_1|=1<e$).
\end{thmbox}

\subsection{Почему не хватает второй теоремы растяжения}

Из формулы Коши и оценки $\max_{|z|=r}|f(z)|\le\frac{r}{(1-r)^2}$ получаем (как при выводе неравенств Коши)
\[
 |a_n|\le\frac1{r^n}\cdot\frac{r}{(1-r)^2}=\frac1{r^{n-1}(1-r)^2},\qquad 0<r<1,\ n\ge2 .
\]
\begin{detail}{Минимизация и результат $|a_n|\le Cn^2$}
$\frac{d}{dr}\log\bigl(r^{n-1}(1-r)^2\bigr)=\frac{n-1}r-\frac2{1-r}=0\iff(n-1)(1-r)=2r\iff r=\frac{n-1}{n+1}$; это точка максимума функции
$r^{n-1}(1-r)^2$ на $(0,1)$ (на концах она равна $0$), т.е. минимума правой части. Подставляя, $1-r=\frac2{n+1}$, $\frac1r=1+\frac2{n-1}$:
\[
 |a_n|\le\frac{(n+1)^2}{4}\Bigl(1+\frac2{n-1}\Bigr)^{n-1}<\frac{e^2}{4}(n+1)^2\le e^2n^2,
\]
так как $(1+\frac2m)^m<e^2$. Эта оценка имеет порядок $n^2$, а не $n$ (выражение не меньше $\frac{(n+1)^2}4$), т.е. существенно грубее гипотезы.
\end{detail}
Причина: мы оценили $|f|$ на окружности её \emph{максимумом} $\frac{r}{(1-r)^2}$, тогда как для формулы Коши достаточно \emph{среднего} значения $|f|$,
а оно растёт медленнее --- как $\frac{r}{1-r}$ (лемма 13.6).

\subsection{Лемма 13.6 (оценка интегрального среднего)}

\begin{thmbox}{Лемма 13.6}
Пусть $f\in\Sc$ и $r\in[0,1)$. Тогда
\[
 \frac1{2\pi}\int_0^{2\pi}|f(re^{it})|\,dt\le\frac{r}{1-r}.
\]
\end{thmbox}

\begin{proof}
При $r=0$ обе части равны $0$. Пусть $h\in\Sc$ --- функция из леммы 13.3: $h(z)^2=f(z^2)$; пусть $h(z)=\sum_{n\ge1}c_nz^n$ ($c_1=1$).

\textbf{Шаг 1: оценка $|h|$.} Для $|z|=\rho<1$ по второй теореме растяжения, применённой к $f$ в точке $z^2$ ($|z^2|=\rho^2$):
\[
 |h(z)|^2=|f(z^2)|\le\frac{\rho^2}{(1-\rho^2)^2},\qquad\text{т.е.}\quad |h(z)|\le\frac{\rho}{1-\rho^2}.
\]
Так как $\rho\mapsto\frac{\rho}{1-\rho^2}$ возрастает, при $|z|<\rho$ также $|h(z)|<R(\rho):=\frac{\rho}{1-\rho^2}$. Итак, $h$ отображает $B_\rho$ конформно
(однолистно) на область $G_\rho\subset B(0,R(\rho))$, и $m_2(G_\rho)\le\pi R(\rho)^2$.

\textbf{Шаг 2: площадь $G_\rho$ через коэффициенты.} Так как $h$ однолистна, в полярных координатах $z=se^{it}$:
\begin{align*}
 m_2(G_\rho)&=\iint_{B_\rho}|h'|^2\,dm_2=\int_0^\rho\Bigl(\int_0^{2\pi}\Bigl|\sum_{n\ge1}nc_ns^{n-1}e^{i(n-1)t}\Bigr|^2dt\Bigr)s\,ds\\
 &=\int_0^\rho2\pi\sum_{n\ge1}n^2|c_n|^2s^{2n-1}\,ds=\pi\sum_{n\ge1}n|c_n|^2\rho^{2n}
\end{align*}
(внутренний интеграл --- по Парсевалю для $h'$; почленное интегрирование по $s$ ряда с неотрицательными членами законно по теореме о монотонной сходимости;
$\int_0^\rho s^{2n-1}ds=\frac{\rho^{2n}}{2n}$). Сравнивая с шагом 1:
\[
 \pi\sum_{n\ge1}n|c_n|^2\rho^{2n}\le\pi\frac{\rho^2}{(1-\rho^2)^2},\qquad\text{т.е.}\qquad \sum_{n\ge1}n|c_n|^2\rho^{2n-1}\le\frac{\rho}{(1-\rho^2)^2}
 \quad(0<\rho<1).
\]

\textbf{Шаг 3: интегрирование по $\rho$.} Проинтегрируем последнее неравенство по $\rho\in[0,r]$. Слева (почленно, члены неотрицательны):
$\sum_n n|c_n|^2\frac{r^{2n}}{2n}=\frac12\sum_n|c_n|^2r^{2n}$. Справа:
\[
 \int_0^r\frac{\rho\,d\rho}{(1-\rho^2)^2}=\Bigl[\frac1{2(1-\rho^2)}\Bigr]_0^r=\frac12\Bigl(\frac1{1-r^2}-1\Bigr)=\frac12\cdot\frac{r^2}{1-r^2}.
\]
Итак, $\sum_{n\ge1}|c_n|^2r^{2n}\le\frac{r^2}{1-r^2}$, и по равенству Парсеваля
\[
 \frac1{2\pi}\int_0^{2\pi}|h(re^{i\vartheta})|^2\,d\vartheta=\sum_{n\ge1}|c_n|^2r^{2n}\le\frac{r^2}{1-r^2}.
\]

\textbf{Шаг 4: возвращение к $f$.} $|h(re^{i\vartheta})|^2=|f(r^2e^{2i\vartheta})|$. Замена $t=2\vartheta$ и $2\pi$-периодичность дают
\[
 \frac1{2\pi}\int_0^{2\pi}|f(r^2e^{2i\vartheta})|\,d\vartheta=\frac1{4\pi}\int_0^{4\pi}|f(r^2e^{it})|\,dt=\frac1{2\pi}\int_0^{2\pi}|f(r^2e^{it})|\,dt.
\]
Следовательно, $\frac1{2\pi}\int_0^{2\pi}|f(r^2e^{it})|\,dt\le\frac{r^2}{1-r^2}$. Так как $r^2$ пробегает всё $[0,1)$, заменяя $r^2$ на $r$, получаем
утверждение леммы.
\end{proof}

\begin{remarkbox}{Сравнение}
Для функции Кёбе оценка леммы 13.6 точна по порядку: $|k(re^{it})|=\frac{r}{|1-re^{it}|^2}$ и $\frac1{2\pi}\int_0^{2\pi}\frac{dt}{|1-re^{it}|^2}=\frac1{1-r^2}$ (ядро Пуассона), так что среднее равно
$\frac{r}{1-r^2}$ --- того же порядка $\frac1{1-r}$, что и в лемме. Максимум же $|k|$ на $C_r$ равен $\frac{r}{(1-r)^2}$, порядка $\frac1{(1-r)^2}$.
\end{remarkbox}

\subsection{Доказательство оценки Литтлвуда}

\begin{proof}[Доказательство предложения 13.11]
Пусть $n\ge2$, $0<r<1$. По формуле Коши для коэффициентов и лемме 13.6
\[
 |a_n|=\Bigl|\frac1{2\pi r^n}\int_0^{2\pi}f(re^{it})e^{-int}\,dt\Bigr|\le\frac1{r^n}\cdot\frac1{2\pi}\int_0^{2\pi}|f(re^{it})|\,dt\le\frac1{r^n}\cdot\frac{r}{1-r}
 =\frac1{r^{n-1}(1-r)}.
\]
\begin{detail}{Выбор $r$}
$\frac{d}{dr}\log\bigl(r^{n-1}(1-r)\bigr)=\frac{n-1}{r}-\frac1{1-r}=0\iff r=1-\frac1n$ --- точка максимума $r^{n-1}(1-r)$ на $(0,1)$, т.е. минимума правой части.
При $r=1-\frac1n=\frac{n-1}n$: $1-r=\frac1n$, $\frac1{r^{n-1}}=\bigl(\frac{n}{n-1}\bigr)^{n-1}=\bigl(1+\frac1{n-1}\bigr)^{n-1}$.
\end{detail}
Подставляя $r=1-\frac1n$:
\[
 |a_n|\le n\Bigl(1+\frac1{n-1}\Bigr)^{n-1}<e\,n,
\]
поскольку последовательность $(1+\frac1m)^m$ строго возрастает и стремится к $e$, так что $(1+\frac1m)^m<e$ при всех $m\ge1$.
(Строгое возрастание: $\log(1+\frac1m)^m=m\log(1+\frac1m)=\int_0^1\frac{dx}{1+x/m}$, а подынтегральная функция строго возрастает по $m$ при $x\in(0,1]$.)
\end{proof}

\begin{remarkbox}{Замечания}
\begin{itemize}
\item При $n=2$ оценка даёт $|a_2|\le2\cdot2=4$, т.е. хуже леммы 13.4 ($|a_2|\le2$); ценность оценки Литтлвуда --- в правильном \emph{линейном} порядке роста по $n$.
\item Ключевая идея: перейти к нечётной функции $h=\sqrt{f(z^2)}$, для которой квадрат модуля $|h|^2=|f(z^2)|$ и $L^2$-норма $h$ контролируется
через площадь образа (теорема площадей <<изнутри>>), а площадь --- через вторую теорему растяжения.
\item Литтлвуд и Пэли (1932) предположили, что для нечётных $h\in\Sc$ коэффициенты ограничены ($|c_n|\le1$, гипотеза Литтлвуда--Пэли); она неверна
(Фекете--Сегё), но Литтлвуд и Пэли доказали оценку $|c_n|<C$ с абсолютной константой. Окончательно гипотеза Бибербаха доказана де Бранжем (1984).
\end{itemize}
\end{remarkbox}

\begin{remarkbox}{Коротко: план ответа}
\begin{enumerate}
\item Формулировка: $f\in\Sc\Rightarrow|a_n|<en$.
\item Неравенство Коши с максимумом $\frac{r}{(1-r)^2}$ даёт лишь $|a_n|\le Cn^2$.
\item Лемма 13.6: среднее $|f|$ на $C_r$ не больше $\frac r{1-r}$.
\item Доказательство: $h=\sqrt{f(z^2)}\in\Sc$, $|h|\le\frac{\rho}{1-\rho^2}$ на $B_\rho$; площадь $\pi\sum n|c_n|^2\rho^{2n}\le\pi\frac{\rho^2}{(1-\rho^2)^2}$.
\item Интегрирование по $\rho$: $\sum|c_n|^2r^{2n}\le\frac{r^2}{1-r^2}$, Парсеваль, $|h|^2=|f(z^2)|$.
\item $|a_n|\le\frac1{r^{n-1}(1-r)}$, $r=1-\frac1n$: $|a_n|\le n(1+\frac1{n-1})^{n-1}<en$.
\end{enumerate}
\end{remarkbox}

% ===== part 3 =====
% ===================================================================
%  Часть 3: билеты 16--20
% ===================================================================

% -------------------------------------------------------------------
\ticket{16}{Существование радиальных пределов почти всюду на единичной окружности у функции $f$, конформно отображающей единичный круг на некоторую ограниченную область}
\label{t16:begin}

\begin{defs}
\textbf{Однолистные функции.} Для области $D\subset\CC$ через $H(D)$ обозначается множество голоморфных в $D$ функций;
$f\in H(D)$ \emph{локально однолистна} ($f\in U_\ell(D)$), если она инъективна в некоторой окрестности каждой точки
(равносильно: $f'\neq0$ всюду в $D$), и \emph{однолистна} ($f\in U(D)$), если она инъективна во всей $D$.
Если $f\in U(D)$, то $f(D)$ --- область, $f^{-1}\colon f(D)\to D$ голоморфна, и $f$ --- \emph{конформное отображение}
$D$ на $f(D)$; множество конформных отображений $D$ на $G$ обозначается $\mathrm{Conf}(D,G)$.

\textbf{Класс $\Sc$:} $\Sc=\{f\in U(\DD):\ f(0)=0,\ f'(0)=1\}$. Если $f\in U(\DD)$, то $F=\dfrac{f-f(0)}{f'(0)}\in\Sc$
(здесь $f'(0)\ne0$ по критерию локальной однолистности). Пример: функция Кёбе $k(z)=\dfrac{z}{(1-z)^2}\in\Sc$ отображает $\DD$
на $\CC\setminus(-\infty,-1/4]$ --- неограниченную область.

\textbf{Радиальный предел.} Пусть $f\colon\DD\to\CC$, $\zeta\in\TT=\{|\zeta|=1\}$. Если существует конечный предел
$\lim_{r\to1-}f(r\zeta)$, то он называется \emph{радиальным пределом} $f$ в точке $\zeta$ и обозначается $f(\zeta)$.
Например, у функции Кёбе $k(r\zeta)\to k(\zeta)$ при $\zeta\neq1$, а $k(r)\to+\infty$ при $r\to1-$, т.\,е. конечного
радиального предела в точке $1$ нет.

\textbf{«Почти всюду на $\TT$».} Мера на $\TT$ --- длина дуги: множество $E\subset\TT$ имеет меру нуль, если множество
$\{t\in[0,2\pi):\ e^{it}\in E\}$ имеет нулевую меру Лебега на прямой. Свойство выполнено \emph{для почти всех $\zeta\in\TT$}, если
множество точек, где оно нарушено, имеет меру нуль.

\textbf{Длина кривой.} Если $\gamma\colon[\alpha,\beta]\to\CC$ кусочно гладкий путь, а $h$ голоморфна в окрестности $[\gamma]$, то
$\ell(h\circ\gamma)=\int_\alpha^\beta|h'(\gamma(t))|\,|\gamma'(t)|\,dt=\int_\gamma|h'(z)|\,|dz|$.
Для луча $t\mapsto re^{it}$, $r\in[0,1)$, при фиксированном $t$: длина образа $h([0,e^{it}))$ равна $\int_0^1|h'(re^{it})|\,dr$
(может быть $+\infty$).

\textbf{Формула площади.} Если $h\in U(\Omega)$ и $E\subset\Omega$ измеримо, то $m_2(h(E))=\int_E|h'(z)|^2\,dm_2(z)$
($m_2$ --- плоская мера Лебега): якобиан голоморфного отображения $h=u+iv$ равен $u_x v_y-u_yv_x=u_x^2+v_x^2=|h'|^2$
(условия Коши--Римана), и работает обычная формула замены переменных для инъективного гладкого отображения.

\textbf{Пространства Харди (для ориентира).} $f\in H(\DD)$ принадлежит $H^p$ ($0<p<\infty$), если
$\sup_{0\le r<1}\frac1{2\pi}\int_0^{2\pi}|f(re^{it})|^p\,dt<\infty$; $H^\infty$ --- ограниченные голоморфные в $\DD$ функции.
Конформное отображение $\DD$ на \emph{ограниченную} область принадлежит $H^\infty$; Предложение 13.12 ниже показывает, что
\emph{любая} $f\in U(\DD)$ лежит в $H^{2/5}$.
\end{defs}

\subsection{Используемые факты}

\begin{thmbox}{Теорема 13.5 (вторая теорема растяжения; см. билет 14)}
Если $f\in U(\DD)$, то для всех $z\in\DD$
\[
|f'(0)|\frac{|z|}{(1+|z|)^2}\le |f(z)-f(0)|\le|f'(0)|\frac{|z|}{(1-|z|)^2}.
\]
В частности, для $f\in\Sc$: $|f(z)|\le |z|/(1-|z|)^2$.
\end{thmbox}

\begin{thmbox}{Лемма А (формула Парсеваля и площадь образа круга)}
Пусть $g(z)=\sum_{n\ge0}b_nz^n$ голоморфна в $\DD$. Тогда при $0\le r<1$
\[
\frac1{2\pi}\int_0^{2\pi}|g(re^{it})|^2\,dt=\sum_{n\ge0}|b_n|^2r^{2n}. \tag{16.1}
\]
Если, кроме того, $g\in U(\DD)$, то при $0<\rho\le1$ (при $\rho=1$ обе части могут быть равны $+\infty$)
\[
m_2\bigl(g(B(0,\rho))\bigr)=\pi\sum_{n\ge1}n|b_n|^2\rho^{2n}. \tag{16.2}
\]
\end{thmbox}
\begin{proof}
Ряд $g(re^{it})=\sum b_nr^ne^{int}$ сходится равномерно по $t$ (так как $\sum|b_n|r^n<\infty$ при $r<1$), поэтому его можно
умножить на сопряжённый ряд и проинтегрировать почленно; из $\frac1{2\pi}\int_0^{2\pi}e^{i(n-m)t}dt=\delta_{nm}$ получаем (16.1).

Для (16.2) применяем формулу площади и полярные координаты: при $\rho<1$
\begin{align*}
m_2(g(B(0,\rho)))&=\int_{B(0,\rho)}|g'|^2dm_2=\int_0^\rho\Bigl(\int_0^{2\pi}|g'(re^{it})|^2dt\Bigr)r\,dr\\
&=\int_0^\rho 2\pi\sum_{n\ge1}n^2|b_n|^2r^{2n-2}\,r\,dr=\pi\sum_{n\ge1}n|b_n|^2\rho^{2n},
\end{align*}
где в третьем равенстве использована (16.1) для $g'(z)=\sum_{n\ge1}nb_nz^{n-1}$, а в последнем --- почленное интегрирование
ряда с неотрицательными членами (теорема о монотонной сходимости). При $\rho=1$ переходим к пределу $\rho\to1-$ (непрерывность меры
снизу слева, монотонная сходимость справа).
\end{proof}

\begin{thmbox}{Лемма Б (конечная длина образа радиуса $\Rightarrow$ радиальный предел)}
Пусть $F\in H(\DD)$, $t\in\RR$ и $\int_0^1|F'(re^{it})|\,dr<\infty$. Тогда существует конечный предел
$\lim_{r\to1-}F(re^{it})$.
\end{thmbox}
\begin{proof}
При $0\le r_1<r_2<1$ по формуле Ньютона--Лейбница $F(r_2e^{it})-F(r_1e^{it})=\int_{r_1}^{r_2}F'(re^{it})e^{it}dr$, откуда
$|F(r_2e^{it})-F(r_1e^{it})|\le\int_{r_1}^{r_2}|F'(re^{it})|dr\to0$ при $r_1,r_2\to1-$ (хвост сходящегося интеграла).
По критерию Коши предел существует и конечен.
\end{proof}

\begin{remarkbox}{Замечание}
На слайде этот шаг объясняется так: образ радиуса под действием однолистной $g$ --- простая полуоткрытая кривая конечной длины,
и потому у неё есть конечный конец. Как видно из Леммы Б, однолистность здесь не нужна: достаточно конечности длины образа
радиуса. Однолистность нужна в другом месте --- для оценки площади (формула (16.2)).
\end{remarkbox}

\subsection{Случай ограниченной области: метод «длина--площадь»}

Для ограниченной области (это и есть формулировка билета) есть очень короткое доказательство, которое полезно знать и которое
показывает идею, лежащую в основе общего случая.

\begin{thmbox}{Теорема В (радиальные пределы для ограниченной области)}
Пусть $f\in\mathrm{Conf}(\DD,G)$, где $G\subset\CC$ --- ограниченная область. Тогда для почти всех $t\in[0,2\pi)$ образ радиуса
$f([0,e^{it}))$ имеет конечную длину; в частности, для почти всех $\zeta\in\TT$ существует конечный радиальный предел
$f(\zeta)=\lim_{r\to1-}f(r\zeta)$, и при этом $f(\zeta)\in\partial G$. Точнее,
\[
\int_0^{2\pi}\Bigl(\int_{1/2}^1|f'(re^{it})|\,dr\Bigr)^2dt\le \ln 2\cdot m_2(G)<\infty .
\]
\end{thmbox}
\begin{proof}
По формуле площади и в полярных координатах
\[
\int_0^{2\pi}\!\!\int_0^1|f'(re^{it})|^2\,r\,dr\,dt=\int_\DD|f'|^2\,dm_2=m_2(f(\DD))=m_2(G)<\infty .
\]
При фиксированном $t$ по неравенству Коши--Буняковского (записываем $|f'|=\frac1{\sqrt r}\cdot|f'|\sqrt r$)
\[
\Bigl(\int_{1/2}^1|f'(re^{it})|\,dr\Bigr)^2\le\int_{1/2}^1\frac{dr}{r}\cdot\int_{1/2}^1|f'(re^{it})|^2\,r\,dr
=\ln2\int_{1/2}^1|f'(re^{it})|^2\,r\,dr .
\]
Интегрируя по $t\in[0,2\pi]$ (функции измеримы по теореме Тонелли), получаем требуемое неравенство. Неотрицательная функция,
интеграл от которой конечен, конечна почти всюду; значит, $\int_{1/2}^1|f'(re^{it})|dr<\infty$ для почти всех $t$, а тогда и
$\int_0^1|f'(re^{it})|dr<\infty$ (на $[0,1/2]$ подынтегральная функция непрерывна). По Лемме Б для таких $t$ существует
конечный предел $w=\lim_{r\to1-}f(re^{it})$.

\emph{Почему $w\in\partial G$.} Ясно, что $w\in\clos G$. Если бы $w\in G$, то, так как $f^{-1}\colon G\to\DD$ непрерывна,
$re^{it}=f^{-1}(f(re^{it}))\to f^{-1}(w)\in\DD$ при $r\to1-$, что невозможно, ведь $|re^{it}|\to1$. Значит, $w\in\clos G\setminus G=\partial G$.
\end{proof}

\begin{remarkbox}{Замечание (почему одной ограниченности $|f|$ здесь мало и что даёт однолистность)}
В доказательстве ограниченность $G$ использовалась только через конечность площади $m_2(G)=\int_\DD|f'|^2dm_2$
(конечность интеграла Дирихле). Для неограниченных $G$ площадь может быть бесконечной (функция Кёбе), и нужен дополнительный приём ---
переход к функции $g(z)=(f(z^5))^{1/5}$, которая растёт медленнее (Предложение 13.12).
\end{remarkbox}

\subsection{Общий случай: Предложение 13.12}

\begin{thmbox}{Лемма Г (преобразование корня)}
Пусть $f\in\Sc$ и $k\in\NN$. Тогда существует единственная функция $g\in H(\DD)$ с $g'(0)=1$ и $g(z)^k=f(z^k)$; она
однолистна, т.\,е. $g\in\Sc$, и $g(\omega z)=\omega g(z)$ для всех $\omega$ с $\omega^k=1$. Обозначение: $g(z)=(f(z^k))^{1/k}$.
\end{thmbox}
\begin{proof}
(Это упражнение со слайда 316; при $k=2$ это Лемма 13.3 из билета 13.)
Так как $f(0)=0$, то $f(z)=z\varphi(z)$, $\varphi\in H(\DD)$, $\varphi(0)=f'(0)=1$. Если $z\neq0$, то $f(z)\ne f(0)=0$ по
однолистности, поэтому $\varphi$ не имеет нулей в $\DD$. Функция $\varphi(z^k)$ голоморфна и не обращается в нуль в односвязной
области $\DD$, поэтому у неё есть голоморфный корень $k$-й степени: существует $q\in H(\DD)$, $q(z)^k=\varphi(z^k)$, $q(0)=1$
(берём $q=\exp(\frac1k\log\varphi(z^k))$ с ветвью логарифма, равной нулю в нуле). Положим $g(z)=zq(z)$. Тогда
$g(z)^k=z^k\varphi(z^k)=f(z^k)$ и $g'(0)=q(0)=1$.

\emph{Симметрия.} Пусть $\omega^k=1$. Тогда $q(\omega z)^k=\varphi(\omega^kz^k)=\varphi(z^k)=q(z)^k$, т.\,е. непрерывная в связной
области $\DD$ функция $q(\omega z)/q(z)$ принимает значения в конечном множестве корней $k$-й степени из $1$; значит, она постоянна и
равна своему значению в нуле, т.\,е. $1$. Отсюда $g(\omega z)=\omega zq(z)=\omega g(z)$.

\emph{Однолистность.} Пусть $g(z_1)=g(z_2)$. Тогда $f(z_1^k)=g(z_1)^k=g(z_2)^k=f(z_2^k)$, и по однолистности $f$ имеем
$z_1^k=z_2^k$, т.\,е. $z_2=\omega z_1$ с $\omega^k=1$. Тогда $g(z_1)=g(z_2)=\omega g(z_1)$. Если $\omega\ne1$, то $g(z_1)=0$,
т.\,е. $f(z_1^k)=0$, откуда $z_1=0$ и $z_2=0$. Если $\omega=1$, то $z_1=z_2$. Итак, $g\in U(\DD)$.

\emph{Единственность.} Если $\tilde g$ --- другая такая функция, то $(\tilde g/g)^k=1$ на $\DD\setminus\{0\}$, отношение
$\tilde g/g$ голоморфно в $\DD$ (у $g$ лишь простой нуль в $0$, у $\tilde g$ тоже нуль в $0$) и постоянно; в нуле оно равно
$\tilde g'(0)/g'(0)=1$.
\end{proof}

\begin{thmbox}{Предложение 13.12}
Пусть $f$ конформно отображает $\DD$ на некоторую область в $\CC$ (не обязательно ограниченную). Тогда для почти всех $\zeta\in\TT$
существует конечный предел $f(\zeta):=\lim_{r\to1-}f(r\zeta)$. Кроме того, при всех $r\in[0,1)$
\[
\frac1{2\pi}\int_0^{2\pi}\bigl|f(re^{it})-f(0)\bigr|^{2/5}\,dt\le5\,|f'(0)|^{2/5}.
\]
\end{thmbox}

\begin{proof}
\textbf{Шаг 0: сведение к классу $\Sc$.} Положим $F=(f-f(0))/f'(0)\in\Sc$. Радиальный предел у $f=f(0)+f'(0)F$ существует
тогда и только тогда, когда он существует у $F$, и $|f(re^{it})-f(0)|^{2/5}=|f'(0)|^{2/5}|F(re^{it})|^{2/5}$. Поэтому
достаточно доказать: \emph{если $f\in\Sc$, то радиальные пределы существуют почти всюду, и
$\frac1{2\pi}\int_0^{2\pi}|f(re^{it})|^{2/5}dt\le5$.} Далее $f\in\Sc$.

\textbf{Шаг 1: вспомогательная функция.} Пусть $g(z)=(f(z^5))^{1/5}$ --- функция из Леммы Г при $k=5$; $g\in\Sc$,
$g(z)=\sum_{n\ge1}b_nz^n$, $b_1=1$. По второй теореме растяжения
\[
|g(z)|=|f(z^5)|^{1/5}\le\Bigl(\frac{|z|^5}{(1-|z|^5)^2}\Bigr)^{1/5}=\frac{|z|}{(1-|z|^5)^{2/5}}. \tag{16.3}
\]
\emph{Смысл приёма:} $|f(z)|$ может расти как $(1-|z|)^{-2}$, а $|g(z)|$ растёт не быстрее $(1-|z|)^{-2/5}$.

\textbf{Шаг 2: оценка интеграла от $|g'|^2$ по окружности.} Пусть $0<r<1$. По (16.1) для $g'$:
\[
\frac{1-r}{2\pi}\int_0^{2\pi}|g'(re^{it})|^2dt=(1-r)\sum_{n\ge1}n^2|b_n|^2r^{2n-2}\le\sum_{n\ge1}n|b_n|^2r^{n-1}, \tag{16.4}
\]
так как $(1-r)n^2r^{2n-2}=n r^{n-1}\cdot\bigl(n(1-r)r^{n-1}\bigr)$ и
\[
n(1-r)r^{n-1}\le(1-r)(1+r+\dots+r^{n-1})=1-r^n\le1
\]
(каждое из $n$ слагаемых $r^j\ge r^{n-1}$). По (16.2) с $\rho=\sqrt r$ последняя сумма в (16.4) равна
$\dfrac{m_2(Q_r)}{\pi r}$, где $Q_r=g(B(0,\sqrt r))$. Так как $Q_r\subset\clos{B(0,M_r)}$, $M_r=\max_{|z|\le\sqrt r}|g(z)|$, то
$m_2(Q_r)\le\pi M_r^2$, и по (16.3)
\[
\sum_{n\ge1}n|b_n|^2r^{n-1}\le\frac1r\Bigl(\frac{\sqrt r}{(1-r^{5/2})^{2/5}}\Bigr)^2=\frac1{(1-r^{5/2})^{4/5}}\le\frac1{(1-r)^{4/5}} \tag{16.5}
\]
(последнее --- так как $r^{5/2}\le r$). Вместе с (16.4):
\[
\frac1{2\pi}\int_0^{2\pi}|g'(re^{it})|^2dt\le\frac1{(1-r)^{9/5}},\qquad 0<r<1. \tag{16.6}
\]

\textbf{Шаг 3: интегрирование по $r$ с весом.} Умножим (16.6) на $(1-r)^{9/10}$ и проинтегрируем по $r\in(0,1)$ (теорема Тонелли):
\[
\frac1{2\pi}\int_0^{2\pi}\!\!\int_0^1(1-r)^{9/10}|g'(re^{it})|^2\,dr\,dt\le\int_0^1\frac{dr}{(1-r)^{9/10}}=10 .
\]

\textbf{Шаг 4: Коши--Буняковский по радиусу.} При каждом $t$, записывая $|g'|=(1-r)^{-9/20}\cdot(1-r)^{9/20}|g'|$,
\[
\Bigl(\int_0^1|g'(re^{it})|\,dr\Bigr)^2\le\int_0^1\frac{dr}{(1-r)^{9/10}}\cdot\int_0^1(1-r)^{9/10}|g'(re^{it})|^2dr
=10\int_0^1(1-r)^{9/10}|g'(re^{it})|^2dr .
\]
Интегрируя по $t$ и используя шаг 3:
\[
\frac1{2\pi}\int_0^{2\pi}\Bigl(\int_0^1|g'(re^{it})|\,dr\Bigr)^2dt\le10\cdot10=100 .
\]
Следовательно, внутренний интеграл $\int_0^1|g'(re^{it})|dr$ (длина образа радиуса $[0,e^{it})$ под действием $g$) конечен для
всех $t$ вне некоторого множества $N\subset[0,2\pi)$ нулевой меры. По Лемме Б при $t\notin N$ существует конечный предел
$\lim_{r\to1-}g(re^{it})$.

\textbf{Шаг 5: возвращение к $f$.} Пусть $\theta\in[0,2\pi)$. Положим $t=\theta/5$. Тогда при $r\to1-$
\[
f(r^5e^{i\theta})=f\bigl((re^{it})^5\bigr)=g(re^{it})^5 .
\]
Если $t\notin N$, то правая часть имеет конечный предел, а так как $r\mapsto r^5$ --- возрастающая биекция $[0,1)$ на себя, то
существует и конечный $\lim_{s\to1-}f(se^{i\theta})$. Значит, плохие $\theta$ лежат в множестве $5N=\{5t:\ t\in N\}$, которое имеет
нулевую меру (линейное растяжение умножает меру на $5$). Первое утверждение доказано.

\textbf{Шаг 6: интегральная оценка.} Из (16.5), интегрируя по $r\in[0,1)$ почленно (ряд с неотрицательными членами):
\[
\sum_{n\ge1}|b_n|^2=\int_0^1\sum_{n\ge1}n|b_n|^2r^{n-1}dr\le\int_0^1\frac{dr}{(1-r)^{4/5}}=5 .
\]
Наконец, при $r\in[0,1)$, делая замену $s=5t$ и пользуясь $2\pi$-периодичностью подынтегральной функции и (16.1),
\[
\frac1{2\pi}\int_0^{2\pi}|f(r^5e^{is})|^{2/5}ds=\frac1{2\pi}\int_0^{2\pi}|f(r^5e^{5it})|^{2/5}dt
=\frac1{2\pi}\int_0^{2\pi}|g(re^{it})|^{2}dt=\sum_{n\ge1}|b_n|^2r^{2n}\le5 .
\]
(В первом равенстве: $\int_0^{2\pi}\phi(5t)dt=\frac15\int_0^{10\pi}\phi(s)ds=\int_0^{2\pi}\phi(s)ds$ для $2\pi$-периодической $\phi$.)
Так как $r^5$ пробегает всё $[0,1)$, неравенство доказано для всех радиусов.
\end{proof}

\begin{detail}{почему именно корень пятой степени}
Повторим рассуждение для $g=(f(z^k))^{1/k}$. Из $|g(z)|\le|z|(1-|z|^k)^{-2/k}$ так же, как в шаге 2, получаем
$\frac1{2\pi}\int_0^{2\pi}|g'(re^{it})|^2dt\le(1-r)^{-1-4/k}$. Шаги 3--4 с весом $(1-r)^{\alpha}$ требуют
одновременно $\int_0^1(1-r)^{-\alpha}dr<\infty$ (т.\,е. $\alpha<1$) и $\int_0^1(1-r)^{\alpha-1-4/k}dr<\infty$ (т.\,е. $\alpha>4/k$).
Это возможно тогда и только тогда, когда $k>4$. При $k=5$ можно взять $\alpha=9/10$ --- это и сделано выше. Для самой $f$
($k=1$) оценка $|f|\lesssim(1-|z|)^{-2}$ слишком слаба, и метод не работает.
При любом $k\ge5$ тем же способом получаем $\frac1{2\pi}\int_0^{2\pi}|f(re^{it})|^{2/k}dt\le\int_0^1(1-r)^{-4/k}dr=\frac{k}{k-4}$
для $f\in\Sc$.
\end{detail}

\begin{remarkbox}{Замечания}
\begin{enumerate}
\item \emph{Показатель $2/5$ не точный.} Аналогичное утверждение верно для любого $p<1/2$: всякая $f\in\Sc$ лежит в $H^p$ при всех
$p<1/2$ (теорема Правица; без доказательства). Значение $1/2$ точное: для функции Кёбе $|k(re^{it})|^{p}\asymp|1-re^{it}|^{-2p}$
вблизи $t=0$, и $\sup_r\int_0^{2\pi}|1-re^{it}|^{-2p}dt<\infty$ лишь при $2p<1$.
\item \emph{Теорема Фату (без доказательства, с идеей).} Если $F\in H^\infty$ (ограниченная голоморфная в $\DD$ функция), то для
почти всех $\zeta\in\TT$ существует $\lim F(z)$ при $z\to\zeta$ по любому углу с вершиной $\zeta$, не касающемуся $\TT$ (в частности,
радиальный предел). \emph{Идея:} функции $F_r(e^{it})=F(re^{it})$ ограничены в $L^\infty(\TT)$; по теореме Банаха--Алаоглу найдётся
последовательность $r_j\to1$ и $\phi\in L^\infty(\TT)$ с $F_{r_j}\to\phi$ слабо-$*$. Переходя к пределу в формуле Пуассона для
$F(r_jz)$, получаем $F=P[\phi]$ --- интеграл Пуассона функции $\phi$. Ядро Пуассона --- аппроксимативная единица, мажорируемая
радиально убывающей интегрируемой функцией; поэтому $P[\phi](re^{it_0})\to\phi(e^{it_0})$ в каждой точке Лебега $t_0$ функции
$\phi$, а по теореме Лебега о дифференцировании почти каждая точка --- точка Лебега. Для ограниченной $G$ теорема В --- частный
случай теоремы Фату; Предложение 13.12 вместе с теоремой Ф.~Рисса (функции из $H^p$ имеют радиальные пределы почти всюду) даёт
другой путь к общему случаю. Наше доказательство обходится без этой теории.
\item \emph{Всюду --- нельзя.} Существуют ограниченные односвязные области (например, «спиральная» область, граница которой
бесконечно наматывается на окружность), для которых у конформного отображения в отдельных точках $\zeta\in\TT$ радиального
предела нет (без доказательства). Как будет видно в билетах 18--19, если $\partial G$ локально связна, то $f$ непрерывна в $\clos\DD$,
и радиальные пределы существуют всюду.
\end{enumerate}
\end{remarkbox}

\subsection{Условия однолистности Беккера и Нехари (доп. материал)}

На слайде 318 приведены (без доказательства) два достаточных условия однолистности. Пусть $F(z)=f''(z)/f'(z)$ (логарифмическая
производная $f'$) и
\[
S_f(z)=F'(z)-\tfrac12F(z)^2=\Bigl(\frac{f''}{f'}\Bigr)'-\frac12\Bigl(\frac{f''}{f'}\Bigr)^2
\]
--- \emph{шварциан} функции $f$.

\begin{thmbox}{Условие Беккера (без доказательства)}
Пусть $f\in U_\ell(\DD)$ и $(1-|z|^2)\,|zF(z)|\le1$ для всех $z\in\DD$. Тогда $f\in U(\DD)$.
\end{thmbox}

\begin{thmbox}{Условие Нехари (без доказательства)}
Пусть $f$ мероморфна и локально однолистна в $\DD$, и $(1-|z|^2)^2|S_f(z)|\le2$ для всех $z\in\DD$. Тогда $f$ однолистна в $\DD$.
\end{thmbox}

\begin{remarkbox}{Комментарии к условиям}
\begin{itemize}
\item Шварциан инвариантен относительно дробно-линейных преобразований: $S_{T\circ f}=S_f$ для ДЛФ $T$, и $S_T\equiv0$; поэтому условие
Нехари естественно для мероморфных функций.
\item Для функции Кёбе $F(z)=\frac{k''(z)}{k'(z)}=\frac{4+2z}{1-z^2}$, откуда $F'=\frac{2+8z+2z^2}{(1-z^2)^2}$,
$\frac12F^2=\frac{8+8z+2z^2}{(1-z^2)^2}$ и $S_k(z)=-\dfrac{6}{(1-z^2)^2}$; при $z=r\to1-$ величина $(1-r^2)^2|S_k(r)|=6$.
Обратно (теорема Крауса), для любой однолистной $f$ выполнено $(1-|z|^2)^2|S_f(z)|\le6$. Таким образом, необходимое условие
имеет константу $6$, а достаточное условие Нехари --- константу $2$ (тоже точную).
\end{itemize}
\end{remarkbox}

\begin{remarkbox}{Коротко: план ответа на экзамене}
\begin{enumerate}
\item Формулировка: $f\in\mathrm{Conf}(\DD,G)$ $\Rightarrow$ для п.в. $\zeta\in\TT$ существует конечный $\lim_{r\to1-}f(r\zeta)$
(Предложение 13.12 --- даже без ограниченности $G$, с оценкой $\frac1{2\pi}\int|f(re^{it})-f(0)|^{2/5}dt\le5|f'(0)|^{2/5}$).
\item Ключевая идея: если образ радиуса имеет конечную длину $\int_0^1|f'(re^{it})|dr$, то радиальный предел есть (критерий Коши).
\item Ограниченная $G$: $\int_\DD|f'|^2=m_2(G)<\infty$, Коши--Буняковский с весом $dr/r$ $\Rightarrow$ $\int_0^{2\pi}(\int_{1/2}^1|f'|dr)^2dt\le\ln2\cdot m_2(G)$.
\item Общий случай: сводим к $f\in\Sc$ и берём $g=(f(z^5))^{1/5}\in\Sc$, $|g(z)|\le|z|(1-|z|^5)^{-2/5}$.
\item $(1-r)\frac1{2\pi}\int|g'|^2dt\le\sum n|b_n|^2r^{n-1}=\frac{m_2(g(B(0,\sqrt r)))}{\pi r}\le(1-r)^{-4/5}$.
\item Вес $(1-r)^{9/10}$, интегрирование по $r$ и КБШ: $\frac1{2\pi}\int(\int_0^1|g'|dr)^2dt\le100$.
\item Пределы $g$ п.в. $\Rightarrow$ пределы $f$ п.в. (так как $f(r^5e^{5it})=g(re^{it})^5$).
\item $\sum|b_n|^2\le5$ и $\frac1{2\pi}\int|f(r^5e^{is})|^{2/5}ds=\sum|b_n|^2r^{2n}\le5$.
\end{enumerate}
\end{remarkbox}

% -------------------------------------------------------------------
\ticket{17}{Теорема Янишевского и ее комплексно-аналитическое доказательство}
\label{t17:begin}

\begin{defs}
\textbf{Расширенная плоскость} $\Ch=\CC\cup\{\infty\}$ с топологией сферы Римана: окрестности $\infty$ --- множества
$\{|z|>R\}\cup\{\infty\}$. $\Ch$ компактна. Множество $E\subset\Ch$ \emph{замкнуто}, если $\Ch\setminus E$ открыто; например,
замкнутое в $\CC$ множество $E$ замкнуто в $\Ch$ тогда и только тогда, когда оно ограничено или к нему добавлена точка $\infty$.

\textbf{Путь (кривая)} в $X\subset\Ch$ --- непрерывное отображение $\gamma\colon[0,1]\to X$; $[\gamma]=\gamma([0,1])$ --- его
носитель. Путь \emph{соединяет} $\gamma(0)$ и $\gamma(1)$.

\textbf{Связность.} Множество связно, если его нельзя представить как объединение двух непустых непересекающихся относительно
открытых подмножеств. \emph{Компонента связности} --- максимальное связное подмножество. Для открытого $O\subset\Ch$ его компоненты
открыты и линейно связны (любые две точки компоненты соединяются путём в ней); в $\CC$ --- даже ломаной.

\textbf{Разделение точек (слайд 328).} Пусть $E\subset\Ch$ замкнуто и $a,b\in\Ch\setminus E$. Говорят, что $E$
\emph{разделяет} $a$ и $b$, если $a$ и $b$ лежат в разных компонентах связности открытого множества $\Ch\setminus E$. Иначе
говоря, $a$ и $b$ \emph{не разделены} множеством $E$, если их можно соединить путём в $\Ch\setminus E$.
\emph{Примеры:} окружность $\TT$ разделяет $0$ и $\infty$; замкнутая дуга $\{e^{it}:0\le t\le\pi\}$ точки $0$ и $\infty$ не
разделяет (их соединяет луч $\{-iy: y\ge0\}\cup\{\infty\}$).

\textbf{Индекс.} Для замкнутого кусочно гладкого пути $\gamma$ в $\CC$ и $a\notin[\gamma]$:
$\Ind(\gamma,a)=\frac1{2\pi i}\int_\gamma\frac{dz}{z-a}$. Это целое число, постоянное на каждой компоненте связности
$\CC\setminus[\gamma]$ и равное нулю на неограниченной компоненте. Для положительно ориентированной границы $\partial^+Q$ квадрата $Q$:
$\Ind(\partial^+Q,a)=1$ при $a\in\operatorname{int}Q$ и $=0$ при $a\notin Q$.

\textbf{Ветвь логарифма.} Пусть $V\subset\CC\setminus\{0\}$ открыто. Функция $f\in H(V)$ называется (однозначной голоморфной)
ветвью $\Ln z$ в $V$, если $e^{f(z)}=z$ при всех $z\in V$. Тогда $f'(z)=1/z$ (дифференцируем: $e^{f}f'=1$). Две ветви на
\emph{связном} $V$ отличаются на постоянную $2\pi ik$, $k\in\ZZ$ (функция $(f_1-f_2)/(2\pi i)$ непрерывна и целочисленна).

\textbf{Цепь из отрезков.} Конечная формальная сумма ориентированных отрезков $\Gamma=\sum e_j$; $\int_\Gamma\phi\,dz:=\sum_j\int_{e_j}\phi\,dz$.
Если в сумме встречаются отрезок и он же с противоположной ориентацией, их вклады в любой интеграл взаимно уничтожаются
(«сокращение»).
\end{defs}

\subsection{Формулировка и примеры}

\begin{thmbox}{Теорема 13.7 (теорема Янишевского)}
Пусть $E_1$ и $E_2$ --- два замкнутых множества в $\Ch$ такие, что множество $E_1\cap E_2$ связно (допускается
$E_1\cap E_2=\varnothing$). Если точки $a,b\in\Ch$ не разделены ни $E_1$, ни $E_2$, то они не разделены и $E_1\cup E_2$.
\end{thmbox}

Эквивалентная (контрапозитивная) форма: \emph{если $E_1\cup E_2$ разделяет $a$ и $b$, а $E_1\cap E_2$ связно, то $a$ и $b$
разделяет уже одно из множеств $E_1$, $E_2$.}

\begin{remarkbox}{Примеры: условие связности пересечения существенно}
\begin{itemize}
\item $E_1=\{e^{it}:0\le t\le\pi\}$, $E_2=\{e^{it}:\pi\le t\le2\pi\}$. Ни одна замкнутая полуокружность не разделяет $0$ и
$\infty$ (их соединяет луч $\{-iy:y\ge0\}$, соответственно $\{iy:y\ge0\}$, вместе с $\infty$). Но $E_1\cup E_2=\TT$ разделяет $0$ и
$\infty$. Здесь $E_1\cap E_2=\{1,-1\}$ несвязно.
\item Если $E_1\cap E_2=\varnothing$, теорема утверждает: объединение двух непересекающихся замкнутых множеств, ни одно из которых не
разделяет $a$ и $b$, тоже их не разделяет. Например, две непересекающиеся замкнутые дуги не разделяют никакие точки вне них.
\item \emph{Применение} (билет 18): в доказательстве теоремы Каратеодори о непрерывном продолжении теорема применяется к
$E_1=Y\cup f(\Gamma)$ (маленькое множество) и $E_2=\Ch\setminus G$.
\end{itemize}
\end{remarkbox}

\subsection{Вспомогательные утверждения}

\begin{thmbox}{Лемма А (выбор хороших путей)}
Пусть $E\subset\Ch$ замкнуто, $0,\infty\notin E$, и $E$ не разделяет $0$ и $\infty$. Тогда существует путь $C$ из $0$ в $\infty$ в
$\Ch\setminus E$, состоящий из ломаной с конечным числом звеньев $[0=p_0,p_1]\cup\dots\cup[p_{s-1},p_s]$ и луча
$\{tp_s:\ t\ge1\}\cup\{\infty\}$, где $p_s\ne0$. Множество $C\cap\CC$ замкнуто в $\CC$.
\end{thmbox}
\begin{proof}
Пусть $V$ --- компонента $\Ch\setminus E$, содержащая $0$ и $\infty$. Так как $\infty\notin E$ и $E$ замкнуто, то
$N=\{|z|>R\}\cup\{\infty\}\subset V$ при некотором $R$.

\emph{$V\setminus\{\infty\}$ связно.} Пусть $V\setminus\{\infty\}=A\sqcup A'$, $A,A'$ открыты и непусты. Связное множество
$\{|z|>R\}$ лежит в одном из них, скажем в $A$. Тогда $A\cup\{\infty\}=A\cup N$ открыто в $\Ch$, и $V=(A\cup\{\infty\})\sqcup A'$
--- разбиение связного $V$ на два непустых открытых множества, противоречие.

Открытое связное множество $V\setminus\{\infty\}\subset\CC$ связно ломаными; соединим ломаной в нём $0$ с точкой $p_s$, $|p_s|>R$.
Луч $\{tp_s:t\ge1\}$ лежит в $\{|z|>R\}\subset V$. Замкнутость $C\cap\CC$ очевидна (конечное объединение отрезков и замкнутого луча).
\end{proof}

\begin{thmbox}{Лемма Б (ветвь логарифма вне пути из $0$ в $\infty$)}
Пусть $C$ --- путь из Леммы А. Тогда на открытом множестве $\CC\setminus C$ существует ветвь $\Ln z$.
\end{thmbox}
\begin{proof}
Пусть $\gamma$ --- замкнутый кусочно гладкий путь в $\CC\setminus C$; так как $[\gamma]$ компактно, $[\gamma]\subset B(0,M)$.
Множество $C\cap\CC$ связно, содержит $0$ и точки с $|z|>M$, не пересекает $[\gamma]$. Значит, $C\cap\CC$ целиком лежит в одной
компоненте $\CC\setminus[\gamma]$, и это неограниченная компонента (она содержит точки вне $B(0,M)$, а такие точки лежат в неограниченной
компоненте). Поэтому $0$ лежит в неограниченной компоненте, и $\int_\gamma\frac{dz}{z}=2\pi i\,\Ind(\gamma,0)=0$.

Итак, форма $dz/z$ имеет нулевые интегралы по всем замкнутым путям в $V=\CC\setminus C$. На каждой компоненте $V_0$ множества $V$
фиксируем $z_0\in V_0$, значение $w_0$ с $e^{w_0}=z_0$ и положим $f(z)=w_0+\int_{z_0}^z\frac{d\zeta}{\zeta}$ (интеграл по любой
ломаной в $V_0$; он не зависит от ломаной). Тогда $f\in H(V_0)$, $f'=1/z$, и $(ze^{-f})'=e^{-f}(1-zf')=0$, т.\,е. $ze^{-f}$ постоянна
на $V_0$; в точке $z_0$ она равна $1$. Значит, $e^f=z$ на $V_0$, а значит, и на $V$.
\end{proof}

\begin{remarkbox}{Замечание к слайду 328}
На слайде существование ветви объясняется тем, что каждая компонента $\CC\setminus C_j$ односвязна и не содержит $0$. Это верно
(компоненты дополнения к связному компакту в $\Ch$ односвязны), но доказательство через индекс короче и не требует теории
односвязных областей. Выбор \emph{ломаных} путей $C_j$ (вместо произвольных кривых) гарантирует, что $\CC\setminus(C_1\cup C_2)\ne\varnothing$
и все множества имеют простое строение; в общем случае кривая в $\Ch$ может даже заполнять всю сферу (кривые Пеано).
\end{remarkbox}

\subsection{Доказательство теоремы Янишевского}

\begin{proof}[Доказательство теоремы 13.7]
\textbf{Шаг 0: нормировка.} Точки $a,b$ не лежат в $E_1\cup E_2$ (по определению разделения $a,b\in\Ch\setminus E_j$). Если $a=b$,
доказывать нечего. Если $a\ne b$, возьмём ДЛФ $T$ с $T(a)=0$, $T(b)=\infty$. $T$ --- гомеоморфизм $\Ch$ на себя, переводящий
замкнутые множества в замкнутые, связные --- в связные, компоненты дополнения --- в компоненты дополнения. Поэтому можно считать
$a=0$, $b=\infty$.

\textbf{Шаг 1: пути и ветви логарифма.} По Лемме А найдём пути $C_1$, $C_2$ из $0$ в $\infty$ («ломаная + луч»), причём
$C_j\cap E_j=\varnothing$, $j=1,2$. По Лемме Б существуют ветви $f_j\in H(\CC\setminus C_j)$ функции $\Ln z$.

\textbf{Шаг 2: согласование ветвей.} Множество $E_1\cap E_2$ связно и не пересекается с $C_1\cup C_2$ (так как $E_1\cap C_1=\varnothing$,
$E_2\cap C_2=\varnothing$), а также не содержит $\infty$. Значит, оно лежит в одной компоненте $B$ открытого множества
$\CC\setminus(C_1\cup C_2)$ (это множество непусто: $C_1\cup C_2$ --- объединение конечного числа отрезков и лучей). Если
$E_1\cap E_2=\varnothing$, в качестве $B$ берём любую компоненту. Связное $B$ лежит в некоторой компоненте $V_1$ множества
$\CC\setminus C_1$ и в некоторой компоненте $V_2$ множества $\CC\setminus C_2$. На $B$ функции $f_1,f_2$ --- ветви логарифма, поэтому
$f_1-f_2\equiv2\pi ik$ на $B$. Заменив $f_2$ на $f_2+2\pi ik$ в компоненте $V_2$ (это снова ветвь логарифма на $\CC\setminus C_2$),
добьёмся равенства $f_1=f_2$ на $B$.

\textbf{Шаг 3: склейка в одну функцию.} Множества $K_j:=E_j\setminus B=E_j\cap(\Ch\setminus B)$ замкнуты в $\Ch$ (пересечение
замкнутых), не содержат $\infty$, значит, компактны в $\CC$. Они не пересекаются: $K_1\cap K_2=(E_1\cap E_2)\setminus B=\varnothing$.
Кроме того, $K_j\cap C_j=\varnothing$. Выберем $\eta>0$ так, чтобы
\[
2\eta<\dist(K_1,K_2),\qquad \eta<\dist(K_j,C_j\cap\CC),\ j=1,2
\]
(расстояния от компакта до непересекающегося с ним замкнутого множества положительны; если $K_j=\varnothing$, условие отсутствует).
Положим $U_j=\{z:\dist(z,K_j)<\eta\}$ (при $K_j=\varnothing$ полагаем $U_j=\varnothing$). Тогда $U_j$ открыты,
\[
K_j\subset U_j\subset\CC\setminus C_j,\qquad U_1\cap U_2=\varnothing .
\]
Положим $W=U_1\cup U_2\cup B$ и определим $f\colon W\to\CC$:
\[
f=f_1\ \text{на } U_1,\qquad f=f_2\ \text{на } U_2,\qquad f=f_1=f_2\ \text{на } B .
\]
Определение корректно: $U_1\cap U_2=\varnothing$, а на пересечениях $U_j\cap B$ оба предписания дают $f_j$. Итак, $f\in H(W)$ и
$e^{f(z)}=z$, т.\,е. $f'(z)=1/z$ на $W$. При этом $E_1\cup E_2\subset K_1\cup K_2\cup B\subset W$, и $0\notin W$ (так как
$0\in C_1\cap C_2$, а $U_j\cap C_j=\varnothing$, $B\cap(C_1\cup C_2)=\varnothing$).

\textbf{Шаг 4: предположение противного.} Пусть $E_1\cup E_2$ разделяет $0$ и $\infty$. Пусть $G$ --- компонента открытого
множества $\Ch\setminus(E_1\cup E_2)$, содержащая $0$; тогда $\infty\notin G$. Множество $G$ ограничено: окрестность
$\{|z|>R\}\cup\{\infty\}$ точки $\infty$, не пересекающаяся с замкнутым множеством $E_1\cup E_2$, связна и потому лежит в компоненте
точки $\infty$, отличной от $G$. Далее, $\partial G\subset E_1\cup E_2$: если $p\in\partial G$ и $p\notin E_1\cup E_2$, то $p$ лежит
в некоторой компоненте $G'$ открытого множества $\Ch\setminus(E_1\cup E_2)$, $G'$ открыта и пересекается с $G$ (так как $p\in\clos G$),
откуда $G'=G\ni p$ --- но $p\notin G$, так как $G$ открыто и $p\in\partial G$. Итак, $\partial G$ --- компакт, лежащий в открытом $W$, и
\[
d:=\dist(\partial G,\CC\setminus W)>0 .
\]

\textbf{Шаг 5: сетка квадратов.} Возьмём $\delta>0$ с $\delta\sqrt2<d$ и покроем плоскость сеткой замкнутых квадратов со стороной
$\delta$ и сторонами, параллельными осям, сдвинув сетку так, чтобы $0$ не лежал на линиях сетки. Пусть $Q_1,\dots,Q_N$ --- все
квадраты сетки, \emph{содержащиеся в $G$} (их конечное число, так как $G$ ограничено). Положим
\[
\Gamma=\sum_{n=1}^N\partial^+Q_n .
\]
\begin{detail}{свойства цепи $\Gamma$}
\emph{(a) $0$ лежит внутри ровно одного $Q_n$.} Пусть $Q_0$ --- квадрат сетки, содержащий $0$ внутри. Так как $0\notin W$, то
$\dist(0,\partial G)\ge d>\delta\sqrt2=\diam Q_0$, поэтому $Q_0\cap\partial G=\varnothing$; связное $Q_0$ содержит точку $0\in G$ и не
пересекает $\partial G$, значит, $Q_0\subset G$, т.\,е. $Q_0$ --- один из $Q_n$. Другие квадраты сетки не содержат $0$ (сетка
выбрана так, что $0$ не на её линиях). Следовательно,
\[
\frac1{2\pi i}\int_\Gamma\frac{dz}{z}=\sum_{n=1}^N\Ind(\partial^+Q_n,0)=1 .
\]
\emph{(b) Сокращения.} Если сторона $e$ является общей для двух квадратов $Q_n,Q_m$ набора, то в $\partial^+Q_n$ и $\partial^+Q_m$
она входит с противоположными ориентациями, и её вклады сокращаются. Пусть $\mathcal E'$ --- множество ориентированных сторон,
которые принадлежат ровно одному квадрату набора («несокращённые» стороны). Тогда для любой функции $\phi$, непрерывной на
объединении сторон, $\int_\Gamma\phi\,dz=\sum_{e\in\mathcal E'}\int_e\phi\,dz$.

\emph{(c) Несокращённые стороны лежат в $W$.} Пусть $e\in\mathcal E'$ --- сторона квадрата $Q_n\subset G$, и $Q'$ --- соседний по
этой стороне квадрат сетки. Тогда $Q'\not\subset G$ (иначе $e$ сократилась бы). Связное $Q'$ пересекает $G$ (по $e$) и
$\CC\setminus G$, значит, содержит точку $p\in\partial G$. Для любой $x\in e\subset Q'$: $|x-p|\le\diam Q'=\delta\sqrt2<d$, поэтому
$x\in W$. Итак, $e\subset W$, и, в частности, $0\notin e$.

\emph{(d) Цепь несокращённых сторон --- цикл.} Для любой функции $\psi$ на вершинах сетки
$\sum_{n}\sum_{e\subset\partial^+Q_n}\bigl(\psi(e^+)-\psi(e^-)\bigr)=0$ ($e^\pm$ --- конец и начало $e$): для каждого квадрата сумма
телескопическая по замкнутой ломаной. Сокращающиеся пары вносят $\psi(e^+)-\psi(e^-)+\psi(e^-)-\psi(e^+)=0$. Значит,
$\sum_{e\in\mathcal E'}(\psi(e^+)-\psi(e^-))=0$.
\end{detail}

\textbf{Шаг 6: противоречие.} Применим (d) к функции $\psi=f$ на вершинах из $W$ и $\psi=0$ на остальных вершинах (концы сторон из
$\mathcal E'$ лежат в $W$ по (c)). Так как на каждой стороне $e\in\mathcal E'\subset W$ имеем $\frac1z=f'(z)$, то по (a), (b)
\[
2\pi i=\int_\Gamma\frac{dz}{z}=\sum_{e\in\mathcal E'}\int_e\frac{dz}{z}=\sum_{e\in\mathcal E'}\int_ef'(z)\,dz
=\sum_{e\in\mathcal E'}\bigl(f(e^+)-f(e^-)\bigr)=0 ,
\]
что невозможно. Значит, $E_1\cup E_2$ не разделяет $0$ и $\infty$.
\end{proof}

\begin{remarkbox}{Замечания к слайдам 329--330}
\begin{itemize}
\item На слайде 329 написано «$e^{f(z)}=z$ при всех $z\in\Omega$»; должно быть: при всех $z\in W$.
\item Открытые множества $U_1,U_2$ нужно выбирать \emph{непересекающимися} (на слайде это не оговорено), иначе $f$ на $U_1\cap U_2$
могла бы быть определена неоднозначно: вне $B$ ветви $f_1$ и $f_2$ не обязаны совпадать.
\item Условие на шаг сетки «$\delta<\dist(\partial G,\partial W)$» следует усилить до $\delta\sqrt2<\dist(\partial G,\CC\setminus W)$:
несокращённая сторона лежит в квадрате, задевающем $\partial G$, т.\,е. на расстоянии не более диагонали $\delta\sqrt2$ от $\partial G$.
\end{itemize}
\end{remarkbox}

\begin{remarkbox}{Почему доказательство «комплексно-аналитическое»}
Чисто топологическое утверждение (о связности дополнений) доказано средствами ТФКП: разделённость точек $0$ и $\infty$ кодируется
существованием ветви логарифма ($\int dz/z=0$ по циклам), и противоречие получается из сравнения $\int_\Gamma dz/z=2\pi i$ (индекс)
с $\int_\Gamma f'dz=0$ (интеграл от производной по циклу).
\end{remarkbox}

\begin{remarkbox}{Коротко: план ответа на экзамене}
\begin{enumerate}
\item Определение: $E$ разделяет $a,b$, если они в разных компонентах $\Ch\setminus E$. Формулировка теоремы; пример двух
полуокружностей (пересечение несвязно $\Rightarrow$ вывод неверен).
\item Сводим к $a=0$, $b=\infty$ (ДЛФ); пути $C_j$ (ломаная + луч) из $0$ в $\infty$ вне $E_j$.
\item На $\CC\setminus C_j$ есть ветвь $f_j=\Ln z$ (индекс любого цикла вокруг $0$ равен $0$).
\item Связное $E_1\cap E_2$ лежит в одной компоненте $B$ множества $\CC\setminus(C_1\cup C_2)$; нормируем $f_1=f_2$ на $B$.
\item $E_1\setminus B$, $E_2\setminus B$ --- непересекающиеся компакты; их непересекающиеся окрестности $U_j\subset\CC\setminus C_j$;
склеиваем $f\in H(W)$, $W=U_1\cup U_2\cup B\supset E_1\cup E_2$, $e^f=z$, $0\notin W$.
\item Если $E_1\cup E_2$ разделяет, то компонента $G\ni0$ ограничена и $\partial G\subset W$.
\item Квадраты сетки внутри $G$: $\int_\Gamma dz/z=2\pi i$, но несокращённые стороны лежат в $W$, где $1/z=f'$, и интеграл равен $0$.
\end{enumerate}
\end{remarkbox}
% -------------------------------------------------------------------
\ticket{18}{Области с локально связными границами и теорема Каратеодори о непрерывном продолжении конформного отображения единичного круга на данную область}
\label{t18:begin}

\begin{defs}
\textbf{Постановка (слайд 322).} Пусть $G\subset\CC$ --- односвязная область и $f\in\mathrm{Conf}(\DD,G)$ (конформное, т.\,е.
голоморфное и инъективное, отображение $\DD$ \emph{на} $G$; существует по теореме Римана, если $G\ne\CC$). Вопрос: когда $f$ можно
непрерывно продолжить на $\clos\DD$, и когда это продолжение инъективно?

\textbf{Континуум} --- непустое связное компактное множество. Примеры: точка, отрезок, окружность, носитель любого пути.

\textbf{Путь (кривая)} --- непрерывное отображение $\gamma\colon[0,1]\to\CC$; его \emph{носитель} $[\gamma]=\gamma([0,1])$ --- континуум.
Фраза «множество $K$ является кривой» означает: $K=[\gamma]$ для некоторого пути $\gamma$ (самопересечения допускаются).

\textbf{Локальная связность (Определение 13.5).} Замкнутое множество $E\subset\CC$ \emph{локально связно в точке} $a\in E$, если
для любого $\eps>0$ найдётся $\delta>0$ такое, что для любых $z_1,z_2\in E\cap B(a,\delta)$ существует континуум $Y$ с
\[
z_1,z_2\in Y\subset E\cap B(a,\eps).
\]
$E$ \emph{локально связно}, если оно локально связно в каждой своей точке. Неформально: близкие точки $E$ соединяются внутри $E$
«маленьким» связным множеством.

\textbf{Равномерная локальная связность.} Замкнутое $E$ \emph{равномерно локально связно на множестве} $K\subset E$, если
$\forall\eps>0\ \exists\delta>0$: для любых $z_1\in K$, $z_2\in E$ с $|z_1-z_2|<\delta$ найдётся континуум $Y\subset E$,
$z_1,z_2\in Y$, $\diam Y<\eps$. При $K=E$ говорим просто «равномерно локально связно».

\textbf{Стандартное (топологическое) определение.} $E$ локально связно в смысле общей топологии, если у каждой точки $a\in E$ и
любого $\eps>0$ есть \emph{связная} относительно открытая в $E$ окрестность $V\ni a$, $V\subset B(a,\eps)$.

\textbf{Примеры.}
1) Отрезок, окружность, любая кривая (Лемма 18.2 ниже), в частности граница жордановой области, локально связны.
2) Граница области $\DD\setminus[0,1)$ --- множество $\TT\cup[0,1]$ --- локально связна (это кривая: пройти отрезок от $0$ до $1$,
затем окружность), хотя область не жорданова.
3) \emph{«Гребёнка» (слайд 322).} $G=\{x+iy:\ |x|<1,\ 0<y<1\}\setminus\bigcup_{n\ge2}[\,i/n,\ 1+i/n\,]$. Её граница не локально
связна в точке $0$ (см. Пример 18.3).
4) Множество $\{0\}\cup\{1/n:n\in\NN\}$ не локально связно в $0$ (континуум, содержащий две различные точки этого множества, ---
связное подмножество счётного множества, т.\,е. точка).
5) Замыкание графика $y=\sin(1/x)$, $0<x\le1$ («синусоида топологов»), не локально связно в точках отрезка $\{0\}\times[-1,1]$.

\textbf{Разделение точек (слайд 328).} Замкнутое $E\subset\Ch$ \emph{разделяет} точки $a,b\in\Ch\setminus E$, если они лежат в разных
компонентах связности $\Ch\setminus E$; иначе (т.\,е. если их можно соединить путём в $\Ch\setminus E$) --- \emph{не разделяет}.

\textbf{Длина кривой.} Для пути, параметризованного дугой окружности, $\ell(h(L))=\int_L|h'(z)|\,|dz|$; если $L$ --- открытое
подмножество окружности (не более чем счётное объединение дуг), длина образа --- сумма длин образов дуг.
\end{defs}

\subsection{Локальная связность: эквивалентные определения и свойства}

\begin{thmbox}{Лемма 18.1 (эквивалентность определений)}
Пусть $E\subset\CC$ замкнуто.
\begin{enumerate}
\item[(а)] $E$ локально связно в смысле Определения 13.5 (в каждой точке) $\iff$ $E$ локально связно в стандартном смысле.
\item[(б)] Если $E$ локально связно в каждой точке компакта $K\subset E$, то $E$ равномерно локально связно на $K$. В частности,
компакт $E$ локально связен $\iff$ он равномерно локально связен.
\end{enumerate}
\end{thmbox}
\begin{proof}
(а) «$\Leftarrow$». Пусть $a\in E$, $\eps>0$. Возьмём связную относительно открытую окрестность $V\ni a$ с $V\subset B(a,\eps/2)$.
Её замыкание $\clos V$ (замыкание в $\CC$; оно лежит в $E$, так как $E$ замкнуто) связно (замыкание связного множества связно),
компактно и лежит в $\clos{B(a,\eps/2)}\subset B(a,\eps)$. Так как $V$ открыто в $E$, то $V\supset E\cap B(a,\delta)$ при
некотором $\delta>0$. Для $z_1,z_2\in E\cap B(a,\delta)$ годится $Y=\clos V$.

«$\Rightarrow$». Пусть $a\in E$, $\eps>0$. Пусть $V$ --- множество тех $z\in E\cap B(a,\eps)$, которые можно соединить с $a$
связным подмножеством $E\cap B(a,\eps)$ (т.\,е. компонента точки $a$ в $E\cap B(a,\eps)$). $V$ связно (объединение связных множеств,
содержащих $a$) и $a\in V\subset B(a,\eps)$. Проверим, что $V$ открыто в $E$. Пусть $z\in V$; положим $\eps'=\eps-|z-a|>0$ и найдём
по Определению 13.5 в точке $z$ число $\delta'$: любая $w\in E\cap B(z,\delta')$ соединяется с $z$ континуумом
$Y\subset E\cap B(z,\eps')\subset E\cap B(a,\eps)$. Объединение $Y$ со связным множеством, соединяющим $z$ с $a$, связно; значит
$w\in V$, т.\,е. $E\cap B(z,\delta')\subset V$.

(б) Пусть $\eps>0$. Для каждой $a\in K$ возьмём $\delta_a>0$ из Определения 13.5 для $\eps/2$. Шары $B(a,\delta_a/2)$, $a\in K$,
покрывают компакт $K$; выберем конечное подпокрытие $B(a_i,\delta_{a_i}/2)$, $i=1,\dots,N$, и положим
$\delta=\min_i\delta_{a_i}/2$. Пусть $z_1\in K$, $z_2\in E$, $|z_1-z_2|<\delta$. Тогда $z_1\in B(a_i,\delta_{a_i}/2)$ при некотором
$i$, и $|z_2-a_i|<\delta_{a_i}/2+\delta\le\delta_{a_i}$. Значит, есть континуум $Y\ni z_1,z_2$, $Y\subset E\cap B(a_i,\eps/2)$.
Так как $Y$ компактно лежит в открытом шаре радиуса $\eps/2$, то $\diam Y<\eps$.

Обратно, если компакт $E$ равномерно локально связен, $a\in E$, $\eps>0$: возьмём $\delta\le\eps$ из равномерного определения для
$\eps/2$; для $z_1,z_2\in E\cap B(a,\delta/2)$ имеем $|z_1-z_2|<\delta$, и континуум $Y\ni z_1$ диаметра $<\eps/2$ лежит в
$B(a,\delta/2+\eps/2)\subset B(a,\eps)$.
\end{proof}

\begin{thmbox}{Лемма 18.2 (кривая локально связна)}
Если $\gamma\colon[0,1]\to\CC$ --- путь, то $K=[\gamma]$ --- локально связный континуум.
\end{thmbox}
\begin{proof}
(Это частный случай утверждения слайда 322 «непрерывный образ компактного локально связного множества локально связен».)
Пусть $a\in K$, $\eps>0$. По равномерной непрерывности $\gamma$ найдём $\eta>0$: $|s-t|\le\eta\Rightarrow|\gamma(s)-\gamma(t)|<\eps/2$.
Множество $T_a=\gamma^{-1}(a)$ компактно и непусто. Его открытая $\eta$-окрестность $O=\{t\in[0,1]:\dist(t,T_a)<\eta\}$ --- открытое
подмножество $[0,1]$; каждая её компонента $J$ --- (относительно открытый) интервал, содержащий точку $s_J\in T_a$ и потому
содержащий $(s_J-\eta,s_J+\eta)\cap[0,1]$, т.\,е. имеющий длину не меньше $\min(\eta,1)$; компоненты не пересекаются, поэтому их
конечное число: $J_1,\dots,J_p$. Для $t\in\clos{J_k}$ имеем $\dist(t,T_a)\le\eta$, значит $|\gamma(t)-a|<\eps/2$. Поэтому
$Y_k:=\gamma(\clos{J_k})$ --- континуум, содержащий $a=\gamma(s_{J_k})$, и $Y_k\subset B(a,\eps/2)$.

Компакт $F=[0,1]\setminus O$ не пересекает $T_a$, поэтому $\delta_0:=\dist(a,\gamma(F))>0$ (при $F=\varnothing$ положим $\delta_0=\infty$).
Пусть $\delta=\min(\delta_0,\eps/2)$ и $z_1,z_2\in K\cap B(a,\delta)$, $z_j=\gamma(t_j)$. Тогда $t_j\notin F$, т.\,е.
$t_j\in J_{k_j}$, и $Y=Y_{k_1}\cup Y_{k_2}$ --- континуум (два связных множества с общей точкой $a$), $z_1,z_2\in Y\subset K\cap B(a,\eps)$.
\end{proof}

\begin{remarkbox}{Пример 18.3 (гребёнка: граница не локально связна)}
Пусть $G=\{x+iy:|x|<1,\,0<y<1\}\setminus\bigcup_{n\ge2}S_n$, $S_n=[\,i/n,\,1+i/n\,]$ --- горизонтальные «зубья», идущие от мнимой
оси до правой стороны прямоугольника. $G$ --- ограниченная односвязная область (левая половина прямоугольника связывает все «щели»
между зубьями). Граница $\partial G$ состоит из границы прямоугольника и всех $S_n$; точка $0$ и точки $i/n\to0$ лежат на $\partial G$.

\emph{Утверждение:} любое связное $Y\subset\partial G$, содержащее $0$ и $i/n$, содержит точку $1+i/n$, поэтому $\diam Y\ge1$; значит
$\partial G$ не локально связна в $0$. \emph{Доказательство.} Положим $S_n'=S_n\setminus\{1+i/n\}=\{x+i/n:0\le x<1\}$. Точки
$S_n'$ не лежат на границе прямоугольника (там $y\in\{0,1\}$ или $x=\pm1$) и находятся на расстоянии $\ge\frac1n-\frac1{n+1}>0$ от
других зубьев. Поэтому $S_n'$ одновременно замкнуто и открыто в $\partial G\setminus\{1+i/n\}$. Если бы $Y\not\ni1+i/n$, то
$Y=(Y\cap S_n')\sqcup(Y\setminus S_n')$ --- разбиение на два непустых ($i/n$ и $0$) относительно открытых множества, противоречие
со связностью. По теореме Каратеодори (ниже) конформное отображение $\DD$ на гребёнку \emph{не} продолжается непрерывно на $\clos\DD$.
\end{remarkbox}

\subsection{Вспомогательные утверждения}

\begin{thmbox}{Лемма 18.4 (соответствие границ)}
Пусть $f\in\mathrm{Conf}(\DD,G)$.
\begin{enumerate}
\item[(а)] Если $\sigma\colon[0,1)\to\DD$ --- путь, $|\sigma(t)|\to1$ при $t\to1$, и существует $w=\lim_{t\to1}f(\sigma(t))\in\CC$, то $w\in\partial G$.
\item[(б)] Если $f$ продолжается до непрерывной функции на $\clos\DD$ (обозначим её тоже $f$), то $f(\TT)=\partial G$.
\item[(в)] Если $f$ равномерно непрерывна в $\DD$, то она продолжается до непрерывной функции на $\clos\DD$.
\end{enumerate}
\end{thmbox}
\begin{proof}
(а) Ясно, что $w\in\clos G$. Если $w\in G$, то по непрерывности $f^{-1}\colon G\to\DD$ имеем $\sigma(t)=f^{-1}(f(\sigma(t)))\to f^{-1}(w)\in\DD$,
что противоречит $|\sigma(t)|\to1$. Значит, $w\in\partial G$.

(б) Для $\zeta\in\TT$ имеем $f(\zeta)=\lim_{r\to1-}f(r\zeta)\in\partial G$ по (а); значит $f(\TT)\subset\partial G$. Обратно, пусть
$w\in\partial G$, $w=\lim f(z_n)$, $z_n\in\DD$. Переходя к подпоследовательности, $z_n\to z_0\in\clos\DD$. Если $z_0\in\DD$, то
$w=f(z_0)\in G$, что невозможно ($G$ открыто, $w\in\partial G$). Значит $z_0\in\TT$ и $w=f(z_0)\in f(\TT)$.

(в) Пусть $\zeta\in\TT$ и $z_n\to\zeta$, $z_n\in\DD$. По равномерной непрерывности последовательность $f(z_n)$ фундаментальна, значит,
сходится; предел не зависит от выбора $(z_n)$ (две последовательности можно перемешать в одну). Положим $f(\zeta)$ равным этому
пределу. Пусть $|z-z'|<\rho\Rightarrow|f(z)-f(z')|<\eps$ для $z,z'\in\DD$. Если $w,w'\in\clos\DD$, $|w-w'|<\rho$, приблизим их точками
$z_n\to w$, $z_n'\to w'$ из $\DD$; при больших $n$ $|z_n-z_n'|<\rho$, откуда $|f(w)-f(w')|=\lim|f(z_n)-f(z_n')|\le\eps$. Значит,
продолжение равномерно непрерывно на $\clos\DD$.
\end{proof}

\begin{thmbox}{Лемма 13.7 (лемма Вольфа)}
Пусть $h$ конформно (голоморфно и инъективно) отображает открытое множество $U\subset\CC$ внутрь круга $B(0,R)$. Пусть $a\in\CC$ и
$L(a,r)=U\cap C(a,r)$, где $C(a,r)=\{|z-a|=r\}$. Тогда для любого $\rho\in(0,1)$
\[
\inf_{\rho<r<\sqrt\rho}\ \ell\bigl(h(L(a,r))\bigr)\le\frac{2\pi R}{\sqrt{\log(1/\rho)}} .
\]
В частности, существует убывающая последовательность $r_n\to0$ такая, что $\ell(h(L(a,r_n)))\to0$.
\end{thmbox}
\begin{proof}
Обозначим $\Theta_r=\{\vartheta\in[0,2\pi]:\ a+re^{i\vartheta}\in U\}$ (открытое множество) и
\[
\lambda(r)=\ell\bigl(h(L(a,r))\bigr)=\int_{L(a,r)}|h'(z)|\,|dz|=\int_{\Theta_r}|h'(a+re^{i\vartheta})|\,r\,d\vartheta\in[0,+\infty].
\]
(Функция $(r,\vartheta)\mapsto|h'(a+re^{i\vartheta})|\,\mathbf 1_U(a+re^{i\vartheta})$ борелевская, поэтому по теореме Тонелли
$\lambda$ измерима по $r$ и все интегралы ниже имеют смысл.)

\textbf{Шаг 1: Коши--Буняковский по окружности.} Запишем $|h'|=1\cdot|h'|$:
\[
\lambda(r)^2=\Bigl(\int_{L(a,r)}|h'(z)|\,|dz|\Bigr)^2\le\int_{L(a,r)}|dz|\cdot\int_{L(a,r)}|h'(z)|^2|dz|
\le2\pi r\int_{\Theta_r}|h'(a+re^{i\vartheta})|^2\,r\,d\vartheta ,
\]
так как длина $L(a,r)$ не больше длины всей окружности $2\pi r$. Делим на $r$:
\[
\frac{\lambda(r)^2}{r}\le2\pi\int_{\Theta_r}|h'(a+re^{i\vartheta})|^2\,r\,d\vartheta .
\]

\textbf{Шаг 2: интегрирование по $r$ (мера $dr/r$).} Интегрируя по $r\in(0,\infty)$ и переходя к полярным координатам с центром $a$
($dm_2=r\,dr\,d\vartheta$):
\[
\int_0^\infty\frac{\lambda(r)^2}{r}\,dr\le2\pi\int_0^\infty\!\!\int_{\Theta_r}|h'(a+re^{i\vartheta})|^2\,r\,d\vartheta\,dr
=2\pi\int_U|h'(z)|^2\,dm_2(z)=2\pi\,m_2(h(U))\le2\pi\cdot\pi R^2 .
\]
Здесь использована формула площади $m_2(h(U))=\int_U|h'|^2dm_2$ для инъективного голоморфного $h$ (якобиан $h$ равен $|h'|^2$ по
условиям Коши--Римана) и включение $h(U)\subset B(0,R)$.

\textbf{Шаг 3: оценка инфимума.} Так как $\int_\rho^{\sqrt\rho}\frac{dr}r=\log\sqrt\rho-\log\rho=\frac12\log\frac1\rho$, то
\[
2\pi^2R^2\ge\int_\rho^{\sqrt\rho}\frac{\lambda(r)^2}{r}\,dr\ge\inf_{\rho<r<\sqrt\rho}\lambda(r)^2\cdot\frac12\log\frac1\rho ,
\]
откуда $\inf_{\rho<r<\sqrt\rho}\lambda(r)^2\le\dfrac{4\pi^2R^2}{\log(1/\rho)}$, что и требовалось.

\textbf{Последовательность $r_n$.} Положим $\rho_1=1/4$ и $\rho_{n+1}=\rho_n^2$; тогда интервалы $(\rho_n,\sqrt{\rho_n})=(\rho_n,\rho_{n-1})$
не пересекаются и идут к нулю. В $n$-м интервале выберем $r_n$ с $\lambda(r_n)<\frac{2\pi R}{\sqrt{\log(1/\rho_n)}}+\frac1n$ (по
определению инфимума). Тогда $r_n$ убывают, $r_n\to0$, $\lambda(r_n)\to0$, так как $\log(1/\rho_n)=2^{n-1}\log4\to\infty$.
\end{proof}

\begin{remarkbox}{Смысл леммы Вольфа (принцип «длина--площадь»)}
Образ $h(U)$ имеет конечную площадь, а площадь $=\int\!\!\int|h'|^2$. Если бы все окружности $C(a,r)$ с $r\in(\rho,\sqrt\rho)$ имели
длинные образы, то, по Коши--Буняковскому, площадь была бы большой: «кольцо» $\rho<|z-a|<\sqrt\rho$ имеет конформный модуль
$\frac12\log\frac1\rho\to\infty$. Ключевое --- мера $dr/r$, относительно которой масса отрезка $(\rho,\sqrt\rho)$ неограниченно растёт.
\end{remarkbox}

\begin{thmbox}{Теорема 13.7 (теорема Янишевского; доказательство --- билет 17)}
Пусть $E_1,E_2\subset\Ch$ замкнуты и $E_1\cap E_2$ связно (возможно, пусто). Если точки $a,b\in\Ch$ не разделены ни $E_1$, ни
$E_2$, то они не разделены и $E_1\cup E_2$.
\end{thmbox}

\subsection{Теорема Каратеодори о непрерывном продолжении}

\begin{thmbox}{Теорема 13.6 (теорема Каратеодори)}
Пусть $f$ отображает $\DD$ конформно на ограниченную область $G$. Следующие условия эквивалентны:
\begin{enumerate}
\item $f$ имеет непрерывное продолжение на $\clos\DD$;
\item множество $\partial G$ является кривой, т.\,е. $\partial G=[\gamma]$ для некоторого пути $\gamma\colon[0,1]\to\CC$;
\item множество $\partial G$ локально связно;
\item множество $\CC\setminus G$ локально связно.
\end{enumerate}
\end{thmbox}

Утверждение 2)$\Rightarrow$1) означает: если $\partial G$ можно задать какой-либо замкнутой кривой (возможно, с самопересечениями), то её
можно задать «конформной параметризацией» $\partial G=\{f(\zeta):\zeta\in\TT\}$.

\begin{proof}
\textbf{1)$\Rightarrow$2).} По Лемме 18.4(б) $\partial G=f(\TT)=[\gamma]$, $\gamma(t)=f(e^{2\pi it})$, $t\in[0,1]$; $\gamma$ непрерывен,
так как продолжение $f$ непрерывно на $\TT$.

\textbf{2)$\Rightarrow$3).} Лемма 18.2.

\textbf{3)$\Rightarrow$4).} $\partial G$ --- компакт ($G$ ограничена), поэтому по Лемме 18.1(б) $\partial G$ равномерно локально связна.
Пусть $\eps>0$; выберем $\delta\in(0,\eps)$ так, что любые $z_1,z_2\in\partial G$ с $|z_1-z_2|<\delta$ лежат в континууме
$Y\subset\partial G$, $\diam Y<\eps$. Пусть $a,b\in\CC\setminus G$, $|a-b|<\delta$.

\begin{detail}{отрезок, не задевающий $\partial G$, лежит вне $G$}
Если связное множество $S$ не пересекает $\partial G$ и содержит точку $p\notin G$, то $S\cap G=\varnothing$: действительно,
$S=(S\cap G)\sqcup(S\cap(\CC\setminus\clos G))$ --- разбиение на два относительно открытых множества, а
$p\in S\setminus G=S\cap(\CC\setminus\clos G)$; по связности $S\cap G=\varnothing$.
\end{detail}

\emph{Случай $[a,b]\cap\partial G=\varnothing$.} По детали $[a,b]\subset\CC\setminus G$; это континуум диаметра $|a-b|<\delta<\eps$.

\emph{Случай $[a,b]\cap\partial G\ne\varnothing$.} Множество $[a,b]\cap\partial G$ компактно; пусть $a'$ и $b'$ --- его ближайшая к $a$
и ближайшая к $b$ точки. Полуинтервал $[a,a')$ связен, не пересекает $\partial G$ и содержит $a\notin G$, поэтому $[a,a')\subset\CC\setminus G$;
и $a'\in\partial G\subset\CC\setminus G$ ($G$ открыто). Итак $[a,a']\subset\CC\setminus G$, аналогично $[b',b]\subset\CC\setminus G$.
Так как $|a'-b'|\le|a-b|<\delta$, найдётся континуум $Y'\subset\partial G$, $a',b'\in Y'$, $\diam Y'<\eps$. Тогда
$Y=[a,a']\cup Y'\cup[b',b]\subset\CC\setminus G$ --- континуум (цепочка связных компактов с общими точками $a'$, $b'$), содержащий $a,b$, и
$\diam Y\le|a-a'|+\diam Y'+|b'-b|<\delta+\eps+\delta<3\eps$.

Таким образом, $\CC\setminus G$ равномерно локально связно (с заменой $\eps$ на $3\eps$), а значит, локально связно в каждой точке
(см. последний абзац доказательства Леммы 18.1(б); он не использует компактность).

\textbf{4)$\Rightarrow$1) --- основная часть.}

\emph{Нормировки.} Сдвигом ($G\mapsto G-f(0)$, при этом $\CC\setminus G$ сдвигается и остаётся локально связным) считаем $f(0)=0$.
Выберем $R_0>0$ и $R>R_0$ так, что $\clos{B(0,R_0)}\subset G$ и $\clos G\subset B(0,R)$.

\emph{Выбор констант.} Пусть $0<\eps<R_0$. Так как $\partial G$ --- компакт в замкнутом множестве $E=\CC\setminus G$, по Лемме 18.1(б)
найдётся $\delta\in(0,\eps)$ такое, что
\[
a\in\partial G,\ b\in\CC\setminus G,\ |a-b|<\delta\ \Longrightarrow\ \exists\ \text{континуум } Y\subset\CC\setminus G:\ a,b\in Y,\ \diam Y<\eps. \tag{18.1}
\]
Выберем $\rho\in(0,\frac14)$ так, что $\dfrac{2\pi R}{\sqrt{\log(1/\rho)}}<\delta$. Заметим, что $\sqrt\rho<\frac12$.

\emph{Цель:} доказать, что для $z_1,z_2\in\DD$ с $\frac12<|z_1|,|z_2|<1$ и $|z_1-z_2|<\rho$ выполнено $|f(z_1)-f(z_2)|<2\eps$.

\begin{detail}{почему цели достаточно}
$f$ равномерно непрерывна на компакте $\clos{B(0,3/4)}$: найдётся $\rho'>0$ с $|f(z)-f(z')|<2\eps$ при $z,z'\in\clos{B(0,3/4)}$,
$|z-z'|<\rho'$. Пусть $\rho_0=\min(\rho,\rho')$ ($\rho<1/4$). Если $z,z'\in\DD$, $|z-z'|<\rho_0$, то либо $|z|,|z'|>1/2$ и работает
цель, либо одна из точек лежит в $\clos{B(0,1/2)}$, и тогда обе лежат в $B(0,3/4)$. Итак, $f$ равномерно непрерывна в $\DD$, и по Лемме
18.4(в) непрерывно продолжается на $\clos\DD$.
\end{detail}

\emph{Выбор окружности.} Пусть $z_1,z_2$ такие, как в цели. По лемме Вольфа (с $U=\DD$, $h=f$, $a=z_1$; $f(\DD)=G\subset B(0,R)$) найдётся
$r\in(\rho,\sqrt\rho)$ с
\[
\lambda:=\ell\bigl(f(\Gamma)\bigr)<\delta,\qquad \Gamma:=L(z_1,r)=\DD\cap C(z_1,r).
\]
Отметим: $r<\frac12<|z_1|$, поэтому $0\notin\clos{B(z_1,r)}$; и $|z_2-z_1|<\rho<r$, т.\,е. $z_2\in B(z_1,r)$.

\begin{detail}{как устроено $\Gamma=\DD\cap C(z_1,r)$}
Пусть $z_1=|z_1|e^{i\varphi}$. Для точки $z=z_1+re^{i(\varphi+\theta)}$ имеем $|z|^2=|z_1|^2+r^2+2r|z_1|\cos\theta$. Положим
$c=\dfrac{1-|z_1|^2-r^2}{2r|z_1|}$; тогда $|z|\ge1\iff\cos\theta\ge c$. Заметим $c>-1$, так как $(|z_1|-r)^2<1$.
\begin{itemize}
\item Если $c>1$ (т.\,е. $|z_1|+r<1$), то $\Gamma=C(z_1,r)$ --- вся окружность, и $\clos{B(z_1,r)}\subset\DD$.
\item Если $c\le1$, положим $\alpha=\arccos c\in[0,\pi)$. Тогда $\Gamma=\{z_1+re^{i(\varphi+\theta)}:\ \alpha<\theta<2\pi-\alpha\}$ ---
\emph{одна} открытая дуга, концы которой $\zeta_\pm=z_1+re^{i(\varphi\pm\alpha)}$ лежат на $\TT$ (при $\alpha=0$ концы совпадают:
окружность касается $\TT$ изнутри).
\end{itemize}
\end{detail}

\emph{Случай 1: $|z_1|+r<1$ (на слайде --- «простой случай, упражнение»).} Тогда $\clos{B(z_1,r)}\subset\DD$, а $f(\Gamma)$ --- замкнутая
кривая длины $\lambda<\delta$. Пусть $\zeta_0\in\Gamma$. Для любой $\zeta\in\Gamma$ точки $f(\zeta)$ и $f(\zeta_0)$ соединены дугой
кривой $f(\Gamma)$ длины $\le\lambda$, поэтому $|f(\zeta)-f(\zeta_0)|\le\lambda<\delta$. По принципу максимума модуля для голоморфной
в окрестности $\clos{B(z_1,r)}$ функции $f-f(\zeta_0)$ имеем $|f(z)-f(\zeta_0)|<\delta$ для всех $z\in B(z_1,r)$, в частности при
$z=z_1,z_2$. Значит $|f(z_1)-f(z_2)|<2\delta<2\eps$.

\emph{Случай 2: $|z_1|+r\ge1$.} Тогда $\Gamma$ --- открытая дуга с концами $\zeta_\pm\in\TT$. Предположим противное:
$|f(z_1)-f(z_2)|\ge2\eps$.

\textit{(i) Концы кривой $f(\Gamma)$.} Положим $\Phi(\theta)=f(z_1+re^{i(\varphi+\theta)})$, $\theta\in(\alpha,2\pi-\alpha)$. Так как
$\int_\alpha^{2\pi-\alpha}|\Phi'(\theta)|\,d\theta=\lambda<\infty$, то при $\theta_1<\theta_2$ имеем
$|\Phi(\theta_2)-\Phi(\theta_1)|\le\int_{\theta_1}^{\theta_2}|\Phi'|\,d\theta$, и по критерию Коши существуют пределы
$a=\lim_{\theta\to\alpha+}\Phi(\theta)$, $b=\lim_{\theta\to(2\pi-\alpha)-}\Phi(\theta)$. По Лемме 18.4(а) (путь $\theta\mapsto z_1+re^{i(\varphi+\theta)}$
стремится к $\zeta_\pm\in\TT$) $a,b\in\partial G$. Переходя к пределу в $|\Phi(\theta_2)-\Phi(\theta_1)|\le\lambda$, получаем
\[
|a-b|\le\lambda<\delta,\qquad |w-a|\le\lambda<\delta<\eps\quad\text{для всех } w\in f(\Gamma). \tag{18.2}
\]
Кроме того, $\clos{f(\Gamma)}=f(\Gamma)\cup\{a,b\}$ (предельная точка последовательности $\Phi(\theta_n)$ --- это либо $\Phi(\theta)$ при
$\theta_n\to\theta$ внутри интервала, либо $a$ или $b$).

\textit{(ii) Маленький континуум.} По (18.1) найдётся континуум $Y\subset\CC\setminus G$, $a,b\in Y$, $\diam Y<\eps$; в частности
$Y\subset B(a,\eps)$. Положим
\[
E_1:=Y\cup f(\Gamma)\quad(\text{компакт, так как } a,b\in Y),\qquad E_2:=\Ch\setminus G .
\]
По (18.2) $E_1\subset B(a,\eps)$. Поскольку $f(\Gamma)\subset G$, а $Y\subset\Ch\setminus G$, имеем $E_1\cap E_2=Y$ --- \emph{связное} множество.

\textit{(iii) Выбор точки.} Так как $|f(z_1)-f(z_2)|\ge2\eps$, обе точки $f(z_1),f(z_2)$ не могут лежать в $B(a,\eps)$; пусть $z_*\in\{z_1,z_2\}$
таково, что $f(z_*)\notin B(a,\eps)$. Заметим: $z_*\in B(z_1,r)$, т.\,е. $z_*\notin\Gamma$.

\textit{(iv) Проверка условий теоремы Янишевского для точек $0$ и $f(z_*)$.}
\begin{itemize}
\item $E_1$ не разделяет $0$ и $f(z_*)$: обе точки лежат в множестве $\Ch\setminus B(a,\eps)$, которое связно (любая его точка соединяется
с $\infty$ радиальным лучом $\{a+t(z-a):t\ge1\}$) и не пересекает $E_1\subset B(a,\eps)$. Здесь $0\notin B(a,\eps)$, так как $|a|\ge R_0>\eps$
($a\in\partial G$, $B(0,R_0)\subset G$).
\item $E_2$ не разделяет $0$ и $f(z_*)$: обе точки лежат в связном множестве $G=\Ch\setminus E_2$.
\end{itemize}
По теореме Янишевского $0$ и $f(z_*)$ не разделены множеством $E_1\cup E_2=f(\Gamma)\cup(\Ch\setminus G)$ (так как $Y\subset E_2$).
Значит, существует путь $\sigma\colon[0,1]\to G\setminus f(\Gamma)$ с $\sigma(0)=0$, $\sigma(1)=f(z_*)$ (компоненты открытого множества
линейно связны).

\textit{(v) Противоречие.} Путь $\tau=f^{-1}\circ\sigma$ лежит в $\DD\setminus\Gamma$ и соединяет $0=f^{-1}(0)$ с $z_*$. Непрерывная функция
$t\mapsto|\tau(t)-z_1|$ равна $|z_1|>r$ при $t=0$ и $|z_*-z_1|<r$ при $t=1$; по теореме о промежуточном значении $|\tau(t_0)-z_1|=r$ при
некотором $t_0$, т.\,е. $\tau(t_0)\in\DD\cap C(z_1,r)=\Gamma$. Противоречие. Значит, $|f(z_1)-f(z_2)|<2\eps$, и цель достигнута.
\end{proof}

\begin{remarkbox}{Замечания}
\begin{itemize}
\item На слайде 326 написано «найдём $r$, $\rho<r<\frac12$»; по лемме Вольфа $r\in(\rho,\sqrt\rho)$, и $\sqrt\rho<\frac12$ при $\rho<\frac14$ ---
это согласовано. Слайд также опускает необходимость брать $\delta$ из \emph{равномерной} локальной связности (Лемма 18.1(б)) ---
иначе $\delta$ зависел бы от точки $a$, которая заранее неизвестна.
\item «Точка $f(z_2)$ (или, возможно, $f(z_1)$)» на слайде --- это наш выбор $z_*$ в шаге (iii).
\item Из теоремы следует нетривиальный топологический факт: граница ограниченной односвязной области, если она локально связна, есть
кривая (частный случай теоремы Хана--Мазуркевича).
\item Ограниченность $G$ существенна в лемме Вольфа (конечная площадь). Для неограниченных областей аналог формулируется на сфере Римана.
\end{itemize}
\end{remarkbox}

\begin{remarkbox}{Коротко: план ответа на экзамене}
\begin{enumerate}
\item Определение 13.5 (локальная связность через континуумы), эквивалентность стандартному определению и равномерной версии для компактов.
\item Примеры: кривая локально связна; $\TT\cup[0,1]$; гребёнка --- не локально связна в $0$.
\item Формулировка теоремы 13.6: 1) продолжение $\iff$ 2) $\partial G$ --- кривая $\iff$ 3) $\partial G$ л.с. $\iff$ 4) $\CC\setminus G$ л.с.
\item 1)$\Rightarrow$2): $\partial G=f(\TT)$; 2)$\Rightarrow$3): кривая л.с.; 3)$\Rightarrow$4): $[a,a']\cup Y'\cup[b',b]$, диаметр $<3\eps$.
\item Лемма Вольфа: КБШ на окружности, $\int_0^\infty\lambda^2dr/r\le2\pi m_2(h(U))\le2\pi^2R^2$, $\int_\rho^{\sqrt\rho}dr/r=\frac12\log\frac1\rho$.
\item 4)$\Rightarrow$1): окружность $C(z_1,r)$ с коротким образом; концы $a,b\in\partial G$ близки; континуум $Y\subset\CC\setminus G$.
\item Янишевский для $E_1=Y\cup f(\Gamma)$, $E_2=\Ch\setminus G$: $0$ и $f(z_*)$ не разделены $f(\Gamma)\cup(\Ch\setminus G)$ $\Rightarrow$
$0$ и $z_*$ не разделены дугой $\Gamma$ в $\DD$ --- противоречие. Значит, $f$ равномерно непрерывна.
\end{enumerate}
\end{remarkbox}
% -------------------------------------------------------------------
\ticket{19}{Точки разреза. Жордановы области. Теорема Каратеодори о гомеоморфном продолжении конформного отображения единичного круга на данную область}
\label{t19:begin}

\begin{defs}
\textbf{Жорданова дуга} --- образ отрезка $[0,1]$ при непрерывном \emph{инъективном} отображении $\gamma\colon[0,1]\to\CC$.
\textbf{Замкнутая жорданова кривая} --- образ окружности $\TT$ при непрерывном инъективном отображении $\TT\to\CC$; эквивалентно ---
носитель пути $\gamma\colon[0,1]\to\CC$ с $\gamma(0)=\gamma(1)$, инъективного на $[0,1)$. Так как $\TT$ компактна, такое
отображение --- гомеоморфизм $\TT$ на кривую. Примеры: окружность, граница квадрата, граница эллипса; «восьмёрка» --- не жорданова
кривая (точка самопересечения).

\textbf{Теорема Жордана} (прошлый семестр, без доказательства). Если $\Gamma$ --- замкнутая жорданова кривая, то $\CC\setminus\Gamma$
состоит ровно из двух компонент связности: ограниченной $\operatorname{Int}\Gamma$ и неограниченной $\operatorname{Ext}\Gamma$, и
$\partial(\operatorname{Int}\Gamma)=\partial(\operatorname{Ext}\Gamma)=\Gamma$.

\textbf{Жорданова область} --- ограниченная область $G\subset\CC$, граница которой --- замкнутая жорданова кривая; тогда
$G=\operatorname{Int}(\partial G)$ (см. деталь ниже), а $G$ односвязна. Примеры: круг, квадрат, внутренность любого многоугольника без
самопересечений. Не жордановы: $\DD\setminus[0,1)$ (граница $\TT\cup[0,1]$ --- не жорданова кривая), гребёнка из билета 18,
кольцо (граница несвязна).

\textbf{Локальная связность} (Определение 13.5, подробнее --- билет 18): замкнутое $E$ локально связно в $a\in E$, если
$\forall\eps>0\ \exists\delta>0$: любые $z_1,z_2\in E\cap B(a,\delta)$ лежат в некотором континууме $Y\subset E\cap B(a,\eps)$.
Любая кривая (в частности, жорданова) локально связна.

\textbf{Точка разреза (Определение 13.6).} Пусть $E$ --- связное локально связное компактное множество, содержащее более одной точки, и
$a\in E$. Точка $a$ называется \emph{точкой разреза} множества $E$, если $E\setminus\{a\}$ несвязно.

\textbf{Примеры.}
1) Замкнутая жорданова кривая $\Gamma$ не имеет точек разреза: $\Gamma\setminus\{a\}=\varphi(\TT\setminus\{\varphi^{-1}(a)\})$ ---
непрерывный образ связной дуги.
2) Каждая внутренняя (не концевая) точка жордановой дуги --- точка разреза: если $a=\gamma(t_0)$, $0<t_0<1$, то
$\gamma([0,1])\setminus\{a\}=\gamma([0,t_0))\sqcup\gamma((t_0,1])$, и оба множества замкнуты в $\gamma([0,1])\setminus\{a\}$ (образы
компактов $\gamma([0,t_0])$, $\gamma([t_0,1])$ без точки $a$). Концы дуги точками разреза не являются.
3) $K=\TT\cup[0,1]$ (граница $\DD\setminus[0,1)$): точки $\zeta\in\TT\setminus\{1\}$ и точка $0$ --- не точки разреза; все точки
$x\in(0,1]$ --- точки разреза. Например, $K\setminus\{1\}=(\TT\setminus\{1\})\sqcup[0,1)$ --- два непересекающихся замкнутых в
$K\setminus\{1\}$ множества; $K\setminus\{x\}$ при $0<x<1$ распадается на $[0,x)$ и $(x,1]\cup\TT$.
4) «Восьмёрка» $C(-1,1)\cup C(1,1)$: единственная точка разреза --- $0$.
5) Отрезок с «усами» --- буква «Т»: точка разветвления --- точка разреза, удаление которой даёт три компоненты.
\end{defs}

\begin{detail}{ограниченная область с жордановой границей есть $\operatorname{Int}$ этой границы}
Пусть $G$ ограниченная область, $\Gamma=\partial G$ --- жорданова кривая. $G$ связна и $G\cap\Gamma=\varnothing$, поэтому $G$ лежит в
одной компоненте $\Omega$ множества $\CC\setminus\Gamma$. $G$ открыто в $\Omega$; $G$ и замкнуто в $\Omega$: $\clos G\cap\Omega=(G\cup\Gamma)\cap\Omega=G$.
По связности $\Omega$ имеем $G=\Omega$. Так как $G$ ограничена, $\Omega=\operatorname{Int}\Gamma$.
\end{detail}

\subsection{Используемые факты}

\begin{thmbox}{Теорема 13.6 (Каратеодори о непрерывном продолжении; билет 18)}
Если $f\in\mathrm{Conf}(\DD,G)$, $G$ ограничена, то $f$ непрерывно продолжается на $\clos\DD$ $\iff$ $\partial G$ --- кривая $\iff$
$\partial G$ локально связна $\iff$ $\CC\setminus G$ локально связно. При этом $f(\TT)=\partial G$ (Лемма 18.4).
\end{thmbox}

\begin{thmbox}{Лемма 19.1 (граничная теорема единственности)}
Пусть $f\in C(\clos\DD)\cap H(\DD)$, $a\in\CC$, и $f\equiv a$ на некоторой дуге $J\subset\TT$ положительной длины. Тогда $f\equiv a$.
\end{thmbox}
\begin{proof}
Пусть дуга $J$ имеет угловую меру $>2\pi/n$, $\omega=e^{2\pi i/n}$. Положим
\[
F(z)=\prod_{k=0}^{n-1}\bigl(f(\omega^kz)-a\bigr).
\]
Тогда
$F\in C(\clos\DD)\cap H(\DD)$, а повёрнутые дуги $\omega^{-k}J$, $k=0,\dots,n-1$, покрывают $\TT$; для $\zeta\in\omega^{-k}J$ множитель с
номером $k$ равен нулю. Значит $F=0$ на $\TT$, и по принципу максимума $F\equiv0$ в $\DD$. Произведение голоморфных функций в связной
области тождественно равно нулю, только если один из множителей $\equiv0$ (иначе нули каждого множителя изолированы, и их объединение
не может быть всей областью). Итак, $f(\omega^kz)\equiv a$ при некотором $k$, т.\,е. $f\equiv a$.
\end{proof}

\subsection{Лемма о точках разреза}

\begin{thmbox}{Лемма 13.8}
Пусть $G\subset\CC$ --- ограниченная область с локально связной границей, $f\in\mathrm{Conf}(\DD,G)$ (по Теореме 13.6 $f\in C(\clos\DD)$).
При $a\in\partial G$ пусть $A=\{\zeta\in\TT:\ f(\zeta)=a\}$, $m=\#A\le\infty$. Тогда $m\ge1$, и точка $a$ является точкой разреза для
$\partial G$ тогда и только тогда, когда $m>1$. При этом
\[
\partial G\setminus\{a\}=\bigcup_jf(I_j),
\]
где $I_j$ --- дуги, на которые окружность $\TT$ разбивается точками множества $A$ (компоненты связности $\TT\setminus A$), и при $m>1$
множества $f(I_j)$ --- это в точности компоненты связности $\partial G\setminus\{a\}$ (в частности, при $m<\infty$ их ровно $m$).
\end{thmbox}

\begin{proof}
\textbf{Предварительно.} $\partial G$ --- связный (как непрерывный образ $f(\TT)$ окружности), локально связный компакт, содержащий более
одной точки (иначе $G=\CC\setminus\{pt\}$ неограничена), так что Определение 13.6 применимо. По Лемме 18.4 $f(\TT)=\partial G$,
поэтому $A\ne\varnothing$, $m\ge1$, и $\partial G\setminus\{a\}=f(\TT\setminus A)=\bigcup_jf(I_j)$. Множество $A=f^{-1}(a)\cap\TT$ замкнуто.
По Лемме 19.1 $A$ не содержит никакой дуги (иначе $f\equiv a$, что противоречит инъективности $f$ в $\DD$).

\textbf{Случай $m=1$.} $A=\{\zeta\}$, и $\partial G\setminus\{a\}=f(\TT\setminus\{\zeta\})$ связно как непрерывный образ связной дуги.
Значит $a$ --- не точка разреза.

\textbf{Случай $m\ge2$.} Тогда $\TT\setminus A$ --- открытое подмножество $\TT$, непустое (так как $A$ не содержит дуг) и отличное от
$\TT\setminus\{pt\}$; его компоненты --- открытые дуги $I_j$ (их не более чем счётно). Возьмём одну из них, $I$, с концами $\alpha,\beta$.
Концы лежат в $A$ (по максимальности $I$ и замкнутости $A$), и $\alpha\ne\beta$: если бы концы совпадали, то $I=\TT\setminus\{\alpha\}$
и $A=\{\alpha\}$, $m=1$.

\emph{Жорданова кривая.} Пусть $C=(\alpha,\beta)$ --- открытая хорда (интервал), $C\subset\DD$. Прямая $\alpha\beta$ делит $\DD\setminus C$
на два выпуклых (значит, связных) открытых множества $W_1,W_2$, причём $\clos{W_1}\cap\TT=\clos I$ и $\clos{W_2}\cap\TT=\TT\setminus I$.
Положим $\sigma(t)=f((1-t)\alpha+t\beta)$, $t\in[0,1]$. Это непрерывный путь, $\sigma(0)=\sigma(1)=a$, на $(0,1)$ инъективный (так как
$f$ инъективна в $\DD$), и $\sigma(t)\in G\not\ni a$ при $0<t<1$. Значит $\Gamma:=[\sigma]=f(C)\cup\{a\}$ --- замкнутая жорданова кривая.
По теореме Жордана $\CC\setminus\Gamma=\Omega_1\sqcup\Omega_2$, области, и $\partial\Omega_1=\partial\Omega_2=\Gamma$.

\emph{$f(W_1)$ и $f(W_2)$ лежат в разных компонентах.} Множество $f(W_k)$ связно и не пересекает $\Gamma$ (оно лежит в $G$, не содержит
точек $f(C)$ по инъективности, и $a\notin G$); значит $f(W_k)$ лежит в одной компоненте $\CC\setminus\Gamma$. Если бы обе лежали в
$\Omega_1$, то $G=f(W_1)\cup f(C)\cup f(W_2)\subset\Omega_1\cup\Gamma$. Возьмём $p\in f(C)\subset G$ и шар $B(p,s)\subset G$. Так как
$p\in\Gamma=\partial\Omega_2$, шар пересекает $\Omega_2$ --- но $\Omega_2\cap(\Omega_1\cup\Gamma)=\varnothing$. Противоречие. Нумеруем так,
что $f(W_1)\subset\Omega_1$, $f(W_2)\subset\Omega_2$.

\emph{Расщепление границы.} По непрерывности $f$ на $\clos\DD$: $f(I)\subset\clos{f(W_1)}\subset\Omega_1\cup\Gamma$. Но $f(I)\subset\partial G\setminus\{a\}$,
а $\Gamma\setminus\{a\}\subset G$; значит $f(I)\subset\Omega_1$. Аналогично $f(\TT\setminus(\clos I\cup A))\subset\Omega_2$
(точки $\TT\setminus I$ являются пределами точек $W_2$). Итак,
\[
\partial G\setminus\{a\}=f(I)\ \sqcup\ f\bigl(\TT\setminus(I\cup A)\bigr),\qquad f(I)\subset\Omega_1,\quad f(\TT\setminus(I\cup A))\subset\Omega_2 ,
\]
и оба множества открыты в $\partial G\setminus\{a\}$ (пересечения с открытыми $\Omega_1$, $\Omega_2$). Оба непусты: $f(I)\ne\varnothing$, а
открытая дуга $\TT\setminus\clos I$ не содержится в $A$ (в $A$ нет дуг), так что $\TT\setminus(I\cup A)\ne\varnothing$. Следовательно,
$\partial G\setminus\{a\}$ несвязно, и $a$ --- точка разреза.

\emph{Компоненты.} Мы показали: для каждой дуги $I=I_j$ множество $f(I_j)$ открыто и замкнуто в $\partial G\setminus\{a\}$ (его дополнение
$f(\TT\setminus(I_j\cup A))$ открыто) и связно (образ дуги). Значит $f(I_j)$ --- компонента связности $\partial G\setminus\{a\}$, а разные
$f(I_j)$, $f(I_k)$ не пересекаются (при $k\ne j$ имеем $I_k\subset\TT\setminus(I_j\cup A)$). При $m<\infty$ точки $A$ делят $\TT$ ровно на
$m$ дуг.
\end{proof}

\begin{remarkbox}{Сравнение со слайдами 332--333}
На слайде дуга $C$ --- дуга окружности в $\DD$, не проходящая через $0$, и используется нормировка $0\in W_2$; мы берём хорду, а то,
что $f(W_1)$ и $f(W_2)$ попадают в разные компоненты, выводим из $\partial\Omega_2=\Gamma$ (на слайде: «непосредственно проверяется»).
Слайд утверждает, что $E_k=\partial G\cap\Omega_k$ непусты, без проверки; непустота $E_2$ требует, чтобы дуга $\TT\setminus\clos I$ не
отображалась целиком в точку $a$ --- это и есть Лемма 19.1. Утверждение слайда 333 «$f(I_j)$ и $f(I_k)$ лежат в разных компонентах
$\CC\setminus\Gamma$» у нас доказано в более сильной форме: $f(I_j)$ --- компоненты $\partial G\setminus\{a\}$.
\end{remarkbox}

\begin{remarkbox}{Пример: $G=\DD\setminus[0,1)$}
Конформное отображение $f$ круга на $\DD\setminus[0,1)$ продолжается непрерывно (граница --- кривая). Каждая точка $x\in(0,1)$ разреза
«видна» с двух берегов, $\#f^{-1}(x)=2$; точка $1$ достигается из верхней и нижней половин окрестности, $\#f^{-1}(1)=2$; кончик $0$
и точки $\TT\setminus\{1\}$ имеют по одному прообразу. Это согласуется с Леммой 13.8 и примером 3 из определений.
\end{remarkbox}

\subsection{Теорема Каратеодори о гомеоморфном продолжении}

\begin{thmbox}{Теорема 13.8 (теорема Каратеодори)}
Пусть $G\subset\CC$ --- ограниченная область, $f\in\mathrm{Conf}(\DD,G)$. Следующие условия эквивалентны:
\begin{enumerate}
\item $f$ имеет непрерывное и инъективное продолжение на $\clos\DD$;
\item $\partial G$ --- замкнутая жорданова кривая (т.\,е. $G$ --- жорданова область);
\item $\partial G$ локально связна и не имеет точек разреза.
\end{enumerate}
При выполнении этих условий продолженное $f$ --- гомеоморфизм $\clos\DD$ на $\clos G$, переводящий $\TT$ на $\partial G$.
\end{thmbox}

\begin{proof}
\textbf{1)$\Rightarrow$2).} По Лемме 18.4(б) $f(\TT)=\partial G$, а $f|_\TT$ непрерывна и инъективна, т.\,е. $\partial G$ ---
замкнутая жорданова кривая.

\textbf{2)$\Rightarrow$3).} Жорданова кривая --- кривая, поэтому локально связна (Лемма 18.2); точек разреза у неё нет (пример 1 в
определениях); она связна и содержит более одной точки.

\textbf{3)$\Rightarrow$1).} Так как $\partial G$ локально связна, по Теореме 13.6 $f\in C(\clos\DD)$ и $f(\TT)=\partial G$. Для каждой
$a\in\partial G$ множество $A=f^{-1}(a)\cap\TT$ непусто, и по Лемме 13.8 $\#A=1$ (иначе $a$ была бы точкой разреза). Итак, $f|_\TT$
инъективна. Далее, $f(\DD)=G$ и $f(\TT)=\partial G$ не пересекаются, а на $\DD$ функция $f$ инъективна; значит, $f$ инъективна на $\clos\DD$.

\textbf{Гомеоморфизм.} $f\colon\clos\DD\to f(\clos\DD)=G\cup\partial G=\clos G$ --- непрерывная биекция компакта на хаусдорфово
пространство, поэтому обратное отображение непрерывно (образ замкнутого, т.\,е. компактного, подмножества компактен, т.\,е. замкнут).
\end{proof}

\begin{remarkbox}{Замечания}
\begin{itemize}
\item В программе сказано «схема доказательства»; ключевые моменты схемы: 1) Теорема 13.6 даёт непрерывность; 2) Лемма 13.8:
число прообразов точки $a$ на $\TT$ равно числу «берегов», и $>1$ ровно для точек разреза; 3) нет точек разреза $\Rightarrow$
инъективность.
\item Для неограниченных областей теорема верна на сфере Римана: если $\partial G\subset\Ch$ --- жорданова кривая в $\Ch$, то
$f$ продолжается до гомеоморфизма $\clos\DD\to\clos G$ (сводится к ограниченному случаю ДЛФ, переводящей точку вне $\clos G$ в $\infty$).
\item Следствия теоремы (граничная нормировка тремя точками, продолжение до гомеоморфизма $\Ch$, теорема Шёнфлиса, произведения Бляшке)
--- билет 20.
\end{itemize}
\end{remarkbox}

\begin{remarkbox}{Коротко: план ответа на экзамене}
\begin{enumerate}
\item Жорданова кривая/область; теорема Жордана; ограниченная область с жордановой границей $=\operatorname{Int}$ границы.
\item Определение 13.6 точки разреза; примеры: жорданова кривая (нет точек разреза), дуга (все внутренние), $\TT\cup[0,1]$.
\item Лемма 13.8: $\#f^{-1}(a)>1\iff a$ --- точка разреза; $\partial G\setminus\{a\}=\bigcup f(I_j)$.
\item Доказательство: $m=1$ --- образ дуги связен; $m\ge2$ --- хорда $C$ между концами дуги $I$, жорданова кривая $f(C)\cup\{a\}$,
по теореме Жордана $f(I)$ и остальная часть границы в разных компонентах.
\item Теорема 13.8: 1)$\Rightarrow$2)$\Rightarrow$3) очевидно; 3)$\Rightarrow$1): Теорема 13.6 + Лемма 13.8.
\item Непрерывная биекция компакта --- гомеоморфизм $\clos\DD\to\clos G$.
\end{enumerate}
\end{remarkbox}
% -------------------------------------------------------------------
\ticket{20}{Теоремы Каратеодори о непрерывном и гомеоморфном продолжении: понятия, формулировки, следствия (граничная нормировка, продолжение до гомеоморфизма $\Ch\to\Ch$, теорема Шёнфлиса)}
\label{t20:begin}

\begin{defs}
\textbf{Конформное отображение.} $\mathrm{Conf}(D,G)$ --- множество голоморфных инъективных отображений области $D$ \emph{на} область $G$;
обратное отображение тоже голоморфно. $\mathrm{Aut}(\DD)=\mathrm{Conf}(\DD,\DD)$; по Предложению 12.2
\[
\mathrm{Aut}(\DD)=\Bigl\{\psi(z)=e^{i\theta}\frac{z-c}{1-\bar cz}:\ \theta\in\RR,\ c\in\DD\Bigr\}
\]
--- это дробно-линейные функции (ДЛФ); они непрерывны на $\clos\DD$ и переводят $\TT$ на $\TT$.

\textbf{Теорема Римана.} Любая односвязная область $G\subsetneq\CC$ конформно эквивалентна $\DD$; более того, для $z_0\in G$
существует единственное $f\in\mathrm{Conf}(\DD,G)$ с $f(0)=z_0$, $f'(0)>0$.

\textbf{Континуум} --- непустой связный компакт. \textbf{Кривая} --- носитель $[\gamma]=\gamma([0,1])$ пути $\gamma$.

\textbf{Локальная связность (Определение 13.5).} Замкнутое $E\subset\CC$ локально связно в $a\in E$, если $\forall\eps>0\ \exists\delta>0$:
любые $z_1,z_2\in E\cap B(a,\delta)$ содержатся в некотором континууме $Y\subset E\cap B(a,\eps)$; $E$ локально связно, если это так в
каждой точке. Для компактов это равносильно стандартному определению (у каждой точки сколь угодно малые связные окрестности) и
равномерной версии ($\delta$ не зависит от $a$) --- см. билет 18, Лемма 18.1. Примеры: любая кривая локально связна; граница
«гребёнки» $\{|x|<1,0<y<1\}\setminus\bigcup_{n\ge2}[i/n,1+i/n]$ не локально связна в $0$.

\textbf{Точка разреза (Определение 13.6).} Для связного локально связного компакта $E$ (из более чем одной точки) точка $a\in E$ ---
точка разреза, если $E\setminus\{a\}$ несвязно. У замкнутой жордановой кривой точек разреза нет; у дуги все внутренние точки --- точки разреза.

\textbf{Жорданова кривая и область.} Замкнутая жорданова кривая --- гомеоморфный образ $\TT$. По теореме Жордана
$\CC\setminus\Gamma=\operatorname{Int}\Gamma\sqcup\operatorname{Ext}\Gamma$ (ограниченная и неограниченная области), обе с границей $\Gamma$.
Жорданова область --- $G=\operatorname{Int}\Gamma$; это в точности ограниченные области с жордановой границей. Жорданова область
односвязна: её дополнение в $\Ch$, равное $\Gamma\cup\operatorname{Ext}\Gamma\cup\{\infty\}=\clos{\operatorname{Ext}\Gamma\cup\{\infty\}}$,
связно (замыкание связного множества). Поэтому к ней применима теорема Римана.

\textbf{Циклический порядок на $\TT$.} Тройка различных точек $(\zeta_1,\zeta_2,\zeta_3)$ на $\TT$ \emph{положительно ориентирована}, если
$\zeta_j=e^{it_j}$ с $t_1<t_2<t_3<t_1+2\pi$ (точки встречаются в этом порядке при обходе против часовой стрелки). Для трёх различных точек
ровно одна из тройек $(\zeta_1,\zeta_2,\zeta_3)$, $(\zeta_1,\zeta_3,\zeta_2)$ положительно ориентирована; свойство не меняется при
циклической перестановке.

\textbf{Циклический порядок на жордановой кривой.} Пусть $G$ --- жорданова область, $h\in\mathrm{Conf}(\DD,G)$, продолженное до гомеоморфизма
$\clos\DD\to\clos G$ (Теорема 13.8). Различные точки $z_1,z_2,z_3\in\partial G$ \emph{имеют положительный циклический порядок}, если
тройка $(h^{-1}(z_1),h^{-1}(z_2),h^{-1}(z_3))$ положительно ориентирована на $\TT$. Корректность (независимость от $h$) --- Лемма 20.2 ниже.
Геометрически это значит: точки встречаются в порядке $z_1,z_2,z_3$ при обходе $\partial G$ так, что $G$ остаётся слева.
\end{defs}

\subsection{Основные теоремы (формулировки)}

\begin{thmbox}{Теорема 13.6 (Каратеодори о непрерывном продолжении)}
Пусть $f$ конформно отображает $\DD$ на ограниченную область $G$. Эквивалентны:
1) $f$ непрерывно продолжается на $\clos\DD$; 2) $\partial G$ --- кривая; 3) $\partial G$ локально связна; 4) $\CC\setminus G$ локально связно.
При этом $f(\TT)=\partial G$.
\end{thmbox}

\begin{thmbox}{Теорема 13.8 (Каратеодори о гомеоморфном продолжении)}
Пусть $G\subset\CC$ --- ограниченная область, $f\in\mathrm{Conf}(\DD,G)$. Эквивалентны:
1) $f$ имеет непрерывное инъективное продолжение на $\clos\DD$; 2) $\partial G$ --- замкнутая жорданова кривая; 3) $\partial G$ локально
связна и не имеет точек разреза. При этом продолжение --- гомеоморфизм $\clos\DD\to\clos G$, $f(\TT)=\partial G$.
\end{thmbox}

Идеи доказательств (подробно --- билеты 18, 19): для Теоремы 13.6 главное --- 4)$\Rightarrow$1): лемма Вольфа (принцип длина--площадь) даёт
окружности $C(z_1,r)$ со сколь угодно коротким $f$-образом; его концы $a,b\in\partial G$ близки и соединяются малым континуумом
$Y\subset\CC\setminus G$; теорема Янишевского показывает, что иначе $0$ и $z_2$ не разделялись бы дугой $\DD\cap C(z_1,r)$; отсюда
равномерная непрерывность $f$. Для Теоремы 13.8: Лемма 13.8 --- точка $a\in\partial G$ является точкой разреза тогда и только тогда, когда у неё
более одного прообраза на $\TT$.

\subsection{Вспомогательные факты о ДЛФ и циклическом порядке}

\begin{thmbox}{Лемма 20.1}
Пусть $(\zeta_1,\zeta_2,\zeta_3)$ --- различные точки $\TT$ и
\[
S(z)=\frac{(z-\zeta_1)(\zeta_2-\zeta_3)}{(z-\zeta_3)(\zeta_2-\zeta_1)}
\]
--- ДЛФ, переводящая $\zeta_1,\zeta_2,\zeta_3$ в $0,1,\infty$. Тогда $S(\TT)=\RR\cup\{\infty\}$, и $S(\DD)=\{\Im w>0\}$, если тройка положительно
ориентирована, и $S(\DD)=\{\Im w<0\}$ в противном случае.
\end{thmbox}
\begin{proof}
ДЛФ переводит окружности в $\Ch$ в окружности; $S(\TT)$ --- окружность через $0,1,\infty$, т.\,е. $\RR\cup\{\infty\}$. Так как $S$ ---
гомеоморфизм $\Ch$, связное $\DD$ переходит в компоненту $\Ch\setminus(\RR\cup\{\infty\})$, т.\,е. в одну из полуплоскостей $H_\pm=\{\pm\Im w>0\}$.

Пусть тройка положительно ориентирована, $\zeta_j=e^{it_j}$, $t_1<t_2<t_3<t_1+2\pi$. Функция $u(t)=S(e^{it})$ на $[t_1,t_3)$ непрерывна,
вещественна (не принимает значения $\infty$, так как $e^{it}\ne\zeta_3$) и инъективна; значит, строго монотонна, а так как $u(t_1)=0<1=u(t_2)$,
она возрастает. Поэтому $u'(t_2)=S'(\zeta_2)\,i\zeta_2=:\kappa\ge0$, а так как $S'(\zeta_2)\ne0$, то $\kappa>0$. Теперь при малых $s>0$
точка $(1-s)\zeta_2\in\DD$ и
\[
S((1-s)\zeta_2)=1-sS'(\zeta_2)\zeta_2+o(s)=1+is\,\kappa+o(s)
\]
(так как $-\zeta_2S'(\zeta_2)=i\cdot i\zeta_2S'(\zeta_2)=i\kappa$). Значит $\Im S((1-s)\zeta_2)>0$, т.\,е. $S(\DD)\cap H_+\ne\varnothing$ и $S(\DD)=H_+$.
Если тройка ориентирована отрицательно, то при движении \emph{по часовой стрелке} от $\zeta_1$ через $\zeta_2$ к $\zeta_3$ то же
рассуждение для $v(t)=S(e^{-it})$ даёт $S'(\zeta_2)(-i\zeta_2)>0$, т.\,е. $\kappa<0$, и $S(\DD)=H_-$.
\end{proof}

\begin{thmbox}{Лемма 20.2 (три точки на окружности)}
(а) Если тройки различных точек $(\zeta_j)$ и $(\eta_j)$ на $\TT$ обе положительно ориентированы, то существует единственный
$\psi\in\mathrm{Aut}(\DD)$ с $\psi(\zeta_j)=\eta_j$, $j=1,2,3$.
(б) Любой $\psi\in\mathrm{Aut}(\DD)$ сохраняет положительную ориентацию троек на $\TT$. Поэтому циклический порядок на жордановой кривой
определён корректно.
(в) ДЛФ, имеющая три различные неподвижные точки, --- тождественное отображение.
\end{thmbox}
\begin{proof}
(в) Если $\frac{az+b}{cz+d}=z$ в трёх точках, то $cz^2+(d-a)z-b=0$ имеет три корня (если один из них $\infty$, то $c=0$ и линейное
уравнение $(d-a)z-b=0$ имеет два корня), значит, все коэффициенты нулевые: $c=b=0$, $a=d$, и отображение тождественно.

(а) Пусть $S_\zeta$, $S_\eta$ --- ДЛФ из Леммы 20.1 для этих троек; обе переводят $\DD$ на $H_+$. Тогда $\psi=S_\eta^{-1}\circ S_\zeta$ ---
ДЛФ, переводящая $\DD$ на $\DD$ и $\zeta_j\mapsto\eta_j$, т.\,е. $\psi\in\mathrm{Aut}(\DD)$. Если $\psi_1$ --- ещё один такой автоморфизм,
то $\psi^{-1}\circ\psi_1$ --- ДЛФ (Предложение 12.2) с тремя неподвижными точками $\zeta_j$, т.\,е. тождество по (в).

(б) Пусть $(\zeta_j)$ положительно ориентирована, а $(\psi(\zeta_j))$ --- нет. Тогда по Лемме 20.1 $S_{\psi(\zeta)}\circ\psi\circ S_\zeta^{-1}$
переводит $H_+$ на $H_-$ и оставляет на месте $0,1,\infty$; по (в) это тождество --- противоречие.
Корректность: если $h,h_1\in\mathrm{Conf}(\DD,G)$, то $\psi=h^{-1}\circ h_1\in\mathrm{Aut}(\DD)$, и на $\TT$ (по непрерывности продолжений)
$h_1^{-1}=\psi^{-1}\circ h^{-1}$; ориентация троек $h^{-1}(z_j)$ и $h_1^{-1}(z_j)$ одинакова по доказанному.
\end{proof}

\subsection{Следствие 1: граничное условие нормировки}

\begin{thmbox}{Следствие (граничное условие нормировки для конформного отображения)}
Пусть $G$ и $B$ --- жордановы области, точки $z_1,z_2,z_3\in\partial G$ различны, точки $w_1,w_2,w_3\in\partial B$ различны, и тройки
имеют одинаковый циклический порядок. Тогда существует единственное $f\in\mathrm{Conf}(G,B)$ (продолженное по непрерывности на $\clos G$)
с условиями $f(z_j)=w_j$, $j=1,2,3$. Это продолжение --- гомеоморфизм $\clos G\to\clos B$.
\end{thmbox}
\begin{proof}
Пусть $h\in\mathrm{Conf}(\DD,G)$ и $g\in\mathrm{Conf}(\DD,B)$ (теорема Римана). По Теореме 13.8 $h$ и $g$ продолжаются до гомеоморфизмов
$\clos\DD\to\clos G$ и $\clos\DD\to\clos B$, переводящих $\TT$ на $\partial G$ и $\partial B$. Точки $\zeta_j=h^{-1}(z_j)$ и $\eta_j=g^{-1}(w_j)$
--- различные точки $\TT$; по определению одинакового циклического порядка тройки $(\zeta_j)$ и $(\eta_j)$ ориентированы одинаково. Если обе
ориентированы положительно, по Лемме 20.2(а) есть $\psi\in\mathrm{Aut}(\DD)$ с $\psi(\zeta_j)=\eta_j$. (Если обе ориентированы
отрицательно, применим Лемму 20.2(а) к тройкам $(\zeta_1,\zeta_3,\zeta_2)$ и $(\eta_1,\eta_3,\eta_2)$.)

\emph{Существование.} $f=g\circ\psi\circ h^{-1}$ --- композиция конформных отображений $G\to\DD\to\DD\to B$, т.\,е. $f\in\mathrm{Conf}(G,B)$;
композиция продолжений --- гомеоморфизм $\clos G\to\clos B$, и $f(z_j)=g(\psi(\zeta_j))=g(\eta_j)=w_j$.

\emph{Единственность.} Пусть $f_1\in\mathrm{Conf}(G,B)$ непрерывно продолжается на $\clos G$ и $f_1(z_j)=w_j$. Функция
$F=\psi^{-1}\circ g^{-1}\circ f_1\circ h$ принадлежит $\mathrm{Aut}(\DD)$, т.\,е. является ДЛФ (Предложение 12.2). На $\TT$ продолжение
$\psi^{-1}\circ g^{-1}\circ f_1\circ h$ (непрерывное на $\clos\DD$) совпадает с этой ДЛФ (по непрерывности обеих на $\clos\DD$), поэтому
$F(\zeta_j)=\psi^{-1}(g^{-1}(w_j))=\psi^{-1}(\eta_j)=\zeta_j$. По Лемме 20.2(в) $F=\mathrm{id}$, откуда $f_1=g\circ\psi\circ h^{-1}=f$.
\end{proof}

\begin{remarkbox}{Замечания}
\begin{itemize}
\item Условие одинакового циклического порядка необходимо: конформное $f$ задаёт $g^{-1}\circ f\circ h\in\mathrm{Aut}(\DD)$, а автоморфизмы сохраняют
ориентацию троек (Лемма 20.2(б)).
\item Сравните с внутренней нормировкой теоремы Римана ($f(z_0)=w_0$, $\arg f'(z_0)$ задан): тоже три вещественных параметра. Нормировку
можно задавать и смешанно: одна внутренняя точка и одна граничная.
\end{itemize}
\end{remarkbox}

\subsection{Следствие 2: продолжение до гомеоморфизма $\Ch\to\Ch$}

\begin{thmbox}{Следствие}
Пусть $G$ --- жорданова область и $f\in\mathrm{Conf}(\DD,G)$. Тогда $f$ продолжается до гомеоморфизма $\Ch\to\Ch$ (с $f(\infty)=\infty$).
\end{thmbox}
\begin{proof}
Пусть $\Gamma=\partial G$, $\DD^*=\{|z|>1\}\cup\{\infty\}$, $G^*=\operatorname{Ext}\Gamma\cup\{\infty\}$.

\textbf{Шаг 1: отображение внешностей.} Фиксируем $p\in G$. ДЛФ $T(w)=\frac1{w-p}$ --- гомеоморфизм $\Ch$, $T(\infty)=0$, $T(p)=\infty$.
Так как $\dist(p,\Gamma)>0$, кривая $\Gamma'=T(\Gamma)$ --- ограниченная жорданова кривая. Множество $G'=T(G^*)$ --- связное открытое, не
содержит $\infty=T(p)$ (так как $p\notin G^*$) и ограничено: $\clos{G^*}=\operatorname{Ext}\Gamma\cup\Gamma\cup\{\infty\}$ --- компакт в $\Ch$, не
содержащий $p$, поэтому $|w-p|\ge d_0>0$ для $w\in G^*\cap\CC$ и $|T|\le1/d_0$ на $G^*$. Граница $\partial G'=T(\partial G^*)=T(\Gamma)=\Gamma'$.
Значит, $G'$ --- жорданова область, $0=T(\infty)\in G'$. По теореме Римана есть $F\in\mathrm{Conf}(\DD,G')$ с $F(0)=0$; по Теореме 13.8
$F$ --- гомеоморфизм $\clos\DD\to\clos{G'}$. Положим
\[
f^*(z)=T^{-1}\bigl(F(1/z)\bigr)=p+\frac1{F(1/z)},\qquad z\in\clos{\DD^*},
\]
($1/\infty=0$, $F(0)=0$, $1/0=\infty$). Тогда $f^*$ --- гомеоморфизм $\clos{\DD^*}=\{|z|\ge1\}\cup\{\infty\}$ на
$\clos{G^*}=G^*\cup\Gamma$, $f^*(\infty)=\infty$, $f^*(\TT)=\Gamma$, и $f^*$ конформно отображает $\{|z|>1\}$ на $\operatorname{Ext}\Gamma$.

\textbf{Шаг 2: склейка на окружности.} По Теореме 13.8 $f$ --- гомеоморфизм $\clos\DD\to\clos G$, $f(\TT)=\Gamma$. Положим
$\varphi=(f^*|_\TT)^{-1}\circ f|_\TT$ --- гомеоморфизм $\TT\to\TT$. Определим
\[
\Phi(z)=\begin{cases}f(z),&|z|\le1,\\ f^*(r\varphi(\zeta)),& z=r\zeta,\ r\ge1,\ \zeta\in\TT,\\ \infty,&z=\infty.\end{cases}
\]
При $r=1$ оба определения дают $f^*(\varphi(\zeta))=f(\zeta)$ --- согласованы.

\textbf{Шаг 3: $\Phi$ --- гомеоморфизм.} Отображение $\Lambda(r\zeta)=r\varphi(\zeta)$ ($r\ge1$), $\Lambda(\infty)=\infty$, --- непрерывная биекция
$\clos{\DD^*}$ на себя (обратное $r\eta\mapsto r\varphi^{-1}(\eta)$), непрерывная и в $\infty$ (так как $|\Lambda(z)|=|z|$). Поэтому
$\Phi|_{\clos{\DD^*}}=f^*\circ\Lambda$ непрерывно, и $\Phi|_{\clos\DD}=f$ непрерывно; так как замкнутые множества $\clos\DD$ и $\clos{\DD^*}$
покрывают $\Ch$, а на пересечении $\TT$ определения совпадают, $\Phi$ непрерывно на $\Ch$ (лемма о склейке). Биективность:
$\Phi$ переводит $\DD$ на $G$, $\TT$ на $\Gamma$, $\{1<|z|<\infty\}$ на $\operatorname{Ext}\Gamma$, $\infty$ в $\infty$, каждый раз
биективно, а $\Ch=G\sqcup\Gamma\sqcup\operatorname{Ext}\Gamma\sqcup\{\infty\}$. Непрерывная биекция компакта $\Ch$ на хаусдорфово $\Ch$ ---
гомеоморфизм.
\end{proof}

\begin{remarkbox}{Замечание}
Продолжение $\Phi$ голоморфно в $\DD$, но в $\{|z|>1\}$, вообще говоря, лишь непрерывно (там оно равно $f^*\circ\Lambda$, а $\Lambda$ не
голоморфно, если $\varphi$ --- не поворот). Конформного продолжения через $\TT$, вообще говоря, нет.
\end{remarkbox}

\subsection{Следствие 3: теорема Шёнфлиса}

\begin{thmbox}{Следствие (теорема Шёнфлиса)}
Пусть $\Gamma$ и $\Upsilon$ --- замкнутые жордановы кривые, а $F\colon\Gamma\to\Upsilon$ --- непрерывное биективное отображение. Тогда $F$
продолжается до гомеоморфизма $\CC\to\CC$ (и $\Ch\to\Ch$), переводящего $\operatorname{Int}\Gamma$ на $\operatorname{Int}\Upsilon$.
\end{thmbox}
\begin{proof}
Пусть $G=\operatorname{Int}\Gamma$, $B=\operatorname{Int}\Upsilon$, $h\in\mathrm{Conf}(\DD,G)$, $g\in\mathrm{Conf}(\DD,B)$. По Следствию 2 $h$ и $g$
продолжаются до гомеоморфизмов $H,\ \widehat g\colon\Ch\to\Ch$ с $H(\infty)=\widehat g(\infty)=\infty$, $H(\TT)=\Gamma$, $\widehat g(\TT)=\Upsilon$,
$H(\DD)=G$, $\widehat g(\DD)=B$.

$F$ --- непрерывная биекция компакта на компакт, т.\,е. гомеоморфизм. Значит, $\varphi=g^{-1}\circ F\circ h\colon\TT\to\TT$ --- гомеоморфизм
(здесь $h,g$ --- продолжения на $\TT$). Положим
\[
\widetilde\varphi(r\zeta)=r\varphi(\zeta)\ \ (r>0,\ \zeta\in\TT),\qquad\widetilde\varphi(0)=0,\qquad\widetilde\varphi(\infty)=\infty .
\]
$\widetilde\varphi$ --- биекция $\Ch$ на себя (обратное строится по $\varphi^{-1}$), непрерывная вне $0,\infty$ очевидно, а в $0$ и $\infty$ ---
так как $|\widetilde\varphi(z)|=|z|$. Значит, $\widetilde\varphi$ --- гомеоморфизм $\Ch$; он переводит $\DD$ на $\DD$ и $\TT$ на $\TT$.

Положим $\Psi=\widehat g\circ\widetilde\varphi\circ H^{-1}\colon\Ch\to\Ch$ --- гомеоморфизм, $\Psi(\infty)=\infty$, поэтому $\Psi|_\CC$ ---
гомеоморфизм $\CC\to\CC$. Для $w\in\Gamma$: $H^{-1}(w)=h^{-1}(w)\in\TT$, и
\[
\Psi(w)=g\bigl(\varphi(h^{-1}(w))\bigr)=g\bigl(g^{-1}(F(h(h^{-1}(w))))\bigr)=F(w).
\]
Наконец, $\Psi(G)=\widehat g(\widetilde\varphi(\DD))=\widehat g(\DD)=B$.
\end{proof}

\begin{remarkbox}{Смысл теоремы Шёнфлиса}
Теорема Жордана утверждает лишь, что жорданова кривая делит плоскость на две области. Теорема Шёнфлиса сильнее: замыкание
внутренности жордановой кривой гомеоморфно замкнутому кругу, и вся конфигурация $(\CC,\Gamma)$ гомеоморфна $(\CC,\TT)$. (В $\RR^3$
аналог неверен: «рогатая сфера Александера».) Здесь она получена средствами ТФКП: теорема Римана $+$ теорема Каратеодори.
\end{remarkbox}

\subsection{Следствие 4: голоморфные функции с граничными значениями на жордановой кривой}

\begin{thmbox}{Следствие (слайды 336--337)}
Пусть $G\subset\CC$ --- жорданова область, $f\in H(\DD)$, $f\not\equiv\mathrm{const}$, и $\dist(f(z),\partial G)\to0$ при $|z|\to1$.
Тогда $f(\DD)=G$, $f$ непрерывно продолжается на $\clos\DD$, и существует $m\in\NN$ такое, что $f$ принимает каждое значение из $G$ ровно
$m$ раз в $\DD$ (с учётом кратности). Именно, $f=\phi\circ B$, где $\phi\in\mathrm{Conf}(\DD,G)$ и
\[
B(z)=c\prod_{k=1}^m\frac{z-w_k}{1-\bar w_kz},\qquad |c|=1,\ w_k\in\DD,
\]
--- конечное произведение Бляшке.
\end{thmbox}
\begin{proof}
\textbf{Шаг 1: $f(\DD)=G$.} Условие $\dist(f(z),\partial G)\to0$ означает: если $z_n\in\DD$, $|z_n|\to1$, то все предельные точки
$f(z_n)$ (в $\Ch$) лежат на $\Gamma=\partial G$ (с компактной $\Gamma$ расстояние $\to0$ исключает и $f(z_n)\to\infty$). Это условие
Предложения 13.3 (билет 3) для $\Omega=\DD$ и жордановой кривой $\Gamma$; по нему $f(\DD)=\operatorname{Int}\Gamma=G$.

\textbf{Шаг 2: $g=\phi^{-1}\circ f$.} Пусть $\phi\in\mathrm{Conf}(\DD,G)$; по Теореме 13.8 $\phi$ --- гомеоморфизм $\clos\DD\to\clos G$.
Положим $g=\phi^{-1}\circ f\in H(\DD)$, $g(\DD)=\DD$, $g\not\equiv$ const. Покажем, что $|g(z)|\to1$ при $|z|\to1$. Пусть $\eps\in(0,1)$;
компакт $\phi(\clos{B(0,1-\eps)})\subset G$ находится на положительном расстоянии $\eta$ от $\partial G$. Если $|z|$ близко к $1$
настолько, что $\dist(f(z),\partial G)<\eta$, то $f(z)\notin\phi(\clos{B(0,1-\eps)})$, т.\,е. $|g(z)|>1-\eps$.

\textbf{Шаг 3: $g$ --- конечное произведение Бляшке (на слайде --- «задача»).} Найдём $r_0<1$ с $|g(z)|>\frac12$ при $r_0\le|z|<1$. Нули
$g$ лежат в компакте $\clos{B(0,r_0)}$ и изолированы ($g\not\equiv0$), поэтому их конечное число; пусть $w_1,\dots,w_m$ --- нули с учётом
кратности. Положим $B_0(z)=\prod_{k=1}^m\frac{z-w_k}{1-\bar w_kz}$ (при $m=0$ $B_0\equiv1$). $B_0$ --- рациональная функция с полюсами
$1/\bar w_k\notin\clos\DD$, $|B_0|=1$ на $\TT$ (так как $|\zeta-w|=|1-\bar w\zeta|$ при $|\zeta|=1$), поэтому $|B_0(z)|\to1$ при $|z|\to1$
равномерно. Функция $q=g/B_0$ голоморфна в $\DD$ (нули сокращаются с учётом кратности), без нулей, и $|q(z)|\to1$ при $|z|\to1$. По
принципу максимума для $q$ и $1/q$: при $|z|\le r$ имеем $|q(z)|\le\max_{|\zeta|=r}|q(\zeta)|\to1$ ($r\to1$), т.\,е. $|q|\le1$; аналогично
$|1/q|\le1$. Итак, $|q|\equiv1$, и по принципу максимума (или открытости) $q\equiv c$, $|c|=1$; $g=cB_0=B$. Наконец, $m\ge1$: иначе $g=c$ ---
константа, а $g$ не постоянна.

\textbf{Шаг 4: $g$ принимает каждое значение $a\in\DD$ ровно $m$ раз.} Возьмём $r<1$ так, что $|g(z)|>|a|$ при $r\le|z|<1$ (возможно,
так как $|g|\to1$ и $|a|<1$). Тогда все нули $g$ и $g-a$ лежат в $B(0,r)$, а на $|z|=r$: $|(g-a)-g|=|a|<|g|$. По теореме Руше $g-a$ и
$g$ имеют в $B(0,r)$ (а значит, и в $\DD$) одинаковое число нулей с учётом кратности, т.\,е. $m$.

\textbf{Шаг 5: выводы для $f$.} $f=\phi\circ g$; $g$ непрерывна на $\clos\DD$ (рациональная функция без полюсов в $\clos\DD$), $g(\clos\DD)\subset\clos\DD$,
$\phi$ непрерывна на $\clos\DD$; значит, $f$ непрерывно продолжается на $\clos\DD$. Для $w\in G$ уравнение $f(z)=w$ равносильно $g(z)=a$,
$a=\phi^{-1}(w)$, а кратности совпадают: $f-w=\phi(g)-\phi(a)=(g-a)\,Q$, где $Q(z)=\frac{\phi(g(z))-\phi(a)}{g(z)-a}$ (при $g(z)=a$ полагаем
$Q=\phi'(a)$) голоморфна и не обращается в нуль ($\phi$ инъективна, $\phi'\ne0$). Итак, $f$ принимает каждое значение из $G$ ровно $m$ раз.
\end{proof}

\begin{remarkbox}{Замечания к слайдам 336--337}
\begin{itemize}
\item В формуле произведения Бляшке на слайде 337 опечатка: знаменатель «$1-zw\,z$» следует читать как $1-\bar w_kz$.
\item «Ровно $m$ раз» --- с учётом кратности: например, $B(z)=z^2$ ($m=2$) принимает значение $0$ один раз, но с кратностью $2$. Различных
прообразов меньше $m$ лишь для конечного числа значений (критических значений $B$).
\item Обратно, при любом $\phi\in\mathrm{Conf}(\DD,G)$ и конечном произведении Бляшке $B$ функция $f=\phi\circ B$ удовлетворяет условию
следствия. Таким образом, такие $f$ --- это в точности $\phi\circ B$; при $m=1$ получаем конформные отображения $\DD$ на $G$.
\end{itemize}
\end{remarkbox}

\begin{remarkbox}{Коротко: план ответа на экзамене}
\begin{enumerate}
\item Понятия: континуум, кривая, локальная связность (Опр. 13.5), точка разреза (Опр. 13.6), жорданова кривая/область, теорема Жордана.
\item Теорема 13.6: продолжение на $\clos\DD$ $\iff$ $\partial G$ --- кривая $\iff$ $\partial G$ л.с. $\iff$ $\CC\setminus G$ л.с.
\item Теорема 13.8: продолжение до гомеоморфизма $\iff$ $\partial G$ жорданова $\iff$ л.с. без точек разреза.
\item Граничная нормировка: $f=g\circ\psi\circ h^{-1}$, $\psi\in\mathrm{Aut}(\DD)$ по трём граничным точкам; единственность --- ДЛФ с тремя неподвижными точками.
\item Продолжение до гомеоморфизма $\Ch$: внешнее отображение $f^*$, $\varphi=f^{*-1}\circ f$ на $\TT$, $\Phi(r\zeta)=f^*(r\varphi(\zeta))$ при $r\ge1$.
\item Шёнфлис: $\Psi=\widehat g\circ\widetilde\varphi\circ H^{-1}$, $\widetilde\varphi(r\zeta)=r\varphi(\zeta)$, $\varphi=g^{-1}\circ F\circ h$.
\item $f\in H(\DD)$, $\dist(f,\partial G)\to0$ $\Rightarrow$ $f=\phi\circ B$, $B$ --- произведение Бляшке из $m$ множителей.
\end{enumerate}
\end{remarkbox}

% ===== part 4 =====
% ====================== Часть 4: билеты 21--24 (гармонические функции) ======================
\providecommand{\pIVHarm}{\operatorname{Harm}}
\providecommand{\pIVConf}{\operatorname{Conf}}

\ticket{21}{Гармонические функции и их основные свойства (теорема о среднем, принцип максимума и минимума)}
\label{t21:start}

\begin{defs}
\textbf{Обозначения.} Плоскость $\RR^2$ отождествляем с $\CC$: $z=x+iy\leftrightarrow(x,y)$.
$B(a,R)=\{z:|z-a|<R\}$ --- открытый круг, $\clos{B(a,R)}=\{|z-a|\le R\}$ --- замкнутый круг.
\emph{Область} --- непустое открытое связное множество $D\subset\CC$; $\clos D$ --- замыкание, $\partial D=\clos D\setminus D$ --- граница.
Если $D$ ограничена, то $\clos D$ и $\partial D$ --- компакты, причём $\partial D\ne\varnothing$ (иначе $D=\clos D$ было бы
непустым открыто-замкнутым собственным подмножеством связного $\CC$). Для функции $h$ и множества $K$ пишем
$\|h\|_K=\sup_{K}|h|$.

\textbf{Классы гладкости.} $C^k(D)$ --- функции $D\to\RR$ (или $\CC$), у которых все частные производные порядка $\le k$
существуют и непрерывны в $D$; $C^\infty(D)=\bigcap_k C^k(D)$. Для $h\in C^2(D)$ смешанные производные равны: $h''_{xy}=h''_{yx}$
(теорема Шварца).

\textbf{Оператор Лапласа:} $\Delta h=h''_{xx}+h''_{yy}$.

\textbf{Определение 14.1.} Пусть $D$ --- область в $\RR^2$, $h\colon D\to\RR$, $h\in C^2(D)$. Если $\Delta h=0$ всюду в $D$, то $h$
называется \emph{гармонической функцией} в $D$. Класс всех гармонических в $D$ функций обозначается $\pIVHarm(D)$.
Для компакта $K\subset\CC$ пишут $h\in\pIVHarm(K)$, если $h\in\pIVHarm(U)$ для некоторого открытого $U\supset K$
(для несвязного $U$ это означает гармоничность на каждой компоненте связности).

\emph{Примеры:} $x$, $y$, $x^2-y^2$, $xy$, $e^x\cos y$, $\ln|z|$ (в $\CC\setminus\{0\}$) гармоничны;
$x^2$, $|z|^2=x^2+y^2$ ($\Delta=4$) --- нет. $\pIVHarm(D)$ --- вещественное векторное пространство (оператор $\Delta$ линеен),
но не алгебра: $x\in\pIVHarm(\CC)$, а $x\cdot x=x^2\notin\pIVHarm(\CC)$.

\textbf{Производные Виртингера.} Для (вещественно) дифференцируемой $F$:
\[
\partial_z F=F'_z=\tfrac12\bigl(F'_x-iF'_y\bigr),\qquad \partial_{\bar z}F=F'_{\bar z}=\tfrac12\bigl(F'_x+iF'_y\bigr).
\]
Если $h$ вещественна, то $h'_{\bar z}=\overline{h'_z}$. Всегда $\partial_{\bar z}\clos F=\overline{\partial_z F}$, $\partial_z\clos F=\overline{\partial_{\bar z}F}$.

\textbf{Голоморфные функции.} $H(D)$ --- класс голоморфных в $D$ функций. \emph{Критерий (условия Коши--Римана):}
функция $f=u+iv$, вещественно дифференцируемая в $D$, голоморфна $\iff$ $f'_{\bar z}=0$ в $D$ $\iff$ $u'_x=v'_y$, $u'_y=-v'_x$.
При этом $f'=f'_z=f'_x=u'_x+iv'_x=u'_x-iu'_y$ и $f'_y=if'$. Голоморфные функции бесконечно дифференцируемы
($f'\in H(D)$), в частности $u,v\in C^\infty(D)$.

\textbf{Односвязная область} --- область, в которой всякий замкнутый путь гомотопен постоянному (для плоских областей
эквивалентно связности $\Ch\setminus D$). Круг односвязен. \emph{Теорема о первообразной:} в односвязной области всякая
$g\in H(D)$ имеет первообразную $F\in H(D)$, $F'=g$ (следствие интегральной теоремы Коши).

\textbf{Теорема о среднем для голоморфных функций:} если $f\in H(U)$, $\clos{B(a,r)}\subset U$, то
$f(a)=\frac1{2\pi}\int_0^{2\pi}f(a+re^{i\vartheta})\,d\vartheta$ (это интегральная формула Коши
$f(a)=\frac1{2\pi i}\oint_{|z-a|=r}\frac{f(z)\,dz}{z-a}$ после подстановки $z=a+re^{i\vartheta}$, $dz=ire^{i\vartheta}d\vartheta$).

\textbf{Теорема единственности для голоморфных функций:} если $g\in H(D)$ и множество нулей $g$ имеет предельную точку в $D$,
то $g\equiv0$ в $D$.

\textbf{Критическая точка} функции $h\in C^1(D)$ --- точка $a$, где $\nabla h(a)=(h'_x(a),h'_y(a))=0$.
\textbf{Интеграл по длине дуги:} $\int_{|z-a|=r}\varphi(z)\,|dz|=\int_0^{2\pi}\varphi(a+re^{i\vartheta})\,r\,d\vartheta$.
\end{defs}

\subsection{Определение и первые примеры}

\begin{thmbox}{Гармоничность $\Re f$ и $\Im f$}
Если $D$ --- область и $f\in H(D)$, то $u=\Re f$ и $v=\Im f$ принадлежат $\pIVHarm(D)$.
\end{thmbox}
\begin{proof}
Так как $f\in C^\infty(D)$, то $u,v\in C^\infty(D)$. Дифференцируя условия Коши--Римана $u'_x=v'_y$, $u'_y=-v'_x$
по $x$ и по $y$ соответственно, получаем $u''_{xx}=v''_{yx}$, $u''_{yy}=-v''_{xy}$. По теореме Шварца $v''_{xy}=v''_{yx}$,
поэтому $\Delta u=0$. Для $v$: $v=\Re(-if)$, $-if\in H(D)$, значит $\Delta v=0$ по уже доказанному.
\end{proof}

\begin{thmbox}{Лемма 21.1 (Лапласиан через производные Виртингера)}
Для любой $h\in C^2(D)$ (вещественной или комплексной) выполнено $4\,\partial_z\partial_{\bar z}h=4\,\partial_{\bar z}\partial_z h=\Delta h$.
\end{thmbox}
\begin{proof}
$4\partial_{\bar z}\partial_z h=(\partial_x+i\partial_y)(h'_x-ih'_y)=h''_{xx}-ih''_{yx}+ih''_{xy}+h''_{yy}=\Delta h$
(мы использовали $h''_{xy}=h''_{yx}$). Для другого порядка вычисление аналогично:
$(\partial_x-i\partial_y)(h'_x+ih'_y)=h''_{xx}+ih''_{yx}-ih''_{xy}+h''_{yy}=\Delta h$.
\end{proof}

\begin{remarkbox}{Пример: требование $h\in C^2(D)$ в определении существенно}
Положим
\[
h_*(x,y)=\frac{xy}{(x^2+y^2)^2}=\Im\Bigl(-\frac1{2z^2}\Bigr)\quad\text{при }(x,y)\ne(0,0),\qquad h_*(0,0)=0 .
\]
\textit{Проверка тождества.} При $z\ne0$: $\dfrac1{z^2}=\dfrac{\bar z^2}{|z|^4}=\dfrac{(x^2-y^2)-2ixy}{(x^2+y^2)^2}$, поэтому
$-\dfrac1{2z^2}=\dfrac{-(x^2-y^2)+2ixy}{2(x^2+y^2)^2}$ и $\Im\bigl(-\frac1{2z^2}\bigr)=\dfrac{xy}{(x^2+y^2)^2}$.
Так как $-1/(2z^2)\in H(\CC\setminus\{0\})$, то $h_*\in\pIVHarm(\CC\setminus\{0\})$.

\textit{Лапласиан в нуле.} Имеем $h_*(t,0)=0$ при всех $t\in\RR$. Значит, $\partial_x h_*(t,0)$ (производная функции
$t\mapsto h_*(t,0)\equiv0$) равна $0$ при всех $t$, и потому $\partial^2_{xx}h_*(0,0)=\frac{d}{dt}\bigl[\partial_xh_*(t,0)\bigr]_{t=0}=0$.
Точно так же $h_*(0,t)\equiv0$ даёт $\partial^2_{yy}h_*(0,0)=0$. Итак, вторые частные производные $h''_{xx},h''_{yy}$ существуют
\emph{во всех} точках $\RR^2$ и $\Delta h_*=0$ всюду в $\CC$.

\textit{Но $h_*\notin\pIVHarm(\{0\})$:} $h_*(t,t)=\dfrac{t^2}{4t^4}=\dfrac1{4t^2}\to+\infty$ при $t\to0$, т.\,е. $h_*$ даже разрывна в нуле.
Если бы в определении требовалось лишь существование вторых производных и $\Delta h=0$, то в класс <<гармонических>>
попали бы разрывные и неограниченные функции, и ни одно из доказываемых ниже свойств (гладкость, принцип максимума и т.\,д.) не было бы верно.
\end{remarkbox}

\begin{remarkbox}{Замечание (опечатка на слайде 339)}
На слайде написано <<$h_*\notin\pIVHarm(0)$, так как $h$ разрывна>>; имеется в виду $h_*\notin\pIVHarm(\{0\})$, так как $h_*$ разрывна в точке $0$.
\end{remarkbox}

\subsection{Сопряжённая гармоническая функция; гладкость}

\begin{thmbox}{Предложение 14.1}
Пусть $D$ --- односвязная область в $\CC$, $h\in\pIVHarm(D)$. Тогда существует $f\in H(D)$ с условием $\Re f=h$.
Функция $f$ определяется единственным образом с точностью до аддитивной чисто мнимой постоянной.
\end{thmbox}
\begin{proof}
\emph{Шаг 1: сопряжённый градиент голоморфен.} Положим $g=2h'_z=h'_x-ih'_y$. Так как $h\in C^2(D)$, то $g\in C^1(D)$.
По лемме 21.1
\[
2g'_{\bar z}=(\partial_x+i\partial_y)(h'_x-ih'_y)=\Delta h=0\quad\text{в }D,
\]
так что $g$ (вещественно дифференцируемая, класса $C^1$) удовлетворяет условию Коши--Римана $g'_{\bar z}=0$, т.\,е. $g\in H(D)$.

\emph{Шаг 2: первообразная.} Так как $D$ односвязна, существует $F\in H(D)$ с $F'=g$. Пусть $F=u+iv$. Тогда
\[
h'_x-ih'_y=g=F'=F'_x=u'_x+iv'_x .
\]
Сравнивая вещественные и мнимые части: $h'_x=u'_x$, $h'_y=-v'_x=u'_y$ (последнее --- условие Коши--Римана для $F$).
Итак, $\nabla(u-h)=0$ в $D$.

\emph{Шаг 3.} Функция с нулевым градиентом в области постоянна: если $\gamma\colon[0,1]\to D$ --- ломаная, соединяющая $z_1$ и $z_2$,
то на каждом звене $\frac{d}{dt}(u-h)(\gamma(t))=\langle\nabla(u-h),\gamma'(t)\rangle=0$; а область линейно связна ломаными.
Следовательно, $u=h+C_0$, $C_0\in\RR$. Функция $f=F-C_0\in H(D)$ удовлетворяет $\Re f=u-C_0=h$.

\emph{Единственность.} Если $f_1,f_2\in H(D)$ и $\Re f_1=\Re f_2=h$, то $\varphi=f_1-f_2\in H(D)$, $\Re\varphi\equiv0$.
Так как $\varphi'=\varphi'_x=(\Re\varphi)'_x+i(\Im\varphi)'_x$ и по Коши--Риману $(\Im\varphi)'_x=-(\Re\varphi)'_y$, то
$\varphi'=(\Re\varphi)'_x-i(\Re\varphi)'_y=0$. Значит, $\varphi\equiv iC$, $C\in\RR$ (по шагу 3, применённому к $\Re\varphi$ и $\Im\varphi$).
Обратно, $f+iC$ при $C\in\RR$, очевидно, тоже подходит.
\end{proof}

\begin{remarkbox}{Сопряжённый градиент}
Функция $g=2h'_z=h'_x-ih'_y$ называется \emph{сопряжённым градиентом} $h$: вектор $(h'_x,-h'_y)$ на плоскости сопряжён
вектору градиента $(h'_x,h'_y)$. Если $h=\Re f$, $f\in H(D)$, то $g=f'$: действительно, $f'=u'_x-iu'_y$ (см. определения).
Функция $v=\Im f$ называется \emph{гармонически сопряжённой} к $h$; она определена с точностью до постоянной.
\end{remarkbox}

\begin{thmbox}{Следствие 21.2 (локальное представление и гладкость)}
Пусть $D$ --- произвольная (не обязательно односвязная) область, $h\in\pIVHarm(D)$. Тогда
\begin{enumerate}
\item[(а)] для любого круга $B(a,\rho)\subset D$ найдётся $f\in H(B(a,\rho))$ с $\Re f=h$ в $B(a,\rho)$;
\item[(б)] $h\in C^\infty(D)$, т.\,е. $\pIVHarm(D)\subset C^\infty(D)$;
\item[(в)] $h$ вещественно аналитична: для любого круга $B(a,\rho)\subset D$, $a=x_0+iy_0$, функция $h$ представляется в
ромбе $|x-x_0|+|y-y_0|<\rho$ абсолютно сходящимся степенным рядом $\sum_{j,k\ge0}\alpha_{jk}(x-x_0)^j(y-y_0)^k$.
\end{enumerate}
\end{thmbox}
\begin{proof}
(а) Сужение $h$ на круг гармонично, а круг односвязен --- применяем предложение 14.1.

(б) Как в шаге 1 доказательства предложения 14.1 (он нигде не использовал односвязность), $g=h'_x-ih'_y\in H(D)\subset C^\infty(D)$.
Значит, $h'_x=\Re g$ и $h'_y=-\Im g$ лежат в $C^\infty(D)$, а тогда все частные производные $h$ любого порядка $\ge1$
(являющиеся производными от $h'_x$ или $h'_y$) существуют и непрерывны.

(в) По (а) $h=\Re f$ в $B(a,\rho)$; ряд Тейлора $f(z)=\sum_n c_n(z-a)^n$ сходится абсолютно в $B(a,\rho)$. Обозначим $s=x-x_0$, $t=y-y_0$.
Тогда $\Re\bigl(c_n(s+it)^n\bigr)=\Re\sum_{k=0}^n c_n\binom nk s^k(it)^{n-k}$ --- однородный многочлен степени $n$ от $s,t$, и сумма модулей его
мономов не превосходит $|c_n|\sum_k\binom nk|s|^k|t|^{n-k}=|c_n|(|s|+|t|)^n$. При $|s|+|t|<\rho$ точка $z=a+(|s|+|t|)$ лежит в круге сходимости, так что
$\sum_n|c_n|(|s|+|t|)^n<\infty$. Значит, двойной ряд из всех мономов сходится абсолютно, и его можно переупорядочить в
степенной ряд по $s,t$, сумма которого равна $\Re f=h$.
\end{proof}

\begin{thmbox}{Упражнение со слайда 341}
$h(z)=\ln|z|\in\pIVHarm(\CC\setminus\{0\})$, но не существует $f\in H(\CC\setminus\{0\})$ с $\Re f=h$.
\end{thmbox}
\begin{proof}
\emph{Гармоничность.} $h=\frac12\ln(x^2+y^2)\in C^\infty(\CC\setminus\{0\})$,
\[
h'_x=\frac{x}{x^2+y^2},\quad h''_{xx}=\frac{y^2-x^2}{(x^2+y^2)^2},\qquad h'_y=\frac{y}{x^2+y^2},\quad h''_{yy}=\frac{x^2-y^2}{(x^2+y^2)^2},
\]
откуда $\Delta h=0$.

\emph{Отсутствие голоморфной $f$.} Пусть $f\in H(\CC\setminus\{0\})$, $\Re f=h$. Тогда (см. замечание о сопряжённом градиенте)
\[
f'=h'_x-ih'_y=\frac{x-iy}{x^2+y^2}=\frac{\bar z}{|z|^2}=\frac1z .
\]
Значит, $1/z$ имеет первообразную $f$ в $\CC\setminus\{0\}$, и интеграл по замкнутой окружности $\oint_{|z|=1}\frac{dz}{z}$ равен нулю
(интеграл от производной по замкнутому пути равен приращению первообразной, т.\,е. $0$). Но прямое вычисление
даёт $\oint_{|z|=1}\frac{dz}z=\int_0^{2\pi}\frac{ie^{i\vartheta}}{e^{i\vartheta}}d\vartheta=2\pi i\ne0$. Противоречие.
\end{proof}

\begin{remarkbox}{Смысл упражнения}
Односвязность в предложении 14.1 существенна. Локально $\ln|z|=\Re(\text{ветвь }\log z)$ (следствие 21.2(а)), но
гармонически сопряжённая функция $\arg z$ не имеет однозначной непрерывной ветви в $\CC\setminus\{0\}$.
\end{remarkbox}

\subsection{Внутренняя теорема единственности}

\begin{thmbox}{Внутренняя теорема единственности для гармонических функций}
Пусть $D$ --- область в $\CC$, $h\in\pIVHarm(D)$, $K_h=\{a\in D:\nabla h(a)=0\}$ --- множество критических точек $h$.
Если $K_h$ имеет в $D$ хотя бы одну предельную точку, то $h=\mathrm{Const}$ в $D$.
\end{thmbox}
\begin{proof}
Пусть $g=2h'_z=h'_x-ih'_y$ --- сопряжённый градиент; по шагу 1 доказательства предложения 14.1 $g\in H(D)$ (односвязность там не нужна).
Для $a\in K_h$ имеем $g(a)=h'_x(a)-ih'_y(a)=0$, т.\,е. $K_h$ содержится в множестве нулей $g$. Если $K_h$ имеет предельную точку в $D$,
то по теореме единственности для голоморфных функций $g\equiv0$ в $D$. Тогда $h'_x=\Re g\equiv0$ и $h'_y=-\Im g\equiv0$, и, как в шаге 3
доказательства предложения 14.1, $h\equiv\mathrm{const}$ в $D$.
\end{proof}

\begin{remarkbox}{Следствия}
(1) Если $h\in\pIVHarm(D)$ постоянна на каком-нибудь непустом открытом множестве $U\subset D$, то $h$ постоянна в $D$
(все точки $U$ --- критические). (2) Если $h_1,h_2\in\pIVHarm(D)$ совпадают на открытом $U\ne\varnothing$, то $h_1\equiv h_2$.
\emph{Предостережение:} в отличие от голоморфных функций, совпадение на множестве с предельной точкой не влечёт
тождества: $h=y$ равна нулю на всей оси $\RR$, но $h\not\equiv0$. Поэтому теорема сформулирована для множества
\emph{критических} точек, а не нулей.
\end{remarkbox}

\subsection{Теорема о среднем}

\begin{thmbox}{Предложение 14.2 (теорема о среднем)}
Пусть $a\in\CC$, $0<R<+\infty$, $h\in\pIVHarm(B(a,R))\cap C(\clos{B(a,R)})$. Тогда при любом $r\in(0,R]$
\[
h(a)=\frac1{2\pi r}\int_{|z-a|=r}h(z)\,|dz|=\frac1{2\pi}\int_0^{2\pi}h(a+re^{i\vartheta})\,d\vartheta .
\]
\end{thmbox}
\begin{proof}
Равенство двух интегралов --- это определение интеграла по длине дуги ($|dz|=r\,d\vartheta$).

\emph{Случай $r\in(0,R)$.} Круг $B(a,R)$ односвязен, поэтому по предложению 14.1 найдётся $f\in H(B(a,R))$ с $\Re f=h$.
Так как $\clos{B(a,r)}\subset B(a,R)$, по теореме о среднем для голоморфных функций
$f(a)=\frac1{2\pi}\int_0^{2\pi}f(a+re^{i\vartheta})\,d\vartheta$. Беря вещественные части обеих частей (вещественная часть
интеграла равна интегралу от вещественной части), получаем требуемое.

\emph{Случай $r=R$.} Функция $(r,\vartheta)\mapsto h(a+re^{i\vartheta})$ непрерывна на компакте $[R/2,R]\times[0,2\pi]$ и потому
равномерно непрерывна; значит, $h(a+re^{i\vartheta})\to h(a+Re^{i\vartheta})$ при $r\to R-$ равномерно по $\vartheta$, и
\[
h(a)=\frac1{2\pi}\int_0^{2\pi}h(a+re^{i\vartheta})\,d\vartheta\xrightarrow[r\to R-]{}\frac1{2\pi}\int_0^{2\pi}h(a+Re^{i\vartheta})\,d\vartheta .\qedhere
\]
\end{proof}

\begin{remarkbox}{Замечание (опечатка на слайде 342)}
В доказательстве написано <<найдём $f\in H(D)$>>; следует читать $f\in H(B(a,R))$ --- именно круг $B(a,R)$ играет роль
односвязной области в предложении 14.1.
\end{remarkbox}

\begin{remarkbox}{Теорема о среднем по площади (полезное дополнение)}
В условиях предложения 14.2: $h(a)=\dfrac1{\pi R^2}\iint_{B(a,R)}h\,dx\,dy$. Действительно, в полярных координатах
$\iint_{B(a,R)}h\,dx\,dy=\int_0^R\Bigl(\int_0^{2\pi}h(a+re^{i\vartheta})d\vartheta\Bigr)r\,dr=\int_0^R2\pi h(a)\,r\,dr=\pi R^2h(a)$.
\end{remarkbox}

\subsection{Принцип максимума и минимума}

\begin{thmbox}{Теорема 14.1 (принцип минимума и максимума)}
Пусть $D$ --- ограниченная область в $\CC$, $h\in\pIVHarm(D)\cap C(\clos D)$. Тогда для любого $z_0\in D$
\[
\min_{z\in\partial D}h(z)\le h(z_0)\le\max_{z\in\partial D}h(z).
\]
Если найдётся $z_0\in D$, для которого любое из этих неравенств обращается в равенство, то $h=\mathrm{Const}$ в $\clos D$.
\end{thmbox}
\begin{proof}
Максимум и минимум по $\partial D$ существуют: $\partial D$ --- непустой компакт (см. определения), $h$ непрерывна.
Положим $M=\max_{\clos D}h$ (существует, так как $\clos D$ компактен).

\emph{Утверждение.} Если существует $z_0\in D$ с $h(z_0)\ge\max_{\partial D}h$, то $h\equiv\mathrm{const}$ в $\clos D$.

Пусть $E=\{z\in D: h(z)=M\}$. Покажем, что $E\ne\varnothing$. Максимум $M$ достигается в некоторой точке $z_1\in\clos D$.
Если $z_1\in D$, то $z_1\in E$. Если же $z_1\in\partial D$, то $M=h(z_1)\le\max_{\partial D}h\le h(z_0)\le M$, откуда $h(z_0)=M$ и $z_0\in E$.

\begin{detail}{<<нетрудно показать, что $h\equiv h(z_0)$ в круге>> (слайд 343)}
Возьмём $z_1\in E$ и $R=\dist(z_1,\partial D)>0$ (расстояние от точки открытого $D$ до замкнутого $\partial D$, не содержащего $z_1$).
Тогда $B=B(z_1,R)\subset D$: круг $B$ связен, не пересекает $\partial D$, поэтому $B=(B\cap D)\sqcup(B\setminus\clos D)$ ---
разбиение на два открытых множества, первое из которых содержит $z_1$; по связности второе пусто.

Для любого $r\in(0,R)$ имеем $\clos{B(z_1,r)}\subset B\subset D$, и $h\in\pIVHarm(B(z_1,r))\cap C(\clos{B(z_1,r)})$, так что по
предложению 14.2
\[
0=M-h(z_1)=\frac1{2\pi}\int_0^{2\pi}\bigl(M-h(z_1+re^{i\vartheta})\bigr)\,d\vartheta .
\]
Подынтегральная функция непрерывна и неотрицательна ($M$ --- максимум). Интеграл от непрерывной неотрицательной функции
равен нулю, только если она тождественно равна нулю (если бы она была $\ge c>0$ на интервале длины $\ell>0$, интеграл был бы $\ge c\ell/(2\pi)>0$).
Значит, $h=M$ на каждой окружности $|z-z_1|=r$, $0<r<R$, т.\,е. $h\equiv M$ в $B$.
\end{detail}

Итак, $h\equiv M$ в круге $B$, значит $\nabla h=0$ в $B$: все точки $B$ --- критические, и любая из них является предельной точкой
множества $K_h$. По внутренней теореме единственности $h\equiv M$ в $D$, а по непрерывности и в $\clos D=$ замыкании $D$.
Утверждение доказано.

\emph{Вывод теоремы.} Если $h\not\equiv\mathrm{const}$, то по утверждению $h(z_0)<\max_{\partial D}h$ для всех $z_0\in D$.
Если $h\equiv\mathrm{const}$, все неравенства --- равенства. В обоих случаях $h(z_0)\le\max_{\partial D}h$, а равенство
в какой-нибудь точке $z_0\in D$ влечёт (по утверждению) постоянство $h$. Для минимума применяем доказанное к $-h\in\pIVHarm(D)\cap C(\clos D)$:
$\max_{\partial D}(-h)=-\min_{\partial D}h$.
\end{proof}

\begin{thmbox}{Следствие 21.3 (сильный принцип максимума, локальная форма)}
Если $D$ --- произвольная (не обязательно ограниченная) область, $h\in\pIVHarm(D)$ и $h$ имеет в точке $z_1\in D$ локальный
максимум или минимум, то $h\equiv\mathrm{const}$ в $D$.
\end{thmbox}
\begin{proof}
Пусть $h(z)\le h(z_1)$ в $B(z_1,R)\subset D$. Тот же аргумент со средними (деталь выше, с $M=h(z_1)$) даёт $h\equiv h(z_1)$ в $B(z_1,R)$,
а внутренняя теорема единственности --- $h\equiv\mathrm{const}$ в $D$. Для минимума --- заменить $h$ на $-h$.
\end{proof}

\begin{remarkbox}{Ограниченность области существенна}
В верхней полуплоскости $\Pi_+=\{\Im z>0\}$ функция $h=y$ гармонична, непрерывна в $\clos{\Pi_+}$, равна $0$ на $\partial\Pi_+=\RR$,
но $h>0$ в $\Pi_+$. Принцип максимума (в форме теоремы 14.1) для неограниченных областей неверен.
\end{remarkbox}

\subsection{Граничная теорема единственности}

\begin{thmbox}{Предложение 14.3 (граничная теорема единственности)}
Пусть $D$ --- ограниченная область в $\CC$, $h_1,h_2\in\pIVHarm(D)\cap C(\clos D)$. Если $h_1=h_2$ на $\partial D$, то $h_1=h_2$ в $\clos D$.
\end{thmbox}
\begin{proof}
$h=h_1-h_2\in\pIVHarm(D)\cap C(\clos D)$ и $h=0$ на $\partial D$. По теореме 14.1 для любого $z_0\in D$:
$0=\min_{\partial D}h\le h(z_0)\le\max_{\partial D}h=0$.
\end{proof}

\begin{remarkbox}{Более сильная оценка}
Тем же рассуждением: $\|h_1-h_2\|_{\clos D}=\|h_1-h_2\|_{\partial D}$ для любых $h_1,h_2\in\pIVHarm(D)\cap C(\clos D)$
(применить теорему 14.1 к $\pm(h_1-h_2)$). Это --- ключевая оценка устойчивости решения задачи Дирихле (билеты 22, 24).
\end{remarkbox}

\begin{remarkbox}{Коротко: план ответа на экзамене}
\begin{enumerate}
\item Определение 14.1; $\Re f,\Im f$ гармоничны (Коши--Риман + Шварц); пример $h_*=xy/(x^2+y^2)^2$: $\Delta h_*=0$ всюду, но $h_*$ разрывна.
\item $\Delta=4\partial_z\partial_{\bar z}$; сопряжённый градиент $g=h'_x-ih'_y$ голоморфен.
\item Предл. 14.1: в односвязной области $h=\Re f$ (первообразная от $g$); единственность с точностью до $iC$.
\item Следствия: $\pIVHarm\subset C^\infty$, вещественная аналитичность; $\ln|z|$ --- контрпример без односвязности ($f'=1/z$).
\item Внутренняя теорема единственности: критические точки с предельной точкой $\Rightarrow$ $h\equiv\mathrm{const}$.
\item Предл. 14.2: среднее по окружности --- вещественная часть формулы среднего для $f$; при $r=R$ --- предельный переход.
\item Теорема 14.1: в точке максимума среднее неотрицательной непрерывной $M-h$ равно $0$ $\Rightarrow$ $h\equiv M$ в круге
$\Rightarrow$ (внутр. единственность) в $D$.
\item Предл. 14.3: граничная единственность; ограниченность области существенна ($h=y$ в полуплоскости).
\end{enumerate}
\end{remarkbox}
\ticket{22}{Постановка задачи Дирихле для гармонических функций. Сохранение гармоничности при голоморфной замене переменных и метод конформного отображения для решения задачи Дирихле. Формулировка теоремы Лебега. Гармоничность непрерывной функции, удовлетворяющей в произвольном круге в данной области теореме о среднем}
\label{t22:start}

\begin{defs}
\textbf{Гармонические функции} (см. билет 21): $h\in\pIVHarm(D)$, если $h\in C^2(D)$ вещественна и $\Delta h=h''_{xx}+h''_{yy}=0$ в $D$.
Локально $h=\Re f$, $f$ голоморфна (следствие 21.2); $\pIVHarm(D)\subset C^\infty(D)$.

\textbf{Задача Дирихле.} Пусть $D$ --- ограниченная область в $\CC$, $h_0\in C(\partial D)$ --- вещественнозначная функция.
Требуется найти $h\in\pIVHarm(D)\cap C(\clos D)$ такую, что $h|_{\partial D}=h_0$. Функция $h_0$ называется \emph{граничными данными}.
Задача \emph{разрешима} для данной $h_0$, если такая $h$ существует.

\textbf{Жорданова кривая} --- гомеоморфный образ окружности в $\CC$. \emph{Жорданова область} --- ограниченная область, граница
которой --- жорданова кривая (по теореме Жордана это внутренность жордановой кривой). Жорданова область односвязна.

\textbf{Конформное отображение.} $\pIVConf(D,G)$ --- множество биголоморфных отображений $D$ на $G$ (голоморфных биекций; обратное
тогда автоматически голоморфно). \emph{Теорема Римана:} всякая односвязная область $D\subsetneq\CC$ конформно эквивалентна
единичному кругу $\DD$. \emph{Теорема Каратеодори} (билет 20): если $D,G$ --- жордановы области и $\varphi\in\pIVConf(D,G)$, то $\varphi$
продолжается до гомеоморфизма $\clos D\to\clos G$, переводящего $\partial D$ на $\partial G$.

\textbf{Цепное правило для производных Виртингера.} Если $F$ дифференцируема (вещественно) в точке $w=\varphi(z)$, а $\varphi$ --- в точке $z$, то
\[
(F\circ\varphi)'_z=F'_w(\varphi)\,\varphi'_z+F'_{\bar w}(\varphi)\,(\clos\varphi)'_z,\qquad
(F\circ\varphi)'_{\bar z}=F'_w(\varphi)\,\varphi'_{\bar z}+F'_{\bar w}(\varphi)\,(\clos\varphi)'_{\bar z}.
\]
(Вывод: из $dx=\frac{dz+d\bar z}2$, $dy=\frac{dz-d\bar z}{2i}$ имеем $dF=F'_xdx+F'_ydy=F'_wdw+F'_{\bar w}d\bar w$; подставляем
$dw=\varphi'_zdz+\varphi'_{\bar z}d\bar z$, $d\bar w=(\clos\varphi)'_zdz+(\clos\varphi)'_{\bar z}d\bar z$ и собираем коэффициенты при $dz$ и $d\bar z$.)
Для голоморфной $\varphi$: $\varphi'_{\bar z}=0$, $\varphi'_z=\varphi'$, $(\clos\varphi)'_z=\overline{\varphi'_{\bar z}}=0$, $(\clos\varphi)'_{\bar z}=\overline{\varphi'}$.

\textbf{Свойство среднего.} Непрерывная функция $h\colon D\to\RR$ обладает \emph{свойством среднего в $D$}, если для любого
замкнутого круга $\clos{B(a,r)}\subset D$
\[
h(a)=\frac1{2\pi}\int_0^{2\pi}h(a+re^{i\vartheta})\,d\vartheta .
\]
\emph{Локальное} свойство среднего: для каждой $a\in D$ найдётся $\rho(a)>0$ с $\clos{B(a,\rho(a))}\subset D$, такое что равенство
выполнено при всех $r\in(0,\rho(a))$. По предложению 14.2 всякая гармоническая функция обладает свойством среднего.

\textbf{Интеграл Пуассона} (билет 24): для $h_0\in C(\partial B(a,R))$
\[
\mathcal P[h_0](z)=\frac1{2\pi R}\int_{|\zeta-a|=R}\frac{R^2-|z-a|^2}{|\zeta-z|^2}\,h_0(\zeta)\,|d\zeta|,\qquad |z-a|<R .
\]
\end{defs}

\subsection{Постановка задачи Дирихле и единственность решения}

Постановка дана в рамке определений. Слово <<гармоническая>> относится к внутренности $D$, а на границе требуется лишь
непрерывное примыкание к заданным значениям.

\begin{thmbox}{Единственность решения задачи Дирихле}
Если задача Дирихле в ограниченной области $D$ разрешима для данной $h_0\in C(\partial D)$, то её решение единственно.
Более того, если $h_1,h_2$ --- решения для данных $h_{0,1},h_{0,2}$, то $\|h_1-h_2\|_{\clos D}=\|h_{0,1}-h_{0,2}\|_{\partial D}$.
\end{thmbox}
\begin{proof}
$h=h_1-h_2\in\pIVHarm(D)\cap C(\clos D)$; по принципу максимума и минимума (теорема 14.1, билет 21) для $z\in D$
\[
-\|h_{0,1}-h_{0,2}\|_{\partial D}\le\min_{\partial D}h\le h(z)\le\max_{\partial D}h\le\|h_{0,1}-h_{0,2}\|_{\partial D},
\]
а на $\partial D$ имеем $h=h_{0,1}-h_{0,2}$. При $h_{0,1}=h_{0,2}$ получаем $h_1=h_2$ (предложение 14.3).
\end{proof}

\begin{remarkbox}{Замечание}
Итак, задача Дирихле <<корректна>>: решение единственно и непрерывно (в равномерной норме) зависит от граничных данных.
Главный вопрос --- \emph{существование}. Его мы решаем для круга (билет 24), а затем переносим на другие области
конформными отображениями.
\end{remarkbox}

\subsection{Сохранение гармоничности при голоморфной замене переменных}

\begin{thmbox}{Предложение 14.4}
Пусть $D$ и $G$ --- области в $\CC$, $\varphi\in H(D)$ и $\varphi(D)\subset G$. Пусть $H\in\pIVHarm(G)$. Тогда $h=H\circ\varphi\in\pIVHarm(D)$.
\end{thmbox}

\begin{proof}[Первое доказательство]
(Через локальное представление $H=\Re F$; слайд 344.) Фиксируем произвольную точку $z_0\in D$, $w_0=\varphi(z_0)\in G$. Так как $G$ открыто, найдётся $\eps>0$ с $B(w_0,\eps)\subset G$;
так как $\varphi$ непрерывна и $D$ открыто, найдётся $\delta>0$ с $B=B(z_0,\delta)\subset D$ и $\varphi(B)\subset B(w_0,\eps)$.
Круг $B(w_0,\eps)$ односвязен, поэтому по предложению 14.1 существует $F\in H(B(w_0,\eps))$ с $\Re F=H$ в $B(w_0,\eps)$.
Тогда $f=F\circ\varphi\in H(B)$ как композиция голоморфных функций, и $h=H\circ\varphi=\Re(F\circ\varphi)=\Re f$ в $B$.
По билету 21 $\Re f\in\pIVHarm(B)$; в частности $h\in C^2(B)$ и $\Delta h=0$ в $B$. Так как $z_0\in D$ произвольна, $h\in C^2(D)$ и
$\Delta h=0$ всюду в $D$, т.\,е. $h\in\pIVHarm(D)$.
\end{proof}

\begin{thmbox}{Лемма 22.1 (закон преобразования лапласиана)}
Пусть $\varphi\in H(D)$, $\varphi(D)\subset G$, $H\in C^2(G)$ (вещественная или комплексная). Тогда $h=H\circ\varphi\in C^2(D)$ и
\[
\Delta(H\circ\varphi)(z)=|\varphi'(z)|^2\,(\Delta H)(\varphi(z)),\qquad z\in D .
\]
\end{thmbox}
\begin{proof}
$h\in C^2(D)$ как композиция $C^2$-отображений ($\varphi\in C^\infty$ как отображение $\RR^2\to\RR^2$). По цепному правилу
(см. определения), так как $(\clos\varphi)'_z=0$,
\[
h'_z=H'_w(\varphi)\,\varphi'+H'_{\bar w}(\varphi)\cdot0=H'_w(\varphi)\,\varphi' .
\]
Функция $H'_w$ класса $C^1(G)$, поэтому к $H'_w\circ\varphi$ снова применимо цепное правило; учитывая $\varphi'_{\bar z}=0$,
$(\clos\varphi)'_{\bar z}=\overline{\varphi'}$ и $\partial_{\bar z}\varphi'=0$ (так как $\varphi'\in H(D)$), получаем
\[
h''_{z\bar z}=\partial_{\bar z}\bigl[H'_w(\varphi)\bigr]\cdot\varphi'+H'_w(\varphi)\cdot\partial_{\bar z}\varphi'
=\bigl[H''_{ww}(\varphi)\cdot0+H''_{w\bar w}(\varphi)\,\overline{\varphi'}\bigr]\varphi'
=|\varphi'|^2\,H''_{w\bar w}(\varphi).
\]
Умножая на $4$ и используя лемму 21.1 ($\Delta=4\partial_{\bar z}\partial_z$, и то же в переменной $w$), получаем
$\Delta h=|\varphi'|^2\,(4H''_{w\bar w})\circ\varphi=|\varphi'|^2\,(\Delta H)\circ\varphi$.
\end{proof}

\begin{proof}[Второе доказательство]
(Прямое вычисление.) По лемме 22.1 $h=H\circ\varphi\in C^2(D)$ и $\Delta h=|\varphi'|^2(\Delta H)\circ\varphi=|\varphi'|^2\cdot0=0$ в $D$.
\end{proof}

\begin{remarkbox}{Следствия леммы 22.1}
(1) Если $\varphi\in\pIVConf(D,G)$ (тогда $\varphi'\ne0$), то $H\in\pIVHarm(G)\iff H\circ\varphi\in\pIVHarm(D)$ (для $\Leftarrow$ применить
предложение 14.4 к $\varphi^{-1}$). (2) Частные случаи: сдвиг $\varphi(z)=z+b$, гомотетия и поворот $\varphi(z)=\lambda z$ сохраняют
гармоничность. Например, если $H\in\pIVHarm(B(0,R))$, то $H(z-a)\in\pIVHarm(B(a,R))$.
(3) Без голоморфности утверждение неверно: $H(w)=\Re w$, $\varphi(z)=|z|^2$ (гладкое, но не голоморфное) дают $h=x^2+y^2$, $\Delta h=4$.
\end{remarkbox}

\subsection{Метод конформного отображения}

\begin{thmbox}{Предложение 22.2 (метод конформного отображения; слайды 344--345)}
Пусть $D$ и $G$ --- жордановы области, $\varphi\in\pIVConf(D,G)$, и $\varphi$ продолжена до гомеоморфизма $\clos D\to\clos G$
(теорема Каратеодори). Пусть $h_0\in C(\partial D)$ и $H_0=h_0\circ\varphi^{-1}\in C(\partial G)$. Если $H$ --- решение задачи Дирихле
в $G$ с данными $H_0$, то $h=H\circ\varphi$ --- решение задачи Дирихле в $D$ с данными $h_0$.
\end{thmbox}
\begin{proof}
Продолженное $\varphi$ --- гомеоморфизм $\clos D\to\clos G$, и так как $\varphi(D)=G$, то $\varphi(\partial D)=\partial G$; поэтому
$\varphi^{-1}|_{\partial G}\colon\partial G\to\partial D$ непрерывно и $H_0\in C(\partial G)$. Далее:
\begin{itemize}
\item $h=H\circ\varphi\in C(\clos D)$ как композиция непрерывных отображений $\clos D\to\clos G\to\RR$;
\item $h|_D\in\pIVHarm(D)$ по предложению 14.4 ($\varphi|_D\in H(D)$, $\varphi(D)=G$, $H\in\pIVHarm(G)$);
\item для $\zeta\in\partial D$: $h(\zeta)=H(\varphi(\zeta))=H_0(\varphi(\zeta))=h_0(\varphi^{-1}(\varphi(\zeta)))=h_0(\zeta)$. \qedhere
\end{itemize}
\end{proof}

\begin{thmbox}{Следствие 22.3 (задача Дирихле в жордановых областях)}
В любой жордановой области $D$ любая задача Дирихле разрешима. Явно: если $\varphi\in\pIVConf(D,\DD)$ (продолженное на $\clos D$), то
\[
h(z)=\frac1{2\pi}\int_{\TT}\frac{1-|\varphi(z)|^2}{|\zeta-\varphi(z)|^2}\,h_0\bigl(\varphi^{-1}(\zeta)\bigr)\,|d\zeta|,\qquad z\in D .
\]
\end{thmbox}
\begin{proof}
Жорданова область односвязна и отлична от $\CC$ (она ограничена), поэтому по теореме Римана существует $\varphi\in\pIVConf(D,\DD)$;
$\DD$ --- жорданова область, и по теореме Каратеодори $\varphi$ продолжается до гомеоморфизма $\clos D\to\clos\DD$. В круге $\DD$
задача Дирихле с любыми непрерывными данными $H_0$ разрешима, и решение даётся интегралом Пуассона
$H(w)=\frac1{2\pi}\int_\TT\frac{1-|w|^2}{|\zeta-w|^2}H_0(\zeta)|d\zeta|$ (предложения 14.8, 14.9, теорема 24.7 в билете 24).
Остаётся применить предложение 22.2 с $H_0=h_0\circ\varphi^{-1}$.
\end{proof}

\begin{remarkbox}{Пример 1: произвольный круг $B(a,R)$}
$\varphi(z)=(z-a)/R\in\pIVConf(B(a,R),\DD)$, $\varphi^{-1}(\omega)=a+R\omega$. Пусть $w=\varphi(z)$, $\omega=\varphi(\zeta)$, $|\zeta-a|=R$. Тогда
$1-|w|^2=\frac{R^2-|z-a|^2}{R^2}$, $|\omega-w|^2=\frac{|\zeta-z|^2}{R^2}$, $|d\omega|=\frac{|d\zeta|}R$, и формула следствия 22.3 даёт
\[
h(z)=\frac1{2\pi R}\int_{|\zeta-a|=R}\frac{R^2-|z-a|^2}{|\zeta-z|^2}\,h_0(\zeta)\,|d\zeta|=\mathcal P[h_0](z),
\]
т.\,е. формулу Пуассона для круга $B(a,R)$.
\end{remarkbox}

\begin{remarkbox}{Пример 2: верхняя полуплоскость (неограниченная область; граница рассматривается в $\Ch$)}
Пусть $\Pi_+=\{z=x+iy:y>0\}$, $h_0\in C(\RR)$ и существует конечный предел $L=\lim_{t\to\pm\infty}h_0(t)$ (один и тот же на $\pm\infty$),
т.\,е. $h_0$ непрерывна на окружности $\RR\cup\{\infty\}\subset\Ch$. \textbf{Утверждение:}
\[
h(x+iy)=\frac1\pi\int_{-\infty}^{+\infty}\frac{y\,h_0(t)}{(t-x)^2+y^2}\,dt
\]
гармонична в $\Pi_+$, непрерывна в $\clos{\Pi_+}$ (при $h:=h_0$ на $\RR$) и $h(z)\to L$ при $z\to\infty$, $z\in\clos{\Pi_+}$.

\textit{Доказательство.} Отображение Кэли $\varphi(z)=\frac{z-i}{z+i}$ --- дробно-линейное, $|\varphi(z)|<1\iff|z-i|<|z+i|\iff y>0$;
значит, $\varphi\in\pIVConf(\Pi_+,\DD)$, а как гомеоморфизм $\Ch\to\Ch$ оно переводит $\clos{\Pi_+}\cup\{\infty\}$ на $\clos\DD$,
$\RR$ --- на $\TT\setminus\{1\}$, $\infty\mapsto1$. Положим $H_0(\zeta)=h_0(\varphi^{-1}(\zeta))$ при $\zeta\in\TT\setminus\{1\}$, $H_0(1)=L$;
по условию $H_0\in C(\TT)$. Пусть $H$ --- решение задачи Дирихле в $\DD$ (интеграл Пуассона). Тогда $h=H\circ\varphi$ гармонична в $\Pi_+$
(предложение 14.4), непрерывна на $\clos{\Pi_+}$, равна $h_0$ на $\RR$ и $h(z)=H(\varphi(z))\to H(1)=L$ при $z\to\infty$.

\textit{Вычисление ядра.} Параметризуем $\TT\setminus\{1\}$ точками $\zeta=\varphi(t)$, $t\in\RR$; $\varphi'(z)=\frac{2i}{(z+i)^2}$, так что
$|d\zeta|=|\varphi'(t)|dt=\frac{2\,dt}{1+t^2}$. Для $w=\varphi(z)$:
\[
1-|w|^2=\frac{|z+i|^2-|z-i|^2}{|z+i|^2}=\frac{4y}{|z+i|^2},\qquad
\zeta-w=\frac{(t-i)(z+i)-(z-i)(t+i)}{(t+i)(z+i)}=\frac{2i(t-z)}{(t+i)(z+i)} .
\]
Поэтому
\[
\frac{1-|w|^2}{|\zeta-w|^2}\,|d\zeta|=\frac{4y}{|z+i|^2}\cdot\frac{(1+t^2)|z+i|^2}{4|t-z|^2}\cdot\frac{2\,dt}{1+t^2}=\frac{2y\,dt}{(t-x)^2+y^2},
\]
и $h(z)=\frac1{2\pi}\int_\TT\frac{1-|w|^2}{|\zeta-w|^2}H_0(\zeta)|d\zeta|=\frac1\pi\int_\RR\frac{y\,h_0(t)\,dt}{(t-x)^2+y^2}$ (точка $\zeta=1$ имеет
нулевую длину). Проверка: при $h_0\equiv1$ получаем $\frac1\pi\int_\RR\frac{y\,dt}{(t-x)^2+y^2}=\frac1\pi\cdot\pi=1$.

\textit{Единственность} в классе функций, гармоничных в $\Pi_+$, непрерывных на $\clos{\Pi_+}$ и имеющих предел $L$ на $\infty$:
для двух таких $h_1,h_2$ функции $h_j\circ\varphi^{-1}$ (доопределённые значением $L$ в точке $1$) непрерывны на $\clos\DD$, гармоничны в $\DD$
и совпадают на $\TT$, значит совпадают (предложение 14.3). Без условия на бесконечности единственности нет: $h=y$ и $h=0$ имеют одинаковые
граничные значения на $\RR$.
\end{remarkbox}

\begin{remarkbox}{Пример 3: полукруг (ограниченная жорданова область)}
$D=\{|z|<1,\ \Im z>0\}$. Отображение $\psi_1(z)=\frac{1+z}{1-z}$ имеет вещественные коэффициенты и
$\Re\psi_1=\frac{1-|z|^2}{|1-z|^2}$, $\Im\psi_1=\frac{2\Im z}{|1-z|^2}$; поэтому оно переводит $D=\DD\cap\Pi_+$ биективно на первый
квадрант $Q=\{\Re w>0,\Im w>0\}$. Далее $w\mapsto w^2$ --- биекция $Q$ на $\Pi_+$ (аргумент удваивается: $(0,\frac\pi2)\to(0,\pi)$), и
$z\mapsto\frac{z-i}{z+i}$ переводит $\Pi_+$ на $\DD$. Итак,
$\varphi(z)=\dfrac{\psi_1(z)^2-i}{\psi_1(z)^2+i}\in\pIVConf(D,\DD)$, и задача Дирихле в полукруге решается формулой следствия 22.3.
Аналогично в секторе $\{0<\arg z<\pi/\alpha,\ |z|<1\}$ используется $z\mapsto z^{\alpha}$, сводящее его к полукругу.
\end{remarkbox}

\subsection{Теорема Лебега}

\begin{thmbox}{Теорема 14.2 (теорема Лебега)}
Пусть $D$ --- произвольная ограниченная односвязная область в $\CC$, а $h_0\in C(\partial D)$ --- произвольные граничные данные.
Тогда соответствующая задача Дирихле всегда разрешима, т.\,е. найдётся $h\in\pIVHarm(D)\cap C(\clos D)$ такая, что $h|_{\partial D}=h_0$.
\end{thmbox}

\begin{remarkbox}{О теореме Лебега}
\begin{itemize}
\item Доказательство в общем случае далеко выходит за рамки курса (оно использует, например, метод Перрона и понятие
\emph{барьера}: точка $\zeta_0\in\partial D$ регулярна, если существует функция, гармоническая в $D\cap B(\zeta_0,r)$, непрерывная
вплоть до границы, положительная всюду кроме $\zeta_0$ и равная $0$ в $\zeta_0$; у односвязной области все граничные точки регулярны,
поскольку каждая лежит на невырожденном связном куске дополнения).
\item \emph{Для жордановых областей} теорема Лебега доказана в курсе: это следствие 22.3 (метод конформного отображения
+ разрешимость в круге, билет 24). Граница произвольной односвязной области может быть устроена очень сложно (не быть кривой),
тогда теорема Каратеодори неприменима.
\item Решение единственно (принцип максимума), так что теорема утверждает существование.
\item \emph{О чём теорема.} Разрешимость задачи Дирихле зависит только от устройства границы: препятствием служат <<слишком тонкие>> куски дополнения, например изолированные точки границы (пример 22.4 ниже). У ограниченной односвязной области дополнение $\Ch\setminus D$ связно, и каждая граничная точка принадлежит связному множеству $\Ch\setminus D$, содержащему более одной точки; теорема Лебега утверждает, что этого достаточно. (В пространстве $\RR^3$ ситуация хуже: известный пример Лебега --- <<шип Лебега>> --- даёт область, гомеоморфную шару, в которой задача Дирихле разрешима не для всех непрерывных данных; в плоском случае такого не бывает.)
\end{itemize}
\end{remarkbox}

\begin{thmbox}{Пример 22.4 (без односвязности теорема неверна)}
Пусть $D=\DD\setminus\{0\}$ (ограниченная, но не односвязная область), $\partial D=\TT\cup\{0\}$, и $h_0=0$ на $\TT$, $h_0(0)=1$.
Тогда задача Дирихле с данными $h_0\in C(\partial D)$ неразрешима.
\end{thmbox}
\begin{proof}
Пусть $h\in\pIVHarm(D)\cap C(\clos\DD)$ --- решение; $M=\max_{\clos\DD}|h|$. Зафиксируем $\eps>0$ и при $0<r<1$ рассмотрим кольцо
$A_r=\{r<|z|<1\}$ и функцию $w_\eps=h+\eps\ln|z|\in\pIVHarm(A_r)\cap C(\clos{A_r})$ (сумма гармонических; $\ln|z|$ гармоничен, билет 21).
На $|z|=1$: $w_\eps=0$. На $|z|=r$: $w_\eps\le M+\eps\ln r<0$, если $r<e^{-M/\eps}$. По принципу максимума (теорема 14.1 в ограниченной
области $A_r$) $w_\eps\le0$ в $A_r$.

Возьмём любое $z$, $0<|z|<1$, и $r<\min(|z|,e^{-M/\eps})$: тогда $z\in A_r$ и $h(z)\le-\eps\ln|z|$. Устремляя $\eps\to0$ (при
фиксированном $z$), получаем $h(z)\le0$. Применяя то же к $-h$ (решение задачи с данными $-h_0$, для которых рассуждение то же: на $\TT$ нуль),
получаем $h(z)\ge0$. Итак, $h\equiv0$ в $D$, и по непрерывности $h(0)=0\ne1=h_0(0)$ --- противоречие.
\end{proof}

\subsection{Гармоничность непрерывной функции со свойством среднего}

Начнём с принципа максимума, в котором вместо гармоничности используется только (локальное) свойство среднего.

\begin{thmbox}{Лемма 22.5 (принцип максимума для функций со свойством среднего)}
Пусть $G$ --- ограниченная область, $u\in C(\clos G)$ вещественна и обладает локальным свойством среднего в $G$. Тогда
\[
\min_{\partial G}u\le u(z)\le\max_{\partial G}u\qquad\text{для всех }z\in G ,
\]
и если $u$ достигает $\max_{\clos G}u$ (или $\min_{\clos G}u$) в точке $G$, то $u\equiv\mathrm{const}$ в $\clos G$.
\end{thmbox}
\begin{proof}
Пусть $M=\max_{\clos G}u$ (существует: $\clos G$ компакт), $E=\{z\in G:u(z)=M\}$.

\emph{$E$ замкнуто в $G$}: это прообраз точки при непрерывном отображении $u|_G$.

\emph{$E$ открыто.} Пусть $a\in E$. Для $r\in(0,\rho(a))$ по свойству среднего
\[
0=M-u(a)=\frac1{2\pi}\int_0^{2\pi}\bigl(M-u(a+re^{i\vartheta})\bigr)\,d\vartheta ,
\]
подынтегральная функция непрерывна и $\ge0$, значит тождественно равна нулю (см. деталь в доказательстве теоремы 14.1). Итак, $u=M$ на всех
окружностях $|z-a|=r$, $0<r<\rho(a)$, т.\,е. $B(a,\rho(a))\subset E$.

Так как $G$ связна, либо $E=\varnothing$, либо $E=G$. Во втором случае $u\equiv M$ на $G$ и, по непрерывности, на $\clos G$. Это доказывает
второе утверждение. Если $u$ не постоянна, то $E=\varnothing$, т.\,е. максимум $M$ достигается только на $\partial G$, и $u(z)<M=\max_{\partial G}u$
при $z\in G$; если постоянна --- неравенства очевидны. Для минимума применяем доказанное к $-u$ (свойство среднего линейно).
\end{proof}

\begin{thmbox}{Теорема 22.6 (гармоничность функций со свойством среднего)}
Пусть $D$ --- область в $\CC$, $h\in C(D)$ вещественна и для любого замкнутого круга $\clos{B(a,R)}\subset D$
\[
h(a)=\frac1{2\pi}\int_0^{2\pi}h(a+Re^{i\vartheta})\,d\vartheta .
\]
Тогда $h\in\pIVHarm(D)$. (Достаточно даже локального свойства среднего.)
\end{thmbox}
\begin{proof}
Фиксируем произвольный круг $\clos{B(a,R)}\subset D$; докажем, что $h\in\pIVHarm(B(a,R))$. Так как каждая точка $D$ является центром
такого круга, это даст $h\in C^2(D)$ и $\Delta h=0$ в $D$.

\emph{Шаг 1: гармоническая функция с теми же граничными значениями.} Функция $h_0=h|_{\partial B(a,R)}$ непрерывна. По теореме о
решении задачи Дирихле в круге интегралом Пуассона (теорема 24.7 в билете 24, перенесённая на круг $B(a,R)$ сдвигом $z\mapsto z-a$,
см. пример~1 выше) функция
\[
H(z)=\begin{cases}\displaystyle\frac1{2\pi R}\int_{|\zeta-a|=R}\frac{R^2-|z-a|^2}{|\zeta-z|^2}\,h(\zeta)\,|d\zeta|,& |z-a|<R,\\[2mm]
h(z),&|z-a|=R,\end{cases}
\]
принадлежит $\pIVHarm(B(a,R))\cap C(\clos{B(a,R)})$. (Доказательство этой теоремы в билете 24 не использует теорему 22.6, так что
порочного круга нет.)

\emph{Шаг 2: разность обладает свойством среднего.} Положим $u=h-H\in C(\clos{B(a,R)})$. Пусть $b\in B(a,R)$ и $\rho(b)=R-|b-a|>0$.
Для $0<r<\rho(b)$ выберем $r'\in(r,\rho(b))$; тогда $\clos{B(b,r')}\subset B(a,R)\subset D$. Свойство среднего на окружности $|z-b|=r$
выполнено для $h$ (по условию, так как $\clos{B(b,r)}\subset D$) и для $H$ (по предложению 14.2, так как
$H\in\pIVHarm(B(b,r'))\cap C(\clos{B(b,r')})$). По линейности оно выполнено для $u$. Итак, $u$ обладает локальным свойством среднего в $B(a,R)$.

\emph{Шаг 3: принцип максимума.} На $\partial B(a,R)$ имеем $u=h-h=0$. По лемме 22.5 (с $G=B(a,R)$)
$0=\min_{\partial B}u\le u\le\max_{\partial B}u=0$ в $B(a,R)$, т.\,е. $h\equiv H$ в $B(a,R)$. Значит, $h\in\pIVHarm(B(a,R))$.

\emph{Локальный вариант.} Если дано лишь локальное свойство среднего, возьмём для $a\in D$ круг $\clos{B(a,R)}\subset D$ любого радиуса;
в шаге 2 свойство среднего для $h$ в точке $b$ используется лишь при малых $r<\min(\rho(b),\rho_h(b))$, чего достаточно для леммы 22.5.
\end{proof}

\begin{remarkbox}{Следствия теоремы 22.6}
\begin{enumerate}
\item \emph{Характеризация:} $\pIVHarm(D)=\{h\in C(D): h$ обладает свойством среднего в $D\}$ (включение $\subset$ --- предложение 14.2).
В частности, гармоничность можно определять без всяких производных, а гладкость $C^\infty$ получается автоматически.
\item \emph{Ещё одно доказательство аналога теоремы Вейерштрасса} (предложение 14.7, билет 23): если $h_n\in\pIVHarm(D)$ и $h_n\to h$ локально
равномерно в $D$, то $h$ непрерывна, и для $\clos{B(a,r)}\subset D$ имеем $h_n(a+re^{i\vartheta})\to h(a+re^{i\vartheta})$ равномерно по
$\vartheta$, так что в равенстве $h_n(a)=\frac1{2\pi}\int_0^{2\pi}h_n(a+re^{i\vartheta})d\vartheta$ можно перейти к пределу. По теореме 22.6
$h\in\pIVHarm(D)$.
\item Непрерывность $h$ используется в доказательстве дважды: чтобы интеграл Пуассона от $h|_{\partial B(a,R)}$ давал функцию, \emph{непрерывно} примыкающую к $h$ на окружности (шаг 1), и в лемме 22.5 (интеграл от непрерывной неотрицательной функции равен нулю лишь для тождественного нуля; существование максимума на компакте).
\end{enumerate}
\end{remarkbox}

\begin{remarkbox}{Коротко: план ответа на экзамене}
\begin{enumerate}
\item Постановка задачи Дирихле; единственность и устойчивость $\|h_1-h_2\|_{\clos D}=\|h_{0,1}-h_{0,2}\|_{\partial D}$ (принцип максимума).
\item Предл. 14.4: локально $H=\Re F$, тогда $H\circ\varphi=\Re(F\circ\varphi)$; второе доказательство: $\Delta(H\circ\varphi)=|\varphi'|^2(\Delta H)\circ\varphi$
через $\Delta=4\partial_z\partial_{\bar z}$ и цепное правило.
\item Метод конформного отображения: $\varphi\in\pIVConf(D,G)$ + Каратеодори, $H_0=h_0\circ\varphi^{-1}$, $h=H\circ\varphi$.
\item Следствие: в жордановой области задача Дирихле разрешима (Риман + Каратеодори + круг); пример: полуплоскость, ядро $\frac{y}{\pi((t-x)^2+y^2)}$.
\item Теорема Лебега (формулировка): ограниченная односвязная область --- всегда разрешима; проколотый круг --- контрпример без односвязности.
\item Принцип максимума для непрерывных функций со свойством среднего (множество точек максимума открыто-замкнуто).
\item Теорема: свойство среднего $\Rightarrow$ гармоничность: в круге сравниваем $h$ с интегралом Пуассона $H$ её граничных значений;
$h-H$ имеет свойство среднего и нуль на границе $\Rightarrow$ $h\equiv H$.
\end{enumerate}
\end{remarkbox}
\ticket{23}{Разложение гармонической функции в круге в ряд по однородным гармоническим многочленам. Аналог теоремы Вейерштрасса для гармонических функций}
\label{t23:start}

\begin{defs}
\textbf{Гармонические функции} (билет 21): $h\in\pIVHarm(D)$, если $h\in C^2(D)$ вещественна и $\Delta h=0$. Если $D$ --- круг, то
$h=\Re f$ для некоторой $f\in H(D)$, единственной с точностью до слагаемого $iC$, $C\in\RR$ (предложение 14.1).

\textbf{Обозначения:} $B_R=B(0,R)$, $C_R=\{|z|=R\}=\partial B_R$; полярные координаты $z=re^{i\varphi}$; $\|h\|_K=\sup_K|h|$.

\textbf{Однородный многочлен степени $n$} от $x,y$: $P(x,y)=\sum_{j+k=n}p_{jk}x^jy^k$; он удовлетворяет $P(tx,ty)=t^nP(x,y)$.
\textbf{Однородный гармонический полином (многочлен) порядка $n$} --- всякая вещественная линейная комбинация функций
\[
\Re(x+iy)^n=\Re z^n=r^n\cos n\varphi,\qquad \Im(x+iy)^n=\Im z^n=r^n\sin n\varphi .
\]
Например, $n=1$: $x,y$; $n=2$: $x^2-y^2$, $2xy$; $n=3$: $x^3-3xy^2$, $3x^2y-y^3$.

\textbf{Тригонометрический ряд:} $\frac{\alpha_0}2+\sum_{n\ge1}(\alpha_n\cos n\varphi+\beta_n\sin n\varphi)$.
\textbf{Коэффициенты Фурье} $2\pi$-периодической непрерывной $F$:
$\alpha_n=\frac1\pi\int_{-\pi}^{\pi}F(\vartheta)\cos n\vartheta\,d\vartheta$ ($n\ge0$), $\beta_n=\frac1\pi\int_{-\pi}^{\pi}F(\vartheta)\sin n\vartheta\,d\vartheta$ ($n\ge1$);
ряд Фурье --- тригонометрический ряд с этими коэффициентами.

\textbf{Ряд Тейлора.} Если $f\in H(B(0,\rho))$, то $f(z)=\sum_{n\ge0}c_nz^n$ в $B(0,\rho)$, причём для $0<\rho'<\rho$ выполнены
\emph{неравенства Коши} $|c_n|\le\|f\|_{C_{\rho'}}/\rho'^{\,n}$; ряд сходится абсолютно и равномерно на каждом $\clos{B_R}$, $R<\rho$.

\textbf{Сходимость.} $h_n\rightrightarrows h$ на $K$ --- равномерная сходимость: $\|h_n-h\|_K\to0$. \emph{Локально равномерная} сходимость
в $D$ --- равномерная на каждом компакте $K\subset D$ (эквивалентно: каждая точка имеет окрестность, где сходимость равномерна).
\emph{Признак Вейерштрасса:} если $|u_n|\le M_n$ на $K$ и $\sum M_n<\infty$, то $\sum u_n$ сходится абсолютно и равномерно на $K$.
\emph{Критерий Коши:} последовательность функций, фундаментальная в норме $\|\cdot\|_K$, сходится равномерно на $K$.

\textbf{Теорема Вейерштрасса для голоморфных функций:} если $f_n\in H(U)$ и $f_n\to f$ локально равномерно в открытом $U$, то $f\in H(U)$
и $f_n^{(m)}\to f^{(m)}$ локально равномерно в $U$ при каждом $m$.
\end{defs}

\subsection{Однородные гармонические многочлены}

\begin{thmbox}{Лемма 23.1 (оправдание термина)}
При $n\ge1$ вещественный однородный многочлен $P(x,y)$ степени $n$ гармоничен в $\CC$ тогда и только тогда, когда
$P=\alpha\Re z^n+\beta\Im z^n$ с $\alpha,\beta\in\RR$. Таким образом, однородные гармонические полиномы порядка $n$ --- это в точности
однородные многочлены степени $n$, являющиеся гармоническими функциями; при $n\ge1$ они образуют двумерное пространство.
\end{thmbox}
\begin{proof}
$\Leftarrow$: $z^n=(x+iy)^n=\sum_k\binom nk x^k(iy)^{n-k}$ --- все мономы степени $n$, так что $\Re z^n,\Im z^n$ однородны степени $n$; они
гармоничны как $\Re$ и $\Im$ голоморфной функции (билет 21).

$\Rightarrow$: Подставив $x=\frac{z+\bar z}2$, $y=\frac{z-\bar z}{2i}$, запишем $P=\sum_{k=0}^n\gamma_kz^k\bar z^{n-k}$, $\gamma_k\in\CC$.
Функции $z^k\bar z^{m-k}$, $0\le k\le m$, линейно независимы над $\CC$: на окружности $z=e^{i\varphi}$ они равны $e^{i(2k-m)\varphi}$ с попарно
различными частотами, а $\int_{-\pi}^{\pi}e^{ip\varphi}\overline{e^{iq\varphi}}\,d\varphi=2\pi\delta_{pq}$ (если $\sum_k\lambda_ke^{i(2k-m)\varphi}\equiv0$,
то, умножая на $e^{-i(2j-m)\varphi}$ и интегрируя, получаем $\lambda_j=0$). Поэтому представление единственно.
По правилам Виртингера ($\partial_zz=1$, $\partial_z\bar z=0$, $\partial_{\bar z}\bar z=1$, $\partial_{\bar z}z=0$, формула Лейбница) и лемме 21.1
\[
\Delta P=4\partial_z\partial_{\bar z}P=4\sum_{k=1}^{n-1}\gamma_k\,k(n-k)\,z^{k-1}\bar z^{\,n-1-k}.
\]
Если $\Delta P\equiv0$, то по линейной независимости (при $m=n-2$) $\gamma_k=0$ для $1\le k\le n-1$, т.\,е. $P=\gamma_nz^n+\gamma_0\bar z^n$.
Так как $P$ вещественен, $P=\clos P=\bar\gamma_n\bar z^n+\bar\gamma_0z^n$, и по единственности $\gamma_0=\bar\gamma_n$. Итак,
$P=\gamma_nz^n+\overline{\gamma_nz^n}=2\Re(\gamma_nz^n)=2(\Re\gamma_n)\Re z^n-2(\Im\gamma_n)\Im z^n$.
Функции $\Re z^n=r^n\cos n\varphi$ и $\Im z^n=r^n\sin n\varphi$ линейно независимы над $\RR$ (на $C_1$ это $\cos n\varphi$ и $\sin n\varphi$), поэтому размерность равна $2$.
\end{proof}

\subsection{Разложение в ряд (предложение 14.5)}

\begin{thmbox}{Лемма 23.2 (ортогональность и коэффициенты равномерно сходящегося ряда)}
(а) При целых $n,m\ge0$:
\begin{gather*}
\int_{-\pi}^{\pi}\cos n\vartheta\cos m\vartheta\,d\vartheta=\begin{cases}0,&n\ne m,\\ \pi,&n=m\ge1,\\ 2\pi,&n=m=0,\end{cases}\\
\int_{-\pi}^{\pi}\sin n\vartheta\sin m\vartheta\,d\vartheta=\pi\delta_{nm}\ (n,m\ge1),\qquad
\int_{-\pi}^{\pi}\sin n\vartheta\cos m\vartheta\,d\vartheta=0 .
\end{gather*}
(б) Если тригонометрический ряд $\frac{\alpha_0}2+\sum_{n\ge1}(\alpha_n\cos n\varphi+\beta_n\sin n\varphi)$ сходится равномерно на $[-\pi,\pi]$
к функции $F$, то $F$ непрерывна и $\alpha_n,\beta_n$ --- коэффициенты Фурье функции $F$.
\end{thmbox}
\begin{proof}
(а) $\int_{-\pi}^{\pi}\cos k\vartheta\,d\vartheta=0$ при целом $k\ne0$ и $=2\pi$ при $k=0$; $\int_{-\pi}^{\pi}\sin k\vartheta\,d\vartheta=0$ при всех $k$.
Остаётся применить формулы $\cos a\cos b=\frac12(\cos(a-b)+\cos(a+b))$, $\sin a\sin b=\frac12(\cos(a-b)-\cos(a+b))$,
$\sin a\cos b=\frac12(\sin(a+b)+\sin(a-b))$ с $a=n\vartheta$, $b=m\vartheta$ (при $n,m\ge1$ число $n+m\ne0$).

(б) $F$ непрерывна как равномерный предел непрерывных частичных сумм. Умножение на ограниченную функцию $\cos m\varphi$ сохраняет
равномерную сходимость, поэтому ряд можно интегрировать почленно:
$\frac1\pi\int_{-\pi}^{\pi}F\cos m\vartheta\,d\vartheta=\frac1\pi\bigl(\frac{\alpha_0}2\int\cos m\vartheta+\sum_n\alpha_n\int\cos n\vartheta\cos m\vartheta+\sum_n\beta_n\int\sin n\vartheta\cos m\vartheta\bigr)$,
что по (а) равно $\alpha_m$ (при $m=0$: $\frac1\pi\cdot\frac{\alpha_0}2\cdot2\pi=\alpha_0$). Аналогично с $\sin m\varphi$ получаем $\beta_m$.
\end{proof}

\begin{thmbox}{Предложение 14.5}
При фиксированных $R>0$ и $\eps>0$ пусть $h\in\pIVHarm(B_{R+\eps})$. Тогда в замкнутом круге $\clos{B_R}$
\begin{equation}\label{t23:e141}
h(re^{i\varphi})=\frac{A_0}2+\sum_{n=1}^\infty r^n\bigl(A_n\cos n\varphi+B_n\sin n\varphi\bigr),\qquad0\le r\le R,\tag{14.1}
\end{equation}
где
\begin{align}
A_n&=\frac1{\pi R^n}\int_{-\pi}^{\pi}h(Re^{i\vartheta})\cos n\vartheta\,d\vartheta,\quad n\in\ZZ_+,\tag{14.2}\\
B_n&=\frac1{\pi R^n}\int_{-\pi}^{\pi}h(Re^{i\vartheta})\sin n\vartheta\,d\vartheta,\quad n\in\NN,\tag{14.3}
\end{align}
а ряд \eqref{t23:e141} сходится абсолютно и равномерно на $\clos{B_R}$. Иными словами, $h=\frac{A_0}2+\sum_{n\ge1}(A_n\Re z^n+B_n\Im z^n)$ ---
разложение в ряд по однородным гармоническим полиномам.
\end{thmbox}
\begin{proof}
\emph{Шаг 1: голоморфная функция.} Круг $B_{R+\eps}$ односвязен, поэтому по предложению 14.1 существует $f_1\in H(B_{R+\eps})$ с $\Re f_1=h$;
все такие функции имеют вид $f_1+iC$. Условие $f(0)=h(0)$ означает $\Im f(0)=0$ и выделяет ровно одну: $f=f_1-i\Im f_1(0)$.
Итак, существует единственная $f\in H(B_{R+\eps})$ с $\Re f=h$, $f(0)=h(0)$.

\emph{Шаг 2: ряд Тейлора.} $f(z)=c_0+\sum_{n\ge1}c_nz^n$ в $B_{R+\eps}$. Положим $a_n=\Re c_n$, $b_n=\Im c_n$; тогда $c_0=f(0)=h(0)\in\RR$, т.\,е. $c_0=a_0$.
Выберем $\rho\in(R,R+\eps)$; по неравенствам Коши $|c_n|\le\|f\|_{C_\rho}\rho^{-n}$, поэтому $\sum_n|c_n|R^n\le\|f\|_{C_\rho}\sum_n(R/\rho)^n<\infty$.

\emph{Шаг 3: вещественная часть.} При $z=re^{i\varphi}$:
$c_nz^n=(a_n+ib_n)r^n(\cos n\varphi+i\sin n\varphi)$, $\Re(c_nz^n)=r^n(a_n\cos n\varphi-b_n\sin n\varphi)$. Так как вещественная часть суммы сходящегося
ряда равна сумме ряда из вещественных частей,
\begin{equation}\label{t23:e144}
h(re^{i\varphi})=a_0+\sum_{n=1}^\infty r^n\bigl(a_n\cos n\varphi-b_n\sin n\varphi\bigr),\qquad 0\le r<R+\eps.\tag{14.4}
\end{equation}
При $r\le R$ модуль $n$-го члена не превосходит $r^n(|a_n|+|b_n|)\le2|c_n|R^n$, и по признаку Вейерштрасса (шаг 2) ряд \eqref{t23:e144}
сходится абсолютно и равномерно на $\clos{B_R}$. Это даёт \eqref{t23:e141} с
\[
A_0=2a_0,\qquad A_n=a_n,\quad B_n=-b_n\quad(n\in\NN).
\]

\emph{Шаг 4: формулы для коэффициентов.} При $r=R$ \eqref{t23:e144} --- равномерно сходящийся по $\varphi\in[-\pi,\pi]$ тригонометрический ряд с
коэффициентами $\alpha_0=2a_0$, $\alpha_n=R^na_n$, $\beta_n=-R^nb_n$, сумма которого равна непрерывной $2\pi$-периодической
функции $\varphi\mapsto h(Re^{i\varphi})$. По лемме 23.2(б) это коэффициенты Фурье этой функции:
\begin{gather*}
2a_0=\frac1\pi\int_{-\pi}^{\pi}h(Re^{i\vartheta})\,d\vartheta,\qquad
R^na_n=\frac1\pi\int_{-\pi}^{\pi}h(Re^{i\vartheta})\cos n\vartheta\,d\vartheta\ \ (n\ge1),\\
-R^nb_n=\frac1\pi\int_{-\pi}^{\pi}h(Re^{i\vartheta})\sin n\vartheta\,d\vartheta\ \ (n\ge1).
\end{gather*}
Деля на $R^n$, получаем ровно $A_n$ и $B_n$ из \eqref{t23:e141} в виде (14.2), (14.3).
\end{proof}

\begin{remarkbox}{Замечание (опечатка на слайде 347)}
На слайде написано <<(14.1) установлено при $A_0=a_0$>>. Правильно $A_0=2a_0$: в (14.1) стоит $A_0/2$, а по (14.2) с $n=0$ и теореме о
среднем $A_0=\frac1\pi\int_{-\pi}^{\pi}h(Re^{i\vartheta})d\vartheta=2h(0)=2a_0$. Именно это соотношение $|a_0|=|A_0|/2$ используется далее на слайде 348.
\end{remarkbox}

\begin{remarkbox}{Дополнения к предложению 14.5}
\begin{enumerate}
\item \emph{Коэффициенты не зависят от $R$:} $A_n=a_n$, $B_n=-b_n$ определяются функцией $h$ (через $f$), поэтому формулы (14.2), (14.3) дают одно и
то же при любом радиусе $R'\in(0,R+\eps)$ в роли $R$. Попутно: $A_n-iB_n=c_n$ --- коэффициенты Тейлора голоморфной $f$ с $\Re f=h$, $f(0)=h(0)$.
\item \emph{Гармоническая функция в открытом круге.} Если $h\in\pIVHarm(B_R)$, то для любого $R'<R$ применимо предложение 14.5 (с $R'$ и
$\eps=R-R'$); с учётом п.~1 ряд \eqref{t23:e141} с коэффициентами, вычисленными по любому радиусу $R'<R$, сходится к $h$ абсолютно и
равномерно на каждом $\clos{B_{R'}}$, т.\,е. локально равномерно в $B_R$.
\item \emph{Единственность разложения.} Если $h=\frac{\alpha_0}2+\sum_n(\alpha_n\Re z^n+\beta_n\Im z^n)$ с рядом, равномерно сходящимся на $C_R$,
то по лемме 23.2(б) $\alpha_n=A_n$, $\beta_n=B_n$.
\item \emph{Оценка коэффициентов.} Из (14.2), (14.3): $|A_n|,|B_n|\le\frac{2}{R^n}\|h\|_{C_R}$.
\end{enumerate}
\end{remarkbox}

\subsection{Оценка голоморфной функции через её вещественную часть (предложение 14.6)}

\begin{thmbox}{Предложение 14.6}
В обозначениях предложения 14.5 и его доказательства ($h\in\pIVHarm(B_{R+\eps})$, $f\in H(B_{R+\eps})$, $\Re f=h$, $f(0)=h(0)$) при $r\in(0,R)$
\[
\|f\|_{\clos{B_r}}\le\frac{R+2r}{R-r}\,\|h\|_{\clos{B_R}} .
\]
\end{thmbox}
\begin{proof}
Положим $M=\|h\|_{\clos{B_R}}$; по принципу максимума и минимума (теорема 14.1) $M=\|h\|_{C_R}$. Из (14.2), (14.3) и $A_0=2a_0$:
\begin{gather*}
|a_0|=\frac{|A_0|}2\le\frac1{2\pi}\int_{-\pi}^{\pi}|h(Re^{i\vartheta})|\,d\vartheta\le M,\\
|a_n|=|A_n|\le\frac{2\pi M}{\pi R^n}=\frac{2M}{R^n},\qquad |b_n|=|B_n|\le\frac{2M}{R^n}\quad(n\in\NN).
\end{gather*}
Отсюда $|c_n|=\sqrt{a_n^2+b_n^2}\le\frac{2\sqrt2\,M}{R^n}\le\frac{3M}{R^n}$. Пусть $|z|=\rho\le r$ и $q=\rho/R<1$. Тогда
\[
|f(z)|\le|c_0|+\sum_{n\ge1}|c_n|\rho^n\le M\Bigl(1+3\sum_{n=1}^\infty q^n\Bigr)=M\Bigl(1+\frac{3q}{1-q}\Bigr)=M\,\frac{1+2q}{1-q}=M\,\frac{R+2\rho}{R-\rho}.
\]
Функция $\rho\mapsto\frac{R+2\rho}{R-\rho}$ возрастает на $[0,R)$ (числитель растёт, знаменатель убывает и положителен), поэтому
$|f(z)|\le M\frac{R+2r}{R-r}$ при всех $|z|\le r$.
\end{proof}

\begin{remarkbox}{Замечания к слайду 348}
(1) На слайде написано $|c_n|<3M/R^n$ (строгое неравенство) --- при $M=0$ оно неверно; правильно $\le$, на результат это не влияет.
Слайд оценивает лишь $|f(re^{i\vartheta})|$; для нормы на всём $\clos{B_r}$ нужна монотонность по $\rho$, отмеченная выше.

(2) \emph{Более точная оценка.} Так как $c_n=a_n+ib_n=A_n-iB_n=\frac1{\pi R^n}\int_{-\pi}^{\pi}h(Re^{i\vartheta})e^{-in\vartheta}d\vartheta$, то
$|c_n|\le2M/R^n$, и тем же вычислением $\|f\|_{\clos{B_r}}\le M\bigl(1+\frac{2q}{1-q}\bigr)=\frac{R+r}{R-r}M$. Для дальнейшего годится любая
оценка вида $C(r,R)\,M$.

(3) \emph{Сдвинутый вариант.} Если $h\in\pIVHarm(B(a,R'))$, $R<R'$, $f\in H(B(a,R'))$, $\Re f=h$, $f(a)=h(a)$, то
$\|f\|_{\clos{B(a,r)}}\le\frac{R+2r}{R-r}\|h\|_{\clos{B(a,R)}}$ при $0<r<R$: применить предложение 14.6 к $h(a+z)$ и $f(a+z)$
(сдвиг сохраняет гармоничность и голоморфность).
\end{remarkbox}

\subsection{Аналог теоремы Вейерштрасса (предложение 14.7)}

\begin{thmbox}{Предложение 14.7 (аналог теоремы Вейерштрасса для гармонических функций)}
Пусть $D$ --- область в $\CC$ и последовательность $h_n\in\pIVHarm(D)$ сходится к функции $h$ локально равномерно в $D$. Тогда $h\in\pIVHarm(D)$.
\end{thmbox}
\begin{proof}
Фиксируем произвольную точку $a\in D$ и $R'>0$ с $\clos{B(a,R')}\subset D$. Пусть $r\in(0,R')$ и $R\in(r,R')$. Докажем, что $h\in\pIVHarm(B(a,r))$.

\emph{Голоморфные функции.} Круг $B(a,R')$ односвязен; по предложению 14.1 (как в шаге 1 доказательства предложения 14.5) для каждого $n$ найдётся
единственная $f_n\in H(B(a,R'))$ с $\Re f_n=h_n$ и $f_n(a)=h_n(a)$.

\emph{Оценка.} Для $n,m\in\NN$ функция $f_n-f_m\in H(B(a,R'))$ имеет вещественную часть $h_n-h_m\in\pIVHarm(B(a,R'))$ и значение
$(f_n-f_m)(a)=(h_n-h_m)(a)$ в центре. По сдвинутому варианту предложения 14.6
\[
\|f_n-f_m\|_{\clos{B(a,r)}}\le\frac{R+2r}{R-r}\,\|h_n-h_m\|_{\clos{B(a,R)}} .
\]

\emph{Предельный переход.} Круг $\clos{B(a,R)}$ --- компакт в $D$, поэтому $h_n\rightrightarrows h$ на нём, и $\{h_n\}$ фундаментальна в
$\|\cdot\|_{\clos{B(a,R)}}$. Из оценки $\{f_n\}$ фундаментальна в $\|\cdot\|_{\clos{B(a,r)}}$ и по критерию Коши $f_n\rightrightarrows f$ на
$\clos{B(a,r)}$ для некоторой функции $f$. По теореме Вейерштрасса для голоморфных функций $f\in H(B(a,r))$. Наконец,
$\Re f=\lim\Re f_n=\lim h_n=h$ в $B(a,r)$, так что $h=\Re f\in\pIVHarm(B(a,r))$ (билет 21).

Так как $a\in D$ произвольна, $h\in C^2(D)$ и $\Delta h=0$ в окрестности каждой точки, т.\,е. $h\in\pIVHarm(D)$.
\end{proof}

\begin{thmbox}{Следствие 23.3 (сходимость производных)}
В условиях предложения 14.7 для любых $j,k\ge0$ имеем $\partial_x^j\partial_y^kh_n\to\partial_x^j\partial_y^kh$ локально равномерно в $D$.
\end{thmbox}
\begin{proof}
Пусть $F$ голоморфна. Так как $F'_x=F'$ и $F'_y=iF'$, индукцией получаем $\partial_x^j\partial_y^kF=i^kF^{(j+k)}$; поскольку дифференцирование по
вещественным переменным перестановочно с взятием $\Re$, для $h=\Re F$
\[
\partial_x^j\partial_y^k h=\Re\bigl(i^kF^{(j+k)}\bigr).
\]
Возьмём $a$, $r$, $f_n$, $f$ из доказательства предложения 14.7. Так как $f_n\rightrightarrows f$ на $\clos{B(a,r)}$, то по теореме Вейерштрасса
$f_n^{(m)}\to f^{(m)}$ равномерно на компакте $\clos{B(a,r/2)}\subset B(a,r)$ при каждом $m$. Следовательно,
\[
\bigl\|\partial_x^j\partial_y^kh_n-\partial_x^j\partial_y^kh\bigr\|_{\clos{B(a,r/2)}}\le\bigl\|f_n^{(j+k)}-f^{(j+k)}\bigr\|_{\clos{B(a,r/2)}}\to0 .
\]
Любой компакт $K\subset D$ покрывается конечным числом таких открытых кругов $B(a,r_a/2)$, $a\in K$, откуда равномерная сходимость на $K$.
\end{proof}

\begin{remarkbox}{Замечания}
\begin{enumerate}
\item \emph{Ряды.} Если ряд $\sum_n u_n$ гармонических в $D$ функций сходится локально равномерно, то его сумма гармонична и ряд можно
почленно дифференцировать любое число раз (применить предложение 14.7 и следствие 23.3 к частичным суммам). Например, этим объясняется,
что ряд (14.1) с коэффициентами $|A_n|,|B_n|\le CR^{-n}$ задаёт гармоническую функцию в $B_R$ (это используется в билете 24).
\item \emph{Другое доказательство} предложения 14.7 --- через характеризацию гармонических функций свойством среднего (теорема 22.6, билет 22).
\item Сравнение с вещественным анализом: равномерный предел гладких функций вещественной переменной может быть не дифференцируем;
гармоничность (как и голоморфность) --- <<жёсткое>> свойство, сохраняющееся при равномерных пределах.
\item Существенна именно равномерная сходимость самих $h_n$: из сходимости $\Re$ без нормировки $f_n(a)=h_n(a)$ сходимость $f_n$ не следует
(к $f_n$ можно прибавлять $in$).
\end{enumerate}
\end{remarkbox}

\begin{remarkbox}{Коротко: план ответа на экзамене}
\begin{enumerate}
\item Однородные гармонические полиномы $r^n\cos n\varphi=\Re z^n$, $r^n\sin n\varphi=\Im z^n$; (лемма) других однородных гармонических нет.
\item Предл. 14.5: $h=\Re f$, $f(0)=h(0)$; ряд Тейлора $f=\sum c_nz^n$; $\Re(c_nz^n)=r^n(a_n\cos n\varphi-b_n\sin n\varphi)$.
\item Абсолютная и равномерная сходимость на $\clos{B_R}$ --- из неравенств Коши на окружности радиуса $\rho\in(R,R+\eps)$.
\item Коэффициенты (14.2), (14.3) --- коэффициенты Фурье $h(Re^{i\vartheta})$ (почленное интегрирование + ортогональность); $A_0=2a_0$.
\item Предл. 14.6: $|c_n|\le3M/R^n$ $\Rightarrow$ $\|f\|_{\clos{B_r}}\le\frac{R+2r}{R-r}\|h\|_{\clos{B_R}}$.
\item Предл. 14.7: нормированные $f_n$ с $\Re f_n=h_n$ фундаментальны по оценке 14.6 $\Rightarrow$ $f_n\rightrightarrows f$, $h=\Re f$.
\item Следствие: производные сходятся локально равномерно, $\partial_x^j\partial_y^kh=\Re(i^kf^{(j+k)})$.
\end{enumerate}
\end{remarkbox}

\ticket{24}{Метод Фурье и формула Пуассона для решения задачи Дирихле для гармонических функций в круге}
\label{t24:start}

\begin{defs}
\textbf{Обозначения.} $B_R=B(0,R)$, $C_R=\{|z|=R\}=\partial B_R$, $z=re^{i\varphi}$ (полярные координаты), $q=r/R$;
$\|u\|_K=\sup_K|u|$; $\pIVHarm(D)$ --- гармонические в $D$ функции (вещественные, класса $C^2$, $\Delta h=0$; билет 21).

\textbf{Задача Дирихле в круге.} Дана вещественная $h_0\in C(C_R)$. Найти $h\in\pIVHarm(B_R)\cap C(\clos{B_R})$ с $h|_{C_R}=h_0$.
Удобно перейти к функции угла: $H_0(\varphi)=h_0(Re^{i\varphi})$. Тогда $H_0\in C([-\pi,\pi])$ и $H_0(-\pi)=H_0(\pi)$; продолжая $H_0$ с периодом
$2\pi$, получаем непрерывную $2\pi$-периодическую функцию на $\RR$ (и обратно, всякая такая функция задаёт $h_0\in C(C_R)$).
Решение, если существует, единственно (граничная теорема единственности, предложение 14.3, билет 21).

\textbf{Тригонометрический ряд и ряд Фурье.} Для $2\pi$-периодической непрерывной (или кусочно непрерывной) $F$ её
\emph{коэффициенты Фурье}
\[
\alpha_n=\frac1\pi\int_{-\pi}^{\pi}F(\vartheta)\cos n\vartheta\,d\vartheta\ (n\ge0),\qquad
\beta_n=\frac1\pi\int_{-\pi}^{\pi}F(\vartheta)\sin n\vartheta\,d\vartheta\ (n\ge1),
\]
\emph{ряд Фурье} $\frac{\alpha_0}2+\sum_{n\ge1}(\alpha_n\cos n\varphi+\beta_n\sin n\varphi)$, его частичные суммы
$S_N(\varphi)=\frac{\alpha_0}2+\sum_{n=1}^N(\alpha_n\cos n\varphi+\beta_n\sin n\varphi)$. В обозначениях (14.2), (14.3) (билет 23) с $h_0$ вместо $h$:
\[
A_n=\frac1{\pi R^n}\int_{-\pi}^{\pi}H_0(\vartheta)\cos n\vartheta\,d\vartheta=\frac{\alpha_n}{R^n},\qquad
B_n=\frac1{\pi R^n}\int_{-\pi}^{\pi}H_0(\vartheta)\sin n\vartheta\,d\vartheta=\frac{\beta_n}{R^n},
\]
где $\alpha_n,\beta_n$ --- коэффициенты Фурье $H_0$. Сразу видно: $|A_n|,|B_n|\le 2\|h_0\|_{C_R}/R^n$.

\textbf{Класс $C^{1+1}$.} $f\in C^{1+1}([-\pi,\pi])$, если $f\in C^1([-\pi,\pi])$ и существует разбиение
$-\pi=t_0<t_1<\dots<t_J=\pi$, такое что на каждом отрезке $I_j=[t_{j-1},t_j]$ сужение $f|_{I_j}$ принадлежит $C^2(I_j)$
(с односторонними производными на концах). Тогда $\int_{-\pi}^{\pi}f''\,g:=\sum_j\int_{I_j}(f|_{I_j})''g$. Пример: $f(\vartheta)=\vartheta|\vartheta|$ (производная $2|\vartheta|$ непрерывна, вторая производная скачет в нуле).

\textbf{Модуль непрерывности} функции $F$ на $\RR$: $\omega_F(\delta)=\sup\{|F(s)-F(t)|:|s-t|\le\delta\}$. Непрерывная периодическая $F$
равномерно непрерывна, т.\,е. $\omega_F(\delta)\to0$ при $\delta\to0+$.

\textbf{Сходимость.} \emph{Признак Вейерштрасса:} если $|u_n|\le M_n$ на $K$ и $\sum M_n<\infty$, то $\sum u_n$ сходится на $K$ абсолютно и равномерно;
равномерно сходящийся ряд непрерывных функций можно почленно интегрировать по отрезку. Равномерный предел непрерывных функций непрерывен.

\textbf{Факты из билетов 21, 23.} \emph{Принцип максимума и минимума} (теорема 14.1): для $u\in\pIVHarm(D)\cap C(\clos D)$, $D$ ограничена,
$\|u\|_{\clos D}=\|u\|_{\partial D}$. \emph{Аналог теоремы Вейерштрасса} (предложение 14.7): локально равномерный предел гармонических функций гармоничен.
Функции $\Re z^n=r^n\cos n\varphi$, $\Im z^n=r^n\sin n\varphi$ гармоничны в $\CC$ (однородные гармонические полиномы).

\textbf{Ядро Пуассона} круга $B_R$:
\[
P(z,\zeta)=\frac1{2\pi R}\cdot\frac{R^2-|z|^2}{|\zeta-z|^2},\qquad z\in B_R,\ \zeta\in C_R;
\]
в угловой записи ($z=re^{i\varphi}$, $\zeta=Re^{i\vartheta}$, $|d\zeta|=R\,d\vartheta$) $P(z,\zeta)\,|d\zeta|=P_r(\vartheta-\varphi)\,d\vartheta$, где
\[
P_r(\psi)=\frac1{2\pi}\cdot\frac{R^2-r^2}{R^2-2Rr\cos\psi+r^2},\qquad 0\le r<R .
\]
\emph{Интеграл Пуассона} функции $h_0\in C(C_R)$: $\mathcal P[h_0](z)=\int_{C_R}P(z,\zeta)h_0(\zeta)\,|d\zeta|$, $z\in B_R$.

\textbf{Аппроксимативная единица.} Семейство непрерывных функций $K_r$ на $[-\pi,\pi]$, $r\to R-$, называется аппроксимативной единицей, если
(i) $K_r\ge0$; (ii) $\int_{-\pi}^{\pi}K_r=1$; (iii) для любого $\delta\in(0,\pi)$: $\max_{\delta\le|\psi|\le\pi}K_r(\psi)\to0$ при $r\to R-$.
\end{defs}

\subsection{Сходимость ряда Фурье гладкой функции}

На слайде 351 используется без доказательства, что для гладких граничных данных ряд Фурье функции $H_0$ сходится \emph{к самой} $H_0$.
Здесь мы это доказываем.

\begin{thmbox}{Лемма 24.1 (убывание коэффициентов Фурье функции класса $C^{1+1}$)}
Пусть $H\in C^{1+1}([-\pi,\pi])$, $H(-\pi)=H(\pi)$, $H'(-\pi)=H'(\pi)$, и $\alpha_n,\beta_n$ --- коэффициенты Фурье $H$. Тогда при $n\ge1$
\[
\alpha_n=-\frac1{\pi n^2}\int_{-\pi}^{\pi}H''(\vartheta)\cos n\vartheta\,d\vartheta,\qquad
\beta_n=-\frac1{\pi n^2}\int_{-\pi}^{\pi}H''(\vartheta)\sin n\vartheta\,d\vartheta,
\]
и потому $|\alpha_n|+|\beta_n|\le\dfrac{M_0}{n^2}$, где $M_0=\dfrac2\pi\displaystyle\int_{-\pi}^{\pi}|H''(\vartheta)|\,d\vartheta$.
\end{thmbox}
\begin{proof}
\emph{Коэффициент $\alpha_n$.} Так как $H\in C^1([-\pi,\pi])$, интегрируем по частям на всём отрезке:
\[
\int_{-\pi}^{\pi}H\cos n\vartheta\,d\vartheta=\Bigl[\frac{H(\vartheta)\sin n\vartheta}{n}\Bigr]_{-\pi}^{\pi}-\frac1n\int_{-\pi}^{\pi}H'\sin n\vartheta\,d\vartheta
=-\frac1n\int_{-\pi}^{\pi}H'\sin n\vartheta\,d\vartheta,
\]
так как $\sin(\pm n\pi)=0$. Функция $H'$ непрерывна на $[-\pi,\pi]$ и класса $C^1$ на каждом $I_j$; интегрируем по частям на каждом $I_j$ и складываем:
\[
\int_{-\pi}^{\pi}H'\sin n\vartheta\,d\vartheta=\sum_{j=1}^J\Bigl[-\frac{H'(\vartheta)\cos n\vartheta}{n}\Bigr]_{t_{j-1}}^{t_j}+\frac1n\int_{-\pi}^{\pi}H''\cos n\vartheta\,d\vartheta .
\]
Так как $H'$ непрерывна в точках $t_1,\dots,t_{J-1}$, внутренние подстановки взаимно уничтожаются (телескопическая сумма), и остаётся
$-\frac1n\bigl(H'(\pi)\cos n\pi-H'(-\pi)\cos n\pi\bigr)=0$ в силу $H'(\pi)=H'(-\pi)$. Итак,
$\pi\alpha_n=-\frac1{n^2}\int H''\cos n\vartheta\,d\vartheta$.

\emph{Коэффициент $\beta_n$.} Аналогично, но теперь используется условие $H(\pi)=H(-\pi)$:
\[
\int_{-\pi}^{\pi}H\sin n\vartheta\,d\vartheta=-\frac{\cos n\pi}{n}\bigl(H(\pi)-H(-\pi)\bigr)+\frac1n\int_{-\pi}^{\pi}H'\cos n\vartheta\,d\vartheta=\frac1n\int_{-\pi}^{\pi}H'\cos n\vartheta\,d\vartheta,
\]
\[
\int_{-\pi}^{\pi}H'\cos n\vartheta\,d\vartheta=\sum_{j=1}^J\Bigl[\frac{H'(\vartheta)\sin n\vartheta}{n}\Bigr]_{t_{j-1}}^{t_j}-\frac1n\int_{-\pi}^{\pi}H''\sin n\vartheta\,d\vartheta
=-\frac1n\int_{-\pi}^{\pi}H''\sin n\vartheta\,d\vartheta
\]
(телескопическая сумма равна \(\frac1n\bigl(H'(\pi)\sin n\pi-H'(-\pi)\sin(-n\pi)\bigr)=0\)).\\ Значит,
$\pi\beta_n=-\frac1{n^2}\int H''\sin n\vartheta\,d\vartheta$.

\emph{Оценка.} $|\alpha_n|\le\frac1{\pi n^2}\int|H''|$, $|\beta_n|\le\frac1{\pi n^2}\int|H''|$; складываем.
\end{proof}

\begin{thmbox}{Лемма 24.2 (ядро Дирихле)}
Пусть $F$ --- $2\pi$-периодическая непрерывная функция, $S_N$ --- частичные суммы её ряда Фурье. Тогда
\[
S_N(\varphi)=\frac1\pi\int_{-\pi}^{\pi}F(\varphi+t)\,D_N(t)\,dt,\qquad
D_N(t)=\frac12+\sum_{n=1}^N\cos nt=\frac{\sin\bigl((N+\frac12)t\bigr)}{2\sin\frac t2}\ \ (t\notin2\pi\ZZ),
\]
и $\frac1\pi\int_{-\pi}^{\pi}D_N(t)\,dt=1$.
\end{thmbox}
\begin{proof}
Подставим формулы для коэффициентов и воспользуемся тождеством
\[\cos n\vartheta\cos n\varphi+\sin n\vartheta\sin n\varphi=\cos n(\vartheta-\varphi):\]
\[
S_N(\varphi)=\frac1\pi\int_{-\pi}^{\pi}F(\vartheta)\Bigl[\frac12+\sum_{n=1}^N\cos n(\vartheta-\varphi)\Bigr]d\vartheta=\frac1\pi\int_{-\pi}^{\pi}F(\vartheta)D_N(\vartheta-\varphi)\,d\vartheta .
\]
Замена $\vartheta=\varphi+t$ даёт интеграл по отрезку $[-\pi-\varphi,\pi-\varphi]$ от $2\pi$-периодической функции $t\mapsto F(\varphi+t)D_N(t)$;
интеграл периодической функции по любому отрезку длины периода один и тот же, так что можно интегрировать по $[-\pi,\pi]$.
Далее, $\int_{-\pi}^{\pi}\cos nt\,dt=0$ при $n\ge1$, откуда $\frac1\pi\int D_N=\frac1\pi\cdot\frac12\cdot2\pi=1$.
Наконец, по формуле $2\sin\frac t2\cos nt=\sin(n+\frac12)t-\sin(n-\frac12)t$
\[
2\sin\tfrac t2\,D_N(t)=\sin\tfrac t2+\sum_{n=1}^N\Bigl(\sin(n+\tfrac12)t-\sin(n-\tfrac12)t\Bigr)=\sin(N+\tfrac12)t .\qedhere
\]
\end{proof}

\begin{thmbox}{Лемма 24.3 (Риман--Лебег для непрерывных функций)}
Если $g\in C([-\pi,\pi])$, то $\displaystyle\int_{-\pi}^{\pi}g(t)\sin\lambda t\,dt\to0$ при $\lambda\to+\infty$.
\end{thmbox}
\begin{proof}
Пусть $\lambda>1$, $\eta=\pi/\lambda<\pi$, $I(\lambda)=\int_{-\pi}^{\pi}g(t)\sin\lambda t\,dt$, $\|g\|=\max|g|$, $\omega_g$ --- модуль непрерывности $g$ на $[-\pi,\pi]$.
Сделаем замену $t=s+\eta$; так как $\sin(\lambda s+\pi)=-\sin\lambda s$,
\[
I(\lambda)=-\int_{-\pi-\eta}^{\pi-\eta}g(s+\eta)\sin\lambda s\,ds=-\int_{-\pi}^{\pi-\eta}g(s+\eta)\sin\lambda s\,ds+\rho_1,\qquad |\rho_1|\le\|g\|\eta .
\]
С другой стороны, $I(\lambda)=\int_{-\pi}^{\pi-\eta}g(s)\sin\lambda s\,ds+\rho_2$, $|\rho_2|\le\|g\|\eta$. Складывая,
\[
2|I(\lambda)|\le\int_{-\pi}^{\pi-\eta}|g(s)-g(s+\eta)|\,ds+2\|g\|\eta\le2\pi\,\omega_g(\eta)+2\|g\|\eta\xrightarrow[\lambda\to\infty]{}0,
\]
так как $\eta=\pi/\lambda\to0$, а $g$ равномерно непрерывна на отрезке.
\end{proof}

\begin{thmbox}{Теорема 24.4 (сходимость ряда Фурье гладкой функции)}
Пусть $F\colon\RR\to\RR$ --- $2\pi$-периодическая функция класса $C^1(\RR)$. Тогда её ряд Фурье сходится к $F(\varphi)$ в каждой точке $\varphi\in\RR$.
Если, кроме того, $\sum_n(|\alpha_n|+|\beta_n|)<\infty$, то сходимость равномерная на $\RR$.
\end{thmbox}
\begin{proof}
Зафиксируем $\varphi$. По лемме 24.2, используя $\frac1\pi\int D_N=1$,
\[
S_N(\varphi)-F(\varphi)=\frac1\pi\int_{-\pi}^{\pi}\bigl(F(\varphi+t)-F(\varphi)\bigr)D_N(t)\,dt=\frac1{2\pi}\int_{-\pi}^{\pi}g(t)\sin\bigl((N+\tfrac12)t\bigr)\,dt,
\]
где $g(t)=\dfrac{F(\varphi+t)-F(\varphi)}{\sin(t/2)}$ при $t\in[-\pi,\pi]\setminus\{0\}$ и $g(0)=2F'(\varphi)$. Функция $g$ непрерывна на $[-\pi,\pi]$:
вне нуля --- как частное непрерывных с ненулевым знаменателем ($\sin\frac t2\ne0$ при $0<|t|\le\pi$), а в нуле
$g(t)=\frac{F(\varphi+t)-F(\varphi)}{t}\cdot\frac{t}{\sin(t/2)}\to F'(\varphi)\cdot2$. (В равенстве под интегралом точка $t=0$ роли не играет.) По лемме 24.3 с
$\lambda=N+\frac12$ получаем $S_N(\varphi)\to F(\varphi)$.

Если $\sum(|\alpha_n|+|\beta_n|)<\infty$, то по признаку Вейерштрасса ряд Фурье сходится равномерно на $\RR$ к некоторой функции $S$; по доказанному
$S(\varphi)=\lim S_N(\varphi)=F(\varphi)$ при каждом $\varphi$, т.\,е. $S_N\rightrightarrows F$.
\end{proof}

\begin{remarkbox}{Применение к $C^{1+1}$}
Если $H_0\in C^{1+1}([-\pi,\pi])$, $H_0(-\pi)=H_0(\pi)$, $H'_0(-\pi)=H'_0(\pi)$, то $2\pi$-периодическое продолжение $F$ функции $H_0$ принадлежит $C^1(\RR)$:
в точках $\pi+2\pi k$ левая производная равна $H_0'(\pi)$, правая --- $H_0'(-\pi)$ (односторонние производные $C^1$-функции на отрезке),
они совпадают, и производная непрерывна. По лемме 24.1 $\sum(|\alpha_n|+|\beta_n|)\le|\alpha_0|+M_0\sum n^{-2}<\infty$, и по теореме 24.4
ряд Фурье $H_0$ сходится к $H_0$ равномерно.
\end{remarkbox}

\subsection{Метод Фурье (предложение 14.8)}

\begin{thmbox}{Лемма 24.5 (упражнение со слайда 351: гладкая аппроксимация граничных данных)}
Пусть $F$ --- непрерывная $2\pi$-периодическая функция на $\RR$. Тогда существуют $2\pi$-периодические функции $F_m\in C^2(\RR)$,
$m\in\NN$, такие что $\|F_m-F\|_{\RR}\to0$. В частности, для любой $h_0\in C(C_R)$ найдутся $h_{0,m}\in C(C_R)$ с
$H_{0,m}(\varphi)=h_{0,m}(Re^{i\varphi})\in C^{1+1}([-\pi,\pi])$ (даже $C^2$), $H_{0,m}(-\pi)=H_{0,m}(\pi)$, $H'_{0,m}(-\pi)=H'_{0,m}(\pi)$ и $h_{0,m}\rightrightarrows h_0$ на $C_R$.
\end{thmbox}
\begin{proof}
Используем \emph{усреднение Стеклова}: для $\delta>0$ и непрерывной периодической $G$ положим
\[
(T_\delta G)(\varphi)=\frac1\delta\int_\varphi^{\varphi+\delta}G(t)\,dt=\frac1\delta\int_0^\delta G(\varphi+s)\,ds .
\]
Свойства (для непрерывной $2\pi$-периодической $G$):
\begin{enumerate}
\item[(а)] $T_\delta G$ $2\pi$-периодична (подынтегральная функция во второй записи периодична по $\varphi$).
\item[(б)] $T_\delta G\in C^1(\RR)$ и $(T_\delta G)'(\varphi)=\frac{G(\varphi+\delta)-G(\varphi)}\delta$ --- по теореме Барроу (производная интеграла с переменным пределом от
непрерывной функции). Если же $G\in C^1(\RR)$, то эта производная --- функция класса $C^1$, т.\,е. $T_\delta G\in C^2(\RR)$.
\item[(в)] $\|T_\delta G-G\|_\RR\le\omega_G(\delta)$: $|(T_\delta G)(\varphi)-G(\varphi)|=\bigl|\frac1\delta\int_0^\delta(G(\varphi+s)-G(\varphi))\,ds\bigr|\le\omega_G(\delta)$.
\item[(г)] $\omega_{T_\delta G}(\eta)\le\omega_G(\eta)$: при $|\varphi-\psi|\le\eta$ имеем
$|(T_\delta G)(\varphi)-(T_\delta G)(\psi)|\le\frac1\delta\int_0^\delta|G(\varphi+s)-G(\psi+s)|\,ds\le\omega_G(\eta)$.
\end{enumerate}
Положим $F_m=T_{1/m}T_{1/m}F$. По (а), (б) $T_{1/m}F$ периодична и класса $C^1$, а $F_m$ периодична и класса $C^2$. По (в), (г)
\[
\|F_m-F\|\le\|T_{1/m}(T_{1/m}F)-T_{1/m}F\|+\|T_{1/m}F-F\|\le\omega_{T_{1/m}F}(\tfrac1m)+\omega_F(\tfrac1m)\le2\,\omega_F(\tfrac1m)\to0,
\]
так как непрерывная периодическая $F$ равномерно непрерывна на $\RR$ (она равномерно непрерывна на компакте $[-2\pi,2\pi]$, а любые две точки
на расстоянии $\le\pi$ сдвигом на кратное $2\pi$ переводятся в этот отрезок).

Для второго утверждения возьмём $F$ --- периодическое продолжение $H_0$, и $h_{0,m}(Re^{i\varphi})=F_m(\varphi)$ (корректно в силу периодичности).
Тогда $H_{0,m}=F_m|_{[-\pi,\pi]}\in C^2([-\pi,\pi])\subset C^{1+1}([-\pi,\pi])$ (разбиение из одного отрезка), а равенства на концах --- следствие
периодичности $F_m$ и $F_m'$. Наконец, $\|h_{0,m}-h_0\|_{C_R}=\|F_m-F\|_\RR\to0$.
\end{proof}

\begin{thmbox}{Предложение 14.8 (метод Фурье)}
В круге $B_R$ всякая задача Дирихле разрешима, т.\,е. для любой $h_0\in C(C_R)$ найдётся $h\in\pIVHarm(B_R)\cap C(\clos{B_R})$ с $h|_{C_R}=h_0$. При
$z=re^{i\varphi}$
\[
h(z)=\begin{cases} h_0(Re^{i\varphi}), & r=R,\\[1mm]
\displaystyle\frac{A_0}2+\sum_{n=1}^\infty r^n\bigl(A_n\cos n\varphi+B_n\sin n\varphi\bigr), & r<R,\end{cases}
\]
где $A_n,B_n$ определяются формулами (14.2), (14.3), в которых $h$ заменено на $h_0$.
\end{thmbox}

\begin{remarkbox}{Откуда берётся формула (эвристика)}
Если решение $h$ существует и гармонично в чуть большем круге, то по предложению 14.5 (билет 23) оно обязано иметь вид (14.1) с коэффициентами,
вычисленными по его значениям на $C_R$, т.\,е. по $h_0$. Метод Фурье --- взять эту формулу за \emph{определение} $h$ и проверить, что она даёт решение.
При $r<R$ ряд всегда сходится: $|r^n(A_n\cos n\varphi+B_n\sin n\varphi)|\le4\|h_0\|_{C_R}q^n$. Трудность --- непрерывность вплоть до границы.
\end{remarkbox}

\begin{proof}
\emph{Шаг 1: гладкие данные.} Пусть $H_0\in C^{1+1}([-\pi,\pi])$, $H_0(-\pi)=H_0(\pi)$, $H'_0(-\pi)=H'_0(\pi)$. Так как $A_n=\alpha_n/R^n$, $B_n=\beta_n/R^n$
($\alpha_n,\beta_n$ --- коэффициенты Фурье $H_0$), по лемме 24.1 при $n\ge1$
\[
|A_n|+|B_n|\le\frac{M_0}{n^2R^n},\qquad M_0=\frac2\pi\int_{-\pi}^{\pi}|H''_0(\vartheta)|\,d\vartheta
\]
(на слайде 350 это получено двукратным интегрированием по частям; детали --- в лемме 24.1). Члены ряда
\[
u_n(z)=r^n\bigl(A_n\cos n\varphi+B_n\sin n\varphi\bigr)=A_n\Re z^n+B_n\Im z^n,\qquad u_0=\frac{A_0}2,
\]
--- однородные гармонические полиномы, непрерывные на $\CC$ и гармонические в $\CC$. На $\clos{B_R}$ имеем $|u_n|\le R^n(|A_n|+|B_n|)\le M_0/n^2$,
и по признаку Вейерштрасса ряд $\sum u_n$ сходится абсолютно и равномерно на $\clos{B_R}$ к некоторой функции $S$. Тогда:
\begin{itemize}
\item $S\in C(\clos{B_R})$ как равномерный предел непрерывных частичных сумм;
\item $S\in\pIVHarm(B_R)$ по аналогу теоремы Вейерштрасса (предложение 14.7): частичные суммы гармоничны в $B_R$ и сходятся равномерно
(тем более локально равномерно);
\item при $r=R$: $S(Re^{i\varphi})=\frac{\alpha_0}2+\sum_n(\alpha_n\cos n\varphi+\beta_n\sin n\varphi)$ --- это ряд Фурье функции $H_0$, и по теореме 24.4 (см. замечание
<<Применение к $C^{1+1}$>>) его сумма равна $H_0(\varphi)=h_0(Re^{i\varphi})$.
\end{itemize}
Итак, функция $h$ из формулировки совпадает с $S$ и является решением задачи Дирихле.

\begin{detail}{почему сумма ряда на границе равна $h_0$}
На слайде 351 сказано лишь, что $h$ непрерывна и гармонична. То, что $h|_{C_R}=h_0$, т.\,е. что ряд Фурье $H_0$ сходится \emph{к $H_0$}, а не к
какой-то другой функции, --- содержательный факт; он доказан в теореме 24.4 (ядро Дирихле + лемма Римана--Лебега). Абсолютная сходимость
из леммы 24.1 сама по себе даёт лишь существование суммы.
\end{detail}

\emph{Шаг 2: общий случай.} Пусть $h_0\in C(C_R)$ произвольна. По лемме 24.5 найдём $h_{0,m}$ с гладкими $H_{0,m}$, удовлетворяющими условиям шага 1, и
$\|h_{0,m}-h_0\|_{C_R}\to0$. По шагу 1 существуют решения $h_m\in\pIVHarm(B_R)\cap C(\clos{B_R})$, $h_m|_{C_R}=h_{0,m}$, задаваемые при $r<R$ рядами
с коэффициентами $A_n^{(m)},B_n^{(m)}$ функции $h_{0,m}$.

\emph{Равномерная фундаментальность.} Для любых $m,k$ функция $h_m-h_k\in\pIVHarm(B_R)\cap C(\clos{B_R})$ равна $h_{0,m}-h_{0,k}$ на $C_R$. По принципу
максимума и минимума (теорема 14.1), применённому к $h_m-h_k$:
\[
\|h_m-h_k\|_{\clos{B_R}}\le\|h_{0,m}-h_{0,k}\|_{C_R}\le\|h_{0,m}-h_0\|_{C_R}+\|h_0-h_{0,k}\|_{C_R}\xrightarrow[m,k\to\infty]{}0 .
\]
По критерию Коши $h_m\rightrightarrows h$ на $\clos{B_R}$ для некоторой функции $h$. Тогда $h\in C(\clos{B_R})$ (равномерный предел непрерывных),
$h\in\pIVHarm(B_R)$ (предложение 14.7) и $h|_{C_R}=\lim h_{0,m}=h_0$. Решение найдено.

\begin{detail}{решение задаётся рядом из формулировки и при $r<R$}
На слайде это утверждается без проверки. Пусть $|z|=r<R$, $q=r/R$, $\eps_m=\|h_{0,m}-h_0\|_{C_R}$. Из (14.2), (14.3) для разности данных
\[
\tfrac12|A_0^{(m)}-A_0|\le\frac1{2\pi}\int_{-\pi}^{\pi}|H_{0,m}-H_0|\,d\vartheta\le\eps_m,\qquad |A_n^{(m)}-A_n|,\ |B_n^{(m)}-B_n|\le\frac{2\eps_m}{R^n}\quad(n\ge1).
\]
Ряд $\Sigma(z)=\frac{A_0}2+\sum_n r^n(A_n\cos n\varphi+B_n\sin n\varphi)$ сходится абсолютно (см. эвристику выше), и, вычитая из него ряд для $h_m(z)$ почленно,
\[
|h_m(z)-\Sigma(z)|\le\eps_m\Bigl(1+\sum_{n\ge1}4q^n\Bigr)=\eps_m\,\frac{1+3q}{1-q}\xrightarrow[m\to\infty]{}0 .
\]
Поскольку $h_m(z)\to h(z)$, получаем $h(z)=\Sigma(z)$.
\end{detail}
Предложение доказано.
\end{proof}

\begin{remarkbox}{Следствия}
(1) Вместе с методом конформного отображения (билет 22, следствие 22.3) это доказывает, что в любой жордановой области любая задача Дирихле разрешима.
(2) Решение задачи Дирихле в круге зависит от данных устойчиво: $\|h-\tilde h\|_{\clos{B_R}}=\|h_0-\tilde h_0\|_{C_R}$ (теорема 14.1).
\end{remarkbox}

\subsection{Ядро Пуассона и его свойства}

\begin{thmbox}{Лемма 24.6 (суммирование ряда; свойства ядра Пуассона)}
Пусть $\zeta=Re^{i\vartheta}\in C_R$, $z=re^{i\varphi}\in B_R$, $q=r/R$, $\psi=\vartheta-\varphi$. Тогда
\begin{enumerate}
\item[(а)] $|\zeta-z|^2=R^2-2Rr\cos\psi+r^2=(R-r)^2+4Rr\sin^2\frac\psi2$;
\item[(б)] $\displaystyle\frac12+\sum_{n=1}^\infty q^n\cos n\psi=\frac12\Re\frac{\zeta+z}{\zeta-z}=\frac12\cdot\frac{R^2-r^2}{R^2-2Rr\cos\psi+r^2}=\pi P_r(\psi)$,
причём ряд сходится равномерно по $\psi\in\RR$ при фиксированном $r<R$;
\item[(в)] $P_r(\psi)>0$; \quad (г) $\displaystyle\int_{-\pi}^{\pi}P_r(\psi)\,d\psi=1$, т.\,е. $\displaystyle\int_{C_R}P(z,\zeta)\,|d\zeta|=1$;
\item[(д)] при $0<\delta\le\pi$ и $\delta\le|\psi|\le\pi$: $P_r(\psi)\le\dfrac{R^2-r^2}{8\pi Rr\sin^2(\delta/2)}\to0$ при $r\to R-$;
\item[(е)] при фиксированном $\zeta\in C_R$ функция $z\mapsto P(z,\zeta)$ гармонична в $B_R$: $2\pi R\,P(z,\zeta)=\Re\dfrac{\zeta+z}{\zeta-z}$.
\end{enumerate}
Таким образом, $\{P_r\}_{r\to R-}$ --- аппроксимативная единица (свойства (в), (г), (д)).
\end{thmbox}
\begin{proof}
(а) Теорема косинусов в треугольнике с вершинами $0,\zeta,z$ (угол при вершине $0$ равен $\psi$ по модулю $2\pi$), или прямо:
$|\zeta-z|^2=|\zeta|^2-2\Re(\zeta\bar z)+|z|^2=R^2-2Rr\cos(\vartheta-\varphi)+r^2$. Далее $-2Rr\cos\psi=-2Rr+4Rr\sin^2\frac\psi2$.

(б) Положим $w=z/\zeta$, $|w|=q<1$, $w^n=q^ne^{-in\psi}$, $\Re w^n=q^n\cos n\psi$. Геометрическая прогрессия:
$\sum_{n\ge1}w^n=\frac{w}{1-w}=\frac{z}{\zeta-z}$, поэтому
\[
\frac12+\sum_{n=1}^\infty q^n\cos n\psi=\Re\Bigl(\frac12+\frac{z}{\zeta-z}\Bigr)=\frac12\Re\frac{\zeta+z}{\zeta-z}
=\frac12\Re\frac{(\zeta+z)(\bar\zeta-\bar z)}{|\zeta-z|^2}=\frac12\cdot\frac{|\zeta|^2-|z|^2}{|\zeta-z|^2},
\]
так как $(\zeta+z)(\bar\zeta-\bar z)=|\zeta|^2-|z|^2+(z\bar\zeta-\zeta\bar z)$, а $z\bar\zeta-\zeta\bar z=2i\Im(z\bar\zeta)$ --- чисто мнимое. С учётом (а) получаем (б).
Равномерная сходимость: $|q^n\cos n\psi|\le q^n$, $\sum q^n<\infty$ (признак Вейерштрасса).

(в) $R^2-r^2>0$, а знаменатель $=|\zeta-z|^2>0$, так как $z\ne\zeta$.

(г) Интегрируя равномерно сходящийся ряд (б) почленно и учитывая $\int_{-\pi}^{\pi}\cos n\psi\,d\psi=0$ ($n\ge1$):
$\pi\int_{-\pi}^{\pi}P_r=\frac12\cdot2\pi$, т.\,е. $\int P_r=1$. В записи по окружности: $P(z,\zeta)|d\zeta|=\frac{1}{2\pi R}\frac{R^2-r^2}{|\zeta-z|^2}R\,d\vartheta=P_r(\vartheta-\varphi)\,d\vartheta$,
а интеграл по $\vartheta\in[-\pi,\pi]$ от $2\pi$-периодической $P_r(\vartheta-\varphi)$ равен $\int_{-\pi}^{\pi}P_r$.

(д) При $\delta\le|\psi|\le\pi$ имеем $\sin^2\frac\psi2\ge\sin^2\frac\delta2$ ($\sin$ возрастает на $[0,\frac\pi2]$), и по (а) знаменатель $\ge4Rr\sin^2\frac\delta2$;
остаётся заметить, что $R^2-r^2\to0$ при $r\to R-$, а $r\ge R/2$ для $r$, близких к $R$.

(е) Функция $z\mapsto\frac{\zeta+z}{\zeta-z}$ голоморфна в $B_R$ (знаменатель не обращается в нуль), её вещественная часть по вычислению в (б) равна
$\frac{R^2-|z|^2}{|\zeta-z|^2}$; вещественная часть голоморфной функции гармонична (билет 21).
\end{proof}

\begin{remarkbox}{Смысл свойств}
$P_r$ --- чётная функция, наибольшая в $\psi=0$: $P_r(0)=\frac1{2\pi}\frac{R+r}{R-r}\to+\infty$, а вне любой окрестности нуля $P_r\to0$ равномерно,
при этом полная масса равна $1$. Поэтому при $z\to\zeta_0\in C_R$ <<масса>> ядра $P(z,\cdot)$ сосредотачивается около $\zeta_0$, и
$\int P(z,\zeta)h_0(\zeta)|d\zeta|\to h_0(\zeta_0)$ --- это и доказывается в теореме 24.7.
\end{remarkbox}

\subsection{Формула Пуассона}

\begin{thmbox}{Предложение 14.9 (формула Пуассона в круге)}
Решение задачи Дирихле в круге $B_R$ с граничной функцией $h_0\in C(C_R)$ имеет вид
\[
h(z)=\begin{cases}h_0(z),&|z|=R,\\[1mm]
\displaystyle\frac1{2\pi R}\int_{C_R}\frac{R^2-|z|^2}{|\zeta-z|^2}\,h_0(\zeta)\,|d\zeta|
=\frac1{2\pi}\int_{-\pi}^{\pi}\frac{(R^2-r^2)\,h_0(Re^{i\vartheta})}{R^2-2Rr\cos(\vartheta-\varphi)+r^2}\,d\vartheta,&|z|=r<R .\end{cases}
\]
\end{thmbox}
\begin{proof}
Равенство двух интегралов --- подстановка $\zeta=Re^{i\vartheta}$, $|d\zeta|=R\,d\vartheta$ и лемма 24.6(а). Пусть $h$ --- решение (оно существует и единственно,
предложение 14.8), $0<r<R$, $q=r/R$. По предложению 14.8, подставляя (14.2), (14.3) с $h_0$:
\begin{align*}
r^n\bigl(A_n\cos n\varphi+B_n\sin n\varphi\bigr)&=\frac{q^n}\pi\int_{-\pi}^{\pi}H_0(\vartheta)\bigl(\cos n\vartheta\cos n\varphi+\sin n\vartheta\sin n\varphi\bigr)d\vartheta\\
&=\frac{1}\pi\int_{-\pi}^{\pi}H_0(\vartheta)\,q^n\cos n(\vartheta-\varphi)\,d\vartheta .
\end{align*}
Ряд $\sum_n H_0(\vartheta)q^n\cos n(\vartheta-\varphi)$ сходится равномерно по $\vartheta$ (мажоранта $\|H_0\|q^n$), поэтому сумму и интеграл можно переставить:
\[
h(re^{i\varphi})=\frac1\pi\int_{-\pi}^{\pi}H_0(\vartheta)\Bigl[\frac12+\sum_{n=1}^\infty q^n\cos n(\vartheta-\varphi)\Bigr]d\vartheta
=\int_{-\pi}^{\pi}H_0(\vartheta)\,P_r(\vartheta-\varphi)\,d\vartheta
\]
по лемме 24.6(б), что и есть требуемая формула. При $r=0$ формула даёт $h(0)=\frac1{2\pi}\int H_0=\frac{A_0}2$, что совпадает с предложением 14.8.
\end{proof}

\begin{remarkbox}{Замечание (опечатка на слайде 353)}
В последней строке вывода на слайде 353 записано $\frac1{2\pi R}\int_{-\pi}^{\pi}h_0(Re^{i\vartheta})\frac{|\zeta|^2-|z|^2}{|\zeta-z|^2}\,d|\zeta|$; правильно либо
$\frac1{2\pi R}\int_{C_R}\dots|d\zeta|$, либо $\frac1{2\pi}\int_{-\pi}^{\pi}\dots d\vartheta$ (так как $|d\zeta|=R\,d\vartheta$).
\end{remarkbox}

Предложение 14.9 опирается на метод Фурье. Дадим теперь \emph{независимое} прямое доказательство того, что интеграл Пуассона решает задачу
Дирихле, --- оно не использует ни леммы 24.1--24.5, ни предложение 14.8 (и именно оно используется в билете 22 при доказательстве теоремы 22.6).

\begin{thmbox}{Теорема 24.7 (интеграл Пуассона решает задачу Дирихле)}
Пусть $h_0\in C(C_R)$. Положим $h(z)=\mathcal P[h_0](z)=\int_{C_R}P(z,\zeta)h_0(\zeta)\,|d\zeta|$ при $z\in B_R$ и $h(z)=h_0(z)$ при $z\in C_R$.
Тогда $h\in\pIVHarm(B_R)\cap C(\clos{B_R})$, т.\,е. $h$ --- решение задачи Дирихле. Кроме того, $\min_{C_R}h_0\le h\le\max_{C_R}h_0$.
\end{thmbox}
\begin{proof}
\emph{Шаг 1: гармоничность в $B_R$.} По лемме 24.6(е), $h=\Re F$, где
\[
F(z)=\frac1{2\pi}\int_{-\pi}^{\pi}\frac{\zeta+z}{\zeta-z}\,H_0(\vartheta)\,d\vartheta,\qquad\zeta=Re^{i\vartheta}
\]
(интеграл Шварца; вещественную часть можно внести под интеграл, так как $H_0$ вещественна). Покажем, что $F\in H(B_R)$. При фиксированном $z\in B_R$
\[
\frac{\zeta+z}{\zeta-z}=1+\frac{2z/\zeta}{1-z/\zeta}=1+2\sum_{n=1}^\infty\frac{z^n}{\zeta^n},\qquad\Bigl|\frac{z^n}{\zeta^n}H_0(\vartheta)\Bigr|\le\|H_0\|\Bigl(\frac{|z|}R\Bigr)^n,
\]
ряд сходится равномерно по $\vartheta$, и почленное интегрирование даёт
\[
F(z)=c_0+\sum_{n=1}^\infty c_nz^n,\qquad c_0=\frac1{2\pi}\int_{-\pi}^{\pi}H_0\,d\vartheta,\quad c_n=\frac1{\pi R^n}\int_{-\pi}^{\pi}H_0(\vartheta)e^{-in\vartheta}\,d\vartheta .
\]
Так как $|c_n|\le2\|H_0\|/R^n$, радиус сходимости этого степенного ряда не меньше $R$ (признак Коши--Адамара), и $F\in H(B_R)$.
Следовательно, $h=\Re F\in\pIVHarm(B_R)$ (билет 21); в частности, $h$ непрерывна во всех точках $B_R$.
(Заметим: $c_n=A_n-iB_n$, так что $\Re(c_nz^n)=r^n(A_n\cos n\varphi+B_n\sin n\varphi)$ --- это снова ряд метода Фурье.)

\emph{Шаг 2: оценка через граничные значения.} При $z\in B_R$ по лемме 24.6(в),(г)
\[
h(z)=\int_{C_R}P(z,\zeta)h_0(\zeta)\,|d\zeta|\le\max_{C_R}h_0\cdot\int_{C_R}P(z,\zeta)\,|d\zeta|=\max_{C_R}h_0,
\]
аналогично $h(z)\ge\min_{C_R}h_0$. На $C_R$ это очевидно.

\emph{Шаг 3: непрерывность в точках окружности ($\eps$--$\delta$).} Фиксируем $\zeta_0\in C_R$ и $\eps>0$. Пусть $M=\max_{C_R}|h_0|$.
Так как $h_0$ непрерывна в $\zeta_0$, найдётся $\delta>0$ такое, что
\[
\zeta\in C_R,\ |\zeta-\zeta_0|<\delta\ \Longrightarrow\ |h_0(\zeta)-h_0(\zeta_0)|<\eps .
\]
Положим
\[
\eta=\min\Bigl(\frac\delta2,\ \frac{\eps\,\delta^2}{16MR+1}\Bigr)>0
\]
и докажем, что $|h(z)-h(\zeta_0)|<2\eps$ для всех $z\in\clos{B_R}$ с $|z-\zeta_0|<\eta$.

\emph{(i) $z\in C_R$.} Тогда $|z-\zeta_0|<\eta\le\delta$ и $|h(z)-h(\zeta_0)|=|h_0(z)-h_0(\zeta_0)|<\eps$.

\emph{(ii) $z\in B_R$.} По нормировке (лемма 24.6(г)) $h_0(\zeta_0)=\int_{C_R}P(z,\zeta)h_0(\zeta_0)|d\zeta|$, поэтому
\[
h(z)-h(\zeta_0)=\int_{C_R}P(z,\zeta)\bigl(h_0(\zeta)-h_0(\zeta_0)\bigr)|d\zeta|=I_1+I_2,
\]
где $I_1$ --- интеграл по дуге $\Gamma_1=\{\zeta\in C_R:|\zeta-\zeta_0|<\delta\}$, а $I_2$ --- по $\Gamma_2=C_R\setminus\Gamma_1$.
\begin{itemize}
\item На $\Gamma_1$: $|h_0(\zeta)-h_0(\zeta_0)|<\eps$ и $P\ge0$, поэтому $|I_1|\le\eps\int_{\Gamma_1}P(z,\zeta)|d\zeta|\le\eps\int_{C_R}P(z,\zeta)|d\zeta|=\eps$.
\item На $\Gamma_2$: $|\zeta-z|\ge|\zeta-\zeta_0|-|z-\zeta_0|\ge\delta-\frac\delta2=\frac\delta2$, поэтому
$P(z,\zeta)\le\frac{1}{2\pi R}\cdot\frac{R^2-|z|^2}{\delta^2/4}$. Кроме того, $|h_0(\zeta)-h_0(\zeta_0)|\le2M$ и длина $\Gamma_2$ не больше $2\pi R$. Значит,
\[
|I_2|\le2M\cdot2\pi R\cdot\frac{4(R^2-|z|^2)}{2\pi R\,\delta^2}=\frac{8M(R^2-|z|^2)}{\delta^2}.
\]
Так как $R-|z|=|\zeta_0|-|z|\le|\zeta_0-z|$ и $R+|z|\le2R$, имеем $R^2-|z|^2\le2R|z-\zeta_0|$, и
\[
|I_2|\le\frac{16MR\,|z-\zeta_0|}{\delta^2}<\frac{16MR}{\delta^2}\cdot\frac{\eps\,\delta^2}{16MR+1}<\eps .
\]
\end{itemize}
Итак, $|h(z)-h(\zeta_0)|\le|I_1|+|I_2|<2\eps$. Поскольку $\eps>0$ произвольно, $h$ непрерывна в $\zeta_0$ (как функция на $\clos{B_R}$).

Вместе с шагом 1 получаем $h\in C(\clos{B_R})$, а $h|_{C_R}=h_0$ по определению.
\end{proof}

\begin{remarkbox}{Следствия теоремы 24.7}
\begin{enumerate}
\item \emph{Второе доказательство предложения 14.8.} Теорема 24.7 даёт решение для любой непрерывной $h_0$ без аппроксимации гладкими данными,
а шаг 1 её доказательства показывает, что при $r<R$ это решение задаётся рядом метода Фурье. По единственности (предложение 14.3) решения,
построенные двумя способами, совпадают.
\item \emph{Интегральное представление.} Если $h\in\pIVHarm(B_R)\cap C(\clos{B_R})$, то $h(z)=\int_{C_R}P(z,\zeta)h(\zeta)|d\zeta|$ при $z\in B_R$: обе части ---
решения задачи Дирихле с данными $h|_{C_R}$. При $z=0$ ($P(0,\zeta)=\frac1{2\pi R}$) получаем теорему о среднем (предложение 14.2).
\item \emph{Неравенство Харнака.} Так как $R-r\le|\zeta-z|\le R+r$ при $|z|=r$, то $\frac{R-r}{R+r}\le\frac{R^2-r^2}{|\zeta-z|^2}\le\frac{R+r}{R-r}$; поэтому для
$h\ge0$ из п.~2: $\frac{R-r}{R+r}h(0)\le h(z)\le\frac{R+r}{R-r}h(0)$.
\item \emph{Равномерная сходимость по радиусам.} Тот же аргумент (с аппроксимативной единицей $P_r$) показывает, что
$\int_{-\pi}^{\pi}P_r(\vartheta-\varphi)H_0(\vartheta)d\vartheta\to H_0(\varphi)$ при $r\to R-$ \emph{равномерно} по $\varphi$; в шаге 3 это заложено в том, что $\eta$ зависит
от $\zeta_0$ только через $\delta$, а $\delta$ можно выбрать общим для всех $\zeta_0$ по равномерной непрерывности $h_0$.
\end{enumerate}
\end{remarkbox}

\begin{remarkbox}{Коротко: план ответа на экзамене}
\begin{enumerate}
\item Постановка в круге; $H_0(\varphi)=h_0(Re^{i\varphi})$; эвристика: решение обязано иметь вид (14.1) с коэффициентами Фурье данных (предл. 14.5).
\item Предл. 14.8, гладкий случай: дважды по частям $|A_n|+|B_n|\le M_0n^{-2}R^{-n}$ $\Rightarrow$ равномерная сходимость на $\clos{B_R}$ $\Rightarrow$ непрерывность и
гармоничность (предл. 14.7); на границе ряд Фурье сходится к $H_0$ (ядро Дирихле, Риман--Лебег).
\item Общий случай: гладкая аппроксимация (усреднение Стеклова дважды), принцип максимума $\|h_m-h_k\|\le\|h_{0,m}-h_{0,k}\|$, предел --- решение.
\item Предл. 14.9: подставить коэффициенты, переставить сумму и интеграл, просуммировать $\frac12+\sum q^n\cos n\psi=\frac12\Re\frac{\zeta+z}{\zeta-z}$.
\item Ядро Пуассона: $P>0$, $\int P=1$, $P\to0$ вне окрестности $\zeta_0$ (аппроксимативная единица); $2\pi RP=\Re\frac{\zeta+z}{\zeta-z}$ гармонично по $z$.
\item Теорема 24.7: $\eps$--$\delta$: разбить окружность на дугу около $\zeta_0$ (вклад $<\eps$) и остальное (вклад $\le16MR|z-\zeta_0|/\delta^2$).
\item Следствия: представление Пуассона, теорема о среднем, неравенство Харнака; с конформными отображениями --- разрешимость в жордановых областях.
\end{enumerate}
\end{remarkbox}

% ===== part 5a =====
% ====================== Часть 5a: билеты 25--27 (аналитическое продолжение) ======================

\ticket{25}{Элементы и их аналитическое продолжение. Аналитическое продолжение по цепочке и вдоль пути, их основные свойства}
\label{t25:start}

\begin{defs}
\textbf{Обозначения и базовые факты прошлого семестра.}
$B(a,R)=\{z:|z-a|<R\}$ при $0<R\le+\infty$ (при $R=+\infty$ это $\CC$);
$V(a,\delta)=B(a,\delta)\setminus\{a\}$ --- проколотая окрестность, $V(\infty,\delta)=\{|z|>1/\delta\}$.
\emph{Область} --- непустое открытое связное множество. $H(D)$ --- множество голоморфных в открытом $D$ функций.
\emph{Путь} --- непрерывное отображение $\gamma\colon[\alpha,\beta]\to\CC$; $\gamma^*=\gamma([\alpha,\beta])$ --- его носитель;
обычно $[\alpha,\beta]=I:=[0,1]$. $\|\gamma-\sigma\|_I=\max_{t\in I}|\gamma(t)-\sigma(t)|$.
Пути $\gamma\sigma$ (сначала $\gamma$, затем $\sigma$; определён, если $\gamma(1)=\sigma(0)$) и $\gamma^-(t)=\gamma(1-t)$ --- произведение и обратный путь.

\smallskip
\textbf{Теорема единственности.} Если $D$ --- область, $f,g\in H(D)$ и $f=g$ на множестве, имеющем предельную точку в $D$
(в частности, на непустом открытом подмножестве), то $f\equiv g$ в $D$.

\textbf{Теорема Коши--Тейлора.} Если $f\in H(B(a,R))$, то ряд Тейлора $\sum_{n\ge0}\frac{f^{(n)}(a)}{n!}(z-a)^n$ сходится к $f$
во всём круге $B(a,R)$. Следовательно, радиус сходимости ряда Тейлора функции $f$ в точке $a$ равен
\emph{точной верхней грани} тех $r$, для которых $f$ (точнее, её росток в $a$) голоморфно продолжается в $B(a,r)$.

\textbf{Пересечение кругов} выпукло, а значит связно. Поэтому две функции, голоморфные в кругах $B_0$ и $B_1$ и совпадающие на
непустом открытом подмножестве $B_0\cap B_1$, совпадают на всём $B_0\cap B_1$ (теорема единственности в области $B_0\cap B_1$).
Этот простой факт используется ниже постоянно.

\smallskip
\textbf{Определение 15.1.} Пусть $D,G$ --- области, $D\cap G$ непусто и связно. Функция $g\in H(G)$ называется
\emph{непосредственным аналитическим продолжением} $f\in H(D)$, если $f=g$ на $D\cap G$.
Функция $g\in H(G)$ --- \emph{аналитическое продолжение $f\in H(D)$ по цепочке областей}, если существуют области
$\Omega_1=D,\Omega_2,\dots,\Omega_N=G$ с непустыми связными $\Omega_n\cap\Omega_{n+1}$ и функции $f_n\in H(\Omega_n)$,
$f_1=f$, $f_N=g$, $f_n=f_{n+1}$ на $\Omega_n\cap\Omega_{n+1}$ ($n=1,\dots,N-1$). Наиболее удобен случай, когда все $\Omega_n$ --- круги
(так это понятие и вводил Вейерштрасс).

\textbf{Определение 15.2.} \emph{(Аналитический) элемент} --- пара $F=(B,f)$, где $B=B(a,R)$, $0<R\le\infty$, $f\in H(B)$;
$a$ --- \emph{центр}, $R$ --- \emph{радиус} элемента. Элемент \emph{канонический}, если $B$ совпадает с кругом сходимости
ряда Тейлора функции $f$ в точке $a$ (т.е. $R$ равно радиусу сходимости этого ряда).

\emph{Обозначение.} Если $f$ голоморфна в точке $a$ (т.е. в некоторой её окрестности), то $E_a(f)$ --- канонический элемент с центром $a$,
определяемый $f$: его круг --- круг сходимости $B(a,R)$ ряда Тейлора $f$ в $a$, функция --- сумма этого ряда.

\emph{Пример.} $(B(1,1),\ln z)$, где $\ln z=\ln|z|+i\arg z$, $\arg z\in(-\pi,\pi)$ --- главная ветвь логарифма в $\CC\setminus\RR_\pi$
($\RR_\alpha=\{te^{i\alpha}:t\ge0\}$). Так как $(\ln z)'=1/z$, то
$\ln z=\sum_{n\ge1}\frac{(-1)^{n-1}}{n}(z-1)^n$ при $|z-1|<1$; радиус сходимости ряда равен $1$, поэтому элемент канонический.

\textbf{Определение 15.3.} Элемент $F_1=(B_1,f_1)$ --- \emph{непосредственное аналитическое продолжение} (НАП) элемента
$F_0=(B_0,f_0)$, если $B_0\cap B_1\neq\varnothing$ и $f_0=f_1$ на $B_0\cap B_1$. (Отношение НАП симметрично и рефлексивно, но не транзитивно.)
Элемент $F_1$ --- \emph{аналитическое продолжение $F_0$ по цепочке}, если существуют элементы $\Phi_1=F_0,\Phi_2,\dots,\Phi_N=F_1$,
в которых $\Phi_{n+1}$ --- НАП $\Phi_n$ ($n=1,\dots,N-1$).

\textbf{Определение 15.4.} Пусть $\gamma\colon I\to\CC$ --- путь, $F_t=(B_t,f_t)$, $B_t=B(a_t,R_t)$, $t\in I$, --- семейство
\emph{канонических} элементов. Оно называется \emph{аналитическим продолжением элемента $F_0$ вдоль пути $\gamma$}, если
\begin{enumerate}
\item $a_t=\gamma(t)$ и $R_t>0$ при всех $t\in I$;
\item для каждого $t\in I$ существует связная (т.е. являющаяся промежутком) окрестность $U_t$ точки $t$ в $I$ такая, что
для всех $\xi\in U_t$: $\gamma(\xi)\in B_t$ и $F_\xi$ --- НАП $F_t$.
\end{enumerate}
Элемент $F_1$ называется \emph{результатом продолжения} $F_0$ вдоль $\gamma$; говорят, что $F_0$ \emph{продолжается вдоль $\gamma$}.
Для пути на $[\alpha,\beta]$ определение то же самое (с заменой $I$ на $[\alpha,\beta]$).
\end{defs}

\subsection{Цепочки элементов: пример неоднозначности}

Результат продолжения по цепочке, вообще говоря, зависит от цепочки --- даже если все элементы канонические.

\begin{remarkbox}{Пример (слайд 357): $\sqrt z$ после обхода нуля меняет знак}
Пусть $F_0=(B(1,1),\sqrt z)$, $\sqrt z=\sqrt{|z|}\,e^{i\arg z/2}$, $\arg z\in(-\pi,\pi)$. Рассмотрим элементы
\begin{align*}
&\Phi_1=F_0,\quad \Phi_2=\bigl(B(i,1),\ \sqrt{|z|}e^{i\arg^*z/2}\bigr),\ \arg^*z\in(-\tfrac\pi2,\tfrac{3\pi}2);\\
&\Phi_3=\bigl(B(-1,1),\ \sqrt{|z|}e^{i\arg^{***}z/2}\bigr),\ \arg^{***}z\in(0,2\pi);\\
&\Phi_4=\bigl(B(-i,1),\ \sqrt{|z|}e^{i\arg^{**}z/2}\bigr),\ \arg^{**}z\in(\tfrac\pi2,\tfrac{5\pi}2);\quad
\Phi_5=\bigl(B(1,1),\ -\sqrt{|z|}e^{i\arg z/2}\bigr).
\end{align*}
\end{remarkbox}

\begin{detail}{$\Phi_1,\dots,\Phi_5$ --- цепочка канонических элементов, $\Phi_5\ne\Phi_1$}
Каждая функция голоморфна в своём круге: например, $\arg^*$ --- непрерывная ветвь аргумента в $\CC\setminus\RR_{-\pi/2}\supset B(i,1)$.
\emph{НАП.} Если $|z-1|<1$ и $|z-i|<1$, то $\Re z>0$ и $\Im z>0$, т.е. $B(1,1)\cap B(i,1)$ лежит в открытом первом квадранте;
там $\arg z\in(0,\pi/2)$ и $\arg^*z$ --- аргументы $z$ из интервалов длины $2\pi$, оба содержащих $(0,\pi/2)$, поэтому $\arg z=\arg^*z$, и функции $\Phi_1,\Phi_2$ совпадают.
Аналогично $B(i,1)\cap B(-1,1)$ лежит во втором квадранте, где $\arg^*z=\arg^{***}z\in(\pi/2,\pi)$;
$B(-1,1)\cap B(-i,1)$ --- в третьем, где $\arg^{***}z=\arg^{**}z\in(\pi,3\pi/2)$;
$B(-i,1)\cap B(1,1)$ --- в четвёртом, где $\arg^{**}z\in(3\pi/2,2\pi)$, $\arg z\in(-\pi/2,0)$, т.е. $\arg^{**}z=\arg z+2\pi$, и
$\sqrt{|z|}e^{i\arg^{**}z/2}=e^{i\pi}\sqrt{|z|}e^{i\arg z/2}$ --- функция элемента $\Phi_5$.

\emph{Каноничность.} Каждая из функций $f$ удовлетворяет $f^2=z$, $f'=1/(2f)$, $|f'(z)|=1/(2\sqrt{|z|})\to\infty$ при $z\to0$;
значит, $f$ не продолжается голоморфно ни в какой круг с центром $a$ ($|a|=1$) радиуса $>1$ (такой круг содержит $0$, а голоморфная
в окрестности нуля функция имеет ограниченную производную около нуля; продолжение совпадало бы с $f$ на $B(a,1)$ по теореме единственности).
По теореме Коши--Тейлора радиус сходимости равен $1$.
\end{detail}

С другой стороны, тривиальная цепочка из одного элемента даёт $F_0$. Итак, \emph{продолжения по разным цепочкам с общими началом и
кругом конца могут различаться}.

\subsection{Простейшие свойства НАП}

\begin{thmbox}{Замечание (слайд 358)}
Если $F_0$ и $F_1$ --- канонические элементы с общим центром $a$ и $F_1$ --- НАП $F_0$, то $F_0=F_1$.
\end{thmbox}
\begin{proof}
$f_0=f_1$ на $B_0\cap B_1$, а это круг с центром $a$; значит, все производные в $a$ совпадают, ряды Тейлора одинаковы.
Канонические круги --- круги сходимости одного и того же ряда, т.е. $B_0=B_1$, а $f_0=f_1$ --- суммы одного ряда (Коши--Тейлор).
\end{proof}

\begin{thmbox}{Предложение 15.1 (простейшие свойства НАП)}
\begin{enumerate}
\item Пусть $F_0=(B_0,f_0)$, $F_1=(B_1,f_1)$ --- элементы, $F_1$ --- НАП $F_0$, и центр $a_1$ элемента $F_1$ лежит в $B_0$. Тогда
\[
f_1(z)=\sum_{n=0}^\infty\frac{f_0^{(n)}(a_1)}{n!}(z-a_1)^n\qquad\forall z\in B_1 .
\]
Обратно: если $a_1\in B_0$ и $B_1$ --- круг сходимости этого ряда, а $f_1$ --- его сумма, то $F_1=(B_1,f_1)=E_{a_1}(f_0)$ --- НАП $F_0$;
его радиус не меньше $R_0-|a_1-a_0|$.
\item Пусть $F_2$ --- НАП $F_1$, $F_3$ --- НАП $F_2$ и $B_1\cap B_2\cap B_3\ne\varnothing$. Тогда $F_3$ --- НАП $F_1$.
\end{enumerate}
\end{thmbox}
\begin{remarkbox}{Замечание}
На слайде в п.~1 пропущено условие «$F_1$ --- НАП $F_0$»; без него утверждение, очевидно, неверно.
\end{remarkbox}
\begin{proof}
1) $a_1\in B_0\cap B_1$ и это множество открыто, поэтому $f_1=f_0$ в окрестности $a_1$ и $f_1^{(n)}(a_1)=f_0^{(n)}(a_1)$. Так как $f_1\in H(B_1)$,
а $B_1$ --- круг с центром $a_1$, по теореме Коши--Тейлора $f_1$ --- сумма своего ряда Тейлора в $a_1$ во всём $B_1$.

Обратно, по теореме Коши--Тейлора для $f_0$ в круге $B(a_1,R_0-|a_1-a_0|)\subset B_0$ ряд сходится там к $f_0$; поэтому радиус сходимости
не меньше $R_0-|a_1-a_0|>0$, а $f_1=f_0$ на этом круге. Итак $f_1=f_0$ на непустом открытом подмножестве выпуклого множества
$B_0\cap B_1$, значит, на всём $B_0\cap B_1$.

2) На непустом открытом множестве $B_1\cap B_2\cap B_3$ имеем $f_1=f_2=f_3$. Множество $B_1\cap B_3$ выпукло, т.е. связно, и содержит
это множество; по теореме единственности $f_1=f_3$ на $B_1\cap B_3$.
\end{proof}

\subsection{Технические операции с продолжениями вдоль пути}

Следующие простые утверждения на слайдах явно не выделены (слайд 381 оставляет их как упражнение), но постоянно используются.

\begin{thmbox}{Лемма 25.A (замена параметра, обращение, склейка, сужение)}
Пусть $\{F_t\}_{t\in I}$ --- продолжение $F_0$ вдоль $\gamma\colon I\to\CC$.
\begin{enumerate}
\item Если $\varphi\colon[\alpha,\beta]\to I$ --- возрастающий гомеоморфизм, то $\{F_{\varphi(s)}\}_{s\in[\alpha,\beta]}$ --- продолжение $F_0$ вдоль
$\gamma\circ\varphi$. В частности, продолжаемость вдоль пути на $[\alpha,\beta]$ равносильна продолжаемости вдоль пути
$t\mapsto\sigma((\beta-\alpha)t+\alpha)$ на $I$ (слайд 381), а результат продолжения не зависит от параметризации.
\item $\{F_{1-t}\}$ --- продолжение $F_1$ вдоль $\gamma^-$; результат --- $F_0$.
\item Для $0\le t_1<t_2\le1$ семейство $\{F_t\}_{t\in[t_1,t_2]}$ --- продолжение $F_{t_1}$ вдоль $\gamma|_{[t_1,t_2]}$.
В частности, \emph{все $F_t$ получаются из $F_0$ продолжением} (вдоль $\gamma|_{[0,t]}$).
\item Если $\{G_t\}$ --- продолжение $G_0=F_1$ вдоль $\sigma$, $\sigma(0)=\gamma(1)$, то семейство $H_t=F_{2t}$ ($t\le\frac12$), $H_t=G_{2t-1}$ ($t\ge\frac12$)
--- продолжение $F_0$ вдоль $\gamma\sigma$ с результатом $G_1$.
\end{enumerate}
\end{thmbox}
\begin{proof}
Условие 1 определения переносится тривиально. Для условия 2 берём образы (прообразы) окрестностей $U_t$: при гомеоморфизме
промежутка связная окрестность переходит в связную окрестность, а условия «$\gamma(\xi)\in B_t$, $F_\xi$ --- НАП $F_t$» не зависят от
параметризации. Это доказывает п.~1, п.~2 (гомеоморфизм $t\mapsto1-t$ убывает, но определение не использует ориентацию) и п.~3
(в качестве окрестности берём $U_t\cap[t_1,t_2]$ --- пересечение промежутков есть промежуток). В п.~4 для $t\ne\frac12$ окрестности
получаются из окрестностей для $\{F_t\}$ или $\{G_t\}$ (пересечённых с $[0,\frac12]$ или $[\frac12,1]$), а для $t=\frac12$ берём объединение
$\{\xi\le\frac12:2\xi\in U_1\}\cup\{\xi\ge\frac12:2\xi-1\in U'_0\}$ ($U'_0$ --- окрестность $0$ для $\{G_t\}$) --- объединение двух промежутков с общей
точкой $\frac12$ есть промежуток.
\end{proof}

\begin{thmbox}{Лемма 25.B (переразложение внутри круга)}
Пусть $F=(B(c,R),f)$ --- элемент, $\lambda\colon I\to B(c,R)$ --- путь. Тогда $\Psi_t:=E_{\lambda(t)}(f)$ --- продолжение $\Psi_0=E_{\lambda(0)}(f)$
вдоль $\lambda$, радиусы $\ge R-|\lambda(t)-c|>0$, и каждый $\Psi_t$ --- НАП $F$.
\end{thmbox}
\begin{proof}
По предложению 15.1(1) (обратная часть) $\Psi_t$ --- НАП $F$ и радиус $\rho_t\ge R-|\lambda(t)-c|>0$. Фиксируем $t$; по непрерывности $\lambda$
найдётся промежуток $U_t\ni t$, на котором $|\lambda(\xi)-\lambda(t)|<\rho_t$, т.е. $\lambda(\xi)\in$ круг $\Psi_t$. Точка $\lambda(\xi)$ лежит
в кругах $\Psi_t$, $\Psi_\xi$ и $B(c,R)$; в её окрестности функции элементов $\Psi_t$ и $\Psi_\xi$ обе равны $f$ (оба --- НАП $F$).
Пересечение кругов $\Psi_t,\Psi_\xi$ выпукло, значит функции совпадают на всём пересечении: $\Psi_\xi$ --- НАП $\Psi_t$.
\end{proof}

\subsection{Продолжение вдоль пути: непрерывность радиуса и единственность}

\begin{remarkbox}{Почему требуется каноничность элементов в определении 15.4}
Без каноничности теряется единственность и непрерывность радиуса: например, семейство $F_t=(B(0,1-t/2),z)$ вдоль постоянного пути
$\gamma\equiv0$ удовлетворяло бы условиям, а радиус можно было бы менять произвольно (и даже разрывно).
Каноничность «жёстко» привязывает элемент к его центру и ростку функции.
\end{remarkbox}

\begin{thmbox}{Лемма 15.1 (о радиусе)}
В условиях определения 15.4 либо $R_t\equiv\infty$ на $I$, либо все $R_t<\infty$, и тогда $R_t$ непрерывна на $I$; точнее,
\[
|R_t-R_\xi|\le|\gamma(t)-\gamma(\xi)|\qquad\text{при всех }t\in I,\ \xi\in U_t .
\]
Как следствие, $R_\gamma:=\min_{t\in I}R_t>0$.
\end{thmbox}
\begin{proof}
\emph{Шаг 1: случай $R_t=\infty$.} Пусть $R_t=\infty$ для некоторого $t$, т.е. $f_t$ --- целая функция. Если $\xi\in U_t$, то $F_\xi$ --- НАП $F_t$ и центр
$\gamma(\xi)\in B_t=\CC$, так что по предложению 15.1(1) $f_\xi$ --- ряд Тейлора целой функции $f_t$ в точке $\gamma(\xi)$; он сходится во
всём $\CC$, и в силу каноничности $R_\xi=\infty$. Итак, множество $E_\infty=\{t:R_t=\infty\}$ \emph{открыто} в $I$ (вместе с $t$ содержит $U_t$).
Оно и \emph{замкнуто}: если $t$ --- предельная точка $E_\infty$, то найдётся $\xi\in U_t\cap E_\infty$. Тогда $F_t$ --- НАП $F_\xi$ (симметрия НАП),
центр $\gamma(t)$ лежит в $B_\xi=\CC$, и, как только что показано (с переставленными ролями), $R_t=\infty$. Так как $I$ связно, $E_\infty=\varnothing$ или $E_\infty=I$.

\emph{Шаг 2: все $R_t$ конечны.} Пусть $\xi\in U_t$, $d=|\gamma(\xi)-\gamma(t)|<R_t$. По предложению 15.1(1) ряд $f_\xi$ сходится в
$B(\gamma(\xi),R_t-d)$, поэтому $R_\xi\ge R_t-d$. Обратно, предположим, что $R_\xi>R_t+d$. Тогда $\gamma(t)\in B_\xi$ и
$B(\gamma(t),R_\xi-d)\subset B_\xi$ --- круг с центром $\gamma(t)$ радиуса $R_\xi-d>R_t$. Функция $f_\xi$ голоморфна в нём и совпадает с $f_t$
на $B_t\cap B_\xi\ni\gamma(\xi)$ (непустое открытое множество, а пересечение кругов связно). По теореме единственности $f_\xi=f_t$ на круге
$B_t\subset B(\gamma(t),R_\xi-d)$, так что ряд Тейлора $f_t$ в $\gamma(t)$ совпадает с рядом $f_\xi$ в $\gamma(t)$ и по Коши--Тейлору сходится в
$B(\gamma(t),R_\xi-d)$ --- противоречие с каноничностью $F_t$. Итак $R_t-d\le R_\xi\le R_t+d$.

\emph{Непрерывность.} $U_t$ --- окрестность $t$ в $I$, и при $\xi\to t$, $\xi\in U_t$, имеем $|R_\xi-R_t|\le|\gamma(\xi)-\gamma(t)|\to0$
по непрерывности $\gamma$. Непрерывная положительная функция на компакте $I$ достигает минимума, поэтому $R_\gamma>0$.
\end{proof}

\begin{remarkbox}{Геометрический смысл (как на слайде 360)}
На слайде то же неравенство получено так: окружности $\partial B_t$ и $\partial B_\xi$ пересекаются (если бы один круг компактно лежал
в другом, то меньший был бы неканоничен --- функция продолжается в больший концентрический круг), и для точки $w$ пересечения
неравенство треугольника в треугольнике $\gamma(t),\gamma(\xi),w$ даёт $|R_t-R_\xi|\le|\gamma(t)-\gamma(\xi)|$.
Наш шаг~2 --- аналитическая запись того же рассуждения.
\end{remarkbox}

\begin{thmbox}{Предложение 15.2 (единственность продолжения вдоль пути)}
Пусть семейства $\{F_t\}$ и $\{\widetilde F_t\}$ реализуют продолжение канонического элемента $F_0=\widetilde F_0$ вдоль пути
$\gamma\colon I\to\CC$. Тогда $F_t=\widetilde F_t$ при всех $t\in I$; в частности, $F_1=\widetilde F_1$.
\end{thmbox}
\begin{proof}
Пусть $E=\{t\in I:F_t=\widetilde F_t\}$; $0\in E$. Обозначим через $U_t,\widetilde U_t$ окрестности из определения для двух семейств;
$U_t\cap\widetilde U_t$ --- окрестность $t$.

\emph{$E$ открыто.} Пусть $t\in E$, $\xi\in U_t\cap\widetilde U_t$. Тогда $F_\xi$ и $\widetilde F_\xi$ --- НАП элемента $F_t=\widetilde F_t$ с центром
$\gamma(\xi)\in B_t$. По предложению 15.1(1) функции обоих --- ряд Тейлора $f_t$ в точке $\gamma(\xi)$, а круги --- его круг сходимости
(каноничность). Значит, $F_\xi=\widetilde F_\xi$, и $U_t\cap\widetilde U_t\subset E$.

\emph{$E$ замкнуто.} Пусть $t$ --- предельная точка $E$. Возьмём $\xi\in E\cap U_t\cap\widetilde U_t$. Тогда $f_t=f_\xi$ на $B_t\cap B_\xi$ и
$\tilde f_t=\tilde f_\xi=f_\xi$ на $\widetilde B_t\cap B_\xi$, поэтому $f_t=\tilde f_t$ на открытом множестве $B_t\cap\widetilde B_t\cap B_\xi$,
непустом, так как оно содержит $\gamma(\xi)$. Круги $B_t$ и $\widetilde B_t$ концентричны (центр $\gamma(t)$), их пересечение --- меньший из них;
по теореме единственности $f_t=\tilde f_t$ в окрестности $\gamma(t)$, т.е. $F_t$ --- НАП $\widetilde F_t$ с общим центром, и по замечанию
выше $F_t=\widetilde F_t$, $t\in E$.

Итак, $E$ --- непустое открыто-замкнутое подмножество связного $I$, т.е. $E=I$.
\end{proof}

\begin{remarkbox}{Замечание}
На слайде при доказательстве замкнутости говорится, что $F_t,\widetilde F_t$ --- НАП элемента $F_\xi$ «с центром $\gamma(t)$», и применяется
замечание об общем центре; но п.~1 предложения 15.1 требует $\gamma(t)\in B_\xi$, что из определения напрямую не следует
(там гарантировано лишь $\gamma(\xi)\in B_t$). Выше этот момент обойден через точку $\gamma(\xi)$.
\end{remarkbox}

\subsection{Связь продолжений по цепочке и вдоль пути}

\begin{thmbox}{Предложение 15.3}
\begin{enumerate}
\item Пусть $\{F_t\}_{t\in I}$ --- продолжение канонического $F_0$ вдоль пути $\gamma$. Тогда найдутся $0=t_0<t_1<\dots<t_N=1$ такие, что
$F_{t_n}$ --- НАП $F_{t_{n-1}}$, $n=1,\dots,N$; т.е. $F_1$ --- продолжение $F_0$ по цепочке $F_{t_0},\dots,F_{t_N}$.
\item Пусть канонический $F_1$ --- продолжение канонического $F_0$ по цепочке канонических элементов $\Phi_1=F_0,\dots,\Phi_N=F_1$
и $\lambda\colon I\to\CC$ --- ломаная, последовательно соединяющая центры $\Phi_1,\dots,\Phi_N$. Тогда существует семейство $\{\Psi_t\}$,
$\Psi_0=F_0$, реализующее продолжение $F_0$ вдоль $\lambda$, и $\Psi_1=F_1$.
\end{enumerate}
\end{thmbox}
\begin{remarkbox}{Замечание}
На слайде во втором пункте написано «вдоль $\gamma$»; должно быть «вдоль $\lambda$». Кроме того, слайд оставляет второй пункт упражнением,
а в первом пункте не объяснено, почему соседние элементы подпокрытия являются НАП друг друга; ниже оба места доказаны полностью
(вместо выбора подпокрытия используется число Лебега).
\end{remarkbox}
\begin{proof}
\emph{1)} По лемме 15.1 найдётся $\eps>0$ с $R_t>\eps$ при всех $t$ (если $R_t\equiv\infty$, берём $\eps=1$). По равномерной непрерывности $\gamma$
на компакте $I$ найдётся $\delta_1>0$: $|s-t|<\delta_1\Rightarrow|\gamma(s)-\gamma(t)|<\eps$. Окрестности $\{\operatorname{int}_I U_t\}_{t\in I}$
(внутренности относительно $I$) образуют открытое покрытие компакта $I$; пусть $\delta_2>0$ --- его число Лебега: любое подмножество $I$
диаметра $<\delta_2$ лежит в одном из $U_\tau$. Возьмём разбиение $0=t_0<\dots<t_N=1$ с шагом $<\min(\delta_1,\delta_2)$.
Для каждого $n$ найдётся $\tau$ с $[t_{n-1},t_n]\subset U_\tau$; тогда $F_{t_{n-1}}$ и $F_{t_n}$ --- НАП $F_\tau$ и $\gamma(t_n)\in B_\tau$.
Кроме того, $\gamma(t_n)\in B_{t_n}$ (центр) и $|\gamma(t_n)-\gamma(t_{n-1})|<\eps<R_{t_{n-1}}$, т.е. $\gamma(t_n)\in B_{t_{n-1}}$. Значит,
$B_{t_{n-1}}\cap B_\tau\cap B_{t_n}\ne\varnothing$, и по предложению 15.1(2) (цепочка $F_{t_{n-1}},F_\tau,F_{t_n}$) $F_{t_n}$ --- НАП $F_{t_{n-1}}$.

\emph{2)} По лемме 25.A(1,4) достаточно разобрать одно звено: $F_1=(B_1,f_1)$ --- НАП $F_0=(B_0,f_0)$, $B_j=B(a_j,R_j)$, $\lambda$ --- отрезок
$[a_0,a_1]$ (при $a_0=a_1$ по замечанию выше $F_1=F_0$ и годится постоянное семейство). Пусть $d=|a_1-a_0|>0$. Так как $B_0\cap B_1\ne\varnothing$,
то $d<R_0+R_1$, поэтому интервал $(d-R_1,R_0)\cap(0,d)$ непуст; возьмём из него $x$ и точку $p=a_0+x\frac{a_1-a_0}{d}$. Тогда
$|p-a_0|=x<R_0$, $|p-a_1|=d-x<R_1$, т.е. $p\in B_0\cap B_1$, а отрезки $[a_0,p]\subset B_0$, $[p,a_1]\subset B_1$ (круги выпуклы).
По лемме 25.B семейство $E_z(f_0)$, $z\in[a_0,p]$, --- продолжение $E_{a_0}(f_0)=F_0$ вдоль $[a_0,p]$ с результатом $E_p(f_0)$; семейство
$E_z(f_1)$, $z\in[p,a_1]$, --- продолжение $E_p(f_1)$ вдоль $[p,a_1]$ с результатом $E_{a_1}(f_1)=F_1$. Но $f_0=f_1$ в окрестности
$p\in B_0\cap B_1$, так что $E_p(f_0)=E_p(f_1)$, и по лемме 25.A(4) (склейка, затем замена параметра) получаем продолжение $F_0$ вдоль
$\lambda=[a_0,a_1]$ с результатом $F_1$. (Здесь использовано, что $F_0,F_1$ канонические: $E_{a_j}(f_j)=F_j$.)
\end{proof}

\begin{remarkbox}{Следствие}
Продолжения вдоль путей --- это ровно продолжения по цепочкам канонических элементов, \emph{но} вдоль пути результат определён
однозначно (предложение 15.2), а по цепочке --- нет (пример с $\sqrt z$). Разница в том, что путь «помнит», как именно мы обходим
особенности: в примере ломаная, соединяющая центры $1,i,-1,-i,1$, обходит $0$, а тривиальная цепочка --- нет.
\end{remarkbox}

\begin{remarkbox}{Коротко: план ответа на экзамене}
\begin{enumerate}
\item Определения 15.1--15.3: непосредственное продолжение, элемент, канонический элемент, НАП, продолжение по цепочке.
\item Пример $\sqrt z$: продолжение по цепочке неоднозначно.
\item Предложение 15.1: переразложение в ряд Тейлора и «транзитивность при непустом тройном пересечении».
\item Определение 15.4 продолжения вдоль пути; роль каноничности.
\item Лемма 15.1: $R_t\equiv\infty$ или $|R_t-R_\xi|\le|\gamma(t)-\gamma(\xi)|$; отсюда $R_\gamma>0$.
\item Предложение 15.2: единственность (множество совпадения открыто-замкнуто).
\item Предложение 15.3: путь $\Rightarrow$ цепочка (компактность, число Лебега); цепочка $\Rightarrow$ ломаная (переразложение вдоль отрезка).
\end{enumerate}
\end{remarkbox}

% =====================================================================
\ticket{26}{Аналитическое продолжение по близким путям. Теорема об аналитическом продолжении по гомотопным путям. Теорема о монодромии}
\label{t26:start}

\begin{defs}
\textbf{Продолжение вдоль пути} (определение 15.4, см. билет 25). Семейство канонических элементов $F_t=(B(\gamma(t),R_t),f_t)$,
$t\in I=[0,1]$, $R_t>0$, такое, что у каждого $t$ есть промежуток-окрестность $U_t\subset I$, для всех $\xi\in U_t$ которой
$\gamma(\xi)\in B_t$ и $F_\xi$ --- НАП $F_t$ (т.е. круги пересекаются и функции совпадают на пересечении).
Результат $F_1$ определён однозначно (предложение 15.2). Всюду ниже $E_a(f)$ --- канонический элемент с центром $a$, задаваемый $f$.

\textbf{Радиус продолжения:} $R_\gamma:=\min_{t\in I}R_t\in(0,+\infty]$ --- положителен по лемме 15.1 ($R_t$ либо $\equiv\infty$, либо
непрерывна и $|R_t-R_\xi|\le|\gamma(t)-\gamma(\xi)|$ при $\xi\in U_t$). По предложению 15.2 $R_\gamma$ не зависит от выбора семейства.

\textbf{Близкие пути:} $\|\gamma-\sigma\|_I=\max_{t\in I}|\gamma(t)-\sigma(t)|$.

\textbf{Гомотопия с закреплёнными концами.} Пути $\gamma_0,\gamma_1\colon I\to D$ с $\gamma_0(0)=\gamma_1(0)=a$, $\gamma_0(1)=\gamma_1(1)=b$
\emph{гомотопны в $D$}, если существует непрерывное $\gamma\colon I\times I\to D$ с $\gamma(0,t)=\gamma_0(t)$, $\gamma(1,t)=\gamma_1(t)$,
$\gamma(s,0)=a$, $\gamma(s,1)=b$ при всех $s,t$. Пишем $\gamma_s(t)=\gamma(s,t)$; каждый $\gamma_s$ --- путь из $a$ в $b$.

\textbf{Односвязная область} $D\subset\CC$: любой замкнутый путь в $D$ гомотопен в $D$ постоянному. Известный топологический факт
(прошлый семестр/курс топологии): \emph{в односвязной области любые два пути с общими началом и концом гомотопны в $D$ с закреплёнными
концами} (фундаментальная группа тривиальна). Примеры односвязных областей: выпуклые и звёздные области, $\CC\setminus\RR_\alpha$,
круг с разрезом по радиусу. Неодносвязны: $\CC^*=\CC\setminus\{0\}$, кольца, проколотые круги.

\textbf{Равномерная непрерывность.} Непрерывное отображение компакта ($I$ или $I\times I$) равномерно непрерывно (теорема Кантора).
\end{defs}

\subsection{Лемма о продолжении по близким путям}

На слайдах это утверждение доказано внутри доказательства теоремы 15.1 (с оценкой $\eps/4$) и затем сформулировано отдельно как
предложение 15.4. Мы докажем его сразу в сильной форме и затем выведем из него теорему 15.1.

\begin{thmbox}{Предложение 15.4 (лемма о продолжении по близким путям)}
Пусть канонический элемент $F=F_0$ продолжается вдоль пути $\gamma\colon I\to\CC$ семейством $\{F_t=(B_t,f_t)\}$, и пусть путь
$\sigma\colon I\to\CC$ имеет те же концы: $\sigma(0)=\gamma(0)$, $\sigma(1)=\gamma(1)$, причём $\|\gamma-\sigma\|_I<R_\gamma$.
Тогда $F$ продолжается вдоль $\sigma$, а именно семейством
\[
\Phi_t:=E_{\sigma(t)}(f_t)\qquad(\text{переразложение $f_t$ в ряд Тейлора с центром }\sigma(t)\in B_t),
\]
и результаты продолжения вдоль $\gamma$ и $\sigma$ совпадают: $\Phi_1=F_1$.
\end{thmbox}
\begin{proof}
Выберем конечное $r$ с $\|\gamma-\sigma\|_I<r\le R_\gamma$ (при $R_\gamma<\infty$ можно $r=R_\gamma$) и положим $\rho:=r-\|\gamma-\sigma\|_I>0$.

\emph{Шаг 1: элементы $\Phi_t$.} Так как $|\sigma(t)-\gamma(t)|<r\le R_t$, то $\sigma(t)\in B_t$, и по предложению 15.1(1) $\Phi_t$ --- НАП $F_t$,
а его радиус не меньше $R_t-|\sigma(t)-\gamma(t)|\ge r-\|\gamma-\sigma\|_I=\rho>0$. Элементы $\Phi_t$ канонические по построению,
центр $\Phi_t$ --- $\sigma(t)$. Далее, $\Phi_0=E_{a}(f_0)=F_0$ и $\Phi_1=E_b(f_1)=F_1$, так как $\sigma(0)=\gamma(0)=a$, $\sigma(1)=\gamma(1)=b$, а $F_0,F_1$ канонические.

\emph{Шаг 2: окрестности.} По равномерной непрерывности $\gamma$ и $\sigma$ на $I$ найдётся $\delta>0$ такое, что при $|\xi-t|<\delta$
\[
|\gamma(\xi)-\gamma(t)|<\rho/2,\qquad |\sigma(\xi)-\sigma(t)|<\rho/2 .
\]
Для $t\in I$ положим $\Omega_t:=U_t\cap(t-\delta,t+\delta)$ --- промежуток, окрестность $t$ в $I$.

\emph{Шаг 3: проверка условия 2 определения.} Пусть $\xi\in\Omega_t$ и $w:=\sigma(\xi)$. Тогда $w$ лежит во всех четырёх кругах:
\begin{itemize}
\item в круге $\Phi_t$: $|w-\sigma(t)|<\rho/2<\rho\le$ радиус $\Phi_t$ (в частности, $\sigma(\xi)\in$ круг $\Phi_t$, как и требует определение);
\item в круге $\Phi_\xi$: это его центр;
\item в $B_\xi$: $|w-\gamma(\xi)|\le\|\sigma-\gamma\|_I<r\le R_\xi$;
\item в $B_t$: $|w-\gamma(t)|\le|\sigma(\xi)-\sigma(t)|+|\sigma(t)-\gamma(t)|<\rho/2+\|\sigma-\gamma\|_I=r-\rho/2<R_t$.
\end{itemize}
В некоторой окрестности $w$ (пересечении четырёх кругов) имеем: функция $\Phi_t$ равна $f_t$ ($\Phi_t$ --- НАП $F_t$), $f_t=f_\xi$
($\xi\in U_t$, $F_\xi$ --- НАП $F_t$), $f_\xi$ равна функции $\Phi_\xi$ ($\Phi_\xi$ --- НАП $F_\xi$). Значит, функции элементов $\Phi_t$ и $\Phi_\xi$ совпадают на
непустом открытом подмножестве пересечения их кругов, а оно выпукло, т.е. связно; по теореме единственности они совпадают на всём
пересечении: $\Phi_\xi$ --- НАП $\Phi_t$.

Итак, $\{\Phi_t\}$ --- продолжение $F$ вдоль $\sigma$, а его результат $\Phi_1=F_1$. По предложению 15.2 любое другое продолжающее семейство
вдоль $\sigma$ даёт тот же результат.
\end{proof}

\begin{remarkbox}{Замечания}
1) Если $R_\gamma=\infty$, то условие $\|\gamma-\sigma\|<R_\gamma$ выполнено для любого $\sigma$; в этом случае все $f_t$ --- одна и та же
целая функция (по предложению 15.3(1) и предложению 15.1(1)), и утверждение тривиально.

2) Из доказательства видно, что радиус продолжения вдоль $\sigma$ не меньше $R_\gamma-\|\gamma-\sigma\|_I$ (шаг 1 и предложение 15.2).

3) Концы путей должны совпадать лишь для того, чтобы $\Phi_0=F_0$ и $\Phi_1=F_1$; без этого условия переразложение даёт продолжение
элемента $E_{\sigma(0)}(f_0)$ вдоль $\sigma$ с результатом $E_{\sigma(1)}(f_1)$.
\end{remarkbox}

\begin{thmbox}{Дополнение к предложению 15.4 (выбор окрестностей $U_t$; слайд 366)}
Пусть $\{F_t\}$ --- продолжение вдоль $\gamma$, $t\in I$, и $W\ni t$ --- промежуток в $I$ такой, что $\gamma(t)\in B_\xi$ при всех $\xi\in W$.
Тогда $F_\xi$ --- НАП $F_t$ при всех $\xi\in W$. В частности, в качестве $U_t$ можно взять компоненту связности точки $t$ множества
$\{\xi\in I:|\gamma(\xi)-\gamma(t)|<R_\gamma\}$ (на ней $\gamma(t)\in B_\xi$ и $\gamma(\xi)\in B_t$).
\end{thmbox}
\begin{proof}
Для $s\in W$ положим $G_s:=E_{\gamma(t)}(f_s)$ (корректно, так как $\gamma(t)\in B_s$). Если $s\in W$ и $s'\in U_s\cap W$, то $F_{s'}$ --- НАП $F_s$ и
$\gamma(t)\in B_s\cap B_{s'}$, откуда $f_s=f_{s'}$ в окрестности $\gamma(t)$ и $G_s=G_{s'}$. Значит, $s\mapsto G_s$ локально постоянно на связном $W$,
т.е. постоянно: $G_\xi=G_t=F_t$. Поэтому $f_\xi=f_t$ около $\gamma(t)$, и, так как $B_\xi\cap B_t$ выпукло, --- на всём $B_\xi\cap B_t$.
Для указанной компоненты $|\gamma(\xi)-\gamma(t)|<R_\gamma\le\min(R_t,R_\xi)$, так что оба включения выполнены.
\end{proof}

\subsection{Теорема о продолжении по гомотопным путям}

\begin{thmbox}{Теорема 15.1 (о продолжении по гомотопным путям)}
Пусть $\gamma_0,\gamma_1\colon I\to\CC$ --- пути из $a$ в $b$, гомотопные (с закреплёнными концами) посредством $\gamma\colon I\times I\to\CC$.
Пусть канонический элемент $F$ с центром $a$ продолжается вдоль \emph{каждого} пути $\gamma_s=\gamma(s,\cdot)$, $s\in I$, и
$\{F_{s,t}\}_{t\in I}$ --- соответствующие продолжающие семейства. Тогда $F_{0,1}=F_{1,1}$; более того, $F_{s,1}$ не зависит от $s$.
\end{thmbox}
\begin{proof}
Пусть $F_{s,t}=(B(\gamma(s,t),R_{s,t}),f_{s,t})$.

\emph{Шаг 1 (положительный радиус для каждого пути).} При фиксированном $s$ по лемме 15.1 $t\mapsto R_{s,t}$ либо $\equiv\infty$, либо непрерывна и
положительна на компакте $I$; поэтому $r_s:=\min\bigl(1,\min_{t\in I}R_{s,t}\bigr)>0$ и $r_s\le R_{\gamma_s}$.

\emph{Шаг 2 (равномерная непрерывность гомотопии).} Отображение $\gamma$ непрерывно на компакте $I\times I$, значит, равномерно непрерывно:
для $r_s>0$ найдётся $\eta_s>0$ такое, что $|(s',t')-(s'',t'')|<\eta_s\Rightarrow|\gamma(s',t')-\gamma(s'',t'')|<r_s$. В частности, при $|s'-s|<\eta_s$
\[
\|\gamma_{s'}-\gamma_s\|_I=\max_{t\in I}|\gamma(s',t)-\gamma(s,t)|<r_s\le R_{\gamma_s}.
\]
(Ключевой момент: оценка \emph{равномерна по $t$} --- именно для этого нужна равномерная непрерывность на квадрате, а не просто
непрерывность по $s$ при каждом $t$.)

\emph{Шаг 3 (локальное постоянство).} Пусть $|s'-s|<\eta_s$. Пути $\gamma_s$ и $\gamma_{s'}$ имеют общие концы $a,b$, и по шагу 2 и предложению 15.4
продолжение $F$ вдоль $\gamma_{s'}$ существует и даёт $F_{s,1}$. По единственности (предложение 15.2) $F_{s',1}=F_{s,1}$.

\emph{Шаг 4 (связность/компактность $I$).} Пусть $E=\{s\in I:F_{s,1}=F_{0,1}\}\ni0$. Если $s\in E$, то $(s-\eta_s,s+\eta_s)\cap I\subset E$ по шагу 3 ---
$E$ открыто. Если $s$ --- предельная точка $E$, возьмём $s'\in E$ с $|s'-s|<\eta_s$; по шагу 3 $F_{s,1}=F_{s',1}=F_{0,1}$ --- $E$ замкнуто.
Так как $I$ связно, $E=I$.
\end{proof}

\begin{detail}{Вариант шага 4 через компактность (как предлагается на слайде 365)}
Интервалы $J_s=(s-\eta_s,s+\eta_s)$, $s\in I$, покрывают компакт $I$; выберем конечное подпокрытие $J_{s_1},\dots,J_{s_m}$, $s_1<\dots<s_m$,
и выбросим лишние, так чтобы никакой интервал не содержался в объединении остальных. Тогда соседние интервалы пересекаются
(иначе между ними была бы непокрытая точка), $0\in J_{s_1}$, $1\in J_{s_m}$. По шагу 3 $F_{\cdot,1}$ постоянно на каждом $J_{s_k}\cap I$ (равно $F_{s_k,1}$),
а на пересечении соседних интервалов значения совпадают; по цепочке $F_{0,1}=F_{s_1,1}=\dots=F_{s_m,1}=F_{1,1}$.
\end{detail}

\begin{remarkbox}{Условие продолжаемости вдоль всех $\gamma_s$ существенно}
Пусть $F=E_1(\sqrt z)$ (главная ветвь), $\gamma_0$ --- верхняя полуокружность из $1$ в $-1$, $\gamma_1$ --- нижняя. Они гомотопны в $\CC$
(линейная гомотопия $\gamma(s,t)=(1-s)\gamma_0(t)+s\gamma_1(t)$), но промежуточный путь при $s=\frac12$ --- отрезок $[1,-1]$ --- проходит через $0$,
и вдоль него $F$ не продолжается. Результаты: вдоль $\gamma_0$ --- элемент $E_{-1}(\sqrt{|z|}e^{i\arg^{***}z/2})$ со значением $i$ в точке $-1$,
вдоль $\gamma_1$ --- элемент со значением $-i$; они различны.
\end{remarkbox}

\subsection{Теорема о монодромии}

\begin{thmbox}{Лемма 26.A (от независимости от пути --- к голоморфной функции)}
Пусть $G$ --- область, $c\in G$, канонический элемент $F$ с центром $c$ продолжается вдоль любого пути в $G$ с началом $c$, и результат
продолжения в любую точку $z\in G$ не зависит от пути; обозначим его $F^z=(B(z,R^z),f^z)$. Положим $\Phi(z):=f^z(z)$, $z\in G$. Тогда
$\Phi\in H(G)$, $\Phi=f^{z_0}$ на $B(z_0,r)$, $r=\min(R^{z_0},\dist(z_0,\partial G))$, и $E_{z}(\Phi)=F^z$ для всех $z\in G$.
\end{thmbox}
\begin{proof}
Фиксируем $z_0\in G$ и путь $\gamma$ в $G$ из $c$ в $z_0$. Пусть $z\in B(z_0,r)$. Отрезок $[z_0,z]$ лежит в $B(z_0,r)\subset G\cap B(z_0,R^{z_0})$.
По лемме 25.B семейство $E_w(f^{z_0})$, $w\in[z_0,z]$, продолжает $E_{z_0}(f^{z_0})=F^{z_0}$ вдоль отрезка, с результатом $E_z(f^{z_0})$.
Склеивая (лемма 25.A(4)), получаем продолжение $F$ вдоль пути $\gamma\cdot[z_0,z]$ в $G$ с результатом $E_z(f^{z_0})$. По независимости от пути
$F^z=E_z(f^{z_0})$; в частности, $\Phi(z)=f^z(z)=f^{z_0}(z)$. Итак $\Phi=f^{z_0}$ на $B(z_0,r)$, т.е. $\Phi$ голоморфна около $z_0$, и
$E_{z_0}(\Phi)=E_{z_0}(f^{z_0})=F^{z_0}$.
\end{proof}

\begin{thmbox}{Теорема 15.2 (о монодромии)}
Пусть $D\subset\CC$ --- односвязная область, $F=(B,f)$ --- канонический элемент с центром $a\in D$, продолжающийся вдоль любого пути в $D$ с
началом $a$. Тогда для любой $b\in D$ результат продолжения $F$ вдоль пути в $D$ из $a$ в $b$ не зависит от пути; обозначим его $\widetilde F_b$.
Существует (и единственна) функция $\Phi\in H(D)$ с $E_b(\Phi)=\widetilde F_b$ для всех $b\in D$; в частности, $\Phi=f$ на $B\cap B(a,\dist(a,\partial D))$,
и в окрестности любой $b\in D$ функция $\Phi$ задаётся рядом Тейлора элемента $\widetilde F_b$.
\end{thmbox}
\begin{proof}
Пусть $\gamma_0,\gamma_1$ --- пути в $D$ из $a$ в $b$. Так как $D$ односвязна, они гомотопны в $D$ с закреплёнными концами; пусть $\gamma$ ---
гомотопия, $\gamma(I\times I)\subset D$. Каждый $\gamma_s$ --- путь в $D$ с началом $a$, поэтому по условию $F$ продолжается вдоль $\gamma_s$.
По теореме 15.1 результаты вдоль $\gamma_0$ и $\gamma_1$ совпадают. Итак, $\widetilde F_b$ корректно определён, и лемма 26.A (с $G=D$, $c=a$)
даёт $\Phi\in H(D)$ с $E_b(\Phi)=\widetilde F_b$. При $b=a$ (постоянный путь) $\widetilde F_a=F$, поэтому $\Phi=f$ около $a$, а значит, и на
выпуклом множестве $B\cap B(a,\dist(a,\partial D))$. Единственность $\Phi$: два таких $\Phi$ совпадают около $a$, а $D$ --- область.
\end{proof}

\begin{remarkbox}{Замечания}
1) Радиус $\widetilde F_b$ может быть больше $\dist(b,\partial D)$ (например, $D=\{\Re z>0\}$, $F=E_1(\ln z)$: радиус $\widetilde F_b$ равен $|b|>\Re b$);
утверждается лишь, что $\Phi$ совпадает с функцией элемента $\widetilde F_b$ на $B(b,\min(R,\dist(b,\partial D)))$.

2) Односвязность существенна: $D=\CC^*$, $F=E_1(\sqrt z)$ продолжается вдоль любого пути в $D$, но после обхода окружности $|z|=1$ переходит в
$E_1(-\sqrt z)$ (пример из билета 25), и голоморфной в $\CC^*$ функции $\Phi$ с $\Phi^2=z$ нет.

3) Следствие (слайды 377--378): если $D$ односвязна, $g\in H(D)$, $g(D)\subset\CC^*$, то существуют $h,q\in H(D)$ с $e^h=g$, $q^n=g$.
Действительно, элемент $E_a(\ln_{(*)}\circ g)$ (с любой ветвью логарифма около $g(a)$) продолжается вдоль любого пути $\gamma$ в $D$:
это композиция элемента $E_a(g)$ и логарифма, продолжаемого вдоль любого пути $g\circ\gamma\subset\CC^*$ (теорема о продолжении сложного
элемента, см. билет 29). По теореме о монодромии получаем $h\in H(D)$, а тождество $e^h=g$, верное около $a$, верно в $D$ по теореме
единственности; $q=e^{h/n}$.
\end{remarkbox}

\begin{remarkbox}{Коротко: план ответа на экзамене}
\begin{enumerate}
\item Напомнить определение продолжения вдоль пути, лемму 15.1 и $R_\gamma>0$.
\item Предложение 15.4: если $\|\gamma-\sigma\|<R_\gamma$ и концы общие, то $\Phi_t=E_{\sigma(t)}(f_t)$ продолжает $F$ вдоль $\sigma$ с тем же результатом;
ключ --- точка $\sigma(\xi)$ лежит в четырёх кругах.
\item Теорема 15.1: $r_s>0$; равномерная непрерывность гомотопии на $I\times I$ даёт $\|\gamma_{s'}-\gamma_s\|<r_s$ при $|s'-s|<\eta_s$;
$s\mapsto F_{s,1}$ локально постоянно; связность (или компактность) $I$.
\item Пример: условие продолжаемости вдоль всех $\gamma_s$ нельзя опустить ($\sqrt z$).
\item Теорема 15.2: в односвязной области пути с общими концами гомотопны, результат не зависит от пути;
лемма 26.A --- склейка с переразложением вдоль отрезка даёт голоморфную $\Phi$.
\item Пример $\sqrt z$ в $\CC^*$; следствие о логарифме и корне.
\end{enumerate}
\end{remarkbox}

% =====================================================================
\ticket{27}{Полная аналитическая функция в смысле Вейерштрасса (П.А.Ф.). Теорема Пуанкаре--Вольтерра. Аналитические и голоморфные ветви П.А.Ф. Особые точки П.А.Ф. и их классификация}
\label{t27:start}

\begin{defs}
\textbf{Напоминание} (билеты 25, 26). Элемент $F=(B(a,R),f)$; канонический --- если $B(a,R)$ --- круг сходимости ряда Тейлора $f$ в $a$;
$E_a(f)$ --- канонический элемент, задаваемый $f$ в точке $a$; НАП; продолжение вдоль пути (определение 15.4) и его единственность
(предложение 15.2); лемма 25.A (склейка, обращение, сужение путей); предложение 15.4 (близкие пути); теорема 15.1 (гомотопные пути);
теорема 15.2 (монодромия) и лемма 26.A.

\textbf{Индекс} замкнутого пути $\gamma$ относительно $a\notin\gamma^*$: $\Ind_a(\gamma)=\frac{\varphi(1)-\varphi(0)}{2\pi}\in\ZZ$, где $\varphi$ --- непрерывная
ветвь $\Arg(\gamma(t)-a)$ на $I$ (она существует и единственна с точностью до $2\pi k$). Для кусочно-гладкого $\gamma$:
$\Ind_a(\gamma)=\frac1{2\pi i}\int_\gamma\frac{dz}{z-a}$. Индекс аддитивен при склейке путей и не меняется при гомотопии в $\CC\setminus\{a\}$.

\textbf{Определение 15.5 (П.А.Ф.).} Пусть $F_a$ --- канонический элемент с центром $a\in\CC$. \emph{Полной аналитической функцией}, определяемой $F_a$,
называется совокупность $\Phi$ всех канонических элементов, получаемых из $F_a$ продолжением вдоль всех путей в $\CC$ с началом $a$,
вдоль которых продолжение возможно. $D_\Phi$ --- объединение кругов всех элементов $\Phi$ --- \emph{область определения} $\Phi$.
$\Phi(b)$ --- множество элементов $\Phi$ с центром $b$ («элементы над $b$»).

\textbf{Определение 15.6 (аналитическая ветвь).} Пусть $D_1\subset D_\Phi$ --- область, $F_b\in\Phi$, $b\in D_1$, и $F_b$ продолжается вдоль любого
пути в $D_1$ с началом $b$. Совокупность $\Phi_1$ всех полученных так элементов --- \emph{(в общем случае многозначная) аналитическая ветвь
$\Phi$ над $D_1$, определяемая $F_b$}.

\textbf{Определение 15.7 (голоморфная ветвь).} Если $D_2\subset D_1$ --- область, $f\in H(D_2)$ и $E_c(f)\in\Phi_1$ для всех $c\in D_2$, то $f$ ---
\emph{голоморфная (однозначная) ветвь} $\Phi_1$ в $D_2$.

\textbf{Определение 15.8 (распадение на ветви).} $\Phi_1$ \emph{распадается над $D_2$ на голоморфные ветви $\{f_j\}_{j\in J}$} (различные функции), если
каждая $f_j$ --- голоморфная ветвь $\Phi_1$ в $D_2$ и для любых $c\in D_2$, $F_c\in\Phi_1$ найдётся $j$ с $F_c=E_c(f_j)$.
Точка $a\in D_1$ --- \emph{правильная} для $\Phi_1$ ($\Phi_1$ аналитична над $a$), если $\Phi_1$ распадается на голоморфные ветви над некоторым кругом $B(a,\delta)$.

\textbf{Определение 15.9 (изолированные особые точки).} Пусть $D_1=V(a,\delta)\subset D_\Phi$ ($a\in\CC$ или $a=\infty$; сама $a$ может лежать
или не лежать в $D_\Phi$), $\Phi_1$ --- аналитическая ветвь $\Phi$ над $D_1$, определяемая $F_b$. Возможны три случая:
\begin{enumerate}
\item после одного обхода $a$ (по пути в $D_1$) $F_b$ переходит в себя --- \emph{изолированная особая точка однозначного характера}; тогда $\Phi_1$ ---
одна голоморфная функция $f_1\in H(D_1)$, и $a$ классифицируется как для $f_1$: устранимая, полюс, существенно особая;
\item при некотором $n\ge2$ элемент $F_b$ после $1,\dots,n-1$ обходов не переходит в себя, а после $n$ обходов переходит ---
\emph{точка ветвления порядка $n$} ($\Phi_1$ имеет над $a$ ветвление порядка $n$);
\item ни при каком числе обходов $F_b$ не переходит в себя --- \emph{точка ветвления бесконечного порядка (логарифмическая)}.
\end{enumerate}
\textbf{«Обход»} точки $a\in\CC$ из $b$ --- замкнутый путь в $D_1$ с началом $b$ и $\Ind_a=1$; для $a=\infty$ --- замкнутый путь в $\{|z|>1/\delta\}$
с $\Ind_0=\pm1$ (направление роли не играет, см. ниже).
\end{defs}

\subsection{П.А.Ф. и её область определения}

\begin{thmbox}{Утверждение 27.1}
1) П.А.Ф. определяется любым своим элементом: если $G\in\Phi$, то совокупность элементов, получаемых продолжением $G$ вдоль путей, есть $\Phi$.

2) $D_\Phi$ --- область, и над каждой точкой $b\in D_\Phi$ есть хотя бы один элемент $\Phi$.
\end{thmbox}
\begin{proof}
1) Пусть $G$ получен из $F_a$ вдоль $\gamma$. Если $H$ получен из $G$ вдоль $\sigma$, то $H$ получен из $F_a$ вдоль $\gamma\sigma$ (лемма 25.A(4)), т.е. $H\in\Phi$.
Обратно, $F_a$ получается из $G$ вдоль $\gamma^-$ (лемма 25.A(2)), поэтому любой $H\in\Phi$, полученный из $F_a$ вдоль $\sigma$, получается из $G$ вдоль $\gamma^-\sigma$.

2) $D_\Phi$ открыто как объединение кругов. Если $b\in D_\Phi$, то $b$ лежит в круге некоторого $G=(B,g)\in\Phi$, и по лемме 25.B (переразложение
вдоль отрезка от центра $G$ до $b$) элемент $E_b(g)$ получается из $G$ продолжением, т.е. $E_b(g)\in\Phi(b)$. Связность: если $G\in\Phi$ получен вдоль $\gamma$,
то по лемме 25.A(3) все $F_t$ ($t\in I$) лежат в $\Phi$, так что $\gamma^*\subset D_\Phi$; любой точке круга элемента $G$ соответствует отрезок до центра
$G$, затем путь $\gamma^-$ до $a$. Значит, $D_\Phi$ линейно связно.
\end{proof}

\subsection{Теорема Пуанкаре--Вольтерра}

\begin{thmbox}{Предложение 15.7 (теорема Пуанкаре--Вольтерра)}
Если $\Phi$ --- П.А.Ф. и $b\in D_\Phi$, то множество $\Phi(b)$ элементов $\Phi$ с центром $b$ не более чем счётно.
\end{thmbox}
\begin{proof}
Фиксируем $F_b\in\Phi(b)$ (утверждение 27.1(2)). По утверждению 27.1(1) любой $\widetilde F_b\in\Phi(b)$ получается из $F_b$ вдоль некоторого
замкнутого пути $\gamma$, $\gamma(0)=\gamma(1)=b$. Пусть $R_\gamma>0$ (лемма 15.1) и $r:=\min(R_\gamma,1)$.

\emph{Построение ломаной.} По равномерной непрерывности $\gamma$ найдётся $N$ такое, что $|\gamma(t)-\gamma(t')|<r/4$ при $|t-t'|\le1/N$.
Положим $t_k=k/N$, выберем точки $q_k\in\QQ+i\QQ$ с $|q_k-\gamma(t_k)|<r/4$ при $0<k<N$, и $q_0=q_N=b$. Пусть $\lambda$ --- ломаная, линейная на каждом
$[t_{k-1},t_k]$, $\lambda(t_k)=q_k$. При $t\in[t_{k-1},t_k]$ точка $\lambda(t)$ --- выпуклая комбинация $q_{k-1},q_k$, поэтому
\[
|\lambda(t)-\gamma(t)|\le\max_{j\in\{k-1,k\}}|q_j-\gamma(t)|\le\max_j\bigl(|q_j-\gamma(t_j)|+|\gamma(t_j)-\gamma(t)|\bigr)<\tfrac r4+\tfrac r4<R_\gamma .
\]
По предложению 15.4 $F_b$ продолжается вдоль $\lambda$ с тем же результатом $\widetilde F_b$.

\emph{Подсчёт.} Ломаная $\lambda$ (с указанной параметризацией) однозначно задаётся конечным набором рациональных точек $q_1,\dots,q_{N-1}$,
а результат продолжения вдоль неё определён однозначно (предложение 15.2). Множество
\[\textstyle\bigcup_{N\ge1}(\QQ+i\QQ)^{N-1}\]
счётно, а $\widetilde F_b$ --- образ
одного из его элементов при частично определённом отображении, сопоставляющем набору результат продолжения $F_b$ вдоль соответствующей ломаной. Образ
не более чем счётного множества не более чем счётен.
\end{proof}

\subsection{Аналитические и голоморфные ветви}

\begin{thmbox}{Утверждение 27.2 (упражнение слайда 374)}
Пусть $\Phi_1$ --- аналитическая ветвь $\Phi$ над областью $D_1$. Тогда: (а) любой элемент $\Phi_1$ продолжается вдоль любого пути в $D_1$, результат лежит в
$\Phi_1$, и $\Phi_1$ определяется любым своим элементом; (б) число $n\in\NN\cup\{\infty\}$ элементов $\Phi_1$ с центром $c$ одно и то же для всех $c\in D_1$.
\end{thmbox}
\begin{proof}
(а) Если $G\in\Phi_1$ получен из $F_b$ вдоль пути $\gamma$ в $D_1$, а $\sigma$ --- путь в $D_1$ из центра $G$, то $\gamma\sigma$ --- путь в $D_1$ из $b$, вдоль него $F_b$
продолжается, и по единственности и лемме 25.A(3) (сужение на вторую половину) результат --- продолжение $G$ вдоль $\sigma$. Он лежит в $\Phi_1$.
Независимость от выбора порождающего элемента --- как в утверждении 27.1(1) (с путями в $D_1$).

(б) Пусть $c,c'\in D_1$, $\sigma$ --- путь в $D_1$ из $c$ в $c'$. По (а) определено отображение $\Lambda_\sigma\colon\Phi_1(c)\to\Phi_1(c')$ --- продолжение вдоль $\sigma$.
По лемме 25.A(2) $\Lambda_{\sigma^-}\circ\Lambda_\sigma=\mathrm{id}$ и $\Lambda_\sigma\circ\Lambda_{\sigma^-}=\mathrm{id}$, т.е. $\Lambda_\sigma$ --- биекция, и мощности $\Phi_1(c)$ и $\Phi_1(c')$ равны.
\end{proof}

\begin{thmbox}{Утверждение 27.3 (ветви над односвязной областью; слайд 374)}
Если $D_1$ односвязна, то существует $f_1\in H(D_1)$ такая, что $\Phi_1=\{E_c(f_1):c\in D_1\}$. Таким образом, $f_1$ --- единственная голоморфная ветвь $\Phi_1$
в $D_1$, и $\Phi_1$ «отождествляется» с $f_1$; над каждой точкой $D_1$ ровно один элемент $\Phi_1$.
\end{thmbox}
\begin{proof}
$F_b$ продолжается вдоль всех путей в $D_1$; по теореме о монодромии (теорема 15.2) результат продолжения в $c$ не зависит от пути и равен $E_c(f_1)$ для
некоторой $f_1\in H(D_1)$. Все элементы $\Phi_1$ получаются именно так. Если $g$ --- голоморфная ветвь $\Phi_1$ в $D_2\subset D_1$, то $E_c(g)\in\Phi_1(c)=\{E_c(f_1)\}$, т.е.
$g=f_1$ в $D_2$.
\end{proof}

\begin{thmbox}{Утверждение 27.4 (распадение; правильные точки; слайд 375)}
1) Если $\Phi_1$ распадается над $D_2$ на ветви $\{f_j\}_{j\in J}$, то $J$ не более чем счётно.

2) Каждая точка $a\in D_1$ правильная: если $B(a,\delta)\subset D_1$, то $\Phi_1$ распадается над $B(a,\delta)$ на голоморфные ветви
$\{g_G\}_{G\in\Phi_1(a)}$, где $g_G\in H(B(a,\delta))$ --- функция, задаваемая элементом $G$ по теореме о монодромии.
\end{thmbox}
\begin{proof}
1) Фиксируем $c\in D_2$. Отображение $j\mapsto E_c(f_j)\in\Phi(c)$ инъективно: если $E_c(f_j)=E_c(f_k)$, то $f_j=f_k$ около $c$, и по теореме единственности ---
в $D_2$. По теореме Пуанкаре--Вольтерра $\Phi(c)$ не более чем счётно.

2) $G\in\Phi_1(a)$ продолжается вдоль всех путей в односвязном $B(a,\delta)$ (утверждение 27.2(а)), монодромия даёт $g_G$. Если $c\in B(a,\delta)$ и $H\in\Phi_1(c)$, то,
продолжив $H$ вдоль отрезка $[c,a]$, получим некоторый $G\in\Phi_1(a)$, и $H$ --- продолжение $G$ вдоль $[a,c]$, т.е. $H=E_c(g_G)$.
\end{proof}

\subsection{Изолированные особые точки: обоснование классификации}

\begin{thmbox}{Лемма 27.5 (замкнутые пути в проколотой окрестности)}
Пусть $D_1=V(a,\delta)$, $a\in\CC$, $0<\delta\le\infty$, $b\in D_1$, $c_b(t)=a+(b-a)e^{2\pi it}$ --- окружность через $b$ (лежит в $D_1$).
Любой замкнутый путь $\gamma$ в $D_1$ с $\gamma(0)=\gamma(1)=b$ гомотопен в $D_1$ (с закреплёнными концами) пути $c_b^k$ (окружность, пройденная $k$ раз;
$c_b^{k}$ при $k<0$ --- $|k|$ раз в обратную сторону, $c_b^0\equiv b$), где $k=\Ind_a(\gamma)$. То же верно для $D_1=\{|z|>R\}$ с $a:=0$.
\end{thmbox}
\begin{proof}
Пусть $\Pi=\{w:\Re w<\ln\delta\}$ ($\Pi=\CC$ при $\delta=\infty$); отображение $w\mapsto a+e^w$ переводит $\Pi$ в $D_1$. Пусть $\varphi$ --- непрерывная ветвь
$\Arg(\gamma(t)-a)$; положим $L(t)=\ln|\gamma(t)-a|+i\varphi(t)\in\Pi$ («поднятие»: $a+e^{L(t)}=\gamma(t)$). Тогда $L(1)-L(0)=2\pi ik$. Пусть $M(t)=L(0)+2\pi ikt$,
так что $a+e^{M(t)}=c_b^k(t)$. Положим
\[
H(s,t)=a+\exp\bigl((1-s)L(t)+sM(t)\bigr).
\]
Так как $\Pi$ выпукла, $(1-s)L(t)+sM(t)\in\Pi$ и $H(s,t)\in D_1$; $H$ непрерывна, $H(0,\cdot)=\gamma$, $H(1,\cdot)=c_b^k$, $H(s,0)=a+e^{L(0)}=b$,
$H(s,1)=a+e^{L(0)+2\pi ik}=b$. Для $\{|z|>R\}$ берём $\Pi=\{\Re w>\ln R\}$, $a=0$.
\end{proof}

\textbf{Оператор обхода.} Пусть $\Phi_1$ --- аналитическая ветвь над $D_1=V(a,\delta)$ (или $\{|z|>R\}$, $a=\infty$; тогда индексы берутся относительно $0$).
Для $c\in D_1$ определим $T_c\colon\Phi_1(c)\to\Phi_1(c)$ --- продолжение вдоль $c_c$ (утверждение 27.2(а)). Это биекция (обратная --- продолжение вдоль $c_c^-$).

\begin{thmbox}{Утверждение 27.6}
1) Продолжение $G\in\Phi_1(c)$ вдоль любого замкнутого пути $\gamma$ в $D_1$ с началом $c$ равно $T_c^{k}G$, $k=\Ind_a(\gamma)$; в частности, «обход» корректно
определён, и $\Phi_1(c)=\{T_c^kG:k\in\ZZ\}$ для любого $G\in\Phi_1(c)$ (одна орбита).

2) Если $\sigma$ --- путь в $D_1$ из $b$ в $c$, то $T_c\circ\Lambda_\sigma=\Lambda_\sigma\circ T_b$ (продолжение вдоль $\sigma$ сопрягает операторы обхода).

3) Следовательно, число $n=|\Phi_1(c)|$ --- длина единственной орбиты --- не зависит ни от $c$, ни от выбора элемента, и три случая определения 15.9
--- это $n=1$, $2\le n<\infty$ (ветвление порядка $n$; $T_c^jG\ne G$ при $0<j<n$, $T_c^nG=G$ для любого $G$), $n=\infty$ (логарифмическое ветвление;
$T_c^jG\ne G$ при всех $j\ne0$). Для $a=\infty$ смена направления обхода заменяет $T$ на $T^{-1}$ и ничего не меняет.
\end{thmbox}
\begin{proof}
1) По лемме 27.5 $\gamma$ гомотопен $c_c^k$ в $D_1$, а $G$ продолжается вдоль всех путей гомотопии (они в $D_1$). По теореме 15.1 результаты совпадают, а
продолжение вдоль $c_c^k$ есть $T_c^kG$ (лемма 25.A(4) и (2)). Все элементы $\Phi_1(c)$ получаются из $G$ вдоль замкнутых путей в $D_1$ (утверждение 27.2(а)).

2) Путь $\sigma c_c\sigma^-$ замкнут, начинается в $b$ и имеет $\Ind_a=\Ind_a(c_c)=1$ (изменения аргумента вдоль $\sigma$ и $\sigma^-$ взаимно уничтожаются).
По п.~1 продолжение $G\in\Phi_1(b)$ вдоль $\sigma c_c\sigma^-$ равно $T_bG$. Продолжив затем ещё вдоль $\sigma$, получим $\Lambda_\sigma T_bG$; с другой стороны,
продолжение вдоль $\sigma c_c\sigma^-\sigma$ равно $T_c\Lambda_\sigma G$, поскольку участок $\sigma^-\sigma$ возвращает элемент (лемма 25.A(2)).

3) Из 1) $\Phi_1(c)$ --- одна орбита биекции $T_c$, и все её элементы имеют одинаковый период $n$ (если $T^nG=G$, то $T^n(T^jG)=T^j(T^nG)=T^jG$).
Из 2) $\Lambda_\sigma$ переводит орбиту $T_b$ в орбиту $T_c$ с тем же периодом.
\end{proof}

\begin{thmbox}{Утверждение 27.7 (случай 1 определения 15.9; слайды 375--376)}
Если $T_bF_b=F_b$, то существует $f_1\in H(D_1)$ с $\Phi_1=\{E_c(f_1):c\in D_1\}$ ($\Phi_1$ распадается на одну голоморфную ветвь).
\end{thmbox}
\begin{proof}
По утверждению 27.6(1) продолжение $F_b$ вдоль любого замкнутого пути в $D_1$ возвращает $F_b$. Значит, результат продолжения $F_b$ в точку $c$ не зависит от
пути: если $\sigma_1,\sigma_2$ --- пути из $b$ в $c$ с результатами $H_1,H_2$, то продолжение $F_b$ вдоль замкнутого $\sigma_1\sigma_2^-$ равно $F_b$, т.е.
$\Lambda_{\sigma_2^-}H_1=F_b$, откуда $H_1=\Lambda_{\sigma_2}F_b=H_2$. Остаётся применить лемму 26.A.
\end{proof}
\begin{remarkbox}{Замечание}
На слайде это получено иначе: теорема о монодромии в двух односвязных областях $D_1\setminus\ell_1$, $D_1\setminus\ell_2$ (разрезы по двум разным радиусам, не
проходящим через $b$) даёт две функции, совпадающие на секторе, содержащем $b$; условие «обход возвращает $F_b$» обеспечивает их совпадение и на втором секторе.
(Проколотый круг с разрезом по радиусу $[a,a+\delta e^{i\theta})$ звёздный относительно точки $a-\frac\delta2e^{i\theta}$, а значит, односвязен.)
\end{remarkbox}

\subsection{Действия над П.А.Ф. и сохранение тождеств}

\begin{thmbox}{Лемма 27.8 (операции коммутируют с продолжением)}
Пусть $\{F_t=(B(\gamma(t),R_t),f_t)\}$ и $\{G_t=(B(\gamma(t),\widetilde R_t),g_t)\}$ --- продолжения элементов $F_0,G_0$ вдоль $\gamma$. Тогда семейства
\[
E_{\gamma(t)}(f_t+g_t),\quad E_{\gamma(t)}(f_tg_t),\quad E_{\gamma(t)}(f_t'),\quad E_{\gamma(t)}(\Psi\circ f_t)\ (\Psi\text{ --- целая функция})
\]
--- продолжения соответствующих элементов вдоль $\gamma$, с радиусами $\ge\min(R_t,\widetilde R_t)$ (для $f_t'$ и $\Psi\circ f_t$ --- $\ge R_t$).
В частности, если $F_0$ переходит в $F_1$, то $\pm F_0+C$ переходит в $\pm F_1+C$ ($C$ --- константа), а $C F_0$ --- в $CF_1$.
\end{thmbox}
\begin{proof}
Пусть $h_t$ --- одна из функций $f_t+g_t,\dots$; она голоморфна в $B(\gamma(t),r_t)$, $r_t=\min(R_t,\widetilde R_t)$, и обладает «локальностью»: если $f_t=f_\xi$ и
$g_t=g_\xi$ в окрестности точки $w$, то $h_t=h_\xi$ около $w$. Пусть $H_t=E_{\gamma(t)}(h_t)$; его радиус $\ge r_t>0$. Для $t\in I$ возьмём промежуток
$W\subset U_t\cap\widetilde U_t$, $W\ni t$, на котором $|\gamma(\xi)-\gamma(t)|<r_t$. При $\xi\in W$ точка $w=\gamma(\xi)$ лежит в $B_t\cap\widetilde B_t\cap B_\xi\cap\widetilde B_\xi$, и
около неё $f_\xi=f_t$, $g_\xi=g_t$, значит $h_\xi=h_t$. Так как $w$ лежит в кругах $H_t$ и $H_\xi$, а их пересечение выпукло, $H_\xi$ --- НАП $H_t$.
\end{proof}

Так, на слайде 377: если $\Phi,\Psi$ --- П.А.Ф. в $D$, $F,G$ --- их элементы с центром $a$, то $E_a(f')$, $E_a(f+g)$, $E_a(fg)$ продолжаются вдоль всех путей в $D$
и определяют совокупности $\Phi'$, $\Phi+\Psi$, $\Phi\Psi$, которые могут распадаться на несколько П.А.Ф. (так как над точкой $a$ можно комбинировать разные пары элементов).

\begin{remarkbox}{Примеры (слайд 377) с проверкой}
1) $\sqrt z+\sqrt z$: над точкой $1$ элементы $\pm\sqrt z$ (главная ветвь); суммы пар: $2\sqrt z$, $-2\sqrt z$ (одна П.А.Ф. $2\sqrt z$) и $\sqrt z-\sqrt z=0$. Итого две П.А.Ф.:
$2\sqrt z$ и $0$.

2) $i\Ln z+\Ln z$: над $1$ элементы $\Ln z$ --- это $\ln z+2\pi ik$, и суммы имеют вид $i(\ln z+2\pi ik)+(\ln z+2\pi im)=(1+i)\ln z+C_{k,m}$,
$C_{k,m}=-2\pi k+2\pi im$. По лемме 27.8 обход нуля прибавляет к первому слагаемому $i\cdot2\pi i=-2\pi$, ко второму $2\pi i$, т.е. переводит $(k,m)$ в $(k+1,m+1)$.
Поэтому элементы с $(k,m)$ и $(k',m')$ лежат в одной П.А.Ф. тогда и только тогда, когда $m-k=m'-k'$. Получаем счётное число различных П.А.Ф.
$(1+i)\Ln z+2\pi ij$, $j=m-k\in\ZZ$ (при $j\ne j'$ разность $2\pi i(j-j')$ не имеет вида $2\pi N(i-1)$).

3) Композиция определяется через теорему о продолжении сложного элемента (билет 29). $w^2\circ\sqrt[4]z$: элементы $(i^k\sqrt[4]z)^2=\pm\sqrt z$ --- П.А.Ф. $\sqrt z$. А произведение
$\sqrt[4]z\cdot\sqrt[4]z$: пары $i^k\sqrt[4]z\cdot i^m\sqrt[4]z=i^{k+m}\sqrt z$ дают две П.А.Ф.: $\sqrt z$ ($k+m$ чётно) и $i\sqrt z$ ($k+m$ нечётно); так что композиция $\ne$ произведение.

4) Упражнение слайда: $\sqrt[4]w\circ z^2$. Внутренняя функция $z^2$ целая; внешний элемент --- любой из $i^k\sqrt[4]w_{(\circ)}$ около $w=1$. Вдоль пути $\gamma$ из $1$ продолжение
$\sqrt[4]w$ вдоль $\Gamma=\gamma^2$ возможно $\iff$ $\Gamma$ не проходит через $0$ $\iff$ $\gamma$ не проходит через $0$, и функция результата около $\gamma(t)$ равна
$i^k\exp\bigl(\frac14\ln^{(2\varphi(t))}(z^2)\bigr)$, где $\varphi$ --- непрерывный аргумент $\gamma$ (а $2\varphi$ --- непрерывный аргумент $\gamma^2$). Около $\gamma(t)$
$\ln^{(2\varphi(t))}(z^2)=2\ln^{(\varphi(t))}z$ (оба --- логарифмы $z^2$ с мнимой частью около $2\varphi(t)$), так что это $i^k\exp(\frac12\ln^{(\varphi(t))}z)=i^k\sqrt z$ (ветвь вдоль $\gamma$).
Итак, $\sqrt[4]w\circ z^2$ распадается на две П.А.Ф.: $\sqrt z$ (при $k=0,2$) и $i\sqrt z$ (при $k=1,3$), с $D=\CC^*$, ветвлениями второго порядка в $0$ и $\infty$.
\end{remarkbox}

\begin{thmbox}{Лемма 27.9 (сохранение тождеств; непродолжаемость через точку)}
Пусть $Q(z,u,v)$ таково, что для любой голоморфной в круге функции $f$ функция $z\mapsto Q(z,f(z),f'(z))$ голоморфна в этом круге (например, $Q$ --- многочлен
от $u,v$ с целыми по $z$ коэффициентами, или $Q=\Psi(u)-z$ с целой $\Psi$; ниже используются $v(1+z^2)-1$, $v(z-w)-1$, $v(z-w)-\lambda u$, $u^2-2zu+1$, $v^2(1-z^2)-1$, $\sin u-z$). Если для элемента $F=(B,f)$ выполнено $Q(z,f(z),f'(z))\equiv0$ в $B$, то это же
тождество выполнено для каждого элемента П.А.Ф., порождённой $F$. Если, кроме того, ни одна функция $g$, голоморфная в окрестности точки $p$, не может
удовлетворять $Q(z,g,g')\equiv0$ около $p$, то $p\notin D_\Phi$, и $F$ не продолжается ни вдоль какого пути, проходящего через $p$; радиус элемента $\Phi$ с центром $c$
не превосходит $|c-p|$.
\end{thmbox}
\begin{proof}
Если $G=(B',g)$ --- НАП $F$, то $g=f$ на непустом открытом $B\cap B'$, так что голоморфная в $B'$ функция $Q(z,g,g')$ обращается в нуль на открытом множестве, а значит,
во всём круге $B'$. По предложению 15.3(1) любой элемент $\Phi$ соединён с $F$ цепочкой НАП. Если $p\in D_\Phi$, то над $p$ есть элемент (утверждение 27.1(2)) с
функцией, голоморфной около $p$ и удовлетворяющей тождеству, --- противоречие; путь через $p$ давал бы элемент с центром $p$. Если радиус элемента с центром $c$
больше $|c-p|$, то $p$ лежит в его круге, т.е. в $D_\Phi$.
\end{proof}

\begin{thmbox}{Лемма 27.10 (продолжение первообразной; ср. предложение 15.5, билет 28)}
Пусть элемент $E_a(f')$ продолжается вдоль $\gamma$ семейством $G_t=(B(\gamma(t),r_t),g_t)$. Тогда $F_0=E_a(f)$ продолжается вдоль $\gamma$ семейством $F_t=(B(\gamma(t),r_t),f_t)$
с $f_t'=g_t$ (те же радиусы). Если $\gamma$ кусочно-гладкий, то $f_1(\gamma(1))-f_0(\gamma(0))=\int_0^1g_t(\gamma(t))\gamma'(t)\,dt$. В частности, если $G_1=G_0$ (путь замкнут), то
$F_1=F_0+C$, $C=\int_\gamma g$.
\end{thmbox}
\begin{proof}
Как в доказательстве предложения 15.3(1), выберем разбиение $0=t_0<\dots<t_N=1$ и точки $\tau_k$ с $[t_{k-1},t_k]\subset U_{\tau_k}$ (окрестности для $\{G_t\}$); тогда
$\gamma([t_{k-1},t_k])\subset B_k:=B(\gamma(\tau_k),r_{\tau_k})$, и при $s\in[t_{k-1},t_k]$ функция $g_{\tau_k}$ совпадает с $g_s$ около $\gamma(s)$. Пусть $P_k$ --- первообразная $g_{\tau_k}$ в круге $B_k$;
константы выберем по индукции: $P_1=f$ около $a$ (обе --- первообразные $g_{\tau_1}=g_0$ около $a$), а $P_k=P_{k-1}$ около $\gamma(t_{k-1})$ (там $g_{\tau_k}=g_{t_{k-1}}=g_{\tau_{k-1}}$, так что
$P_k-P_{k-1}$ --- константа около $\gamma(t_{k-1})$). При $s\in[t_{k-1},t_k]$ положим $F_s=E_{\gamma(s)}(P_k)$ (на стыках определения согласованы). По лемме 25.B на каждом отрезке
это продолжение; склейка (лемма 25.A(4)) даёт продолжение вдоль $\gamma$. Ряд производной $P_k'=g_{\tau_k}$ в $\gamma(s)$ --- это $g_s$ (единственность продолжения $\{G_t\}$, лемма 25.B),
а ряд функции и ряд её производной имеют один радиус сходимости; поэтому радиус $F_s$ равен $r_s$ и $f_s'=g_s$. Формула для приращения:
$P_k(\gamma(t_k))-P_k(\gamma(t_{k-1}))=\int_{t_{k-1}}^{t_k}g_{\tau_k}(\gamma(t))\gamma'(t)\,dt$ и суммирование по $k$. Если $G_1=G_0$, то $f_1'=f_0'$, т.е. $f_1-f_0$ --- константа.
\end{proof}

\subsection{Поведение П.А.Ф. около точки ветвления конечного порядка}

Нам понадобится (без доказательства, см. билет 29) \textbf{теорема о продолжении сложного элемента (предложение 15.6)}: если $F_a$ продолжается вдоль $\gamma$
семейством $\{F_t=(B_t,f_t)\}$, $\Gamma(t)=f_t(\gamma(t))$, и элемент $G_A$, $A=\Gamma(0)$, продолжается вдоль пути $\Gamma$ семейством $\{G_t=(\widetilde B_t,g_t)\}$, то $E_a(g_A\circ f_a)$
продолжается вдоль $\gamma$ семейством $E_{\gamma(t)}(g_t\circ f_t)$. Мы применяем её с целой внутренней функцией $\phi$ (тогда $F_t=E_{\gamma(t)}(\phi)$, $\Gamma=\phi\circ\gamma$).

Пусть $\Phi_1$ --- аналитическая ветвь над $V=V(a,\eps)$ с ветвлением порядка $n\ge2$ в $a\in\CC$. Положим $W=V(0,\eps^{1/n})$, $\phi(w)=a+w^n$, так что $\phi\colon W\to V$,
$\omega_n=e^{2\pi i/n}$.

\begin{thmbox}{Предложение 15.8}
$\Phi_1\circ\phi$ распадается над $W$ на $n$ различных голоморфных функций $f_0,\dots,f_{n-1}$, причём $f_j(w)=f_0(\omega_n^jw)$. Точнее: для любого $w\in W$ и любого
$G=(B,g)\in\Phi_1(\phi(w))$ элемент $E_w(g\circ\phi)$ равен $E_w(f_j)$ при некотором $j$. Обратно, $\Phi_1$ восстанавливается по $f_0$: $\Phi_1=f_0\circ\phi^{-1}$, т.е.
$\Phi_1(z)=\{E_z(f_0\circ\psi)\}$, где $\psi$ пробегает $n$ голоморфных ветвей $\phi^{-1}=a+(z-a)^{1/n}$ около $z$.
\end{thmbox}
\begin{proof}
\emph{Подготовка.} Фиксируем $w_0\in W$, $z_0=\phi(w_0)$, $F_0\in\Phi_1(z_0)$ и $F_j:=T_{z_0}^jF_0$, $j=0,\dots,n-1$; по утверждению 27.6 это все $n$ (различных) элементов над $z_0$,
$F_j=(B_j,h_j)$. Пусть $G_j=E_{w_0}(h_j\circ\phi)$.

\emph{Индексы.} Если $\theta$ --- непрерывный аргумент пути $\sigma$ в $W$, то $n\theta$ --- непрерывный аргумент $\phi(\sigma(t))-a=\sigma(t)^n$. Поэтому для замкнутого $\sigma$
$\Ind_a(\phi\circ\sigma)=n\Ind_0(\sigma)$, а для дуги $\sigma_j$ окружности $|w|=|w_0|$ от $w_0$ до $\omega_n^jw_0$ (против часовой стрелки) путь $\phi\circ\sigma_j$ замкнут и $\Ind_a(\phi\circ\sigma_j)=j$.

\emph{$G_j$ задаёт функцию $f_j\in H(W)$.} Для пути $\sigma$ в $W$ из $w_0$ путь $\phi\circ\sigma$ лежит в $V$, и $F_j$ продолжается вдоль него (утверждение 27.2(а)); по теореме о сложном элементе
$G_j$ продолжается вдоль $\sigma$. Если $\sigma$ замкнут, $\Ind_0\sigma=k$, то $F_j$ вдоль $\phi\circ\sigma$ переходит в $T^{nk}F_j=F_j$, и $G_j$ переходит в себя. Как в доказательстве
утверждения 27.7, результат продолжения $G_j$ не зависит от пути, и лемма 26.A даёт $f_j\in H(W)$ с $f_j=h_j\circ\phi$ около $w_0$.

\emph{Различность.} Если $f_j=f_k$, то $h_j\circ\phi=h_k\circ\phi$ около $w_0$; так как $\phi'(w_0)=nw_0^{n-1}\ne0$, $\phi$ отображает окрестность $w_0$ на открытую окрестность $z_0$ (теорема
об открытости), и $h_j=h_k$ около $z_0$, т.е. $F_j=F_k$, $j=k$.

\emph{Все элементы $\Phi_1\circ\phi$.} Пусть $w\in W$, $G\in\Phi_1(\phi(w))$, $\sigma$ --- путь в $W$ из $w_0$ в $w$. Продолжив $G$ вдоль $(\phi\circ\sigma)^-$, получим некоторый $F_j$; значит, $G$ --- продолжение
$F_j$ вдоль $\phi\circ\sigma$, и по теореме о сложном элементе $E_w(g\circ\phi)$ --- продолжение $G_j$ вдоль $\sigma$, т.е. $E_w(f_j)$.

\emph{Формула $f_j(w)=f_0(\omega_n^jw)$.} Продолжение $F_0$ вдоль $\phi\circ\sigma_j$ (индекс $j$) есть $F_j$; поэтому продолжение $G_0$ вдоль $\sigma_j$ равно $E_{\omega_n^jw_0}(h_j\circ\phi)$, а с другой
стороны, $E_{\omega_n^jw_0}(f_0)$. Значит, $f_0=h_j\circ\phi$ около $\omega_n^jw_0$. Функция $\tilde f(w):=f_0(\omega_n^jw)$ голоморфна в $W$, и около $w_0$:
$\tilde f(w)=h_j(\phi(\omega_n^jw))=h_j(\phi(w))=f_j(w)$ (так как $\phi(\omega_n^jw)=\phi(w)$). По теореме единственности $\tilde f=f_j$ в $W$.

\emph{$\Phi_1=f_0\circ\phi^{-1}$.} Пусть $z\in V$, $\psi$ --- ветвь $\phi^{-1}$ около $z$, $w=\psi(z)$. По доказанному $f_0=g\circ\phi$ около $w$ для некоторого $G=(B,g)\in\Phi_1(z)$ (возьмём $G$ ---
продолжение $F_0$ вдоль $\phi\circ\sigma$, $\sigma$ из $w_0$ в $w$); тогда $f_0\circ\psi=g$ около $z$, т.е. $E_z(f_0\circ\psi)\in\Phi_1(z)$. Различные ветви $\psi$ (их $n$: $\omega_n^j\psi$) дают
различные элементы, так как $f_0(\omega_n^j\psi)=f_j\circ\psi$, а $f_j$ различны; поскольку $|\Phi_1(z)|=n$, получены все элементы.
\end{proof}

\begin{remarkbox}{Следствие (ряды Пюизё)}
Разложив $f_0$ в ряд Лорана в $W$: $f_0(w)=\sum_{k\in\ZZ}c_kw^k$, получаем, что все элементы $\Phi_1$ над $V$ имеют вид
$\sum_{k\in\ZZ}c_k\bigl((z-a)^{1/n}_{(*)}\bigr)^k$ с разными ветвями корня. Особенность $f_0$ в нуле (устранимая/полюс/существенная) характеризует поведение $\Phi_1$ около $a$.
\end{remarkbox}

\subsection{Схема полного описания П.А.Ф. (слайд 380)}
\begin{enumerate}
\item Выбрать исходный элемент $F\in\Phi$ и обосновать его каноничность (найти радиус). Помнить, что «формула» может задавать не одну П.А.Ф.
(например, $\sqrt{z^2}$ задаёт две: $z$ и $-z$).
\item Определить, вдоль каких путей $F$ продолжается, а вдоль каких --- нет; для каждого допустимого пути указать продолжающее семейство; найти $D_\Phi$.
\item Найти число элементов $\Phi$ над произвольной точкой $D_\Phi$ и связь между ними.
\item Найти и классифицировать все изолированные особые точки всех ветвей $\Phi$, указав, сколько ветвей имеют особенность данного типа над данной точкой.
\end{enumerate}

\subsection{Примеры описания П.А.Ф. (слайды 382--399)}

Обозначения: $\ln^{(\alpha)}\zeta=\ln|\zeta|+i\arg^{(\alpha)}\zeta$, $\arg^{(\alpha)}\zeta\in(\alpha-\pi,\alpha+\pi)$ --- ветвь логарифма, голоморфная в $\CC\setminus\RR_{\alpha+\pi}$;
$(z-w)^\lambda_{(\alpha)}=\exp(\lambda\ln^{(\alpha)}(z-w))$; $\ln=\ln^{(0)}$ --- главная ветвь; $\CC^*_w=\CC\setminus\{w\}$.

\subsubsection*{1. П.А.Ф. $\Ln(z-w)$}

\emph{Исходный элемент.} Для $a\in\CC^*_w$ пусть $\alpha(a)=\arg(a-w)\in(-\pi,\pi]$ и $L_w(a)=E_a\bigl(\ln^{(\alpha(a))}(z-w)\bigr)$. Его радиус равен $|a-w|$:
круг $B(a,|a-w|)$ не пересекает луч $w+\RR_{\alpha+\pi}=\{w-se^{i\alpha}:s\ge0\}$ (расстояние от $a=w+|a-w|e^{i\alpha}$ до $w-se^{i\alpha}$ равно $|a-w|+s$), так что радиус
$\ge|a-w|$; а по лемме 27.9 с $Q=v(z-w)-1$ (функция $g$, голоморфная около $w$, дала бы $0=g'(w)\cdot0=1$) радиус $\le|a-w|$ и $w\notin D_\Phi$.

\emph{Продолжение.} Пусть $\gamma\colon I\to\CC^*_w$, $\gamma(0)=a$, $\varphi$ --- непрерывная ветвь $\Arg(\gamma(t)-w)$ с $\varphi(0)=\alpha(a)$. Тогда продолжение $L_w(a)$ вдоль $\gamma$ ---
\[
L_{w,\gamma,t}=\bigl(B(\gamma(t),|\gamma(t)-w|),\ \ell_{w,\gamma,t}\bigr),\qquad \ell_{w,\gamma,t}(z)=\ln^{(\varphi(t))}(z-w).
\]
\begin{detail}{проверка определения 15.4}
Каноничность и радиус --- как для исходного элемента ($\varphi(t)$ --- аргумент $\gamma(t)-w$). Для $t\in I$ возьмём промежуток $U_t$, на котором
$|\gamma(\xi)-\gamma(t)|<|\gamma(t)-w|$ и $|\varphi(\xi)-\varphi(t)|<\pi$. Тогда $\gamma(\xi)\in B_t$. На выпуклом множестве $B_t\cap B_\xi\not\ni w$ функции $\ell_{w,\gamma,t}$ и $\ell_{w,\gamma,\xi}$ ---
логарифмы $z-w$, их разность непрерывна и принимает значения в $2\pi i\ZZ$, т.е. постоянна. В точке $\gamma(\xi)$: $\Im\ell_{w,\gamma,\xi}=\varphi(\xi)$, а $\Im\ell_{w,\gamma,t}(\gamma(\xi))$ ---
аргумент $\gamma(\xi)-w$ из $(\varphi(t)-\pi,\varphi(t)+\pi)$, и это тоже $\varphi(\xi)$. Разность равна нулю: $L_{w,\gamma,\xi}$ --- НАП $L_{w,\gamma,t}$.
\end{detail}
Итак, $R_\gamma=\min_t|\gamma(t)-w|$; вдоль путей через $w$ продолжения нет; $D_\Phi=\CC^*_w$, и $\Phi$ --- своя (единственная) аналитическая ветвь над $\CC^*_w$.

\emph{Элементы над точкой.} Для замкнутого $\gamma$ с началом $b$: $\varphi(1)=\varphi(0)+2\pi\Ind_w(\gamma)$, поэтому $L^k_w(b):=L_w(b)+2\pi ik$ переходит в
$L^k_w(b)+2\pi i\Ind_w(\gamma)$. Над каждой $b\ne w$ ровно счётное множество элементов $\{L_w(b)+2\pi ik:k\in\ZZ\}$ (все различны).

\emph{Особые точки.} $D_1=\CC^*_w=V(w,\infty)$, а также $\{|z|>|w|\}$ --- проколотая окрестность $\infty$ (обход $\infty$ имеет $\Ind_w=\Ind_0=\pm1$). Обход прибавляет $\pm2\pi i\ne0$, ни одна степень
$T$ не тождественна: $w$ и $\infty$ --- логарифмические точки ветвления.

\subsubsection*{2. П.А.Ф. $(z-w)^\lambda$, $\lambda\in\CC\setminus\ZZ$}

Исходный элемент $P^\lambda_w(a)=E_a\bigl(\exp(\lambda\ln^{(\alpha(a))}(z-w))\bigr)$. По лемме 27.8 (с целой $\Psi(u)=e^{\lambda u}$) продолжение вдоль $\gamma\subset\CC^*_w$ ---
семейство $P^\lambda_{w,\gamma,t}=E_{\gamma(t)}\bigl(\exp(\lambda\ell_{w,\gamma,t})\bigr)$, причём окрестности $U_t$ можно брать те же, что для $\Ln$, а радиусы $\ge|\gamma(t)-w|$.

\begin{detail}{радиус равен $|\gamma(t)-w|$ и через $w$ продолжения нет}
Все элементы удовлетворяют $f'(z)(z-w)=\lambda f(z)$ (лемма 27.9) и $f\not\equiv0$. Пусть $g\not\equiv0$ голоморфна около $w$ и удовлетворяет этому тождеству; $g=(z-w)^mh$, $h(w)\ne0$, $m\ge0$.
Подстановка: $(z-w)^m\bigl[(m-\lambda)h+(z-w)h'\bigr]\equiv0$, откуда $(m-\lambda)h(w)=0$, $\lambda=m\in\ZZ$ --- противоречие. (Функция элемента с центром $w$ или с кругом,
содержащим $w$, была бы такой $g$: она $\not\equiv0$, так как совпадает на открытом множестве с экспонентой.)
\end{detail}

\emph{Элементы над точкой.} Обход по замкнутому $\gamma$ с $N=\Ind_w(\gamma)$ умножает функцию на $e^{2\pi i\lambda N}$ (так как $\ell$ увеличивается на $2\pi iN$), круг не меняется.
Над $b\ne w$ элементы: $\{P^\lambda_w(b)e^{2\pi i\lambda k}:k\in\ZZ\}$, различные при различных значениях множителя.
\begin{itemize}
\item $\lambda=p/q$ несократимая дробь, $q\ge2$: $e^{2\pi ipk/q}$ зависит лишь от $k\bmod q$ и принимает $q$ различных значений ($p$ взаимно просто с $q$). Над каждой точкой $q$ элементов;
$w$ и $\infty$ --- точки ветвления порядка $q$. В частности, при $\lambda=1/n$ ($\sqrt[n]{z-w}$) --- ровно $n$ элементов и ветвления порядка $n$ в $w$ и $\infty$.
\item $\lambda\notin\QQ$ (в том числе невещественное): $e^{2\pi i\lambda k}=e^{2\pi i\lambda k'}\iff\lambda(k-k')\in\ZZ\iff k=k'$. Над каждой точкой счётное множество элементов; $w$ и $\infty$ ---
логарифмические точки ветвления.
\item Упражнение слайда 386: $\dfrac1{(z-w)^\pi}=(z-w)^{-\pi}$, $\lambda=-\pi\notin\QQ$: $w$ и $\infty$ --- логарифмические точки ветвления ($D_\Phi=\CC^*_w$, над каждой точкой счётно много элементов,
отличающихся множителями $e^{-2\pi^2ik}$).
\end{itemize}

\subsubsection*{3. П.А.Ф. $\Arctg z$}

\emph{Исходный элемент} $F_0=E_0(\operatorname{arctg}z)$, $\operatorname{arctg}z=\sum_{n\ge0}\frac{(-1)^n}{2n+1}z^{2n+1}$, радиус $1$. Так как
$(\operatorname{arctg}z)'=\frac1{1+z^2}=\frac i2\bigl(\frac1{z+i}-\frac1{z-i}\bigr)$, по теореме единственности в $B(0,1)$
\[
f_0(z)=\tfrac i2\ln^{(\pi/2)}(z+i)-\tfrac i2\ln^{(-\pi/2)}(z-i)+\tfrac\pi2
\]
(производные совпадают, и при $z=0$: $\frac i2\cdot\frac{i\pi}2-\frac i2\cdot\bigl(-\frac{i\pi}2\bigr)+\frac\pi2=0$). Здесь $\ln^{(\pi/2)}(z+i)$ и $\ln^{(-\pi/2)}(z-i)$ --- функции элементов $L_{-i}(0)$, $L_i(0)$.

\emph{Продолжение.} По лемме 27.8 вдоль любого $\gamma\subset\CC\setminus\{\pm i\}$, $\gamma(0)=0$, $F_0$ продолжается семейством с функциями
$f_t=\frac i2\ell_{-i,\gamma,t}-\frac i2\ell_{i,\gamma,t}+\frac\pi2$ и радиусами $r_t=\min\{|\gamma(t)+i|,|\gamma(t)-i|\}$: радиус $\ge r_t$ по лемме 27.8 и $\le r_t$ по лемме 27.9
с $Q=v(1+z^2)-1$ (в точке $\pm i$ получилось бы $0=1$). По той же лемме $F_0$ \emph{не продолжается вдоль путей через $\pm i$} (упражнение слайда 387). $D_\Phi=\CC\setminus\{\pm i\}$,
$\Phi$ --- аналитическая ветвь над $D_\Phi$.

\emph{Элементы над $b$.} По формуле для $\Ln$ обход по замкнутому $\gamma$ с $N_{\pm i}=\Ind_{\pm i}(\gamma)$ переводит $F_b$ в
$F_b+\frac i2\cdot2\pi iN_{-i}-\frac i2\cdot2\pi iN_i=F_b-\pi N_{-i}+\pi N_i$. Так как $N_{\pm i}$ принимают любые значения (например, обход одной $i$), над $b$ ровно
счётное множество элементов $F_b+\pi k$, $k\in\ZZ$, с общим кругом $B(b,r_b)$.

\emph{Особые точки.} Над $V(\pm i,2)$ (там $\Ind_{\mp i}=0$ для замкнутых путей, так как $\mp i\notin B(\pm i,2)$) обход сдвигает на $\pm\pi$ --- логарифмическое ветвление, по одному над каждой
точкой ($\Phi|_{V(\pm i,2)}$ --- одна аналитическая ветвь, т.к. все сдвиги $\pi k$ достигаются). Над $W=\{|z|>1\}$: обход $\infty$ обходит $i$ и $-i$ по одному разу в одну сторону, сдвиг
$-\pi+\pi=0$ --- случай 1; поэтому $\Phi$ над $W$ распадается на счётное число голоморфных ветвей $f_j$, $f_j-f_s=\pi(j-s)$ (каждый элемент над $b\in W$ порождает свою
аналитическую ветвь, и это голоморфная функция по утверждению 27.7). Каждая $f_j$ имеет в $\infty$ устранимую особенность: $|f_j'(z)|=\frac1{|z^2+1|}\le\frac1{|z|^2-1}$, и,
интегрируя $f_j'$ от точки $2$ по дуге окружности $|\zeta|=2$ (длина $\le4\pi$, $|f'_j|\le\frac13$) и затем по радиусу ($\int_2^\infty\frac{dr}{r^2-1}<\infty$), получаем ограниченность
$f_j$ в $\{|z|\ge2\}$; по теореме Римана особенность устранима.
\begin{remarkbox}{Замечание}
Главная ветвь $\operatorname{arctg}x\to\pm\frac\pi2$ при $x\to\pm\infty$ не противоречит устранимости: главная ветвь голоморфна в $\CC\setminus i\bigl((-\infty,-1]\cup[1,\infty)\bigr)$, и на $W$ это
не одна голоморфная ветвь --- при переходе через разрез $i[1,\infty)$ голоморфная в $W$ ветвь отличается от главной на $\pi$; для неё значения при $x\to+\infty$ и $x\to-\infty$ (через верхнюю
полуплоскость) равны $\frac\pi2$ и $-\frac\pi2+\pi=\frac\pi2$.
\end{remarkbox}

\subsubsection*{4. П.А.Ф. $\Zh^{-1}(z)$}

\emph{Исходный элемент} $F_0=(B(0,1),f)$, $f=\Zh^{-1}_+|_{B(0,1)}$, где $\Zh^{-1}_+$ --- обратная к $\Zh(w)=\frac12(w+\frac1w)$ на верхней полуплоскости, голоморфная в
$\Omega_0=\CC\setminus((-\infty,-1]\cup[1,\infty))$, $\Zh^{-1}_+(0)=i$. Уравнение $\Zh(w)=z$ равносильно $w^2-2zw+1=0$, $w\in z+\sqrt{z^2-1}$, и
\[
f(z)=z+(z+1)^{1/2}_{(0)}(z-1)^{1/2}_{(\pi)}\qquad(z\in\Omega_0).
\]
\begin{detail}{формула для $f$}
Правая часть $\hat f$ голоморфна в $\Omega_0$ (первый корень --- вне $(-\infty,-1]$, второй --- вне $[1,\infty)$), $(\hat f-z)^2=z^2-1$, т.е. $\hat f$ --- корень уравнения, и $\hat f(0)=1\cdot e^{i\pi/2}=i$.
Множество $\{z\in\Omega_0:\hat f(z)=\Zh^{-1}_+(z)\}$ замкнуто в $\Omega_0$ и непусто; оно открыто: корни уравнения различны при $z^2\ne1$, т.е. отделены друг от друга, и две непрерывные
функции-корни, совпавшие в точке, совпадают и рядом. По связности $\hat f=\Zh^{-1}_+$ в $\Omega_0$.
\end{detail}
Радиус $F_0$ равен $1$; более общо --- лемма 27.9 с $Q=u^2-2zu+1$: если $g$ голоморфна около $1$ и $g^2-2zg+1\equiv0$, то $g(1)=1$, а дифференцирование даёт
$2gg'-2g-2zg'\equiv0$, в точке $1$: $2g'(1)-2-2g'(1)=-2\ne0$. Аналогично в $-1$ ($g(-1)=-1$, получаем $2\ne0$). Значит, $\pm1\notin D_\Phi$, \emph{через $\pm1$ продолжать нельзя}
(упражнение слайда 389), радиус элемента с центром $c$ не больше $\min|c\mp1|$.

\emph{Продолжение.} Вдоль $\gamma\subset D=\CC\setminus\{\pm1\}$, $\gamma(0)=0$, по лемме 27.8 (сумма и произведение продолжений $P^{1/2}_{-1}(0)$, $P^{1/2}_1(0)$ и тождественной функции)
$F_0$ продолжается семейством с функциями $f_t(z)=z+(z+1)^{1/2}_{\gamma,t}(z-1)^{1/2}_{\gamma,t}$ и радиусами $\ge R_t=\min\{|\gamma(t)-1|,|\gamma(t)+1|\}$, т.е. ровно $R_t$
(проверка деталей слайда 389). $D_\Phi=\CC\setminus\{\pm1\}$, $\Phi$ --- аналитическая ветвь над $D_\Phi$.

\emph{Элементы над $b$.} Обход по замкнутому $\gamma$ умножает $(z\mp1)^{1/2}$ на $(-1)^{N_{\pm1}}$, $N_{\pm1}=\Ind_{\pm1}\gamma$. Пусть $F_b=(B(b,R_b),f_b)$, $f_b=z+s_b$, $s_b^2=z^2-1$.
После обхода --- $z+(-1)^{N_1+N_{-1}}s_b$. Итак, над $b$ ровно два элемента: $F_b$ и $\widetilde F_b=(B(b,R_b),z-s_b)$ (различны, так как $s_b\not\equiv0$), и
\[
\tilde f_b\,f_b=(z-s_b)(z+s_b)=z^2-(z^2-1)=1,\qquad\text{т.е. }\tilde f_b=1/f_b\quad\text{(упражнение слайда 390).}
\]
Однократный обход только одной из точек $\pm1$ меняет $F_b\leftrightarrow\widetilde F_b$.

\emph{Особые точки.} Над $V(1,2)$ (там $\Ind_{-1}=0$) обход меняет элементы местами, два обхода возвращают: \emph{ветвление второго порядка} (одно; $\Phi|_{V(1,2)}$ --- одна аналитическая ветвь).
Аналогично над $V(-1,2)$. Над $D_2=\{|z|>1\}$: обход $\infty$ обходит обе точки, множитель $(-1)^2=1$ --- случай 1. Функция $S(z)=z\sqrt{1-1/z^2}_{(\circ)}$ (главная ветвь; $\Re(1-1/z^2)>0$ при $|z|>1$)
голоморфна в $D_2$ и $S^2=z^2-1$; на любом круге $B(b,R_b)$ с $|b|>1$ два корня $s_b$ и $S$ из $z^2-1$ связаны $s_b=\pm S$ (произведение $(s_b-S)(s_b+S)\equiv0$, а кольцо
голоморфных в круге функций без делителей нуля). Поэтому $\Phi$ над $D_2$ \emph{распадается на две голоморфные ветви}
\[
g_+(z)=z+z\sqrt{1-1/z^2}_{(\circ)}=2z-\tfrac1{2z}+O(z^{-3}),\qquad g_-(z)=\frac1{g_+(z)}=\tfrac1{2z}+O(z^{-3}),
\]
$g_+$ имеет в $\infty$ \emph{полюс первого порядка}, $g_-$ --- \emph{устранимую особенность} (нуль первого порядка) (упражнения слайда 390).

\subsubsection*{5. П.А.Ф. $\Arcsin z$}

\emph{Исходный элемент} $F_0=(B(0,1),f_0)$, $f_0=\arcsin z$ --- главная ветвь, голоморфная в $\Omega_0=\CC\setminus((-\infty,-1]\cup[1,\infty))$. По теореме единственности
\[
f_0'(z)=\frac{1}{\sqrt{1-z^2}_{(\circ)}}=\frac{i}{(z+1)^{1/2}_{(0)}(z-1)^{1/2}_{(\pi)}}\quad\text{в }\Omega_0
\]
(все три функции голоморфны в $\Omega_0$ и совпадают на $(-1,1)$; в нуле $i/(1\cdot i)=1$). Тождество $(f')^2(1-z^2)=1$ сохраняется (лемма 27.9), и в точках $\pm1$ даёт $0=1$:
радиус $F_0$ равен $1$ (а $\ge1$, так как $B(0,1)\subset\Omega_0$), $\pm1\notin D_\Phi$.

\emph{Продолжение.} Элемент $F_0'$ продолжается (лемма 27.8 для $i(z+1)^{-1/2}(z-1)^{-1/2}$) вдоль любого пути в $\CC\setminus\{\pm1\}$ с радиусами $\min|\gamma(t)\mp1|$ --- это П.А.Ф. $\Phi'$
(над каждой из областей $\Omega_0$, $\Omega_1=\CC\setminus[-1,\infty)$, $\Omega_2=\CC\setminus(-\infty,1]$ она распадается на две голоморфные ветви, отличающиеся знаком; по одному ветвлению второго порядка над
$V(\pm1,2)$; над $\{|z|>1\}$ --- две голоморфные ветви, $\pm i/(z\sqrt{1-1/z^2}_{(\circ)})$, с устранимой особенностью в $\infty$). По лемме 27.10 $F_0$ продолжается вдоль тех же путей с теми же радиусами.
(На слайде это сделано разбиением пути на куски в $\Omega_j$; лемма 27.10 делает то же самое.) Итак, $D_\Phi=\CC\setminus\{\pm1\}$, и $\Phi$ --- аналитическая ветвь над $D_\Phi$.

\begin{thmbox}{Лемма 27.11 (граничные значения $\arcsin$; упражнение слайда 398)}
Для $x\in\RR$, $|x|>1$ пределы $\arcsin(x\pm i0)=\lim_{y\to0+}\arcsin(x\pm iy)$ существуют, и
\[
\arcsin(x\pm i0)=\frac\pi2\pm i\operatorname{arch}x\ \ (x>1),\qquad
\arcsin(x+i0)+\arcsin(x-i0)=\begin{cases}\pi,&x>1,\\-\pi,&x<-1,\end{cases}
\]
где $\operatorname{arch}x=\ln(x+\sqrt{x^2-1})=\int_1^x\frac{dt}{\sqrt{t^2-1}}$.
\end{thmbox}
\begin{proof}
$\Omega_0$ звёздна относительно $0$, а $\arcsin z=\int_{[0,z]}g(\zeta)\,d\zeta$, $g=1/\sqrt{1-\zeta^2}_{(\circ)}$ (производные совпадают, $\arcsin0=0$). Пусть $x>1$, $y>0$. Замкнутый
треугольник с вершинами $0,iy,x+iy$ пересекает $\RR$ только в $0$, т.е. лежит в $\Omega_0$, и по теореме Коши
$\arcsin(x+iy)=\int_{[0,iy]}g+\int_0^xg(t+iy)\,dt$. Первый интеграл $\to0$ ($g$ ограничена около $0$). Во втором: $1-(t+iy)^2=(1-t^2+y^2)-2ity$; при $t<1$ это стремится к $1-t^2>0$,
при $t>1$ --- к $1-t^2<0$ \emph{снизу} ($\Im<0$), и главный корень стремится к $-i\sqrt{t^2-1}$. Значит, $g(t+iy)\to1/\sqrt{1-t^2}$ при $t<1$ и $\to i/\sqrt{t^2-1}$ при $t>1$.
Мажоранта: $|1-\zeta^2|=|1-\zeta||1+\zeta|\ge|1-t|$ (при $\Re\zeta=t\ge0$), т.е. $|g(t+iy)|\le|1-t|^{-1/2}$ --- интегрируема на $[0,x]$. По теореме Лебега
$\arcsin(x+i0)=\int_0^1\frac{dt}{\sqrt{1-t^2}}+i\int_1^x\frac{dt}{\sqrt{t^2-1}}=\frac\pi2+i\operatorname{arch}x$. Для $x-iy$ треугольник $0,-iy,x-iy$, и $1-(t-iy)^2$ стремится к $1-t^2$ сверху,
корень --- к $+i\sqrt{t^2-1}$, $g\to-i/\sqrt{t^2-1}$: $\arcsin(x-i0)=\frac\pi2-i\operatorname{arch}x$. Сумма равна $\pi$. Для $x<-1$ используем нечётность $\arcsin(-z)=-\arcsin z$ (обе части
голоморфны в $\Omega_0$ и совпадают на $(-1,1)$): $\arcsin(x\pm i0)=-\arcsin(-x\mp i0)$, сумма равна $-\pi$.
\end{proof}

\emph{Обходы точек $\pm1$.} Пусть $\gamma_\pm$ --- полуокружности $|z-1|=1$ из $0$ в $2$ через верхнюю/нижнюю полуплоскость, $F_2^\pm=(B(2,1),f_2^\pm)$ --- результаты продолжения $F_0$ вдоль них.
Внутренние точки $\gamma_\pm$ лежат в $\Omega_0$, где продолжение --- переразложение $\arcsin$ (семейство $E_{\gamma(t)}(\arcsin)$, лемма 25.B локально); поэтому $f_2^\pm=\arcsin$ на полукруге
$B(2,1)\cap\{\pm\Im z>0\}$ (по теореме единственности на этом связном полукруге), и $f_2^\pm(x)=\arcsin(x\pm i0)$ при $x\in(1,3)$. По лемме 27.11 $f_2^++f_2^-=\pi$ на $(1,3)$, значит, в $B(2,1)$:
$F_2^-=\pi-F_2^+$.

Положительный обход $\gamma_1=\{1+e^{it}:t\in[-\pi,\pi]\}$ --- это $\gamma_-$, затем $\gamma_+^-$. Вдоль $\gamma_-$: $F_0\mapsto F_2^-=\pi-F_2^+$; вдоль $\gamma_+^-$ элемент $F_2^+$ переходит в $F_0$
(лемма 25.A(2)), а потому (лемма 27.8) $\pi-F_2^+\mapsto\pi-F_0$. Итак, \textbf{обход $1$: $F_0\mapsto-F_0+\pi$}. Аналогично для $\gamma_{-1}=\{-1+e^{it}:t\in[0,2\pi]\}$ (сначала верхняя
полуокружность до $-2$, затем нижняя) и суммы $-\pi$: \textbf{обход $-1$: $F_0\mapsto-F_0-\pi$}. Двукратный обход одной точки возвращает $F_0$. По лемме 27.8 операторы обхода
перестановочны с $G\mapsto\pm G+C$: $T_1(\pm F_0+C)=\pm(-F_0+\pi)+C$, $T_{-1}(\pm F_0+C)=\pm(-F_0-\pi)+C$. Поэтому
\[
F_0\xrightarrow{\gamma_1}-F_0+\pi\xrightarrow{\gamma_{-1}}-(-F_0-\pi)+\pi=F_0+2\pi,\qquad
F_0\xrightarrow{\gamma_{-1}}-F_0-\pi\xrightarrow{\gamma_{1}}-(-F_0+\pi)-\pi=F_0-2\pi .
\]
\begin{remarkbox}{Замечание (опечатка слайда 398)}
На слайде сказано, что при обходе по $\gamma_1$, а потом по $\gamma_{-1}$ элемент $F_0$ переходит в $F_0-2\pi$. Как показывает вычисление выше (и численная проверка),
в этом порядке получается $F_0+2\pi$, а $F_0-2\pi$ --- при обратном порядке обходов. На дальнейшие выводы это не влияет.
\end{remarkbox}

\emph{Элементы над точкой $b$.} Все элементы $\Phi$ удовлетворяют $\sin f(z)=z$ (лемма 27.9 с $Q=\sin u-z$). Если $f,g$ --- функции элементов над $b$, то
$0=\sin f-\sin g=2\cos\frac{f+g}2\sin\frac{f-g}2$ в круге; так как кольцо голоморфных в круге функций без делителей нуля, $\cos\frac{f+g}2\equiv0$ или $\sin\frac{f-g}2\equiv0$, и по
непрерывности (функция со значениями в дискретном множестве постоянна) $g=-f+\pi+2\pi k$ или $g=f+2\pi k$. Обратно, над $0$ все такие элементы достигаются: $F_0+2\pi k$ --- $k$-кратным
(или $|k|$-кратным обратным) обходом «сначала $1$, затем $-1$», а $-F_0+\pi+2\pi k=T_1(F_0-2\pi k)$. Переносом вдоль пути $\sigma$ из $0$ в $b$ (лемма 27.8: $\Lambda_\sigma(\pm F_0+C)=\pm\Lambda_\sigma F_0+C$)
получаем: \textbf{над каждой $b\in\CC\setminus\{\pm1\}$ ровно счётное множество элементов $F_b+2\pi k$, $-F_b+\pi+2\pi k$, $k\in\ZZ$} ($F_b$ --- фиксированный элемент, например $E_b(\arcsin)$ при $b\in\Omega_0$).

\emph{Особые точки $\pm1$.} Над $V(1,2)$ (там $\Ind_{-1}=0$; $0\in V(1,2)$ и $\gamma_1\subset V(1,2)$): $T_1(\varepsilon F_0+C)=-\varepsilon F_0+\varepsilon\pi+C\ne\varepsilon F_0+C$, $T_1^2=\mathrm{id}$. Значит,
$\Phi$ над $V(1,2)$ распадается на счётное число аналитических ветвей (орбиты $\{F_0+2\pi k,-F_0+\pi-2\pi k\}$), у каждой --- \emph{ветвление второго порядка}. Аналогично над $V(-1,2)$.

\emph{Точка $\infty$.} Пусть $h(z)=\dfrac{i}{z\sqrt{1-1/z^2}_{(\circ)}}\in H(\{|z|>1\})$. На верхнем полукруге $B(2,1)\cap\{\Im z>0\}$ имеем $(f_2^+)'=1/\sqrt{1-z^2}_{(\circ)}$, $h^2=1/(1-z^2)$, так что
$(f_2^+)'=\pm h$; при $z\to x+i0$, $x\in(1,3)$, обе стороны стремятся к $i/\sqrt{x^2-1}$ (см. доказательство леммы 27.11), поэтому знак «$+$», и $(f_2^+)'=h$ в $B(2,1)$. Вдоль окружности $|z|=2$
производная продолжается семейством $E_{\gamma(t)}(h)$ (лемма 25.B локально) и возвращается. По лемме 27.10 $F_2^+$ переходит в $F_2^++C$,
\[
C=\oint_{|z|=2}h(z)\,dz=\oint_{|z|=2}\frac iz\Bigl(1+\frac1{2z^2}+\dots\Bigr)dz=i\cdot2\pi i=-2\pi
\]
(ряд сходится равномерно на окружности). Итак, при обходе $\infty$ $F_2^+\mapsto F_2^+-2\pi$ (слайд 399), а $-F_2^++\pi+2\pi k\mapsto-F_2^++\pi+2\pi(k+1)$. Обе орбиты бесконечны:
\textbf{над $\{1<|z|<\infty\}$ у $\Arcsin z$ два логарифмических ветвления} (две аналитические ветви $\{F_2^++2\pi k\}$ и $\{-F_2^++\pi+2\pi k\}$).

\begin{remarkbox}{Итог для $\Arcsin z$}
$D_\Phi=\CC\setminus\{\pm1\}$; над каждой точкой --- элементы $F_b+2\pi k$, $-F_b+\pi+2\pi k$; над $V(\pm1,2)$ --- счётное число ветвлений второго порядка; над окрестностью $\infty$ --- два
логарифмических ветвления. Сравните с $\Arctg$: там над $\infty$ счётное число однозначных ветвей.
\end{remarkbox}

\begin{remarkbox}{Коротко: план ответа на экзамене}
\begin{enumerate}
\item Определение П.А.Ф., её область определения (область), П.А.Ф. определяется любым элементом.
\item Теорема Пуанкаре--Вольтерра: замкнутый путь заменяется близкой ломаной с рациональными вершинами (лемма о близких путях); таких ломаных счётно.
\item Аналитическая ветвь над $D_1$; число элементов над точками $D_1$ одинаково (продолжение вдоль пути --- биекция). Голоморфная ветвь, распадение; над односвязной $D_1$ --- одна
голоморфная ветвь (монодромия); правильные точки.
\item Изолированные особые точки: в $V(a,\delta)$ любой замкнутый путь гомотопен $k$-кратной окружности, поэтому всё определяется оператором обхода $T$; $\Phi_1(c)$ --- одна орбита $T$, её длина
$n$ не зависит от точки. $n=1$ --- однозначный характер (устранимая, полюс, существенная), $2\le n<\infty$ --- ветвление порядка $n$, $n=\infty$ --- логарифмическое.
\item Предложение 15.8: около точки ветвления порядка $n$ $\Phi_1=f_0\circ\phi^{-1}$, $\phi(w)=a+w^n$; ряды Пюизё.
\item Действия над П.А.Ф.; сохранение тождеств и непродолжаемость через точку.
\item Примеры: $\Ln(z-w)$ и $(z-w)^\lambda$ (логарифмическое ветвление / порядок $q$ при $\lambda=p/q$); $\Arctg$ (логарифмические в $\pm i$, однозначные ветви у $\infty$); $\Zh^{-1}$
(ветвления второго порядка в $\pm1$, у $\infty$ полюс и нуль, $\tilde f=1/f$); $\Arcsin$ (элементы $\pm F+\ldots$, ветвления второго порядка в $\pm1$, логарифмические у $\infty$).
\end{enumerate}
\end{remarkbox}

% ===== part 5b =====
% ===================================================================
%  Часть 5b: билеты 28--29
% ===================================================================

% -------------------------------------------------------------------
\ticket{28}{Теорема об аналитическом продолжении первообразного элемента. Пример её применения для описания продолжения элемента $E_1\bigl(\ln(\sqrt z_{(\circ)}\,\Zh(z))\bigr)$}
\label{t28:begin}

\begin{defs}
\textbf{Пути.} Путь --- непрерывное отображение $\gamma\colon I=[0,1]\to\CC$, $\gamma^*=\gamma(I)$ --- его носитель (компакт).
Путь \emph{спрямляем}, если он имеет конечную длину (например, кусочно гладкий путь или ломаная);
для спрямляемого $\gamma$ и непрерывной на $\gamma^*$ функции $h$ определён интеграл $\int_\gamma h\,dz$.
$\gamma_t:=\gamma|_{[0,t]}$, $\gamma_{t,\xi}:=\gamma|_{[t,\xi]}$ (при $\xi<t$ --- проходимый в обратном направлении),
$[z_1,z_2]$ --- отрезок от $z_1$ к $z_2$. $\|\gamma-\sigma\|_I:=\max_{t\in I}|\gamma(t)-\sigma(t)|$.
$U_\eps(\gamma^*):=\bigcup_{t\in I}B(\gamma(t),\eps)$ --- $\eps$-окрестность носителя.
Для замкнутого спрямляемого пути $\gamma$ и $p\notin\gamma^*$ \emph{индекс}
$\Ind_p(\gamma)=\frac1{2\pi i}\int_\gamma\frac{dz}{z-p}\in\ZZ$; он постоянен на компонентах связности $\CC\setminus\gamma^*$ и равен $0$
на неограниченной компоненте.

\textbf{Первообразная.} $F$ --- (комплексная) первообразная $h$ в открытом $U$, если $F\in H(U)$ и $F'=h$.
В круге (и вообще в односвязной области) любая $h\in H$ имеет первообразную, единственную с точностью до постоянной
(в связном множестве), например $F(z)=\int_{[a,z]}h\,d\zeta$ в круге $B(a,r)$.
\emph{Интегральная теорема Коши для круга:} если $h\in H(B)$, $B$ --- круг, то $\int_\sigma h\,dz=0$ для любого замкнутого спрямляемого пути $\sigma\subset B$.

\textbf{Элементы (Опр.~15.2, 15.3).} Элемент --- пара $F=(B,f)$, $B=B(a,R)$, $f\in H(B)$; $a$ --- центр, $R\in(0,+\infty]$ --- радиус.
$F$ \emph{канонический}, если $B$ --- круг сходимости ряда Тейлора $f$ в точке $a$.
Если $f$ голоморфна в окрестности $a$, то $E_a(f)$ --- канонический элемент с центром $a$, заданный $f$.
По теореме Коши--Тейлора, если $f$ голоморфна в круге $B(a,r)$, то радиус $E_a(f)$ не меньше $r$.
$F_1=(B_1,f_1)$ --- \emph{непосредственное аналитическое продолжение} (НАП) $F_0=(B_0,f_0)$, если $B_0\cap B_1\neq\varnothing$ и $f_0=f_1$ на $B_0\cap B_1$
(пересечение двух кругов выпукло, значит связно; поэтому по теореме единственности достаточно совпадения на непустом открытом подмножестве $B_0\cap B_1$).
Два канонических элемента с общим центром, один из которых --- НАП другого, совпадают.

\textbf{Продолжение вдоль пути (Опр.~15.4).} Семейство канонических элементов $F_t=(B(a_t,R_t),f_t)$, $t\in I$, ---
аналитическое продолжение $F_0$ вдоль $\gamma$, если: 1) $a_t=\gamma(t)$, $R_t>0$; 2) для каждого $t$ есть связная окрестность $U_t\subset I$ точки $t$,
такая что при $\xi\in U_t$: $\gamma(\xi)\in B_t$ и $F_\xi$ --- НАП $F_t$. Говорят, что $F_0$ \emph{переходит} в $F_1$ при продолжении вдоль $\gamma$.
\emph{Лемма 15.1:} либо $R_t\equiv\infty$, либо $R_t$ непрерывна на $I$; в частности, $R_\gamma:=\min_{t}R_t>0$.
\emph{Предложение 15.2 (единственность):} результат продолжения вдоль пути единствен; более того, два продолжающих семейства совпадают при всех $t$
(доказательство показывает, что множество $\{t:F_t=\widetilde F_t\}$ есть весь $I$). Применяя это к путям $\gamma_\tau$, получаем:
\emph{сужение продолжающего семейства на $[0,\tau]$ однозначно определено}.
\emph{Склейка и обращение:} если $F$ переходит в $G$ вдоль $\gamma_1$, а $G$ в $K$ вдоль $\gamma_2$ ($\gamma_1(1)=\gamma_2(0)$), то $F$ переходит в $K$ вдоль
$\gamma_1\gamma_2$; если $F$ переходит в $G$ вдоль $\gamma$, то $G$ переходит в $F$ вдоль обратного пути $\gamma^-(t)=\gamma(1-t)$
(берём семейство $F_{1-t}$).

\textbf{Предложение 15.4 (лемма о близких путях; билет 26).} Пусть $F$ продолжается вдоль $\gamma$, $0<r\le R_\gamma$, $r<\infty$, и путь $\sigma$ с теми же концами
удовлетворяет $\|\gamma-\sigma\|_I<r$. Тогда $F$ продолжается вдоль $\sigma$ и результаты продолжения вдоль $\gamma$ и $\sigma$ совпадают.
(На слайде $r=R_\gamma<\infty$; доказательство дословно проходит для любого конечного $r\in(0,R_\gamma]$.)
\emph{Теорема 15.1:} если $\gamma_0,\gamma_1$ гомотопны с закреплёнными концами (гомотопия $\gamma(s,t)$) и $F$ продолжается вдоль всех $\gamma(s,\cdot)$, то результаты продолжения вдоль $\gamma_0$ и $\gamma_1$ совпадают.

\textbf{П.А.Ф. и ветви (Опр.~15.5--15.9; билет 27).} П.А.Ф. $\Phi$, порождённая каноническим $F_a$, --- совокупность всех канонических элементов,
получаемых из $F_a$ продолжением вдоль путей с началом $a$; $D_\Phi$ --- объединение кругов её элементов (область определения).
Все элементы $\Phi$ с центром $b$ получаются из одного из них, $F_b$, продолжением вдоль замкнутых путей с началом и концом $b$
(если $F_a\to F_b$ вдоль $\alpha$ и $F_a\to\widetilde F_b$ вдоль $\beta$, то $F_b\to\widetilde F_b$ вдоль петли $\alpha^-\beta$).
\emph{Аналитическая ветвь} $\Phi_1$ над областью $D_1\subset D_\Phi$, определяемая $F_b\in\Phi$, $b\in D_1$, --- совокупность результатов продолжения $F_b$
вдоль всех путей в $D_1$ (предполагается, что $F_b$ продолжается вдоль \emph{любого} пути в $D_1$ с началом $b$).
\emph{Голоморфная ветвь} в области $D_2\subset D_1$ --- функция $f\in H(D_2)$ с $E_c(f)\in\Phi_1$ для всех $c\in D_2$; $\Phi_1$ \emph{распадается} над $D_2$ на голоморфные ветви $\{f_j\}$,
если каждый элемент $\Phi_1$ с центром в $D_2$ имеет вид $E_c(f_j)$. По теореме о монодромии над односвязной $D_2$ аналитическая ветвь всегда распадается на голоморфные ветви.
\emph{Изолированная особая точка} (Опр.~15.9): $D_1=V(a,\delta)$ --- проколотая окрестность $a$ ($a\in\Ch$; $V(\infty,\delta)=\{|z|>1/\delta\}$), $V(a,\delta)\subset D_\Phi$, $\Phi_1$ --- аналитическая ветвь над $V(a,\delta)$.
Если после одного обхода $a$ элемент $F_b$ переходит в себя --- $a$ \emph{точка однозначного характера} для $\Phi_1$ (классифицируется как у голоморфной в $V(a,\delta)$ функции: устранимая, полюс, существенно особая; если голоморфная ветвь продолжается в $a$, то $a$ --- правильная точка);
если $F_b$ возвращается в себя впервые после $n\ge2$ обходов --- \emph{ветвление порядка $n$};
если ни после какого числа обходов не возвращается --- \emph{логарифмическое ветвление} (ветвление бесконечного порядка).
\end{defs}

\subsection{Вспомогательные утверждения}

Сначала докажем два простых общих факта о продолжении элементов, у которых известна производная; они используются и в доказательстве
Предложения 15.5, и в примере.

\begin{thmbox}{Лемма 28.1 (производная сохраняется при продолжении)}
Пусть $U\subset\CC$ открыто, $h\in H(U)$, путь $\gamma$ лежит в $U$, $\{F_t=(B_t,f_t)\}$ --- аналитическое продолжение элемента $F_0$ вдоль $\gamma$,
и $f_0'=h$ в некоторой окрестности точки $\gamma(0)$. Положим $d_t:=\dist(\gamma(t),\CC\setminus U)$ ($d_t=+\infty$ при $U=\CC$). Тогда для каждого $t\in I$:
\begin{enumerate}
\item $f_t'=h$ на $B_t\cap B(\gamma(t),d_t)$;
\item $R_t\ge d_t$.
\end{enumerate}
\end{thmbox}
\begin{proof}
Заметим, что $B(\gamma(t),d_t)\subset U$, $d_t>0$, и множество $K_t:=B_t\cap B(\gamma(t),d_t)$ выпукло (пересечение двух кругов), содержит $\gamma(t)$,
и на нём голоморфны обе функции $f_t'$ и $h$. Поэтому, по теореме единственности, утверждение 1 равносильно тому, что $f_t'=h$ в \emph{некоторой}
окрестности $\gamma(t)$. Пусть $E$ --- множество тех $t$, для которых это так. $0\in E$ по условию.

\emph{$E$ открыто в $I$.} Пусть $t\in E$. По непрерывности $\gamma$ найдётся связная окрестность $W\subset U_t$ точки $t$ с $\gamma(\xi)\in B(\gamma(t),d_t)$ при $\xi\in W$.
Для $\xi\in W$: $\gamma(\xi)\in B_t\cap B_\xi$, и $f_\xi=f_t$ на $B_t\cap B_\xi$ (НАП), т.е. в окрестности $\gamma(\xi)$; а так как $\gamma(\xi)\in K_t$ и $f_t'=h$ на $K_t$,
получаем $f_\xi'=h$ в окрестности $\gamma(\xi)$, т.е. $\xi\in E$.

\emph{$E$ замкнуто в $I$.} Пусть $t$ --- предельная точка $E$. Выберем $\xi\in E\cap U_t$ так, чтобы $|\gamma(\xi)-\gamma(t)|<d_t$ (возможно по непрерывности $\gamma$).
Тогда $\gamma(\xi)\in B_t\cap B_\xi$, $f_t=f_\xi$ около $\gamma(\xi)$, а $f_\xi'=h$ около $\gamma(\xi)$. Значит $f_t'=h$ на непустом открытом подмножестве выпуклого
множества $K_t$ (ведь $\gamma(\xi)\in K_t$), а потому на всём $K_t$, т.е. $t\in E$.

Так как $I$ связно, $E=I$; утверждение 1 доказано. Для 2: в круге $B(\gamma(t),d_t)\subset U$ функция $h$ имеет первообразную $H$; в окрестности $\gamma(t)$
имеем $(f_t-H)'=0$, т.е. $f_t=H+c$. Ряд Тейлора $H+c$ в точке $\gamma(t)$ сходится в $B(\gamma(t),d_t)$ (теорема Коши--Тейлора), а он же --- ряд Тейлора $f_t$.
Значит радиус канонического элемента $F_t$ не меньше $d_t$.
\end{proof}

\begin{thmbox}{Лемма 28.2 (непродолжаемость через неустранимую особенность производной)}
Пусть $P\subset\CC$ --- замкнутое дискретное множество, $U=\CC\setminus P$, $h\in H(U)$, и каждая точка $p\in P$ --- неустранимая изолированная особая точка $h$
(полюс или существенно особая). Пусть $F_0=(B_0,f_0)$ --- канонический элемент с центром $a\in U$ и $f_0'=h$ около $a$. Тогда $F_0$ продолжается вдоль пути $\gamma$, $\gamma(0)=a$,
в том и только том случае, когда $\gamma^*\cap P=\varnothing$.
\end{thmbox}
\begin{proof}
Если $\gamma^*\subset U$, продолжаемость будет доказана в Предложении 15.5 ниже (при $\eps=\dist(\gamma^*,P)>0$; $U_\eps(\gamma^*)\subset U$) --- здесь мы
используем лемму только в одну сторону, а «достаточность» выводится из Предложения 15.5, независимого от этой леммы.

Пусть $\gamma^*\cap P\ne\varnothing$ и, от противного, $\{F_t\}$ продолжает $F_0$ вдоль $\gamma$. Множество $\{t:\gamma(t)\in P\}$ замкнуто
(прообраз замкнутого) и непусто; пусть $\tau$ --- его минимум, $p=\gamma(\tau)$. Так как $\gamma(0)=a\notin P$, то $\tau>0$, и $\gamma([0,\tau))\subset U$.
Выберем $r\in(0,R_\tau]$ так, чтобы $B(p,r)\cap P=\{p\}$. Возьмём $\xi<\tau$, $\xi\in U_\tau$, с $|\gamma(\xi)-p|<r$. Сужение семейства на $[0,\xi]$ --- продолжение
$F_0$ вдоль пути $\gamma_\xi\subset U$, и по Лемме 28.1 $f_\xi'=h$ в окрестности $\gamma(\xi)$. Так как $F_\xi$ --- НАП $F_\tau$ и $\gamma(\xi)\in B_\tau\cap B_\xi$, то $f_\tau'=f_\xi'=h$
в окрестности $\gamma(\xi)$. Функции $f_\tau'$ и $h$ голоморфны в связной проколотой окрестности $B(p,r)\setminus\{p\}$ и совпадают на её открытом подмножестве, значит совпадают всюду
в ней. Но $f_\tau'$ голоморфна в $B(p,r)$, т.е. $h$ имеет в $p$ устранимую особенность --- противоречие.
\end{proof}

\subsection{Теорема о продолжении первообразного элемента}

\begin{thmbox}{Предложение 15.5 (о продолжении первообразного элемента)}
Пусть $\gamma\colon[0,1]\to\CC$ --- путь, $\gamma(0)=a$, $\gamma(1)=b$, $\eps>0$, $U_\eps:=U_\eps(\gamma^*)$. Пусть $h\in H(U_\eps)$, а $f$ --- первообразная $h$
в $B(a,\eps)$ с $f(a)=0$, и $F_0:=E_a(f)$. Тогда $F_0$ продолжается вдоль $\gamma$, причём для продолжающего семейства $\{F_{\gamma,t}\}$ выполнено
$R_\gamma=\min_tR_{\gamma,t}\ge\eps$.
\end{thmbox}
Элемент $F_0$ называется \emph{первообразным каноническим элементом для $h$ в точке $a$}.

\begin{remarkbox}{Замечание}
На слайде в формулировке написано «$R_\gamma\ge\eps$» --- именно нестрогое неравенство и верно: например, при $h(z)=1/z$, $\gamma\equiv1$ (постоянный путь), $\eps=1$
радиус элемента $E_1(\ln z)$ равен ровно $1$.
\end{remarkbox}

\begin{proof}
\textbf{Шаг 1: $\gamma$ спрямляем.} Для $t\in I$ и $z\in B(\gamma(t),\eps)$ положим
\begin{equation}\label{t28:prim}
f_{\gamma,t}(z)=\int_{\gamma_t}h(\zeta)\,d\zeta+\int_{[\gamma(t),z]}h(\zeta)\,d\zeta .
\end{equation}
Интегралы определены: $\gamma_t^*\subset\gamma^*\subset U_\eps$, $[\gamma(t),z]\subset B(\gamma(t),\eps)\subset U_\eps$. Первое слагаемое --- константа, второе ---
первообразная $h$ в круге $B(\gamma(t),\eps)$, поэтому $f_{\gamma,t}\in H(B(\gamma(t),\eps))$ и $f_{\gamma,t}'=h$ там. Положим $F_{\gamma,t}:=E_{\gamma(t)}(f_{\gamma,t})$; по теореме Коши--Тейлора
$R_{\gamma,t}\ge\eps$. При $t=0$: $f_{\gamma,0}(z)=\int_{[a,z]}h\,d\zeta=f(z)$ (обе функции --- первообразные $h$ в $B(a,\eps)$, равные $0$ в $a$), т.е. $F_{\gamma,0}=F_0$.

Проверим условие 2) Определения 15.4. Фиксируем $t$. По непрерывности $\gamma$ найдётся $\delta>0$ такое, что $|\gamma(s)-\gamma(t)|<\eps$ для всех $s\in U_t:=(t-\delta,t+\delta)\cap I$
(связная окрестность). Пусть $\xi\in U_t$. Тогда $\gamma(\xi)\in B(\gamma(t),\eps)\subset B_{\gamma,t}$. Пусть $z\in K:=B(\gamma(t),\eps)\cap B(\gamma(\xi),\eps)$ (выпуклое множество, содержит $\gamma(\xi)$). Тогда
\[
f_{\gamma,\xi}(z)-f_{\gamma,t}(z)=\int_{\gamma_{t,\xi}}h\,d\zeta+\int_{[\gamma(\xi),z]}h\,d\zeta+\int_{[z,\gamma(t)]}h\,d\zeta=\int_{\sigma}h\,d\zeta,
\]
где $\sigma$ --- замкнутый путь: $\gamma$ от $\gamma(t)$ до $\gamma(\xi)$, затем отрезки $[\gamma(\xi),z]$ и $[z,\gamma(t)]$ (мы использовали $\int_{\gamma_\xi}=\int_{\gamma_t}+\int_{\gamma_{t,\xi}}$,
что верно и при $\xi<t$ с учётом ориентации). Весь путь $\sigma$ лежит в круге $B(\gamma(t),\eps)$: дуга $\gamma|_{[t,\xi]}$ --- по выбору $\delta$, отрезки --- так как их концы
$\gamma(\xi),z,\gamma(t)$ лежат в этом выпуклом круге. По интегральной теореме Коши для круга ($h\in H(B(\gamma(t),\eps))$) $\int_\sigma h=0$.
Итак, $f_{\gamma,\xi}=f_{\gamma,t}$ на $K$ --- непустом открытом подмножестве $B_{\gamma,t}\cap B_{\gamma,\xi}$; по теореме единственности они совпадают на всём (выпуклом) $B_{\gamma,t}\cap B_{\gamma,\xi}$,
т.е. $F_{\gamma,\xi}$ --- НАП $F_{\gamma,t}$. Семейство $\{F_{\gamma,t}\}$ продолжает $F_0$ вдоль $\gamma$, и $R_\gamma\ge\eps$.

\begin{remarkbox}{Замечание}
На слайде НАП обосновывается равенством $f_{\gamma,\xi}(\gamma(t))=f_{\gamma,t}(\gamma(t))$ и совпадением производных в круге $D_\xi$ с центром $\gamma(\xi)$;
но точка $\gamma(t)$ может не лежать в $D_\xi$. Приведённая выше выкладка с замкнутым путём $\sigma$ в круге $B(\gamma(t),\eps)$ устраняет эту неточность.
\end{remarkbox}

\textbf{Шаг 2: произвольный (не обязательно спрямляемый) путь.} Спрямляемость нужна была только для интеграла по $\gamma$.
По равномерной непрерывности $\gamma$ на $I$ выберем разбиение $0=s_0<s_1<\dots<s_m=1$, для которого колебание $\gamma$ на каждом $[s_{j-1},s_j]$ меньше $\eps/6$,
и пусть $\lambda$ --- вписанная ломаная: $\lambda(s_j)=\gamma(s_j)$, на $[s_{j-1},s_j]$ $\lambda$ линейна.
\begin{detail}{$\|\gamma-\lambda\|_I<\eps/3$}
При $t\in[s_{j-1},s_j]$ точка $\lambda(t)$ --- выпуклая комбинация $\gamma(s_{j-1})$ и $\gamma(s_j)$, а обе эти точки лежат в круге $B(\gamma(t),\eps/6)$ (колебание $<\eps/6$).
Круг выпуклый, поэтому $|\lambda(t)-\gamma(t)|<\eps/6<\eps/3$.
\end{detail}
Далее, $U_{2\eps/3}(\lambda^*)\subset U_\eps$: если $|z-\lambda(t)|<2\eps/3$, то $|z-\gamma(t)|<2\eps/3+\eps/3=\eps$. Значит $h\in H(U_{2\eps/3}(\lambda^*))$, $\lambda(0)=a$,
и $f$ --- первообразная $h$ в $B(a,2\eps/3)$ с $f(a)=0$. Ломаная спрямляема, поэтому по шагу 1 (с $2\eps/3$ вместо $\eps$) $F_0$ продолжается вдоль $\lambda$ и
$R_\lambda\ge2\eps/3$. Возьмём $r:=\min(R_\lambda,2\eps/3)$, конечное, $0<r\le R_\lambda$; так как $\|\gamma-\lambda\|_I<\eps/3<r$, а концы $\gamma$ и $\lambda$ совпадают, то по Предложению 15.4
$F_0$ продолжается вдоль $\gamma$ (и результат тот же, что вдоль $\lambda$). Наконец, путь $\gamma$ лежит в открытом множестве $U_\eps$, $f_0'=h$ около $a$, и
$\dist(\gamma(t),\CC\setminus U_\eps)\ge\eps$ (так как $B(\gamma(t),\eps)\subset U_\eps$). По Лемме 28.1 $R_{\gamma,t}\ge\eps$ при всех $t$, т.е. $R_\gamma\ge\eps$.
\end{proof}

\begin{remarkbox}{Сравнение со слайдом}
На слайде оценка $R_\gamma\ge\eps$ в шаге 2 получена так: $f_{\gamma,t}$ --- переразложение $f_{\lambda,t}$ в точке $\gamma(t)$ (так устроено семейство в доказательстве
Предложения 15.4), $f_{\lambda,t}'=h$, и по теореме единственности $f_{\gamma,t}$ голоморфна там же, где её производная $h$, т.е. в $B(\gamma(t),\eps)$.
Лемма 28.1 --- аккуратное оформление именно этого рассуждения, не зависящее от устройства доказательства Предложения 15.4.
\end{remarkbox}

\subsection{Следствия: интеграл по пути, И.Т.К. для гомотопных путей, логарифм}

\begin{thmbox}{Следствие 28.3}
В условиях Предложения 15.5, если $\gamma$ спрямляем, то
\[
\int_\gamma h\,dz=f_{\gamma,1}(b)-f_{\gamma,0}(a)=f_{\gamma,1}(b).
\]
Для произвольного (неспрямляемого) пути правую часть естественно принять за \emph{определение} $\int_\gamma h\,dz$. Более общо: если $F=(B,g)$ --- \emph{любой} элемент
с центром $a$ и $g'=h$ около $a$, то $F$ продолжается вдоль $\gamma$ и результат продолжения --- элемент $E_b(g_1)$, где
$g_1(z)=g(a)+\int_\gamma h\,d\zeta+\int_{[b,z]}h\,d\zeta$ около $b$. В частности, если $\gamma$ замкнут ($b=a$), то $F$ переходит в элемент $(B,\,g+\int_\gamma h\,dz)$.
\end{thmbox}
\begin{proof}
Первое --- формула \eqref{t28:prim} при $t=1$, $z=b$ (и $f_{\gamma,0}(a)=f(a)=0$). Для второго: $g-g(a)$ --- первообразная $h$ около $a$, равная $0$ в $a$, т.е. $E_a(g)=F_0+g(a)$
(прибавление константы ко всем элементам семейства, очевидно, сохраняет свойства Определения 15.4); остаётся применить \eqref{t28:prim} при $t=1$.
При $b=a$: $g_1-g$ --- постоянная $\int_\gamma h$ (обе --- первообразные $h$ около $a$, их значения в $a$ различаются на $\int_\gamma h$), и у канонических элементов $E_a(g)$ и $E_a(g+c)$ один и тот же круг.
\end{proof}

\begin{thmbox}{Следствие 28.4 (И.Т.К. для гомотопных путей)}
Пусть $D$ --- область, $h\in H(D)$, $\gamma_0,\gamma_1$ --- спрямляемые пути в $D$ с общими концами, гомотопные в $D$ с закреплёнными концами. Тогда $\int_{\gamma_0}h\,dz=\int_{\gamma_1}h\,dz$.
\end{thmbox}
\begin{proof}
Пусть $\gamma(s,t)$ --- гомотопия в $D$, $\gamma_s:=\gamma(s,\cdot)$. Для каждого $s$ носитель $\gamma_s^*$ компактен и лежит в $D$, поэтому $\eps_s:=\dist(\gamma_s^*,\CC\setminus D)>0$
(при $D=\CC$ берём $\eps_s=1$) и $U_{\eps_s}(\gamma_s^*)\subset D$. Пусть $f$ --- первообразная $h$ в круге $B(a,r)\subset D$ с $f(a)=0$, $F_0=E_a(f)$. По Предложению 15.5 $F_0$
продолжается вдоль каждого $\gamma_s$, и по Теореме 15.1 результаты продолжения вдоль $\gamma_0$ и $\gamma_1$ совпадают: $f_{\gamma_0,1}=f_{\gamma_1,1}$. По Следствию 28.3
$\int_{\gamma_0}h=f_{\gamma_0,1}(b)=f_{\gamma_1,1}(b)=\int_{\gamma_1}h$.
\end{proof}

\begin{thmbox}{Следствие 28.5 (продолжение логарифма функции)}
Пусть $U$ --- открытое множество, $g\in H(U)$, $g\ne0$ в $U$, $a\in U$, и $\ell$ --- какая-нибудь голоморфная ветвь логарифма $g$ около $a$ (т.е. $e^{\ell}=g$). Тогда элемент
$E_a(\ell)$ продолжается вдоль любого пути $\gamma\subset U$ с началом $a$; для продолжающего семейства $\{(B_t,\ell_t)\}$ выполнено $\ell_t'=g'/g$ и $e^{\ell_t}=g$ около $\gamma(t)$,
$R_t\ge\dist(\gamma(t),\CC\setminus U)$, а если $\gamma$ --- спрямляемая петля, то $E_a(\ell)$ переходит в $E_a\bigl(\ell+\int_\gamma\frac{g'}{g}dz\bigr)$.
\end{thmbox}
\begin{proof}
$\ell'=g'/g=:h\in H(U)$. С $\eps:=\dist(\gamma^*,\CC\setminus U)>0$ имеем $U_\eps(\gamma^*)\subset U$, и $\ell-\ell(a)$ --- первообразный элемент для $h$; по Предложению 15.5 и Следствию 28.3
элемент $E_a(\ell)$ продолжается, по Лемме 28.1 $\ell_t'=h$ и $R_t\ge\dist(\gamma(t),\CC\setminus U)$. Функция $\varphi_t=g\,e^{-\ell_t}$ около $\gamma(t)$ имеет производную $g'e^{-\ell_t}-g\,h\,e^{-\ell_t}=0$,
т.е. $\varphi_t\equiv c_t$ около $\gamma(t)$. При $\xi\in U_t$ функции $\ell_\xi$ и $\ell_t$ совпадают около $\gamma(\xi)$, поэтому $c_\xi=c_t$: функция $t\mapsto c_t$ локально постоянна на связном $I$,
значит $c_t=c_0=1$. Формула для петли --- из Следствия 28.3.
\end{proof}

\subsection{Пример: П.А.Ф. $\Ln\bigl(\sqrt z\,\Zh(z)\bigr)$}

Здесь $\sqrt z_{(\circ)}=\sqrt{|z|}\,e^{i\arg z/2}$, $\arg z\in(-\pi,\pi)$, --- главная ветвь корня в $\CC\setminus R_\pi$, $R_\pi=(-\infty,0]$; $\Zh(z)=\frac12\bigl(z+\frac1z\bigr)$ --- функция Жуковского;
$\ln$ --- главная ветвь логарифма. Опишем П.А.Ф. $\Phi$, порождённую элементом $F_0:=E_1\bigl(\ln(\sqrt z_{(\circ)}\Zh(z))\bigr)$, по схеме полного описания П.А.Ф. (слайд 380).

\subsubsection*{Шаг 0. Подготовка: функция под логарифмом и логарифмическая производная}
Положим $g(z):=\sqrt z_{(\circ)}\,\Zh(z)=\dfrac{z^2+1}{2\sqrt z_{(\circ)}}$, $g\in H(\CC\setminus R_\pi)$. Нули $g$ --- точки $\pm i$ (простые), $g$ не определена в $0$ и не продолжается туда
(при $z\to0$ $|g(z)|\sim\frac1{2\sqrt{|z|}}\to\infty$). Логарифмическая производная (одна и та же для любой ветви корня, так как ветви отличаются знаком):
\[
\frac{g'}{g}=\frac{(\sqrt z)'}{\sqrt z}+\frac{\Zh'}{\Zh}=\frac1{2z}+\frac{1-1/z^2}{z+1/z}=\frac1{2z}+\frac{z^2-1}{z(z^2+1)}.
\]
\begin{detail}{разложение на простейшие дроби}
$\dfrac{z^2-1}{z(z^2+1)}=\dfrac Az+\dfrac{Bz+C}{z^2+1}$: при $z\to0$ получаем $A=-1$; тогда $z^2-1=-(z^2+1)+(Bz+C)z$, откуда $Bz^2+Cz=2z^2$, $B=2$, $C=0$. Так как
$\frac{2z}{z^2+1}=\frac1{z-i}+\frac1{z+i}$, то
\[
h(z):=\frac{g'(z)}{g(z)}=\frac1{2z}-\frac1z+\frac1{z-i}+\frac1{z+i}=-\frac1{2z}+\frac1{z-i}+\frac1{z+i}.
\]
(Проверка: $\ln g=\ln(z^2+1)-\ln2-\frac12\ln z$ локально, и производная этого выражения та же.)
\end{detail}
Итак, $h\in H(\CC\setminus P)$, $P:=\{0,i,-i\}$, и каждая точка $P$ --- простой полюс $h$ с вычетами $\Res_0h=-\frac12$, $\Res_{\pm i}h=1$.

\subsubsection*{Шаг 1. Исходный элемент, его каноничность и радиус}
$g(1)=1\cdot\Zh(1)=1$, поэтому функция $f(z):=\ln g(z)$ (главный логарифм) голоморфна в окрестности $1$ (там $g$ близка к $1$), $f(1)=0$, $f'=h$. Значит $F_0=E_1(f)$ ---
\emph{первообразный канонический элемент для $h$ в точке $1$}.

\emph{Радиус $F_0$ равен $1$.} Действительно, $h\in H(B(1,1))$ ($\dist(1,P)=\min(1,\sqrt2)=1$), поэтому по Лемме 28.1 (с $U=\CC\setminus P$, путь --- точка $1$) радиус $\ge1$.
Если бы радиус был $r>1$, то $f$ (точнее, сумма её ряда Тейлора) была бы голоморфна в $B(1,r)\ni0$, а её производная совпадала бы с $h$ на $B(1,1)$, значит (теорема единственности
в связном $B(1,\min(r,\sqrt2))\setminus\{0\}$) $h$ имела бы устранимую особенность в $0$ --- но это полюс. Итак, $F_0=(B(1,1),f)$.

\begin{remarkbox}{Явная формула}
В $B(1,1)$ $f(z)=\ln(z-i)+\ln(z+i)-\ln2-\tfrac12\ln z$ с главными логарифмами: все три слагаемых голоморфны в $B(1,1)$ (круг не пересекает лучи $\{x\pm i:x\le0\}$ и $R_\pi$),
производная равна $h$, а значение в $1$ равно $\ln\sqrt2-\frac{i\pi}4+\ln\sqrt2+\frac{i\pi}4-\ln2-0=0$. Так что $\Phi$ --- это «сумма» П.А.Ф. $\Ln(z-i)+\Ln(z+i)-\frac12\Ln z-\ln2$,
что даёт независимую проверку приведённых ниже ответов.
\end{remarkbox}

\subsubsection*{Шаг 2. Вдоль каких путей продолжается $F_0$; продолжающие семейства; $D_\Phi$}
\begin{itemize}
\item \emph{$F_0$ продолжается вдоль любого пути $\gamma$, $\gamma(0)=1$, $\gamma^*\subset\CC\setminus\{0,i,-i\}$.} Действительно, $\eps:=\dist(\gamma^*,P)>0$ (расстояние от компакта до конечного
множества, не пересекающего его), $U_\eps(\gamma^*)\subset\CC\setminus P$, $h\in H(U_\eps(\gamma^*))$, и по Предложению 15.5 $F_0$ продолжается вдоль $\gamma$.
Продолжающее семейство (для спрямляемого $\gamma$) явно:
\[
F_{\gamma,t}=\Bigl(B\bigl(\gamma(t),R(\gamma(t))\bigr),\ f_{\gamma,t}\Bigr),\quad f_{\gamma,t}(z)=\int_{\gamma_t}h\,d\zeta+\int_{[\gamma(t),z]}h\,d\zeta,\quad R(c):=\min\{|c|,|c-i|,|c+i|\}.
\]
Радиус ровно $R(\gamma(t))=\dist(\gamma(t),P)$: $\ge$ по Лемме 28.1, а $\le$ --- так как $f_{\gamma,t}'=h$ на пересечении круга элемента с $B(\gamma(t),R(\gamma(t)))$, и если бы круг элемента содержал
точку $p\in P$, то $h$ продолжалась бы голоморфно в $p$ (рассуждение шага 1).

\emph{Что за функции $f_{\gamma,t}$.} Функция $G(z):=g(z)^2=z\,\Zh(z)^2=\frac{(z^2+1)^2}{4z}$ однозначна, голоморфна и не обращается в нуль в $\CC\setminus P$, и $e^{2f}=G$ около $1$.
По Следствию 28.5 (применённому к $G$ и элементу $E_1(2f)$, продолжающемуся семейством $\{E_{\gamma(t)}(2f_{\gamma,t})\}$) $e^{2f_{\gamma,t}}=G$ около $\gamma(t)$. Значит $e^{f_{\gamma,t}}$ ---
одна из двух ветвей $\pm\sqrt{G}=\pm\sqrt z\,\Zh(z)$ около $\gamma(t)$, т.е. каждый элемент $\Phi$ --- ветвь логарифма функции $\sqrt z\,\Zh(z)$ при какой-то ветви корня.
Явно: $f_{\gamma,t}=\ell_{i,\gamma,t}+\ell_{-i,\gamma,t}-\frac12\ell_{0,\gamma,t}-\ln2$, где $\ell_{w,\gamma,t}(z)=\ln_{(\varphi_w(t))}(z-w)$, $\varphi_w$ --- непрерывная ветвь $\Arg(\gamma(t)-w)$ с $\varphi_w(0)=\arg(1-w)$
(описание П.А.Ф. $\Ln(z-w)$).
\item \emph{$F_0$ не продолжается вдоль путей, проходящих через $0$, $i$ или $-i$} --- по Лемме 28.2 ($P$ --- полюсы $h$).
\item \emph{$D_\Phi=\CC\setminus\{0,i,-i\}$.} Любая точка $c\notin P$ соединяется с $1$ путём в связной области $\CC\setminus P$, значит является центром элемента $\Phi$; точки $P$ не лежат ни в одном
круге, так как радиус элемента с центром $c$ равен $\dist(c,P)$.
\end{itemize}
Поскольку все пути продолжения лежат в $D_\Phi$ и вдоль \emph{любого} пути в $D_\Phi$ продолжение возможно, $\Phi$ является своей (единственной) аналитической ветвью над $D_\Phi$.
Тем самым $\Phi$ --- совокупность элементов логарифмов функций $\sqrt z\,\Zh(z)$ при \emph{обеих} ветвях корня (см. шаг 3).

\subsubsection*{Шаг 3. Элементы над точкой и связь между ними}
Фиксируем $b\in D_\Phi$ и элемент $F_b=(B(b,R(b)),f_b)\in\Phi$. Все элементы $\Phi$ с центром $b$ получаются из $F_b$ продолжением вдоль петель $\gamma$ в $D_\Phi$ с началом и концом $b$.
По Предложению 15.4 каждую петлю можно заменить вписанной ломаной (с тем же результатом продолжения и тем же индексом относительно точек $P$, если $\|\gamma-\lambda\|<\dist(\gamma^*,P)$), так что
можно считать $\gamma$ спрямляемой. По Следствию 28.3 $F_b$ переходит в $\bigl(B(b,R(b)),\,f_b+\int_\gamma h\,dz\bigr)$, и по теореме о вычетах (или прямо по определению индекса)
\[
\boxed{\ \int_\gamma h\,dz=2\pi i\bigl(\Ind_i(\gamma)+\Ind_{-i}(\gamma)\bigr)-\pi i\,\Ind_0(\gamma)\ }
\]
\begin{remarkbox}{О подсказке в задании}
Вклад точки $0$ складывается из $+\frac12\cdot2\pi i\,\Ind_0$ (от $\sqrt z$) и $-2\pi i\,\Ind_0$ (от $\Zh$, у которой в $0$ полюс), итого $-\pi i\,\Ind_0$; точки $\pm i$ (нули $\Zh$) дают $+2\pi i\,\Ind_{\pm i}$.
Это совпадает с формулой выше.
\end{remarkbox}
Значит каждый элемент над $b$ имеет вид $F_b+\pi in:=(B(b,R(b)),f_b+\pi in)$, $n\in\ZZ$ (правая часть формулы кратна $\pi i$). Обратно, все такие элементы получаются:
пусть $\beta$ --- спрямляемый путь в $D_\Phi$ из $b$ в точку $c$ с $0<|c|<\frac12$, $\kappa$ --- окружность $|z|=|c|$ с началом $c$, проходимая $|k|$ раз (против часовой стрелки при $k>0$, по ней при $k<0$).
Для петли $\beta\kappa\beta^-$ интегралы по $\beta$ и $\beta^-$ сокращаются, а $\int_\kappa h=-\pi ik$ ($\Ind_0\kappa=k$, $\Ind_{\pm i}\kappa=0$, так как $\pm i$ лежат вне круга $|z|\le|c|$). Поэтому эта петля
переводит $F_b$ в $F_b-\pi ik$, и все $n\in\ZZ$ реализуются.
\begin{thmbox}{Итог шага 3}
Над каждой точкой $b\in\CC\setminus\{0,\pm i\}$ имеется счётное множество элементов $\{F_b+\pi in:\ n\in\ZZ\}$ с одним и тем же кругом $B(b,\min\{|b|,|b-i|,|b+i|\})$; они попарно различны.
Петля $\gamma$ переводит $F_b+\pi in$ в $F_b+\pi i\bigl(n+2\Ind_i\gamma+2\Ind_{-i}\gamma-\Ind_0\gamma\bigr)$. Элементы с чётным $n$ --- ветви $\Ln(g_b)$, с нечётным --- ветви $\Ln(-g_b)$,
где $g_b=e^{f_b}$ --- одна из ветвей $\sqrt z\,\Zh(z)$ около $b$.
\end{thmbox}
\emph{Голоморфные ветви.} Область $\Omega_0:=\CC\setminus\bigl((-\infty,0]\cup\{x+i:x\le0\}\cup\{x-i:x\le0\}\bigr)$ односвязна (её дополнение в $\Ch$ --- три луча, сходящихся в $\infty$, т.е. связно), и
$\Phi$ распадается над $\Omega_0$ на голоморфные ветви
\[
f_n(z)=\ln(z-i)+\ln(z+i)-\ln2-\tfrac12\ln z+\pi in,\qquad n\in\ZZ
\]
(главные логарифмы; $f_0=f$ в $B(1,1)$, а по монодромии элементы над $\Omega_0$ --- в точности $E_c(f_n)$).

\subsubsection*{Шаг 4. Особые точки и их классификация для всех ветвей}
Особые точки --- $0$, $i$, $-i$ и $\infty$. В качестве проколотых окрестностей берём $V(0,1)$, $V(\pm i,1)$ (они лежат в $D_\Phi$: $|i-0|=1$, $|i-(-i)|=2$) и $V(\infty)=\{|z|>1\}$.
Для петли в каждой из этих окрестностей индексы относительно «чужих» точек $P$ находятся так: если петля лежит в круге $B(p,1)$, то точки вне этого круга (в частности, на его границе) лежат в
неограниченной компоненте дополнения к петле, и индекс относительно них $0$; если петля лежит в $\{|z|>1\}$, то связное множество $\{|z|\le1\}\supset P$ содержится в одной компоненте дополнения,
и $\Ind_0=\Ind_i=\Ind_{-i}=:N$.

\begin{itemize}
\item \textbf{Точка $0$.} Петля в $V(0,1)$ с $\Ind_0=N$ даёт сдвиг $-\pi iN$. Элемент $F_b+\pi in$ после $N$ обходов переходит в $F_b+\pi i(n-N)\ne F_b+\pi in$ при $N\ne0$, и все элементы над $b$
получаются из одного обходами вокруг $0$. Значит над $V(0,1)$ \emph{одна} аналитическая ветвь (совпадает со всей $\Phi|_{V(0,1)}$), и $0$ --- \emph{логарифмическая точка ветвления}.
(Около $0$: $f\approx-\frac12\Ln z$ + голоморфная функция, что согласуется с ответом.)
\item \textbf{Точки $\pm i$.} Петля в $V(i,1)$ с $\Ind_i=N$ даёт сдвиг $2\pi iN$. Поэтому элементы над $b\in V(i,1)$ разбиваются на \emph{две} аналитические ветви над $V(i,1)$:
$\{F_b+2\pi ik\}$ и $\{F_b+\pi i+2\pi ik\}$ (чётные и нечётные $n$); обходами внутри $V(i,1)$ чётность $n$ не меняется, а внутри класса все элементы достижимы. В каждой из двух ветвей $i$ --- \emph{логарифмическая точка ветвления}.
Для $-i$ --- всё то же самое: две аналитические ветви над $V(-i,1)$, у каждой логарифмическое ветвление.
(Около $\pm i$: $f\approx\Ln(z\mp i)$ + голоморфная функция.)
\item \textbf{Точка $\infty$.} Петля в $\{|z|>1\}$ с $\Ind_0=\Ind_{\pm i}=N$ даёт сдвиг $\pi i(2N+2N-N)=3\pi iN$. Поэтому над $\{|z|>1\}$ элементы разбиваются на \emph{три} аналитические ветви
$\{F_b+\pi i(j+3k):k\in\ZZ\}$, $j=0,1,2$, и у каждой $\infty$ --- \emph{логарифмическая точка ветвления} ($3\pi iN\ne0$ при $N\ne0$).
(Проверка: при $z\to\infty$ $\sqrt z\,\Zh(z)\sim\frac12z^{3/2}$, так что $f\approx\frac32\Ln z$, а обход $\infty$ меняет $\frac32\Ln z$ на $\pm3\pi i$.)
\end{itemize}

\begin{remarkbox}{Сводка по примеру}
$D_\Phi=\CC\setminus\{0,\pm i\}$; над каждой точкой --- счётно много элементов $F_b+\pi in$, $n\in\ZZ$; радиус $\dist(b,\{0,\pm i\})$.
Над $0$: одно логарифмическое ветвление. Над $i$ и над $-i$: по два логарифмических ветвления. Над $\infty$: три логарифмических ветвления.
\end{remarkbox}

\begin{remarkbox}{Коротко: план ответа на экзамене}
\begin{enumerate}
\item Определения: элемент, канонический элемент, НАП, продолжение вдоль пути; лемма о близких путях.
\item Формулировка Предложения 15.5: $h\in H(U_\eps(\gamma^*))$, первообразный элемент продолжается вдоль $\gamma$, $R_\gamma\ge\eps$.
\item Спрямляемый случай: формула $f_{\gamma,t}(z)=\int_{\gamma_t}h+\int_{[\gamma(t),z]}h$; НАП --- по теореме Коши в круге $B(\gamma(t),\eps)$.
\item Общий случай: вписанная ломаная $\lambda$, $\|\gamma-\lambda\|<\eps/3$, продолжение вдоль $\lambda$ с $R_\lambda\ge2\eps/3$, лемма о близких путях; оценка $R_\gamma\ge\eps$ через $f_{\gamma,t}'=h$.
\item Следствия: $\int_\gamma h=f_{\gamma,1}(b)$; И.Т.К. для гомотопных путей; продолжение $\ln g$ --- первообразный элемент для $g'/g$.
\item Пример: $h=(\ln(\sqrt z\Zh))'=-\frac1{2z}+\frac1{z-i}+\frac1{z+i}$; $F_0=(B(1,1),f)$; продолжение вдоль путей в $\CC\setminus\{0,\pm i\}$, и только вдоль них.
\item Приращение вдоль петли $2\pi i(\Ind_i+\Ind_{-i})-\pi i\Ind_0$; над точкой элементы $F_b+\pi in$.
\item Особые точки: $0$ --- одно логарифмическое ветвление; $\pm i$ --- по два; $\infty$ --- три.
\end{enumerate}
\end{remarkbox}

% -------------------------------------------------------------------
\ticket{29}{Теорема об аналитическом продолжении сложного элемента. Пример её применения для описания продолжения элемента $E_1\bigl(\ln(1+\sqrt z_{(\circ)})\bigr)$}
\label{t29:begin}

\begin{defs}
\textbf{Элементы и продолжение (Опр.~15.2--15.4).} Элемент $F=(B,f)$: $B=B(a,R)$, $f\in H(B)$; канонический, если $B$ --- круг сходимости ряда Тейлора $f$ в $a$;
$E_a(f)$ --- канонический элемент с центром $a$, заданный голоморфной около $a$ функцией $f$ (если $f$ голоморфна в $B(a,r)$, то радиус $E_a(f)$ не меньше $r$).
$F_1$ --- НАП $F_0$, если круги пересекаются и функции совпадают на пересечении (пересечение кругов выпукло, поэтому достаточно совпадения на непустом открытом его подмножестве).
Семейство канонических $F_t=(B(\gamma(t),R_t),f_t)$, $t\in I=[0,1]$, --- продолжение $F_0$ вдоль пути $\gamma$, если у каждого $t$ есть связная окрестность $U_t\subset I$,
для $\xi\in U_t$ которой $\gamma(\xi)\in B_t$ и $F_\xi$ --- НАП $F_t$. \emph{Единственность (Предл.~15.2):} продолжающее семейство однозначно определено (при всех $t$).
\emph{Склейка/обращение:} продолжение вдоль $\gamma_1\gamma_2$ --- последовательное продолжение; вдоль $\gamma^-$ --- обратный переход.

\textbf{Простейший случай.} Если $G\in H(\Omega)$, $\Omega$ --- область, $\Gamma\subset\Omega$ --- путь, то семейство $E_{\Gamma(t)}(G)$ продолжает $E_{\Gamma(0)}(G)$ вдоль $\Gamma$
(круг $E_{\Gamma(t)}(G)$ содержит $B(\Gamma(t),\dist(\Gamma(t),\partial\Omega))$, где ряд равен $G$; для $\xi$ близких к $t$ функции элементов совпадают с $G$ около $\Gamma(\xi)$).
В частности, продолжение вдоль петли в $\Omega$ возвращает элемент в себя.

\textbf{П.А.Ф., область определения, ветви, особые точки (Опр.~15.5--15.9, билет 27).} П.А.Ф. $\Phi$ --- все канонические элементы, получаемые из данного продолжением вдоль путей;
$D_\Phi$ --- объединение их кругов. Все элементы над точкой $b$ получаются из одного $F_b$ продолжением вдоль петель с началом $b$.
Аналитическая ветвь над областью $D_1\subset D_\Phi$ --- всё, что получается из $F_b$ продолжением по путям в $D_1$ (если $F_b$ продолжается вдоль всех таких путей);
голоморфная ветвь --- $f\in H(D_2)$ с $E_c(f)$ в ветви при всех $c\in D_2$; распадение на голоморфные ветви.
Изолированная особая точка $a$ ветви над $V(a,\delta)\subset D_\Phi$: однозначного характера (при обходе элемент возвращается в себя; затем --- правильная/устранимая, полюс, существенно особая),
ветвление порядка $n\ge2$ (возврат впервые после $n$ обходов), логарифмическое ветвление (возврата нет).

\textbf{Композиция П.А.Ф. (слайд 377).} Если $\Phi$ и $\Psi$ --- П.А.Ф., то «$\Psi\circ\Phi$» --- совокупность элементов $E_a(g\circ f)$, где $(B,f)\in\Phi$ с центром $a$, $(\widetilde B,g)\in\Psi$ с центром $f(a)$,
и их продолжений; корректность обеспечивает Предложение 15.6. Например, $w^2\circ\sqrt[4]{z}$ --- это П.А.Ф. $\sqrt z$ (а произведение $\sqrt[4]z\cdot\sqrt[4]z$ даёт две П.А.Ф. $\sqrt z$ и $i\sqrt z$).

\textbf{Обозначения.} $\sqrt z_{(\circ)}$ --- главная ветвь корня в $\CC\setminus R_\pi$, $R_\pi=(-\infty,0]$; $\ln$ --- главный логарифм ($\arg\in(-\pi,\pi)$);
$\ln_{(\alpha)}w=\ln|w|+i\arg_{(\alpha)}w$, $\arg_{(\alpha)}w\in(\alpha-\pi,\alpha+\pi)$ --- ветвь логарифма в $\CC\setminus\{-te^{i\alpha}:t\ge0\}$.
Продолжение $E_a(\sqrt z_{(\circ)})$ ($a>0$) вдоль пути $\gamma\subset\CC^*$: элементы $\bigl(B(\gamma(t),|\gamma(t)|),\ \sqrt{|z|}e^{i\arg_{(\varphi(t))}z/2}\bigr)$, $\varphi$ --- непрерывная ветвь $\Arg\gamma(t)$, $\varphi(0)=0$;
продолжение $E_A(\ln(1+w))$ ($A>-1$) вдоль пути $\Gamma\subset\CC\setminus\{-1\}$: элементы $\bigl(B(\Gamma(t),|\Gamma(t)+1|),\ \ln_{(\psi(t))}(w+1)\bigr)$, $\psi$ --- непрерывная ветвь $\Arg(\Gamma(t)+1)$, $\psi(0)=0$
(описание П.А.Ф. $(z-w)^{1/2}$ и $\Ln(z-w)$ из билета 27).
\end{defs}

\subsection{Теорема о продолжении сложного элемента}

\begin{thmbox}{Предложение 15.6 (теорема о продолжении сложного элемента)}
Пусть $F_a=(B_a,f_a)$ --- канонический элемент с центром $a$, который продолжается вдоль пути $\gamma\colon[0,1]\to\CC$, $\gamma(0)=a$, $\gamma(1)=b$, семейством
$\{F_t\}$, $F_t=(B_t,f_t)$, $B_t=B(\gamma(t),R_t)$. Положим
\[
\Gamma(t):=f_t(\gamma(t)),\qquad t\in[0,1].
\]
Пусть элемент $G_A=(B_A,g_A)$ с центром $A=\Gamma(0)$ продолжается вдоль $\Gamma$ семейством $\{G_t\}$, $G_t=(\widetilde B_t,g_t)$, $\widetilde B_t=B(\Gamma(t),\widetilde R_t)$.
Тогда элемент $H_a:=E_a(g_A\circ f_a)$ продолжается вдоль $\gamma$, и продолжающее семейство имеет вид
\[
H_t=E_{\gamma(t)}(g_t\circ f_t),\qquad t\in[0,1].
\]
\end{thmbox}

\begin{proof}
\textbf{(а) $\Gamma$ --- путь.} Фиксируем $t$. Для $\xi\in U_t$ (окрестность из определения продолжения для $\{F_t\}$) имеем $\gamma(\xi)\in B_t\cap B_\xi$ и $f_\xi=f_t$ на $B_t\cap B_\xi$, поэтому
$\Gamma(\xi)=f_\xi(\gamma(\xi))=f_t(\gamma(\xi))$. Так как $f_t$ непрерывна в $B_t$, а $\gamma$ непрерывна, то $\Gamma(\xi)\to f_t(\gamma(t))=\Gamma(t)$ при $\xi\to t$. Итак, $\Gamma$ непрерывен на $I$.
В частности, условие теоремы (продолжение $G_A$ вдоль $\Gamma$) корректно.

\textbf{(б) Элементы $H_t$ определены.} $f_t(\gamma(t))=\Gamma(t)$ --- центр круга $\widetilde B_t$; по непрерывности $f_t$ найдётся круг $B(\gamma(t),\delta_t)\subset B_t$ с $f_t(B(\gamma(t),\delta_t))\subset\widetilde B_t$,
и $g_t\circ f_t$ голоморфна в нём. Значит $H_t=E_{\gamma(t)}(g_t\circ f_t)$ --- канонический элемент с центром $\gamma(t)$ и радиусом $\ge\delta_t>0$. При $t=0$: $H_0=E_a(g_A\circ f_a)=H_a$.

\textbf{(в) Условие НАП.} Фиксируем $t$; пусть $z_0=\gamma(t)$, $w_0=\Gamma(t)=f_t(z_0)$. Выберем $\eps>0$ с $\widetilde B^0:=B(w_0,\eps)\subset\widetilde B_t$, затем $\delta>0$ с
$B^0:=B(z_0,\delta)\subset B_t$ и $f_t(B^0)\subset\widetilde B^0$ (непрерывность $f_t$). Тогда $h:=g_t\circ f_t\in H(B^0)$ и круг элемента $H_t$ содержит $B^0$.
По определению продолжения и непрерывности $\gamma,\Gamma$ найдётся связная окрестность $W\subset U_t\cap\widetilde U_t$ точки $t$ (здесь $\widetilde U_t$ --- окрестность для семейства $\{G_t\}$) такая, что
при $\xi\in W$: $\gamma(\xi)\in B^0$ и $\Gamma(\xi)\in\widetilde B^0$. Пусть $\xi\in W$. Тогда:
\begin{itemize}
\item $\gamma(\xi)\in B^0\subset$ круг $H_t$ --- первое требование выполнено;
\item $f_\xi=f_t$ на $B_\xi\cap B_t\supset B_\xi\cap B^0$, а $g_\xi=g_t$ на $\widetilde B_\xi\cap\widetilde B_t\supset\widetilde B_\xi\cap\widetilde B^0$;
\item множество $\widetilde B_\xi\cap\widetilde B^0$ открыто и содержит $\Gamma(\xi)=f_\xi(\gamma(\xi))$; по непрерывности $f_\xi$ найдётся круг $D_\xi=B(\gamma(\xi),\rho)\subset B_\xi\cap B^0$ с $f_\xi(D_\xi)\subset\widetilde B_\xi\cap\widetilde B^0$,
причём $\rho$ можно взять меньше радиуса $H_\xi$.
\end{itemize}
Для $z\in D_\xi$: $f_\xi(z)=f_t(z)=:w\in\widetilde B_\xi\cap\widetilde B^0$, и $g_\xi(w)=g_t(w)$, т.е. $g_\xi(f_\xi(z))=g_t(f_t(z))=h(z)$. Итак, функции элементов $H_\xi$ и $H_t$ совпадают на непустом
открытом множестве $D_\xi$, лежащем в пересечении их кругов, а значит (выпуклость пересечения, теорема единственности) на всём пересечении: $H_\xi$ --- НАП $H_t$.
Окрестность $W$ связна. Все условия Определения 15.4 выполнены.
\end{proof}

\begin{remarkbox}{Замечание к слайду}
На слайде сказано «$F_\xi$ --- НАП $F_t$, откуда $f_\xi=f_t$ всюду в $B_0$»; но $f_\xi$ определена лишь в $B_\xi$, а $B_0=B(\gamma(t),\delta)$ может не содержаться в $B_\xi$. Это
исправляется (как выше) рассмотрением малого круга $D_\xi$ около $\gamma(\xi)$; либо выбором $\delta<R_t/3$ и использованием непрерывности радиуса (Лемма 15.1). Аналогично для $g_\xi$.
\end{remarkbox}

\begin{remarkbox}{Следствия (слайд 378)}
Если $D$ односвязна, $f\in H(D)$, $f(D)\subset\CC^*$, то композиции $\sqrt w\circ f$ и $\Ln w\circ f$ (по Предложению 15.6 элементы продолжаются вдоль любых путей в $D$, так как $f\circ\gamma$
не проходит через $0$) по теореме о монодромии распадаются над $D$ на голоморфные ветви $g$ и $h$, для которых $g^2=f$, $e^h=f$ (тождества сохраняются при продолжении по теореме единственности).
Упражнение со слайда: композиция $\sqrt[4]w\circ z^2$, начатая с $E_1(\sqrt[4]{z^2})$ (значение $1$ в $1$), --- это П.А.Ф. $\sqrt z$: для её элементов $h^4=z^2$, т.е. $(h^2-z)(h^2+z)=0$
на круге, значит один из множителей $\equiv0$, и тождество $h^2=z$, верное в начале, сохраняется при НАП. Если же брать все элементы $\sqrt[4]w$ над $w=1$, то получатся две П.А.Ф.: $\sqrt z$ и $i\sqrt z$.
\end{remarkbox}

\subsection{Пример: П.А.Ф. $\Ln(1+\sqrt z)$}

Пусть $h(z):=\ln(1+\sqrt z_{(\circ)})$. При $z\in\CC\setminus R_\pi$ имеем $\Re\sqrt z_{(\circ)}>0$, $\Re(1+\sqrt z_{(\circ)})>1$, поэтому $h\in H(\CC\setminus R_\pi)$ --- голоморфная функция.
Опишем П.А.Ф. $\Phi$, порождённую $H_1:=E_1(h)$, по схеме слайда 380.

\subsubsection*{Шаг 0. Алгебраическое тождество}
Для $h$ верно $(e^{h(z)}-1)^2=z$. Это тождество сохраняется при НАП (если оно верно для функции элемента на открытом подмножестве круга, то по теореме единственности --- на всём круге),
а значит выполнено для функции $q$ \emph{любого} элемента $(B,q)\in\Phi$: $s:=e^{q}-1$ голоморфна в $B$ и $s^2=z$, т.е. $s$ --- голоморфная ветвь корня в $B$, а $q$ --- ветвь логарифма $1+s$.
Отсюда два ограничения на элементы:
\begin{enumerate}
\item[(О1)] круг элемента не содержит $0$ (в круге, содержащем $0$, нет голоморфной ветви $\sqrt z$: из $s^2=z$ следует $s(0)=0$, и $s^2$ имела бы в $0$ нуль чётного порядка, а $z$ --- простой);
\item[(О2)] если круг элемента содержит точку $1$, то $s(1)=1$ (иначе $s(1)=-1$ и $e^{q(1)}=1+s(1)=0$, что невозможно).
\end{enumerate}

\subsubsection*{Шаг 1. Исходный элемент: каноничность и радиус}
$H_1=E_1(h)$, $h(1)=\ln2$. Радиус $\ge1$, так как $h$ голоморфна в $B(1,1)\subset\CC\setminus R_\pi$; и $\le1$ по (О1) (круг $B(1,r)$ при $r>1$ содержит $0$). Итак, $H_1=(B(1,1),h)$
(здесь $h=\ln(1+\sqrt z_{(\circ)})$ --- сумма ряда Тейлора, т.е. элемент канонический по построению).

\subsubsection*{Шаг 2. Вдоль каких путей продолжается $H_1$}
Представим $h=g\circ f$ около $1$, где $f=\sqrt z_{(\circ)}$ (внутренний элемент $F_1=E_1(\sqrt z_{(\circ)})$, радиус $1$) и $g(w)=\ln(1+w)$ (внешний элемент $G_1=E_1(\ln(1+w))$, радиус $2$; $A=f(1)=1$).

\emph{Пусть $\gamma\colon I\to\CC^*$, $\gamma(0)=1$.} Внутренний элемент продолжается вдоль $\gamma$: $F_t=\bigl(B(\gamma(t),|\gamma(t)|),f_t\bigr)$, $f_t(z)=\sqrt{|z|}e^{i\arg_{(\varphi(t))}z/2}$, где $\varphi$ ---
непрерывная ветвь $\Arg\gamma(t)$, $\varphi(0)=0$. Путь
\[
\Gamma(t)=f_t(\gamma(t))=\sqrt{|\gamma(t)|}\,e^{i\varphi(t)/2}.
\]
Внешний элемент $G_1$ продолжается вдоль $\Gamma$ тогда и только тогда, когда $\Gamma$ не проходит через $-1$ (П.А.Ф. $\Ln(w+1)$), семейством
$G_t=\bigl(B(\Gamma(t),|\Gamma(t)+1|),\ln_{(\psi(t))}(w+1)\bigr)$, $\psi$ --- непрерывная ветвь $\Arg(\Gamma(t)+1)$, $\psi(0)=0$. Условие $\Gamma(t)=-1$ равносильно
$|\gamma(t)|=1$ и $e^{i\varphi(t)/2}=-1$, т.е.
\[
\gamma(t)=1\quad\text{и}\quad\varphi(t)\in2\pi+4\pi\ZZ
\]
--- путь приходит в точку $1$ «по второму листу» корня (значение корня, продолженного вдоль $\gamma_t$, равно $-1$).

\begin{thmbox}{Утверждение 29.1 (продолжаемость $H_1$)}
$H_1$ продолжается вдоль пути $\gamma$, $\gamma(0)=1$, тогда и только тогда, когда (i) $\gamma$ не проходит через $0$ и (ii) ни при каком $t$ не выполнено одновременно $\gamma(t)=1$ и
$\varphi(t)\equiv2\pi\pmod{4\pi}$. В этом случае продолжающее семейство --- $H_t=E_{\gamma(t)}\bigl(\ln_{(\psi(t))}(1+f_t(z))\bigr)$.
\end{thmbox}
\begin{proof}
\emph{Достаточность} --- Предложение 15.6 с описанными выше семействами $\{F_t\}$, $\{G_t\}$.

\emph{Необходимость.} Пусть $\{H_t=(B^H_t,q_t)\}$ продолжает $H_1$ вдоль $\gamma$. Если при некотором $t$ $\gamma(t)=0$, то круг $B^H_t$ содержит $0$ --- противоречие с (О1).
Значит $\gamma\subset\CC^*$, и непрерывная ветвь $\varphi$ определена на всём $I$. Предположим, что множество «плохих» моментов
$T:=\{t:\ \gamma(t)=1,\ \varphi(t)\in2\pi+4\pi\ZZ\}$ непусто. Оно замкнуто (прообраз замкнутого множества $\{1\}\times(2\pi+4\pi\ZZ)$ при непрерывном $t\mapsto(\gamma(t),\varphi(t))$);
пусть $\tau=\min T$; $\tau>0$, так как $\varphi(0)=0$. Для $\xi<\tau$ путь $\gamma_\xi$ удовлетворяет (i), (ii), поэтому по достаточности и единственности (сужение семейства $\{H_t\}$ на $[0,\xi]$)
$H_\xi=E_{\gamma(\xi)}(g_\xi\circ f_\xi)$, т.е. $q_\xi=g_\xi\circ f_\xi$ около $\gamma(\xi)$ и $s_\xi:=e^{q_\xi}-1=f_\xi$ около $\gamma(\xi)$.
Круг $B^H_\tau$ содержит $\gamma(\tau)=1$, и по (О2) $s_\tau(1)=1$, где $s_\tau=e^{q_\tau}-1\in H(B^H_\tau)$. С другой стороны, для $\xi<\tau$, $\xi\in U_\tau$, элемент $H_\xi$ --- НАП $H_\tau$ и
$\gamma(\xi)\in B^H_\tau\cap B^H_\xi$, поэтому $s_\tau=s_\xi=f_\xi$ около $\gamma(\xi)$ и $s_\tau(\gamma(\xi))=\Gamma(\xi)$. Устремляя $\xi\to\tau-0$ (непрерывность $s_\tau$, $\gamma$, $\Gamma$), получаем
$s_\tau(1)=\Gamma(\tau)=e^{i\varphi(\tau)/2}=-1$ --- противоречие. 
\end{proof}

\subsubsection*{Шаг 3. Радиусы, область определения, элементы над точкой}
\textbf{Радиусы.} Пусть $(B(c,R),q)\in\Phi$, $s=e^q-1$ --- соответствующая ветвь корня около $c$; она голоморфно продолжается на $B(c,|c|)$ (односвязный круг без нуля). Тогда
\[
R=\begin{cases}|c-1|,& \text{если }1\in B(c,|c|)\ (\Leftrightarrow\Re c>\tfrac12)\text{ и }s(1)=-1,\\ |c|,&\text{иначе.}\end{cases}
\]
Действительно, $R\le|c|$ по (О1) и $R\le|c-1|$ в первом случае по (О2). Обратно, в круге $B^*$ указанного радиуса функция $1+s$ голоморфна и не обращается в нуль
($1+s(z)=0\Leftrightarrow s(z)=-1\Rightarrow z=1$ и $s(1)=-1$, а такой точки в $B^*$ нет), поэтому имеет голоморфный логарифм в круге $B^*$, отличающийся от $q$ около $c$ на константу $2\pi ik$;
значит $q$ голоморфно продолжается в $B^*$ и $R\ge$ радиуса $B^*$.

\textbf{Область определения.} $D_\Phi=\CC^*=\CC\setminus\{0\}$: каждая точка $c\neq0$ --- центр элемента (путь из $1$ в $c$ в $\CC^*$, не проходящий через $1$ повторно, удовлетворяет условиям
Утверждения 29.1), а $0$ ни в один круг не попадает (О1). Заметим, что \emph{$\Phi$ не является аналитической ветвью над $D_\Phi$} (не вдоль всякого пути в $\CC^*$ элементы продолжаются),
но является аналитической ветвью над $\CC\setminus\{0,1\}$: любой элемент с центром $b\ne0,1$ продолжается вдоль любого пути в $\CC\setminus\{0,1\}$ (по Предл.~15.6 --- путь $\Gamma$ не попадает в $-1$, так как $\Gamma^2=\gamma\ne1$).

\textbf{Элементы над точкой $b\neq0,1$.} Пусть $r$ --- какая-нибудь ветвь $\sqrt z$ в $B(b,|b|)$, $L_\pm$ --- какие-нибудь ветви $\ln(1\pm r)$ около $b$ ($1\pm r(b)\neq0$, т.к. $b\ne1$).
По шагу 0 каждый элемент над $b$ имеет вид
\[
E_b(L_++2\pi ik)\quad\text{или}\quad E_b(L_-+2\pi ik),\qquad k\in\ZZ,
\]
(функция элемента --- логарифм $1+r$ или $1-r$, а два логарифма одной функции на круге отличаются на $2\pi ik$). Все эти элементы попарно различны, и \emph{все они принадлежат $\Phi$} (см. ниже, шаг 4, пункт «$\infty$»).
\emph{Над точкой $b=1$} элементы только вида $E_1(h+2\pi ik)$, $k\in\ZZ$: вариант $1-\sqrt z$ запрещён (О2). Итак, над каждой точкой $D_\Phi$ --- счётное множество элементов
(в согласии с теоремой Пуанкаре--Вольтерра), над $1$ --- «вдвое меньше».

\emph{Голоморфные ветви.} Над $\CC\setminus R_\pi$ --- голоморфные ветви $\ln(1+\sqrt z_{(\circ)})+2\pi ik$. В односвязной области $\Omega:=\CC\setminus\bigl((-\infty,0]\cup[1,+\infty)\bigr)$ функция $1-\sqrt z_{(\circ)}$
не принимает значений из $(-\infty,0]$ ($1-\sqrt z\le0\Leftrightarrow\sqrt z\ge1\Leftrightarrow z\in[1,\infty)$), так что $\Phi$ распадается над $\Omega$ на голоморфные ветви
\[
\ln\bigl(1+\sqrt z_{(\circ)}\bigr)+2\pi ik,\qquad \ln\bigl(1-\sqrt z_{(\circ)}\bigr)+2\pi ik,\qquad k\in\ZZ .
\]

\subsubsection*{Шаг 4. Особые точки $0$, $1$, $\infty$ и их классификация для всех ветвей}

\textbf{Точка $0$; окрестность $V(0,1)$.} При $0<|z|<1$ и любой ветви корня $|\pm\sqrt z|<1$, т.е. $\Gamma$ лежит в $\DD$, где $G(w):=\ln(1+w)$ голоморфна. Поэтому для элемента
$E_b(G(\sigma r)+2\pi ik)$, $b\in V(0,1)$, $\sigma=\pm1$, продолжение вдоль петли $\gamma\subset V(0,1)$ с $\Ind_0\gamma=N$ по Предл.~15.6 (внешнее семейство --- переразложения однозначной функции $G+2\pi ik$ в $\DD$)
даёт $E_b\bigl(G((-1)^N\sigma r)+2\pi ik\bigr)$: внутренний корень меняет знак при каждом обходе, а $k$ не меняется. Значит над $V(0,1)$ $\Phi$ распадается на \emph{счётное число} аналитических ветвей
\[
A_k=\bigl\{\ln(1+\sqrt z)+2\pi ik,\ \ln(1-\sqrt z)+2\pi ik\bigr\},\quad k\in\ZZ,
\]
и у каждой из них $0$ --- \emph{точка ветвления второго порядка}. (По Предложению 15.8: при $z=w^2$ ветвь $A_k$ превращается в однозначную функцию $\ln(1+w)+2\pi ik$, $|w|<1$.)

\textbf{Точка $1$; окрестность $V(1,1)$.} Круг $B(1,1)$ лежит в правой полуплоскости, $\sqrt z_{(\circ)}$ голоморфна в нём, внутренний корень вдоль петель в $V(1,1)$ не меняется.
\begin{itemize}
\item Элементы $E_b(\ln(1+\sqrt z_{(\circ)})+2\pi ik)$: функции $\ln(1+\sqrt z_{(\circ)})+2\pi ik$ голоморфны во всём $B(1,1)$. Получаем \emph{счётное число однозначных (голоморфных) ветвей}, для которых $1$ ---
правильная точка (принадлежит области голоморфности).
\item Элементы с $1-\sqrt z_{(\circ)}$: их производная
\[
h_-(z)=\frac{d}{dz}\ln\bigl(1-\sqrt z_{(\circ)}\bigr)=\frac{-1}{2\sqrt z(1-\sqrt z)}=\frac1{z-1}-\frac1{2\sqrt z(1+\sqrt z)}
\]
однозначна и голоморфна в $V(1,1)$ (проверка: $1-\sqrt z=\frac{1-z}{1+\sqrt z}$, логарифмическая производная $\frac{-1}{1-z}-\frac{1/(2\sqrt z)}{1+\sqrt z}$). Значит эти элементы --- первообразные для $h_-$ и по Следствию 28.3
петля $\gamma\subset V(1,1)$ сдвигает их на $\int_\gamma h_-=2\pi i\,\Ind_1\gamma$ (второе слагаемое голоморфно в $B(1,1)$). Все они образуют \emph{одну} аналитическую ветвь
$\{\ln(1-\sqrt z_{(\circ)})+2\pi ik:k\in\ZZ\}$ над $V(1,1)$, и у неё $1$ --- \emph{логарифмическая точка ветвления}.
\end{itemize}
Итак, над $1$: счётно много голоморфных ветвей (правильная точка) и одно логарифмическое ветвление.

\textbf{Точка $\infty$; окрестность $\{|z|>1\}$.} Для $z\in\{|z|>1\}\setminus R_\pi$ и $m\in\ZZ$ положим
\[
\Psi_m(z):=\tfrac12\bigl(\ln z+2\pi im\bigr)+\ln\Bigl(1+\frac{(-1)^m}{\sqrt z_{(\circ)}}\Bigr)
\]
(второй логарифм главный: $|1/\sqrt z|<1$). Тогда $e^{\Psi_m}=(-1)^m\sqrt z\,\bigl(1+(-1)^m/\sqrt z\bigr)=1+(-1)^m\sqrt z$, т.е. $\Psi_m$ --- ветвь $\ln(1\pm\sqrt z)$, и $\Psi_{m+2}=\Psi_m+2\pi i$.
Значит $\{\Psi_m:m\in\ZZ\}$ --- это в точности все функции $L_\pm+2\pi ik$ из шага 3 (чётные $m$ --- знак $+$, нечётные --- знак $-$), и они попарно различны.
Монодромия: $\Psi_m=W\circ\ell_m$, где $\ell_m=\ln z+2\pi im$ --- ветвь $\Ln z$, а $W(\zeta):=\frac\zeta2+\ln\bigl(1+e^{-\zeta/2}\bigr)$ голоморфна (однозначна) в полуплоскости $\{\Re\zeta>0\}$
(там $|e^{-\zeta/2}|<1$). Для петли $\gamma\subset\{|z|>1\}$ с $\Ind_0\gamma=N$ путь $\ell(\gamma(t))$ лежит в $\{\Re\zeta=\ln|\gamma(t)|>0\}$, элемент $\ln$ переходит в $\ell_m+2\pi iN$ (П.А.Ф. $\Ln z$), а внешнее
семейство --- переразложения однозначной $W$. По Предложению 15.6 $E_b(\Psi_m)$ переходит в $E_b(\Psi_{m+N})$. Следовательно:
\begin{itemize}
\item над $\{|z|>1\}$ \emph{все} элементы $\Phi$ образуют \emph{одну} аналитическую ветвь, и $\infty$ для неё --- \emph{логарифмическая точка ветвления} ($\Psi_{m+N}\ne\Psi_m$ при $N\ne0$);
при этом каждый обход $\infty$ меняет лист корня, а два обхода прибавляют $2\pi i$;
\item в частности, все элементы $E_b(L_\pm+2\pi ik)$ при $|b|>1$ лежат в $\Phi$ (так как $E_b(\Psi_0)=E_b(\ln(1+\sqrt z_{(\circ)}))\in\Phi$ при $b>1$). Для произвольного $b\ne0,1$ соединим $b$ путём $\beta\subset\CC\setminus\{0,1\}$
с точкой $b'$, $|b'|>1$: любой кандидат $E_b(L_\pm+2\pi ik)$ продолжается вдоль $\beta$ (Предл.~15.6) в некоторый элемент над $b'$, который лежит в $\Phi$; продолжая обратно вдоль $\beta^-$, получаем,
что и сам кандидат лежит в $\Phi$. Над $b=1$ элементы $E_1(h+2\pi ik)$ получаются так: из $1$ по отрезку в $2$, $2k$ раз по окружности $|z|=2$, обратно в $1$.
\end{itemize}

\begin{remarkbox}{О подсказке в задании}
Над $\infty$ получилось \emph{одно} логарифмическое ветвление (а не несколько): обход $\infty$ переставляет листы корня, поэтому элементы $\ln(1+\sqrt z)+2\pi ik$ и $\ln(1-\sqrt z)+2\pi ik$ попадают в одну
аналитическую ветвь над $\{|z|>1\}$. Это согласуется с асимптотикой $\ln(1\pm\sqrt z)=\frac12\Ln z+O(|z|^{-1/2})$: над $\infty$ «работает» $\frac12\Ln z$, у которого одно логарифмическое ветвление.
\end{remarkbox}

\begin{remarkbox}{Сводка по примеру $\Ln(1+\sqrt z)$}
$D_\Phi=\CC^*$; над $b\ne0,1$ элементы $\ln(1\pm\sqrt z)+2\pi ik$ ($k\in\ZZ$), радиус $|b|$, кроме элементов с $1-\sqrt z_{(\circ)}$ при $\Re b>\frac12$, у которых радиус $|b-1|$; над $1$ --- только $\ln(1+\sqrt z)+2\pi ik$.
Над $0$: счётно много ветвлений второго порядка. Над $1$: счётно много однозначных (голоморфных) ветвей, одно логарифмическое ветвление. Над $\infty$: одно логарифмическое ветвление.
\end{remarkbox}

\subsection{Образец из лекций: П.А.Ф. $\sqrt{1+\sqrt z}$ (кратко)}
Исходный элемент $E_1(h)$, $h=\sqrt{1+\sqrt z_{(\circ)}}_{(\circ)}\in H(\CC\setminus R_\pi)$; тождество $(h^2-1)^2=z$ сохраняется при продолжении, поэтому над $b\ne0,1$ не больше четырёх элементов
($h_b\in\{\pm\sqrt{1+r},\pm\sqrt{1-r}\}$). По Предложению 15.6 (внутренний $\sqrt z$, внешний $\sqrt{1+w}$) элемент продолжается вдоль путей в $\CC^*$, не приходящих в $1$ по второму листу корня.
Ответ: особые точки $0,1,\infty$; над каждой $b\ne0,1$ --- четыре элемента, попарно отличающиеся знаком;
над $0$ --- два ветвления 2-го порядка ($\{\sqrt{1+\sqrt z},\sqrt{1-\sqrt z}\}$ и с противоположным внешним знаком);
над $1$ --- две однозначные ветви $\pm h$ и одно ветвление 2-го порядка ($\pm\sqrt{1-\sqrt z}=\pm\sqrt{1-z}/\sqrt{1+\sqrt z}$, причём $\sqrt{1-z}$ ветвится в $1$ с порядком $2$);
над $\infty$ --- одно ветвление 4-го порядка ($h=\sqrt[4]z\sqrt{1+1/\sqrt z}$, обход $\infty$ последовательно даёт $i\sqrt[4]z\sqrt{1-1/\sqrt z}$, $-h$, $-i\sqrt[4]z\sqrt{1-1/\sqrt z}$, $h$).
Сравнение: в примере $\Ln(1+\sqrt z)$ внешняя функция --- логарифм, поэтому «ветвление 2-го порядка» над $1$ становится логарифмическим, а ветвление порядка $4$ над $\infty$ --- логарифмическим.

\begin{remarkbox}{Коротко: план ответа на экзамене}
\begin{enumerate}
\item Определения: элемент, канонический элемент, НАП, продолжение вдоль пути, единственность.
\item Формулировка Предложения 15.6: $\Gamma(t)=f_t(\gamma(t))$, если $G_A$ продолжается вдоль $\Gamma$, то $E_a(g_A\circ f_a)$ продолжается вдоль $\gamma$ семейством $E_{\gamma(t)}(g_t\circ f_t)$.
\item Доказательство: $\Gamma$ непрерывен ($\Gamma(\xi)=f_t(\gamma(\xi))$ при $\xi$ близких к $t$); выбор кругов $B(\Gamma(t),\eps)$, $B(\gamma(t),\delta)$ с $f_t(B^0)\subset\widetilde B^0$; совпадение $g_\xi\circ f_\xi=g_t\circ f_t$ около $\gamma(\xi)$.
\item Применения: композиция П.А.Ф.; $\sqrt f$, $\ln f$ в односвязной области.
\item Пример: $\Ln(1+\sqrt z)=\Ln w\circ(1+\sqrt z)$; тождество $(e^q-1)^2=z$; продолжение вдоль путей в $\CC^*$, не приходящих в $1$ со значением корня $-1$.
\item $D_\Phi=\CC^*$; над точкой --- элементы $\ln(1\pm\sqrt z)+2\pi ik$; над $1$ --- только знак $+$.
\item $0$: счётно много ветвлений 2-го порядка; $1$: счётно много голоморфных ветвей и одно логарифмическое ветвление; $\infty$: одно логарифмическое ветвление.
\end{enumerate}
\end{remarkbox}

% ===== part 6 =====
% ===================== Часть 6: билеты 30--32 (алгебраические П.А.Ф., многообразия, римановы поверхности) =====================
\ticket{30}{Понятие об алгебраической особой точке П.А.Ф. и об алгебраических П.А.Ф.}

\begin{defs}
\begin{itemize}
\item \textbf{Обозначения.} $B(a,r)=\{|z-a|<r\}$ ($r\le+\infty$), $V(a,\delta)=B(a,\delta)\setminus\{a\}$ --- проколотый круг;
  для $a=\infty$ полагаем $V(\infty,\delta)=\{|z|>1/\delta\}$. Для функции $f$, голоморфной в окрестности точки $a$,
  $E_a(f)$ --- канонический элемент с центром $a$, порождённый рядом Тейлора $f$ в точке $a$.
\item \textbf{Элемент} (Опр.~15.2): пара $F=(B,f)$, $B=B(a,R)$, $f\in H(B)$; $a$ --- центр, $R$ --- радиус. $F$ \emph{канонический},
  если $B$ --- круг сходимости ряда Тейлора $f$ в точке $a$. По теореме Коши--Тейлора, если $f\in H(B(a,R))$, то радиус
  канонического элемента $E_a(f)$ не меньше $R$.
\item \textbf{НАП} (Опр.~15.3): $F_1=(B_1,f_1)$ --- непосредственное аналитическое продолжение $F_0=(B_0,f_0)$, если
  $B_0\cap B_1\ne\varnothing$ и $f_0=f_1$ на $B_0\cap B_1$. Пересечение двух кругов выпукло, значит связно; поэтому
  \emph{по теореме единственности для проверки НАП достаточно совпадения $f_0$ и $f_1$ на непустом открытом подмножестве $B_0\cap B_1$.}
\item \textbf{Продолжение вдоль пути} (Опр.~15.4): семейство канонических элементов $F_t=(B(\gamma(t),R_t),f_t)$, $t\in[0,1]$,
  такое, что для каждого $t$ найдётся связная окрестность $U_t\ni t$ в $[0,1]$ со свойством: при $\xi\in U_t$
  точка $\gamma(\xi)\in B(\gamma(t),R_t)$ и $F_\xi$ --- НАП $F_t$. Результат продолжения единственен; продолжения вдоль
  гомотопных (с закреплёнными концами) путей дают один и тот же результат (см.\ билеты 25--27).
\item \textbf{П.А.Ф.} (Опр.~15.5): совокупность $\Phi$ всех канонических элементов, получаемых из данного элемента продолжением
  вдоль всех путей, вдоль которых это возможно; $D_\Phi$ --- объединение кругов элементов $\Phi$ (область определения).
  П.А.Ф.\ определяется любым своим элементом.
\item \textbf{Аналитическая ветвь} (Опр.~15.6): если $D_1\subset D_\Phi$ --- область, $F_b\in\Phi$, $b\in D_1$, и $F_b$ продолжается вдоль
  любого пути в $D_1$ с началом $b$, то совокупность $\Phi_1$ всех таких продолжений --- \emph{аналитическая ветвь} $\Phi$ над $D_1$.
  Число элементов $\Phi_1$ с центром в $c\in D_1$ не зависит от $c$; если оно равно $N<\infty$, ветвь называется
  \emph{$N$-значной} (конечнозначной). \textbf{Теорема о монодромии}: если $D_1$ односвязна, то найдётся $g\in H(D_1)$ с $\Phi_1=\{E_c(g):c\in D_1\}$.
\item \textbf{Изолированные особые точки ветви} (Опр.~15.9). Пусть $\Psi$ --- ветвь $\Phi$ над $V(a,\delta)$. Если элемент после
  одного обхода $a$ переходит в себя --- $a$ точка \emph{однозначного характера} (устранимая, полюс или существенно особая точка
  голоморфной функции в $V(a,\delta)$); если впервые возвращается в себя после $n\ge2$ обходов --- \emph{точка ветвления порядка $n$};
  если не возвращается никогда --- \emph{логарифмическая} точка ветвления.
  \emph{Пример:} для $\sqrt z$ точка $0$ --- точка ветвления порядка $2$, для $\Ln z$ --- логарифмическая.
\item \textbf{Многочлен от $w$ над $\CC[z]$.} $P(z,w)=\sum_{k=0}^N P_k(z)w^k$, $P_k\in\CC[z]$, $P_N\not\equiv0$; $N=\deg_wP$.
  $\CC(z)$ --- поле рациональных функций. $P$ \emph{неприводим над $\CC(z)$}, если его нельзя записать как произведение двух многочленов
  от $w$ положительной степени с коэффициентами из $\CC(z)$.
\item \textbf{Алгебраический элемент} (слайд 400): канонический элемент $F=(B,f)$ с $P(z,f(z))\equiv0$ на $B$; $P$ --- \emph{представляющий}
  многочлен. Представляющий многочлен минимальной степени по $w$ --- \emph{определяющий}.
  \emph{Пример:} любой элемент $\sqrt z$ алгебраический, $P=w^2-z$ --- представляющий и определяющий.
\item \textbf{Сингулярная точка многочлена $P$}: точка $z\in\CC$, для которой $P_N(z)=0$, либо существует $w$ с
  $P(z,w)=0$, $P'_w(z,w)=0$ (т.е.\ у $P(z,\cdot)$ есть кратный корень). $\mathcal S=\mathcal S_P$ --- множество сингулярных точек.
  \emph{Примеры:} для $P=w^2$ все $z\in\CC$ сингулярны; для (приводимого) $P=w^2-z^2$ сингулярна лишь $z=0$; для $w^2-z$ --- тоже лишь $z=0$.
\item \textbf{Алгебраическая особая точка ветви} (Опр.~15.10), \textbf{алгебраическая П.А.Ф.} (Опр.~15.11) --- см.\ п.~\ref{t30:defalg}.
\item \textbf{Группа монодромии (Галуа)} --- см.\ п.~\ref{t30:galois}.
\end{itemize}
\end{defs}

\subsection{Алгебраические элементы. Определяющий многочлен}

\begin{thmbox}{Предложение 30.1 (свойства определяющего многочлена; слайд 400)}
Пусть $F=(B,f)$ --- алгебраический элемент, $Q$ --- его определяющий многочлен, $n=\deg_wQ$. Тогда
\begin{enumerate}
\item[(а)] $n\ge1$, и любой представляющий многочлен $P$ элемента $F$ делится на $Q$ в кольце $\CC(z)[w]$;
\item[(б)] определяющий многочлен единствен с точностью до ненулевого множителя из $\CC(z)$: если $\widetilde Q$ --- тоже
  определяющий, то $v(z)\widetilde Q=u(z)Q$ для некоторых ненулевых $u,v\in\CC[z]$;
\item[(в)] $Q$ неприводим над $\CC(z)$.
\end{enumerate}
\end{thmbox}
\begin{proof}
Будем пользоваться двумя фактами.

\emph{(i) Если $R\in\CC(z)[w]$ и $R(z,f(z))\equiv0$ на $B$ вне конечного числа точек (полюсов коэффициентов), то после умножения на
общий знаменатель $d\in\CC[z]$ коэффициентов $R$ получаем $dR\in\CC[z][w]$ --- представляющий многочлен для $F$.} Действительно,
$d(z)R(z,f(z))$ --- голоморфная в $B$ функция, равная нулю вне конечного множества, значит, $\equiv0$ по непрерывности.

\emph{(ii) Многочлен $R\in\CC[z][w]$ степени $0$ по $w$, т.е.\ $R=r(z)$, представляет $F$ лишь при $r\equiv0$}: ведь $r(z)\equiv0$ на круге
означает $r=0$.

(а) Из (ii) $n\ge1$. Разделим $P$ на $Q$ с остатком в кольце $\CC(z)[w]$ (это возможно, так как $\CC(z)$ --- поле и старший
коэффициент $Q_n\not\equiv0$): $P=TQ+R$, $\deg_wR<n$. Подставляя $w=f(z)$, получаем $R(z,f(z))=0$ на $B$ вне полюсов коэффициентов $T$, $R$.
По (i) $dR$ --- представляющий многочлен степени $<n$; по минимальности $n$ он нулевой, т.е.\ $R=0$, и $Q\mid P$.

(б) Если $\widetilde Q$ --- другой определяющий, то $\deg_w\widetilde Q=n$, и по (а) $\widetilde Q=TQ$ с $\deg_wT=0$, т.е.\
$T=u/v\in\CC(z)$. (Подробнее: многочлен $\widetilde Q_nQ-Q_n\widetilde Q$ представляющий и имеет степень $<n$, значит, равен нулю.)

(в) Пусть $Q=AB$, $A,B\in\CC(z)[w]$, $\deg_wA,\deg_wB\ge1$. Умножив на общий знаменатель, получим $dQ=A_1B_1$ с
$A_1,B_1\in\CC[z][w]$, $d\in\CC[z]\setminus\{0\}$, $\deg_wA_1=\deg_wA<n$, $\deg_wB_1=\deg_wB<n$. Тогда
$A_1(z,f(z))\cdot B_1(z,f(z))=d(z)Q(z,f(z))\equiv0$ в круге $B$. Произведение двух голоморфных в круге функций тождественно равно
нулю, только если одна из них $\equiv0$ (иначе нули обеих изолированы, и вне счётного множества произведение ненулевое).
Значит, $A_1$ или $B_1$ --- ненулевой представляющий многочлен степени $<n$ --- противоречие.
\end{proof}

\begin{remarkbox}{Замечание (неточность слайда 400)}
На слайде сказано «единственный с точностью до ненулевого множителя из $\CC[z]$». Точнее --- из $\CC(z)$: например, $(z+1)(w^2-z)$ и
$z(w^2-z)$ оба определяющие для $\sqrt z$, но ни один не равен другому, умноженному на многочлен. Если требовать, чтобы
коэффициенты $Q$ не имели общего множителя в $\CC[z]$ (примитивность), то $Q$ определён однозначно с точностью до постоянного множителя.
\end{remarkbox}

\begin{thmbox}{Лемма 30.2 (Гаусс)}
Если $P\in\CC[z,w]$ неприводим в кольце $\CC[z,w]$ и $\deg_wP\ge1$, то $P$ неприводим над $\CC(z)$ (как многочлен от $w$).
\end{thmbox}
\begin{proof}
Пусть $P=AB$, $A,B\in\CC(z)[w]$, $\deg_wA,\deg_wB\ge1$. Умножая на общий знаменатель, получаем $dP=A_1B_1$, $A_1,B_1\in\CC[z][w]$,
$d\in\CC[z]\setminus\{0\}$. Выберем такое представление с минимальной степенью $d$. Если $\deg d\ge1$, пусть $c$ --- корень $d$.
Подставим $z=c$: $A_1(c,w)B_1(c,w)=d(c)P(c,w)\equiv0$ в $\CC[w]$; так как $\CC[w]$ --- область целостности, то, скажем,
$A_1(c,w)\equiv0$, т.е.\ все коэффициенты $A_1$ делятся на $z-c$. Тогда $\frac{d}{z-c}P=\frac{A_1}{z-c}B_1$ --- представление с
меньшей степенью $d$, противоречие. Значит, $d$ --- константа, и $P=(A_1/d)B_1$ --- разложение в $\CC[z,w]$ на множители
положительной степени по $w$, что противоречит неприводимости.
\end{proof}

\subsection{Сингулярные точки и конечность сингулярного множества}

\begin{thmbox}{Предложение 30.3 (слайд 401)}
Если $P$ и $P'_w$ не имеют общего множителя положительной степени по $w$ в $\CC(z)[w]$ (в частности, если $P$ неприводим над
$\CC(z)$), то множество $\mathcal S_P$ конечно. Для $z_0\notin\mathcal S_P$ уравнение $P(z_0,w)=0$ имеет ровно $N$ различных корней.
\end{thmbox}
\begin{proof}
\emph{Неприводимый случай сводится к общему.} Если $P$ неприводим над $\CC(z)$ и $N\ge1$, то наибольший общий делитель $P$ и
$P'_w$ в $\CC(z)[w]$ --- делитель $P$, т.е.\ (с точностью до множителя из $\CC(z)$) либо $1$, либо $P$. Но $P'_w=\sum_{k\ge1}kP_kw^{k-1}$
--- ненулевой многочлен (его коэффициент при $w^{N-1}$ равен $NP_N\ne0$) степени $N-1<N$, на $P$ он делиться не может. Значит,
НОД $=1$.

\emph{Общий случай.} Кольцо $\CC(z)[w]$ евклидово, поэтому из взаимной простоты $P$ и $P'_w$ следует соотношение Безу:
найдутся $U,V\in\CC(z)[w]$ с $UP+VP'_w=1$. Умножив на общий знаменатель $h\in\CC[z]\setminus\{0\}$ коэффициентов $U,V$, получим
\[
U_1(z,w)P(z,w)+V_1(z,w)P'_w(z,w)=h(z),\qquad U_1,V_1\in\CC[z][w].
\]
Если $z_0$ сингулярна и $P_N(z_0)\ne0$, то существует $w_0$ с $P(z_0,w_0)=P'_w(z_0,w_0)=0$, и подстановка даёт $h(z_0)=0$.
Следовательно, $\mathcal S_P\subset\{P_N=0\}\cup\{h=0\}$ --- конечное множество (нули двух ненулевых многочленов).

Если $z_0\notin\mathcal S_P$, то $P_N(z_0)\ne0$, т.е.\ $P(z_0,\cdot)$ --- многочлен степени ровно $N$, и у него нет кратных корней
(кратный корень $w_0$ удовлетворяет $P(z_0,w_0)=P'_w(z_0,w_0)=0$). Значит, корней ровно $N$, и они различны.
\end{proof}

Всюду далее $\mathcal S=\mathcal S_P$ конечно.

\subsection{Алгебраические элементы над несингулярными точками}

\begin{thmbox}{Предложение 30.4 (неявная функция; слайды 401--402)}
Пусть $z_0\notin\mathcal S$, $w_1,\dots,w_N$ --- корни $P(z_0,\cdot)$. Тогда существуют круг $B_0=B(z_0,\rho)\subset\CC\setminus\mathcal S$,
попарно непересекающиеся круги $B_k\ni w_k$ и функции $f_k\in H(B_0)$, $f_k(B_0)\subset B_k$, такие, что:
\begin{enumerate}
\item[(а)] при каждом $z\in B_0$ корни $P(z,\cdot)$ --- это ровно $f_1(z),\dots,f_N(z)$; $\;f_k'=-P'_z(z,f_k)/P'_w(z,f_k)$;
\item[(б)] если $G=(B_G,g)$ --- алгебраический элемент с представляющим многочленом $P$ и центром $c\in B_0$, то $g=f_k$
  в окрестности $c$ для некоторого $k$, т.е.\ $G=E_c(f_k)$, и радиус $G$ не меньше $\rho-|c-z_0|$.
\end{enumerate}
В частности, с центром $z_0$ имеется ровно $N$ алгебраических элементов $E_{z_0}(f_k)$ с представляющим многочленом $P$.
\end{thmbox}
\begin{proof}
Отождествим $\CC^2$ с $\RR^4$ и рассмотрим гладкое отображение $(z,w)\mapsto\zeta=P(z,w)$ из $\RR^4$ в $\RR^2$. Так как $P$
голоморфен по $w$, его вещественная матрица Якоби по переменным $\Re w$, $\Im w$ ---
это матрица умножения на $P'_w$, её
определитель равен $|P'_w(z,w)|^2$. В точке $(z_0,w_k)$ он положителен: $z_0\notin\mathcal S$ и $P(z_0,w_k)=0$, поэтому
$P'_w(z_0,w_k)\ne0$. Выберем попарно непересекающиеся круги $B_k\ni w_k$ настолько малые, чтобы по вещественной теореме
о неявной функции нашлись круги $B_0^{(k)}\ni z_0$ и функции $f_k\in C^1(B_0^{(k)})$ с $f_k(z_0)=w_k$,
$f_k(B^{(k)}_0)\subset B_k$, такие, что
\[
\text{для } (z,w)\in B_0^{(k)}\times B_k:\qquad P(z,w)=0\iff w=f_k(z).
\]
Пусть $B_0=B(z_0,\rho)$ --- круг, содержащийся во всех $B_0^{(k)}$ и в открытом множестве $\CC\setminus\mathcal S$ ($\mathcal S$ конечно).

\emph{Голоморфность.} Положим $\Phi_k(z)=P(z,f_k(z))\equiv0$. По цепному правилу для производных Виртингера (учитывая, что $P$
голоморфен по обеим переменным, т.е.\ $\partial P/\partial\bar z=\partial P/\partial\bar w=0$)
\[
0=\partial_{\bar z}\Phi_k=P'_w(z,f_k(z))\,\partial_{\bar z}f_k(z).
\]
При $z\in B_0$ имеем $P(z,f_k(z))=0$, а $z\notin\mathcal S$, значит $P'_w(z,f_k(z))\ne0$; поэтому $\partial_{\bar z}f_k\equiv0$, т.е.\ $f_k\in H(B_0)$.
Теперь $\partial_z f_k=f_k'$ и
$0=\partial_z\Phi_k=P'_z(z,f_k)+P'_w(z,f_k)f_k'$, откуда $f_k'=-P'_z(z,f_k)/P'_w(z,f_k)$.

(а) Числа $f_1(z),\dots,f_N(z)$ лежат в попарно непересекающихся кругах $B_k$, значит, различны, и все являются корнями
$P(z,\cdot)$; так как $P_N(z)\ne0$, корней ровно $N$, и других нет.

(б) Число $g(c)$ --- корень $P(c,\cdot)$, значит $g(c)=f_k(c)\in B_k$ при некотором $k$. По непрерывности $g(z)\in B_k$ для $z$ из
некоторого круга $B'\subset B_0\cap B_G$ с центром $c$; так как $P(z,g(z))=0$, из единственности в теореме о неявной функции
$g=f_k$ на $B'$. Поэтому $G=E_c(f_k)$, а так как $f_k\in H(B(c,\rho-|c-z_0|))$, радиус $G$ не меньше $\rho-|c-z_0|$.
\end{proof}

\begin{remarkbox}{Замечание (опечатки слайда 402)}
На слайде в формуле $0=F'_z(b)=P'_w\,\partial_zf_k$ должна стоять производная $\partial_{\bar z}$ (именно из неё выводится
голоморфность); в формуле для $f_k'$ вместо $\varphi(b)$ следует читать $f_k(b)$; «$B_0\times B_n$» --- это $B_0\times B_k$.
\end{remarkbox}

\begin{thmbox}{Предложение 30.5 (продолжаемость и радиусы; слайд 402)}
Пусть $F=(B(a,r_a),f)$ --- алгебраический элемент с представляющим многочленом $P$ и $a\notin\mathcal S$. Тогда $F$ продолжается вдоль
любого пути в $\CC\setminus\mathcal S$, все элементы продолжения алгебраические с тем же $P$, и $r_a\ge\dist(a,\mathcal S)$.
\end{thmbox}
\begin{proof}
\emph{Сохранение уравнения.} Если $F_1=(B_1,f_1)$ --- НАП алгебраического $F_0=(B_0,f_0)$, то $P(z,f_1(z))$ голоморфна в $B_1$ и
равна $P(z,f_0(z))=0$ на непустом открытом $B_0\cap B_1$, значит, $\equiv0$ в $B_1$. Поэтому при любом продолжении (вдоль пути, по
цепочке) уравнение $P(z,f(z))\equiv0$ сохраняется.

\emph{Продолжаемость.} Пусть $\gamma:[0,1]\to\CC\setminus\mathcal S$, $\gamma(0)=a$. Для каждой точки $z\in\gamma^*$ возьмём круг
$B_0(z)$ из предложения 30.4. Прообразы $\gamma^{-1}(B_0(z))$ образуют открытое покрытие компакта $[0,1]$; пусть $\lambda$ --- его
число Лебега и $0=t_0<t_1<\dots<t_n=1$ --- разбиение с шагом $<\lambda$; тогда $\gamma([t_{j-1},t_j])\subset\Delta_j:=B_0(z^{(j)})$.
Строим продолжение по индукции. На $[t_0,t_1]$: по 30.4(б) $F=E_a(f^{(1)}_{k_1})$ для одной из функций $f^{(1)}_{k}\in H(\Delta_1)$;
положим $F_t=E_{\gamma(t)}(f^{(1)}_{k_1})$. Если $F_t$ уже построено на $[0,t_{j-1}]$ и алгебраично, то $\gamma(t_{j-1})\in\Delta_j$, и по 30.4(б)
$F_{t_{j-1}}=E_{\gamma(t_{j-1})}(f^{(j)}_{k_j})$ для некоторой $f^{(j)}_{k_j}\in H(\Delta_j)$; положим $F_t=E_{\gamma(t)}(f^{(j)}_{k_j})$ при $t\in[t_{j-1},t_j]$
(в точке $t_{j-1}$ оба определения совпадают).
Условие Опр.~15.4 проверяется так: если $t,\xi\in[t_{j-1},t_j]$ и $|\gamma(\xi)-\gamma(t)|<\dist(\gamma(t),\partial\Delta_j)$, то $\gamma(\xi)$ лежит в
круге $F_t$, а функции $F_t$ и $F_\xi$ совпадают с $f^{(j)}_{k_j}$ около $\gamma(\xi)$, т.е.\ $F_\xi$ --- НАП $F_t$; по непрерывности $\gamma$ это
выполнено для $\xi$ из связной окрестности $t$ (в точке $t=t_{j-1}$ используем $\Delta_{j-1}$ слева и $\Delta_j$ справа).

\emph{Радиус.} Пусть $d=\dist(a,\mathcal S)$ ($d=\infty$, если $\mathcal S=\varnothing$). Круг $B(a,d)\subset\CC\setminus\mathcal S$ односвязен, и $F$
продолжается вдоль любого пути в нём. По теореме о монодромии найдётся $g\in H(B(a,d))$ с $E_a(g)=F$; значит, ряд Тейлора $f$ в
точке $a$ сходится в $B(a,d)$, и $r_a\ge d$.
\end{proof}

\subsection{Лемма 15.2: оценки вблизи сингулярной точки и бесконечности}

Ключ к лемме --- элементарная оценка корней многочлена.

\begin{thmbox}{Лемма 30.6 (оценка Коши для корней)}
Если $a_N\ne0$ и $a_Nw^N+a_{N-1}w^{N-1}+\dots+a_0=0$, то $|w|\le\max\Bigl(1,\ \sum_{k=0}^{N-1}|a_k|/|a_N|\Bigr)$.
\end{thmbox}
\begin{proof}
Если $|w|>1$, то $|a_N||w|^N\le\sum_{k<N}|a_k||w|^k\le|w|^{N-1}\sum_{k<N}|a_k|$, откуда $|w|\le\sum_{k<N}|a_k|/|a_N|$.
\end{proof}

\begin{thmbox}{Лемма 15.2}
Пусть $z_1\in\mathcal S$ и $z_1$ --- нуль $P_N$ порядка $m\in\ZZ_+$ ($m=0$, если $P_N(z_1)\ne0$). Тогда найдутся $A>0$ и $r>0$ такие,
что для любого алгебраического элемента $(B(a,r_a),f_a)$ с представляющим многочленом $P$ и условием $0<|z_1-a|\le r$ выполнено
$|z_1-a|\le r_a$ и при всех $z\in B(a,|z_1-a|)$
\[
|f_a(z)|\le\frac{A}{|z-z_1|^m}.
\]
Кроме того, найдутся $M>0$, $s\ge0$, $R>0$ такие, что для любого алгебраического элемента $(B(a,r_a),f_a)$ с представляющим
многочленом $P$ при всех $z\in B(a,r_a)\setminus B(0,R)$ справедлива оценка $|f_a(z)|\le M|z|^s$.
\end{thmbox}
\begin{proof}
Без ограничения общности $z_1=0$.

\emph{Шаг 1 (замена $g=z^mf$).} Если $(B,f)$ --- элемент с представляющим $P$, то $g(z)=z^mf(z)$ удовлетворяет $P_m(z,g(z))\equiv0$, где
\[
P_m(z,w)=\frac{P_N(z)}{z^m}w^N+\sum_{k=0}^{N-1}z^{m(N-1-k)}P_k(z)w^k=:\sum_{k=0}^NQ_k(z)w^k .
\]
Проверка: $P_m(z,z^mf)=\frac{P_N}{z^m}z^{mN}f^N+\sum_{k<N}z^{m(N-1-k)}P_kz^{mk}f^k=z^{m(N-1)}P(z,f)=0$. Все $Q_k$ --- многочлены, и $Q_N(0)\ne0$.

\emph{Шаг 2 (выбор констант).} Выберем $\rho>0$ так, что $\rho\le\dist(0,\mathcal S\setminus\{0\})$ (любое $\rho$, если $\mathcal S=\{0\}$),
$|Q_N(z)|\ge C_1>0$ и $|Q_k(z)|\le C_2$ ($C_2>0$) при $z\in B(0,\rho)$, $k<N$ (возможно по непрерывности, так как $Q_N(0)\ne0$).
Положим $r=\rho/2$.

\emph{Шаг 3 (радиус).} Пусть $0<|a|\le r$. Тогда $a\notin\mathcal S$ и $\dist(a,\mathcal S\setminus\{0\})\ge\rho-|a|\ge|a|$, т.е.\ $\dist(a,\mathcal S)=|a|$.
По предложению 30.5 $r_a\ge|a|$.

\emph{Шаг 4 (оценка; «оставшиеся рассуждения» слайда).} Если $z\in B(a,|a|)$, то $|z|<2|a|\le\rho$ и $z\ne0$. Число $g_a(z)=z^mf_a(z)$ ---
корень многочлена $\sum_kQ_k(z)w^k$ с $|Q_N(z)|\ge C_1$. Слайд оценивает так: из $Q_Ng_a^N=-\sum_{k<N}Q_kg_a^k$
\[
|g_a|^N\le\frac{C_2}{C_1}\sum_{k=0}^{N-1}|g_a|^k=\frac{C_2}{C_1}\cdot\frac{|g_a|^N-1}{|g_a|-1}\quad(|g_a|\ne1),
\]
и при $|g_a|\ge2$ имеем $\frac{|g_a|^N-1}{|g_a|-1}\le\frac{|g_a|^N}{|g_a|/2}=2|g_a|^{N-1}$, откуда $|g_a|\le2C_2/C_1$.
Итак, \emph{в каждой точке} $z\in B(a,|a|)$ либо $|g_a(z)|<2$, либо $|g_a(z)|\le2C_2/C_1$; т.е.\ всегда
$|g_a(z)|\le A:=\max(2,2C_2/C_1)$. Значит, $|f_a(z)|=|g_a(z)|/|z|^m\le A/|z|^m$. (То же даёт лемма 30.6: $|g_a|\le\max(1,NC_2/C_1)$.)
Заметим, что $A,r$ не зависят от элемента: оценка --- это оценка \emph{всех корней} $P(z,\cdot)$ при $0<|z|<\rho$.

\emph{Шаг 5 (второе утверждение).} Пусть $d=\deg P_N$, $p\ne0$ --- старший коэффициент $P_N$, $D=\max_{k<N}\deg P_k$ (если все $P_k\equiv0$
при $k<N$, то $f_a\equiv0$ и доказывать нечего). Так как $|P_N(z)|/|z|^d\to|p|$ при $z\to\infty$, найдётся $R\ge1$ такое, что
$|P_N(z)|\ge\frac{|p|}{2}|z|^d$ при $|z|\ge R$. Для многочлена $P_k=\sum_j c_{kj}z^j$ степени $\le D$ при $|z|\ge1$ имеем
$|P_k(z)|\le C|z|^D$, где $C=\max_k\sum_j|c_{kj}|$. Пусть $z\in B(a,r_a)$, $|z|\ge R$. Тогда $f_a(z)$ --- корень $P(z,\cdot)$, и по лемме 30.6
\[
|f_a(z)|\le\max\Bigl(1,\frac{NC|z|^D}{\frac{|p|}{2}|z|^d}\Bigr)\le M|z|^s,\qquad s=\max(D-d,0),\quad M=\max\Bigl(1,\frac{2NC}{|p|}\Bigr),
\]
так как $|z|\ge1$ и $|z|^{D-d}\le|z|^s$.
\end{proof}

\begin{remarkbox}{Замечание}
Лемма 15.2 по существу --- утверждение о \emph{корнях} уравнения $P(z,w)=0$: при $0<|z-z_1|<\rho$ все корни по модулю $\le A|z-z_1|^{-m}$,
при $|z|\ge R$ все корни $\le M|z|^s$. На слайде в формулировке оценка $|f_a(z)|\le A/|z-z_1|^m$; именно это и доказано.
\end{remarkbox}

\subsection{Теорема 15.3: многочлен порождает П.А.Ф.}

Нам понадобится

\begin{thmbox}{Лемма 30.7 (рациональность)}
Пусть $T\subset\CC$ конечно, $s\in H(\CC\setminus T)$, и каждая точка $T\cup\{\infty\}$ --- устранимая особая точка или полюс $s$.
Тогда $s$ --- рациональная функция. Достаточное условие: $|s(z)|\le C|z-t|^{-k}$ около каждой $t\in T$ и $|s(z)|\le C|z|^k$ при больших $|z|$.
\end{thmbox}
\begin{proof}
Пусть $g_t(z)$ --- главная часть ряда Лорана $s$ в точке $t\in T$ (многочлен от $1/(z-t)$). Тогда $h=s-\sum_{t\in T}g_t$ имеет в точках $T$
устранимые особенности, т.е.\ продолжается до целой функции; в $\infty$ функция $h$ имеет не более чем полюс (так как $g_t(z)\to0$),
т.е.\ $|h(z)|\le C'|z|^k$ при больших $|z|$. По обобщённой теореме Лиувилля (неравенства Коши: $|c_n|\le\max_{|z|=\varrho}|h|/\varrho^n\to0$ при $n>k$)
$h$ --- многочлен, и $s=h+\sum g_t$ рациональна. Достаточное условие: если $|s|\le C|z-t|^{-k}$, то $(z-t)^ks$ ограничена около $t$, т.е.\ $t$ ---
устранимая точка для $(z-t)^ks$ (теорема Римана), и $s$ имеет в $t$ не более чем полюс; аналогично в $\infty$.
\end{proof}

\begin{thmbox}{Теорема 15.3}
Пусть $P\in\CC[z,w]$ неприводим, $\deg_wP=N\ge1$, $\mathcal S$ --- его сингулярное множество, $D=\CC\setminus\mathcal S$.
Совокупность всех алгебраических элементов с представляющим многочленом $P$ образует П.А.Ф.\ $\Phi$, для которой
$\Psi=\Phi|_D$ (элементы с центрами в $D$) является $N$-значной аналитической ветвью над $D$. При этом $P$ --- определяющий для всех
этих элементов, а все изолированные особые точки всех ветвей $\Phi$ (в $\Ch$) являются либо точками ветвления конечного порядка
($\le N$), либо точками однозначного характера, не являющимися существенно особыми (т.е.\ полюсами или устранимыми).
\end{thmbox}
\begin{proof}
По лемме 30.2 $P$ неприводим над $\CC(z)$, поэтому $\mathcal S$ конечно (предл.~30.3) и $D$ --- область (плоскость без конечного множества).
Обозначим через $\mathcal A$ множество всех (канонических) алгебраических элементов с представляющим $P$.

\emph{Шаг 1: ветвь над $D$.} Зафиксируем $a\in D$ и $F\in\mathcal A$ с центром $a$ (они есть по предл.~30.4). По предл.~30.5 $F$
продолжается вдоль любого пути в $D$; пусть $\Psi$ --- аналитическая ветвь над $D$, порождённая $F$; $\Psi\subset\mathcal A$. Пусть
$F_1,\dots,F_J$ --- все различные элементы $\Psi$ с центром $a$, $F_j=(B_j,f_j)$. По предл.~30.4 элементов $\mathcal A$ с центром $a$ ровно $N$,
поэтому $J\le N$ (на слайде здесь ссылка на лемму 15.2 --- описка: нужно предл.~30.4), и значения $f_j(a)$ попарно различны.

\emph{Шаг 2: перестановки.} Пусть $\gamma$ --- замкнутый путь в $D$ с началом $a$. Продолжение $F_j$ вдоль $\gamma$ --- снова элемент $\Psi$
с центром $a$, т.е.\ $F_{\sigma(j)}$. Отображение $\sigma=\sigma_\gamma$ инъективно: если $F_i$ и $F_j$ переходят в один элемент, то продолжая его
вдоль обратного пути $\gamma^{-}$ получаем (по единственности продолжения) и $F_i$, и $F_j$, т.е.\ $i=j$. Итак, $\sigma_\gamma$ --- перестановка.

\emph{Шаг 3: симметрические функции рациональны.} Пусть $\tilde S_1=\sum q_i,\ \dots,\ \tilde S_J=q_1\cdots q_J$ --- элементарные
симметрические многочлены и $S_n:=E_a(\tilde S_n(f_1,\dots,f_J))$. Если $F_j^t$ --- продолжения $F_j$ вдоль пути $\gamma$, то
$E_{\gamma(t)}(\tilde S_n(f^t_1,\dots,f^t_J))$ --- продолжение $S_n$ вдоль $\gamma$ (сумма и произведение элементов, совпадающих на пересечениях кругов,
снова совпадают). Для замкнутого $\gamma$ набор $\{F_j\}$ переходит в себя с перестановкой, а $\tilde S_n$ симметричен, поэтому $S_n$
переходит в себя. Значит, формула $s_n(z):=$ значение в $z$ результата продолжения $S_n$ вдоль любого пути из $a$ в $z$ в $D$
корректно определяет функцию на $D$ (два пути $\gamma_1,\gamma_2$ дают одно и то же, так как $\gamma_1\gamma_2^{-}$ замкнут), и $s_n\in H(D)$
(локально $s_n$ совпадает с функцией элемента). Значение $s_n(z)$ --- элементарная симметрическая функция от $J$ различных
корней $P(z,\cdot)$, поэтому по лемме 15.2 около каждой $t\in\mathcal S$: $|s_n(z)|\le\binom Jn\max(1,A)^n|z-t|^{-mn}$ (при $|z-t|\le1$),
а при $|z|\ge R$: $|s_n(z)|\le\binom Jn M^n|z|^{sn}$. По лемме 30.7 все $s_n$ рациональны.

\emph{Шаг 4: $J=N$.} Пусть $h\in\CC[z]$ --- общий знаменатель $s_1,\dots,s_J$ и
\[
Q(z,w)=h(z)\bigl(w^J-s_1(z)w^{J-1}+\dots+(-1)^Js_J(z)\bigr)\in\CC[z][w].
\]
По формулам Виета $Q(z,w)=h(z)\prod_{j}(w-f_j(z))$ около $a$, поэтому $Q$ представляет все $F_j$. Пусть $J<N$. Разделим $P$ на
$Q$ в $\CC(z)[w]$: $P=TQ+R$, $\deg_wR<J$. Если $R=0$, то $P=TQ$ с $\deg_wT=N-J\ge1$, $\deg_wQ=J\ge1$ --- $P$ приводим над $\CC(z)$,
противоречие. Если $R\ne0$, то $R(z,f_j(z))=0$ для всех $j=1,\dots,J$ около $a$. Возьмём $z$ около $a$, не являющуюся полюсом
коэффициентов $R$ и нулём какого-нибудь ненулевого коэффициента $R$ (таких точек --- все, кроме конечного числа). Тогда
$R(z,\cdot)$ --- ненулевой многочлен степени $<J$, имеющий $J$ различных корней $f_1(z),\dots,f_J(z)$ (различны при $z\notin\mathcal S$,
предл.~30.4) --- противоречие. Значит, $J=N$: \emph{все} $N$ элементов $\mathcal A$ с центром $a$ лежат в $\Psi$. Точка $a\in D$ была
произвольной, а ветвь, порождённая любым своим элементом, совпадает с $\Psi$ (продолжения обратимы); поэтому $\Psi=\mathcal A|_D$
--- $N$-значная ветвь над $D$.

\emph{Шаг 5: $\mathcal A$ --- П.А.Ф.} Пусть $\Phi$ --- П.А.Ф., порождённая $F$. Все её элементы --- продолжения $F$, значит, $\Phi\subset\mathcal A$
(уравнение сохраняется). Обратно, $\mathcal A|_D=\Psi\subset\Phi$. Если $G=(B(c,r),g)\in\mathcal A$ и $c\in\mathcal S$, возьмём $c'\in B(c,r)\setminus\mathcal S$;
$E_{c'}(g)\in\mathcal A|_D\subset\Phi$, а $G$ получается из $E_{c'}(g)$ продолжением вдоль отрезка $[c',c]\subset B(c,r)$ (семейство $E_{\gamma(t)}(g)$),
поэтому $G\in\Phi$. Итак, $\Phi=\mathcal A$.

\emph{Шаг 6: $P$ определяющий.} Пусть $Q'\ne0$ представляет некоторый $G\in\Phi$ и $\deg_wQ'<N$. Так как все элементы $\Phi$ --- продолжения $G$,
$Q'$ представляет все элементы $\Phi$, в частности все $N$ элементов с центром $a$; как в шаге 4 получаем противоречие.

\emph{Шаг 7: особые точки.} Пусть $t\in\mathcal S\cup\{\infty\}$, $\delta>0$ таково, что $V(t,\delta)\subset D$, и $\Psi_1$ --- ветвь над $V(t,\delta)$. Петля
$\gamma$, однократно обходящая $t$, действует перестановкой на $N$ элементах над точкой $b\in V(t,\delta)$; орбита элемента $F_b$ под действием
степеней $\sigma_\gamma$ имеет длину $k\le N$, т.е.\ $F_b$ впервые возвращается в себя после $k$ обходов. При $k\ge2$ --- точка ветвления порядка
$k\le N$ (логарифмическое ветвление невозможно). При $k=1$ ветвь однозначна: задаётся функцией $f_1\in H(V(t,\delta))$, и по лемме 15.2
$|f_1(z)|\le A|z-t|^{-m}$ (соотв.\ $\le M|z|^s$ при $t=\infty$), значит, $t$ --- устранимая точка или полюс (как в лемме 30.7).
\end{proof}

\subsection{Алгебраические особые точки и алгебраические П.А.Ф.}\label{t30:defalg}

Пусть $\Phi$ --- П.А.Ф., $z_1\in\Ch$, $D_1=V(z_1,\delta)\subset D_\Phi$ ($z_1$ может как принадлежать $D_\Phi$, так и нет), $b\in D_1$,
$F_b\in\Phi$, $\Psi$ --- ветвь $\Phi$ над $D_1$, порождённая $F_b$. Пусть $\Psi$ имеет в $z_1$ ветвление конечного порядка $m\ge2$ или
однозначный характер ($m=1$).

\begin{defs}[Определения 15.10, 15.11]
\begin{itemize}
\item \textbf{Алгебраическая особая точка.} Точка $z_1\in\CC$ --- \emph{алгебраическая} О.Т.\ для $\Psi$, если найдутся $A>0$, $s\ge0$,
  $r\in(0,\delta)$ такие, что для любого $c\in V(z_1,r)$ и всякого элемента $F_c=(B(c,r_c),f_c)\in\Psi$ при всех
  $z\in B(c,r_c)\cap V(z_1,r)$
  \[
  |f_c(z)|\le\frac{A}{|z-z_1|^s}.
  \]
  Точка $z_1=\infty$ --- алгебраическая, если (при тех же кванторах, с $V(\infty,r)=\{|z|>1/r\}$) $|f_c(z)|\le A|z|^s$.
  Иначе говоря, это точка ветвления конечного порядка или точка однозначного характера, около которой функция растёт не быстрее
  степени. \emph{Пример:} $0$ для $\sqrt z$, $1/\sqrt z$, $z^{-3/2}$; не алгебраические: $0$ для $e^{1/\sqrt z}$ (ветвление порядка $2$, но рост
  экспоненциальный), $0$ для $e^{1/z}$ (существенно особая).
\item \textbf{Алгебраическая П.А.Ф.} $\Phi$ называется алгебраической, если
  (1) найдётся конечное $\mathcal S_\Phi\subset\Ch$, $\infty\in\mathcal S_\Phi$, такое, что $D=\Ch\setminus\mathcal S_\Phi\subset D_\Phi$ и совокупность
  элементов $\Phi|_D$ есть конечнозначная аналитическая ветвь $\Phi$ над $D$ (берётся минимальное такое $\mathcal S_\Phi$);
  (2) для каждой $a\in\mathcal S_\Phi$ найдётся $\delta>0$ такое, что для любой аналитической ветви $\Phi_1$ П.А.Ф.\ $\Phi$ над $V(a,\delta)$
  точка $a$ --- алгебраическая для $\Phi_1$.
\end{itemize}
\end{defs}

\begin{remarkbox}{Связь с теоремой 15.3}
П.А.Ф.\ из теоремы 15.3 алгебраическая: условие (1) выполнено с $\mathcal S_\Phi\subset\mathcal S\cup\{\infty\}$ (над $\CC\setminus\mathcal S$ имеем $N$-значную ветвь),
а (2) --- это лемма 15.2 (при $\delta<\dist(a,\mathcal S\setminus\{a\})$ все значения элементов ветвей --- корни $P(z,\cdot)$, для которых верны оценки
леммы, со степенью $s=m$ в конечных точках). Теорема 15.4 утверждает обратное.
\end{remarkbox}

\begin{thmbox}{Теорема 15.4}
Пусть $\Phi$ --- алгебраическая П.А.Ф., $N$-значная над $D=\Ch\setminus\mathcal S_\Phi$, $N\in\NN$. Тогда найдутся $P_k\in\CC[z]$, $k=0,\dots,N$,
такие, что многочлен $P(z,w)=\sum_{n=0}^NP_n(z)w^n$ степени $N$ по $w$ неприводим над $\CC(z)$, и для каждого элемента $F=(B,f)\in\Phi$
$P(z,f(z))\equiv0$ на $B$, т.е.\ каждый элемент $\Phi$ алгебраический с представляющим (и определяющим) многочленом $P$.
\end{thmbox}
\begin{proof}
Заметим, что $D\subset\CC$ ($\infty\in\mathcal S_\Phi$) --- область. Зафиксируем $b\in D$ и все (различные) элементы $F_k=(B_k,f_k)$, $k=1,\dots,N$,
из $\Phi$ с центром $b$. Как в шаге 2 доказательства теоремы 15.3, продолжение вдоль замкнутого пути $\gamma\subset D$ с началом $b$ действует
перестановкой $\sigma_\gamma$: $F_k\overset{\gamma}{\leadsto}F_{\sigma(k)}$. Поэтому элементарные симметрические элементы
$S_n=E_b(\tilde S_n(f_1,\dots,f_N))$ переходят в себя и, как в шаге 3, задают функции $s_n\in H(D)$.

\begin{detail}{почему $s_n$ рациональны («сразу следует из определений 15.10, 15.11»)}
Пусть $a\in\mathcal S_\Phi$. Уменьшив $\delta$ из условия (2), можно считать $V(a,\delta)\cap\mathcal S_\Phi=\varnothing$, т.е.\ $V(a,\delta)\subset D$
(при уменьшении $\delta$ элементы ветви над меньшим проколотым кругом являются элементами ветви над большим, и оценки сохраняются).
Над $V(a,\delta)$ элементы $\Phi$ (их по $N$ над каждой точкой) распадаются на аналитические ветви; две ветви либо совпадают, либо не
имеют общих элементов (ветвь порождается любым своим элементом), и каждая ветвь имеет элемент над фиксированной точкой $c_0\in V(a,\delta)$,
поэтому ветвей не более $N$. Для каждой ветви $\Phi_i$ по определению 15.10 есть свои $A_i,s_i,r_i$; положим $A=\max(1,\max A_i)$,
$s=\max s_i$, $r=\min(r_i,1)$. Для $z\in V(a,r)$ каждое из $N$ значений в точке $z$ элементов $\Phi$ с центром $z$ есть $f_c(z)$ для элемента
$F_c$ ($c=z$) некоторой ветви $\Phi_i$, поэтому оно по модулю $\le A|z-a|^{-s}$. Значит,
$|s_n(z)|\le\binom Nn A^n|z-a|^{-ns}$, и $a$ --- не более чем полюс $s_n$. Аналогично в $\infty$: $|s_n(z)|\le\binom NnA^n|z|^{ns}$.
По лемме 30.7 (с $T=\mathcal S_\Phi\setminus\{\infty\}$) $s_n$ рациональна. Она не зависит от $b$: при другом $b'$ получаются продолжения
тех же элементов.
\end{detail}

Пусть $h\in\CC[z]$ --- общий знаменатель $s_1,\dots,s_N$, $P(z,w):=h(z)\bigl(w^N-s_1w^{N-1}+\dots+(-1)^Ns_N\bigr)$. По Виета при
$z\in\bigcap_kB_k$ числа $f_k(z)$ --- корни $P(z,\cdot)$, т.е.\ $P$ представляет $F_1$. Каждый элемент $\Phi$ --- продолжение $F_1$, а
уравнение $P(z,f(z))\equiv0$ сохраняется при продолжении (см.\ предл.~30.5), поэтому $P$ представляет все элементы $\Phi$.

\emph{Неприводимость.} Пусть $P=AB$ в $\CC(z)[w]$, $\deg_wA,\deg_wB\ge1$. Тогда (как в предл.~30.1(в)) около $b$ одна из функций
$A(z,f_1(z))$, $B(z,f_1(z))$ (мероморфных) тождественно равна нулю, скажем первая; после умножения на знаменатель многочлен
$A_1\in\CC[z][w]$, $\deg_wA_1<N$, представляет $F_1$, а значит (продолжение) и все $F_1,\dots,F_N$. Функции $f_k$ попарно различны около $b$,
поэтому вне счётного (дискретного) множества точек значения $f_1(z),\dots,f_N(z)$ различны, и $A_1(z,\cdot)$ (ненулевой при всех $z$, кроме
конечного числа) имеет $N>\deg_wA_1$ корней --- противоречие. Тем же рассуждением многочлен степени $<N$ не может представлять элемент $\Phi$,
т.е.\ $P$ --- определяющий.
\end{proof}

\begin{remarkbox}{Итог}
Теоремы 15.3 и 15.4 вместе: \emph{П.А.Ф.\ алгебраическая (в смысле Опр.~15.11) тогда и только тогда, когда её элементы удовлетворяют
неприводимому алгебраическому уравнению $P(z,w)=0$}. Уравнение $P(z,w)=0$ называется алгебраическим уравнением степени $N$ относительно $w$.
\end{remarkbox}

\subsection{Группа монодромии (группа Галуа)}\label{t30:galois}

В обозначениях теоремы 15.4 множество $G_b=\{\sigma_\gamma\}$ всех перестановок по всем замкнутым путям $\gamma\subset D$ с началом $b$ ---
подгруппа симметрической группы $S_N$: тождественная перестановка даёт постоянный путь; $\sigma_{\gamma^-}=\sigma_\gamma^{-1}$; если $\gamma_1\gamma_2$ ---
последовательное прохождение, то $\sigma_{\gamma_1\gamma_2}=\sigma_{\gamma_2}\circ\sigma_{\gamma_1}$. По теореме о продолжении вдоль гомотопных путей
$\sigma_\gamma$ зависит лишь от гомотопического класса $[\gamma]\in\pi_1(D,b)$, поэтому $[\gamma]\mapsto\sigma_\gamma$ --- (анти)гомоморфизм
$\pi_1(D,b)$ \emph{на} $G_b$ (гомоморфизм, если перестановки компоновать «слева направо»). Если $b'\in D$ и $\delta$ --- путь из $b$ в $b'$, то
продолжение вдоль $\delta$ задаёт биекцию элементов над $b$ и над $b'$, и $\sigma_{\gamma'}$ для $\gamma'=\delta^-\gamma\delta$ соответствует $\sigma_\gamma$;
поэтому $G_{b'}\cong G_b$ (сопряжены). Эта группа $G$ --- \emph{группа Галуа} (\emph{монодромии}) П.А.Ф.\ $\Phi$.
Неприводимость $P$ $\iff$ транзитивность $G$ (все $N$ элементов над $b$ получаются друг из друга продолжением), см.\ шаг 4 теоремы 15.3.

\subsection{\texorpdfstring{Подробный пример: $\Phi=\sqrt{1+\sqrt z}$}{Подробный пример: sqrt(1+sqrt z)}}

\textbf{Многочлен.} Если $w^2=1+\sqrt z$, то $(w^2-1)^2=z$. Положим
\[
P(z,w)=(w^2-1)^2-z=w^4-2w^2+1-z,\qquad N=4,\quad P_4\equiv1.
\]
\emph{Неприводимость.} $P$ имеет степень $1$ по $z$. Если $P=AB$ в $\CC[z,w]$, то одна из степеней по $z$ равна $0$, скажем $A=A(w)$;
тогда $A(w)$ делит коэффициенты $P$ при $z^1$ и $z^0$, т.е.\ $-1$ и $(w^2-1)^2$, значит $A$ --- константа. По лемме 30.2 $P$ неприводим над $\CC(z)$.
По теореме 15.3 все элементы с представляющим $P$ образуют одну П.А.Ф.\ (это и есть $\sqrt{1+\sqrt z}$, так как её начальный элемент
алгебраический с этим $P$), а $P$ --- её определяющий многочлен.

\textbf{Сингулярные точки.} $P_4\equiv1$ нулей не имеет. $P'_w=4w^3-4w=4w(w-1)(w+1)$. При $w=0$: $P=1-z=0\Rightarrow z=1$;
при $w=\pm1$: $P=-z=0\Rightarrow z=0$. Итак, $\mathcal S=\{0,1\}$. Над $z=0$: $(w^2-1)^2=0$, корни $\pm1$ двукратные; над $z=1$:
$w^2-1=\pm1$, корни $\pm\sqrt2$ (простые) и $0$ (двукратный).

\textbf{Монодромия.} Возьмём $b=1/4$ и пронумеруем элементы с центром $b$ (главные ветви корней):
$w_1=\sqrt{1+\sqrt z}$, $w_2=-w_1$, $w_3=\sqrt{1-\sqrt z}$, $w_4=-w_3$ (при $z=1/4$: $\sqrt z=1/2$, значения $\pm\sqrt{3/2}$, $\pm\sqrt{1/2}$).
\begin{itemize}
\item \emph{Обход $0$} по окружности $\gamma_0=\{|z|=1/4\}$: $\sqrt z=\frac12e^{it/2}$ переходит в $-\sqrt z$, а $1\pm\sqrt z$ всё время остаётся в
  $\{\Re\zeta>0\}$, где главная ветвь корня голоморфна. Поэтому $w_1\leftrightarrow w_3$, $w_2\leftrightarrow w_4$: $\sigma_0=(13)(24)$.
\item \emph{Обход $1$} по окружности $\gamma_1=\{|z-1|=3/4\}$ (лежит в $\Re z\ge1/4$, $0$ не обходит): $\sqrt z$ (главная) возвращается к себе,
  $\Re(1+\sqrt z)>1$, так что $w_1,w_2$ неподвижны. Далее $1-\sqrt z=\frac{1-z}{1+\sqrt z}$; знаменатель лежит в $\{\Re>1\}$ и не меняет аргумента,
  а $\arg(1-z)$ возрастает на $2\pi$; значит, $\arg(1-\sqrt z)$ возрастает на $2\pi$, и $\sqrt{1-\sqrt z}$ меняет знак: $\sigma_1=(34)$.
\item \emph{Обход $\infty$} (большая окружность, гомотопная $\gamma_0\gamma_1$ в $\CC\setminus\{0,1\}$): $\sigma_\infty=\sigma_1\circ\sigma_0$:
  $1\mapsto3\mapsto4$, $4\mapsto2\mapsto2$, $2\mapsto4\mapsto3$, $3\mapsto1\mapsto1$, т.е.\ $4$-цикл $(1\,4\,2\,3)$.
\end{itemize}

\textbf{Тип особенностей.}
\begin{itemize}
\item $z=0$: $\sigma_0=(13)(24)$ --- над $V(0,\delta)$ две ветви, каждая с ветвлением порядка $2$ (точки ветвления второго порядка).
  Значения ограничены (корни около $\pm1$), $m=0$ в лемме 15.2.
\item $z=1$: $\sigma_1=(34)$ --- две однозначные голоморфные ветви $\pm\sqrt{1+\sqrt z}$ (значения $\pm\sqrt2$ в точке $1$; голоморфны в $B(1,1)$) и
  одно ветвление второго порядка ($\pm\sqrt{1-\sqrt z}\approx\pm\sqrt{(1-z)/2}$).
\item $z=\infty$: $\sigma_\infty$ --- $4$-цикл: одно ветвление порядка $4$; $|w|\sim|z|^{1/4}$ --- алгебраическая особая точка.
\end{itemize}
Все особые точки алгебраические, $\mathcal S_\Phi=\{0,1,\infty\}$ (точку $1$ нельзя выкинуть: элементы $\pm\sqrt{1-\sqrt z}$ не продолжаются
вдоль путей через $1$). Однозначных ветвей над окрестностью точки $1$ ровно две, т.е.\ над $z=1$ лежат только \emph{два} элемента $\Phi$
(корни $\pm\sqrt2$ простые, $P'_w(1,\pm\sqrt2)\ne0$, и по теореме о неявной функции им отвечают два элемента; голоморфного решения с $w(1)=0$ нет: для него
$w^2-1$ голоморфна, равна $-1$ в точке $1$ и $(w^2-1)^2=z$, значит $w^2-1=-\sqrt z$ (главная ветвь), т.е.\ $w^2=1-\sqrt z$ --- функция с
\emph{простым} нулём в точке $1$, а квадрат голоморфной функции имеет нуль чётного порядка).

\textbf{Группа Галуа.} $\pi_1(\CC\setminus\{0,1\},b)$ порождается классами $\gamma_0,\gamma_1$, поэтому $G=\langle(13)(24),(34)\rangle$.
Имеем $\sigma_0\sigma_1\sigma_0^{-1}=(\sigma_0(3)\,\sigma_0(4))=(12)$, так что $G\ni(12),(34),(13)(24)$. Все элементы $G$ сохраняют разбиение
$\{\{1,2\},\{3,4\}\}$ (пары $\pm w$), а группа перестановок, сохраняющих это разбиение, имеет порядок $2\cdot2\cdot2=8$ и порождается
указанными тремя элементами; итак, $G$ --- группа порядка $8$ (диэдральная группа $D_4$ симметрий квадрата с вершинами $1,3,2,4$).
\emph{Некоммутативность:} $\sigma_0\sigma_1\sigma_0^{-1}=(12)\ne(34)=\sigma_1$. \emph{Разрешимость:} $H=\{e,(12),(34),(12)(34)\}\cong\ZZ_2\times\ZZ_2$
--- нормальная подгруппа индекса $2$ (ядро действия на двух блоках), $H$ и $G/H\cong\ZZ_2$ абелевы, значит $G$ разрешима.

\begin{remarkbox}{Проверка (формула Римана--Гурвица)}
Сумма $(\text{порядок ветвления}-1)$: $0\mapsto1+1$, $1\mapsto1$, $\infty\mapsto3$, всего $6$; $2g-2=4\cdot(-2)+6=-2$, т.е.\ род $g=0$. Согласуется
с тем, что $w$ --- глобальный параметр: $z=(w^2-1)^2$.
\end{remarkbox}

\subsection{\texorpdfstring{Второй пример: $\Zh^{-1}(z)$}{Второй пример: обратная функция Жуковского}}

$\Zh(w)=\frac12(w+w^{-1})=z\iff w^2-2zw+1=0$. Многочлен $P=w^2-2zw+1$ неприводим (степень $1$ по $z$, коэффициенты $-2w$ и $w^2+1$ взаимно
просты), $P_2\equiv1$; $P'_w=2(w-z)=0\Rightarrow w=z$, $P(z,z)=1-z^2=0\Rightarrow\mathcal S=\{1,-1\}$. Корни $w=z\pm\sqrt{z^2-1}$, их произведение $1$.
Обход точки $1$ (или $-1$) меняет знак $\sqrt{z^2-1}$ --- транспозиция; обход $\infty$ (обе точки) --- тождественная перестановка: над $V(\infty,\delta)$
две однозначные ветви $w\approx2z$ (полюс) и $w\approx1/(2z)$ (нуль). Группа Галуа $\cong\ZZ_2$ --- коммутативна. Определяющий многочлен:
$w^2-2zw+1$.

\subsection{О теореме Абеля (доп. материал)}

Класс $\mathcal R$ элементов, «выражающихся в радикалах», определяется индуктивно: $\mathcal R_0$ --- элементы констант $E(c)$ и $E(z)$; $\mathcal R_n$
получается из $\mathcal R_{n-1}$ операциями $F\pm G$, $FG$, $F/G$ (при $g(a)\ne0$) и $E_a\bigl(\sqrt[m]{g}\,^{(j)}\bigr)$ (голоморфные ветви корня около
$g(a)\ne0$); $\mathcal R=\bigcup\mathcal R_n$. Можно показать (задачи), что элементы $\mathcal R$ алгебраичны и группа Галуа их П.А.Ф.\ разрешима (используются:
подгруппы и прямые произведения разрешимых групп разрешимы; если $N\triangleleft G$ и $N$, $G/N$ разрешимы, то $G$ разрешима).

Пример: $P_5(z,w)=3w^5-25w^3+60w-z$ (неприводим: степень $1$ по $z$). $\partial_wP_5=15(w^4-5w^2+4)=15(w^2-1)(w^2-4)$, корни $\pm1,\pm2$;
$z=3w^5-25w^3+60w$ даёт $\mathcal S=\{\pm38,\pm16\}$ ($w=1\mapsto38$, $w=2\mapsto16$). При $z=z_j$ у $P_5(z_j,\cdot)$ один двойной корень $w_0$ и три простых;
так как $\partial_zP_5=-1\ne0$ и $\partial^2_wP_5(z_j,w_0)\ne0$ ($w_0$ --- простой корень $\partial_wP_5$), около $w_0$: $z-z_j\approx c(w-w_0)^2$, $c\ne0$ ---
два элемента меняются местами при обходе $z_j$, остальные три однозначны. Значит, $G$ содержит транспозиции (по одной на каждую $z_j$, после
сопряжения путями), порождается ими (петли вокруг $z_j$ порождают $\pi_1$) и транзитивна (неприводимость). Транзитивная подгруппа $S_5$, порождённая
транспозициями, есть $S_5$ (граф на $\{1,\dots,5\}$ с рёбрами-транспозициями связен, а транспозиции связного графа порождают $S_5$). Группа $S_5$
неразрешима, поэтому никакой элемент с определяющим многочленом $P_5$ не лежит в $\mathcal R$: уравнение $P_5=0$ неразрешимо в радикалах.
Над $\infty$ --- одно ветвление порядка $5$ ($w\sim(z/3)^{1/5}$).

\begin{remarkbox}{Коротко: план ответа на экзамене}
\begin{enumerate}
\item Алгебраический элемент, представляющий и определяющий многочлен; определяющий единствен с точностью до множителя из $\CC(z)$,
  неприводим, делит любой представляющий (деление с остатком + теорема единственности).
\item Сингулярные точки; при $\gcd(P,P'_w)=1$ их конечное число (соотношение Безу).
\item Над несингулярной точкой --- ровно $N$ голоморфных решений (теорема о неявной функции, $\partial_{\bar z}f=0$); элементы продолжаются вдоль
  путей в $\CC\setminus\mathcal S$, радиусы $\ge\dist(a,\mathcal S)$ (монодромия).
\item Лемма 15.2: степенные оценки корней около $z_1\in\mathcal S$ и около $\infty$.
\item Теорема 15.3: неприводимый $P$ задаёт П.А.Ф., $N$-значную над $\CC\setminus\mathcal S$; особенности --- конечные ветвления и полюсы.
\item Опр.~15.10--15.11: алгебраическая О.Т.\ (конечное ветвление + степенной рост), алгебраическая П.А.Ф.
\item Теорема 15.4: симметрические функции однозначны и имеют не более чем полюсы $\Rightarrow$ рациональны; Виета даёт $P$.
\item Группа монодромии; пример $\sqrt{1+\sqrt z}$: $\mathcal S=\{0,1\}$, $G\cong D_4$; $\Zh^{-1}$: $G\cong\ZZ_2$.
\end{enumerate}
\end{remarkbox}

% =====================================================================
\ticket{31}{Одномерные комплексные многообразия и их голоморфные отображения. Неразветвлённые голоморфные накрытия. Теорема о поднятии путей}

\begin{defs}
\begin{itemize}
\item \textbf{Хаусдорфово пространство}: топологическое пространство, в котором любые две различные точки имеют непересекающиеся
  окрестности. \textbf{Линейно связное}: любые две точки соединяются путём (непрерывным $\gamma:[0,1]\to X$).
\item \textbf{Одномерное комплексное многообразие} (слайд 418). Пусть $\mathcal X$ --- хаусдорфово топологическое пространство,
  $\{U_\alpha\}_{\alpha\in\mathcal A}$ --- покрытие $\mathcal X$ открытыми множествами, $\varphi_\alpha:U_\alpha\to B_\alpha$ --- гомеоморфизмы на круги
  $B_\alpha\subset\CC$, причём при $U_{\alpha\beta}:=U_\alpha\cap U_\beta\ne\varnothing$ \emph{отображения перехода (замены карт)}
  \[
  \varphi_\beta\circ\varphi_\alpha^{-1}:\varphi_\alpha(U_{\alpha\beta})\to\varphi_\beta(U_{\alpha\beta}),\qquad
  \varphi_\alpha\circ\varphi_\beta^{-1}:\varphi_\beta(U_{\alpha\beta})\to\varphi_\alpha(U_{\alpha\beta})
  \]
  голоморфны (их области определения открыты, так как $\varphi_\alpha$ --- гомеоморфизмы). Тогда $(\mathcal X,\{U_\alpha\},\{\varphi_\alpha\})$ ---
  одномерное комплексное многообразие; $\{U_\alpha\}$ (вместе с $\varphi_\alpha$) --- \emph{атлас}, $\varphi_\alpha$ --- \emph{локальные карты}
  (локальные координаты). Отображения перехода --- взаимно обратные биекции, т.е.\ конформные отображения открытых множеств.
\item \textbf{Голоморфное отображение.} Пусть $(\mathcal X,\{U_\alpha\},\{\varphi_\alpha\})$, $(\mathcal Y,\{V_\beta\},\{\psi_\beta\})$ --- многообразия. Непрерывное
  отображение $f:\mathcal X\to\mathcal Y$ \emph{голоморфно}, если все функции $\psi_\beta\circ f\circ\varphi_\alpha^{-1}$ голоморфны на своих
  (открытых в силу непрерывности $f$) областях определения $\varphi_\alpha(U_\alpha\cap f^{-1}(V_\beta))$.
  Функция $f:\mathcal X\to\CC$ голоморфна, если голоморфны все $f\circ\varphi_\alpha^{-1}:B_\alpha\to\CC$.
\item \textbf{Биголоморфное отображение}: голоморфная биекция $f:\mathcal X\to\mathcal Y$, у которой $f^{-1}$ голоморфно.
\item \textbf{Неразветвлённое голоморфное накрытие} (слайд 420): отображение $\pi:\mathcal X\to\mathcal Y$ многообразий такое, что
  (1) $\mathcal X$ линейно связно; (2) $\pi$ голоморфно и сюръективно; (3) у любой $z\in\mathcal Y$ есть окрестность $V$ (\emph{правильно накрываемая})
  и набор попарно непересекающихся открытых $U_j\subset\mathcal X$, $j\in J$, с $\pi^{-1}(V)=\bigsqcup_{j\in J}U_j$, причём каждое
  $\pi|_{U_j}:U_j\to V$ биголоморфно. Множества $U_j$ --- \emph{листы} над $V$; $\pi^{-1}(z)$ --- \emph{слой}; мощность слоя (одна и та же для всех $z$,
  см.\ следствие леммы 15.3) --- \emph{число листов}.
\item \textbf{Поднятие пути}: для пути $\gamma:[0,1]\to\mathcal Y$ --- путь $\tilde\gamma:[0,1]\to\mathcal X$ с $\pi\circ\tilde\gamma=\gamma$.
\item \textbf{Послойный изоморфизм} (слайд 421) накрытий $\pi_1:\mathcal X_1\to\mathcal Y$ и $\pi_2:\mathcal X_2\to\mathcal Y$: биголоморфное
  $f:\mathcal X_1\to\mathcal X_2$ с $\pi_2\circ f=\pi_1$ (слои переходят в слои).
\item \textbf{Число Лебега}: для открытого покрытия компактного метрического пространства $K$ найдётся $\lambda>0$ такое, что любое
  подмножество $K$ диаметра $<\lambda$ целиком лежит в одном из элементов покрытия.
\end{itemize}
\end{defs}

\begin{remarkbox}{Замечание (о непрерывности)}
На слайде в определении голоморфного отображения не оговорена непрерывность $f$. Она нужна, чтобы множества
$\varphi_\alpha(U_\alpha\cap f^{-1}(V_\beta))$ были открытыми и голоморфность имела смысл (для функций $f:\mathcal X\to\CC$ непрерывность
следует из голоморфности $f\circ\varphi_\alpha^{-1}$). Голоморфность --- локальное свойство, его достаточно проверять в одной карте около каждой
точки: если $f\circ\varphi_\alpha^{-1}$ голоморфна, то $f\circ\varphi_{\alpha'}^{-1}=(f\circ\varphi_\alpha^{-1})\circ(\varphi_\alpha\circ\varphi_{\alpha'}^{-1})$ на
пересечении --- композиция голоморфных.
\end{remarkbox}

\subsection{Примеры многообразий и голоморфных отображений}

\textbf{1. Область $D\subset\CC$.} Атлас: любое семейство открытых кругов $U_\alpha\subset D$, покрывающих $D$, карты $\varphi_\alpha=\mathrm{id}$.
Отображения перехода --- тождественные, голоморфны. $f:D\to\CC$ голоморфна в смысле многообразий $\iff$ $f|_{U_\alpha}$ голоморфна при всех $\alpha$
$\iff$ $f\in H(D)$ (голоморфность локальна). Для областей $D,G$ отображение $f:D\to G$ голоморфно в новом смысле $\iff$ $f\in H(D)$;
биголоморфное отображение областей --- это конформное отображение $D$ на $G$ (для голоморфной биекции обратное отображение голоморфно
автоматически, так как $f'\ne0$ у однолистной функции).

\textbf{2. Расширенная плоскость $\Ch$.} Топология $\Ch$ --- топология сферы Римана (окрестности $\infty$ --- $\{|z|>R\}\cup\{\infty\}$),
она хаусдорфова. Атлас из двух карт:
\[
U_1=\{|z|<2\},\ \varphi_1(z)=z;\qquad U_2=\{|z|>1/2\}\cup\{\infty\},\ \varphi_2(z)=1/z\ (\varphi_2(\infty)=0).
\]
$\varphi_1:U_1\to B(0,2)$ и $\varphi_2:U_2\to B(0,2)$ --- гомеоморфизмы. $U_1\cap U_2=\{1/2<|z|<2\}$, $\varphi_1(U_{12})=\varphi_2(U_{12})=\{1/2<|\zeta|<2\}$,
и $\varphi_2\circ\varphi_1^{-1}(\zeta)=\varphi_1\circ\varphi_2^{-1}(\zeta)=1/\zeta$ --- голоморфна в кольце. Значит, $\Ch$ --- одномерное комплексное многообразие.
Функция $f:\Ch\to\CC$ голоморфна $\iff$ $f$ голоморфна в $\CC$ и $\zeta\mapsto f(1/\zeta)$ голоморфна около $0$ --- ровно определение голоморфности
в $\infty$ из начала курса. (По теореме Лиувилля такие $f$ --- константы.) Отображение $f:\Ch\to\Ch$ голоморфно, например, если $f$ рациональна
(около полюса $a$ в карте $1/w$ имеем $1/f$ --- голоморфную функцию).

\subsection{Неразветвлённые голоморфные накрытия: примеры}

\begin{thmbox}{Лемма 31.1 (критерий накрытия для областей)}
Пусть $G,D\subset\CC$ --- области, $g\in H(G)$, $g(G)=D$. Пусть у каждой точки $z_0\in D$ есть круг $V\ni z_0$, $V\subset D$, такой, что для любых
$z'\in V$ и $w'\in G$ с $g(w')=z'$ найдётся $h\in H(V)$ с $h(V)\subset G$, $g\circ h=\mathrm{id}_V$ и $h(z')=w'$ (\emph{локальное обращение через $w'$}).
Тогда $g:G\to D$ --- неразветвлённое голоморфное накрытие, и $V$ --- правильно накрываемая окрестность.
\end{thmbox}
\begin{proof}
(1), (2) очевидны ($G$ --- область, значит линейно связна). Пусть $\{h_j\}_{j\in J}$ --- все голоморфные $h:V\to G$ с $g\circ h=\mathrm{id}_V$.
Каждое $h_j$ инъективно (имеет левый обратный $g$), значит, $h_j(V)$ открыто (принцип открытости) и $g|_{h_j(V)}:h_j(V)\to V$ --- биекция
с голоморфным обратным $h_j$, т.е.\ биголоморфна. Объединение $h_j(V)$ равно $g^{-1}(V)$: включение «$\subset$» очевидно, а «$\supset$» --- это условие леммы.
\emph{Дизъюнктность:} пусть $h_j(z_1)=h_k(z_2)$. Применяя $g$, получаем $z_1=z_2$. Множество $E=\{z\in V:h_j(z)=h_k(z)\}$ непусто и замкнуто в $V$.
Оно открыто: если $h_j(z)=h_k(z)=w$, то $g'(w)\ne0$ (так как $g'(h_j(z))h_j'(z)=1$), и по теореме об обратной функции $g$ инъективна в некоторой
окрестности $O\ni w$; при $z'$ близких к $z$ оба значения $h_j(z'),h_k(z')$ лежат в $O$ и $g(h_j(z'))=z'=g(h_k(z'))$, откуда $h_j(z')=h_k(z')$.
Круг $V$ связен, значит $E=V$, т.е.\ $h_j=h_k$. Итак, различные $h_j$ имеют непересекающиеся образы.
\end{proof}

\textbf{Пример 1: $z\mapsto e^z$, $\CC\to\CC^*$.} Сюръективно. Для $z_0\ne0$ возьмём круг $V\ni z_0$, $0\notin V$; в нём есть голоморфная
ветвь логарифма $L$, и все решения $e^w=z'$ ($z'\in V$) --- это $L(z')+2\pi ik$, $k\in\ZZ$; соответствующие обращения $h_k=L+2\pi ik$. По лемме 31.1 ---
накрытие; листы $h_k(V)$, число листов счётно.

\textbf{Пример 2: $z\mapsto z^n$, $\CC^*\to\CC^*$, $n\ge2$.} Для круга $V\not\ni0$ и ветви $r_0=\exp(L/n)$ корня $n$-й степени все решения $w^n=z'$ --- это
$\omega_n^kr_0(z')$, $\omega_n=e^{2\pi i/n}$, $k=0,\dots,n-1$; $h_k=\omega_n^kr_0$. $\CC^*$ линейно связна. Накрытие $n$-листное.

\textbf{Пример 3: $h(z)=\exp(\exp z)$.}
\emph{Как отображение $\CC\to\CC^*$ оно не накрытие}, хотя сюръективно ($e^w$ при $w\in\CC^*$ принимает все значения из $\CC^*$, в том числе
$1=e^{2\pi i}$) и локально биголоморфно ($h'=h\,e^z\ne0$). Действительно, у накрытия любой путь поднимается (лемма 15.3 ниже). Рассмотрим
$\gamma(t)=e^{1-t}$, $t\in[0,1]$, $\gamma(0)=e=h(0)$, $\gamma(1)=1$. Если бы $\tilde\gamma$ --- поднятие с $\tilde\gamma(0)=0$, то $u(t)=e^{\tilde\gamma(t)}$ ---
непрерывный путь с $e^{u(t)}=e^{1-t}$, т.е.\ $u(t)-(1-t)\in2\pi i\ZZ$; непрерывная функция со значениями в дискретном множестве постоянна,
и $u(0)=1$, значит $u(t)=1-t$, $u(1)=0$ --- но $e^{\tilde\gamma(1)}\ne0$. Противоречие.

\emph{Как отображение $\mathcal X=\CC\setminus h^{-1}(1)\to\CC\setminus\{0,1\}$ это накрытие.} Здесь $h^{-1}(1)=\{z:e^z\in2\pi i\ZZ\}=\{z:e^z\in2\pi i(\ZZ\setminus\{0\})\}$,
а $h^{-1}(0)=\varnothing$; значит $h(\mathcal X)\subset\CC\setminus\{0,1\}$.
\begin{itemize}
\item \emph{$\mathcal X$ линейно связно}: $h^{-1}(1)=\{\ln(2\pi|k|)+i(\pm\frac\pi2+2\pi j)\}$ --- замкнутое дискретное (в каждом круге конечное) множество. Для
  $p,q\in\mathcal X$ ломаные $p\to m\to q$, где $m$ пробегает серединный перпендикуляр к $[p,q]$, попарно пересекаются лишь в $p,q$; их
  несчётно, а точек $h^{-1}(1)$ счётно, так что одна из ломаных лежит в $\mathcal X$.
\item \emph{Сюръективность}: при $c\ne0,1$ возьмём $w$ с $e^w=c$; тогда $w\notin2\pi i\ZZ$ (иначе $c=1$), в частности $w\ne0$; возьмём $z$ с $e^z=w$; тогда
  $z\in\mathcal X$ и $h(z)=c$.
\item \emph{Локальные обращения}: пусть $V\subset\CC\setminus\{0,1\}$ --- круг, $L_k=L+2\pi ik$ --- ветви логарифма на $V$. Образ $L_k(V)$ --- односвязная
  область (конформный образ круга), не содержащая $0$ (так как $1\notin V$); на ней есть ветви логарифма $M_{k,j}$, $j\in\ZZ$. Функции
  $h_{k,j}=M_{k,j}\circ L_k$ голоморфны на $V$, $h\circ h_{k,j}=\mathrm{id}$, значения в $\mathcal X$ ($e^{h_{k,j}}=L_k\notin2\pi i\ZZ$, так как $1\notin V$). Если $h(z')=\zeta\in V$,
  то $e^{z'}$ --- логарифм $\zeta$, т.е.\ $e^{z'}=L_k(\zeta)$ при некотором $k$, а $z'$ --- логарифм $L_k(\zeta)$, т.е.\ $z'=h_{k,j}(\zeta)$ при некотором $j$.
\end{itemize}
По лемме 31.1 $h:\mathcal X\to\CC\setminus\{0,1\}$ --- (счётнолистное) неразветвлённое голоморфное накрытие.

\begin{remarkbox}{Примеры послойных изоморфизмов (слайд 421)}
$f_k(z)=z+2\pi ik$ --- послойный изоморфизм накрытия $e^z:\CC\to\CC^*$ на себя: $f_k$ биголоморфно ($f_k^{-1}=f_{-k}$) и $e^{f_k(z)}=e^z$.
$z\mapsto\omega_n^kz$ ($k=0,\dots,n-1$) --- послойный изоморфизм накрытия $z^n:\CC^*\to\CC^*$: биголоморфно $\CC^*\to\CC^*$ и $(\omega_n^kz)^n=z^n$.
\end{remarkbox}

\subsection{Лемма о поднятии путей}

\begin{thmbox}{Лемма 15.3 (о поднятии путей)}
Пусть $\pi:\mathcal X\to\mathcal Y$ --- неразветвлённое голоморфное накрытие. Для любой $z_0\in\mathcal Y$, любого пути $\gamma:[0,1]\to\mathcal Y$ с
$\gamma(0)=z_0$ и любой $\zeta_0\in\pi^{-1}(z_0)$ существует единственный путь $\tilde\gamma:[0,1]\to\mathcal X$ с $\tilde\gamma(0)=\zeta_0$ и
$\pi(\tilde\gamma(t))=\gamma(t)$ при всех $t\in[0,1]$.
\end{thmbox}
\begin{proof}
\emph{Разбиение.} Для каждой точки $z\in\gamma^*$ возьмём правильно накрываемую окрестность $V(z)$. Множества $\gamma^{-1}(V(z))$ открыты в $[0,1]$
(непрерывность $\gamma$) и покрывают $[0,1]$. Пусть $\lambda>0$ --- число Лебега этого покрытия и $T=(0=t_0<t_1<\dots<t_N=1)$ --- разбиение с
$\max(t_j-t_{j-1})<\lambda$. Тогда $[t_{j-1},t_j]\subset\gamma^{-1}(V_j)$ для некоторой правильно накрываемой $V_j$, т.е.\ $\gamma([t_{j-1},t_j])\subset V_j$.
(На слайде говорится о конечном подпокрытии $\gamma^*$; само по себе оно не даёт разбиения, нужное разбиение даёт именно число Лебега.)

\emph{Существование.} Строим $\tilde\gamma$ на $[0,t_j]$ индукцией по $j$. Пусть $\tilde\gamma$ уже построен на $[0,t_{j-1}]$ (при $j=1$ --- только
$\tilde\gamma(0)=\zeta_0$), непрерывен и $\pi\circ\tilde\gamma=\gamma$. Точка $\tilde\gamma(t_{j-1})$ лежит в $\pi^{-1}(V_j)=\bigsqcup_iU_{j,i}$, т.е.\ ровно в одном листе $U_{j,i_j}$.
Положим
\[
\tilde\gamma(t)=(\pi|_{U_{j,i_j}})^{-1}(\gamma(t)),\qquad t\in[t_{j-1},t_j].
\]
Это непрерывно (композиция непрерывных), $\pi\circ\tilde\gamma=\gamma$, и в точке $t_{j-1}$ новое значение равно старому:
$(\pi|_{U_{j,i_j}})^{-1}(\gamma(t_{j-1}))$ --- единственная точка $U_{j,i_j}$ над $\gamma(t_{j-1})$, а $\tilde\gamma(t_{j-1})$ такова. Склейка двух непрерывных
отображений на отрезках с общим концом непрерывна. За $N$ шагов получаем $\tilde\gamma$ на $[0,1]$.

\emph{Единственность.} Пусть $\tilde\gamma_1,\tilde\gamma_2$ --- два поднятия с началом $\zeta_0$, $E=\{t\in[0,1]:\tilde\gamma_1(t)=\tilde\gamma_2(t)\}$.
$0\in E$. $E$ \emph{замкнуто}: если $t\notin E$, то точки $\tilde\gamma_1(t)\ne\tilde\gamma_2(t)$ имеют непересекающиеся окрестности (хаусдорфовость
$\mathcal X$), и по непрерывности при $\xi$, близких к $t$, $\tilde\gamma_1(\xi)$ и $\tilde\gamma_2(\xi)$ лежат в этих окрестностях, т.е.\ различны; дополнение
$E$ открыто. $E$ \emph{открыто}: пусть $t\in E$, $V$ --- правильно накрываемая окрестность $\gamma(t)$, $U$ --- лист, содержащий
$\tilde\gamma_1(t)=\tilde\gamma_2(t)$. По непрерывности найдётся интервал $I\ni t$ (открытый в $[0,1]$) с $\tilde\gamma_1(I),\tilde\gamma_2(I)\subset U$. При $\xi\in I$
обе точки $\tilde\gamma_i(\xi)$ лежат в $U$ над одной точкой $\gamma(\xi)$, а $\pi|_U$ инъективно; значит, $\tilde\gamma_1(\xi)=\tilde\gamma_2(\xi)=(\pi|_U)^{-1}(\gamma(\xi))$.
Так как $[0,1]$ связен, $E=[0,1]$.
\end{proof}

\begin{thmbox}{Следствие (о мощности слоя)}
Если $\pi:\mathcal X\to\mathcal Y$ --- неразветвлённое голоморфное накрытие, то мощность $\pi^{-1}(z)$ одинакова для всех $z\in\mathcal Y$.
\end{thmbox}
\begin{proof}
$\mathcal Y=\pi(\mathcal X)$ --- непрерывный образ линейно связного пространства, значит линейно связно. Пусть $z_1,z_2\in\mathcal Y$, $\gamma$ --- путь из
$z_1$ в $z_2$, $\gamma^-(t)=\gamma(1-t)$. Определим $\Gamma:\pi^{-1}(z_1)\to\pi^{-1}(z_2)$: $\Gamma(\zeta)=\tilde\gamma_\zeta(1)$, где $\tilde\gamma_\zeta$ --- поднятие $\gamma$ с
началом $\zeta$; аналогично $\Gamma^-:\pi^{-1}(z_2)\to\pi^{-1}(z_1)$ через $\gamma^-$. Если $\tilde\gamma_\zeta$ --- поднятие $\gamma$ из $\zeta$, то $t\mapsto\tilde\gamma_\zeta(1-t)$ ---
поднятие $\gamma^-$ из $\Gamma(\zeta)$, и по единственности $\Gamma^-(\Gamma(\zeta))=\zeta$; симметрично $\Gamma(\Gamma^-(\eta))=\eta$. Итак, $\Gamma$ --- биекция (на
слайде инъективность выводится из единственности, а «очевидная» сюръективность --- это как раз существование $\Gamma^-$).
\end{proof}

\begin{remarkbox}{Замечание}
В примере 3 слои $h:\CC\to\CC^*$ над $1$ и над $c\ne1$ оба счётны, так что следствие не различает их --- поэтому для доказательства того, что это не
накрытие, мы использовали саму лемму о поднятии путей. Локальная биголоморфность ($h'\ne0$) не влечёт свойства накрытия.
\end{remarkbox}

\subsection{\texorpdfstring{П.А.Ф. $f\circ\pi^{-1}$, порождённая голоморфной функцией на накрытии}{П.А.Ф., порождённая голоморфной функцией на накрытии}}

Пусть $D\subset\CC$ --- область, $\pi:\mathcal X\to D$ --- неразветвлённое голоморфное накрытие, $f:\mathcal X\to\CC$ --- голоморфная функция.
Если $V$ --- правильно накрываемая окрестность, то и любой открытый круг $V'\subset V$ правильно накрываемый (листы $U_j\cap\pi^{-1}(V')$). Поэтому
для $\zeta\in\mathcal X$ возьмём правильно накрываемый круг $V$ \emph{с центром} $\pi(\zeta)$, лист $U(\zeta)\ni\zeta$ над ним и \emph{локальное сечение}
$s_\zeta=(\pi|_{U(\zeta)})^{-1}:V\to U(\zeta)$ (голоморфно). Положим $h_\zeta=f\circ s_\zeta\in H(V)$ (голоморфна как композиция голоморфных отображений;
в карте $\varphi$ около $\zeta$: $h_\zeta=(f\circ\varphi^{-1})\circ(\varphi\circ s_\zeta)$) и
\[
E(\zeta):=E_{\pi(\zeta)}(h_\zeta)\quad\text{--- канонический элемент с центром } \pi(\zeta).
\]

\begin{detail}{$E(\zeta)$ не зависит от выбора $V$; совпадение сечений}
Если $s,s'$ --- два локальных сечения (над окрестностями $V,V'$ точки $z$) с $s(z)=s'(z)=\zeta$, $U=s(V)$, $U'=s'(V')$, то $W=U\cap U'$ открыто,
содержит $\zeta$, $\pi(W)$ открыто ($\pi|_U$ --- гомеоморфизм на $V$), и для $y\in\pi(W)$, $y=\pi(x)$, $x\in W$: $s(y)=x=s'(y)$ (инъективность $\pi$ на $U$ и на $U'$).
Значит, $s=s'$ около $z$, функции $f\circ s$, $f\circ s'$ совпадают около $z$ и имеют один ряд Тейлора.
\end{detail}

\begin{thmbox}{Предложение 31.2 (слайды 423--424)}
(а) Если $\gamma:[0,1]\to D$ --- путь, $\tilde\gamma$ --- его поднятие, то $\{E(\tilde\gamma(t))\}_{t\in[0,1]}$ --- аналитическое продолжение $E(\tilde\gamma(0))$ вдоль $\gamma$.
(б) Для любой $\zeta_0\in\mathcal X$ П.А.Ф., порождённая $E(\zeta_0)$, содержит все $E(\zeta)$, $\zeta\in\mathcal X$, а продолжения $E(\zeta_0)$ вдоль путей в $D$
--- это в точности элементы $E(\zeta)$. Значит, эта П.А.Ф.\ не зависит от $\zeta_0$; она обозначается $f\circ\pi^{-1}$, а $\{E(\zeta)\}_{\zeta\in\mathcal X}$ --- её
аналитическая ветвь над $D$.
\end{thmbox}
\begin{proof}
(а) Пусть $t\in[0,1]$, $\zeta=\tilde\gamma(t)$, $V$ --- правильно накрываемый круг с центром $\gamma(t)$, $U$ --- лист над $V$, содержащий $\zeta$. Круг
элемента $E(\zeta)$ содержит $V$ (функция $h_\zeta$ голоморфна в $V$). По непрерывности $\tilde\gamma$ найдётся связная окрестность $I_t\ni t$ с
$\tilde\gamma(I_t)\subset U$. Для $\xi\in I_t$: $\gamma(\xi)=\pi(\tilde\gamma(\xi))\in V$, т.е.\ $\gamma(\xi)$ лежит в круге $E(\zeta)$; сечения $s_\zeta$ и $s_{\tilde\gamma(\xi)}$ оба
переводят $\gamma(\xi)$ в $\tilde\gamma(\xi)$ и по детали выше совпадают около $\gamma(\xi)$. Значит, функции элементов $E(\zeta)$ и $E(\tilde\gamma(\xi))$ совпадают
около $\gamma(\xi)$, т.е.\ на непустом открытом подмножестве (выпуклого) пересечения кругов; по теореме единственности $E(\tilde\gamma(\xi))$ --- НАП
$E(\zeta)$. Это и есть условие Опр.~15.4.

(б) Пусть $\zeta_1\in\mathcal X$. $\mathcal X$ линейно связно: возьмём путь $\alpha$ в $\mathcal X$ из $\zeta_0$ в $\zeta_1$; $\gamma=\pi\circ\alpha$ --- путь в $D$, а $\alpha$ --- его
поднятие. По (а) $E(\zeta_1)$ --- продолжение $E(\zeta_0)$ вдоль $\gamma$. Обратно, если $\gamma$ --- путь в $D$ из $\pi(\zeta_0)$, то по лемме 15.3 он поднимается в
$\tilde\gamma$ с $\tilde\gamma(0)=\zeta_0$, и по (а) и единственности продолжения результат продолжения $E(\zeta_0)$ вдоль $\gamma$ равен $E(\tilde\gamma(1))$. В частности,
$E(\zeta_0)$ продолжается вдоль любого пути в $D$, и совокупность продолжений --- $\{E(\zeta)\}$, т.е.\ это ветвь над $D$ (Опр.~15.6).
\end{proof}

\begin{remarkbox}{Замечание}
П.А.Ф.\ $f\circ\pi^{-1}$ может «выходить» за пределы $D$: например, для $\mathcal X=D=\CC^*$, $\pi=\mathrm{id}$, $f(z)=z$ получаем П.А.Ф.\ $z$ с
областью определения $\CC$. Число элементов ветви над точкой $z\in D$ не больше числа листов (разные $\zeta$ из слоя могут давать один элемент,
например при $f\equiv\mathrm{const}$).
\end{remarkbox}

\begin{remarkbox}{Коротко: план ответа на экзамене}
\begin{enumerate}
\item Определение многообразия: хаусдорфово пространство, атлас, карты-гомеоморфизмы на круги, голоморфные замены карт.
\item Примеры: область (тождественные карты), $\Ch$ (карты $z$ и $1/z$, замена $1/\zeta$ в кольце).
\item Голоморфные отображения и функции (в картах), биголоморфизмы; для областей --- обычные понятия, для $\Ch$ --- голоморфность в $\infty$.
\item Неразветвлённое накрытие: связность, сюръективность, листы над правильно накрываемыми окрестностями; примеры $e^z$, $z^n$; $\exp(\exp z)$ ---
  не накрытие $\CC\to\CC^*$ (путь $e^{1-t}$ не поднимается), но накрытие $\CC\setminus h^{-1}(1)\to\CC\setminus\{0,1\}$.
\item Послойные изоморфизмы: $z+2\pi ik$, $\omega_n^kz$.
\item Лемма о поднятии: число Лебега $\to$ разбиение $\to$ подъём по листам; единственность: множество совпадения открыто-замкнуто (хаусдорфовость).
\item Следствие: слои равномощны (биекция через поднятия $\gamma$ и $\gamma^-$).
\item П.А.Ф.\ $f\circ\pi^{-1}$: элементы $E_{\pi(\zeta)}(f\circ s_\zeta)$, продолжение вдоль $\gamma$ --- вдоль поднятия.
\end{enumerate}
\end{remarkbox}

% =====================================================================
\ticket{32}{Понятие о римановой поверхности П.А.Ф.}

\begin{defs}
\begin{itemize}
\item \textbf{Элемент, канонический элемент, НАП, П.А.Ф.\ $\Phi$, область определения $D_\Phi$, аналитическая ветвь} --- см.\ рамку билета 30
  (Опр.~15.2--15.6). $E_a(f)$ --- канонический элемент функции $f$ с центром $a$. Для элемента $F=(B(a,r),f)$ пишем $r_F=r$ (радиус, $\le+\infty$).
\item \textbf{Одномерное комплексное многообразие, голоморфные отображения и функции, биголоморфизм, неразветвлённое голоморфное накрытие,
  послойный изоморфизм, лемма о поднятии путей} --- см.\ билет 31.
\item \textbf{Риманова поверхность П.А.Ф.} $\Phi$ (слайды 424--425): множество $\mathcal R=\mathcal R(\Phi)$, \emph{точками которого являются элементы $\Phi$};
  \emph{проекция} $\pi:\mathcal R\to D_\Phi$, $\pi(F)=a$ для $F=(B(a,r),f)$ (элемент проектируется в свой центр);
  \emph{$\eps$-окрестность} элемента $F_0$:
  \[
  W_\eps(F_0)=\{F\in\Phi:\ |\pi(F)-\pi(F_0)|<\eps,\ F \text{ --- НАП } F_0\};
  \]
  топология --- порождённая базой $\{W_\eps(F)\}$; атлас --- $U_\alpha=W_\eps(F)$, $\varphi_\alpha=\pi|_{U_\alpha}$; \emph{функция} $f_\Phi:\mathcal R\to\CC$,
  $f_\Phi(F)=f(\pi(F))$ --- значение функции элемента в его центре.
\item \textbf{Пример-идея.} Для $\Phi=\sqrt z$ над каждой точкой $z\ne0$ два элемента, и $\mathcal R$ --- «двулистная поверхность» над $\CC^*$; на ней
  $\sqrt z$ становится обычной однозначной функцией $f_\Phi$.
\item \textbf{Берега разреза.} Для области вида $\Omega=\CC\setminus L$ ($L$ --- луч или объединение лучей на $\RR$) и $x\in L$: $x^+$ --- «верхний берег»
  (предел из $\Im z>0$), $x^-$ --- «нижний берег» (предел из $\Im z<0$). $x^{\pm}(k)$ --- берег на $k$-м экземпляре («листе»). $\Lambda_{\alpha,\beta}=\{\alpha<\arg w<\beta\}$,
  $\RR_+=[0,+\infty)$.
\end{itemize}
\end{defs}

\subsection{Идея Римана}

Конструкция Вейерштрасса (элементы, продолжение, П.А.Ф.) --- рабочий инструмент, но П.А.Ф.\ не функция в обычном смысле. Риман предложил
трактовать П.А.Ф.\ $\Phi$ над областью $D$ как \emph{однозначную} функцию $\tilde\Phi$ на поверхности, многолистно лежащей над $D$: над каждой точкой
$z$ лежит столько точек поверхности, сколько у $\Phi$ различных элементов с центром $z$. Сначала --- примеры, затем формальная конструкция и
предложение 15.9, и в конце (предл.~32.6) --- доказательство того, что для примеров формальная конструкция даёт именно описанные поверхности.

Во всех примерах схема одна: есть область $G\subset\CC$ и голоморфная $g:G\to D$, являющаяся неразветвлённым накрытием (проверяем по лемме 31.1),
и П.А.Ф.\ $\Phi$ --- обратная к $g$, т.е.\ $\Phi=\iota\circ g^{-1}$, $\iota(w)=w$ (предл.~31.2). Область $G$ разбивается на «фундаментальные области» $G_k$,
каждая из которых $g$ конформно отображает на одну и ту же область $\Omega\subset D$; образ $g(G_k)$ --- «$k$-й лист». Общие граничные кривые соседних
$G_k$ переходят в разрезы, и то, \emph{какие берега каких листов} они склеивают, вычисляется по знаку $\Im g$ около границы. «Склеенная поверхность»
--- это просто $G$, а $g^{-1}$ --- тождественная функция на ней, т.е.\ однозначна.

\subsection{Примеры римановых поверхностей}

\textbf{1. $\sqrt z$.} Отображение $w\mapsto w^2$ конформно переводит $\Lambda_{0,\alpha}$ на $\Lambda_{0,2\alpha}$ при $\alpha\in(0,\pi]$; при $\alpha>\pi$ оно остаётся
локально конформным ($2w\ne0$), но теряет инъективность. $G=\CC^*$, $g(w)=w^2$ --- двулистное накрытие $\CC^*\to\CC^*$ (билет 31, пример 2).
Листы: $G_1=\{\Im w>0\}\to\Omega=\CC\setminus\RR_+$ (лист 1), $G_2=\{\Im w<0\}\to\Omega$ (лист 2).
\emph{Склейка.} Луч $w=x>0$: подход из $G_1$ ($w=\sqrt x\,e^{i\theta}$, $\theta\to0+$): $\arg w^2=2\theta\to0+$ --- верхний берег листа 1; из $G_2$ ($\theta\to2\pi-$):
$\arg w^2\to4\pi-\equiv2\pi-$ --- нижний берег листа 2. Луч $w<0$: из $G_1$ ($\theta\to\pi-$): $\arg w^2\to2\pi-$ --- нижний берег листа 1; из $G_2$ ($\theta\to\pi+$):
$\arg w^2\to2\pi+\equiv0+$ --- верхний берег листа 2. Итак, \emph{$x^+(1)\sim x^-(2)$, $x^-(1)\sim x^+(2)$} ($x>0$), как на слайде 414. Окрестность склеенной
точки --- это образ малого круга с центром на вещественной оси плоскости $w$: верхний полукруг на одном листе плюс нижний на другом.
Отображение $z\mapsto\sqrt z$ (обратное к $w^2$) взаимно однозначно отображает $X_2$ на $\CC^*$ и однозначно на $X_2$.

\textbf{2. $\sqrt[n]z$, $n\ge2$.} $G=\CC^*$, $g=w^n$ ($n$-листное накрытие). Листы $G_k=\Lambda_{2\pi(k-1)/n,\,2\pi k/n}\to\Omega=\CC\setminus\RR_+$, $k=1,\dots,n$.
Общая граница $G_k$ и $G_{k+1}$ --- луч $\arg w=2\pi k/n$, $g$ переводит его в $(0,+\infty)$; из $G_k$ ($\arg w\to2\pi k/n-$) $\arg w^n\to2\pi k-$: нижний берег
листа $k$; из $G_{k+1}$: $\arg w^n\to2\pi k+$: верхний берег листа $k+1$. Граница $G_n$ и $G_1$ --- луч $\arg w=0$. Склейка:
\[
x^-(k)\sim x^+(k+1),\ k=1,\dots,n-1;\qquad x^-(n)\sim x^+(1)\qquad(x>0).
\]

\begin{remarkbox}{Замечание (слайд 415)}
На слайде: «$x^+(k)$ отождествляется с $x^-(k+1)$ при $k=1,2,\dots,k-1$; $x^+(n)$ --- с $x^-(1)$». Во-первых, опечатка: должно быть $k=1,\dots,n-1$.
Во-вторых, при указанной на слайде нумерации листов (лист $k$ --- образ сектора $\Lambda_{2\pi(k-1)/n,2\pi k/n}$) верна склейка $x^-(k)\sim x^+(k+1)$;
вариант слайда соответствует обратной нумерации листов ($k\mapsto n+1-k$) и даёт изоморфную поверхность. При $n=2$ оба варианта совпадают.
То же для $\Ln z$ ниже.
\end{remarkbox}

\textbf{3. $\Ln z$.} $G=\CC$, $g=e^w$ (счётнолистное накрытие $\CC\to\CC^*$). Полосы $\Pi_k=\{2\pi k<\Im w<2\pi(k+1)\}$, $k\in\ZZ$, конформно отображаются
на $\Omega=\CC\setminus\RR_+$. Общая граница $\Pi_k$ и $\Pi_{k+1}$ --- прямая $w=s+2\pi i(k+1)$, $e^w=e^s>0$; подход из $\Pi_k$ (снизу):
$e^{s+2\pi i(k+1)-i\delta}=e^se^{-i\delta}$ --- нижний берег листа $k$; из $\Pi_{k+1}$ --- верхний берег листа $k+1$. Склейка: $x^-(k)\sim x^+(k+1)$, $x>0$,
$k\in\ZZ$ (на слайде --- в обратной нумерации, см.\ замечание выше). Получается бесконечнолистная «винтовая лестница» $X_L$; $\Ln$ --- однозначная
функция на ней.

\textbf{4. $\Zh^{-1}(z)$.} $\Zh(w)=\frac12(w+w^{-1})$. \emph{Накрытие:} $G=\CC\setminus\{0,1,-1\}$, $D=\CC\setminus\{\pm1\}$. Действительно, $\Zh'(w)=\frac12(1-w^{-2})\ne0$
на $G$; $\Zh(w)=\pm1\iff w=\pm1$, поэтому $\Zh(G)\subset D$. Для круга $V\subset D$ и голоморфной ветви $q$ корня $\sqrt{z^2-1}$ на $V$ все решения
$\Zh(w)=z'$ --- это $z'\pm q(z')$ (корни $w^2-2z'w+1=0$), они $\ne0$ (произведение $1$) и $\ne\pm1$; обращения $h_\pm=z\pm q$. По лемме 31.1 ---
двулистное накрытие.
\emph{Листы:} $\Zh$ конформно отображает верхнюю полуплоскость $G_1$ и нижнюю $G_2$ на $\Omega=\CC\setminus((-\infty,-1]\cup[1,+\infty))$
(инъективность: $\Zh(w_1)=\Zh(w_2)\iff w_2\in\{w_1,1/w_1\}$, а $1/w_1$ лежит в другой полуплоскости; образ: вещественная ось переходит в
$(-\infty,-1]\cup[1,\infty)$, полуокружность $|w|=1$ --- в $[-1,1]$).
\emph{Склейка.} Для $w=x+i\delta$ ($x\in\RR\setminus\{0\}$, $\delta\to0+$): $\Im\Zh(w)=\frac\delta2\bigl(1-\frac1{x^2+\delta^2}\bigr)$. При $|x|>1$ это $>0$: точка $x$, подходя из
$G_1$, даёт верхний берег (листа 1) луча, содержащего $\Zh(x)$, а из $G_2$ ($\delta\to0-$) --- нижний берег листа 2. При $0<|x|<1$ --- наоборот.
Каждая точка $y\in(1,\infty)$ (или $(-\infty,-1)$) имеет два прообраза $x$ и $1/x$, $|x|>1$; поэтому
$y^+(1)\sim y^-(2)$ и $y^-(1)\sim y^+(2)$ для $y\in(-\infty,-1)\cup(1,+\infty)$ --- как на слайде 416.

\begin{remarkbox}{Замечание (слайд 416)}
На слайде: «$\Zh$ локально конформно и взаимно однозначно отображает $\CC\setminus\{\pm1\}$ на $X_{\Zh}$». Точку $w=0$ надо исключить: $\Zh(0)=\infty$, а
$X_{\Zh}$ лежит над $\CC\setminus\{\pm1\}$. Правильно: $\Zh:\CC\setminus\{0,\pm1\}\to X_{\Zh}$ (или, добавив над $\infty$ две точки, $\Ch\setminus\{\pm1\}\to\Ch\setminus\{\pm1\}$).
\end{remarkbox}

\textbf{5. $\Arcsin z$.} $G=\CC\setminus(\frac\pi2+\pi\ZZ)$, $g=\sin$, $D=\CC\setminus\{\pm1\}$.
\emph{Накрытие:} $\sin'=\cos\ne0$ на $G$, $\sin w=\pm1\iff w\in\frac\pi2+\pi\ZZ$. Пусть $V\subset D$ --- круг, $q$ --- голоморфная ветвь $\sqrt{1-z^2}$ на $V$. Так как
$(iz+q)(-iz+q)=q^2+z^2=1$, функция $u_\pm=iz\pm q$ не обращается в нуль, и на $V$ есть ветви $\ell$ логарифма $u_\pm$. Тогда $h=-i\ell$ удовлетворяет
$\sin h=\frac{e^{ih}-e^{-ih}}{2i}=\frac{u_\pm-1/u_\pm}{2i}=\frac{(iz\pm q)-(-iz\pm q)}{2i}=z$. Все решения $\sin w=z'$: $\zeta=e^{iw}$ удовлетворяет
$\zeta^2-2iz'\zeta-1=0$, т.е.\ $\zeta=iz'\pm q(z')=u_\pm(z')$, а $w$ --- одно из значений $-i\log\zeta$; значит, $w=h(z')$ для подходящих знака и ветви $\ell$. По лемме 31.1
--- счётнолистное накрытие.
\emph{Листы:} полосы $\Pi_k=\{-\frac\pi2+\pi k<\Re w<\frac\pi2+\pi k\}$ конформно отображаются на $\Omega=\CC\setminus((-\infty,-1]\cup[1,\infty))$ (для $k=0$ --- известный
факт; $\sin(w+\pi k)=(-1)^k\sin w$, а $\Omega$ симметрична относительно $z\mapsto-z$).
\emph{Соответствие границ.} $\sin(a+ib)=\sin a\ch b+i\cos a\sh b$. Для $k=0$: при $a=-\frac\pi2+\delta$ ($\delta\to0+$) $\Im\sin=\sin\delta\,\sh b$, $\Re\to-\ch b$: луч
$-\frac\pi2+i\RR_+$ ($b>0$) переходит в верхний берег $(-\infty,-1)$, луч $-\frac\pi2-i\RR_+$ --- в нижний; при $a=\frac\pi2-\delta$: $\Im\sin=\sin\delta\,\sh b$, $\Re\to\ch b$: луч
$\frac\pi2+i\RR_+$ --- в верхний берег $(1,\infty)$, $\frac\pi2-i\RR_+$ --- в нижний. Для чётных $k$ так же; \emph{для нечётных $k$ из-за множителя $-1$}: луч
$-\frac\pi2+\pi k+i\RR_+$ переходит в нижний берег $(1,+\infty)$, $-\frac\pi2+\pi k-i\RR_+$ --- в верхний берег $(1,\infty)$, $\frac\pi2+\pi k\pm i\RR_+$ --- в нижний/верхний берег
$(-\infty,-1)$.
\emph{Склейка.} Общая граница $\Pi_{2k-1}$ и $\Pi_{2k}$ --- прямая $\Re w=-\frac\pi2+2\pi k$, образ $(-\infty,-1]$; со стороны $\Pi_{2k}$ точка с $b>0$ даёт верхний берег,
со стороны $\Pi_{2k-1}$ ($a=-\frac\pi2-\delta$, $\cos a=-\sin\delta$) --- нижний. Общая граница $\Pi_{2k}$ и $\Pi_{2k+1}$ --- $\Re w=\frac\pi2+2\pi k$, образ $[1,\infty)$; со стороны
$\Pi_{2k}$ при $b>0$ --- верхний берег, со стороны $\Pi_{2k+1}$ --- нижний. Итог (как на слайде 417):
\begin{align*}
x<-1&:\quad x^+(2k)\sim x^-(2k-1),\quad x^-(2k)\sim x^+(2k-1);\\
x>1&:\quad x^+(2k)\sim x^-(2k+1),\quad x^-(2k)\sim x^+(2k+1),\qquad k\in\ZZ.
\end{align*}
$\Arcsin$ --- однозначная функция на этой поверхности.

\begin{remarkbox}{Замечание (слайд 416)}
Соответствие границ на слайде 416 («луч $-\pi/2+\pi k+i\RR_+$ переходит в верхний берег луча $(-\infty,-1)$» и т.д.) верно лишь для \emph{чётных} $k$; для нечётных
$k$ лучи на левой границе полосы переходят в берега $(1,+\infty)$, на правой --- в берега $(-\infty,-1)$, с переменой верхнего и нижнего (см.\ выше).
Итоговая схема склейки на слайде 417 при этом верна.
\end{remarkbox}

\textbf{Общий взгляд (слайды 417--418).} Будем «склеивать» соседние элементы $F_1=(B_1,f_1)$, $F_2=(B_2,f_2)$ П.А.Ф.\ $\Phi$ ($B_1\cap B_2\ne\varnothing$): если $F_1$ ---
НАП $F_2$, они представляют один и тот же лист и склеиваются; если нет --- представляют разные листы. Над окрестностью точки получается столько
листов, сколько у $\Phi$ различных элементов над ней. Формализуем это.

\subsection{Вспомогательные факты о канонических элементах}

\begin{thmbox}{Лемма 32.1}
(а) Если $G=(B(b,r),g)$ --- канонический элемент и $c\in B(b,r)$, то $E_c(g)$ имеет радиус $\ge r-|c-b|$ и является НАП элемента $G$; если $G\in\Phi$, то
$E_c(g)\in\Phi$.\\
(б) Если канонические элементы $F_1$, $F_2$ имеют общий центр и $F_1$ --- НАП $F_2$, то $F_1=F_2$.\\
(в) Если $G_1,G_2$ имеют общий центр и оба являются НАП элемента $F$, то $G_1=G_2$.\\
(г) Если $F=(B(a,r),f)$ и $0<\eps\le r$, то $W_\eps(F)=\{E_c(f):c\in B(a,\eps)\}$, и $\pi:W_\eps(F)\to B(a,\eps)$ --- биекция с обратным $c\mapsto E_c(f)$.
\end{thmbox}
\begin{proof}
(а) $g$ голоморфна в $B(c,r-|c-b|)\subset B(b,r)$, поэтому ряд Тейлора $g$ в точке $c$ сходится в этом круге к $g$ (теорема Коши--Тейлора). Функции $E_c(g)$
и $G$ совпадают с $g$ около $c$, т.е.\ на непустом открытом подмножестве выпуклого пересечения кругов, значит, на всём пересечении. Если $G\in\Phi$, то $E_c(g)$
получается из $G$ продолжением вдоль отрезка $[b,c]$ (семейство $E_{\gamma(t)}(g)$ удовлетворяет Опр.~15.4 по только что доказанному), т.е.\ $E_c(g)\in\Phi$.

(б) Круги концентричны, пересечение --- меньший из них; функции совпадают около центра, значит, имеют один ряд Тейлора, а канонический элемент
определяется рядом Тейлора.

(в) Пусть $r_{G_1}\ge r_{G_2}$; тогда $B_{G_2}\subset B_{G_1}$ и $\varnothing\ne B_{G_2}\cap B_F\subset B_{G_1}\cap B_F$; на $B_{G_2}\cap B_F$ имеем $g_1=f=g_2$. Значит, $G_1$ --- НАП $G_2$, и по
(б) $G_1=G_2$.

(г) Если $c\in B(a,\eps)\subset B(a,r)$, то $E_c(f)$ --- НАП $F$ по (а) и $|c-a|<\eps$, т.е.\ $E_c(f)\in W_\eps(F)$. Обратно, если $G=(B_G,g)\in W_\eps(F)$ с центром $c$, то
$c\in B(a,\eps)\subset B_F$ и $c\in B_G$, $g=f$ на $B_G\cap B_F\ni c$, т.е.\ около $c$; значит, $G=E_c(f)$. Биективность $\pi$ очевидна из этого описания.
\end{proof}

\subsection{Топология на $\mathcal R$}

\begin{thmbox}{Лемма 32.2 (окрестности задают базу топологии)}
Если $G\in W_\eps(F)$, то найдётся $\delta\in(0,r_G]$ с $W_\delta(G)\subset W_\eps(F)$. Следовательно, семейство $\{W_\eps(F):F\in\Phi,\eps>0\}$ --- база топологии на
$\mathcal R$, причём подсемейство с $\eps\le r_F$ --- база той же топологии.
\end{thmbox}
\begin{proof}
Пусть $F=(B(a,r),f)$, $G=(B(c,\rho),g)$, $|c-a|<\eps$, и $g=f$ на непустом $B(c,\rho)\cap B(a,r)$. Возьмём точку $p$ в этом пересечении и $\delta>0$ так, чтобы
$\delta\le\eps-|c-a|$, $\delta\le\rho$ и $2\delta<\rho-|p-c|$. По лемме 32.1(г) любой $H\in W_\delta(G)$ имеет вид $E_{c'}(g)$, $|c'-c|<\delta$. Тогда
$|c'-a|<\delta+|c-a|\le\eps$; радиус $H$ не меньше $\rho-|c'-c|>\rho-\delta$ (лемма 32.1(а)), а $|p-c'|<|p-c|+\delta<\rho-\delta$, т.е.\ $p\in B_H$. Итак, $p$ лежит в
открытом множестве $B_H\cap B(c,\rho)\cap B(a,r)$, на котором (по лемме 32.1(а) и условию) функция $H$ равна $g$, а $g=f$. Значит, $H$ и $F$ совпадают на непустом
открытом подмножестве $B_H\cap B_F$ и $H$ --- НАП $F$: $H\in W_\eps(F)$.

\emph{База.} (1) $F\in W_\eps(F)$ ($F$ --- НАП самого себя), так что множества покрывают $\mathcal R$. (2) Если $G\in W_{\eps_1}(F_1)\cap W_{\eps_2}(F_2)$, то по доказанному
$W_{\delta_i}(G)\subset W_{\eps_i}(F_i)$, и $W_{\min\delta_i}(G)$ лежит в пересечении (очевидно, $W_\eps(G)$ возрастает по $\eps$). Это и есть аксиомы базы.
Так как в доказательстве $\delta\le r_G$, любое $W_\eps(F)$ --- объединение множеств $W_\delta(G)$ с $\delta\le r_G$.
\end{proof}

\subsection{Предложение 15.9}

\begin{remarkbox}{Замечание (уточнение атласа слайда 424)}
На слайде атлас: $(F,\eps)$ с $B(\pi(F),\eps)\subset D_\Phi$. Этого мало: если $\eps>r_F$, то $\pi(W_\eps(F))$ может не совпадать с кругом $B(\pi(F),\eps)$ (над частью круга у $F$ нет
НАП). Пример: $\Phi=\sqrt{1+\sqrt z}$ (билет 30), $D_\Phi=\CC^*$, $F$ --- элемент типа $\sqrt{1-\sqrt z}$ с центром $a=1+0{,}1i$, $r_F=0{,}1$; $B(a,0{,}5)\subset\CC^*$, но над
$c=1-0{,}1i$ нет НАП $F$: элементы типа $\pm\sqrt{1-\sqrt z}$ с центром $c$ имеют радиус $0{,}1$, и их круги не пересекают $B(a,0{,}1)$, а элементы
$\pm\sqrt{1+\sqrt z}$ имеют значения около $\pm\sqrt2$, а не около $0$. Поэтому берём $\mathcal A=\{(F,\eps):F\in\Phi,\ 0<\eps\le r_F\}$ (тогда автоматически
$B(\pi(F),\eps)\subset B_F\subset D_\Phi$).
\end{remarkbox}

\begin{thmbox}{Предложение 15.9 (в уточнённой форме)}
Пусть $\Phi$ --- П.А.Ф., $\mathcal R=\mathcal R(\Phi)$ с топологией из леммы 32.2, $U_\alpha=W_\eps(F)$, $\varphi_\alpha=\pi|_{U_\alpha}$ для $\alpha=(F,\eps)\in\mathcal A$.
\begin{enumerate}
\item[(1)] $\mathcal R$ хаусдорфово и линейно связно; $(\mathcal R,\{U_\alpha\},\{\varphi_\alpha\})$ --- одномерное комплексное многообразие (\emph{риманова поверхность $\Phi$}).
\item[(2)] $\pi:\mathcal R\to D_\Phi$ голоморфно, сюръективно и локально биголоморфно.
\item[(3)] Если $D'\subset D_\Phi$ --- область и $\Psi\subset\Phi$ --- аналитическая ветвь $\Phi$ над $D'$, то $\Psi$ --- открытое линейно связное подмножество $\mathcal R$ и
  $\pi:\Psi\to D'$ --- неразветвлённое голоморфное накрытие. В частности, если каждый элемент $\Phi$ продолжается вдоль любого пути в $D_\Phi$ (как для $\sqrt[n]z$,
  $\Ln z$, $\Zh^{-1}$, $\Arcsin$), то $\pi:\mathcal R\to D_\Phi$ --- неразветвлённое голоморфное накрытие.
\item[(4)] $f_\Phi$ голоморфна на $\mathcal R$; каждый элемент $F\in\Phi$ равен $E_{\pi(F)}(f_\Phi\circ s)$, где $s$ --- локальное обращение $\pi$ с $s(\pi(F))=F$.
  Если $\pi$ --- накрытие (как в (3)), то $\Phi=f_\Phi\circ\pi^{-1}$ в смысле предл.~31.2.
\end{enumerate}
\end{thmbox}
\begin{proof}
\emph{(1) Хаусдорфовость} (слайд 425). Пусть $F_1=(B_1,h_1)\ne F_2=(B_2,h_2)$. Если $\pi(F_1)\ne\pi(F_2)$, возьмём $\eps<|\pi(F_1)-\pi(F_2)|/2$: проекции
$W_\eps(F_i)$ лежат в непересекающихся кругах $B(\pi(F_i),\eps)$, значит $W_\eps(F_1)\cap W_\eps(F_2)=\varnothing$. Пусть $\pi(F_1)=\pi(F_2)=a$; выберем $\eps$ с
$B(a,\eps)\subset B_1\cap B_2$. Если $F=(B,h)$ с центром $b$ лежит в обоих $W_\eps(F_i)$, то $b\in B(a,\eps)\subset B_1\cap B_2$ и $b\in B$; около $b$ (в открытом
$B\cap B_1\cap B_2$) $h=h_1$ и $h=h_2$, т.е.\ $h_1=h_2$ на непустом открытом подмножестве $B_1\cap B_2$, и $F_1$ --- НАП $F_2$; по лемме 32.1(б) $F_1=F_2$ ---
противоречие.

\emph{Карты.} Пусть $\alpha=(F,\eps)$, $F=(B(a,r),f)$, $\eps\le r$. По лемме 32.1(г) $\varphi_\alpha:U_\alpha\to B(a,\eps)$ --- биекция. \emph{Непрерывность:}
$\pi(W_\delta(G))\subset B(\pi(G),\delta)$, т.е.\ $\pi$ непрерывна на всём $\mathcal R$. \emph{Открытость:} любое открытое подмножество $U_\alpha$ --- объединение базисных
$W_\delta(G)\subset U_\alpha$ с $\delta\le r_G$ (лемма 32.2), а $\pi(W_\delta(G))=B(\pi(G),\delta)$ (лемма 32.1(г)) --- открытый круг. Значит, $\varphi_\alpha$ --- гомеоморфизм на круг.
\emph{Замены карт:} если $c\in\varphi_\alpha(U_\alpha\cap U_\beta)$, то $\varphi_\alpha^{-1}(c)$ --- элемент с центром $c$, лежащий в $U_\beta$, и $\varphi_\beta(\varphi_\alpha^{-1}(c))=c$. Итак,
$\varphi_\beta\circ\varphi_\alpha^{-1}=\mathrm{id}$ на $\varphi_\alpha(U_\alpha\cap U_\beta)$ --- голоморфно. Множества $U_\alpha$ покрывают $\mathcal R$ ($F\in W_{r_F}(F)$; при $r_F=\infty$
берём любое $\eps$).

\emph{Линейная связность.} Зафиксируем $F_0\in\Phi$. Любой $F\in\Phi$ --- результат продолжения $\{F_t\}$ элемента $F_0$ вдоль пути $\gamma$. Отображение $t\mapsto F_t$
непрерывно в $\mathcal R$: по Опр.~15.4 для $\xi$ из окрестности $t$ элемент $F_\xi$ --- НАП $F_t$ и $|\gamma(\xi)-\gamma(t)|<\eps$ при $\xi$, близких к $t$, т.е.\ $F_\xi\in W_\eps(F_t)$
(а $W_\eps(F_t)$, $\eps>0$, образуют базу окрестностей точки $F_t$ по лемме 32.2). Итак, $t\mapsto F_t$ --- путь в $\mathcal R$ из $F_0$ в $F$.

\emph{(2)} В картах $\pi\circ\varphi_\alpha^{-1}=\mathrm{id}$, т.е.\ $\pi$ голоморфно, а $\pi|_{U_\alpha}=\varphi_\alpha$ --- биголоморфизм $U_\alpha$ на круг. Сюръективность: если
$z\in D_\Phi$, то $z\in B_F$ для некоторого $F\in\Phi$, и $E_z(f)\in\Phi$ по лемме 32.1(а).

\emph{(3) Накрытие.} \emph{Шаг 1: элементы ветви продолжаются по путям в $D'$.} Элементы $\Psi$ --- продолжения $F_b$ вдоль путей в $D'$; продолжение такого
элемента вдоль пути в $D'$ --- это продолжение $F_b$ вдоль составного пути, оно существует по определению ветви и лежит в $\Psi$.
\emph{Шаг 2: радиусы.} Если $F\in\Psi$ имеет центр $c$, то $r_F\ge d(c):=\dist(c,\CC\setminus D')$: круг $B(c,d(c))\subset D'$ односвязен, и по теореме о монодромии $F$
задаётся функцией, голоморфной в нём.
\emph{Шаг 3: открытость и связность.} Если $F\in\Psi$ и $\eps\le\min(r_F,d(\pi(F)))$, то $W_\eps(F)=\{E_c(f):c\in B(\pi(F),\eps)\}\subset\Psi$ (продолжения $F$ вдоль отрезков в
$D'$); значит, $\Psi$ открыто. Линейная связность --- как в (1), пути $t\mapsto F_t$ для продолжений вдоль путей в $D'$ лежат в $\Psi$. Сюръективность на $D'$:
$F_b$ продолжается в любую точку $D'$.
\emph{Шаг 4: правильно накрываемые окрестности} («упражнение» слайда 426). Пусть $z_0\in D'$, $\eps=d(z_0)/2$ ($\eps=1$, если $D'=\CC$), $V=B(z_0,\eps)$, $\{F_j\}_{j\in J}$
--- все элементы $\Psi$ с центром $z_0$, $F_j=(B_j,f_j)$; по шагу 2 $r_{F_j}\ge2\eps$. Положим $U_j=W_\eps(F_j)=\{E_c(f_j):c\in V\}$ (лемма 32.1(г)); $U_j\subset\Psi$ открыты.
\begin{itemize}
\item \emph{$\pi|_{U_j}:U_j\to V$ биголоморфно}: это карта $\varphi_{(F_j,\eps)}$, а карта --- биголоморфизм на свой круг (в картах она тождественна).
\item \emph{Дизъюнктность:} если $E_c(f_j)=E_c(f_k)$, то $f_j=f_k$ около $c$; обе функции голоморфны в $B(z_0,2\eps)$, значит $f_j\equiv f_k$ там, и $F_j=F_k$.
\item \emph{$\pi^{-1}(V)\cap\Psi=\bigcup_jU_j$:} пусть $F=(B_F,f)\in\Psi$, центр $c\in V$. По шагу 2 $r_F\ge d(c)\ge2\eps-|c-z_0|>|c-z_0|$, т.е.\ $z_0\in B_F$. По
  лемме 32.1(а) $E_{z_0}(f)\in\Psi$ (продолжение вдоль отрезка в $V\subset D'$), т.е.\ $E_{z_0}(f)=F_j$ для некоторого $j$: $f_j=f$ около $z_0$, значит и на выпуклом
  $B_F\cap B(z_0,2\eps)\ni c$; поэтому $F=E_c(f)=E_c(f_j)\in U_j$.
\end{itemize}
Все условия определения накрытия выполнены.

\emph{(4)} В карте $\varphi=\pi|_{W_\eps(F)}$, $F=(B,h)$, $\eps\le r_F$: $f_\Phi(\varphi^{-1}(c))=f_\Phi(E_c(h))=h(c)$, т.е.\ $f_\Phi\circ\varphi^{-1}=h|_{B(\pi(F),\eps)}$ --- голоморфна. Значит,
$f_\Phi$ голоморфна, и $E_{\pi(F)}(f_\Phi\circ\varphi^{-1})=E_{\pi(F)}(h)=F$. Если $\pi$ --- накрытие, то в обозначениях предл.~31.2 $s_F=\varphi^{-1}$ около $\pi(F)$ и $E(F)=F$
для всех $F$; по предл.~31.2 П.А.Ф.\ $f_\Phi\circ\pi^{-1}$ порождается любым $E(F)=F$, т.е.\ совпадает с $\Phi$.
\end{proof}

\begin{remarkbox}{Замечание (ошибка слайдов 425--426: $\pi$ не всегда накрытие над всей $D_\Phi$)}
На слайде утверждается, что $\pi:\mathcal R\to D_\Phi$ всегда неразветвлённое накрытие, и в доказательстве: «пусть $F=(B,h)\in\pi^{-1}(V)$. Тогда $z_0\in B$» ---
это неверно, если радиусы элементов над $V$ не ограничены снизу числом $|c-z_0|$. Контрпример: $\Phi=\sqrt{1+\sqrt z}$, $D_\Phi=\CC^*$. Над точкой $z=1$ лежат
только \emph{два} элемента ($\pm\sqrt{1+\sqrt z}$, значения $\pm\sqrt2$), а над близкими точками --- \emph{четыре} (билет 30). По следствию леммы 15.3 у накрытия все
слои равномощны, так что $\pi:\mathcal R(\Phi)\to\CC^*$ --- не накрытие (это локальный гомеоморфизм, но путь, ведущий к $1$ вдоль ветви $\sqrt{1-\sqrt z}$, не
поднимается). Верная версия --- пункт (3): накрытием является проекция любой аналитической ветви над областью. Например, для алгебраической П.А.Ф.\ теоремы
15.3 $\pi:\mathcal R(\Phi|_{\CC\setminus\mathcal S})\to\CC\setminus\mathcal S$ --- $N$-листное накрытие.
\end{remarkbox}

\subsection{Примеры как римановы поверхности в смысле формальной конструкции}

\begin{thmbox}{Предложение 32.6}
Пусть $G,D\subset\CC$ --- области, $g:G\to D$ --- неразветвлённое голоморфное накрытие, $w_0\in G$, и $\Phi$ --- П.А.Ф.\ $\iota\circ g^{-1}$ ($\iota(w)=w$), т.е.\ П.А.Ф.,
порождённая элементом $E(w_0)=E_{g(w_0)}(s_{w_0})$ локального обращения $g$. Пусть $D_\Phi=D$. Тогда $\Theta:G\to\mathcal R(\Phi)$, $\Theta(w)=E_{g(w)}(s_w)$, ---
послойный изоморфизм накрытий $g:G\to D$ и $\pi:\mathcal R(\Phi)\to D$; в частности, $\mathcal R(\Phi)$ биголоморфна $G$, и $f_\Phi\circ\Theta=\iota$.
Условие $D_\Phi=D$ выполнено для $\sqrt[n]z$, $\Ln z$ ($D=\CC^*$), $\Zh^{-1}$, $\Arcsin$ ($D=\CC\setminus\{\pm1\}$).
\end{thmbox}
\begin{proof}
Каждый элемент $\Phi$ --- продолжение $E(w_0)$ вдоль пути, лежащего в $D_\Phi=D$; по предл.~31.2(б) такие продолжения --- это в точности $\Theta(w)$, $w\in G$; так что $\Theta$ сюръективно и
$\pi\circ\Theta=g$. Инъективность: функция элемента $\Theta(w)$ в центре $g(w)$ равна $s_w(g(w))=w$, т.е.\ $f_\Phi(\Theta(w))=w$ --- по элементу восстанавливается $w$.
Локально: пусть $V$ --- правильно накрываемый круг с центром $g(w_1)$, $U$ --- лист через $w_1$, $s=(g|_U)^{-1}$, $\eps=$ радиус $V$ ($\le r_{\Theta(w_1)}$). Для $w\in U$:
$\Theta(w)=E_{g(w)}(s)$ (сечения совпадают около $g(w)$), т.е.\ $\Theta(U)=W_\eps(\Theta(w_1))$ по лемме 32.1(г), и в карте $\varphi=\pi|_{W_\eps}$:
$\varphi\circ\Theta|_U=g|_U$ --- биголоморфизм $U\to V$. Значит, $\Theta$ --- биективное локально биголоморфное отображение, т.е.\ биголоморфизм, и $\pi\circ\Theta=g$.

\emph{Проверка $D_\Phi=D$.} Иначе $\Phi$ имела бы элемент $(B(t,\rho),h)$ с центром в точке $t\notin D$, т.е.\ (по построению $\Phi$) голоморфное решение $g(h(z))=z$
около $t$: для $\sqrt[n]z$: $h^n=z$ около $0$ --- невозможно (левая часть имеет нуль кратности, делящейся на $n$, или не имеет нуля, правая --- простой нуль);
для $\Ln$: $e^h=z$ около $0$ --- невозможно ($e^h\ne0$); для $\Zh^{-1}$: $h^2-2zh+1=0$ около $\pm1$ $\Rightarrow$ $(h-z)^2=z^2-1$ --- квадрат голоморфной функции с простым нулём,
невозможно; для $\Arcsin$: $\sin h=z$ около $\pm1$ $\Rightarrow$ $\cos^2h=1-z^2$ --- то же противоречие.
\end{proof}

Таким образом, описанные «склейки листов» --- это (с точностью до послойного изоморфизма) римановы поверхности $\mathcal R(\Phi)$ П.А.Ф.\ $\sqrt[n]z$, $\Ln z$,
$\Zh^{-1}(z)$, $\Arcsin z$, а П.А.Ф.\ становится на них однозначной голоморфной функцией $f_\Phi$ (в модели $G$ --- просто $w$).

\begin{remarkbox}{Коротко: план ответа на экзамене}
\begin{enumerate}
\item Идея: П.А.Ф.\ --- однозначная функция на поверхности, многолистно лежащей над областью; листов над точкой --- столько, сколько элементов.
\item Примеры: $\sqrt z$ (две полуплоскости $\to$ два листа $\CC\setminus\RR_+$, склейка $x^\pm(1)\sim x^\mp(2)$), $\sqrt[n]z$ (циклическая склейка $n$ листов), $\Ln z$
  (бесконечная лестница), $\Zh^{-1}$ (два листа $\CC\setminus((-\infty,-1]\cup[1,\infty))$, склейка крест-накрест по обоим лучам), $\Arcsin$ (счётно много листов, чередование лучей).
\item Формально: точки $\mathcal R$ --- элементы $\Phi$; $\pi(F)$ --- центр; $W_\eps(F)$ --- элементы с центром ближе $\eps$, являющиеся НАП $F$; они образуют базу.
\item Хаусдорфовость (разные центры --- разнести; общий центр --- иначе $F_1$ НАП $F_2$).
\item Карты $\pi|_{W_\eps(F)}$, $\eps\le r_F$; замены карт тождественны.
\item $\mathcal R$ линейно связно (путь продолжения --- путь в $\mathcal R$); $\pi$ локально биголоморфно; над ветвью --- накрытие (радиусы $\ge\dist$ до границы).
\item $f_\Phi(F)=$ значение в центре --- голоморфна; $\Phi=f_\Phi\circ\pi^{-1}$.
\end{enumerate}
\end{remarkbox}

% ===== part 7 =====
% ===================== Часть 7: билеты 33--37 (целые функции) =====================
\ticket{33}{Порядок и тип целой функции. Выражение порядка (с доказательством) и типа через коэффициенты Тейлора}

\begin{defs}
\begin{itemize}
\item \textbf{Целая функция} --- функция $f\in H(\CC)$, т.е. голоморфная во всей плоскости. Она
  представляется всюду сходящимся рядом Тейлора $f(z)=\sum_{n=0}^\infty c_nz^n$,
  $c_n=f^{(n)}(0)/n!$. По формуле Коши--Адамара радиус сходимости $R$ ряда $\sum c_nz^n$ равен
  $1/\limsup_{n\to\infty}|c_n|^{1/n}$; для целой функции $R=\infty$, т.е.
  $\lim_{n\to\infty}|c_n|^{1/n}=0$. Обратно, если $|c_n|^{1/n}\to0$, то сумма ряда --- целая функция.
\item \textbf{Максимум модуля}: $M(r)=M_f(r)=\max_{|z|=r}|f(z)|$, $r>0$. По принципу максимума
  $M_f(r)=\max_{|z|\le r}|f(z)|$, поэтому $M_f$ не убывает. Если $f\ne\mathrm{const}$, то
  $M_f(r)\to\infty$ при $r\to\infty$ (теорема Лиувилля).
\item Запись $A\lesssim B$ означает: существует константа $C$, что $A\le CB$ при всех допустимых значениях
  параметров.
\item \textbf{Порядок} (Определение 16.1). Целая $f$ --- функция \emph{конечного порядка}, если существуют
  $\alpha>0$ и $R>0$ такие, что
  \begin{equation}\label{t33:ord}
  M_f(r)<\exp(r^\alpha)\qquad\text{при всех } r>R. \tag{16.1}
  \end{equation}
  Нижняя грань $\rho=\rho_f=\rho(f)$ таких $\alpha$ называется \emph{порядком} $f$. Если
  \eqref{t33:ord} не выполнено ни при каком $\alpha$, то $f$ имеет \emph{бесконечный порядок}, $\rho=\infty$.
  \emph{Пример:} $e^z$ имеет порядок $1$, многочлен --- порядок $0$.
\item \textbf{Тип} (Определение 16.2). Пусть $0<\rho<\infty$ --- порядок $f$. Функция $f$ имеет
  \emph{конечный тип} относительно порядка $\rho$, если найдутся $s>0$, $R_s>0$ с
  \begin{equation}\label{t33:type}
  M_f(r)<\exp(sr^\rho)\qquad\text{при всех } r>R_s. \tag{16.2}
  \end{equation}
  Нижняя грань $\sigma=\sigma_f$ таких $s$ --- \emph{тип} $f$. Если \eqref{t33:type} не выполнено ни при каком $s$,
  то $\sigma=\infty$. При $\sigma=0$ тип называют \emph{минимальным} (нулевым), при $0<\sigma<\infty$ ---
  \emph{нормальным}, при $\sigma=\infty$ --- \emph{максимальным} (бесконечным).
  \emph{Пример:} $e^{5z}$ имеет порядок $1$ и тип $5$.
\item \textbf{Верхний предел}: $\limsup_{n\to\infty}x_n=\lim_{n\to\infty}\sup_{k\ge n}x_k\in[-\infty,+\infty]$.
  Удобная характеристика: $\limsup x_n\le L$ $\iff$ для любого $L'>L$ неравенство $x_n<L'$ верно при
  всех достаточно больших $n$; $\limsup x_n\ge L$ $\iff$ для любого $L'<L$ неравенство $x_n>L'$ верно
  для бесконечно многих $n$. То же для $\limsup_{r\to\infty}$.
\item \textbf{Соглашение} для формулы порядка через коэффициенты: если $c_n=0$, то считаем
  $\dfrac{n\ln n}{\ln(1/|c_n|)}=0$ (это согласуется с $\ln(1/0)=+\infty$). Так как $|c_n|^{1/n}\to0$, при
  больших $n$ имеем $|c_n|<1$, т.е. $\ln(1/|c_n|)>0$, и выражение корректно определено.
\end{itemize}
\end{defs}

\subsection{Предварительные факты: неравенства Коши и рост многочленов}

\begin{thmbox}{Неравенства Коши}
Если $f(z)=\sum_{n\ge0}c_nz^n$ --- целая функция, то $|c_n|\le M_f(r)/r^n$ для всех $r>0$, $n\in\ZZ_+$.
\end{thmbox}
\begin{proof}
По формуле Коши для коэффициентов $c_n=\frac1{2\pi i}\oint_{|\zeta|=r}\frac{f(\zeta)}{\zeta^{n+1}}\,d\zeta$,
откуда $|c_n|\le\frac1{2\pi}\cdot2\pi r\cdot\frac{M_f(r)}{r^{n+1}}=\frac{M_f(r)}{r^n}$.
\end{proof}

\begin{thmbox}{Утверждение 33.1 (рост многочленов и трансцендентных функций)}
\begin{enumerate}
\item[(а)] Если $|f(z)|\lesssim |z|^k$ при $|z|\ge r_0$ ($k\in\ZZ_+$), то $f$ --- многочлен степени $\le k$.
\item[(б)] Если целая $f$ не многочлен, то $M_f(r)/r^m\to\infty$ при $r\to\infty$ для любого $m\ge 0$.
\end{enumerate}
\end{thmbox}
\begin{proof}
(а) По неравенствам Коши $|c_n|\le Cr^k/r^n=Cr^{k-n}$ при $r\ge r_0$; при $n>k$, устремляя $r\to\infty$,
получаем $c_n=0$.

(б) Если $f$ не многочлен, то $c_n\ne0$ для некоторого $n>m$. По неравенствам Коши
$M_f(r)\ge|c_n|r^n$, поэтому $M_f(r)/r^m\ge|c_n|r^{n-m}\to\infty$.
\end{proof}

Таким образом, у целой функции, не являющейся многочленом, $M_f(r)$ растёт быстрее любой степени $r$,
и для измерения роста естественно сравнивать $M_f(r)$ с функциями $\exp(r^\alpha)$.

\subsection{Порядок и тип через $M_f(r)$}

\begin{thmbox}{Утверждение 33.2 (формулы для порядка и типа через $M_f$)}
Пусть $f$ --- непостоянная целая функция. Тогда
\[
\rho_f=\limsup_{r\to\infty}\frac{\ln\ln M_f(r)}{\ln r},
\qquad\text{а при }0<\rho_f<\infty\qquad
\sigma_f=\limsup_{r\to\infty}\frac{\ln M_f(r)}{r^{\rho}}.
\]
В частности, если $0\le\rho<\infty$, то для любого $\eps>0$ найдётся $R_\eps$ с $M_f(r)<\exp(r^{\rho+\eps})$
при $r>R_\eps$, и найдётся последовательность $r_k\uparrow\infty$ с $M_f(r_k)>\exp(r_k^{\rho-\eps})$.
\end{thmbox}
\begin{proof}
Так как $M_f(r)\to\infty$, при больших $r$ величина $\ln\ln M_f(r)$ определена. Обозначим
$\mu=\limsup_{r\to\infty}\frac{\ln\ln M_f(r)}{\ln r}\in[-\infty,\infty]$.

\emph{$\mu\le\rho$.} Если $\rho=\infty$, доказывать нечего. Пусть $\alpha$ удовлетворяет \eqref{t33:ord}. Тогда
$\ln\ln M_f(r)<\alpha\ln r$ при $r>R$, т.е. $\frac{\ln\ln M_f(r)}{\ln r}<\alpha$ (при $r>\max(R,1)$), значит
$\mu\le\alpha$. Переходя к нижней грани по $\alpha$, получаем $\mu\le\rho$.

\emph{$\rho\le\mu$.} Если $\mu=\infty$, доказывать нечего. Пусть $\alpha>\max(\mu,0)$. По характеристике верхнего
предела $\frac{\ln\ln M_f(r)}{\ln r}<\alpha$ при всех больших $r$, т.е. $M_f(r)<\exp(r^\alpha)$; значит $\alpha$
удовлетворяет \eqref{t33:ord} и $\rho\le\alpha$. Так как $\alpha>\max(\mu,0)$ произвольно, $\rho\le\max(\mu,0)$.
Наконец, $\mu\ge0$: иначе $\ln\ln M_f(r)<0$, т.е. $M_f(r)<e$ при больших $r$, и $f$ ограничена, что
противоречит теореме Лиувилля. Итак, $\rho=\mu$.

Формула для типа доказывается дословно так же с заменой $\exp(r^\alpha)$ на $\exp(sr^\rho)$ и величины
$\frac{\ln\ln M}{\ln r}$ на $\frac{\ln M}{r^\rho}$ (здесь $\limsup\ge0$, так как $M_f(r)\ge1$ при больших $r$).
Последнее утверждение --- просто расшифровка равенства $\rho=\mu$ по характеристике $\limsup$.
\end{proof}

\subsection{Примеры вычисления порядка и типа по определению}

\begin{remarkbox}{Примеры}
\begin{enumerate}
\item $f=e^z$: $|e^z|=e^{\Re z}\le e^r$, равенство при $z=r$; $M(r)=e^r$, $\rho=1$, $\sigma=1$.
\item $f=e^{z^2}$: $|e^{z^2}|=e^{\Re z^2}\le e^{r^2}$, равенство при $z=r$; $M(r)=e^{r^2}$, $\rho=2$, $\sigma=1$.
\item $f=e^{P(z)}$, $P(z)=p_0z^m+p_1z^{m-1}+\dots+p_m$, $p_0\ne0$, $m\ge1$: $\rho=m$, $\sigma=|p_0|$.
  Действительно, пусть $K=|p_1|+\dots+|p_m|$. При $|z|=r\ge1$:
  $\Re P(z)\le|P(z)|\le|p_0|r^m+Kr^{m-1}$, значит $\ln M(r)\le|p_0|r^m+Kr^{m-1}$. С другой стороны, на луче
  $\arg z=-\frac{\arg p_0}{m}$ имеем $p_0z^m=|p_0|r^m>0$, и $\Re P(z)\ge|p_0|r^m-Kr^{m-1}$. Итак,
  $\ln M(r)=|p_0|r^m+O(r^{m-1})$, откуда по Утверждению 33.2 $\rho=m$ и $\sigma=|p_0|$.
\item $f=\sin z$: $\rho=1$, $\sigma=1$. Действительно,
  $|\sin z|=\frac{|e^{iz}-e^{-iz}|}{2}\le\frac{e^{-\Im z}+e^{\Im z}}2=\operatorname{ch}(\Im z)\le\operatorname{ch}r$
  при $|z|=r$, и $|\sin(ir)|=\operatorname{sh}r$. Значит $\operatorname{sh}r\le M(r)\le\operatorname{ch}r$, т.е.
  $\ln M(r)=r+O(1)$. Так же $\cos z$: $\rho=1$, $\sigma=1$.
\item $f(z)=\cos\sqrt z:=\sum_{n\ge0}\frac{(-1)^nz^n}{(2n)!}$ (корректно: $\cos$ чётен, для любого $w$ с
  $w^2=z$ имеем $\cos w=f(z)$). При $|z|=r$, $|w|=\sqrt r$: $|f(z)|=|\cos w|\le\operatorname{ch}(\Im w)\le e^{\sqrt r}$,
  а при $z=-r$, $w=i\sqrt r$: $f(-r)=\operatorname{ch}\sqrt r\ge\frac12e^{\sqrt r}$. Поэтому
  $\ln M(r)=\sqrt r+O(1)$: $\rho=\frac12$, $\sigma=1$. Аналогично $\frac{\sin\sqrt z}{\sqrt z}=\sum\frac{(-1)^nz^n}{(2n+1)!}$
  имеет $\rho=\frac12$, $\sigma=1$.
\item $f=\exp(e^z)$: $|f(z)|=e^{\Re e^z}\le e^{|e^z|}=e^{e^{\Re z}}\le e^{e^r}$, равенство при $z=r$.
  Значит $\frac{\ln\ln M(r)}{\ln r}=\frac{r}{\ln r}\to\infty$: \emph{бесконечный порядок}.
  Аналогично $\exp(\cos z)$: $M(r)\ge|\exp(\cos ir)|=e^{\operatorname{ch}r}\ge e^{e^r/2}$, $\ln\ln M(r)\ge r-\ln2$, $\rho=\infty$.
\item $1/\Gamma(z)$ (без доказательства): целая функция порядка $1$ \emph{максимального} типа; в действительности
  $\ln M_{1/\Gamma}(r)=r\ln r\,(1+o(1))$.
\end{enumerate}
\end{remarkbox}

\begin{thmbox}{Задача (со слайда 432)}
Если $P\not\equiv0$ --- многочлен, $f$ --- целая функция порядка $\rho\in[0,\infty]$, то $fP$ имеет порядок $\rho$;
при $0<\rho<\infty$ типы $f$ и $fP$ совпадают. Аналогично для $f+a$, $a\in\CC$ ($f\ne\mathrm{const}$).
\end{thmbox}
\begin{proof}
Пусть $P(z)=p_0z^d+\dots$, $p_0\ne0$. При больших $r$ на $|z|=r$ имеем
$\frac12|p_0|r^d\le|P(z)|\le2|p_0|r^d$, поэтому
$\frac12|p_0|r^dM_f(r)\le M_{fP}(r)\le2|p_0|r^dM_f(r)$, т.е.
$\ln M_{fP}(r)=\ln M_f(r)+O(\ln r)$. Если $f$ --- многочлен, то $fP$ --- многочлен, оба порядка равны $0$.
Иначе по Утверждению 33.1(б) $\ln M_f(r)/\ln r\to\infty$, так что
$\ln M_{fP}(r)=\ln M_f(r)\,(1+o(1))$ и $\ln\ln M_{fP}(r)=\ln\ln M_f(r)+o(1)$; по Утверждению 33.2 порядки равны.
При $0<\rho<\infty$: $\frac{\ln M_{fP}(r)}{r^\rho}=\frac{\ln M_f(r)}{r^\rho}+O\big(\frac{\ln r}{r^\rho}\big)$, и
$\frac{\ln r}{r^\rho}\to0$, так что типы равны. Для $f+a$: $M_f(r)-|a|\le M_{f+a}(r)\le M_f(r)+|a|$ и
$M_f(r)\to\infty$, поэтому $\ln M_{f+a}(r)=\ln M_f(r)+o(1)$.
\end{proof}

\subsection{Две леммы о связи роста и коэффициентов}

\begin{thmbox}{Лемма 16.1}
Пусть $f(z)=\sum_{n\ge0}c_nz^n$ --- целая функция и $M_f(r)<\exp(Ar^\alpha)$ при всех $r>R$, где $A>0$,
$\alpha>0$. Тогда найдётся $N$ такое, что при всех $n>N$
\begin{equation}\label{t33:163}
|c_n|^{1/n}<\Big(\frac{A\alpha e}{n}\Big)^{1/\alpha}.\tag{16.3}
\end{equation}
\end{thmbox}
\begin{proof}
По неравенствам Коши при $r>R$: $|c_n|\le M_f(r)r^{-n}<\varphi_n(r):=\exp(Ar^\alpha)\,r^{-n}$. Имеем
$(\ln\varphi_n)'(r)=A\alpha r^{\alpha-1}-\frac nr$, что обращается в нуль при
$r_n=\big(\frac n{A\alpha}\big)^{1/\alpha}$ (это точка минимума $\varphi_n$, но нам достаточно просто подставить её).
Так как $r_n\to\infty$, найдётся $N$ такое, что $r_n>R$ при $n>N$. Для таких $n$, подставляя $r=r_n$:
\[
\ln|c_n|<Ar_n^\alpha-n\ln r_n=\frac n\alpha-\frac n\alpha\ln\frac n{A\alpha}=\frac n\alpha\ln\frac{A\alpha e}{n},
\]
т.е. $|c_n|<\big(\frac{A\alpha e}n\big)^{n/\alpha}$, что равносильно \eqref{t33:163}.
\end{proof}

\begin{thmbox}{Лемма 16.2}
Пусть $f(z)=\sum_{n\ge0}c_nz^n$ --- степенной ряд, и при всех $n>N$ выполнено \eqref{t33:163} с некоторыми
$A>0$, $\alpha>0$. Тогда $f$ --- целая функция и для любого $\eps>0$ найдётся $R_\eps$ такое, что
\[
M_f(r)<\exp\big((A+\eps)r^\alpha\big)\qquad\text{при всех }r>R_\eps.
\]
\end{thmbox}
\begin{proof}
Из \eqref{t33:163} $|c_n|^{1/n}\to0$, поэтому ряд сходится во всей плоскости и $f$ --- целая функция.

Фиксируем $r>0$. При $n>N$ из \eqref{t33:163}
\[
|c_n|r^n<\Big(\frac{A\alpha e\,r^\alpha}{n}\Big)^{n/\alpha}=\psi_r(n),\qquad
\psi_r(t):=\exp\Big(\frac t\alpha\ln\frac{A\alpha e\,r^\alpha}{t}\Big),\ t>0.
\]
\begin{detail}{максимум $\psi_r$ (на слайде --- «непосредственное вычисление»)}
$\frac{d}{dt}\Big[\frac t\alpha\big(\ln(A\alpha e r^\alpha)-\ln t\big)\Big]=\frac1\alpha\big(\ln(A\alpha er^\alpha)-\ln t-1\big)
=\frac1\alpha\ln\frac{A\alpha r^\alpha}{t}$. Эта производная положительна при $t<t^*:=A\alpha r^\alpha$ и
отрицательна при $t>t^*$, так что $\max_{t>0}\psi_r(t)=\psi_r(t^*)=\exp\big(\frac{A\alpha r^\alpha}{\alpha}\ln e\big)=\exp(Ar^\alpha)$.
\end{detail}
Итак, (а) $|c_n|r^n<\exp(Ar^\alpha)$ при всех $n>N$.
(б) Пусть $T(r)=2^\alpha A\alpha e\,r^\alpha$. Если $n\ge T(r)$, то $\frac{A\alpha er^\alpha}{n}\le2^{-\alpha}$, и
при $n>N$ получаем $|c_n|r^n<(2^{-\alpha})^{n/\alpha}=2^{-n}$.

Разобьём ряд для $|f(z)|$, $|z|=r$, на три части: $n\le N$; $N<n<T(r)$ (таких $n$ меньше $T(r)$; если
$T(r)\le N+1$, их нет); $n>N$, $n\ge T(r)$. Получаем
\[
M_f(r)\le\sum_{n=0}^{N}|c_n|r^n+T(r)\,e^{Ar^\alpha}+\sum_{n\ge1}2^{-n}=Q(r)+2^\alpha A\alpha e\,r^\alpha e^{Ar^\alpha}+1,
\]
где $Q(r)=\sum_{n\le N}|c_n|r^n$ --- многочлен. Умножим правую часть на $e^{-(A+\eps)r^\alpha}$:
\[
\big(Q(r)+1\big)e^{-(A+\eps)r^\alpha}+2^\alpha A\alpha e\cdot r^\alpha e^{-\eps r^\alpha}\xrightarrow[r\to\infty]{}0,
\]
так как экспонента $e^{-\eps r^\alpha}$ убывает быстрее любой степени $r$. Значит при $r>R_\eps$ это выражение
меньше $1$, т.е. $M_f(r)<\exp((A+\eps)r^\alpha)$.
\end{proof}

\begin{remarkbox}{Замечание}
На слайдах 433--435 та же идея изложена через $\mu(r)=\max_n|c_n|r^n$ и неравенство
$M_f(r)\le N_1(r)\mu(r)+1$; приведённое выше рассуждение --- то же самое, но без необходимости следить за
номером $m(r)$ максимального члена.
\end{remarkbox}

\subsection{Теорема 16.1: порядок и тип через коэффициенты Тейлора}

\begin{thmbox}{Теорема 16.1}
Пусть $f(z)=\sum_{n\ge0}c_nz^n$ --- целая функция порядка $\rho$. Тогда
\begin{equation}\label{t33:164}
\rho=\limsup_{n\to\infty}\frac{n\ln n}{\ln(1/|c_n|)}.\tag{16.4}
\end{equation}
Если $0<\rho<\infty$ и $\sigma$ --- тип $f$ относительно порядка $\rho$, то
\begin{equation}\label{t33:165}
(e\sigma\rho)^{1/\rho}=\limsup_{n\to\infty}n^{1/\rho}|c_n|^{1/n},\qquad\text{т.е.}\qquad
\sigma=\frac1{e\rho}\limsup_{n\to\infty}n\,|c_n|^{\rho/n}.\tag{16.5}
\end{equation}
\end{thmbox}

На слайдах формула \eqref{t33:164} доказывается для $\rho<\infty$, а случай $\rho=\infty$ оставлен как упражнение;
мы докажем её сразу для любого $\rho\in[0,\infty]$. Обозначим $L=\limsup_{n\to\infty}\frac{n\ln n}{\ln(1/|c_n|)}\in[0,\infty]$
(с принятым соглашением при $c_n=0$).

\begin{proof}[Доказательство \eqref{t33:164}, шаг 1: $L\le\rho$]
Если $\rho=\infty$, доказывать нечего. Пусть $\rho<\infty$ и $\eps>0$. По Утверждению 33.2 найдётся $R_0$ с
$M_f(r)<\exp(r^{\rho+\eps})$ при $r>R_0$. Лемма 16.1 с $A=1$, $\alpha=\rho+\eps$ даёт $N_0$ такое, что при $n>N_0$
\[
|c_n|^{1/n}<\Big(\frac{e(\rho+\eps)}{n}\Big)^{1/(\rho+\eps)}.
\]
Если $c_n\ne0$, логарифмируем: $\frac1n\ln\frac1{|c_n|}>\frac{\ln n-\ln(e(\rho+\eps))}{\rho+\eps}$, т.е.
\[
\ln\frac1{|c_n|}>\frac{n\ln n}{\rho+\eps}-\lambda n,\qquad \lambda=\frac{\ln\big(e(\rho+\eps)\big)}{\rho+\eps}.
\]
При $n$ таких, что $\ln n>2|\lambda|(\rho+\eps)$, правая часть положительна, и
\[
\frac{n\ln n}{\ln(1/|c_n|)}<\frac{n\ln n}{\frac{n\ln n}{\rho+\eps}-\lambda n}=\frac{\rho+\eps}{1-\frac{\lambda(\rho+\eps)}{\ln n}}
\xrightarrow[n\to\infty]{}\rho+\eps .
\]
Поэтому найдётся $N_1$, что при $n>N_1$ выражение $\frac{n\ln n}{\ln(1/|c_n|)}$ меньше $\rho+2\eps$ (при $c_n=0$ оно равно $0$
по соглашению). Следовательно, $L\le\rho+2\eps$ для любого $\eps>0$, т.е. $L\le\rho$.
\end{proof}

\begin{proof}[Шаг 2: $\rho\le L$]
Если $L=\infty$, доказывать нечего. Пусть $L<\infty$ и $\eps>0$. Найдётся $N_2$ такое, что при $n>N_2$
$\frac{n\ln n}{\ln(1/|c_n|)}<L+\eps$. Если $c_n\ne0$ (и $n$ так велико, что $|c_n|<1$), это означает
$\ln\frac1{|c_n|}>\frac{n\ln n}{L+\eps}$, т.е. $|c_n|^{1/n}<\big(\frac1n\big)^{1/(L+\eps)}$; при $c_n=0$ это неравенство
тривиально. Оно имеет вид \eqref{t33:163} с
\[
\alpha=L+\eps>0,\qquad A=\frac1{e(L+\eps)}\quad(\text{тогда }A\alpha e=1).
\]
По Лемме 16.2 $M_f(r)<\exp\big((A+\eps)r^{L+\eps}\big)$ при $r>R_2$. Так как
$(A+\eps)r^{L+\eps}<r^{L+2\eps}$ при больших $r$, то $\alpha'=L+2\eps$ удовлетворяет \eqref{t33:ord} и
$\rho\le L+2\eps$. В силу произвольности $\eps$: $\rho\le L$.

Шаг 2 не использует конечность $\rho$; поэтому если $\rho=\infty$, то $L=\infty$ (иначе $\rho\le L<\infty$). Это и есть
«упражнение» со слайда 437. Вместе с шагом 1 получаем \eqref{t33:164} при всех $\rho\in[0,\infty]$.
\end{proof}

По программе формула \eqref{t33:165} для типа даётся без доказательства. Однако при наличии Лемм 16.1 и 16.2 она
доказывается в несколько строк --- приводим это доказательство как дополнительный материал.

\begin{proof}[Доказательство \eqref{t33:165} (доп. материал)]
Обозначим $\nu=\limsup_{n\to\infty}n^{1/\rho}|c_n|^{1/n}\in[0,\infty]$.

\emph{$\nu\le(e\sigma\rho)^{1/\rho}$.} При $\sigma=\infty$ доказывать нечего. Пусть $s>\sigma$; тогда
$M_f(r)<\exp(sr^\rho)$ при больших $r$, и по Лемме 16.1 (с $A=s$, $\alpha=\rho$) при больших $n$
$|c_n|^{1/n}<(s\rho e/n)^{1/\rho}$, т.е. $n^{1/\rho}|c_n|^{1/n}<(e\rho s)^{1/\rho}$. Значит $\nu\le(e\rho s)^{1/\rho}$ для
всех $s>\sigma$, откуда $\nu\le(e\rho\sigma)^{1/\rho}$.

\emph{$(e\sigma\rho)^{1/\rho}\le\nu$.} При $\nu=\infty$ доказывать нечего. Пусть $t>\nu$ ($t>0$). При больших $n$
$n^{1/\rho}|c_n|^{1/n}<t$, т.е. $|c_n|^{1/n}<\big(\frac{t^\rho}{n}\big)^{1/\rho}=\big(\frac{A\rho e}n\big)^{1/\rho}$ с
$A=\frac{t^\rho}{e\rho}$. По Лемме 16.2 (с $\alpha=\rho$) для любого $\eps>0$ $M_f(r)<\exp((A+\eps)r^\rho)$ при больших $r$,
значит $\sigma\le A+\eps$, и $\sigma\le A=\frac{t^\rho}{e\rho}$. Устремляя $t\downarrow\nu$, получаем $e\rho\sigma\le\nu^\rho$.
\end{proof}

\begin{thmbox}{Следствие (слайд 437)}
Целая функция $f$ и её производная $f'$ имеют одинаковый порядок, а при $0<\rho<\infty$ --- и одинаковый тип.
\end{thmbox}
\begin{proof}
На слайде предлагается применить \eqref{t33:164}, \eqref{t33:165} к коэффициентам $(n+1)c_{n+1}$ ряда $f'$; приведём
более прозрачное рассуждение через $M$. Если $f$ --- многочлен, утверждение очевидно (оба порядка равны $0$). Пусть $f$ не
многочлен; тогда и $f'$ не многочлен.

(1) По неравенству Коши для $f'(z)$ на окружности $|\zeta-z|=1$: $|f'(z)|\le\max_{|\zeta-z|=1}|f(\zeta)|\le M_f(|z|+1)$,
т.е. $M_{f'}(r)\le M_f(r+1)$. (2) $f(z)=f(0)+\int_{[0,z]}f'(\zeta)\,d\zeta$, поэтому $M_f(r)\le|f(0)|+rM_{f'}(r)$.

Из (1): $\frac{\ln\ln M_{f'}(r)}{\ln r}\le\frac{\ln\ln M_f(r+1)}{\ln(r+1)}\cdot\frac{\ln(r+1)}{\ln r}$, и так как
$\frac{\ln(r+1)}{\ln r}\to1$, имеем $\rho(f')\le\rho(f)$. Из (2) и Утверждения 33.1(б) (для $f'$: $M_{f'}(r)\ge 2r$ и
$rM_{f'}(r)\ge|f(0)|$ при больших $r$): $M_f(r)\le2rM_{f'}(r)\le M_{f'}(r)^2$, т.е.
$\ln\ln M_f(r)\le\ln2+\ln\ln M_{f'}(r)$, откуда $\rho(f)\le\rho(f')$.

Типы при общем порядке $\rho\in(0,\infty)$: из (1)
$\frac{\ln M_{f'}(r)}{r^\rho}\le\frac{\ln M_f(r+1)}{(r+1)^\rho}\big(\frac{r+1}{r}\big)^\rho$, значит $\sigma(f')\le\sigma(f)$;
из (2) $\ln M_f(r)\le\ln(2r)+\ln M_{f'}(r)$ и $\frac{\ln(2r)}{r^\rho}\to0$, значит $\sigma(f)\le\sigma(f')$.
\end{proof}

\subsection{Примеры вычисления по коэффициентам}

\begin{detail}{асимптотика $n!$: $\ (n!)^{1/n}\sim n/e$, т.е. $\frac{n}{(n!)^{1/n}}\to e$}
Так как $\ln$ возрастает, $\int_{k-1}^k\ln t\,dt\le\ln k\le\int_k^{k+1}\ln t\,dt$. Суммируя,
\[
n\ln n-n+1=\int_1^n\ln t\,dt\le\ln n!\le\int_1^{n+1}\ln t\,dt=(n+1)\ln(n+1)-n.
\]
Поэтому $\frac{\ln n!}{n}-\ln n+1$ заключено между $\frac1n$ и $\ln\frac{n+1}{n}+\frac{\ln(n+1)}{n}$, т.е. стремится к $0$.
Отсюда $\ln\frac{(n!)^{1/n}e}{n}\to0$. Кроме того, $\ln n!=n\ln n\,(1+o(1))$.
\end{detail}

\begin{remarkbox}{Примеры}
\begin{enumerate}
\item $e^z=\sum\frac{z^n}{n!}$: $\frac{n\ln n}{\ln n!}\to1$, так что $\rho=1$; $n|c_n|^{1/n}=\frac{n}{(n!)^{1/n}}\to e$, так что
  $\sigma=\frac{e}{e\cdot1}=1$. Совпадает с прямым вычислением.
\item $e^{z^2}=\sum\frac{z^{2k}}{k!}$: при $n=2k$ $\frac{n\ln n}{\ln(1/|c_n|)}=\frac{2k\ln(2k)}{\ln k!}\to2$ (при нечётных $n$
  $c_n=0$), $\rho=2$; $n|c_n|^{2/n}=\frac{2k}{(k!)^{1/k}}\to2e$, $\sigma=\frac{2e}{2e}=1$.
\item $\sin z$: $|c_{2k+1}|=\frac1{(2k+1)!}$, $c_{2k}=0$ --- как у $e^z$ по нечётным номерам: $\rho=1$, $\sigma=1$.
\item $\cos\sqrt z=\sum\frac{(-1)^nz^n}{(2n)!}$: $\frac{n\ln n}{\ln(2n)!}=\frac{n\ln n}{2n\ln(2n)(1+o(1))}\to\frac12$, $\rho=\frac12$;
  $n|c_n|^{1/(2n)}=\frac{n}{((2n)!)^{1/(2n)}}\to\frac e2$, и $\sigma=\frac{e/2}{e\cdot\frac12}=1$.
\item $f(z)=\sum_{n\ge0}\frac{z^n}{(n!)^{1/\rho}}$, $0<\rho<\infty$: $\ln\frac1{|c_n|}=\frac{\ln n!}{\rho}=\frac{n\ln n}{\rho}(1+o(1))$,
  значит порядок равен $\rho$; $n|c_n|^{\rho/n}=\frac{n}{(n!)^{1/n}}\to e$, тип $\sigma=\frac{e}{e\rho}=\frac1\rho$.
  (При $\rho=1$ это $e^z$.)
\item (слайд 438, а) $c_n=(A/n)^{n/\alpha}$: $\ln\frac1{c_n}=\frac n\alpha(\ln n-\ln A)$, и
  $\frac{n\ln n}{\ln(1/c_n)}\to\alpha$: порядок $\alpha$. (б) $c_n=(e\alpha\beta/n)^{n/\alpha}$: порядок $\alpha$ и
  $n^{1/\alpha}c_n^{1/n}=(e\alpha\beta)^{1/\alpha}$, т.е. по \eqref{t33:165} $e\sigma\alpha=e\alpha\beta$, тип $\sigma=\beta$.
\item (слайд 438, в) $|c_n|=n^{-n/\alpha_n}$, $\alpha_n>0$: $\frac{n\ln n}{\ln(1/|c_n|)}=\alpha_n$, так что
  (если ряд задаёт целую функцию) порядок равен $\limsup\alpha_n$. При $\alpha_n=1/n$ ($|c_n|=n^{-n^2}$) получаем порядок $0$;
  при $\alpha_n=(\ln n)^\gamma$, $\gamma\in(0,1)$, имеем $|c_n|^{1/n}=e^{-(\ln n)^{1-\gamma}}\to0$ (функция целая), а порядок
  бесконечен.
\item (слайд 439) $\sum_{n\ge2}\big(\frac{\ln n}{n}\big)^{n/\rho}z^n$: $\ln\frac1{|c_n|}=\frac n\rho(\ln n-\ln\ln n)$, порядок $\rho$;
  $n^{1/\rho}|c_n|^{1/n}=(\ln n)^{1/\rho}\to\infty$ --- максимальный тип. Ряд $\sum_{n\ge2}\big(\frac1{n\ln n}\big)^{n/\rho}z^n$:
  порядок $\rho$, $n^{1/\rho}|c_n|^{1/n}=(\ln n)^{-1/\rho}\to0$ --- минимальный тип.
\end{enumerate}
\end{remarkbox}

\begin{remarkbox}{Замечание к слайду 436}
На слайде $\lambda=\ln\big((e(\rho+\eps))\big/(\rho+\eps)\big)$ записано с двойными скобками; понимать нужно как
$\lambda=\ln\big(e(\rho+\eps)\big)/(\rho+\eps)$ (именно это получается при логарифмировании оценки Леммы 16.1).
Величина $\lambda$ может быть отрицательной (при $\rho+\eps<1/e$), что не мешает рассуждению.
\end{remarkbox}

\begin{remarkbox}{Коротко: план ответа на экзамене}
\begin{enumerate}
\item $M_f(r)$; неравенства Коши; рост $\lesssim r^k$ $\Rightarrow$ многочлен, поэтому сравниваем с $\exp(r^\alpha)$.
\item Определения порядка и типа; формулы $\rho=\limsup\frac{\ln\ln M}{\ln r}$, $\sigma=\limsup\frac{\ln M}{r^\rho}$.
\item Примеры: $e^{P}$ (порядок $\deg P$, тип $|p_0|$), $\sin z$, $\cos\sqrt z$, $\exp(e^z)$.
\item Лемма 16.1: $M<e^{Ar^\alpha}$ $\Rightarrow$ $|c_n|^{1/n}<(A\alpha e/n)^{1/\alpha}$ (Коши + минимизация по $r$).
\item Лемма 16.2: обратно (оценка $|c_n|r^n\le e^{Ar^\alpha}$ и хвост $\le\sum2^{-n}$ при $n\ge2^\alpha A\alpha er^\alpha$).
\item Теорема 16.1: $\rho=\limsup\frac{n\ln n}{\ln(1/|c_n|)}$: «$\le$» через Лемму 16.1, «$\ge$» через Лемму 16.2.
\item Тип: $(e\sigma\rho)^{1/\rho}=\limsup n^{1/\rho}|c_n|^{1/n}$ (формулировка; доказывается теми же леммами).
\item Следствие: $f$ и $f'$ одного порядка и типа; примеры $\sum z^n/(n!)^{1/\rho}$.
\end{enumerate}
\end{remarkbox}

% ======================================================================
\ticket{34}{Теорема Вейерштрасса о существовании целой функции с заданными нулями}

\begin{defs}
\begin{itemize}
\item \textbf{Целая функция} --- функция, голоморфная во всей плоскости $\CC$. Нули непостоянной (не тождественно
  нулевой) целой функции изолированы, поэтому в каждом круге их конечное число; если их бесконечно много, то они
  стремятся к $\infty$.
\item \textbf{Кратность нуля}: $a$ --- нуль кратности $k\ge1$ функции $f$, если $f(z)=(z-a)^kh(z)$, $h$ голоморфна около $a$,
  $h(a)\ne0$; эквивалентно $f(a)=\dots=f^{(k-1)}(a)=0\ne f^{(k)}(a)$.
\item \textbf{Стандартная запись последовательности нулей.} Нули целой $f\not\equiv0$ записываем в виде
  $\{0,\dots,0,a_1,a_2,\dots\}$: первые $m\in\ZZ_+$ членов --- нули (кратность нуля $f$ в точке $0$), далее
  $0<|a_1|\le|a_2|\le\dots$, каждый нуль повторён столько раз, какова его кратность, и (если нулей бесконечно много)
  $|a_n|\to\infty$. Обратно, назовём \emph{допустимой} любую последовательность такого вида.
\item \textbf{Бесконечное произведение} $\prod_{n=1}^\infty f_n(z)$ функций на множестве $K$ \emph{сходится равномерно на $K$},
  если частичные произведения $\Pi_N=\prod_{n=1}^Nf_n$ сходятся равномерно на $K$. Для голоморфных $f_n$ в области
  $G$ \emph{локально равномерная} сходимость (равномерная на каждом компакте) влечёт голоморфность предела
  (теорема Вейерштрасса о равномерных пределах голоморфных функций).
\item \textbf{Главная ветвь логарифма}: $\operatorname{Log}w=\ln|w|+i\arg w$, $\arg w\in(-\pi,\pi)$, голоморфна в
  $\CC\setminus(-\infty,0]$; при $|u|<1$ $\operatorname{Log}(1-u)=-\sum_{k\ge1}\frac{u^k}{k}$.
  \emph{Пример:} $\operatorname{Log}(1-\tfrac12)=-\ln2=-\sum_k\frac1{k2^k}$.
\item \textbf{Первичные множители Вейерштрасса}: $E_0(u)=1-u$, и при $p\in\NN$
  \[
  E_p(u)=(1-u)\exp\Big(u+\frac{u^2}2+\dots+\frac{u^p}p\Big).
  \]
  $E_p$ --- целая функция с единственным (простым) нулём в точке $u=1$. Идея: при малых $|u|$ показатель
  $u+\dots+u^p/p$ есть начальный отрезок ряда $-\operatorname{Log}(1-u)$, поэтому $E_p(u)$ очень близко к $1$:
  $E_p(u)=\exp\big(-\sum_{k>p}u^k/k\big)$.
\end{itemize}
\end{defs}

\subsection{Мотивировка: функции с конечным числом нулей}

\begin{thmbox}{Лемма 12.3 (логарифм функции без нулей)}
Если $G$ --- односвязная область, $f\in H(G)$ и $f\ne0$ в $G$, то существует $g\in H(G)$ с $f=e^g$.
\end{thmbox}
\begin{proof}
Функция $f'/f$ голоморфна в $G$; в односвязной области у неё есть первообразная $h$ (интеграл не зависит от пути по
теореме Коши). Тогда $(fe^{-h})'=e^{-h}(f'-fh')=0$, т.е. $fe^{-h}\equiv c\ne0$. Выбрав $b$ с $e^b=c$, получаем
$f=e^{h+b}$.
\end{proof}

\begin{thmbox}{Утверждение 34.1}
(а) Целая функция без нулей имеет вид $e^{g}$, $g$ --- целая. (б) Целая функция, у которой ровно конечное множество
нулей $a_1,\dots,a_n$ (с учётом кратности), имеет вид $f(z)=e^{g(z)}\prod_{k=1}^n(z-a_k)$, $g$ --- целая.
\end{thmbox}
\begin{proof}
(а) --- Лемма 12.3 для $G=\CC$. (б) Функция $h=f/\prod_k(z-a_k)$ имеет в точках $a_k$ устранимые особенности
(в $a$ кратности $s$ и числитель, и знаменатель имеют нуль ровно порядка $s$), после доопределения $h$ --- целая без нулей,
и $h=e^g$ по (а).
\end{proof}

Для бесконечного множества нулей произведение $\prod(z-a_n)$ расходится; нужно нормировать множители (рассматривать
$1-z/a_n$) и, кроме того, домножать их на «сходящие» экспоненциальные множители.

\subsection{Бесконечные произведения голоморфных функций}

\begin{thmbox}{Лемма 34.2 (сходимость бесконечных произведений)}
Пусть $G\subset\CC$ --- область, $h_n\in H(G)$ и ряд $\sum_n h_n$ сходится равномерно на компакте $K\subset G$.
Тогда $\prod_{n}e^{h_n}$ сходится равномерно на $K$ к $e^{H}$, $H=\sum_n h_n$; в частности, предел нигде на $K$ не обращается в нуль.
\end{thmbox}
\begin{proof}
Пусть $H_N=\sum_{n\le N}h_n$, тогда $\prod_{n\le N}e^{h_n}=e^{H_N}$. Функция $H$ непрерывна на $K$, поэтому
$|H|\le B$ на $K$; при больших $N$ $|H_N-H|\le1$, т.е. $|H_N|\le B+1$. Для $|w|,|w'|\le B+1$
$|e^w-e^{w'}|=\big|\int_{[w',w]}e^\zeta d\zeta\big|\le e^{B+1}|w-w'|$. Значит $\sup_K|e^{H_N}-e^H|\le e^{B+1}\sup_K|H_N-H|\to0$.
\end{proof}

\begin{thmbox}{Следствие 34.3}
Пусть $f_n\in H(G)$, и для каждого компакта $K\subset G$ найдётся $N_K$ такое, что при $n>N_K$ на $K$ определены
голоморфные (в окрестности $K$) $h_n$ с $f_n=e^{h_n}$ и ряд $\sum_{n>N_K}h_n$ сходится равномерно на $K$. Тогда
$f=\prod_{n\ge1}f_n$ сходится локально равномерно в $G$, $f\in H(G)$, и нули $f$ --- это объединение нулей всех $f_n$,
причём кратность нуля $a$ функции $f$ равна сумме кратностей нуля $a$ у всех $f_n$ (сумма конечна).
\end{thmbox}
\begin{proof}
На компакте $K$: $\prod_{n\le N}f_n=\Big(\prod_{n\le N_K}f_n\Big)\cdot\prod_{N_K<n\le N}e^{h_n}$. Первый множитель не зависит
от $N$ и ограничен на $K$, второй по Лемме 34.2 сходится равномерно на $K$ к $e^{H_K}$, $H_K=\sum_{n>N_K}h_n$; значит частичные
произведения сходятся равномерно на $K$ к $f=\big(\prod_{n\le N_K}f_n\big)e^{H_K}$. Так как это верно для любого компакта,
$f\in H(G)$ (теорема Вейерштрасса). Для точки $a$ возьмём $K$ --- замкнутый круг с центром $a$ в $G$; на $K$ (в окрестности $a$)
$f=\big(\prod_{n\le N_K}f_n\big)e^{H_K}$ и $e^{H_K}\ne0$; кратность нуля конечного произведения равна сумме кратностей.
\end{proof}

\subsection{Теорема Вейерштрасса}

\begin{thmbox}{Теорема 16.2 (Вейерштрасс)}
Для любой допустимой последовательности $\{0,\dots,0,a_1,a_2,\dots\}$ ($m$ нулей, затем $0<|a_n|\le|a_{n+1}|\to\infty$)
существует целая функция, нули которой (с учётом кратности) в точности совпадают с этой последовательностью. А именно,
такова функция
\begin{equation}\label{t34:166}
\Phi(z)=z^m\prod_{n=1}^\infty\Big(1-\frac z{a_n}\Big)\exp\big(P_n(z)\big),\qquad
P_n(z)=\sum_{k=1}^n\frac{z^k}{k\,a_n^k}.\tag{16.6}
\end{equation}
\end{thmbox}
\begin{proof}
(Если последовательность $\{a_n\}$ конечна, годится $z^m\prod_n(1-z/a_n)$; далее она бесконечна.)
Положим $f_n(z)=(1-z/a_n)e^{P_n(z)}=E_n(z/a_n)$ --- целые функции, единственный нуль $f_n$ --- простой нуль в точке $a_n$.
Применим Следствие 34.3 с $G=\CC$. Достаточно рассмотреть компакты $K=\clos{D_R}$, $D_R=\{|z|<R\}$, $R>0$ (любой
компакт лежит в таком круге).

Фиксируем $R$ и выберем $N$ так, что $|a_n|>2R$ при $n>N$ (возможно, так как $|a_n|\to\infty$). Тогда при
$|z|\le R$, $n>N$: $|z/a_n|<\frac12$, и определено
\[
h_n(z):=\operatorname{Log}\Big(1-\frac z{a_n}\Big)+P_n(z)=-\sum_{k=1}^\infty\frac{z^k}{ka_n^k}+\sum_{k=1}^n\frac{z^k}{ka_n^k}
=-\sum_{k=n+1}^\infty\frac{z^k}{ka_n^k},
\]
голоморфная в круге $|z|<|a_n|$ (содержащем $\clos{D_R}$), причём $e^{h_n}=f_n$. Оценка:
\[
|h_n(z)|\le\sum_{k=n+1}^\infty\frac1k\Big|\frac z{a_n}\Big|^k\le\sum_{k=n+1}^\infty\frac1{2^k}=\frac1{2^n},\qquad |z|\le R,\ n>N.
\]
По признаку Вейерштрасса ряд $\sum_{n>N}h_n$ сходится равномерно на $\clos{D_R}$. По Следствию 34.3 произведение
$\prod_nf_n$ сходится локально равномерно в $\CC$ к целой функции, нули которой --- точки $a_n$, причём кратность нуля
в точке $a$ равна числу индексов $n$ с $a_n=a$ (каждый $f_n$ имеет простой нуль только в $a_n$). Множитель $z^m$
добавляет нуль кратности $m$ в нуле (а $f_n(0)=1\ne0$). Итак, $\Phi$ --- искомая функция.
\end{proof}

\begin{remarkbox}{Замечание (зачем экспоненты)}
Без множителей $e^{P_n}$ произведение $\prod(1-z/a_n)$ может расходиться: так, для $a_n=n$ имеем
$\prod_{n\le N}(1+\frac1n)=N+1\to\infty$ при $z=-1$. Множитель $e^{P_n}$ «гасит» начальные члены ряда $\operatorname{Log}(1-z/a_n)$, и
остаток $h_n$ имеет порядок $|z/a_n|^{n+1}$. Степень $n$ многочлена $P_n$ выбрана с запасом: она годится для любой
допустимой последовательности, но растёт с $n$, и поэтому рост $\Phi$ трудно контролировать (слайд 442). Для функций
конечного порядка достаточно многочленов фиксированной степени $p$ (см. билет 36).
\end{remarkbox}

\begin{thmbox}{Следствие 34.4 (общий вид целой функции с данными нулями)}
Если $f\not\equiv0$ --- целая функция с нулями $\{0,\dots,0,a_1,a_2,\dots\}$ ($m$ нулей в начале координат), то
\[
f(z)=z^me^{g(z)}\prod_{n=1}^\infty\Big(1-\frac z{a_n}\Big)e^{P_n(z)},\qquad g\ \text{--- целая функция.}
\]
\end{thmbox}
\begin{proof}
$f/\Phi$ имеет лишь устранимые особенности (в каждой точке $a$ нули $f$ и $\Phi$ одной кратности), после доопределения это
целая функция без нулей, по Лемме 12.3 она равна $e^g$.
\end{proof}

\subsection{Первичные множители и их оценка}

\begin{thmbox}{Лемма 34.5 (оценка первичного множителя)}
Для $p\in\ZZ_+$ и $|u|\le1$ выполнено $|1-E_p(u)|\le|u|^{p+1}$.
\end{thmbox}
\begin{proof}
При $p=0$: $1-E_0(u)=u$, и $|u|\le|u|^1$. Пусть $p\ge1$, $s(u)=u+\dots+u^p/p$. Дифференцируем:
\[
E_p'(u)=-e^{s(u)}+(1-u)(1+u+\dots+u^{p-1})e^{s(u)}=-e^{s(u)}+(1-u^p)e^{s(u)}=-u^pe^{s(u)}.
\]
Пусть $1-E_p(u)=\sum_{k\ge0}b_ku^k$ (целая функция). Так как $E_p(0)=1$, $b_0=0$. Далее
$\big(1-E_p\big)'(u)=u^pe^{s(u)}$, а у $e^{s(u)}=\sum_j\frac{s(u)^j}{j!}$ все тейлоровские коэффициенты \emph{неотрицательны}
(коэффициенты $s$ неотрицательны, а произведения и суммы рядов с неотрицательными коэффициентами имеют неотрицательные
коэффициенты). Значит $kb_k\ge0$ для всех $k$, и $b_k=0$ при $1\le k\le p$ (производная начинается с $u^p$). Итак,
\[
1-E_p(u)=u^{p+1}\varphi(u),\qquad \varphi(u)=\sum_{j\ge0}b_{p+1+j}u^j,\qquad b_{p+1+j}\ge0.
\]
При $|u|\le1$: $|\varphi(u)|\le\sum_jb_{p+1+j}|u|^j\le\sum_jb_{p+1+j}=\varphi(1)=1-E_p(1)=1$, так как $E_p(1)=0$.
Значит $|1-E_p(u)|\le|u|^{p+1}$.
\end{proof}

\begin{thmbox}{Теорема 34.6 (Вейерштрасс, вариант с первичными множителями)}
Пусть $\{a_n\}$ --- допустимая последовательность ненулевых чисел, $|a_n|\to\infty$, и $p_n\in\ZZ_+$ таковы, что
\begin{equation}\label{t34:pn}\tag{34.1}
\sum_{n=1}^\infty\Big(\frac r{|a_n|}\Big)^{p_n+1}<\infty\qquad\text{для всех } r>0.
\end{equation}
(Например, $p_n=n-1$: при $n$ таких, что $|a_n|>2r$, член ряда $\le2^{-n}$.) Тогда
$P(z)=\prod_{n=1}^\infty E_{p_n}(z/a_n)$ сходится локально равномерно в $\CC$ к целой функции, нули которой --- в точности $a_n$
(с учётом кратности).
\end{thmbox}
\begin{proof}
Фиксируем $R>0$; по \eqref{t34:pn} $(R/|a_n|)^{p_n+1}\to0$, и при $n>N$ имеем $|a_n|>R$ (так как $|a_n|\to\infty$).
Положим $u_n(z)=E_{p_n}(z/a_n)-1$. При $|z|\le R$, $n>N$ имеем $|z/a_n|\le1$, и по Лемме 34.5
$|u_n(z)|\le(R/|a_n|)^{p_n+1}$; по \eqref{t34:pn} ряд $\sum_n\sup_{|z|\le R}|u_n(z)|$ сходится. Выберем $N'\ge N$ так, что
$|u_n|\le\frac12$ на $\clos{D_R}$ при $n>N'$. Там определено $h_n=\operatorname{Log}(1+u_n)$, и
$|h_n|=\big|\sum_{k\ge1}\frac{(-1)^{k+1}u_n^k}{k}\big|\le|u_n|\sum_{k\ge1}2^{1-k}=2|u_n|$, так что $\sum_{n>N'}h_n$ сходится
равномерно на $\clos{D_R}$. Остаётся применить Следствие 34.3 (каждый $E_{p_n}(z/a_n)$ имеет единственный простой нуль $a_n$).
\end{proof}

\begin{remarkbox}{Замечание}
Произведение \eqref{t34:166} со слайдов есть частный случай: $(1-\frac z{a_n})e^{P_n(z)}=E_n(z/a_n)$, т.е. $p_n=n$.
Если же $\sum|a_n|^{-(p+1)}<\infty$ для некоторого фиксированного $p$, то \eqref{t34:pn} выполнено с $p_n\equiv p$, и получается
произведение $\prod E_p(z/a_n)$ --- каноническое произведение рода $p$ (билет 36).
\end{remarkbox}

\begin{thmbox}{Следствие 34.7 (мероморфные функции)}
Всякая функция $h$, мероморфная в $\CC$, есть отношение двух целых функций.
\end{thmbox}
\begin{proof}
Полюсы $h$ образуют дискретное в $\CC$ множество (конечное или стремящееся к $\infty$). По Теореме 16.2 существует целая $\psi$,
нули которой --- полюсы $h$ с кратностями, равными порядкам полюсов. Тогда $\varphi=h\psi$ имеет во всех полюсах $h$ устранимые
особенности, т.е. $\varphi$ --- целая, и $h=\varphi/\psi$.
\end{proof}

\begin{remarkbox}{Коротко: план ответа на экзамене}
\begin{enumerate}
\item Без нулей: $f=e^g$ (логарифм в односвязной области); конечное число нулей: $f=e^g\prod(z-a_k)$.
\item Формулировка: для любой последовательности $0<|a_n|\le|a_{n+1}|\to\infty$ (и $m$ нулей в $0$) есть целая функция с этими нулями.
\item Конструкция $z^m\prod(1-z/a_n)e^{P_n(z)}$, $P_n=\sum_{k\le n}z^k/(ka_n^k)$.
\item Сходимость произведения $\Leftarrow$ равномерная сходимость ряда логарифмов (Лемма 34.2).
\item В круге $|z|\le R$ при $|a_n|>2R$: $|\operatorname{Log}(1-z/a_n)+P_n|=|\sum_{k>n}z^k/(ka_n^k)|\le2^{-n}$.
\item Нули предела = нули множителей (хвост произведения $=e^{H}\ne0$).
\item Следствие: всякая целая функция с нулями $\{a_n\}$ равна $z^me^g\prod(\dots)$; мероморфная = отношение целых.
\item Первичные множители $E_p$ и оценка $|1-E_p(u)|\le|u|^{p+1}$ при $|u|\le1$ (через $E_p'=-u^pe^{s}$).
\end{enumerate}
\end{remarkbox}

% ======================================================================
\ticket{35}{Связь показателя сходимости последовательности нулей и порядка целой функции конечного порядка}

\begin{defs}
\begin{itemize}
\item \textbf{Порядок} целой $f$ (Определение 16.1): нижняя грань $\rho$ чисел $\alpha>0$ таких, что $M_f(r)<\exp(r^\alpha)$
  при всех больших $r$; здесь $M_f(r)=\max_{|z|=r}|f(z)|$. Для непостоянной $f$:
  $\rho=\limsup_{r\to\infty}\frac{\ln\ln M_f(r)}{\ln r}$. Если $\rho<\infty$, то для всякого $\eps>0$
  $\ln M_f(r)<r^{\rho+\eps}$ при $r>R_\eps$. \textbf{Тип} $\sigma$ (Определение 16.2) при $0<\rho<\infty$: нижняя грань $s>0$ с
  $M_f(r)<\exp(sr^\rho)$ при больших $r$. (Подробно --- билет 33.)
\item \textbf{Последовательность нулей} $f\not\equiv0$: ненулевые нули $a_1,a_2,\dots$, повторённые с учётом кратности,
  $0<|a_n|\le|a_{n+1}|$, $|a_n|\to\infty$ (если их бесконечно много); $m$ --- кратность нуля в точке $0$.
\item \textbf{Считающая функция нулей}: $n_f(r)=\#\{z\in B_r:\ f(z)=0\}$ (с учётом кратности), $B_r=B(0,r)=\{|z|<r\}$.
  Она не убывает, непрерывна слева и кусочно постоянна. \emph{Пример:} для $\sin z$ $n(r)=2\lceil r/\pi\rceil-1$
  (нули $0,\pm\pi,\dots$ в открытом круге).
\item \textbf{Показатель сходимости} (Определение 16.3) последовательности $0<|a_n|\le|a_{n+1}|\to\infty$:
  $\tau=\inf\{\lambda>0:\ \sum_n|a_n|^{-\lambda}<\infty\}$; если ряд расходится при всех $\lambda>0$, то $\tau=\infty$.
  Множество $\lambda$, при которых ряд сходится, --- луч: если ряд сходится при $\lambda$, то и при $\lambda'>\lambda$
  (так как $|a_n|^{-\lambda'}\le|a_n|^{-\lambda}$ при $|a_n|\ge1$). Поэтому ряд сходится при $\lambda>\tau$ и расходится при
  $\lambda<\tau$; при $\lambda=\tau$ возможно и то, и другое. Для конечной последовательности (или пустой) полагаем $\tau=0$.
  \emph{Примеры:} $a_n=n$: $\tau=1$, ряд при $\lambda=\tau$ расходится; $a_n=n\ln^2(n+1)$: $\tau=1$, ряд при $\lambda=1$ сходится.
\item \textbf{Род последовательности} $\{a_n\}$ --- наименьшее $p\in\ZZ_+$ с $\sum|a_n|^{-(p+1)}<\infty$ (если такое есть).
  Выполнено $p\le\tau\le p+1$ (см. билет 36).
\item \textbf{Гармоническая функция и теорема о среднем}: если $F\in H(B(0,r'))$, то для $u=\Re F$ и $r<r'$
  $u(0)=\frac1{2\pi}\int_0^{2\pi}u(re^{i\vartheta})\,d\vartheta$ (вещественная часть формулы Коши
  $F(0)=\frac1{2\pi}\int_0^{2\pi}F(re^{i\vartheta})d\vartheta$).
\end{itemize}
\end{defs}

\subsection{Формула Йенсена}

\begin{thmbox}{Лемма 35.1}
$\displaystyle\int_0^{2\pi}\ln|1-e^{i\vartheta}|\,d\vartheta=0$ (интеграл сходится абсолютно); то же для $\ln|1-e^{i(\vartheta-\beta)}|$, $\beta\in\RR$.
\end{thmbox}
\begin{proof}
При $0<s<1$ функция $\operatorname{Log}(1-sz)$ голоморфна в круге $|z|<1/s$, содержащем $\clos\DD$, её вещественная часть
$\ln|1-sz|$ по теореме о среднем даёт $\int_0^{2\pi}\ln|1-se^{i\vartheta}|\,d\vartheta=2\pi\ln|1-0|=0$. Перейдём к пределу при
$s\to1-$. При $s\in[\frac12,1)$
\[
|1-se^{i\vartheta}|^2=(1-s)^2+2s(1-\cos\vartheta)\ge 4s\sin^2\tfrac\vartheta2\ge2\sin^2\tfrac\vartheta2,
\qquad |1-se^{i\vartheta}|\le2,
\]
так что $\big|\ln|1-se^{i\vartheta}|\big|\le\ln2+\big|\ln|\sin\tfrac\vartheta2|\big|$. Мажоранта интегрируема на $[0,2\pi]$: при
$x\in[0,\frac\pi2]$ $\sin x\ge\frac{2x}\pi$ (вогнутость синуса), так что около концов $|\ln|\sin\frac\vartheta2||\le|\ln\frac{\vartheta}{\pi}|$
(и симметрично у $2\pi$), а $\ln$ интегрируем в нуле. При $\vartheta\ne0,2\pi$ $\ln|1-se^{i\vartheta}|\to\ln|1-e^{i\vartheta}|$.
По теореме Лебега о мажорированной сходимости интеграл предельной функции равен $0$. Сдвиг на $\beta$ не меняет интеграла
по периоду.
\end{proof}

\begin{thmbox}{Лемма 12.2 (формула Йенсена)}
Пусть $f\in H(B(0,R))$, $f(0)\ne0$, $0<r<R$, и $a_1,a_2,\dots$ --- нули $f$ в $B(0,R)$ с учётом кратности. Тогда
\begin{equation}\label{t35:jensen}
\frac1{2\pi}\int_0^{2\pi}\ln|f(re^{i\vartheta})|\,d\vartheta=\ln|f(0)|+\sum_{n:\,|a_n|<r}\ln\frac r{|a_n|}.\tag{12.1}
\end{equation}
\end{thmbox}
\begin{proof}
В замкнутом круге $|z|\le r$ (компакт в $B(0,R)$) лишь конечное число нулей: $a_1,\dots,a_k$ с $|a_j|<r$ и $b_1,\dots,b_l$ с
$|b_j|=r$, $b_j=re^{i\beta_j}$. Положим
\[
F(z)=f(z)\prod_{j=1}^k\frac{r^2-\clos{a_j}z}{r(a_j-z)}\prod_{j=1}^l\frac{b_j}{b_j-z}.
\]
Нули знаменателей сокращаются с нулями $f$ (с учётом кратности), числители $r^2-\clos{a_j}z$ обращаются в нуль лишь в точках
$r^2/\clos{a_j}$ с модулем $r^2/|a_j|>r$. Так как нули $f$ дискретны, найдётся $r'\in(r,R)$ такое, что в кольце $r<|z|<r'$ нет нулей $f$
и $r'<\min_j r^2/|a_j|$. Тогда $F$ голоморфна и не имеет нулей в $B(0,r')$, по Лемме 12.3 $F=e^{G}$, $G\in H(B(0,r'))$, и
$\ln|F|=\Re G$. По теореме о среднем
\[
\ln|F(0)|=\frac1{2\pi}\int_0^{2\pi}\ln|F(re^{i\vartheta})|\,d\vartheta.
\]
Вычислим обе части. $F(0)=f(0)\prod_{j\le k}\frac{r^2}{ra_j}\cdot1$, т.е. $\ln|F(0)|=\ln|f(0)|+\sum_{j\le k}\ln\frac r{|a_j|}$.
При $|z|=r$: $|r^2-\clos{a_j}z|=|z\bar z-\clos{a_j}z|=r|\bar z-\clos{a_j}|=r|z-a_j|$, так что множители первого произведения имеют
модуль $1$; а $\big|\frac{b_j}{b_j-re^{i\vartheta}}\big|=\frac1{|1-e^{i(\vartheta-\beta_j)}|}$. Поэтому (при $re^{i\vartheta}\ne b_j$)
\[
\ln|F(re^{i\vartheta})|=\ln|f(re^{i\vartheta})|-\sum_{j=1}^l\ln|1-e^{i(\vartheta-\beta_j)}|.
\]
По Лемме 35.1 интегралы от слагаемых суммы равны нулю (и, значит, $\ln|f(re^{i\vartheta})|$ интегрируема). Подставляя, получаем \eqref{t35:jensen}.
\end{proof}

\begin{remarkbox}{Замечание}
На слайде 264 формула доказывается для $r\ne|a_n|$, а общий случай оставлен как упражнение; приведённое доказательство
(с множителями $b_j/(b_j-z)$ и Леммой 35.1) закрывает и этот случай.
\end{remarkbox}

\begin{thmbox}{Следствие 35.2 (формула Йенсена через считающую функцию)}
Если $f$ --- целая, $f(0)\ne0$, то при всех $r>0$
\[
\frac1{2\pi}\int_0^{2\pi}\ln|f(re^{i\vartheta})|\,d\vartheta=\ln|f(0)|+\int_0^r\frac{n_f(t)}t\,dt,
\qquad\text{и потому}\qquad
\int_0^r\frac{n_f(t)}t\,dt\le\ln\frac{M_f(r)}{|f(0)|}.
\]
\end{thmbox}
\begin{proof}
Для каждого нуля $a_n$ с $|a_n|<r$: $\ln\frac r{|a_n|}=\int_{|a_n|}^r\frac{dt}t=\int_0^r\mathbf 1_{\{|a_n|<t\}}\frac{dt}t$. Суммируя по
таким $n$ и замечая, что при $0<t\le r$ $\sum_{n:|a_n|<r}\mathbf 1_{\{|a_n|<t\}}=\#\{n:|a_n|<t\}=n_f(t)$, получаем
$\sum_{|a_n|<r}\ln\frac r{|a_n|}=\int_0^r\frac{n_f(t)}t\,dt$ (интеграл собственный: $n_f(t)=0$ при $t<|a_1|$). Неравенство: подынтегральная
функция слева не превосходит $\ln M_f(r)$. (На слайде 444 то же получено разбиением $[0,r]$ точками $|a_k|$.)
\end{proof}

\subsection{Оценка считающей функции через рост}

\begin{thmbox}{Утверждение 35.3 (оценки (16.8)--(16.10))}
Пусть $f$ --- целая, $f(0)\ne0$.
\begin{enumerate}
\item[(а)] Для $\delta\in(0,1)$, $r>0$: $\displaystyle n_f(\delta r)\le\frac1{\ln(1/\delta)}\ln\frac{M_f(r)}{|f(0)|}$. \hfill(16.8)
\item[(б)] Если $\ln M_f(r)<Br^\alpha$ при $r>R_0$ ($B,\alpha>0$), то $\displaystyle\limsup_{r\to\infty}\frac{n_f(r)}{r^\alpha}\le e\alpha B$.
\item[(в)] Если $f$ имеет конечный порядок $\rho$, то $\displaystyle\limsup_{r\to\infty}\frac{n_f(r)}{r^{\rho+\eps}}\le e(\rho+\eps)$ для всех $\eps>0$ \hfill(16.9)\\
  и, если $0<\rho<\infty$ и тип $\sigma<\infty$, то $\displaystyle\limsup_{r\to\infty}\frac{n_f(r)}{r^{\rho}}\le e\rho\sigma$. \hfill(16.10)
\end{enumerate}
\end{thmbox}
\begin{proof}
(а) $n_f$ не убывает, поэтому
$\int_0^r\frac{n_f(t)}t\,dt\ge\int_{\delta r}^r\frac{n_f(t)}t\,dt\ge n_f(\delta r)\int_{\delta r}^r\frac{dt}t=n_f(\delta r)\ln\frac1\delta$; применяем Следствие 35.2.

(б) Фиксируем $\delta\in(0,1)$. По (а) при $r>R_0$: $n_f(\delta r)\le\frac{Br^\alpha-\ln|f(0)|}{\ln(1/\delta)}$. Полагая $t=\delta r$
(т.е. $r=t/\delta$, $t>\delta R_0$):
\[
\frac{n_f(t)}{t^\alpha}\le\frac{B\delta^{-\alpha}}{\ln(1/\delta)}-\frac{\ln|f(0)|}{t^\alpha\ln(1/\delta)},\qquad\text{откуда}\qquad
\limsup_{t\to\infty}\frac{n_f(t)}{t^\alpha}\le\frac{B\delta^{-\alpha}}{\ln(1/\delta)}.
\]
Это верно при любом $\delta\in(0,1)$. Положим $s=\ln(1/\delta)\in(0,\infty)$; правая часть равна $\psi(s)=Be^{\alpha s}/s$,
$\psi'(s)=Be^{\alpha s}(\alpha s-1)/s^2$, так что минимум достигается при $s=1/\alpha$ (т.е. $\delta=e^{-1/\alpha}$) и равен
$\psi(1/\alpha)=Be\alpha$.

(в) Для (16.9): по определению порядка $\ln M_f(r)<r^{\rho+\eps}$ при больших $r$; это (б) с $B=1$, $\alpha=\rho+\eps$.
Для (16.10): при $\eps>0$ $\ln M_f(r)<(\sigma+\eps)r^\rho$ при больших $r$ (определение типа); по (б) с $B=\sigma+\eps$, $\alpha=\rho$
верхний предел $\le e\rho(\sigma+\eps)$; так как $\eps$ произвольно, он $\le e\rho\sigma$.
\end{proof}

\begin{remarkbox}{Замечание}
На слайде 447 вывод (16.10) назван выходящим за рамки курса; как видно, он дословно повторяет вывод (16.9) (пункт (б) выше
доказан сразу для произвольных $B$ и $\alpha$). Отметим также, что из (16.9), применённого с $\eps/2$, следует даже
$n_f(r)=o(r^{\rho+\eps})$.
\end{remarkbox}

\begin{thmbox}{Лемма 16.4}
Пусть $f\not\equiv0$ --- целая функция конечного порядка $\rho$, $\{a_n\}$ --- её последовательность (ненулевых) нулей,
бесконечная. Тогда для любого $\eps>0$
\[
\limsup_{n\to\infty}\frac n{|a_n|^{\rho+\eps}}\le e(\rho+\eps),\tag{16.11}
\]
а если $0<\rho<\infty$ и $f$ имеет конечный тип $\sigma$, то
$\displaystyle\limsup_{n\to\infty}\frac n{|a_n|^{\rho}}\le e\rho\sigma$. \hfill(16.12)
\end{thmbox}
\begin{proof}
\emph{Сведение к $f(0)\ne0$.} Пусть $f(z)=z^mh(z)$, $h(0)\ne0$; ненулевые нули $f$ и $h$ совпадают, и $n_h(t)=\#\{n:|a_n|<t\}$.
По принципу максимума $M_h(r)=\max_{|z|=r}|f(z)|/r^m=M_f(r)/r^m\le M_f(r)$ при $r\ge1$. Поэтому оценки роста
$\ln M_f(r)<r^{\rho+\eps}$, $\ln M_f(r)<(\sigma+\eps)r^\rho$ при больших $r$ верны и для $h$, и по Утверждению 35.3(б)
\[
\limsup_{t\to\infty}\frac{n_h(t)}{t^{\rho+\eps}}\le e(\rho+\eps),\qquad \limsup_{t\to\infty}\frac{n_h(t)}{t^{\rho}}\le e\rho(\sigma+\eps).
\]
\emph{Переход к $a_n$.} Положим $t_n=|a_n|+1$. Так как $|a_1|\le\dots\le|a_n|<t_n$, точки $a_1,\dots,a_n$ лежат в $B_{t_n}$, т.е.
$n\le n_h(t_n)$. Поэтому для $\alpha>0$
\[
\frac n{|a_n|^\alpha}\le\frac{n_h(t_n)}{t_n^\alpha}\Big(\frac{t_n}{|a_n|}\Big)^\alpha .
\]
При $n\to\infty$: $t_n\to\infty$ и $t_n/|a_n|\to1$, поэтому $\limsup_n\frac n{|a_n|^\alpha}\le\limsup_{t\to\infty}\frac{n_h(t)}{t^\alpha}$.
При $\alpha=\rho+\eps$ получаем (16.11), при $\alpha=\rho$ --- оценку $e\rho(\sigma+\eps)$ для всех $\eps>0$, т.е. (16.12).
\end{proof}

\subsection{Показатель сходимости: формула и примеры}

\begin{thmbox}{Лемма 16.3}
Показатель сходимости последовательности $0<|a_n|\le|a_{n+1}|$, $|a_n|\to\infty$, равен
$\displaystyle\tau=\limsup_{n\to\infty}\frac{\ln n}{\ln|a_n|}$.
\end{thmbox}
\begin{proof}
(На слайде 446 --- задача.) Обозначим верхний предел через $\mu\in[0,\infty]$ (при больших $n$ $|a_n|>1$, так что
выражение определено и неотрицательно).

\emph{$\tau\le\mu$.} Пусть $\mu<\infty$ и $\lambda>\mu$; возьмём $\lambda'\in(\mu,\lambda)$, $\lambda'>0$. При больших $n$
$\frac{\ln n}{\ln|a_n|}<\lambda'$, т.е. $|a_n|>n^{1/\lambda'}$, и $|a_n|^{-\lambda}<n^{-\lambda/\lambda'}$, где $\lambda/\lambda'>1$. Ряд
$\sum|a_n|^{-\lambda}$ сходится, значит $\tau\le\lambda$; так как $\lambda>\mu$ любое, $\tau\le\mu$.

\emph{$\mu\le\tau$.} Пусть $\tau<\infty$ и $\lambda>\tau$, так что $S=\sum_k|a_k|^{-\lambda}<\infty$. Так как $|a_k|\le|a_n|$ при $k\le n$,
\[
\frac n{|a_n|^\lambda}\le\sum_{k=1}^n\frac1{|a_k|^\lambda}\le S\quad\Longrightarrow\quad \lambda\ln|a_n|\ge\ln n-\ln S
\quad\Longrightarrow\quad \frac{\ln n}{\ln|a_n|}\le\frac{\lambda\ln n}{\ln n-\ln S}\xrightarrow[n\to\infty]{}\lambda
\]
(при больших $n$, когда $\ln n>\ln S$). Значит $\mu\le\lambda$ для любого $\lambda>\tau$, т.е. $\mu\le\tau$.
Случаи $\mu=\infty$ или $\tau=\infty$ покрываются этими же неравенствами.
\end{proof}

\begin{remarkbox}{Примеры (слайд 446)}
\begin{enumerate}
\item $a_n=e^n$: $\frac{\ln n}{n}\to0$, $\tau=0$ (ряд $\sum e^{-\lambda n}$ сходится при всех $\lambda>0$).
\item $a_n=n^{1/\beta}$, $\beta>0$: $\frac{\ln n}{\ln n/\beta}=\beta$, $\tau=\beta$ (ряд $\sum n^{-\lambda/\beta}$ сходится $\iff\lambda>\beta$).
\item $a_n=\ln n$ ($n\ge2$): $\frac{\ln n}{\ln\ln n}\to\infty$, $\tau=\infty$ (ряд $\sum(\ln n)^{-\lambda}$ расходится при любом $\lambda$, так как
  $(\ln n)^\lambda<n$ при больших $n$).
\end{enumerate}
\end{remarkbox}

\subsection{Лемма Адамара о нулях}

\begin{thmbox}{Предложение 16.1 (лемма Адамара о нулях)}
Пусть $f\not\equiv0$ --- целая функция конечного порядка $\rho$. Тогда последовательность её нулей имеет конечный
показатель сходимости $\tau$ и $\tau\le\rho$.
\end{thmbox}
\begin{proof}
Если ненулевых нулей конечное число, $\tau=0\le\rho$ по соглашению. Пусть их бесконечно много, $\eps>0$. По (16.11)
$\limsup_n n|a_n|^{-(\rho+\eps)}\le e(\rho+\eps)<e(\rho+2\eps)$, значит найдётся $N$ такое, что при $n>N$
\[
\frac n{|a_n|^{\rho+\eps}}<e(\rho+2\eps),\qquad\text{т.е.}\qquad \frac1{|a_n|}<\Big(\frac{e(\rho+2\eps)}n\Big)^{1/(\rho+\eps)}.
\]
Для $\lambda>0$: $|a_n|^{-\lambda}<\big(e(\rho+2\eps)\big)^{\lambda/(\rho+\eps)}n^{-\lambda/(\rho+\eps)}$. При $\lambda>\rho+\eps$ показатель
$\lambda/(\rho+\eps)>1$, и ряд $\sum|a_n|^{-\lambda}$ сходится. Значит $\tau\le\rho+\eps$ для любого $\eps>0$, т.е. $\tau\le\rho$.
\end{proof}

\begin{remarkbox}{Другой вариант через Лемму 16.3}
Из (16.11): $n<e(\rho+2\eps)|a_n|^{\rho+\eps}$ при больших $n$, т.е. $\ln n<\ln\big(e(\rho+2\eps)\big)+(\rho+\eps)\ln|a_n|$, и
$\limsup\frac{\ln n}{\ln|a_n|}\le\rho+\eps$. По Лемме 16.3 $\tau\le\rho+\eps$.
\end{remarkbox}

\begin{thmbox}{Следствие (слайд 449)}
Если порядок целой $f$ не превосходит $\rho$, а показатель сходимости её нулей больше $\rho$, то $f\equiv0$.
\end{thmbox}
\begin{proof}
Если $f\not\equiv0$, то по Предложению 16.1 показатель сходимости нулей не больше порядка $f$, т.е. не больше $\rho$, --- противоречие.
\end{proof}

\begin{remarkbox}{Примеры: когда $\tau<\rho$ и когда $\tau=\rho$}
\begin{enumerate}
\item \emph{Строгое неравенство.} $e^z$ (порядок $1$) нулей не имеет: $\tau=0<1=\rho$. Вообще $e^{P(z)}$, $\deg P=k\ge1$:
  $\tau=0<k=\rho$. С бесконечным числом нулей: $f(z)=e^{z^2}\sin z$. Нули --- $\pi k$, $\tau=1$ (см. п.~2). Порядок $2$: сверху
  $M_f(r)\le e^{r^2}\operatorname{ch}r\le e^{r^2+r}$, а в точках $r_k=\pi k+\frac\pi2$ $|f(r_k)|=e^{r_k^2}$, так что
  $\limsup\frac{\ln\ln M_f(r)}{\ln r}=2$. Итак, $\tau=1<2=\rho$.
\item \emph{Равенство.} $\sin z$: $\rho=1$; ненулевые нули $\pm\pi n$, т.е. $|a_{2n-1}|=|a_{2n}|=\pi n$, и
  $\sum|a_n|^{-\lambda}=2\pi^{-\lambda}\sum n^{-\lambda}<\infty\iff\lambda>1$: $\tau=1=\rho$. Так же $\cos z$, $e^z-1$ (нули $2\pi ik$):
  $\tau=\rho=1$ --- хотя $e^z-1$ растёт так же, как $e^z$, у которой нулей нет.
\item $\cos\sqrt z$: $\rho=\frac12$ (билет 33); нули $z_k=\pi^2(k-\frac12)^2$, $k\in\NN$ ($\cos w=0\iff w=\pi(k-\frac12)$,
  $k\in\ZZ$, а $w$ и $-w$ дают одно $z=w^2$). $\sum z_k^{-\lambda}<\infty\iff2\lambda>1$: $\tau=\frac12=\rho$.
\item $1/\Gamma(z)$: порядок $1$, нули $0,-1,-2,\dots$, $\tau=1=\rho$ (без доказательства роста).
\item Как будет показано в билете 37 (теорема Бореля), если $\rho\notin\ZZ$, то всегда $\tau=\rho$; строгое неравенство $\tau<\rho$
  возможно лишь при целом $\rho$, и тогда $\rho=\deg g$ в факторизации Адамара.
\end{enumerate}
\end{remarkbox}

\begin{remarkbox}{Коротко: план ответа на экзамене}
\begin{enumerate}
\item Определения: $n_f(r)$, показатель сходимости $\tau$; формула $\tau=\limsup\frac{\ln n}{\ln|a_n|}$ (Лемма 16.3), примеры.
\item Формула Йенсена; переписать как $\frac1{2\pi}\int\ln|f|=\ln|f(0)|+\int_0^r\frac{n(t)}t dt$, откуда $\int_0^r\frac{n(t)}tdt\le\ln\frac{M(r)}{|f(0)|}$.
\item $n(\delta r)\ln\frac1\delta\le\ln\frac{M(r)}{|f(0)|}$ (16.8).
\item Подставить $\ln M(r)<r^{\rho+\eps}$, оптимизировать по $\delta$: $\limsup n(r)/r^{\rho+\eps}\le e(\rho+\eps)$ (16.9); аналогично с типом (16.10).
\item $n\le n(|a_n|+1)$ $\Rightarrow$ $\limsup n/|a_n|^{\rho+\eps}\le e(\rho+\eps)$ (Лемма 16.4).
\item Отсюда $|a_n|^{-\lambda}\lesssim n^{-\lambda/(\rho+\eps)}$, ряд сходится при $\lambda>\rho+\eps$: $\tau\le\rho$.
\item Примеры: $e^z$, $e^{z^2}\sin z$ ($\tau<\rho$); $\sin z$, $\cos\sqrt z$ ($\tau=\rho$).
\end{enumerate}
\end{remarkbox}

% ======================================================================
\ticket{36}{Теорема Адамара о факторизации. Канонические произведения}

\begin{defs}
\begin{itemize}
\item \textbf{Порядок} $\rho$ целой функции $f$: нижняя грань $\alpha>0$ с $M_f(r)<\exp(r^\alpha)$ при больших $r$,
  $M_f(r)=\max_{|z|=r}|f|$; при $\rho<\infty$ для любого $\eps>0$ $\ln M_f(r)<r^{\rho+\eps}$ при $r>R_\eps$ (билет 33).
\item \textbf{Последовательность нулей} $f\not\equiv0$: $m$ --- кратность нуля в $0$; $a_1,a_2,\dots$ --- остальные нули с учётом
  кратности, $0<|a_n|\le|a_{n+1}|$, $|a_n|\to\infty$.
\item \textbf{Показатель сходимости} $\tau$ (Определение 16.3): $\inf\{\lambda>0:\sum|a_n|^{-\lambda}<\infty\}$; для конечной
  последовательности $\tau=0$. Ряд сходится при $\lambda>\tau$ и расходится при $0<\lambda<\tau$ (билет 35).
  \emph{Лемма Адамара (Предложение 16.1, билет 35):} если $f$ конечного порядка $\rho$, то $\tau\le\rho$.
\item \textbf{Род последовательности}: наименьшее $p\in\ZZ_+$, для которого $\sum_n|a_n|^{-(p+1)}<\infty$. \emph{Пример:} для
  $a_n=n$ ряд $\sum n^{-1}$ расходится, $\sum n^{-2}$ сходится, $p=1$; для $a_n=n^2$ $p=0$; для $a_n=2^n$ $p=0$.
\item \textbf{Многочлены} $P^p_n(z)=\sum_{k=1}^p\frac{z^k}{k\,a_n^k}$ при $p\ge1$, $P^0_n\equiv0$. Тогда
  $(1-\frac z{a_n})e^{P_n^p(z)}=E_p(z/a_n)$, где $E_p(u)=(1-u)\exp(u+\dots+u^p/p)$ --- первичный множитель Вейерштрасса (билет 34).
\item \textbf{Каноническое произведение} (Определение 16.4) рода $p$:
  \[
  F(z)=z^m\prod_{n=1}^\infty\Big(1-\frac z{a_n}\Big)\exp\big(P^p_n(z)\big)=z^m\prod_{n=1}^\infty E_p\Big(\frac z{a_n}\Big),\tag{16.17}
  \]
  где $p$ --- род последовательности $\{a_n\}$. \emph{Пример:} для нулей $\{\pm\pi n\}$ функции $\sin z/z$ род $p=1$, и
  каноническое произведение есть $\prod_{n\ge1}(1-\frac z{\pi n})e^{z/(\pi n)}(1+\frac z{\pi n})e^{-z/(\pi n)}=\prod_{n\ge1}(1-\frac{z^2}{\pi^2n^2})$.
\item \textbf{Тейлоровские коэффициенты}: $c_j(h)=h^{(j)}(0)/j!$ для $h$, голоморфной около $0$.
\item \textbf{Род целой функции} (доп.): если $f=z^me^{g}F$, $F$ --- каноническое произведение рода $p$, $g$ --- многочлен, то
  $\max(p,\deg g)$ называют родом $f$.
\end{itemize}
\end{defs}

\subsection{Род и показатель сходимости}

\begin{thmbox}{Лемма 36.1}
Пусть $\{a_n\}$ бесконечна, $\tau<\infty$, $p$ --- её род. Тогда $p\le\tau\le p+1$. Если ряд $\sum|a_n|^{-\tau}$ сходится, то $p<\tau\le p+1$;
если расходится, то $p\le\tau<p+1$.
\end{thmbox}
\begin{proof}
Род существует: возьмём целое $q\ge0$ с $q+1>\tau$, тогда $\sum|a_n|^{-(q+1)}<\infty$. Ряд сходится при $\lambda=p+1$, значит
$\tau\le p+1$. Если $p\ge1$, то ряд при $\lambda=p$ расходится (иначе род был бы $\le p-1$); если бы $\tau<p$, то ряд сходился бы при
$\lambda=p>\tau$, значит $\tau\ge p$. При $p=0$ неравенство $\tau\ge0$ тривиально. Если ряд при $\lambda=\tau$ сходится,
то $\tau>0$ (при бесконечной последовательности $\sum|a_n|^0=\infty$) и $\tau\ne p$ (при $p\ge1$ ряд с показателем $p$ расходится), т.е. $p<\tau$.
Если ряд при $\lambda=\tau$ расходится, то $\tau\ne p+1$, т.е. $\tau<p+1$.
\end{proof}

\subsection{Сходимость канонических произведений}

\begin{thmbox}{Лемма 16.5}
Пусть $a_n\ne0$ и существует целое $p\ge0$ с
\begin{equation}\label{t36:1613}
\sum_{n=1}^\infty\frac1{|a_n|^{p+1}}<\infty.\tag{16.13}
\end{equation}
Тогда произведение $f(z)=\prod_{n=1}^\infty\big(1-\frac z{a_n}\big)\exp\big(P^p_n(z)\big)$ \hfill(16.14)\\
сходится локально равномерно в $\CC$, $f$ --- целая, и её нули (с учётом кратности) --- в точности элементы $\{a_n\}$.
\end{thmbox}
\begin{proof}
Из \eqref{t36:1613} $|a_n|\to\infty$. Фиксируем $R>0$ и $N$ с $|a_n|>2R$ при $n>N$. При $|z|\le R$, $n>N$ имеем $|z/a_n|\le\frac12$, и
\[
h_n(z):=\operatorname{Log}\Big(1-\frac z{a_n}\Big)+P^p_n(z)=-\sum_{k=p+1}^\infty\frac{z^k}{ka_n^k},
\]
\[
|h_n(z)|\le\sum_{k=p+1}^\infty\Big|\frac z{a_n}\Big|^k=\frac{|z|^{p+1}}{|a_n|^{p+1}}\sum_{j=0}^\infty\Big|\frac z{a_n}\Big|^j
\le\frac{R^{p+1}}{|a_n|^{p+1}}\sum_{j=0}^\infty\frac1{2^j}=\frac{2R^{p+1}}{|a_n|^{p+1}}.
\]
По \eqref{t36:1613} и признаку Вейерштрасса $\sum_{n>N}h_n$ сходится равномерно на $|z|\le R$, причём $e^{h_n}=(1-\frac z{a_n})e^{P^p_n}$.
По Следствию 34.3 (билет 34: если хвост произведения есть $\prod e^{h_n}$ с равномерно сходящимся $\sum h_n$, то произведение сходится
равномерно к $(\text{начальное конечное произведение})\cdot e^{\sum_{n>N}h_n}$) получаем локально равномерную сходимость, голоморфность
предела, и нули предела --- нули множителей; $n$-й множитель имеет единственный простой нуль $a_n$.
\end{proof}

\subsection{Оценка тейлоровских коэффициентов через вещественную часть}

\begin{thmbox}{Лемма 16.6}
Пусть $f\in H(B_R)$, $f(z)=\sum_{n\ge0}b_nz^n$, и $\Re f(z)\le A$ при $z\in B_R$. Тогда
$|b_n|\le\dfrac{2(A-\Re f(0))}{R^n}$ при всех $n\in\NN$.
\end{thmbox}
\begin{proof}
Пусть $g=A-f$; тогда $\Re g\ge0$ в $B_R$, $g(z)=(A-b_0)-\sum_{n\ge1}b_nz^n$. Фиксируем $n\ge1$, $r\in(0,R)$. Ряд $g(re^{i\vartheta})=\sum_k c_kr^ke^{ik\vartheta}$
($c_0=A-b_0$, $c_k=-b_k$) сходится равномерно по $\vartheta$, и почленное интегрирование даёт
\[
\int_0^{2\pi}g(re^{i\vartheta})e^{-in\vartheta}d\vartheta=2\pi c_nr^n=-2\pi b_nr^n,\qquad
\int_0^{2\pi}g(re^{i\vartheta})e^{in\vartheta}d\vartheta=\sum_kc_kr^k\int_0^{2\pi}e^{i(k+n)\vartheta}d\vartheta=0
\]
(во втором $k+n\ge1$). Сопрягая второе равенство: $\int_0^{2\pi}\clos{g(re^{i\vartheta})}e^{-in\vartheta}d\vartheta=0$. Складывая с первым и
учитывая $g+\bar g=2\Re g$:
\[
-b_n=\frac1{\pi r^n}\int_0^{2\pi}\Re g(re^{i\vartheta})\,e^{-in\vartheta}\,d\vartheta.
\]
Так как $\Re g\ge0$, $|b_n|\le\frac1{\pi r^n}\int_0^{2\pi}\Re g(re^{i\vartheta})d\vartheta=\frac{2\Re g(0)}{r^n}=\frac{2(A-\Re f(0))}{r^n}$ (теорема о среднем для
$\Re g$). Устремляя $r\to R$, получаем утверждение.
\end{proof}

\begin{thmbox}{Предложение 12.5 (лемма Бореля--Каратеодори)}
Пусть $f$ голоморфна в окрестности замкнутого круга $\clos{B(0,R)}$, $M(r)=\max_{|z|=r}|f|$, $A(r)=\max_{|z|=r}\Re f$. Тогда при
$0<r<R$
\[
M(r)\le\frac{R+r}{R-r}|f(0)|+\frac{2r}{R-r}A(R).
\]
\end{thmbox}
\begin{proof}
\emph{Случай $f\equiv c$:} $\frac{R+r}{R-r}|c|+\frac{2r}{R-r}\Re c\ge\frac{R+r-2r}{R-r}|c|=|c|=M(r)$.

\emph{Случай $f(0)=0$, $f\ne\mathrm{const}$.} (1) $A(R)>0$: $u=\Re f$ гармонична, $\frac1{2\pi}\int_0^{2\pi}u(Re^{i\vartheta})d\vartheta=u(0)=0$; если бы
$u\le0$ на окружности $|z|=R$, то $u\equiv0$ на ней (непрерывная неположительная функция с нулевым средним), и по принципу максимума
для гармонических функций (применённому к $u$ и $-u$) $u\equiv0$ в круге, откуда по условиям Коши--Римана $f\equiv\mathrm{const}$ ---
противоречие. Значит $A:=A(R)>0$. (2) По принципу максимума для $u$: $u\le A$ во всём круге $|z|\le R$, поэтому
$\Re(2A-f)\ge A>0$, знаменатель ниже не обращается в нуль, и
\[
\varphi(z)=\frac{f(z)}{2A-f(z)}\in H\big(\clos{B(0,R)}\big),\qquad\varphi(0)=0.
\]
(3) При $f=u+iv$: $|\varphi|^2=\frac{u^2+v^2}{(2A-u)^2+v^2}\le1$, так как $(2A-u)^2-u^2=4A(A-u)\ge0$. (4) По лемме Шварца (для круга
радиуса $R$, т.е. для $\varphi(Rz)$ в $\DD$) $|\varphi(z)|\le|z|/R$. (5) Из $\varphi(2A-f)=f$ находим $f=\frac{2A\varphi}{1+\varphi}$; при $|z|=r$
\[
|f(z)|\le\frac{2A|\varphi(z)|}{1-|\varphi(z)|}\le\frac{2A\,r/R}{1-r/R}=\frac{2Ar}{R-r}
\]
(функция $x\mapsto\frac x{1-x}$ возрастает на $[0,1)$).

\emph{Общий случай.} $h=f-f(0)$, $h(0)=0$, $A_h(R)=A(R)-\Re f(0)\le A(R)+|f(0)|$. По предыдущему (если $h\ne\mathrm{const}$)
$M(r)\le M_h(r)+|f(0)|\le\frac{2r}{R-r}\big(A(R)+|f(0)|\big)+|f(0)|=\frac{2r}{R-r}A(R)+\frac{R+r}{R-r}|f(0)|$.
\end{proof}

\subsection{Теорема Адамара}

\begin{thmbox}{Теорема 16.3 (Адамар)}
Пусть $f\not\equiv0$ --- целая функция конечного порядка $\rho$, $m$ --- кратность нуля $f$ в точке $0$, $\{a_n\}$ --- её ненулевые
нули, $p\in\ZZ_+$ --- наименьшее целое, для которого ряд $\sum_n|a_n|^{-(p+1)}$ сходится. Тогда
\begin{equation}\label{t36:1616}
f(z)=z^me^{g(z)}\prod_{n=1}^\infty\Big(1-\frac z{a_n}\Big)\exp\big(P^p_n(z)\big),\tag{16.16}
\end{equation}
где $g$ --- многочлен степени не выше $\rho$. Иначе говоря, $f=e^gF$, где $F$ --- каноническое произведение, построенное по нулям $f$.
\end{thmbox}

\begin{proof}
\emph{Шаг 1: существование $p$ и представление $f=e^gF$.} По лемме Адамара о нулях (Предложение 16.1) $\tau\le\rho<\infty$; по Лемме 36.1
род $p$ существует и $p\le\tau\le\rho$. (Если ненулевых нулей конечное число, $p=0$ и произведение конечно.) По Лемме 16.5
$F(z)=z^m\prod_n(1-\frac z{a_n})e^{P^p_n(z)}$ --- целая функция с теми же нулями и кратностями, что и $f$. Значит $f/F$ после устранения
особенностей --- целая функция без нулей, и по Лемме 12.3 $f/F=e^{g}$, $g$ --- целая. Остаётся доказать, что $c_k(g)=0$ при всех целых $k>\rho$.

\emph{Шаг 2: разложение $f$ при данном $R$.} Возьмём $R\ge\frac12$ и пусть $N=N(R)$ --- число $n$ с $|a_n|\le R$; тогда $|a_n|\le R$ при $n\le N$
и $|a_n|>R$ при $n>N$. Запишем
\[
f(z)=z^m\prod_{n=1}^N\Big(1-\frac z{a_n}\Big)\cdot F_N(z),\qquad
F_N(z)=\exp\Big(g(z)+\sum_{n=1}^NP^p_n(z)\Big)\prod_{n=N+1}^\infty\Big(1-\frac z{a_n}\Big)e^{P^p_n(z)}.
\]
$F_N$ --- целая функция (по Лемме 16.5 хвост произведения сходится локально равномерно).

\emph{Шаг 3: оценка $|F_N|$.} При $|z|=2R$ и $n\le N$: $\big|1-\frac z{a_n}\big|\ge\frac{|z|}{|a_n|}-1\ge\frac{2R}R-1=1$, и $|z|^m=(2R)^m\ge1$
(здесь используется $R\ge\frac12$). Поэтому $|F_N(z)|\le|f(z)|\le M_f(2R)$ при $|z|=2R$, и по принципу максимума модуля
\[
|F_N(z)|\le M_f(2R)\qquad\text{при }|z|\le2R.
\]

\emph{Шаг 4: логарифм $F_N$ в $B_R$.} При $z\in B_R$ и $n>N$ $|z/a_n|<1$, и определена голоморфная в $B_R$ функция
$h_n(z)=\operatorname{Log}(1-\frac z{a_n})+P^p_n(z)=-\sum_{k>p}\frac{z^k}{ka_n^k}$, $e^{h_n}=(1-\frac z{a_n})e^{P^p_n}$. Ряд $\sum_{n>N}h_n$
сходится равномерно в $B_R$: для $n$ с $|a_n|>2R$ (все, кроме конечного числа) $|h_n|\le2R^{p+1}|a_n|^{-(p+1)}$ (как в Лемме 16.5).
По Лемме 34.2 хвост произведения равен $\exp\big(\sum_{n>N}h_n\big)$, т.е.
\[
F_N=e^{g_N}\ \text{в } B_R,\qquad g_N(z)=g(z)+\sum_{n=1}^NP^p_n(z)+\sum_{n=N+1}^\infty h_n(z)\in H(B_R).
\]
Из Шага 3: $\Re g_N(z)=\ln|F_N(z)|\le\ln M_f(2R)$ при $z\in B_R$. Кроме того, $P^p_n(0)=0$, $h_n(0)=0$, так что
\[
g_N(0)=g(0)\qquad\text{(не зависит от $R$)}.
\]

\emph{Шаг 5: коэффициенты.} Пусть $k>p$ --- целое. Тогда $c_k(P^p_n)=0$ ($\deg P^p_n\le p$), $c_k(h_n)=-\frac1{ka_n^k}$, а по теореме
Вейерштрасса о почленном дифференцировании равномерно сходящегося ряда голоморфных функций $c_k\big(\sum_{n>N}h_n\big)=\sum_{n>N}c_k(h_n)$.
Итак,
\[
c_k(g_N)=c_k(g)-\sum_{n=N+1}^\infty\frac1{k\,a_n^k}.
\]
По Лемме 16.6 (в круге $B_R$ с $A=\ln M_f(2R)$) $|c_k(g_N)|\le2\big(\ln M_f(2R)-\Re g(0)\big)R^{-k}$. Следовательно,
\begin{equation}\label{t36:ck}\tag{36.1}
|c_k(g)|\le\frac{2\ln M_f(2R)+2|\Re g(0)|}{R^k}+\frac1k\sum_{n>N(R)}\frac1{|a_n|^k}.
\end{equation}

\emph{Шаг 6: предельный переход $R\to\infty$.} Пусть $k>\rho$ --- целое (тогда $k>\rho\ge p$, так что Шаг 5 применим). Выберем $\eps>0$ с
$\rho+\eps<k$. При больших $R$ $\ln M_f(2R)<(2R)^{\rho+\eps}$, поэтому первое слагаемое в \eqref{t36:ck} не больше
$2^{1+\rho+\eps}R^{\rho+\eps-k}+2|\Re g(0)|R^{-k}\to0$. Второе: так как $k>\rho\ge\tau$, ряд $\sum_n|a_n|^{-k}$ сходится; $N(R)\to\infty$ при
$R\to\infty$ (если нулей бесконечно много; иначе при больших $R$ хвост пуст), поэтому его остаток $\sum_{n>N(R)}|a_n|^{-k}\to0$.
Левая часть \eqref{t36:ck} от $R$ не зависит, значит $c_k(g)=0$. Итак, $c_k(g)=0$ при всех целых $k>\rho$: $g$ --- многочлен степени
$\le[\rho]\le\rho$.
\end{proof}

\begin{remarkbox}{Замечания к доказательству на слайдах 453--456}
\begin{enumerate}
\item На слайде 456 в оценке остаётся величина $\Re g_N(0)$, зависящая, вообще говоря, от $R$; для предельного перехода нужно, чтобы она
  не уходила в $-\infty$. Это так, потому что $g_N(0)=g(0)$ (Шаг 4); при $m=0$ $\Re g(0)=\ln|f(0)|$.
\item Условие $R>\frac12$ на слайде нужно ровно для $(2R)^m\ge1$ в Шаге 3.
\item Сходимость ряда $\sum_{n>N}h_n$ в $B_R$: для $n$ с $R<|a_n|\le2R$ (их конечное число) оценка Леммы 16.5 не работает, но
  каждая такая $h_n$ голоморфна в $B_R$, а конечное число слагаемых на сходимость не влияет.
\end{enumerate}
\end{remarkbox}

\begin{remarkbox}{Вариант Шага 5 через лемму Бореля--Каратеодори}
Вместо Леммы 16.6 можно оценить сам $g_N$. Для $R'<R$ функция $g_N$ голоморфна около $\clos{B(0,R')}$ и $A_{g_N}(R')\le A:=\ln M_f(2R)$
(при больших $R$ $A\ge0$). По Предложению 12.5 с $r=R/2$: $M_{g_N}(R/2)\le\frac{2r}{R'-r}A+\frac{R'+r}{R'-r}|g(0)|$; при $R'\to R$ получаем
$M_{g_N}(R/2)\le2A+3|g(0)|$. По неравенствам Коши $|c_k(g_N)|\le(2A+3|g(0)|)(2/R)^k$, и при $k>\rho+\eps$
$(2(2R)^{\rho+\eps}+3|g(0)|)2^kR^{-k}\to0$; дальше --- как в Шаге 6.
\end{remarkbox}

\begin{remarkbox}{Единственность представления}
Многочлен $g$ в \eqref{t36:1616} определён однозначно с точностью до слагаемого $2\pi i\ell$, $\ell\in\ZZ$: если $z^me^{g}F=z^me^{\tilde g}F$, то
$e^{g-\tilde g}=1$ вне дискретного множества нулей, по непрерывности --- везде; непрерывная функция $g-\tilde g$ со значениями в $2\pi i\ZZ$
на связном $\CC$ постоянна. В частности, $\deg g$ определяется функцией $f$.
\end{remarkbox}

\begin{remarkbox}{Примеры и следствия}
\begin{enumerate}
\item Если $f$ конечного порядка не имеет нулей, то $f=e^{g}$, $g$ --- многочлен степени $\le\rho$; значит порядок такой $f$ ---
  целое число $\deg g$ (билет 33: $e^{P}$ имеет порядок $\deg P$). Например, не существует целой функции порядка $\frac12$ без нулей.
\item Если $0\le\rho<1$, то $\tau<1$, $p=0$, $\deg g=0$: $f(z)=Cz^m\prod(1-z/a_n)$. Например,
  $\cos\sqrt z=\prod_{k\ge1}\big(1-\frac{z}{\pi^2(k-\frac12)^2}\big)$ (константа $C=f(0)=1$).
\item $\sin z=z\prod_{n\ge1}(1-\frac{z^2}{\pi^2n^2})$ (билет 37).
\item (без доказательства) $\dfrac1{\Gamma(z)}=ze^{\gamma z}\prod_{n\ge1}\big(1+\frac zn\big)e^{-z/n}$, $\gamma$ --- постоянная Эйлера: здесь
  $a_n=-n$, $p=1$, $\deg g=1$.
\end{enumerate}
\end{remarkbox}

\begin{remarkbox}{Коротко: план ответа на экзамене}
\begin{enumerate}
\item Определения: показатель сходимости, род $p$ ($p\le\tau\le p+1$), $P^p_n$, каноническое произведение (Опр.~16.4).
\item Лемма 16.5: при $\sum|a_n|^{-(p+1)}<\infty$ произведение сходится (оценка $|\operatorname{Log}(1-z/a_n)+P_n^p|\le2R^{p+1}/|a_n|^{p+1}$).
\item Формулировка теоремы Адамара: $f=z^me^{g}\prod(\dots)$, $\deg g\le\rho$.
\item $\tau\le\rho$ (лемма Адамара) $\Rightarrow$ $p$ существует; $f/F=e^g$.
\item Отщепить нули в $|a_n|\le R$: $|F_N|\le M_f(2R)$ в круге $B_{2R}$; $F_N=e^{g_N}$ в $B_R$, $\Re g_N\le\ln M_f(2R)$.
\item Лемма 16.6: $|c_k(g_N)|\le2(A-\Re g_N(0))/R^k$ (или Борель--Каратеодори + Коши).
\item $c_k(g_N)=c_k(g)-\sum_{n>N}\frac1{ka_n^k}$ при $k>p$; при $k>\rho$ оба слагаемых оценки $\to0$, значит $c_k(g)=0$.
\end{enumerate}
\end{remarkbox}

% ======================================================================
\ticket{37}{Теорема Бореля. Разложение функции $\sin z$ в бесконечное произведение}

\begin{defs}
\begin{itemize}
\item \textbf{Порядок} $\rho$ и \textbf{тип} $\sigma$ целой функции (Определения 16.1, 16.2; билет 33):
  $\rho=\inf\{\alpha>0:\ \ln M_f(r)<r^\alpha\ \text{при больших }r\}$, при $0<\rho<\infty$
  $\sigma=\inf\{s>0:\ \ln M_f(r)<sr^\rho\ \text{при больших }r\}$. Если $\ln M_f(r)\le Br^\kappa$ при больших $r$ ($B,\kappa>0$),
  то $\rho\le\kappa$ (так как $Br^\kappa<r^{\kappa+\eps}$ при больших $r$), а если $\rho=\kappa$, то $\sigma\le B<\infty$.
\item \textbf{Показатель сходимости} $\tau$ последовательности $0<|a_n|\le|a_{n+1}|\to\infty$: $\inf\{\lambda>0:\sum|a_n|^{-\lambda}<\infty\}$
  ($\tau=0$ для конечной последовательности). \textbf{Род} $p$ --- наименьшее $p\in\ZZ_+$ с $\sum|a_n|^{-(p+1)}<\infty$;
  $p\le\tau\le p+1$, причём $p<\tau$, если $\sum|a_n|^{-\tau}<\infty$, и $\tau<p+1$, если этот ряд расходится (Лемма 36.1, билет 36).
\item \textbf{Первичный множитель}: $E_p(u)=(1-u)\exp\big(\sum_{j=1}^pu^j/j\big)$ ($E_0(u)=1-u$).
\item \textbf{Каноническое произведение} рода $p$ (Определение 16.4): $F(z)=z^m\prod_{n\ge1}E_p(z/a_n)
  =z^m\prod_{n\ge1}(1-\frac z{a_n})\exp(P^p_n(z))$, $P^p_n(z)=\sum_{j=1}^p\frac{z^j}{ja_n^j}$; оно сходится и задаёт целую функцию с нулями
  $0$ (кратности $m$) и $a_n$ (Лемма 16.5).
\item \textbf{Теорема Адамара} (Теорема 16.3, билет 36): целая $f$ конечного порядка $\rho$ представляется в виде
  \[
  f(z)=z^me^{g(z)}\prod_{n=1}^\infty E_p\Big(\frac z{a_n}\Big),\tag{16.16}
  \]
  где $p$ --- род последовательности нулей, $g$ --- многочлен, $\deg g\le\rho$; $g$ определён однозначно с точностью до $2\pi i\ell$.
  \textbf{Лемма Адамара} (Предложение 16.1, билет 35): $\tau\le\rho$.
\end{itemize}
\end{defs}

\subsection{Оценка первичного множителя}

\begin{thmbox}{Лемма 37.1}
Пусть $p\in\ZZ_+$, $\lambda\in(p,p+1]$. Существует $C=C(p,\lambda)>0$ такая, что $|E_p(u)|\le\exp(C|u|^\lambda)$ при всех $u\in\CC$.
Можно взять $C=\max\{2,(p+1)2^p\}$ при $p\ge1$ и $C=\max\{2,1/\lambda\}$ при $p=0$.
\end{thmbox}
\begin{proof}
\emph{$|u|\le\frac12$.} Тогда $\operatorname{Log}(1-u)=-\sum_{j\ge1}u^j/j$ и $E_p(u)=\exp\big(-\sum_{j>p}u^j/j\big)$, так что
\[
|E_p(u)|\le\exp\Big(\sum_{j>p}|u|^j\Big)=\exp\Big(\frac{|u|^{p+1}}{1-|u|}\Big)\le\exp\big(2|u|^{p+1}\big)\le\exp\big(2|u|^\lambda\big),
\]
так как $|u|<1$ и $\lambda\le p+1$.

\emph{$|u|\ge\frac12$, $p\ge1$.} $|E_p(u)|\le(1+|u|)\exp\big(\sum_{j=1}^p|u|^j\big)$, $\ln(1+|u|)\le|u|$. При $1\le j\le p$
$|u|^j=|u|^p|u|^{j-p}\le2^{p-j}|u|^p\le2^{p-1}|u|^p$. Значит $\ln|E_p(u)|\le|u|+\sum_{j=1}^p|u|^j\le(p+1)2^{p-1}|u|^p$. Наконец,
$|u|^p=|u|^\lambda|u|^{p-\lambda}\le2^{\lambda-p}|u|^\lambda\le2|u|^\lambda$. Итого $\ln|E_p(u)|\le(p+1)2^p|u|^\lambda$.

\emph{$|u|\ge\frac12$, $p=0$.} $|E_0(u)|\le1+|u|$, и $\ln(1+x)\le x^\lambda/\lambda$ при $x\ge0$, $\lambda\in(0,1]$: действительно,
$t^{1-\lambda}\le\max(1,t)\le1+t$, т.е. $\frac1{1+t}\le t^{\lambda-1}$, и $\ln(1+x)=\int_0^x\frac{dt}{1+t}\le\int_0^xt^{\lambda-1}dt=\frac{x^\lambda}\lambda$.
\end{proof}

\begin{remarkbox}{Замечание к слайду 458}
На слайде оценка при $|u|\ge\frac12$ проведена в виде $\ln(1+|u|)+|u|^p\sum_{j<p}|u|^{-j}\le C_12^p|u|^p$; при $p=0$ это неверно
($\ln(1+|u|)$ не ограничен), поэтому случай $p=0$ разобран отдельно, и для этого существенно, что $\lambda>p$. Это условие
всегда выполнимо: см. выбор $\lambda$ ниже.
\end{remarkbox}

\subsection{Теорема Бореля}

\begin{thmbox}{Теорема 16.4 (Борель)}
Пусть $f$ имеет вид (16.16): $f(z)=z^me^{g(z)}\prod_{n\ge1}E_p(z/a_n)$, где $\{a_n\}$, $a_n\ne0$, --- последовательность с показателем
сходимости $\tau$ и родом $p$, а $g$ --- многочлен степени $k$. Тогда $f$ имеет конечный порядок
\[
\rho=\max\{\tau,k\}.
\]
Более того, если $\tau<k$ или если $\sum_n|a_n|^{-\tau}<\infty$ (и $\rho>0$), то $f$ имеет конечный тип относительно порядка $\rho$.
\end{thmbox}

\begin{proof}
Пусть $g(z)=q_0+q_1z+\dots+q_kz^k$, $q_k\ne0$ (при $k=0$ $g\equiv q_0$), $K=\sum_{j<k}|q_j|$. При $|z|=r\ge1$
\begin{equation}\label{t37:g}\tag{37.1}
\ln|z^me^{g(z)}|=m\ln r+\Re g(z)\le m\ln r+|q_k|r^k+Kr^{k-1}.
\end{equation}

\emph{Случай конечного числа нулей.} Тогда $\tau=0$, $\prod E_p(z/a_n)$ --- многочлен $Q$, $f=z^mQ\,e^{g}$. Функция $e^{g}$ имеет порядок $k$ и тип
$|q_k|$ (билет 33), а умножение на многочлен не меняет порядка и типа (задача со слайда 432, билет 33). Итак, $\rho=k=\max\{\tau,k\}$, и при
$k\ge1$ тип конечен. Далее нулей бесконечно много; тогда $\tau<\infty$ (род $p$ существует, $\tau\le p+1$).

\emph{Шаг 1: выбор $\lambda$.} Если $\sum|a_n|^{-\tau}<\infty$, то по Лемме 36.1 $p<\tau\le p+1$; положим $\lambda=\tau$. Если ряд расходится, то
$p\le\tau<p+1$; положим $\lambda=\tau+\eps$, где $\eps\in(0,p+1-\tau]$ --- любое. В обоих случаях $\lambda\in(p,p+1]$ и
$S_\lambda:=\sum_n|a_n|^{-\lambda}<\infty$ (при $\lambda=\tau+\eps>\tau$ --- по определению $\tau$).

\emph{Шаг 2: оценка сверху.} По Лемме 37.1 $|E_p(z/a_n)|\le\exp\big(C|z|^\lambda/|a_n|^\lambda\big)$ для всех $z$ и $n$, поэтому для частичных
произведений $\prod_{n\le N}|E_p(z/a_n)|\le\exp(CS_\lambda|z|^\lambda)$, и то же верно для предела. Вместе с \eqref{t37:g}: при $r\ge1$
\begin{equation}\label{t37:1618}
\ln M_f(r)\le|q_k|r^k+Kr^{k-1}+m\ln r+CS_\lambda r^\lambda.\tag{16.18}
\end{equation}
Пусть $\kappa=\max\{k,\lambda\}>0$. Так как $\ln r\le\frac{r^\kappa}{\kappa}$ (из $\ln x\le x-1$ при $x=r^\kappa$), то
$\ln M_f(r)\le\big(|q_k|+K+\frac m\kappa+CS_\lambda\big)r^\kappa$ при $r\ge1$. Значит $f$ имеет конечный порядок и $\rho\le\max\{k,\lambda\}$.
Если $\lambda=\tau$, это $\rho\le\max\{k,\tau\}$; если $\lambda=\tau+\eps$, то $\rho\le\max\{k,\tau+\eps\}$ при всех малых $\eps>0$, т.е. снова
$\rho\le\max\{k,\tau\}$.

\emph{Шаг 3: оценка снизу.} Раз $f$ конечного порядка, по лемме Адамара о нулях $\tau\le\rho$. По теореме Адамара
$f=z^me^{\tilde g}\prod E_p(z/a_n)$ с многочленом $\deg\tilde g\le\rho$ (каноническое произведение то же: оно определяется нулями $f$).
Сравнивая с исходным представлением (16.16), получаем $e^{g-\tilde g}\equiv1$, откуда $g-\tilde g$ --- константа (замечание о единственности в билете 36),
и $k=\deg g=\deg\tilde g\le\rho$ (при $k=0$ доказывать нечего). Итак, $\rho=\max\{k,\tau\}$.

\emph{Шаг 4: тип при $\tau<k$.} Тогда $\rho=k\ge1$. Выберем $\lambda$ как в Шаге 1, но дополнительно $\lambda<k$: если $\sum|a_n|^{-\tau}<\infty$,
то $\lambda=\tau<k$; иначе $\eps<k-\tau$. В \eqref{t37:1618} слагаемые $Kr^{k-1}$, $m\ln r$, $CS_\lambda r^\lambda$ суть $o(r^k)$, поэтому при больших
$r$ $\ln M_f(r)<(|q_k|+\eta)r^k$ для любого $\eta>0$: тип $\sigma\le|q_k|<\infty$.

\emph{Шаг 5: тип при $k\le\tau$ и $\sum|a_n|^{-\tau}<\infty$.} Тогда $\rho=\tau=\lambda>0$, и \eqref{t37:1618} с $k\le\tau$ даёт при $r\ge1$
$\ln M_f(r)\le\big(|q_k|+K+\frac m\tau+CS_\tau\big)r^\tau$, т.е. $\sigma\le|q_k|+K+\frac m\tau+CS_\tau<\infty$.
\end{proof}

\begin{remarkbox}{Замечания}
\begin{enumerate}
\item На слайде 459 оценка $|z^me^{g(z)}|<\exp\big((|q_k|+\eps)|z|^k\big)$ при больших $|z|$ верна лишь при $k\ge1$; при $k=0$ множитель $|z|^m$
  не ограничен. В доказательстве выше это учтено слагаемым $m\ln r$ (оно $o(r^\kappa)$, так как $\kappa\ge\lambda>0$).
\item Условия конечности типа достаточны, но не необходимы: у $\sin z$ $\tau=k=1$, ряд $\sum|a_n|^{-1}$ расходится, а тип равен $1$. Если же
  они нарушены, тип может быть бесконечным: $1/\Gamma(z)=ze^{\gamma z}\prod(1+\frac zn)e^{-z/n}$ имеет порядок $1$ и максимальный тип
  (без доказательства).
\end{enumerate}
\end{remarkbox}

\subsection{Следствия теоремы Бореля}

\begin{thmbox}{Следствие 37.2 (порядок канонического произведения)}
Порядок канонического произведения равен показателю сходимости его нулей; при $\sum|a_n|^{-\tau}<\infty$ (и $\tau>0$) тип конечен.
\end{thmbox}
\begin{proof}
Теорема 16.4 с $g\equiv0$, $k=0$: $\rho=\max\{\tau,0\}=\tau$.
\end{proof}

\begin{thmbox}{Следствие 37.3 (слайд 460)}
Если $f$ --- целая функция конечного порядка $\rho\notin\ZZ$, то $\rho=\tau$.
\end{thmbox}
\begin{proof}
По теореме Адамара $f$ имеет вид (16.16) с многочленом $g$ степени $k\in\ZZ_+$; по теореме Бореля $\rho=\max\{\tau,k\}$. Так как $\rho\ne k$,
то $\rho=\tau$.
\end{proof}

\begin{thmbox}{Следствие 37.4 (слайд 460)}
Всякая целая функция $f\not\equiv0$ порядка $\rho\in[0,1)$ представляется в виде $f(z)=Cz^m\prod_{n\ge1}\big(1-\frac z{a_n}\big)$, $C\ne0$.
\end{thmbox}
\begin{proof}
По лемме Адамара $\tau\le\rho<1$, поэтому $\sum|a_n|^{-1}<\infty$ (так как $1>\tau$) и род $p=0$, $P^0_n\equiv0$. По теореме Адамара $\deg g\le\rho<1$, т.е.
$g\equiv\mathrm{const}$, $C=e^{g}$.
\end{proof}

\begin{thmbox}{Следствие 37.5}
(а) Целая функция нецелого порядка имеет бесконечно много нулей; более того, для любого $a\in\CC$ уравнение $f(z)=a$ имеет бесконечно много
решений. (б) Если у целой функции конечного порядка $\tau<\rho$, то $\rho$ --- целое число, равное $\deg g$ в (16.16).
\end{thmbox}
\begin{proof}
(а) Если бы нулей было конечное число, то $\tau=0$ и по Следствию 37.3 $\rho=0\in\ZZ$ --- противоречие. Функция $f-a$ имеет тот же порядок, что
$f$ (задача со слайда 432, билет 33; $f\ne\mathrm{const}$, так как порядок нецелый), поэтому и у неё бесконечно много нулей.
(б) $\rho=\max\{\tau,k\}$ и $\rho\ne\tau$, значит $\rho=k$.
\end{proof}

\subsection{Разложение $\sin z$ и $\cos z$ в бесконечные произведения}

\begin{thmbox}{Теорема 37.6}
При всех $z\in\CC$
\[
\sin z=z\prod_{n=1}^\infty\Big(1-\frac{z^2}{\pi^2n^2}\Big),\qquad
\cos z=\prod_{n=1}^\infty\Big(1-\frac{z^2}{\pi^2(n-\frac12)^2}\Big),
\]
произведения сходятся локально равномерно в $\CC$.
\end{thmbox}
\begin{proof}
\emph{1. Порядок.} $\sin z$ имеет порядок $1$: $\operatorname{sh}r=|\sin(ir)|\le M(r)\le\operatorname{ch}r$ (билет 33).

\emph{2. Нули.} $\sin z=\frac{e^{iz}-e^{-iz}}{2i}=0\iff e^{2iz}=1\iff2iz\in2\pi i\ZZ\iff z=\pi\ell$, $\ell\in\ZZ$. Все нули простые: $(\sin)'(\pi\ell)=\cos\pi\ell=(-1)^\ell\ne0$.
Итак, $m=1$, а ненулевые нули в стандартном порядке: $a_{2n-1}=\pi n$, $a_{2n}=-\pi n$, $n\in\NN$.

\emph{3. Показатель сходимости и род.} $\sum_j|a_j|^{-\lambda}=2\pi^{-\lambda}\sum_{n\ge1}n^{-\lambda}$ сходится $\iff\lambda>1$. Значит $\tau=1$, ряд при $\lambda=1$
расходится, при $\lambda=2$ сходится: род $p=1$, $P^1_j(z)=z/a_j$.

\emph{4. Теорема Адамара.} $\sin z=ze^{g(z)}\prod_{j\ge1}\big(1-\frac z{a_j}\big)e^{z/a_j}$, где $\deg g\le\rho=1$, т.е. $g(z)=az+b$.

\emph{5. Группировка множителей.} Частичное произведение с номером $2N$ равно
\[
\prod_{n=1}^N\Big(1-\frac z{\pi n}\Big)e^{\frac z{\pi n}}\Big(1+\frac z{\pi n}\Big)e^{-\frac z{\pi n}}=\prod_{n=1}^N\Big(1-\frac{z^2}{\pi^2n^2}\Big).
\]
Последовательность частичных произведений сходится (локально равномерно), значит и её подпоследовательность с чётными номерами сходится к тому же
пределу. Поэтому
\[
\sin z=ze^{az+b}\,\Pi(z),\qquad \Pi(z)=\prod_{n=1}^\infty\Big(1-\frac{z^2}{\pi^2n^2}\Big),
\]
и $\Pi$ --- целая функция (локально равномерный предел многочленов). Она \emph{чётна} (все частичные произведения чётны) и $\Pi(0)=1$.

\emph{6. $e^b=1$.} При $z\to0$: $\frac{\sin z}z\to1$ (так как $\sin z=z-\frac{z^3}6+\dots$), а $e^{az+b}\Pi(z)\to e^b\Pi(0)=e^b$ (непрерывность). Значит $e^b=1$.

\emph{7. $a=0$.} Так как $\sin(-z)=-\sin z$: $-ze^{-az}\Pi(-z)=-ze^{az}\Pi(z)$, т.е. $z\Pi(z)\big(e^{az}-e^{-az}\big)=0$. При $z\notin\pi\ZZ$ $z\Pi(z)\ne0$
(это $\sin z\,e^{-az-b}$), поэтому $e^{2az}=1$ на $\CC\setminus\pi\ZZ$. По теореме единственности $e^{2az}\equiv1$; дифференцируя в нуле, $2a=0$.

Итак, $\sin z=z\Pi(z)$.

\emph{8. Косинус.} Аналогично (проводим все шаги): порядок $\cos z$ равен $1$ ($\operatorname{sh}r\le|\cos(ir)|=\operatorname{ch}r$, $M(r)\le\operatorname{ch}r$). Нули:
$\cos z=0\iff e^{2iz}=-1\iff z=\pi(\ell-\frac12)$, $\ell\in\ZZ$; простые, так как $(\cos)'=-\sin\ne0$ в этих точках; $m=0$; ненулевые нули $\pm\pi(n-\frac12)$,
$n\in\NN$. $\sum|a_j|^{-\lambda}=2\pi^{-\lambda}\sum_n(n-\frac12)^{-\lambda}<\infty\iff\lambda>1$, так что $\tau=1$, $p=1$. По теореме Адамара и группировке
\[
\cos z=e^{az+b}\prod_{n\ge1}\Big(1-\frac{z^2}{\pi^2(n-\frac12)^2}\Big)=e^{az+b}\Pi_c(z),\qquad\Pi_c\ \text{чётна},\ \Pi_c(0)=1.
\]
$z=0$: $1=e^b$. Чётность $\cos$: $e^{-az}\Pi_c(z)=e^{az}\Pi_c(z)$, $\Pi_c\ne0$ вне нулей $\cos$, откуда $e^{2az}\equiv1$ и $a=0$.
\end{proof}

\begin{remarkbox}{Следствия разложения (доп.)}
\begin{enumerate}
\item \emph{Формула Валлиса.} При $z=\frac\pi2$: $1=\frac\pi2\prod_{n\ge1}\big(1-\frac1{4n^2}\big)=\frac\pi2\prod_{n\ge1}\frac{(2n-1)(2n+1)}{(2n)^2}$, т.е.
  $\dfrac\pi2=\prod_{n\ge1}\dfrac{(2n)^2}{(2n-1)(2n+1)}$.
\item \emph{$\sum1/n^2=\pi^2/6$.} При $|z|\le1$ имеем $|z^2/(\pi^2n^2)|\le\pi^{-2}<\frac12$, и (Лемма 34.2, билет 34) $\Pi=e^{H}$,
  $H(z)=\sum_n\operatorname{Log}\big(1-\frac{z^2}{\pi^2n^2}\big)$, ряд сходится равномерно ($|\operatorname{Log}(1-w)|\le2|w|$ при $|w|\le\frac12$). По теореме
  Вейерштрасса коэффициенты $H$ получаются почленно: $H(z)=-\Big(\sum_n\frac1{\pi^2n^2}\Big)z^2+O(z^4)$. Тогда $\Pi(z)=e^{H(z)}=1-\Big(\sum_n\frac1{\pi^2n^2}\Big)z^2+O(z^4)$, а
  $\frac{\sin z}z=1-\frac{z^2}6+O(z^4)$. Сравнивая коэффициенты при $z^2$: $\sum_{n\ge1}\frac1{n^2}=\frac{\pi^2}6$.
\item \emph{Гиперболические функции и $e^z-1$.} $\operatorname{sh}w=-i\sin(iw)=w\prod_n\big(1+\frac{w^2}{\pi^2n^2}\big)$, и
  $e^z-1=2e^{z/2}\operatorname{sh}\frac z2=ze^{z/2}\prod_{n\ge1}\big(1+\frac{z^2}{4\pi^2n^2}\big)$ --- здесь $g(z)=z/2$, $k=1=\tau=\rho$.
\end{enumerate}
\end{remarkbox}

\begin{remarkbox}{Коротко: план ответа на экзамене}
\begin{enumerate}
\item Формулировка: для $f=z^me^{g}\prod E_p(z/a_n)$, $\deg g=k$: $\rho=\max\{\tau,k\}$; тип конечен, если $\tau<k$ или $\sum|a_n|^{-\tau}<\infty$.
\item Оценка $|E_p(u)|\le e^{C|u|^\lambda}$, $\lambda\in(p,p+1]$: при $|u|\le\frac12$ через $\exp(-\sum_{j>p}u^j/j)$, при $|u|\ge\frac12$ грубо.
\item Выбор $\lambda$: $\lambda=\tau$, если $\sum|a_n|^{-\tau}<\infty$, иначе $\lambda=\tau+\eps$.
\item $\ln M_f(r)\le|q_k|r^k+o(r^k)+CS_\lambda r^\lambda$ $\Rightarrow$ $\rho\le\max\{k,\tau\}$; обратно $\tau\le\rho$ (Адамар о нулях), $k\le\rho$ (теорема Адамара).
\item Следствия: порядок канонического произведения $=\tau$; $\rho\notin\ZZ\Rightarrow\rho=\tau$; $\rho<1\Rightarrow f=Cz^m\prod(1-z/a_n)$.
\item $\sin z$: порядок 1, нули $\pi\ell$ простые, $\tau=1$, $p=1$, $g=az+b$; группировка пар $\pm\pi n$; $e^b=1$ из $\frac{\sin z}z\to1$, $a=0$ из нечётности.
\item Ответ: $\sin z=z\prod(1-\frac{z^2}{\pi^2n^2})$; аналогично $\cos z=\prod(1-\frac{z^2}{\pi^2(n-1/2)^2})$.
\end{enumerate}
\end{remarkbox}

% ===== part 8 =====
% ===================== Часть 8: билеты 38, 39, 40 и вопрос 43* =====================

\ticket{38}{Плотность системы $\{e^{i\lambda_n t}\}$ в пространстве {$C[-\pi,\pi]$} в случае последовательности $\Lambda=\{\lambda_n\}_{n=1}^\infty$ с условием $\liminf_{t\to\infty} n(t)/t>2$, где $n(\cdot)$ --- считающая функция $\Lambda$}

\begin{defs}
\textbf{Максимум модуля, порядок, тип} (см. билет 33). Для целой функции (ц.ф.) $f$ полагаем $M_f(r)=\max_{|z|=r}|f(z)|$.
$f$ имеет \emph{конечный порядок}, если $M_f(r)<\exp(r^\alpha)$ при всех $r>R$ для некоторых $\alpha>0$, $R>0$; \emph{порядок}
$\rho$ --- нижняя грань таких $\alpha$; эквивалентно $\rho=\limsup_{r\to\infty}\frac{\ln\ln M_f(r)}{\ln r}$.
При $0<\rho<\infty$ \emph{тип} $\sigma$ относительно порядка $\rho$ --- нижняя грань чисел $s>0$ таких, что $M_f(r)<\exp(sr^\rho)$ при
всех больших $r$; эквивалентно $\sigma=\limsup_{r\to\infty}\frac{\ln M_f(r)}{r^\rho}$. В частности, \emph{для любого $\eps>0$ при всех
достаточно больших $r$} выполнено $\ln M_f(r)<(\sigma+\eps)r^\rho$ (если $\sigma<\infty$) и $\ln M_f(r)<r^{\rho+\eps}$.
Примеры: $e^z$ --- порядок $1$, тип $1$; $\sin z$ --- порядок $1$, тип $1$; $e^{z^2}$ --- порядок $2$, тип $1$.

\textbf{Последовательность нулей.} Нули ц.ф. $f\not\equiv0$, отличные от $0$, записываются с учётом кратности в последовательность
$\Lambda=\{a_n\}$, упорядоченную «стандартно»: $0<|a_1|\le|a_2|\le\dots$, $|a_n|\to\infty$.

\textbf{Показатель сходимости} (см. билет 35, Определение 16.3) последовательности $0<|a_n|\le|a_{n+1}|\to\infty$:
$\tau=\inf\{\lambda>0:\ \sum_n|a_n|^{-\lambda}<\infty\}$; по Лемме 16.3 $\tau=\limsup_{n}\frac{\ln n}{\ln|a_n|}$.
Для конечной последовательности полагаем $\tau=0$. Пример: $\{\pi k\}_{k\in\ZZ\setminus\{0\}}$ имеет $\tau=1$.

\textbf{Считающая функция} последовательности $\Lambda=\{a_n\}$: $n_\Lambda(t)=\#\{n:\ |a_n|<t\}$, $t>0$ (число членов в открытом
круге радиуса $t$, с учётом повторений). Для ц.ф. $f$ через $n_f(t)$ обозначается число нулей $f$ в круге $|z|<t$ (с учётом кратности).
В Предложении 16.4 для вещественной последовательности используется $n(t)=\#\{n:\ |\lambda_n|\le t\}$ --- число точек
$\Lambda$ на отрезке $[-t,t]$; разница между «$<$» и «$\le$» на асимптотику не влияет.

\textbf{Плотности} (Определение 16.5). Пусть $\rho>0$. Верхняя и нижняя $\rho$-плотности последовательности $\Lambda$:
\[
\Delta^\sharp_\rho(\Lambda)=\limsup_{n\to\infty}\frac{n}{|a_n|^\rho},\qquad \Delta^\flat_\rho(\Lambda)=\liminf_{n\to\infty}\frac{n}{|a_n|^\rho};
\]
если существует $\lim_n n/|a_n|^\rho$, он называется $\rho$-плотностью $\Delta_\rho(\Lambda)$.
Пример: для $a_n=n$ имеем $\Delta_1=1$; для нулей $\sin z$ ($\pm\pi,\pm2\pi,\dots$; $a_{2k-1},a_{2k}=\pm\pi k$) $\Delta_1=2/\pi$.

\textbf{Полнота системы.} Система $\{x_n\}$ элементов линейного топологического пространства $E$ называется \emph{полной}, если
замыкание её линейной оболочки совпадает с $E$, т.е. любой элемент $E$ сколь угодно хорошо приближается конечными линейными
комбинациями $x_n$. В названии билета «плотность системы» означает именно полноту (плотность линейной оболочки).

\textbf{Пространство $C[-\pi,\pi]$} --- комплекснозначные непрерывные функции на $[-\pi,\pi]$ с нормой $\|h\|=\max|h|$; это банахово
пространство. $E^*$ --- пространство непрерывных линейных функционалов.

\textbf{Комплексная борелевская мера} $\mu$ на $[-\pi,\pi]$ --- счётно-аддитивная функция на борелевских множествах со значениями в $\CC$;
её \emph{полная вариация} $|\mu|$ --- наименьшая положительная мера с $|\mu(B)|\le|\mu|(B)$; $\|\mu\|=|\mu|([-\pi,\pi])<\infty$.
Для ограниченной измеримой $h$: $\big|\int h\,d\mu\big|\le\int|h|\,d|\mu|\le\|\mu\|\sup|h|$.
\end{defs}

\subsection{Плотности последовательностей}

\begin{thmbox}{Лемма 16.7}
Пусть $\Lambda=\{a_n\}$, $0<|a_n|\le|a_{n+1}|$, $|a_n|\to\infty$, $\rho>0$, $n_\Lambda(t)=\#\{n:|a_n|<t\}$. Тогда
\[
\Delta^\sharp_\rho(\Lambda)=\limsup_{t\to\infty}\frac{n_\Lambda(t)}{t^\rho},\qquad
\Delta^\flat_\rho(\Lambda)=\liminf_{t\to\infty}\frac{n_\Lambda(t)}{t^\rho},\qquad
\Delta_\rho(\Lambda)=\lim_{t\to\infty}\frac{n_\Lambda(t)}{t^\rho},
\]
причём в последнем равенстве пределы существуют или не существуют одновременно (значения $+\infty$ допускаются).
\end{thmbox}
На слайдах доказательство оставлено как задача. Приведём его.
\begin{proof}
Положим $x_n=n/|a_n|^\rho$, $\phi(t)=n_\Lambda(t)/t^\rho$. Заметим, что $n_\Lambda(t)<\infty$ при всех $t$ (так как $|a_n|\to\infty$),
$n_\Lambda$ не убывает и $n_\Lambda(t)\to\infty$ при $t\to\infty$.

\emph{Два элементарных наблюдения.} (1) Если $t>|a_n|$, то $|a_1|\le\dots\le|a_n|<t$, поэтому $n_\Lambda(t)\ge n$.
(2) Если $N=n_\Lambda(t)$, то в силу монотонности $|a_k|$ члены с $|a_k|<t$ --- это ровно $a_1,\dots,a_N$; значит,
$|a_N|<t\le|a_{N+1}|$ (при $N\ge1$).

\textbf{$\limsup_t\phi\ge\limsup_n x_n$.} Для каждого $n$ по (1) при $t\downarrow|a_n|$ имеем $\phi(t)\ge n/t^\rho\to x_n$; выберем
$t_n\in(|a_n|,|a_n|+1)$ с $\phi(t_n)\ge x_n-\tfrac1n$ (если $x_n$ конечно, это возможно по непрерывности $n/t^\rho$). Так как
$t_n\to\infty$, получаем $\limsup_t\phi\ge\limsup_n(x_n-\tfrac1n)=\limsup_n x_n$.

\textbf{$\limsup_t\phi\le\limsup_n x_n$.} При $t>|a_1|$ имеем $N=n_\Lambda(t)\ge1$ и по (2) $|a_N|<t$, откуда
$\phi(t)=N/t^\rho\le N/|a_N|^\rho=x_{N}$. При $t\to\infty$ индекс $N=n_\Lambda(t)\to\infty$, поэтому
$\limsup_t\phi(t)\le\limsup_{N\to\infty}x_N$.

\textbf{$\liminf_t\phi\ge\liminf_n x_n$.} При $t>|a_1|$, $N=n_\Lambda(t)$, по (2) $t\le|a_{N+1}|$, значит
\[
\phi(t)=\frac{N}{t^\rho}\ge\frac{N}{|a_{N+1}|^\rho}=\frac{N}{N+1}\,x_{N+1}.
\]
При $t\to\infty$ имеем $N\to\infty$, $\frac N{N+1}\to1$, а $x_{N+1}$ пробегает члены последовательности со сколь угодно большими
номерами; поэтому $\liminf_t\phi\ge\liminf_n x_n$ (если $\liminf_n x_n=+\infty$, то $\phi(t)\ge\frac12x_{N+1}\to\infty$).

\textbf{$\liminf_t\phi\le\liminf_n x_n$.} При $t=|a_n|$ в круг $|z|<t$ не попадает $a_n$ и все следующие члены, поэтому
$n_\Lambda(|a_n|)\le n-1<n$ и $\phi(|a_n|)\le x_n$. Беря подпоследовательность $x_{n_j}\to\liminf_n x_n$, получаем
$\liminf_t\phi\le\lim_j\phi(|a_{n_j}|)\le\liminf_n x_n$.

Итак, верхние и нижние пределы совпадают. Предел $\lim_n x_n$ существует $\iff$ $\liminf_n x_n=\limsup_n x_n$ $\iff$
$\liminf_t\phi=\limsup_t\phi$ $\iff$ существует $\lim_t\phi(t)$, и тогда пределы равны.
\end{proof}

Следующие два утверждения на лекциях приведены без доказательства (последовательность нулей ц.ф. всегда упорядочена стандартным образом).

\begin{thmbox}{Предложение 16.2}
Пусть $f$ --- ц.ф. конечного порядка $\rho>0$ и типа $\sigma$ относительно этого порядка, $\Lambda$ --- последовательность
её нулей. Тогда
\[
\Delta^\sharp_\rho(\Lambda)\le e\rho\sigma,\qquad \Delta^\flat_\rho(\Lambda)\le\rho\sigma,
\]
и эти неравенства точные (не могут быть улучшены в классе всех таких $f$).
\end{thmbox}

Неравенства (без утверждения о точности) доказываются той же формулой Йенсена, что и в билете 35; приведём это
рассуждение (сверх лекций), так как оно хорошо иллюстрирует связь роста и плотности нулей. Сначала --- нужная версия формулы Йенсена.

\begin{thmbox}{Формула Йенсена (см. раздел 0) и её интегральная форма}
Пусть $G$ голоморфна в окрестности замкнутого круга $\{|z|\le r\}$, $G(0)\ne0$, на окружности $|z|=r$ нулей нет, и
$a_1,\dots,a_N$ --- нули $G$ в круге $|z|<r$ с учётом кратности. Тогда
\[
\ln|G(0)|=\sum_{k=1}^N\ln\frac{|a_k|}{r}+\frac1{2\pi}\int_0^{2\pi}\ln|G(re^{i\theta})|\,d\theta .
\]
Эквивалентно, если $n_G(s)$ --- число нулей $G$ в круге $|z|<s$, то
\begin{equation}\label{t38:jensen}
\int_0^r\frac{n_G(s)}{s}\,ds=\frac1{2\pi}\int_0^{2\pi}\ln|G(re^{i\theta})|\,d\theta-\ln|G(0)|.
\end{equation}
\end{thmbox}
\begin{detail}{переход к интегральной форме}
Так как $G(0)\ne0$, в некотором круге $|z|<s_0$ нулей нет и $n_G(s)=0$ при $s<s_0$; интеграл в \eqref{t38:jensen} собственный.
Для каждого нуля $a_k$, $|a_k|<r$: $\ln\frac{r}{|a_k|}=\int_{|a_k|}^r\frac{ds}{s}=\int_0^r\mathbf 1_{\{|a_k|<s\}}\frac{ds}s$.
Суммируя по $k=1,\dots,N$ и учитывая $\sum_k\mathbf 1_{\{|a_k|<s\}}=n_G(s)$ при $0<s\le r$, получаем
$\sum_k\ln\frac r{|a_k|}=\int_0^r\frac{n_G(s)}s\,ds$, что вместе с формулой Йенсена даёт \eqref{t38:jensen}.
Если вместо $n_G(s)$ взять число нулей в замкнутом круге $|z|\le s$, интеграл не изменится: эти функции различаются лишь в
конечном числе точек $s=|a_k|$.
\end{detail}

\begin{proof}[Доказательство неравенств Предложения 16.2]
Можно считать $\sigma<\infty$. Пусть $f$ имеет в нуле нуль кратности $m\ge0$; положим $G(z)=f(z)/z^m$. Тогда $G$ --- ц.ф.,
$G(0)\ne0$, её нули --- в точности $\Lambda$, т.е. $n_G=n_\Lambda$, и при $|z|=R\ge1$ имеем $|G(z)|\le|f(z)|$. По
\eqref{t38:jensen} (для $R\ge1$, не равных модулям нулей)
\begin{equation}\label{t38:J2}
\int_0^R\frac{n_\Lambda(s)}{s}\,ds\le\ln M_f(R)-\ln|G(0)|.
\end{equation}
Обе части \eqref{t38:J2} непрерывны по $R$ (левая --- как интеграл с переменным верхним пределом, правая --- так как $M_f$
непрерывна), а исключённые $R$ образуют дискретное множество, поэтому \eqref{t38:J2} верно при всех $R\ge1$.
Фиксируем $\eps>0$; при всех больших $R$ имеем $\ln M_f(R)<(\sigma+\eps)R^\rho$.

\emph{Верхняя плотность.} Пусть $k>1$. Так как $n_\Lambda$ не убывает,
\[
n_\Lambda(r)\ln k=\int_r^{kr}\frac{n_\Lambda(r)}{s}\,ds\le\int_r^{kr}\frac{n_\Lambda(s)}{s}\,ds\le\int_0^{kr}\frac{n_\Lambda(s)}{s}\,ds
\le(\sigma+\eps)k^\rho r^\rho-\ln|G(0)|
\]
при больших $r$. Делим на $r^\rho\ln k$ и переходим к $\limsup$: $\limsup_r n_\Lambda(r)/r^\rho\le(\sigma+\eps)\,k^\rho/\ln k$.
Минимизируем $\psi(k)=k^\rho/\ln k$ по $k>1$: $\psi'(k)=k^{\rho-1}(\rho\ln k-1)/\ln^2k=0$ при $\ln k=1/\rho$, т.е. $k=e^{1/\rho}$,
$\psi(e^{1/\rho})=e\rho$. Устремляя $\eps\to0$ и применяя Лемму 16.7, получаем $\Delta^\sharp_\rho(\Lambda)\le e\rho\sigma$.

\emph{Нижняя плотность.} Допустим, $\Delta^\flat_\rho(\Lambda)>\rho\sigma$. По Лемме 16.7 найдутся $\delta>0$ и $T$ такие, что
$n_\Lambda(s)\ge(\rho\sigma+\delta)s^\rho$ при $s\ge T$. Тогда при больших $R$
\[
\int_0^R\frac{n_\Lambda(s)}{s}\,ds\ge(\rho\sigma+\delta)\int_T^R s^{\rho-1}ds=\Big(\sigma+\frac\delta\rho\Big)(R^\rho-T^\rho),
\]
а по \eqref{t38:J2} левая часть не больше $(\sigma+\eps)R^\rho-\ln|G(0)|$. Взяв $\eps<\delta/\rho$ и устремив $R\to\infty$
(разделив на $R^\rho$), получаем $\sigma+\delta/\rho\le\sigma+\eps$ --- противоречие. Значит, $\Delta^\flat_\rho(\Lambda)\le\rho\sigma$.
\end{proof}

\begin{remarkbox}{Замечание}
Первое неравенство --- это оценка (16.12) Леммы 16.4 (см. билет 35); на слайде 447 её вывод назван выходящим за рамки
курса, но, как видно, он занимает несколько строк. Точность неравенств (существование функций, на которых достигаются
равенства) на лекциях не обсуждается; её мы не доказываем. Для $\sin z$: $\rho=\sigma=1$, $\Delta_1(\Lambda)=2/\pi<1=\rho\sigma<e$.
\end{remarkbox}

\begin{thmbox}{Предложение 16.3}
Пусть $f$ --- ц.ф. конечного порядка $\rho\notin\ZZ$, $\Lambda$ --- последовательность её нулей, $\sigma$ --- тип $f$
относительно порядка $\rho$. Тогда величины $\Delta^\sharp_\rho(\Lambda)$ и $\sigma$ одновременно равны нулю, или
одновременно бесконечны, или одновременно конечны и положительны.
\end{thmbox}
\textit{Идея доказательства (без доказательства на лекциях).} В одну сторону это Предложение 16.2: $\Delta^\sharp_\rho\le e\rho\sigma$,
так что $\sigma=0\Rightarrow\Delta^\sharp_\rho=0$ и $\Delta^\sharp_\rho=\infty\Rightarrow\sigma=\infty$. В другую сторону используется
теорема Адамара: так как $\rho\notin\ZZ$, $f=z^me^{g}\Pi$, где $\deg g\le[\rho]<\rho$, а $\Pi$ --- каноническое произведение рода
$p=[\rho]$; множитель $e^{g}$ имеет нулевой тип относительно порядка $\rho$, а для канонического произведения справедлива
оценка вида $\ln M_\Pi(r)\le C\big(r^p\!\int_0^r\!\frac{n(t)}{t^{p+1}}dt+r^{p+1}\!\int_r^\infty\!\frac{n(t)}{t^{p+2}}dt\big)$;
при $n(t)\le At^\rho$ и $p<\rho<p+1$ оба интеграла сходятся и дают $\ln M_\Pi(r)\le C'Ar^\rho$, т.е. $\sigma\le C''\Delta^\sharp_\rho$.
При целом $\rho$ утверждение неверно: $e^z$ имеет тип $1$ и вообще не имеет нулей.

\subsection{Полнота системы экспонент}

Напомним два факта функционального анализа.

\begin{thmbox}{Следствие теоремы Хана--Банаха}
Пусть $E$ --- нормированное пространство, $E_0\subset E$ --- замкнутое подпространство, $x_0\in E\setminus E_0$. Тогда существует
$\varphi\in E^*$ такой, что $\varphi|_{E_0}=0$ и $\varphi(x_0)\ne0$. В частности, система $\{x_n\}$ полна в $E$ тогда и только
тогда, когда не существует ненулевого $\varphi\in E^*$, обращающегося в нуль на всех $x_n$.
\end{thmbox}
\begin{detail}{почему «в частности»}
Пусть $E_0$ --- замыкание линейной оболочки $\{x_n\}$. Если система не полна, то $E_0\ne E$; беря $x_0\notin E_0$, получаем
$\varphi\ne0$ с $\varphi(x_n)=0$. Обратно, если $\varphi\ne0$ и $\varphi(x_n)=0$ при всех $n$, то $\varphi$ равен нулю на
линейной оболочке, а по непрерывности --- и на $E_0$; так как $\varphi\ne0$, то $E_0\ne E$.
\end{detail}

\begin{thmbox}{Теорема Рисса (об общем виде функционала на $C[a,b]$)}
Для каждого $\varphi\in C[a,b]^*$ (комплексный случай) существует единственная комплексная регулярная борелевская мера
$\mu$ на $[a,b]$ такая, что $\varphi(h)=\int_{[a,b]}h\,d\mu$ при всех $h\in C[a,b]$; при этом $\|\varphi\|=\|\mu\|$.
\end{thmbox}

\begin{thmbox}{Предложение 16.4}
Пусть $\Lambda=\{\lambda_n\}_{n=1}^\infty$ --- последовательность \emph{попарно различных} вещественных чисел,
$n(t)=\#\{n:|\lambda_n|\le t\}$ --- её считающая функция. Если
\[
\liminf_{t\to\infty}\frac{n(t)}{t}>2,
\]
то система $\{e^{i\lambda_nt}\}_{n=1}^\infty$ полна в $C[-\pi,\pi]$.
\end{thmbox}

\begin{remarkbox}{Замечание (о попарной различности)}
На слайдах условие попарной различности $\lambda_n$ не выписано, но оно необходимо: повторение одного и того же $\lambda$
не добавляет к системе новых функций, но увеличивает $n(t)$. Например, если каждое целое число взять три раза, то
$n(t)\sim6t$, но система $\{e^{ikt}\}_{k\in\ZZ}$ \emph{не} полна в $C[-\pi,\pi]$ (см. ниже). В доказательстве это
используется так: различные $\lambda_n$ дают \emph{различные} нули функции $F$.
\end{remarkbox}

\begin{proof}
Предположим противное: система не полна. Проведём доказательство в несколько шагов.

\textbf{Шаг 1 (функционал и мера).} По следствию теоремы Хана--Банаха существует ненулевой $\varphi\in C[-\pi,\pi]^*$ с
$\varphi(e^{i\lambda_nt})=0$ для всех $n$. По теореме Рисса $\varphi(h)=\int_{[-\pi,\pi]}h(t)\,d\mu(t)$ для некоторой
комплексной борелевской меры $\mu$, $0<\|\mu\|<\infty$. Итак,
\[
\int_{[-\pi,\pi]}e^{i\lambda_nt}\,d\mu(t)=0,\qquad n=1,2,\dots
\]
(обратно, существование такой ненулевой меры, очевидно, означает неполноту).

\textbf{Шаг 2 (функция $F$ --- целая).} Положим
\[
F(z)=\int_{[-\pi,\pi]}e^{izt}\,d\mu(t),\qquad z\in\CC.
\]
При фиксированном $z$ ряд $e^{izt}=\sum_{k\ge0}\frac{(iz)^kt^k}{k!}$ сходится равномерно по $t\in[-\pi,\pi]$
(мажорируется рядом $\sum_k\frac{(|z|\pi)^k}{k!}$), поэтому его можно почленно интегрировать по конечной мере:
\[
F(z)=\sum_{k=0}^\infty c_kz^k,\qquad c_k=\frac{i^k}{k!}\int_{[-\pi,\pi]}t^k\,d\mu(t),\qquad |c_k|\le\frac{\|\mu\|\pi^k}{k!}.
\]
Из оценки коэффициентов следует, что степенной ряд сходится во всей плоскости; значит, $F$ --- целая функция.
(Попутно: $M_F(r)\le\|\mu\|e^{\pi r}$, т.е. $F$ --- функция экспоненциального типа $\le\pi$.)

\textbf{Шаг 3 (оценка $|F(z)|\le\|\mu\|e^{\pi|\Im z|}$).} Пусть $z=x+iy$. При $t\in[-\pi,\pi]$
\[
|e^{izt}|=|e^{ixt}e^{-yt}|=e^{-yt}\le e^{|y|\,|t|}\le e^{\pi|y|},
\]
следовательно, $|F(z)|\le\int|e^{izt}|\,d|\mu|(t)\le\|\mu\|e^{\pi|\Im z|}$. В частности, $F$ ограничена на $\RR$.

\textbf{Шаг 4 ($F\not\equiv0$).} Если $F\equiv0$, то все коэффициенты $c_k=0$, т.е. $\int t^k\,d\mu=0$ при всех $k\ge0$, откуда
$\int p\,d\mu=0$ для любого многочлена $p$. По теореме Вейерштрасса многочлены плотны в $C[-\pi,\pi]$, а функционал
$h\mapsto\int h\,d\mu$ непрерывен; значит, $\varphi(h)=\int h\,d\mu=0$ для всех $h$, т.е. $\varphi=0$ --- противоречие.

\textbf{Шаг 5 (нули $F$).} По Шагу 1 $F(\lambda_n)=0$ при всех $n$; так как $\lambda_n$ попарно различны, это различные нули $F$.
Пусть $F$ имеет в нуле нуль кратности $m\ge0$, и $G(z)=F(z)/z^m$ --- ц.ф., $G(0)\ne0$, нули $G$ --- ненулевые нули $F$.
Обозначим $n^*(t)=\#\{n:\ 0<|\lambda_n|\le t\}$; так как ноль может встретиться в $\Lambda$ не более одного раза,
$n^*(t)\ge n(t)-1$, и $n^*(t)\le n_G^{\le}(t)$, где $n_G^{\le}(t)$ --- число нулей $G$ в замкнутом круге $|z|\le t$.

\textbf{Шаг 6 (формула Йенсена).} Пусть $r\ge1$ и на окружности $|z|=r$ нет нулей $G$ (это верно для всех $r$, кроме дискретного
множества). По \eqref{t38:jensen} (с $n^{\le}_G$ вместо $n_G$ --- интеграл тот же)
\[
\int_0^r\frac{n_G^{\le}(t)}{t}\,dt=\frac1{2\pi}\int_0^{2\pi}\ln|G(re^{i\theta})|\,d\theta-\ln|G(0)|.
\]
При $|z|=r\ge1$: $\ln|G(z)|=\ln|F(z)|-m\ln r\le\ln|F(z)|\le\ln\|\mu\|+\pi r|\sin\theta|$ (Шаг 3, $\Im(re^{i\theta})=r\sin\theta$). Так как
$\int_0^{2\pi}|\sin\theta|\,d\theta=4$, получаем
\[
\int_0^r\frac{n_G^{\le}(t)}{t}\,dt\le\frac{\pi r}{2\pi}\int_0^{2\pi}|\sin\theta|\,d\theta+\ln\|\mu\|-\ln|G(0)|=2r+C_1 .
\]
Поэтому для фиксированного $\delta\in(0,1)$
\begin{equation}\label{t38:upper}
\int_\delta^r\frac{n^*(t)}{t}\,dt\le\int_0^r\frac{n_G^{\le}(t)}{t}\,dt\le2r+C_1 \qquad(=2r+O(1),\ r\to\infty).
\end{equation}
Это и есть оценка (16.19) слайдов.

\textbf{Шаг 7 (противоречие).} Пусть $L=\liminf_{t\to\infty}n(t)/t>2$. Положим $\eps=\min\{(L-2)/2,\,1\}>0$; найдётся
$T_0\ge1$ такое, что $n(t)\ge(2+\eps)t$ при $t\ge T_0$, а увеличив $T_0$, получим и
$n^*(t)\ge(2+\eps)t-1\ge(2+\tfrac\eps2)t$ при $t\ge T_0$. Тогда при $r>T_0$
\[
\int_\delta^r\frac{n^*(t)}{t}\,dt\ge\int_{T_0}^r\Big(2+\frac\eps2\Big)dt=\Big(2+\frac\eps2\Big)(r-T_0).
\]
Сравнивая с \eqref{t38:upper}: $\frac\eps2 r\le C_1+(2+\frac\eps2)T_0$ для всех достаточно больших $r$, не совпадающих с
модулями нулей $G$; таких $r$ сколь угодно много, и при $r\to\infty$ получаем противоречие. Следовательно, система полна.
\end{proof}

\begin{remarkbox}{Замечания и примеры}
\begin{enumerate}
\item \emph{Откуда константа $2$.} Если пользоваться лишь грубой оценкой $|F(z)|\le\|\mu\|e^{\pi|z|}$, то Йенсен дал бы
$\int^r n/t\le\pi r+O(1)$ и порог $\pi$. Выигрыш даёт оценка через $|\Im z|$: среднее значение $\pi|\sin\theta|$ по окружности равно $2$.
\item \emph{Условие нельзя ослабить до «$\ge2$».} Для $\Lambda=\ZZ$ (занумерованного в последовательность) $n(t)=2[t]+1$,
$n(t)/t\to2$, но система $\{e^{ikt}\}_{k\in\ZZ}$ не полна в $C[-\pi,\pi]$: все её линейные комбинации $h$ удовлетворяют
$h(-\pi)=h(\pi)$, и это равенство сохраняется при равномерном пределе; поэтому, например, $h(t)=t$ не приближается.
Соответствующий функционал $\varphi(h)=h(\pi)-h(-\pi)$ (мера $\mu=\delta_\pi-\delta_{-\pi}$) даёт $F(z)=2i\sin\pi z$ ---
функцию, нули которой суть ровно $\ZZ$.
\item \emph{Пример полной системы.} При любом $\delta>0$ система $\{e^{ikt/(1+\delta)}\}_{k\in\ZZ}$ полна в $C[-\pi,\pi]$:
здесь $n(t)=2[(1+\delta)t]+1$, $n(t)/t\to2(1+\delta)>2$.
\item Вещественность $\lambda_n$ в доказательстве фактически не используется: при $\lambda_n\in\CC$ рассуждение дословно то же
(если $n(t)$ --- число $\lambda_n$ в круге $|\lambda|\le t$).
\end{enumerate}
\end{remarkbox}

\begin{remarkbox}{Коротко: план ответа на экзамене}
\begin{enumerate}
\item Определения: считающая функция $n_\Lambda(t)$, верхняя/нижняя $\rho$-плотности; Лемма 16.7 (плотности через $n_\Lambda(t)/t^\rho$).
\item Предложения 16.2, 16.3 (формулировки; идея $\Delta^\sharp\le e\rho\sigma$ --- Йенсен с $k=e^{1/\rho}$).
\item Полнота системы; критерий через Хана--Банаха: неполна $\iff$ есть $\varphi\ne0$, $\varphi(x_n)=0$.
\item Теорема Рисса: $\varphi(h)=\int h\,d\mu$; неполнота $\iff$ $\int e^{i\lambda_nt}d\mu=0$ для ненулевой меры $\mu$.
\item $F(z)=\int e^{izt}d\mu$ --- целая (степенной ряд), $|F(z)|\le\|\mu\|e^{\pi|\Im z|}$; $F\not\equiv0$ по Вейерштрассу.
\item $\lambda_n$ --- различные нули $F$; формула Йенсена: $\int_\delta^r n(t)/t\,dt\le\frac r2\int_0^{2\pi}|\sin\theta|d\theta+O(1)=2r+O(1)$.
\item Если $n(t)\ge(2+\eps)t$ при $t\ge T_0$, то левая часть $\ge(2+\eps)(r-T_0)$ --- противоречие.
\item Пример $\Lambda=\ZZ$: $n(t)/t\to2$, система не полна ($F=2i\sin\pi z$); константа $2$ точная.
\end{enumerate}
\end{remarkbox}

% =====================================================================
\ticket{39}{$A$-точки целых функций. Исключительные значения в смысле Бореля целых функций}

\begin{defs}
\textbf{Порядок} $\rho$ ц.ф. $f$: $\rho=\inf\{\alpha>0:\ M_f(r)<e^{r^\alpha}\text{ при всех больших }r\}=\limsup_{r\to\infty}\frac{\ln\ln M_f(r)}{\ln r}$
(см. билет 33). Порядок многочлена (и константы) равен $0$.

\textbf{$A$-точка.} Пусть $A\in\CC$. Точка $z\in\CC$ называется \emph{$A$-точкой} функции $f$, если $f(z)=A$; кратность
$A$-точки --- кратность нуля функции $f-A$ в ней. $A$-точки --- это нули функции $f-A$, и \emph{последовательность $A$-точек}
--- ненулевые нули $f-A$, выписанные с учётом кратности в стандартном порядке $0<|z_1|\le|z_2|\le\dots$.

\textbf{Показатель сходимости} последовательности $\{a_n\}$, $0<|a_n|\le|a_{n+1}|\to\infty$:
$\tau=\inf\{\lambda>0:\sum|a_n|^{-\lambda}<\infty\}=\limsup\frac{\ln n}{\ln|a_n|}$ (см. билет 35); $\tau=0$ для конечной
последовательности. Через $\tau_A=\tau_A(f)$ обозначаем показатель сходимости последовательности $A$-точек $f$.

\textbf{Каноническое произведение} рода $p$ (см. билет 36): $\Pi(z)=\prod_n\big(1-\frac z{a_n}\big)\exp\big(\sum_{k=1}^p\frac{z^k}{k a_n^k}\big)$,
где $p\in\ZZ_+$ --- наименьшее число с $\sum|a_n|^{-(p+1)}<\infty$ (род однозначно определяется последовательностью).

\textbf{Исключительное значение по Борелю} (Определение 16.6). Пусть $f$ --- ц.ф. конечного положительного порядка $\rho$.
Значение $A\in\CC$ называется \emph{исключительным борелевским значением} $f$, если $\tau_A<\rho$ (т.е. $A$-точек «существенно
меньше», чем у типичного значения).

\textbf{Исключительное значение по Пикару}: $A$ называется пикаровским исключительным значением $f$, если $f$ принимает его
лишь конечное число раз (или ни разу). Пример: $0$ для $e^z$.
\end{defs}

\subsection{Нужные факты из билетов 33--37}

\begin{thmbox}{Теорема 16.3 (Адамар; см. билет 36)}
Пусть $f\not\equiv0$ --- ц.ф. конечного порядка $\rho$, имеющая в нуле нуль кратности $m\ge0$, $\{a_n\}$ --- её ненулевые нули,
$p$ --- род последовательности $\{a_n\}$. Тогда $f(z)=z^me^{g(z)}\Pi(z)$, где $\Pi$ --- каноническое произведение рода $p$,
а $g$ --- многочлен степени $\le\rho$.
\end{thmbox}
\begin{thmbox}{Теорема 16.4 (Борель; см. билет 37) и Предложение 16.1}
В обозначениях теоремы Адамара, если $\tau$ --- показатель сходимости $\{a_n\}$ и $\deg g=k$, то порядок $f$ равен $\max\{\tau,k\}$.
В частности: (а) каноническое произведение (умноженное на $z^m$) имеет порядок $\tau$; (б) если $\rho\notin\ZZ$, то $\tau=\rho$.
Всегда $\tau\le\rho$ (лемма Адамара о нулях, Предложение 16.1).
\end{thmbox}
\begin{detail}{(а) и (б) из теоремы Бореля}
(а) Для $z^m\Pi$ «$g\equiv0$», $k=0$, порядок $\max\{\tau,0\}=\tau$. (б) $\rho=\max\{\tau,k\}$, причём $k\in\ZZ$; если $\rho\notin\ZZ$,
то $\rho\ne k$, значит $\rho=\tau$.
\end{detail}

\begin{thmbox}{Лемма А (элементарные свойства порядка)}
Обозначим через $\ord f$ порядок ц.ф. $f$. Тогда:
\begin{enumerate}
\item $\ord(f_1f_2)\le\max\{\ord f_1,\ord f_2\}$, $\ord(f_1+f_2)\le\max\{\ord f_1,\ord f_2\}$;
\item $\ord f'\le\ord f$; $\ord(Pf)\le\ord f$ для многочлена $P$;
\item $\ord(f-A)=\ord f$ для любой константы $A$;
\item если $R$ --- многочлен степени $d\ge1$, то $\ord e^{R}=d$.
\end{enumerate}
\end{thmbox}
\begin{proof}
Пусть $\alpha>\max\{\ord f_1,\ord f_2\}$; при больших $r$ $M_{f_j}(r)<\exp(r^{\alpha'})$, где $\max\ord f_j<\alpha'<\alpha$.
Тогда $M_{f_1f_2}(r)\le M_{f_1}M_{f_2}<\exp(2r^{\alpha'})<\exp(r^\alpha)$ и $M_{f_1+f_2}(r)<2\exp(r^{\alpha'})<\exp(r^\alpha)$ при больших $r$.
Это даёт 1. Для 2: по неравенству Коши на окружности $|\zeta-z|=1$: $|f'(z)|\le M_f(|z|+1)$, т.е. $M_{f'}(r)\le M_f(r+1)<\exp((r+1)^{\alpha'})<\exp(r^\alpha)$
при больших $r$; второе утверждение следует из 1, так как $\ord P=0$. Для 3: если $f$ --- константа, обе части равны $0$; иначе
$M_f(r)\to\infty$ (теорема Лиувилля) и $\frac12M_f(r)\le M_f(r)-|A|\le M_{f-A}(r)\le M_f(r)+|A|\le2M_f(r)$ при больших $r$, а
$\ln\ln(2M)/\ln r$ и $\ln\ln M/\ln r$ имеют одинаковый верхний предел. Для 4: $R(z)=cz^d+\dots$, $c\ne0$; при $|z|=r$
$|e^{R}|\le e^{|R(z)|}\le e^{C r^d}$, $r\ge1$, так что $\ord e^R\le d$. На луче, где $cz^d=|c|r^d>0$ (т.е. $\arg z=-\arg c/d$),
$\Re R(z)\ge|c|r^d-C'r^{d-1}$, откуда $M_{e^R}(r)\ge\exp(\frac{|c|}2r^d)$ при больших $r$ и $\ord e^R\ge d$.
\end{proof}

\begin{thmbox}{Лемма Б (точный порядок произведения $e^Ru$)}
Пусть $R$ --- многочлен степени $\rho\ge1$, $u\not\equiv0$ --- ц.ф. порядка $\ord u<\rho$. Тогда $\ord(e^Ru)=\rho$.
\end{thmbox}
\begin{proof}
По Лемме А $h=e^Ru$ имеет порядок $\le\rho<\infty$. Предположим, что $\ord h<\rho$. Функции $h$ и $u$ имеют одни и те же нули
с одинаковыми кратностями (так как $e^R\ne0$): кратность $k$ нуля в начале координат, одну и ту же последовательность $\{a_n\}$
ненулевых нулей, а значит, и один и тот же род $p$ и одно и то же каноническое произведение $\Pi$. По теореме Адамара
$u=z^ke^S\Pi$, $h=z^ke^T\Pi$, где $\deg S\le\ord u<\rho$, $\deg T\le\ord h<\rho$. Тогда
$z^ke^{R+S}\Pi=e^Ru=h=z^ke^T\Pi$, откуда $e^{R+S-T}=1$ вне дискретного множества нулей, а по непрерывности --- везде.
Целая функция $R+S-T$ принимает значения в дискретном множестве $2\pi i\ZZ$, поэтому постоянна (связность $\CC$). Но
$\deg(R+S-T)=\rho\ge1$, так как $\deg S,\deg T<\rho=\deg R$. Противоречие.
\end{proof}

\subsection{$A$-точки ц.ф. конечного порядка}

Во многих ситуациях $A$-точки голоморфных функций ведут себя как нули. Но для ц.ф. есть интересная особенность: величина
$\tau_A$ не зависит от $A$, за исключением, быть может, одного значения.

\begin{thmbox}{Предложение 16.5}
Пусть $f$ --- ц.ф. конечного ненулевого порядка $\rho$. Если $\rho\notin\ZZ$, то для любого $A\in\CC$ последовательность
$A$-точек $f$ имеет показатель сходимости $\tau_A=\rho$. Если же $\rho\in\ZZ$ (т.е. $\rho\in\NN$), то $\tau_A=\rho$ для всех
$A\in\CC$, кроме, быть может, одного.
\end{thmbox}

\begin{proof}
Так как $\rho>0$, $f$ не константа, и при любом $A$ функция $f-A\not\equiv0$ имеет порядок $\rho$ (Лемма А, п.~3).

\textbf{Случай $\rho\notin\ZZ$.} $A$-точки --- нули ц.ф. $f-A$ нецелого порядка $\rho$; по следствию (б) теоремы Бореля $\tau_A=\rho$.

\textbf{Случай $\rho\in\NN$.} По Предложению 16.1 всегда $\tau_A\le\rho$. Значит, достаточно доказать, что неравенство $\tau_A<\rho$
не может выполняться для двух различных значений. Пусть $\tau_A<\rho$ и $\tau_B<\rho$, $A\ne B$; придём к противоречию.

\emph{Шаг 1 (представление; «простое упражнение» слайдов).} Покажем, что
\begin{equation}\label{t39:rep}
f(z)-A=e^{P(z)}\phi(z),\qquad f(z)-B=e^{Q(z)}\psi(z),
\end{equation}
где $P,Q$ --- многочлены степени ровно $\rho$, а $\phi,\psi\not\equiv0$ --- ц.ф. порядка $<\rho$.
\begin{detail}{доказательство представления \eqref{t39:rep}}
По теореме Адамара $f-A=z^me^{P}\Pi_A$, где $\Pi_A$ --- каноническое произведение по $A$-точкам, $\deg P\le\rho$. Положим
$\phi=z^m\Pi_A$; по следствию (а) теоремы Бореля $\ord\phi=\tau_A<\rho$. По теореме Бореля $\rho=\ord(f-A)=\max\{\tau_A,\deg P\}$,
а так как $\tau_A<\rho$, то $\deg P=\rho$. Для $B$ --- аналогично. (Если $A$-точек конечное число, $\Pi_A$ --- многочлен,
$\ord\phi=0<\rho$, и рассуждение то же.)
\end{detail}

\emph{Шаг 2.} Вычитая равенства \eqref{t39:rep}, получаем $A-B=e^{Q}\psi-e^{P}\phi$. Дифференцируем:
\[
0=e^{Q}\big(Q'\psi+\psi'\big)-e^{P}\big(P'\phi+\phi'\big).
\]
\emph{Выражения в скобках не равны тождественно нулю.} Пусть, например, $P'\phi+\phi'\equiv0$. Тогда
$(\phi e^{P})'=e^P(\phi'+P'\phi)\equiv0$, т.е. $\phi e^P\equiv c$ и $\phi=c\,e^{-P}$, причём $c\ne0$, так как $\phi\not\equiv0$.
Но по Лемме А, п.~4, $\ord(ce^{-P})=\deg P=\rho$, а $\ord\phi<\rho$ --- противоречие. Аналогично $Q'\psi+\psi'\not\equiv0$.
Значит,
\begin{equation}\label{t39:eq}
e^{Q-P}\big(Q'\psi+\psi'\big)=P'\phi+\phi'.
\end{equation}
По Лемме А (пп.~1, 2) $\ord(Q'\psi+\psi')\le\ord\psi<\rho$ и $\ord(P'\phi+\phi')\le\ord\phi<\rho$.

\emph{Шаг 3 (старшие коэффициенты совпадают).} Пусть $P(z)=az^\rho+\dots$, $Q(z)=bz^\rho+\dots$ ($a,b\ne0$). Если $a\ne b$, то
$\deg(Q-P)=\rho$, и по Лемме Б левая часть \eqref{t39:eq} имеет порядок ровно $\rho$, а правая --- порядок $<\rho$. Противоречие;
следовательно, $a=b$.

\emph{Шаг 4.} При $a=b$ положим $P_1=P-az^\rho$, $Q_1=Q-az^\rho$ --- многочлены степени $<\rho$. Умножая $A-B=e^Q\psi-e^P\phi$ на
$e^{-az^\rho}$, получаем
\[
(A-B)\,e^{-az^\rho}=e^{Q_1}\psi-e^{P_1}\phi .
\]
Правая часть имеет порядок $\le\max\{\deg Q_1,\deg P_1,\ord\psi,\ord\phi\}<\rho$ (Лемма А; $\ord e^{Q_1}=\deg Q_1$ или $0$).
Если $A\ne B$, левая часть имеет порядок $\rho$ (Лемма А, п.~4). Противоречие. Следовательно, $A=B$, что противоречит выбору.
\end{proof}

\begin{remarkbox}{Замечание (опечатка на слайде 467)}
На слайде сказано: из $P'\phi+\phi'\equiv0$ следует $\phi(z)=\exp(P(z)+C)$. Правильно: $\phi(z)=\exp(-P(z)+C)$ (решение уравнения
$\phi'=-P'\phi$). На вывод это не влияет: порядок $e^{\pm P+C}$ равен $\deg P=\rho$. Кроме того, слова «в левой части стоит ц.ф. порядка
$\rho$» требуют обоснования --- это Лемма Б выше (произведение $e^{R}$ на ненулевую функцию меньшего порядка не может понизить порядок).
\end{remarkbox}

\subsection{Исключительные значения по Борелю}

Из Предложения 16.5 непосредственно получаем:

\begin{thmbox}{Следствие (о борелевских исключительных значениях)}
Ц.ф. конечного нецелого порядка не имеет исключительных борелевских значений. Ц.ф. конечного ненулевого целого порядка может
их иметь, но не более одного.
\end{thmbox}

\begin{thmbox}{Следствие (пикаровские значения для функций конечного порядка)}
Ц.ф. конечного положительного порядка имеет не более одного пикаровского исключительного значения, а при нецелом порядке
--- ни одного: каждое значение принимается бесконечно много раз.
\end{thmbox}
\begin{proof}
Если $A$ принимается конечное число раз, то $\tau_A=0<\rho$, т.е. $A$ --- борелевское исключительное значение. Остаётся
применить предыдущее следствие. (Это частный случай малой теоремы Пикара, полученный элементарными средствами; см. вопрос 42*.)
\end{proof}

\begin{remarkbox}{Примеры}
\begin{enumerate}
\item $f(z)=e^z$, $\rho=1$. Значение $0$ не принимается: $\tau_0=0<1$ --- исключительное (и по Пикару, и по Борелю).
При $A\ne0$ $A$-точки --- $z_k=\ln|A|+i\arg A+2\pi ik$, $k\in\ZZ$; $\sum_k|z_k|^{-\lambda}<\infty\iff\lambda>1$, т.е. $\tau_A=1=\rho$.
\item $f(z)=e^{z^2}\sin z$. По Лемме Б ($R=z^2$, $u=\sin z$ порядка $1<2$) порядок $\rho=2$. Нули --- $\pi k$, $k\in\ZZ$; ненулевые
$\pi k$ имеют $\tau_0=1$ (ряд $\sum_{k\ne0}|\pi k|^{-\lambda}$ сходится $\iff\lambda>1$). Значит, $0$ --- борелевское исключительное
значение, хотя принимается бесконечно много раз (не пикаровское). По Предложению 16.5 $\tau_A=2$ для всех $A\ne0$.
\item $\cos z$, $\rho=1$: уравнение $\cos z=A$ имеет решения $\pm w_A+2\pi k$ ($w_A$ --- одно из решений), так что
$\tau_A=1$ при всех $A$: исключительных значений нет, хотя порядок целый. Так что при целом $\rho$ исключительное значение
может быть, а может и не быть.
\item $e^{z^\rho}$ при $\rho\in\NN$: $0$ --- исключительное, $\tau_A=\rho$ при $A\ne0$ (решения $z^\rho=\ln|A|+i\arg A+2\pi ik$ --- это
$\rho$ корней из чисел $w_k$ с $|w_k|\asymp|k|$, так что $|z|\asymp|k|^{1/\rho}$, $\#\{|z|<t\}\asymp t^\rho$).
\item Порядок $\rho=0$ исключён: определение требует $\tau_A<\rho$, что невозможно при $\rho=0$ (например, у многочлена каждое значение
принимается конечное число раз).
\end{enumerate}
\end{remarkbox}

\begin{remarkbox}{Коротко: план ответа на экзамене}
\begin{enumerate}
\item $A$-точки $=$ нули $f-A$; $\tau_A$ --- показатель сходимости; $\ord(f-A)=\ord f$, поэтому $\tau_A\le\rho$.
\item Нецелый $\rho$: по теореме Бореля $\rho=\max\{\tau_A,\deg g\}$, $\deg g\in\ZZ$ $\Rightarrow$ $\tau_A=\rho$ для всех $A$.
\item Целый $\rho$: пусть $\tau_A,\tau_B<\rho$, $A\ne B$. По Адамару--Борелю $f-A=e^P\phi$, $f-B=e^Q\psi$, $\deg P=\deg Q=\rho$, $\ord\phi,\ord\psi<\rho$.
\item Дифференцируем $A-B=e^Q\psi-e^P\phi$: $e^{Q-P}(Q'\psi+\psi')=P'\phi+\phi'$, скобки $\not\equiv0$.
\item Старшие коэффициенты $P,Q$ равны (иначе слева порядок $\rho$, справа $<\rho$).
\item Тогда $(A-B)e^{-az^\rho}=e^{Q_1}\psi-e^{P_1}\phi$: слева порядок $\rho$, справа $<\rho$ $\Rightarrow$ $A=B$.
\item Определение 16.6, следствия; примеры $e^z$, $e^{z^2}\sin z$ ($\tau_0=1<2$), $\cos z$.
\end{enumerate}
\end{remarkbox}

% =====================================================================
\ticket{40}{Принцип Фрагмена--Линделёфа для целых функций конечного порядка}

\begin{defs}
\textbf{Порядок} ц.ф. $f$ (см. билет 33): $\rho=\inf\{\alpha>0:\ M_f(r)<e^{r^\alpha}$ при всех достаточно больших $r\}$, где
$M_f(r)=\max_{|z|=r}|f(z)|$. Если порядок $f$ \emph{не выше} $\rho$ и $\beta>\rho$, то $M_f(r)<\exp(r^\beta)$ при всех $r>R(\beta)$.

\textbf{Сектор (угол)} раствора $\theta_0\in(0,2\pi)$ с вершиной $0$ --- множество $\{z\ne0:\ \vartheta_1<\arg z<\vartheta_1+\theta_0\}$;
его граница --- два луча (стороны угла) и вершина $0$. Поворотом $z\mapsto e^{i\vartheta}z$ любой такой сектор переводится в
симметричный $S=\{z\ne0:|\arg z|<\theta_0/2\}$, причём поворот не меняет ни порядка ц.ф., ни величин $|f|$ на соответствующих множествах.

\textbf{Главная ветвь степени.} Для $\gamma>0$ и $z=re^{i\vartheta}$, $r>0$, $|\vartheta|<\pi$, положим $z^\gamma=r^\gamma e^{i\gamma\vartheta}$
($=e^{\gamma\ln z}$ с главной ветвью логарифма); это голоморфная функция в $\CC\setminus(-\infty,0]$, и
$\Re z^\gamma=r^\gamma\cos(\gamma\vartheta)$, $|e^{-\eps z^\gamma}|=\exp(-\eps r^\gamma\cos\gamma\vartheta)$. Положив $0^\gamma=0$, получаем
функцию, непрерывную на $\{z:|\arg z|\le\vartheta_0\}\cup\{0\}$ при любом $\vartheta_0<\pi$.

\textbf{Принцип максимума модуля} (для ограниченной области). Если $D$ --- ограниченная область, $g\in H(D)\cap C(\clos D)$, то
$\max_{\clos D}|g|=\max_{\partial D}|g|$.

\textbf{Принцип Фрагмена--Линделёфа} --- общее название утверждений вида «если голоморфная функция в неограниченной области
ограничена на границе и растёт не слишком быстро внутри, то она ограничена и внутри» (для ограниченных областей это
просто принцип максимума; для неограниченных без ограничения роста --- неверно, см. пример $e^z$ ниже).
\end{defs}

\subsection{Мотивировка: рост по направлениям}

Изучая рост ц.ф., мы до сих пор имели дело только с $M_f(r)$. Эта функция не содержит всей информации о росте $f$.
Например, для $f(z)=e^z$ имеем $M_f(r)=e^r$, а при $z=re^{i\vartheta}$ $|e^z|=e^{r\cos\vartheta}$. Поэтому:
\begin{itemize}
\item если $|\arg z|<\pi/2-\eps$ ($\eps>0$), то $\cos\vartheta>\cos(\pi/2-\eps)=\sin\eps$ и $|e^z|>e^{|z|\sin\eps}\to\infty$ при $|z|\to\infty$;
\item если $\arg z\in(\pi/2+\eps,3\pi/2-\eps)$, то $\cos\vartheta<-\sin\eps$ и $|e^z|<e^{-|z|\sin\eps}\to0$;
\item на мнимой оси $|e^{iy}|=1$.
\end{itemize}
В частности, $e^z$ ограничена на границе правой полуплоскости (угла раствора $\pi$), но не ограничена внутри неё.
Принцип Фрагмена--Линделёфа показывает, что это возможно лишь при достаточно большом порядке функции.

\subsection{Теорема Фрагмена--Линделёфа}

\begin{thmbox}{Теорема 16.5 (принцип Фрагмена--Линделёфа)}
Пусть $f$ --- ц.ф. порядка не выше $\rho$, и пусть $0<\alpha<1/\rho$ (при $\rho=0$ --- любое $\alpha$), причём $\alpha<2$.
Пусть $S$ --- сектор раствора $\pi\alpha$ с вершиной в $0$. Если $|f(z)|\le C$ при $z\in\partial S$, где $C$ --- константа, то
$|f(z)|\le C$ и при $z\in S$.
\end{thmbox}
\begin{remarkbox}{Замечание к формулировке}
Условие $\alpha<2$ на слайдах не выписано, но подразумевается: угол раствора $\pi\alpha$ должен быть меньше $2\pi$
(иначе «сектор» не является углом в плоскости). Случай «угла раствора $2\pi$» (плоскость с разрезом по лучу) разобран
отдельно в Следствии 2 ниже.
\end{remarkbox}

\begin{proof}
\textbf{Сведение.} Поворотом сведём задачу к сектору $S=\{z\ne0:\ |\arg z|<\pi\alpha/2\}$: если $S=e^{i\vartheta}S_0$, то функция
$f_1(z)=f(e^{i\vartheta}z)$ имеет тот же порядок, $M_{f_1}=M_f$, и $|f_1|\le C$ на $\partial S_0$. Если $C=0$, то $f=0$ на луче,
и по теореме единственности $f\equiv0$; поэтому считаем $C>0$.

\textbf{Выбор параметров.} Так как $\alpha<1/\rho$, можно выбрать числа $\beta,\gamma$ с
\[
\rho<\beta<\gamma<\frac1\alpha .
\]
Поскольку порядок $f$ не выше $\rho<\beta$, найдётся $R_0$ такое, что $M_f(r)<\exp(r^\beta)$ при $r>R_0$. Положим
$R_n=R_0+n$; тогда $R_n\uparrow\infty$ и
\begin{equation}\label{t40:arc}
|f(z)|<\exp(|z|^\beta)\qquad\text{при }|z|=R_n,\ n=1,2,\dots
\end{equation}
(для доказательства достаточно, чтобы \eqref{t40:arc} выполнялось на какой-нибудь последовательности $R_n\to\infty$).

\textbf{Вспомогательная функция.} Фиксируем $\eps>0$ и положим
\[
g(z)=g_\eps(z)=f(z)\exp(-\eps z^\gamma),
\]
где $z^\gamma$ --- главная ветвь. Так как $\pi\alpha/2<\pi$, замкнутый сектор $\clos S=\{|\arg z|\le\pi\alpha/2\}\cup\{0\}$
лежит в $\big(\CC\setminus(-\infty,0]\big)\cup\{0\}$, так что $g$ голоморфна в $S$ и непрерывна на $\clos S$, $g(0)=f(0)$.

\textbf{Ключевая оценка.} При $z=re^{i\vartheta}\in\clos S$, $|\vartheta|\le\pi\alpha/2$, имеем $\gamma|\vartheta|\le\gamma\pi\alpha/2<\pi/2$,
так как $\gamma<1/\alpha$. Косинус убывает на $[0,\pi/2]$, поэтому
\[
\cos(\gamma\vartheta)\ge c_0:=\cos\Big(\frac{\gamma\pi\alpha}2\Big)>0,\qquad
|\exp(-\eps z^\gamma)|=\exp\big(-\eps r^\gamma\cos\gamma\vartheta\big)\le\exp(-\eps c_0r^\gamma)\le1.
\]

\textbf{Усечённые секторы.} Пусть $S_n=S\cap\{|z|<R_n\}$ --- ограниченная область; $\partial S_n=\Gamma_n\cup\Sigma_n$, где
$\Gamma_n=\{R_ne^{i\vartheta}:|\vartheta|\le\pi\alpha/2\}$ --- дуга, а $\Sigma_n$ --- два отрезка сторон угла (вместе с вершиной $0$).
\begin{itemize}
\item На $\Sigma_n\subset\partial S$: $|g(z)|\le|f(z)|\cdot1\le C$.
\item На $\Gamma_n$: по \eqref{t40:arc} и ключевой оценке $|g(z)|<\exp\big(R_n^\beta-\eps c_0R_n^\gamma\big)$. Так как $\gamma>\beta$,
показатель $R_n^\beta-\eps c_0R_n^\gamma=R_n^\gamma(R_n^{\beta-\gamma}-\eps c_0)\to-\infty$; значит, найдётся $N=N(\eps)$ с
$|g|\le C$ на $\Gamma_n$ при $n\ge N$.
\end{itemize}
Итак, при $n\ge N$ имеем $|g|\le C$ на $\partial S_n$, и по принципу максимума модуля $|g|\le C$ на $S_n$.

\textbf{Предельные переходы.} Каждая точка $z\in S$ лежит в $S_n$ при достаточно большом $n$ (так как $R_n\to\infty$); значит,
$|g_\eps(z)|\le C$ для всех $z\in S$, т.е.
\[
|f(z)|\le C\,|\exp(\eps z^\gamma)|=C\exp(\eps r^\gamma\cos\gamma\vartheta)\le Ce^{\eps|z|^\gamma},\qquad z\in S.
\]
Фиксируя $z$ и устремляя $\eps\to0$, получаем $|f(z)|\le C$.
\end{proof}

\begin{remarkbox}{Смысл вспомогательной функции}
Множитель $e^{-\eps z^\gamma}$ по модулю не превосходит $1$ в замкнутом секторе (на сторонах ограничение $|f|\le C$ не портится),
но на больших дугах убывает как $e^{-\eps c_0r^\gamma}$, т.е. быстрее, чем может расти $f$ (не быстрее $e^{r^\beta}$, $\beta<\gamma$).
Положительность $\cos(\gamma\vartheta)$ во всём секторе возможна только при $\gamma\cdot\pi\alpha/2<\pi/2$, т.е. $\gamma<1/\alpha$;
а «обогнать» $f$ можно лишь при $\gamma>\rho$. Отсюда и условие $\rho<1/\alpha$.
\end{remarkbox}

\begin{remarkbox}{Обобщение (тем же доказательством)}
В доказательстве использовались только: голоморфность $f$ в $S$, непрерывность на $\clos S$, оценка $|f|\le C$ на $\partial S$ и
оценка $|f(z)|<\exp(|z|^\beta)$ на дугах $|z|=R_n\to\infty$ внутри $S$ с некоторым $\beta<1/\alpha$. Поэтому теорема верна для любой
такой функции (не обязательно целой), в том числе для ц.ф., \emph{нижний} порядок которых меньше $1/\alpha$.
\end{remarkbox}

\subsection{Следствия}

\begin{thmbox}{Следствие 1 (переформулировка)}
Пусть ц.ф. $f$ имеет порядок $\rho<\alpha$, $\alpha>1/2$, и ограничена на сторонах угла раствора $\pi/\alpha$ с вершиной в $0$:
$|f|\le C$. Тогда $|f|\le C$ во всём угле. В частности: ц.ф. порядка $<1$, ограниченная на сторонах полуплоскости
(т.е. на прямой, проходящей через $0$), ограничена в этой полуплоскости; ц.ф. порядка $<2$, ограниченная на сторонах прямого угла,
ограничена в нём, и т.д.
\end{thmbox}
\begin{proof}
Это Теорема 16.5 с $\alpha'=1/\alpha$ вместо $\alpha$: раствор $\pi\alpha'=\pi/\alpha<2\pi$, условие $\alpha'<1/\rho$ равносильно $\rho<\alpha$.
\end{proof}

\begin{thmbox}{Следствие 2 (простое следствие из лекций)}
Пусть $f$ --- ц.ф. порядка $\rho<1/2$, ограниченная на некотором луче $\{z:\arg z=\vartheta_0\}$. Тогда $f\equiv\mathrm{const}$.
\end{thmbox}
\begin{proof}
Поворотом можно считать, что луч --- отрицательная вещественная полуось: $|f(-s)|\le C$ при $s\ge0$. Рассмотрим ц.ф.
$h(w)=f(w^2)$. Её порядок не выше $2\rho<1$: $M_h(r)=\max_{|w|=r}|f(w^2)|=M_f(r^2)<\exp(r^{2\beta})$ при любом $\beta>\rho$ и
больших $r$. На мнимой оси $w=iy$ имеем $w^2=-y^2\le0$, поэтому $|h(iy)|\le C$.
Мнимая ось --- граница правой и левой полуплоскостей, каждая из которых --- угол раствора $\pi$ (т.е. $\alpha=1<1/(2\rho)$).
По Теореме 16.5 $|h|\le C$ в обеих полуплоскостях, т.е. во всей плоскости. По теореме Лиувилля $h\equiv\mathrm{const}$,
значит, $f(\zeta)=h(\sqrt\zeta)$ постоянна (каждое $\zeta\in\CC$ есть $w^2$ для некоторого $w$).
\end{proof}

\begin{remarkbox}{Замечание (к Следствию 2)}
Следствие 2 можно получить и прямо из доказательства Теоремы 16.5, применённого к «углу раствора $2\pi$» $\CC\setminus\{\arg z=\vartheta_0\}$
($\alpha=2<1/\rho$), но тогда $g$ не непрерывна на разрезе (значения $z^\gamma$ на двух берегах различны) и принцип максимума
приходится применять аккуратно. Приём $w\mapsto w^2$ («раскрытие» разреза) позволяет этого избежать.
\end{remarkbox}

\begin{remarkbox}{Точность условий: примеры}
\begin{enumerate}
\item \emph{Следствие 2 неверно при $\rho=1/2$.} Функция
\[
f(z)=\frac{\sin\sqrt z}{\sqrt z}=\sum_{k=0}^\infty\frac{(-1)^kz^k}{(2k+1)!}
\]
--- целая (ряд не зависит от выбора ветви $\sqrt z$). При $z=x>0$: $|f(x)|=|\sin\sqrt x|/\sqrt x\le1$ (так как $|\sin u|\le|u|$), т.е.
$f$ ограничена на положительной полуоси. Её порядок равен $1/2$: $M_f(r)\le\sum_k\frac{r^k}{(2k+1)!}=\frac{\operatorname{sh}\sqrt r}{\sqrt r}\le e^{\sqrt r}$,
а при $z=-x<0$ $f(-x)=\frac{\operatorname{sh}\sqrt x}{\sqrt x}\ge\frac{e^{\sqrt x}}{4\sqrt x}$ при $x\ge1$, так что
$\ln\ln M_f(r)/\ln r\to1/2$. При этом $f\ne\mathrm{const}$.
\item \emph{Теорема 16.5 неверна при $\rho=1/\alpha$.} Для угла раствора $\pi/n$ ($\alpha=1/n$, $n\in\NN$) и $f(z)=\exp(z^n)$ порядка $n=1/\alpha$:
на сторонах $\arg z=\pm\frac\pi{2n}$ имеем $z^n=\pm ir^n$, $|f|=1$, но на биссектрисе $f(r)=e^{r^n}\to\infty$. При $n=1$ это пример
$e^z$ в правой полуплоскости.
\item \emph{Следствие 1 для полуплоскостей неверно при $\rho=1$:} $\sin z$ ограничена на $\RR$ ($|\sin x|\le1$), но не ограничена
(и не постоянна); а для порядка $<1$ из Следствия 1 (в обеих полуплоскостях) и теоремы Лиувилля получаем: \emph{ц.ф. порядка $<1$,
ограниченная на прямой, постоянна.}
\end{enumerate}
\end{remarkbox}

\begin{remarkbox}{Коротко: план ответа на экзамене}
\begin{enumerate}
\item $M_f(r)$ не описывает рост по направлениям: пример $e^z$ ($e^{r\cos\vartheta}$).
\item Формулировка Теоремы 16.5: порядок $\le\rho$, угол раствора $\pi\alpha$, $\alpha<1/\rho$, $|f|\le C$ на сторонах $\Rightarrow$ внутри.
\item $S=\{|\arg z|<\pi\alpha/2\}$; $\rho<\beta<\gamma<1/\alpha$; $|f|<e^{r^\beta}$ на $|z|=R_n$.
\item $g=fe^{-\eps z^\gamma}$ (главная ветвь); $|e^{-\eps z^\gamma}|=e^{-\eps r^\gamma\cos\gamma\vartheta}\le e^{-\eps c_0 r^\gamma}$, $c_0=\cos(\gamma\pi\alpha/2)>0$.
\item На сторонах $|g|\le C$, на дугах $|g|\le e^{R_n^\beta-\eps c_0R_n^\gamma}\le C$ при больших $n$; принцип максимума в $S_n$.
\item $|f|\le Ce^{\eps|z|^\gamma}$, $\eps\to0$.
\item Следствие: порядок $<1/2$ и ограниченность на луче $\Rightarrow$ константа (через $h(w)=f(w^2)$ и Лиувилля); пример $\sin\sqrt z/\sqrt z$.
\item Точность: $e^{z^n}$ в угле раствора $\pi/n$.
\end{enumerate}
\end{remarkbox}

% =====================================================================
\ticketx{43*}{Целые функции экспоненциального типа. Преобразование Бореля. Сопряжённая диаграмма, индикатор и индикаторная диаграмма. Теорема Пойа и теорема об индикаторе производной. Применение преобразования Бореля к решению линейных ОДУ}

\textit{(Дополнительный материал: вопрос со звёздочкой, в билеты не входит.)}

\begin{defs}
\textbf{Ц.ф.э.т.} Целая функция $f$ называется \emph{целой функцией экспоненциального типа}, если её порядок $\rho<1$, или
$\rho=1$ и тип $\sigma$ относительно порядка $1$ конечен. Эквивалентно: $|f(z)|\le Ae^{B|z|}$ для некоторых $A,B>0$ (проверка ниже).

\textbf{Степень (тип) ц.ф.э.т.} $\displaystyle\kappa=\kappa_f=\limsup_{r\to\infty}\frac{\ln M_f(r)}{r}$. Это тип $f$ относительно порядка $1$;
величина определена для любой ц.ф.: при $\rho<1$ она равна $0$, при $\rho>1$ --- $+\infty$.

\textbf{Запись ряда Тейлора с факториалами:} $f(z)=\sum_{n\ge0}\frac{b_n}{n!}z^n$, т.е. $b_n=f^{(n)}(0)$.

\textbf{Преобразование Бореля} ц.ф.э.т. $f$: $\displaystyle\gamma_f(t)=\sum_{n=0}^\infty\frac{b_n}{t^{n+1}}$ (ряд Лорана в окрестности $\infty$).

\textbf{Выпуклое множество, выпуклая оболочка.} $G\subset\CC$ выпукло, если вместе с любыми двумя точками содержит соединяющий их
отрезок; выпуклая оболочка $\operatorname{conv}E$ --- наименьшее выпуклое множество, содержащее $E$ (для компакта $E$ это компакт).

\textbf{Опорная функция} выпуклого компакта $G$: $K_G(\alpha)=\max_{z\in G}\Re(ze^{-i\alpha})$, $\alpha\in[0,2\pi]$ (продолжаем
$2\pi$-периодически). $\Re(ze^{-i\alpha})=x\cos\alpha+y\sin\alpha$ --- ортогональная проекция $z=x+iy$ на направление луча
$L_\alpha=\{\arg z=\alpha\}$. Прямая $\{x\cos\alpha+y\sin\alpha=K_G(\alpha)\}$ называется \emph{опорной}: $G$ лежит по одну сторону от неё и
имеет с ней общую точку.

\textbf{Сопряжённая диаграмма} $D^*$ ц.ф.э.т. $f$ --- выпуклая оболочка множества особых точек $\gamma_f$; точнее, наименьший выпуклый
компакт, вне которого $\gamma_f$ продолжается до голоморфной функции (см. ниже). $K(\alpha)=K_{D^*}(\alpha)$.

\textbf{Индикатор} ц.ф.э.т. $f$: $\displaystyle h_f(\alpha)=\limsup_{r\to\infty}\frac{\ln|f(re^{i\alpha})|}{r}$, $\alpha\in[0,2\pi]$ --- характеристика роста
$f$ вдоль луча $L_\alpha$ (на слайдах --- $\ln M_f(re^{i\alpha})$ в том же смысле).

\textbf{Индикаторная диаграмма} $D$ --- выпуклый компакт, опорная функция которого есть $h_f$; по теореме Пойа $D=\clos{D^*}:=\{\bar z:z\in D^*\}$.
\end{defs}

\subsection{Степень ц.ф.э.т. через коэффициенты}

\begin{thmbox}{Утверждение (упражнение слайда 472)}
Для любой ц.ф. $f(z)=\sum b_nz^n/n!$ выполнено $\kappa_f=\limsup_{n\to\infty}|b_n|^{1/n}$ (обе части могут быть $+\infty$).
В частности, $f$ --- ц.ф.э.т. $\iff$ $\kappa_f<\infty$ $\iff$ $|f(z)|\le Ae^{B|z|}$.
\end{thmbox}
\begin{proof}
Положим $L=\limsup|b_n|^{1/n}$. Используем элементарные оценки $\big(\frac ne\big)^n\le n!\le en\big(\frac ne\big)^n$, $n\ge1$
(следуют из $\ln n!=\sum_{k\le n}\ln k$ и сравнения с $\int_1^n\ln x\,dx=n\ln n-n+1$), откуда $\frac{(n!)^{1/n}e}{n}\to1$.

\emph{$L\le\kappa$.} Пусть $\kappa<\infty$, $\eps>0$; тогда $M_f(r)\le Ae^{(\kappa+\eps)r}$ при всех $r>0$ (при больших $r$ --- по
определению $\kappa$, при малых --- за счёт выбора $A$). Неравенство Коши: $|b_n|/n!\le M_f(r)r^{-n}\le Ae^{(\kappa+\eps)r}r^{-n}$.
При $r=n/(\kappa+\eps)$ (минимум правой части) получаем $|b_n|\le A\,n!\,\big(\frac{e(\kappa+\eps)}n\big)^n$, откуда
$|b_n|^{1/n}\le A^{1/n}\frac{(n!)^{1/n}e}{n}(\kappa+\eps)\to\kappa+\eps$. Значит, $L\le\kappa+\eps$ для любого $\eps>0$.

\emph{$\kappa\le L$.} Пусть $L<\infty$, $\eps>0$; тогда $|b_n|\le B(L+\eps)^n$ при всех $n$ (с подходящим $B$), и
$M_f(r)\le\sum_nB\frac{(L+\eps)^nr^n}{n!}=Be^{(L+\eps)r}$, т.е. $\kappa\le L+\eps$.

Последнее утверждение: если $|f|\le Ae^{B|z|}$, то $\kappa\le B$; обратно, при $\kappa<\infty$ выше получено $M_f(r)\le Ae^{(\kappa+1)r}$.
Наконец, $\kappa<\infty$ $\iff$ ($\rho<1$, или $\rho=1$ и $\sigma<\infty$): при $\rho<1$ $\ln M_f(r)<r^{\rho+\eps}=o(r)$, $\kappa=0$; при $\rho=1$ $\kappa=\sigma$;
при $\rho>1$ $\ln M_f(r)>r^{1+\eps}$ для сколь угодно больших $r$, $\kappa=\infty$.
\end{proof}
\begin{remarkbox}{Связь с формулой для типа}
Это согласуется с формулой (16.5) билета 33 при $\rho=1$: $e\sigma=\limsup n|c_n|^{1/n}$, $c_n=b_n/n!$; так как
$n(n!)^{-1/n}\to e$ (формула Стирлинга), то $\limsup n|c_n|^{1/n}=e\limsup|b_n|^{1/n}$.
\end{remarkbox}

\subsection{Преобразование Бореля}

\begin{thmbox}{Сходимость преобразования Бореля}
Пусть $f$ --- ц.ф.э.т. степени $\kappa$. Ряд $\gamma_f(t)=\sum b_nt^{-n-1}$ сходится абсолютно и локально равномерно при $|t|>\kappa$,
расходится при $|t|<\kappa$; на окружности $|t|=\kappa$ у суммы ряда есть хотя бы одна особая точка (т.е. $\gamma_f$ не продолжается
голоморфно ни в какую окрестность $\{|t|\ge\kappa\}$). Кроме того, $\gamma_f(\infty)=0$ и $\gamma_f\equiv0\iff f\equiv0$.
\end{thmbox}
\begin{proof}
Подстановка $w=1/t$ даёт степенной ряд $w\sum b_nw^n$, радиус сходимости которого по формуле Коши--Адамара равен
$1/\limsup|b_n|^{1/n}=1/\kappa$. Значит, ряд сходится при $|t|>\kappa$ и расходится при $|t|<\kappa$, а так как на границе круга
сходимости степенного ряда всегда есть особая точка суммы, то на $|t|=\kappa$ есть особая точка $\gamma_f$ (при $\kappa=0$ особая точка
--- $t=0$, если $f\not\equiv0$: иначе $\gamma_f$ была бы голоморфна в $\Ch$, равна $0$ в $\infty$ и, по теореме Лиувилля, $\equiv0$).
Коэффициенты ряда Лорана определяются функцией однозначно, отсюда последнее утверждение.
\end{proof}

Преобразование Бореля линейно: $\gamma_{c_1f_1+c_2f_2}=c_1\gamma_{f_1}+c_2\gamma_{f_2}$ (почленное сложение рядов при $|t|>\max\kappa_j$).

\begin{remarkbox}{Примеры преобразования Бореля (проверка непосредственно по определению)}
\begin{enumerate}
\item $f=Ae^{az}$: $b_n=Aa^n$, $\gamma_f(t)=\sum_n\frac{Aa^n}{t^{n+1}}=\frac At\cdot\frac1{1-a/t}=\frac A{t-a}$, $|t|>|a|=\kappa$.
\item $f=\sin z$: $b_{2k+1}=(-1)^k$, $b_{2k}=0$; $\gamma(t)=\sum_k\frac{(-1)^k}{t^{2k+2}}=\frac1{t^2}\cdot\frac1{1+t^{-2}}=\frac1{t^2+1}$.
То же из примера 1: $\sin z=\frac{e^{iz}-e^{-iz}}{2i}$, $\gamma=\frac1{2i}\big(\frac1{t-i}-\frac1{t+i}\big)=\frac1{t^2+1}$.
\item $f=\cos z=\frac{e^{iz}+e^{-iz}}2$: $\gamma=\frac12\big(\frac1{t-i}+\frac1{t+i}\big)=\frac t{t^2+1}$.
\item $f=1+e^z$: $\gamma=\frac1t+\frac1{t-1}$.
\item $f=\operatorname{ch}z+\sin z$ (упражнение слайда 475): $\operatorname{ch}z=\frac{e^z+e^{-z}}2$ даёт $\frac12\big(\frac1{t-1}+\frac1{t+1}\big)=\frac t{t^2-1}$, итого
$\gamma(t)=\frac t{t^2-1}+\frac1{t^2+1}$. Особые точки --- простые полюсы $\pm1$ (вычеты $\frac12$) и $\pm i$ (вычеты $\pm\frac1{2i}$).
\item $f=z^m$: $b_m=m!$, остальные $0$; $\gamma=m!/t^{m+1}$ --- единственная особая точка $0$.
\end{enumerate}
\end{remarkbox}

\begin{remarkbox}{Интегральная формула Бореля (сверх лекций)}
При $\Re t>\kappa$ имеем $\gamma_f(t)=\int_0^\infty f(s)e^{-ts}\,ds$, т.е. преобразование Бореля --- это преобразование Лапласа
сужения $f$ на $[0,+\infty)$. \emph{Доказательство:} $\int_0^\infty s^ne^{-ts}ds=n!/t^{n+1}$ при $\Re t>0$; при $\Re t>\kappa+\eps$
имеем $\sum_n\int_0^\infty\frac{|b_n|s^n}{n!}|e^{-ts}|ds\le\int_0^\infty Be^{(\kappa+\eps)s}e^{-s\Re t}ds<\infty$ (оценка $|b_n|\le B(\kappa+\eps)^n$
из предыдущего пункта), так что интеграл от ряда равен сумме интегралов. Это объясняет связь метода ниже с операционным исчислением.
\end{remarkbox}

\subsection{Опорная функция выпуклого множества}

\begin{remarkbox}{Примеры (упражнения слайда 474)}
Максимум линейной функции $z\mapsto\Re(ze^{-i\alpha})$ на выпуклой оболочке конечного множества достигается в одной из его точек.
\begin{enumerate}
\item Круг $\{|z|\le R\}$: $\Re(ze^{-i\alpha})\le|z|\le R$, равенство при $z=Re^{i\alpha}$; $K(\alpha)\equiv R$.
\item Отрезок $[-id,id]$: $z=isd$, $s\in[-1,1]$, $\Re(isd(\cos\alpha-i\sin\alpha))=sd\sin\alpha$; $K(\alpha)=d|\sin\alpha|$.
\item Квадрат с центром $0$ и стороной $2d$ (вершины $d(\pm1\pm i)$): $\Re\big(d(\pm1\pm i)e^{-i\alpha}\big)=\pm d\cos\alpha\pm d\sin\alpha$, так что
$K(\alpha)=d(|\cos\alpha|+|\sin\alpha|)$; при $|\alpha|\le\pi/4$ это $d(\cos\alpha+\sin|\alpha|)=d\sqrt2\cos(\pi/4-|\alpha|)$, и $K$ $\frac\pi2$-периодична.
\item Квадрат с вершинами $\pm1,\pm i$: $K(\alpha)=\max\{|\cos\alpha|,|\sin\alpha|\}$.
\item Точка $\{a\}$, $a=a_1+ia_2$: $K(\alpha)=\Re(ae^{-i\alpha})=a_1\cos\alpha+a_2\sin\alpha$.
\end{enumerate}
\end{remarkbox}

\begin{thmbox}{Свойства опорной функции}
Пусть $G,G_1,G_2$ --- выпуклые компакты. Тогда:
(1) $K_G$ непрерывна (даже липшицева); (2) если $K_{G_1}\equiv K_{G_2}$, то $G_1=G_2$; точнее,
$G=\bigcap_\alpha\{z:\Re(ze^{-i\alpha})\le K_G(\alpha)\}$; (3) для $G_\eps=\{z:\dist(z,G)\le\eps\}$ имеем $K_{G_\eps}=K_G+\eps$;
(4) для $\clos G=\{\bar z:z\in G\}$ имеем $K_{\clos G}(\alpha)=K_G(-\alpha)$.
\end{thmbox}
\begin{proof}
(1) Пусть $G\subset\{|z|\le R\}$. Для $z\in G$: $\Re(ze^{-i\alpha})\le\Re(ze^{-i\beta})+|z||e^{-i\alpha}-e^{-i\beta}|\le K_G(\beta)+R|\alpha-\beta|$;
берём максимум по $z$ и меняем $\alpha,\beta$ местами: $|K_G(\alpha)-K_G(\beta)|\le R|\alpha-\beta|$.

(2) Включение «$\subset$» очевидно. Пусть $w\notin G$. Так как $G$ компакт, есть ближайшая к $w$ точка $w_0\in G$. Для $z\in G$ и
$s\in(0,1]$ точка $w_0+s(z-w_0)\in G$, поэтому $|w-w_0-s(z-w_0)|^2\ge|w-w_0|^2$, т.е. $-2s\Re\big((z-w_0)\overline{(w-w_0)}\big)+s^2|z-w_0|^2\ge0$;
деля на $s$ и устремляя $s\to0$, получаем $\Re\big((z-w_0)\overline{(w-w_0)}\big)\le0$. Положим $e^{i\alpha}=\frac{w-w_0}{|w-w_0|}$; тогда
$\Re(ze^{-i\alpha})\le\Re(w_0e^{-i\alpha})<\Re(w_0e^{-i\alpha})+|w-w_0|=\Re(we^{-i\alpha})$ при всех $z\in G$, т.е. $K_G(\alpha)<\Re(we^{-i\alpha})$,
и $w$ не лежит в пересечении полуплоскостей. Значит, $G$ восстанавливается по $K_G$.

(3) $G_\eps=G+\clos{B(0,\eps)}=\{z+u:z\in G,|u|\le\eps\}$ (для компакта $G$ расстояние достигается), и
$\max\Re((z+u)e^{-i\alpha})=K_G(\alpha)+\max_{|u|\le\eps}\Re(ue^{-i\alpha})=K_G(\alpha)+\eps$.

(4) $\Re(\bar ze^{-i\alpha})=\Re\big(\overline{ze^{i\alpha}}\big)=\Re(ze^{i\alpha})$, и максимум по $z\in G$ равен $K_G(-\alpha)$.
\end{proof}

\subsection{Сопряжённая диаграмма. Представление ц.ф.э.т. интегралом}

\begin{detail}{корректность определения $D^*$}
Пусть $\mathcal K$ --- семейство выпуклых компактов $K$, таких что $\gamma=\gamma_f$ (заданная рядом при $|t|>\kappa$) продолжается
до голоморфной функции $\gamma_K$ в $\CC\setminus K$. Круг $\{|t|\le\kappa\}\in\mathcal K$.
\emph{(а) Для $K_1,K_2\in\mathcal K$ множество $\CC\setminus(K_1\cup K_2)$ связно.} Если $p\notin K_j$, то $K_j$ отделяется от $p$ прямой (п.~(2) выше),
поэтому направления лучей из $p$, пересекающих $K_j$, образуют дугу окружности направлений длины $<\pi$; две такие дуги не покрывают
окружность, и некоторый луч из $p$ уходит в $\infty$, не задевая $K_1\cup K_2$. \emph{(б)} Поэтому $\gamma_{K_1}=\gamma_{K_2}$ на
$\CC\setminus(K_1\cup K_2)$ (совпадают около $\infty$, теорема единственности), и они склеиваются в голоморфную функцию на
$\CC\setminus(K_1\cap K_2)$. \emph{(в)} Тем же склеиванием $\gamma$ продолжается на дополнение к $D^*:=\bigcap_{K\in\mathcal K}K$; это наименьший
элемент $\mathcal K$, и он совпадает с выпуклой оболочкой множества особых точек (точек, в которые продолжение невозможно).
\end{detail}

\begin{thmbox}{Свойства $D^*$ (слайд 475)}
$D^*\subset\{|t|\le\kappa\}$, $D^*\cap\{|t|=\kappa\}\ne\varnothing$; $\max_\alpha K(\alpha)=\kappa$ (и если $\kappa e^{i\alpha_0}\in D^*$, то максимум
достигается при $\alpha_0$); $K(\alpha)\ge-\kappa$ при всех $\alpha$.
\end{thmbox}
\begin{proof}
Включение --- так как круг $\{|t|\le\kappa\}\in\mathcal K$. Если бы $D^*\subset\{|t|\le\kappa'\}$, $\kappa'<\kappa$, то $\gamma$ была бы голоморфна
при $|t|>\kappa'$, её ряд Лорана (тот же самый по единственности) сходился бы там, и $\limsup|b_n|^{1/n}\le\kappa'$ --- противоречие.
Далее, $K(\alpha)\le\max_{|t|\le\kappa}\Re(te^{-i\alpha})=\kappa$; при $\zeta_0=\kappa e^{i\alpha_0}\in D^*$: $K(\alpha_0)\ge\Re(\zeta_0e^{-i\alpha_0})=\kappa$.
Наконец, $\Re(te^{-i\alpha})\ge-|t|\ge-\kappa$ на $D^*$, откуда $K(\alpha)\ge-\kappa$.
\end{proof}

\begin{remarkbox}{Примеры сопряжённых диаграмм}
$Ae^{az}$: $D^*=\{a\}$. $\sin z$, $\cos z$: особые точки $\pm i$, $D^*=[-i,i]$. $1+e^z$: $D^*=[0,1]$. $\operatorname{ch}z+\sin z$: особые
точки $\pm1,\pm i$, $D^*$ --- квадрат с вершинами $\pm1,\pm i$. Многочлен: $D^*=\{0\}$.
\end{remarkbox}

\begin{thmbox}{Предложение 16.6}
Пусть $f$ --- ц.ф.э.т., $\gamma$ --- её преобразование Бореля, $D^*$ --- сопряжённая диаграмма. Тогда
\begin{equation}\label{t43:repr}
f(z)=\frac1{2\pi i}\int_\Gamma\gamma(t)e^{zt}\,dt,
\end{equation}
где $\Gamma$ --- любая замкнутая жорданова спрямляемая кривая (положительно ориентированная) такая, что $D^*\subset\operatorname{Int}\Gamma$.
\end{thmbox}
\begin{proof}
Пусть $R>\kappa$ и $R>\max_{t\in\Gamma}|t|$, $C_R=\{|t|=R\}$. Функция $\gamma(t)e^{zt}$ голоморфна в $\Omega=\CC\setminus D^*$. Циклы
$\Gamma$ и $C_R$ гомологичны в $\Omega$: для любой точки $a\in D^*=\CC\setminus\Omega$ имеем $\Ind_\Gamma(a)=1$ ($a\in\operatorname{Int}\Gamma$) и
$\Ind_{C_R}(a)=1$ ($|a|\le\kappa<R$). По общей интегральной теореме Коши $\int_\Gamma=\int_{C_R}$. На $C_R$ ряд $\sum b_nt^{-n-1}$
сходится равномерно, а $e^{zt}$ ограничена, поэтому можно интегрировать почленно:
\[
\int_{C_R}\gamma(t)e^{zt}dt=\sum_{n=0}^\infty b_n\int_{C_R}\frac{e^{zt}}{t^{n+1}}dt=2\pi i\sum_{n=0}^\infty b_n\frac{z^n}{n!}=2\pi if(z),
\]
так как по теореме о вычетах $\int_{C_R}e^{zt}t^{-n-1}dt=2\pi i\cdot[\text{коэффициент при }t^n\text{ в }e^{zt}]=2\pi i\,z^n/n!$.
\end{proof}

\begin{thmbox}{Следствие (оценка роста через опорную функцию)}
Пусть $K(\alpha)$ --- опорная функция $D^*$. Для любого $\eps>0$ существует $A(\eps)$, не зависящее от $r$ и $\alpha$, такое, что
\[
|f(re^{i\alpha})|\le A(\eps)\exp\big(r(K(-\alpha)+\eps)\big),\qquad r>0,\ \alpha\in[0,2\pi].
\]
\end{thmbox}
\begin{proof}
Пусть $D_\eps=\{t:\dist(t,D^*)\le\eps\}$ --- выпуклый компакт с непустой внутренностью, содержащей $D^*$; его граница $\Gamma=\partial D_\eps$
--- спрямляемая жорданова кривая (граница выпуклой области; её длина не больше $2\pi(\kappa+\eps)$), и $D^*\subset\operatorname{Int}\Gamma$.
По \eqref{t43:repr}, при $z=re^{i\alpha}$ и $t\in\Gamma$: $|e^{zt}|=e^{\Re(tz)}$ и $\Re(tre^{i\alpha})=r\Re(te^{-i(-\alpha)})\le rK_{D_\eps}(-\alpha)=r(K(-\alpha)+\eps)$
(свойство (3) опорной функции). Отсюда
\[
|f(z)|\le\frac{\ell(\Gamma)}{2\pi}\max_{\Gamma}|\gamma|\cdot e^{r(K(-\alpha)+\eps)},\qquad A(\eps)=\frac{\ell(\Gamma)}{2\pi}\max_\Gamma|\gamma|,
\]
причём $\max_\Gamma|\gamma|<\infty$, так как $\gamma$ непрерывна на компакте $\Gamma\subset\Omega$.
\end{proof}

\subsection{Индикатор и индикаторная диаграмма. Теорема Пойа}

Из следствия сразу получаем: $h_f(\alpha)\le K(-\alpha)+\eps$ для всех $\eps>0$, т.е. $h_f(\alpha)\le K(-\alpha)$. Теорема Пойа утверждает,
что на самом деле здесь равенство.

\begin{thmbox}{Теорема Пойа (без доказательства)}
Пусть $f$ --- ц.ф.э.т., $K(\cdot)$ --- опорная функция её сопряжённой диаграммы $D^*$, $h(\cdot)$ --- индикатор $f$. Тогда
$h(\alpha)=K(-\alpha)$ при всех $\alpha\in[0,2\pi]$.
\end{thmbox}

\begin{thmbox}{Следствия теоремы Пойа}
Индикатор $h$ ц.ф.э.т. --- непрерывная $2\pi$-периодическая функция; $\max_\alpha h(\alpha)=\kappa$; $h(\alpha)\ge-\kappa$; $h$ --- опорная функция
множества $D=\clos{D^*}=\{\bar z:z\in D^*\}$ (индикаторной диаграммы).
\end{thmbox}
\begin{proof}
Всё следует из свойств опорной функции и свойств $D^*$: $h(\alpha)=K(-\alpha)$ непрерывна; $\max_\alpha h=\max_\alpha K=\kappa$; $h\ge-\kappa$;
по свойству (4) $K_{\clos{D^*}}(\alpha)=K_{D^*}(-\alpha)=h(\alpha)$.
\end{proof}

\begin{remarkbox}{Примеры индикаторов (непосредственная проверка согласуется с теоремой Пойа)}
\begin{enumerate}
\item $f=e^{az}$, $a=a_1+ia_2$: при $z=re^{i\vartheta}$ $|e^{az}|=e^{\Re(az)}=\exp\big(r(a_1\cos\vartheta-a_2\sin\vartheta)\big)$, поэтому
$h_f(\alpha)=a_1\cos\alpha-a_2\sin\alpha$. По Пойа: $D^*=\{a\}$, $K(\alpha)=a_1\cos\alpha+a_2\sin\alpha$, $K(-\alpha)=a_1\cos\alpha-a_2\sin\alpha$ ---
совпадает. Индикаторная диаграмма $D=\{\bar a\}$. Порядок $1$ (при $a\ne0$), тип $\kappa=|a|$.
\item $f=\sin z$: порядок $\rho=1$, тип $\sigma=\kappa=1$. При $z=re^{i\alpha}$, $y=r\sin\alpha$: $\sin z=\frac{e^{iz}-e^{-iz}}{2i}$,
$|e^{\pm iz}|=e^{\mp y}$, поэтому $\frac{e^{|y|}-e^{-|y|}}2\le|\sin z|\le e^{|y|}$; при $\sin\alpha\ne0$ отсюда $h(\alpha)=|\sin\alpha|$.
При $\alpha=0,\pi$: $|\sin(\pm r)|\le1$ и $=1$ при $r=\pi/2+2\pi k$, так что $h=0=|\sin\alpha|$. По Пойа: $D^*=[-i,i]$,
$K(\alpha)=|\sin\alpha|=K(-\alpha)=h(\alpha)$; $D=\clos{[-i,i]}=[-i,i]$.
\item $f=1+e^z$: $D^*=D=[0,1]$, $K(\alpha)=\max\{0,\cos\alpha\}$, $h(\alpha)=\max\{0,\cos\alpha\}$. Непосредственно: при $\cos\alpha>0$
$|1+e^z|\sim e^{r\cos\alpha}$, при $\cos\alpha<0$ $|1+e^z|\to1$, при $\cos\alpha=0$ $|1+e^{\pm ir}|\le2$ и $=2$ при $r\in2\pi\NN$.
\item $f=\operatorname{ch}z+\sin z$: $h(\alpha)=\max\{|\cos\alpha|,|\sin\alpha|\}$, $D$ --- квадрат с вершинами $\pm1,\pm i$.
\end{enumerate}
\end{remarkbox}

\subsection{Преобразование Бореля производной. Индикатор производной}

\begin{thmbox}{Утверждение (упражнение слайда 478) и формула (16.21)}
Если $f$ --- ц.ф.э.т., $\gamma$ и $\gamma_1$ --- преобразования Бореля $f$ и $f'$, то $f'$ --- ц.ф.э.т. той же степени и
$\gamma_1(t)=t\gamma(t)-f(0)$. Для $\gamma_m$ --- преобразования Бореля $f^{(m)}$:
\[
\gamma_m(t)=t^m\gamma(t)-t^{m-1}f(0)-t^{m-2}f'(0)-\dots-f^{(m-1)}(0).
\]
\end{thmbox}
\begin{proof}
$f'(z)=\sum_{n\ge0}b_{n+1}z^n/n!$, т.е. коэффициенты $f'$ суть $b_{n+1}$, и $\limsup|b_{n+1}|^{1/n}=\limsup|b_n|^{1/n}=\kappa$ (так как
$|b_{n+1}|^{1/n}=(|b_{n+1}|^{1/(n+1)})^{(n+1)/n}$). Далее
\[
\gamma_1(t)=\sum_{n\ge0}\frac{b_{n+1}}{t^{n+1}}=t\sum_{m\ge1}\frac{b_m}{t^{m+1}}=t\Big(\gamma(t)-\frac{b_0}t\Big)=t\gamma(t)-f(0).
\]
Формула для $\gamma_m$ --- индукция: $\gamma_m=t\gamma_{m-1}-f^{(m-1)}(0)=t\big(t^{m-1}\gamma-\dots-f^{(m-2)}(0)\big)-f^{(m-1)}(0)$.
\end{proof}

\begin{thmbox}{Теорема об индикаторе производной (без доказательства)}
Пусть $f$ --- ц.ф.э.т., $D$ --- её индикаторная диаграмма. Пусть выполнено одно из условий: (i) $0$ --- внутренняя точка $D$;
(ii) $0\notin D$; (iii) $0$ --- внутренняя точка одного из отрезков, составляющих $\partial D$. Тогда $h_{f'}=h_f$.
\end{thmbox}

\begin{remarkbox}{Пример: $0$ в «угловой» точке диаграммы}
$f=1+e^z$, $f'=e^z$. $\gamma(t)=\frac1t+\frac1{t-1}$, $\gamma_1(t)=\frac1{t-1}$; проверка формулы: $t\gamma-f(0)=1+\frac t{t-1}-2=\frac1{t-1}$.
Диаграммы: у $f$ --- отрезок $[0,1]$ (точка $0$ --- его конец), у $f'$ --- точка $\{1\}$. Индикаторы различны:
$h_f(\alpha)=\max\{0,\cos\alpha\}$, $h_{f'}(\alpha)=\cos\alpha$ (они различаются при $\cos\alpha<0$). Причина видна из формулы
$\gamma_1=t\gamma-f(0)$: умножение на $t$ может «убить» особенность в точке $0$.
\end{remarkbox}

\subsection{Применение к линейным ОДУ с постоянными коэффициентами}

Рассмотрим задачу Коши
\begin{equation}\label{t43:ode}
a_0f^{(n)}+a_1f^{(n-1)}+\dots+a_{n-1}f'+a_nf=\phi,\qquad f^{(m)}(0)=b_m\ (m=0,\dots,n-1),
\end{equation}
где $a_j\in\CC$, $a_0\ne0$, $\phi$ --- данная ц.ф.э.т. с преобразованием Бореля $\gamma$. Ищем решение $f$ в классе ц.ф.э.т.; пусть $\mu$ ---
преобразование Бореля $f$. По формуле (16.21) преобразование Бореля левой части равно
\[
\sum_{j=0}^na_{n-j}\gamma_j(t)=P(t)\mu(t)-Q(t),\qquad P(t)=a_0t^n+a_1t^{n-1}+\dots+a_n,\quad
Q(t)=\sum_{j=1}^na_{n-j}\sum_{i=0}^{j-1}b_it^{j-1-i}.
\]
Коэффициент $Q$ при $t^{n-1-l}$ ($l=0,\dots,n-1$) равен $a_0b_l+a_1b_{l-1}+\dots+a_lb_0$; $\deg Q\le n-1$. (На слайдах $Q$ взят с
противоположным знаком: $P\mu+Q=\gamma$.)

\begin{thmbox}{Утверждение (метод преобразования Бореля)}
Функция $\mu(t)=\dfrac{\gamma(t)+Q(t)}{P(t)}$ голоморфна в окрестности $\infty$ и $\mu(\infty)=0$; если $\mu(t)=\sum_{k\ge0}\mu_kt^{-k-1}$ ---
её ряд Лорана в $\infty$, то $f(z)=\sum_{k\ge0}\frac{\mu_k}{k!}z^k$ --- решение задачи \eqref{t43:ode}, и это единственное решение в
классе ц.ф.э.т. Кроме того, $f(z)=\frac1{2\pi i}\int_{|t|=R}\mu(t)e^{zt}dt$ при достаточно большом $R$ (формула \eqref{t43:repr}).
\end{thmbox}
\begin{proof}
\emph{Единственность.} Если $f$ --- решение класса ц.ф.э.т., то преобразования Бореля обеих частей \eqref{t43:ode} совпадают:
$P\mu-Q=\gamma$, откуда $\mu=(\gamma+Q)/P$ однозначно, а $\mu$ определяет все $b_n=f^{(n)}(0)$, т.е. $f$.

\emph{Существование.} Пусть $R_0$ больше $\kappa_\phi$ и модулей всех корней $P$. Тогда $\mu=(\gamma+Q)/P$ голоморфна при $|t|>R_0$, и
$\mu(t)\to0$ при $t\to\infty$, так как $\gamma(\infty)=0$, а $\deg Q<\deg P$. Значит, $\mu(t)=\sum_k\mu_kt^{-k-1}$ при $|t|>R_0$, и по
Коши--Адамару $\limsup|\mu_k|^{1/k}\le R_0$; поэтому $f=\sum\mu_kz^k/k!$ --- ц.ф.э.т. с преобразованием Бореля $\mu$. Проверим начальные
условия: выделим в равенстве $P\mu=\gamma+Q$ «полиномиальные части» (неотрицательные степени $t$) рядов Лорана в $\infty$. У $\gamma$
их нет, так что полиномиальная часть $P\mu$ равна $Q=Q_b$. С другой стороны, полиномиальная часть $P(t)\sum_k\mu_kt^{-k-1}$ равна
$\sum_ja_{n-j}\sum_{k\le j-1}\mu_kt^{j-1-k}=Q_\mu$ --- тот же многочлен, что $Q$, но с $\mu_k$ вместо $b_k$. Из $Q_\mu=Q_b$ и треугольности
(коэффициент при $t^{n-1-l}$ равен $a_0c_l+\dots+a_lc_0$, $a_0\ne0$) по индукции по $l$ получаем $\mu_l=b_l$, т.е. $f^{(l)}(0)=b_l$, $l<n$.
Тогда преобразование Бореля левой части \eqref{t43:ode} для $f$ равно $P\mu-Q_b=\gamma$, т.е. левая часть равна $\phi$ (преобразование
Бореля инъективно). Последняя формула --- Предложение 16.6 (при $R>R_0$ окружность охватывает все особенности $\mu$).
\end{proof}
Предложенный метод составляет основу т.н. операционного исчисления. Для нахождения $f$ можно либо разложить $\mu$ в ряд
Лорана, либо вычислить интеграл \eqref{t43:repr} по теореме о вычетах: $f(z)=\sum\Res\big(\mu(t)e^{zt}\big)$ по всем особым точкам $\mu$.

\begin{remarkbox}{Пример 1 (слайд 480): $f''-2f'+f=0$, $f(0)=0$, $f'(0)=1$}
Преобразования Бореля: $f\mapsto\mu$, $f'\mapsto t\mu-f(0)=t\mu$, $f''\mapsto t^2\mu-tf(0)-f'(0)=t^2\mu-1$. Уравнение:
$t^2\mu-1-2t\mu+\mu=0$, $\mu(t)=\frac1{(t-1)^2}$. Сопряжённая диаграмма --- $\{1\}$; берём $\Gamma=\{|t-1|=\delta\}$; по формуле Коши для производной
\[
f(z)=\frac1{2\pi i}\int_{|t-1|=\delta}\frac{e^{zt}}{(t-1)^2}dt=\frac{d}{dt}e^{zt}\Big|_{t=1}=ze^z.
\]
Через ряд Лорана: $\frac1{(t-1)^2}=\frac1{t^2}(1-t^{-1})^{-2}=\sum_{k\ge1}\frac k{t^{k+1}}$, $\mu_k=k$, $f=\sum_{k\ge1}\frac{kz^k}{k!}=z\sum_{k\ge1}\frac{z^{k-1}}{(k-1)!}=ze^z$.
\emph{Проверка:} $f'=(1+z)e^z$, $f''=(2+z)e^z$, $f''-2f'+f=e^z(2+z-2-2z+z)=0$; $f(0)=0$, $f'(0)=1$.
\end{remarkbox}

\begin{remarkbox}{Пример 2 (слайд 481): $f''+f'+f=\cos z$, $f(0)=0$, $f'(0)=1$}
Преобразование Бореля $\cos z$ равно $\frac t{t^2+1}$. Уравнение: $(t^2\mu-1)+t\mu+\mu=\frac t{t^2+1}$, т.е. $(t^2+t+1)\mu=1+\frac t{t^2+1}=\frac{t^2+t+1}{t^2+1}$,
\[
\mu(t)=\frac1{t^2+1}.
\]
(Корни $t^2+t+1$ сократились --- «свободных колебаний» в решении нет.) Сопряжённая диаграмма --- $[-i,i]$; $\Gamma=\{|t|=2\}$. По теореме о вычетах
\[
f(z)=\frac1{2\pi i}\int_{|t|=2}\frac{e^{zt}}{t^2+1}dt=\Res_{t=i}+\Res_{t=-i}=\frac{e^{iz}}{2i}-\frac{e^{-iz}}{2i}=\sin z
\]
(это и пример 2 преобразований Бореля: $\gamma_{\sin}=\frac1{t^2+1}$).
\emph{Проверка:} $(\sin z)''+(\sin z)'+\sin z=-\sin z+\cos z+\sin z=\cos z$; $\sin0=0$, $\cos0=1$.
\end{remarkbox}

\begin{remarkbox}{Пример 3 (резонанс, сверх лекций): $f''+f=\cos z$, $f(0)=f'(0)=0$}
$t^2\mu+\mu=\frac t{t^2+1}$, $\mu=\frac t{(t^2+1)^2}$ --- двукратные полюсы $\pm i$. Ряд Лорана:
$\frac t{(t^2+1)^2}=t^{-3}(1+t^{-2})^{-2}=\sum_{k\ge0}(-1)^k(k+1)t^{-2k-3}$, т.е. $\mu_{2k+2}=(-1)^k(k+1)$, остальные $\mu_j=0$, и
\[
f(z)=\sum_{k\ge0}\frac{(-1)^k(k+1)}{(2k+2)!}z^{2k+2}=\frac z2\sum_{k\ge0}\frac{(-1)^kz^{2k+1}}{(2k+1)!}=\frac{z\sin z}2 .
\]
Через вычеты: $\Res_{t=i}\frac{te^{zt}}{(t-i)^2(t+i)^2}=\frac d{dt}\frac{te^{zt}}{(t+i)^2}\Big|_{t=i}=e^{iz}\Big(\frac{iz}{-4}+\frac1{-4}-\frac{2i}{-8i}\Big)=-\frac{iz}4e^{iz}$,
аналогично $\Res_{t=-i}=\frac{iz}4e^{-iz}$; сумма $\frac{iz}4(e^{-iz}-e^{iz})=\frac{z\sin z}2$.
\emph{Проверка:} $f''=\cos z-\frac z2\sin z$, $f''+f=\cos z$; $f(0)=0$, $f'(0)=\big(\frac{\sin z+z\cos z}2\big)\big|_0=0$.
\end{remarkbox}

\begin{remarkbox}{Коротко: план ответа}
\begin{enumerate}
\item Ц.ф.э.т.: порядок $<1$ или $=1$ с конечным типом; степень $\kappa=\limsup\frac{\ln M_f(r)}r=\limsup|b_n|^{1/n}$, $f=\sum b_nz^n/n!$.
\item Преобразование Бореля $\gamma_f(t)=\sum b_nt^{-n-1}$, $|t|>\kappa$; особая точка на $|t|=\kappa$; примеры $e^{az}$, $\sin$, $\cos$, $1+e^z$.
\item Опорная функция $K_G(\alpha)=\max_G\Re(ze^{-i\alpha})$; примеры (круг, отрезок, квадрат); $G$ определяется по $K_G$.
\item Сопряжённая диаграмма $D^*$; $f=\frac1{2\pi i}\int_\Gamma\gamma(t)e^{zt}dt$; оценка $|f(re^{i\alpha})|\le A(\eps)e^{r(K(-\alpha)+\eps)}$.
\item Индикатор $h_f(\alpha)$; теорема Пойа $h(\alpha)=K(-\alpha)$; индикаторная диаграмма $D=\clos{D^*}$; пример $\sin z$: $h=|\sin\alpha|$, $D=D^*=[-i,i]$.
\item $\gamma_{f'}=t\gamma_f-f(0)$, (16.21); теорема об индикаторе производной; контрпример $1+e^z$.
\item ОДУ: $P\mu-Q=\gamma$, $\mu=(\gamma+Q)/P$, $f=\frac1{2\pi i}\int\mu e^{zt}dt$; примеры $ze^z$ и $\sin z$.
\end{enumerate}
\end{remarkbox}

% ===== part 9 =====
% ===================== Часть 9: вопрос 42* =====================
\ticketx{42*}{Модулярная функция. Доказательство малой теоремы Пикара в общем случае}
\label{t42:start}

\textbf{Идея.} Мы построим голоморфную функцию $\mu\colon\DD\to\Omega:=\CC\setminus\{0,1\}$ (\emph{модулярную функцию}),
которая устроена как <<накрытие>>: каждый круг $B\subset\Omega$ имеет прообраз, распадающийся на счётное число
непересекающихся областей, каждая из которых конформно отображается на $B$. Тогда обратная функция $\mu^{-1}$ --- это
П.А.Ф., продолжаемая вдоль \emph{любого} пути в $\Omega$, со значениями в $\DD$. Если целая $f$ не принимает значений $0$ и $1$,
то $f(\CC)\subset\Omega$, и композиция $\mu^{-1}\circ f$ по теореме о монодромии даёт \emph{однозначную} целую функцию,
ограниченную единицей; по теореме Лиувилля она постоянна, а значит, постоянна и $f$.
Материал --- слайды 482--491 (глава 17 <<Дополнения>>).

\begin{defs}
\begin{itemize}
\item \textbf{Обобщённая окружность} в $\Ch$ --- окружность в $\CC$ или прямая вместе с точкой $\infty$.
\textbf{ДЛО} (дробно-линейное отображение) $L(z)=\frac{az+b}{cz+d}$, $ad-bc\ne0$, --- гомеоморфизм $\Ch$ на себя,
конформный, переводящий обобщённые окружности в обобщённые окружности. ДЛО $L$ с вещественными $a,b,c,d$ и $ad-bc>0$
--- это в точности конформные автоморфизмы верхней полуплоскости $\CC_+=\{\Im w>0\}$; $\CC_-=\{\Im w<0\}$, $\widehat\RR=\RR\cup\{\infty\}$.
\item \textbf{Симметрия относительно обобщённой окружности $\Gamma$.} Для $\Gamma=\widehat\RR$ это $z\mapsto\bar z$; для
$\Gamma=\{|z-c|=R\}$ --- инверсия $\sigma_\Gamma(z)=c+\dfrac{R^2}{\bar z-\bar c}$ ($c\leftrightarrow\infty$); для прямой --- зеркальное отражение.
Точки $z$ и $z^*=\sigma_\Gamma(z)$ называются \emph{симметричными} относительно $\Gamma$. Свойства: $\sigma_\Gamma$ --- инволюция $\Ch$,
неподвижна ровно на $\Gamma$, меняет местами две компоненты $\Ch\setminus\Gamma$; если $K$ --- ДЛО с $K(\Gamma)=\widehat\RR$, то
$\sigma_\Gamma=K^{-1}\circ\overline{(\,\cdot\,)}\circ K$ (результат не зависит от выбора $K$). Пример: $\sigma_{\TT}(z)=1/\bar z$.
\item Окружности $\Gamma_1=\{|z-c|=R\}$ и $\TT$ \textbf{ортогональны} ($\Gamma_1\perp\TT$) iff $|c|^2=1+R^2$. Если дуга окружности
$\Gamma\perp\TT$ соединяет точки $e^{i\alpha},e^{i\beta}\in\TT$ ($0<\beta-\alpha<\pi$), то центр $\Gamma$ --- $e^{i(\alpha+\beta)/2}/\cos\frac{\beta-\alpha}2$,
радиус --- $\operatorname{tg}\frac{\beta-\alpha}2$. В $\CC_+$ обобщённые окружности, ортогональные $\widehat\RR$, --- это окружности с центром на $\RR$
и вертикальные прямые; их части в $\CC_+$ называют \emph{геодезическими}.
\item $\mathrm{Conf}(D,G)$ --- множество конформных (голоморфных и биективных) отображений области $D$ на область $G$.
Функция $f\in H(D)$ \textbf{локально конформна}, если $f'(z)\ne0$ при всех $z\in D$ (эквивалентно: $f$ локально инъективна, билет~1).
\item \textbf{Круговой треугольник с нулевыми углами} (\emph{идеальный треугольник}) в $\DD$ --- область, ограниченная тремя дугами окружностей,
ортогональных $\TT$, попарно соединяющими три точки $\TT$ (\emph{вершины}); соседние стороны в вершине касаются (обе ортогональны $\TT$),
т.\,е. углы треугольника равны $0$. Аналогично в $\CC_+$: например,
$T_0'=\{0<\Re w<1,\ |w-\tfrac12|>\tfrac12\}$ с вершинами $0,1,\infty$.
\item \textbf{Жорданова область} --- ограниченная область, граница которой --- жорданова (простая замкнутая) кривая.
\item \textbf{Элемент} $F=(B,f)$: круг $B=B(a,R)$ и $f\in H(B)$; канонический, если $B$ --- круг сходимости ряда Тейлора $f$ в $a$;
$E_a(f)$ --- канонический элемент функции $f$ в точке $a$. \textbf{Продолжение вдоль пути} $\gamma\colon[0,1]\to\CC$ --- семейство
канонических элементов $F_t=(B(\gamma(t),R_t),f_t)$ такое, что для любого $t_0$ найдётся $\delta>0$: при $|t-t_0|<\delta$
$\gamma(t)\in B_{t_0}$ и $f_t=f_{t_0}$ на $B_t\cap B_{t_0}$ (билет 25). \textbf{П.А.Ф.}, порождённая элементом $F_a$, --- совокупность
всех его продолжений вдоль всех путей, вдоль которых продолжение возможно (билет 27).
\item Пусть $\Phi$ --- П.А.Ф., $V(a,\delta)=\{0<|z-a|<\delta\}$, и элемент $F_b\in\Phi$, $b\in V(a,\delta)$, продолжается вдоль любого пути
в $V(a,\delta)$. Точка $a$ --- \textbf{логарифмическая точка ветвления} (для аналитической ветви, порождённой $F_b$), если после любого
числа $n\ge1$ обходов вокруг $a$ в $V(a,\delta)$ элемент $F_b$ не переходит в себя (билет 27, Определение 15.9).
\item Пусть $f\in H(\DD)$. Точка $\zeta\in\TT$ \textbf{правильная} для $f$, если существуют $r>0$ и $F\in H(B(\zeta,r))$ с $F=f$
на $B(\zeta,r)\cap\DD$. Если ни одна точка $\TT$ не правильная, $\TT$ называется \textbf{естественной границей} $f$.
\item $f$ \textbf{мероморфна в $\CC$}, если она голоморфна в $\CC$ вне локально конечного множества точек, каждая из которых --- полюс;
в полюсе полагаем $f=\infty$. Значение $a\in\Ch$, не принимаемое непостоянной функцией $f$, называется \textbf{исключительным
пикаровским значением} (ИПЗ). Пример: $e^z$ имеет ИПЗ $0$ (как целая) и ИПЗ $0,\infty$ (как мероморфная); $\operatorname{tg}z$ имеет ИПЗ $\pm i$.
\end{itemize}
\end{defs}

\subsection{Симметрия относительно окружностей и принцип симметрии}

\begin{thmbox}{Лемма 42.1 (симметрия и ДЛО)}
\begin{enumerate}[label=\alph*)]
\item Если $L$ --- ДЛО, то $L\circ\sigma_\Gamma=\sigma_{L(\Gamma)}\circ L$, т.\,е. точки, симметричные относительно $\Gamma$, переходят в точки,
симметричные относительно $L(\Gamma)$.
\item $\sigma_\Gamma$ переводит обобщённые окружности в обобщённые окружности и сохраняет величины углов между кривыми
(меняя ориентацию); в частности, сохраняет ортогональность.
\item Если $\Gamma'\perp\Gamma$, то $\sigma_\Gamma(\Gamma')=\Gamma'$. В частности, если $\Gamma\perp\TT$, то $\sigma_\Gamma(\TT)=\TT$ и $\sigma_\Gamma(\DD)=\DD$;
если $\Gamma\perp\widehat\RR$, то $\sigma_\Gamma(\CC_+)=\CC_+$.
\end{enumerate}
\end{thmbox}
\begin{proof}
Сначала проверим корректность представления $\sigma_\Gamma=K^{-1}\circ\bar{\ }\circ K$. Если $K_1,K_2$ --- ДЛО, переводящие $\Gamma$ в $\widehat\RR$,
то $N=K_2\circ K_1^{-1}$ переводит $\widehat\RR$ в $\widehat\RR$, поэтому (после деления на общий множитель) имеет вещественные
коэффициенты и $N(\bar z)=\overline{N(z)}$. Отсюда $K_2^{-1}\bar{\ }K_2=K_1^{-1}N^{-1}\,\bar{\ }\,NK_1=K_1^{-1}\bar{\ }K_1$.
(Для окружности $\{|z-c|=R\}$ и $K(z)=\frac{z-c-R}{z-c+R}\cdot(-i)$ прямое вычисление даёт формулу инверсии из рамки определений.)

a) Если $K(\Gamma)=\widehat\RR$, то $K\circ L^{-1}$ переводит $L(\Gamma)$ в $\widehat\RR$, поэтому
$\sigma_{L(\Gamma)}=(KL^{-1})^{-1}\,\bar{\ }\,(KL^{-1})=L\,\sigma_\Gamma\,L^{-1}$.

b) $\sigma_\Gamma$ --- композиция ДЛО $K$, сопряжения и ДЛО $K^{-1}$; каждое из них переводит обобщённые окружности в обобщённые окружности
и сохраняет величины углов (сопряжение меняет их знак).

c) $K(\Gamma')\perp\widehat\RR$, т.\,е. $K(\Gamma')$ --- окружность с центром на $\RR$ или вертикальная прямая; такие множества
инвариантны относительно сопряжения. Значит, $\sigma_\Gamma(\Gamma')=K^{-1}(\overline{K(\Gamma')})=\Gamma'$. Далее, $\sigma_\Gamma$ переводит
$\Ch\setminus\TT$ в себя, точки $\Gamma\cap\DD$ неподвижны, значит компонента $\DD$ переходит в компоненту, содержащую эти точки, т.\,е. в $\DD$.
Аналогично для $\CC_+$.
\end{proof}

Нам понадобится принцип симметрии в форме лекций (Теорема 13.1, билет 5).

\begin{thmbox}{Теорема 13.1 (принцип симметрии Римана--Шварца)}
Пусть $D$ и $W$ --- области в $\Ch$; $\partial D$ и $\partial W$ содержат дуги $\gamma$ и $\sigma$ обобщённых окружностей $\Gamma$ и $\Sigma$
(непустые, открытые и связные в $\Gamma$, $\Sigma$), причём $D$ и $W$ лежат по одну сторону от $\Gamma$ и $\Sigma$ соответственно.
Пусть $\widetilde D=D\cup\gamma\cup D^*$ и $\widetilde W=W\cup\sigma\cup W^*$ --- области, где $D^*=\sigma_\Gamma(D)$, $W^*=\sigma_\Sigma(W)$.
Пусть $f\in\mathrm{Conf}(D,W)$, $f\in C(D\cup\gamma)$ и $f$ гомеоморфно отображает $\gamma$ на $\sigma$. Тогда существует единственная
$\tilde f\in\mathrm{Conf}(\widetilde D,\widetilde W)$ с $\tilde f|_D=f$; при этом $\tilde f(\sigma_\Gamma(z))=\sigma_\Sigma(\tilde f(z))$, в частности
$\tilde f(D^*)=W^*$.
\end{thmbox}
\begin{proof}
\emph{Сведение к прямой.} Пусть $\gamma\ne\Gamma$ (в нашем приложении это так; случай $\gamma=\Gamma$ разбирается аналогично после
разбиения $\Gamma$ на две дуги). Возьмём ДЛО $\varphi$ с $\varphi(\Gamma)=\widehat\RR$, переводящее какую-нибудь точку $\Gamma\setminus\gamma$ в $\infty$;
тогда $\varphi(\gamma)\subset\RR$ --- интервал, а связная область $\varphi(D)$, не пересекающая $\widehat\RR$, лежит в $\CC_+$ или в $\CC_-$; заменив
при необходимости $\varphi$ на $-\varphi$, получим $\varphi(D)\subset\CC_+$. Аналогично строим $\psi$ для $\Sigma,\sigma,W$. По Лемме~42.1\,a)
$\varphi(D^*)=\overline{\varphi(D)}$, $\psi(W^*)=\overline{\psi(W)}$ (черта --- сопряжение). Функция $h=\psi\circ f\circ\varphi^{-1}$ удовлетворяет условиям
теоремы для $\Gamma=\Sigma=\widehat\RR$; если $\tilde h$ --- нужное продолжение $h$, то $\tilde f=\psi^{-1}\circ\tilde h\circ\varphi$ --- нужное
продолжение $f$, и $\tilde f\circ\sigma_\Gamma=\psi^{-1}\tilde h\,\bar{\ }\,\varphi=\psi^{-1}\,\bar{\ }\,\tilde h\varphi=\sigma_\Sigma\circ\tilde f$ (снова Лемма 42.1\,a).

\emph{Случай прямой}: $D\subset\CC_+$, $\gamma\subset\RR$, $W\subset\CC_+$, $f(\gamma)=\sigma\subset\RR$. Положим $\tilde f=f$ на $D\cup\gamma$ и
$\tilde f(z)=\overline{f(\bar z)}$ на $D^*$. В круге $B(\bar a,r)\subset D^*$ имеем $f(\bar z)=\sum c_k(\bar z-a)^k$, откуда
$\tilde f(z)=\sum\bar c_k(z-\bar a)^k$ --- голоморфна в $D^*$. Функция $\tilde f$ непрерывна в $\widetilde D$, так как $f$ непрерывна вплоть до $\gamma$
и вещественна на $\gamma$. Голоморфность в точках $\gamma$ --- по теореме Морера: для треугольника $\Delta\subset\widetilde D$, пересекающего $\RR$,
отрезаем от него полоску $\{|\Im z|<\delta\}$; интегралы по частям вне полоски равны нулю по теореме Коши, а интеграл по границе части
внутри полоски стремится к нулю при $\delta\to0$ по равномерной непрерывности $\tilde f$ на компакте (подробно --- билет 5).
Однолистность: при $z_1,z_2\in D$ или $z_1,z_2\in D^*$ --- из однолистности $f$; при $z_1\in D\cup\gamma$, $z_2\in D^*$ значения лежат в
$W\cup\sigma$ и $W^*$, которые не пересекаются. Наконец, $\tilde f(D^*)=\overline{f(D)}=W^*$, т.\,е. $\tilde f(\widetilde D)=\widetilde W$.
Единственность --- по теореме единственности.
\end{proof}

\subsection{Треугольник $T_0$ и отображение $\mu_0$}

Следуя слайду 483, положим $\omega=e^{2\pi i/3}$ и
\[
a_0^1=e^{-i\pi/6},\qquad a_0^2=i=\omega a_0^1,\qquad a_0^3=e^{7\pi i/6}=\omega^2a_0^1
\]
--- точки, делящие $\TT$ на три равные дуги (они идут против часовой стрелки). Через $a_0^1,a_0^2$ проведём окружность $\Gamma_0^1\perp\TT$:
её центр $c_1=2e^{i\pi/6}$, радиус $\sqrt3$ (действительно, $|c_1|^2=4=1+3$, $|a_0^1-c_1|=|a_0^2-c_1|=\sqrt3$). Положим
$\Gamma_0^2=\omega\Gamma_0^1$ (через $a_0^2,a_0^3$), $\Gamma_0^3=\omega^2\Gamma_0^1$ (через $a_0^3,a_0^1$), $\gamma_0^k=\Gamma_0^k\cap\DD$,
$B_0^k$ --- открытый круг, ограниченный $\Gamma_0^k$, и
\[
T_0=\DD\setminus\bigl(\clos{B_0^1}\cup\clos{B_0^2}\cup\clos{B_0^3}\bigr)\ni 0 .
\]
Это и есть открытый круговой треугольник, ограниченный дугами $\gamma_0^1,\gamma_0^2,\gamma_0^3$. Очевидно, $\omega T_0=T_0$.

\textbf{Переход в полуплоскость.} Рассмотрим ДЛО
\[
\Phi(z)=\frac{(z-a_0^1)(a_0^2-a_0^3)}{(z-a_0^3)(a_0^2-a_0^1)},\qquad \Phi(a_0^1)=0,\ \ \Phi(a_0^2)=1,\ \ \Phi(a_0^3)=\infty .
\]
\begin{detail}{свойства $\Phi$}
Прямое вычисление: $\dfrac{a_0^2-a_0^3}{a_0^2-a_0^1}=\dfrac{\sqrt3+3i}{-\sqrt3+3i}=\dfrac{1-i\sqrt3}{2}=e^{-i\pi/3}$ и
$a_0^1/a_0^3=e^{-4\pi i/3}=e^{2\pi i/3}$, поэтому $\Phi(0)=e^{2\pi i/3}e^{-i\pi/3}=e^{i\pi/3}\in\CC_+$, $\Phi(\infty)=e^{-i\pi/3}$.
Так как $\Phi$ переводит три точки $\TT$ в три точки $\widehat\RR$, то $\Phi(\TT)=\widehat\RR$, и связная область $\DD$ переходит в компоненту
$\Ch\setminus\widehat\RR$, содержащую $\Phi(0)$: $\Phi(\DD)=\CC_+$. Окружность $\Phi(\Gamma_0^1)$ проходит через $0,1$ и ортогональна $\widehat\RR$
(конформность), т.\,е. это $\{|w-\frac12|=\frac12\}$; аналогично $\Phi(\Gamma_0^2)=\{\Re w=1\}\cup\{\infty\}$, $\Phi(\Gamma_0^3)=\{\Re w=0\}\cup\{\infty\}$.
Круг $B_0^k$ не содержит $0$ ($|c_k|=2>\sqrt3$), поэтому $\Phi(B_0^k)$ --- та из двух областей, ограниченных $\Phi(\Gamma_0^k)$, которая не содержит
$e^{i\pi/3}=\frac12+\frac{\sqrt3}2i$: $\Phi(B_0^1)=\{|w-\frac12|<\frac12\}$, $\Phi(B_0^2)=\{\Re w>1\}$, $\Phi(B_0^3)=\{\Re w<0\}$. Так как $\Phi$ --- биекция $\Ch$,
\begin{gather*}
\Phi(T_0)=\CC_+\setminus\bigl(\{|w-\tfrac12|\le\tfrac12\}\cup\{\Re w\ge1\}\cup\{\Re w\le0\}\bigr)=T_0',\\
\Phi(\gamma_0^1)=S(0,1),\quad \Phi(\gamma_0^2)=\{\Re w=1\}\cap\CC_+,\quad \Phi(\gamma_0^3)=\{\Re w=0\}\cap\CC_+,
\end{gather*}
где $S(0,1)=\{|w-\frac12|=\frac12,\ \Im w>0\}$. В частности, $\partial T_0$ --- жорданова кривая (три дуги, пересекающиеся лишь в концах),
и при обходе $\partial T_0$ против часовой стрелки вершины идут в порядке $a_0^1,a_0^2,a_0^3$.
\end{detail}

\begin{thmbox}{Предложение 42.2 (нормированное отображение $\mu_0$; слайд 485)}
Существует единственный гомеоморфизм $\mu_0\colon\clos{T_0}\to\clos{\CC_+}=\CC_+\cup\widehat\RR$ (замыкание в $\Ch$), конформно отображающий $T_0$
на $\CC_+$, с условиями $\mu_0(a_0^1)=0$, $\mu_0(a_0^2)=1$, $\mu_0(a_0^3)=\infty$. При этом
\[
\mu_0(\gamma_0^1)=(0,1)=:I_0^1,\qquad \mu_0(\gamma_0^2)=(1,+\infty)=:I_0^2,\qquad \mu_0(\gamma_0^3)=(-\infty,0)=:I_0^3 .
\]
\end{thmbox}
\begin{proof}
По теореме Римана (билеты 10--11) существует $F_0\in\mathrm{Conf}(T_0,\DD)$ с $F_0(0)=0$, $F_0'(0)>0$, и она единственна.
Функция $G(z)=\omega^{-1}F_0(\omega z)$ также лежит в $\mathrm{Conf}(T_0,\DD)$ (так как $\omega T_0=T_0$), $G(0)=0$, $G'(0)=F_0'(0)>0$;
по единственности $G=F_0$, т.\,е. $F_0(\omega z)=\omega F_0(z)$.
Так как $\partial T_0$ --- жорданова кривая, по теореме Каратеодори (Теорема 13.8, билеты 19--20) обратное отображение $F_0^{-1}\in\mathrm{Conf}(\DD,T_0)$
продолжается до непрерывного инъективного отображения $\clos\DD\to\clos{T_0}$, т.\,е. (компактность) до гомеоморфизма; значит, $F_0$
продолжается до гомеоморфизма $\clos{T_0}\to\clos\DD$, $F_0(\partial T_0)=\TT$. По непрерывности $F_0(a_0^{k+1})=F_0(\omega a_0^k)=\omega F_0(a_0^k)$.
Пусть $F_0(a_0^1)=e^{i\alpha}a_0^1$; тогда $F_0(a_0^k)=e^{i\alpha}a_0^k$ при $k=1,2,3$. Отображение $F=e^{-i\alpha}F_0$ --- гомеоморфизм
$\clos{T_0}\to\clos\DD$, конформный внутри, с $F(a_0^k)=a_0^k$. Положим $\mu_0=\Phi\circ F$: это гомеоморфизм $\clos{T_0}\to\Phi(\clos\DD)=\clos{\CC_+}$,
конформный на $T_0$, $\mu_0(a_0^1)=\Phi(a_0^1)=0$ и т.\,д.

Гомеоморфизм $\partial T_0\to\widehat\RR$ переводит открытую дугу $\gamma_0^1$ границы между $a_0^1$ и $a_0^2$, не содержащую $a_0^3$, в открытую
дугу $\widehat\RR$ между $0$ и $1$, не содержащую $\infty$, т.\,е. в $(0,1)$; аналогично для $\gamma_0^2$, $\gamma_0^3$.

Единственность: если $\mu_0'$ --- другое такое отображение, то $\mu_0'\circ\mu_0^{-1}$ --- конформный автоморфизм $\CC_+$, т.\,е. вещественное ДЛО,
непрерывно продолжающееся на $\widehat\RR$ с тремя неподвижными точками $0,1,\infty$; ДЛО с тремя неподвижными точками тождественно.
\end{proof}

\begin{remarkbox}{Замечание}
Можно сослаться и на <<граничное условие нормировки>> (следствие теоремы Каратеодори, билет 20): для жордановых областей и троек граничных
точек с одинаковым циклическим порядком существует единственное конформное отображение с заданным соответствием трёх точек.
Использованная выше симметрия $T_0$ позволяет не проверять отдельно совпадение циклических порядков.
\end{remarkbox}

\subsection{Треугольники всех поколений. Исчерпание круга}

\textbf{Конструкция (слайды 483--486).} При инверсии относительно $\Gamma_0^1$ дуга $\gamma_0^1$ остаётся на месте (поточечно), $\TT$ и $\DD$
переходят в себя (Лемма 42.1\,c), точка $a_0^3$ переходит в точку $a_1^1\in\TT$, а $T_0$ --- в треугольник $T_1^1=\sigma_{\Gamma_0^1}(T_0)$ с вершинами
$a_0^1,a_1^1,a_0^2$, ограниченный дугами $\gamma_0^1$, $\gamma_1^1=\sigma_{\Gamma_0^1}(\gamma_0^3)$ и $\gamma_1^2=\sigma_{\Gamma_0^1}(\gamma_0^2)$, ортогональными $\TT$
(Лемма 42.1\,b). Аналогично $T_1^2=\sigma_{\Gamma_0^2}(T_0)$, $T_1^3=\sigma_{\Gamma_0^3}(T_0)$. Стороны треугольников $T_1^j$, отличные от $\gamma_0^j$,
называются \emph{свободными}; это шесть дуг $\gamma_1^k$, $k=1,\dots,6$ (нумерация --- против часовой стрелки от $a_0^1$).

Пусть сделан $n$-й шаг: построены треугольники $T_n^j$, $j=1,\dots,3\cdot2^{n-1}$, со свободными сторонами $\gamma_n^{2j-1},\gamma_n^{2j}$
(дуги окружностей $\Gamma_n^{2j-1},\Gamma_n^{2j}\perp\TT$). Положим
\[
T_{n+1}^{2j-1}=\sigma_{\Gamma_n^{2j-1}}(T_n^j),\qquad T_{n+1}^{2j}=\sigma_{\Gamma_n^{2j}}(T_n^j);
\]
свободные стороны $T_{n+1}^i$ --- две его стороны, отличные от дуги, по которой производилось отражение. Пусть $P_1^j=T_0\sqcup\gamma_0^j\sqcup T_1^j$,
$D_1=\bigcup_{j=1}^3P_1^j$ и
\[
P_{n+1}^{2j-1}=T_n^j\sqcup\gamma_n^{2j-1}\sqcup T_{n+1}^{2j-1},\quad P_{n+1}^{2j}=T_n^j\sqcup\gamma_n^{2j}\sqcup T_{n+1}^{2j},\quad
D_{n+1}=D_n\cup\bigcup_{k=1}^{3\cdot2^n}P_{n+1}^k .
\]
Треугольник $T_n^j$ называем треугольником \emph{$n$-го поколения}; $T_0$ --- нулевого. Для свободной дуги $\gamma_n^k$ \emph{шапочкой}
$Q_n^k$ назовём компоненту $\DD\setminus\gamma_n^k$, не содержащую $T_0$ (это часть $\DD$, отсекаемая дугой $\gamma_n^k$).

Слайды молчаливо предполагают, что треугольники не налегают друг на друга, а <<пары>> $P_n^k$ --- области; кроме того, в качестве
упражнения оставлено равенство $\bigcup_nD_n=\DD$. Всё это --- содержание следующей леммы (самое тонкое место конструкции).

\begin{thmbox}{Лемма 42.3 (геометрия разбиения; упражнение <<$D=\DD$>>)}
\begin{enumerate}[label=\alph*)]
\item Все треугольники всех поколений попарно не пересекаются. Треугольник $T_n^j$ лежит вне круга, ограниченного окружностью любой его
свободной стороны $\gamma_n^k$, а $T_{n+1}^k\subset Q_n^k$. Шапочки одного поколения попарно не пересекаются.
\item При каждом $n$: $\ \DD=\bigcup_{m\le n}\bigcup_j\clos{T_m^j}^{\,\DD}\ \cup\ \bigsqcup_kQ_n^k$, где $\clos{T}^{\,\DD}=T\cup(\text{три стороны }T)$ ---
замыкание в $\DD$; шапочки $Q_n^k$ не пересекаются с треугольниками поколений $\le n$ и их сторонами.
\item Существует абсолютная постоянная $L$ такая, что $\diam Q_n^k\le L/(n+1)$ для всех $n,k$.
\item Каждая точка $\zeta\in\TT$ при каждом $n$ лежит в замыкании некоторой шапочки $Q_n^k$.
\item Следовательно, каждая точка $\DD$ лежит в некотором треугольнике или на некоторой его стороне; каждый компакт $K\subset\DD$ пересекается
лишь с конечным числом треугольников; каждая пара $P_n^k$ --- область, и $D:=\bigcup_nD_n=\DD$.
\end{enumerate}
\end{thmbox}

\begin{proof}
\textbf{Шаг 1: сведение к трети круга.} Поворот $\rho(z)=\omega z$ --- ДЛО, $\rho(T_0)=T_0$, $\rho(\Gamma_0^k)=\Gamma_0^{k+1}$ (индексы по модулю 3; окружность,
ортогональная $\TT$ и проходящая через две данные точки $\TT$, единственна). По Лемме~42.1\,a) $\rho\,\sigma_\Gamma\,\rho^{-1}=\sigma_{\rho(\Gamma)}$, поэтому по
индукции $\rho$ переводит треугольники $n$-го поколения, происходящие от $T_1^k$ (<<потомки>> $T_1^k$), и их свободные стороны в потомков $T_1^{k+1}$
того же поколения и их свободные стороны; шапочки --- в шапочки. Поскольку $\rho$ --- евклидова изометрия, достаточно изучить потомков $T_1^1$,
т.\,е. то, что происходит в шапочке $Q_0^1$; потомки $T_1^2$, $T_1^3$ лежат в $Q_0^2=\rho(Q_0^1)$, $Q_0^3=\rho^2(Q_0^1)$.

\textbf{Шаг 2: переход в $\CC_+$ и обозначения.} По Лемме 42.1\,a) $\Phi\circ\sigma_\Gamma=\sigma_{\Phi(\Gamma)}\circ\Phi$, так что $\Phi$ переводит всю конструкцию
в аналогичную конструкцию в $\CC_+$, начинающуюся с $T_0'=\Phi(T_0)$ и отражений относительно геодезических. По детали предыдущего пункта
$\Phi(Q_0^1)=C(0,1)$, $\Phi(Q_0^2)=\{\Re w>1\}\cap\CC_+$, $\Phi(Q_0^3)=\{\Re w<0\}\cap\CC_+$, где для $x<y$
\[
S(x,y)=\{w\in\CC_+:|w-\tfrac{x+y}2|=\tfrac{y-x}2\},\qquad C(x,y)=\{w\in\CC_+:|w-\tfrac{x+y}2|<\tfrac{y-x}2\},
\]
$\clos B_{xy}$ --- замкнутый круг с диаметром $[x,y]$. Для геодезической $g$ с концами $\alpha,\beta\in\widehat\RR$ множество $\CC_+\setminus g$ состоит из двух
<<полуплоскостей>> (внутренность и внешность полукруга, либо левая и правая части); их замыкания в $\Ch$ пересекают $\widehat\RR$ по двум замкнутым дугам
между $\alpha$ и $\beta$. Через $\Pi(g;\xi)$ ($\xi\in\widehat\RR\setminus\{\alpha,\beta\}$) обозначим ту из них, чья дуга содержит $\xi$. \emph{Идеальным
треугольником} с вершинами $\alpha,\beta,\xi$ назовём
\[
\Delta\langle\alpha,\beta,\xi\rangle=\Pi(g_{\alpha\beta};\xi)\cap\Pi(g_{\beta\xi};\alpha)\cap\Pi(g_{\xi\alpha};\beta)
\]
($g_{\alpha\beta}$ --- геодезическая с концами $\alpha,\beta$). Например, $T_0'=\Delta\langle0,1,\infty\rangle$, а при $x<c<y$
\[
\Delta\langle x,c,y\rangle=C(x,y)\setminus\bigl(\clos B_{xc}\cup\clos B_{cy}\bigr).
\]
Если $A$ --- конформный автоморфизм $\CC_+$ (вещественное ДЛО) или отражение $\sigma_g$ относительно геодезической, то $A$ --- гомеоморфизм $\Ch$,
переводящий $\CC_+$ в $\CC_+$ и $\widehat\RR$ в $\widehat\RR$, поэтому $A$ переводит компоненты $\CC_+\setminus g$ в компоненты $\CC_+\setminus A(g)$ вместе с их
дугами на бесконечности: $A(\Pi(g;\xi))=\Pi(A(g);A(\xi))$, а значит,
\begin{equation}\label{t42:inv}
A\bigl(\Delta\langle\alpha,\beta,\xi\rangle\bigr)=\Delta\langle A\alpha,A\beta,A\xi\rangle .
\end{equation}
\emph{Элементарный факт}: при $x<c<y$ круги $\clos B_{xc}$, $\clos B_{cy}$ касаются изнутри круга $\clos B_{xy}$ (в точках $x$ и $y$) и касаются друг друга в
точке $c$ (расстояния между центрами равны разности и сумме радиусов). Поэтому
\begin{equation}\label{t42:dec}
C(x,y)\cup S(x,y)=\bigl(\Delta\langle x,c,y\rangle\cup S(x,y)\cup S(x,c)\cup S(c,y)\bigr)\sqcup C(x,c)\sqcup C(c,y),
\end{equation}
где первое слагаемое --- замыкание $\Delta\langle x,c,y\rangle$ в $\CC_+$.

\textbf{Шаг 3: индукция «по Фарею».} Докажем индукцией по $m\ge0$ утверждение $(\mathrm F_m)$: \emph{существует набор $\mathcal I_m$ из $2^m$ отрезков,
покрывающих $[0,1]$ и имеющих попарно непересекающиеся внутренности, такой, что}
\begin{enumerate}[label=(\arabic*)]
\item образы при $\Phi$ свободных сторон $m$-го поколения, лежащих в $\clos{Q_0^1}$ (при $m=0$ --- одной дуги $\gamma_0^1$), --- это ровно $S(I)$, $I\in\mathcal I_m$;
\item каждый $I=[x,y]\in\mathcal I_m$ имеет вид $x=p/q$, $y=r/s$, где $p,r\in\ZZ$, $q,s\in\NN$, $rq-ps=1$, $q+s\ge m+2$, и треугольник $m$-го поколения со
стороной $\Phi^{-1}(S(I))$ имеет $\Phi$-образ $\Delta\langle x,y,v\rangle$, где $v=\frac{p-r}{q-s}$ ($v=\infty$ при $q=s$);
\item $C(0,1)\cup S(0,1)=S(0,1)\cup\bigcup\{\Phi(\clos T^{\,\DD})\}\cup\bigsqcup_{I\in\mathcal I_m}C(I)$, где объединение в середине берётся по потомкам $T_1^1$
поколений $1,\dots,m$; эти треугольники попарно не пересекаются и не пересекаются с шапочками $C(I)$, $I\in\mathcal I_m$.
\end{enumerate}
\emph{База} $m=0$: $\mathcal I_0=\{[0,1]\}$, $0=\frac01$, $1=\frac11$, $rq-ps=1$, $q+s=2$, $v=\frac{0-1}{1-1}=\infty$, и $\Phi(T_0)=\Delta\langle0,1,\infty\rangle$; пункт (3) тривиален.

\emph{Переход} $m\to m+1$. Пусть $I=[x,y]\in\mathcal I_m$, $x=p/q$, $y=r/s$, $T$ --- соответствующий треугольник, $\Phi(T)=\Delta\langle x,y,v\rangle$.
Рассмотрим $M(t)=\dfrac{qt-p}{-st+r}$; его определитель $qr-ps=1>0$, так что $M\in\operatorname{Aut}(\CC_+)$, и
\[
\begin{gathered}
M(x)=0,\qquad M(y)=\infty,\\ M(v)=\frac{q(p-r)-p(q-s)}{-s(p-r)+r(q-s)}=\frac{-1}{1}=-1,\qquad M(c)=\frac{q(p+r)-p(q+s)}{-s(p+r)+r(q+s)}=1,
\end{gathered}
\]
где $c=\frac{p+r}{q+s}$ (\emph{медианта}). Геодезическая $S(x,y)$ переходит в геодезическую с концами $0,\infty$, т.\,е. в мнимую полуось, и по Лемме 42.1\,a)
$\sigma_{S(x,y)}=M^{-1}\circ s\circ M$, где $s(w)=-\bar w$. Следовательно, $\sigma_{S(x,y)}(v)=M^{-1}(1)=c$, $\sigma_{S(x,y)}$ неподвижна в $x,y$, и по \eqref{t42:inv}
\[
\Phi(T')=\sigma_{S(x,y)}(\Delta\langle x,y,v\rangle)=\Delta\langle x,c,y\rangle\qquad(T'\ \text{--- ребёнок } T \text{ через сторону }\Phi^{-1}S(I)).
\]
Так как $M$ возрастает на $\RR\setminus\{r/s\}$ и переводит $x\mapsto0$, $c\mapsto1$, $y\mapsto\infty$, то $x<c<y$, а $v\notin[x,y]$ (ибо $M(v)=-1\notin[0,+\infty]$).
Свободные стороны $T'$ --- это $\Phi^{-1}S(x,c)$ и $\Phi^{-1}S(c,y)$; положим $\mathcal I_{m+1}=\{[x,c],[c,y]:[x,y]\in\mathcal I_m\}$. Проверим (2) для $[x,c]$:
$(p+r)q-p(q+s)=rq-ps=1$, $q+(q+s)\ge(q+s)+1\ge m+3$, а третья вершина $T'$ относительно стороны $S(x,c)$ есть $y=\frac{-r}{-s}=\frac{p-(p+r)}{q-(q+s)}$ --- как и
требуется. Для $[c,y]$: $r(q+s)-(p+r)s=1$, $(q+s)+s\ge m+3$, третья вершина $x=\frac{(p+r)-r}{(q+s)-s}$. Пункт (3) следует из \eqref{t42:dec}: шапочка $C(I)$
вместе с $S(I)$ разбивается на замыкание нового треугольника $\Delta\langle x,c,y\rangle$ и две новые шапочки; новые треугольники лежат в старых
шапочках $C(I)$, которые попарно не пересекаются и не пересекаются со старыми треугольниками. Индукция завершена.

\textbf{Шаг 4: оценка шапочек.} Для $I=[p/q,r/s]\in\mathcal I_m$
\[
|I|=\frac rs-\frac pq=\frac{rq-ps}{qs}=\frac1{qs}\le\frac1{q+s-1}\le\frac1{m+1}\qquad\bigl(qs-q-s+1=(q-1)(s-1)\ge0\bigr),
\]
и $\diam C(I)=|I|\le\frac1{m+1}$. Функция $\Phi^{-1}$ голоморфна в окрестности выпуклого компакта $K_0=\clos{C(0,1)}$ (её единственный полюс $\Phi(\infty)=e^{-i\pi/3}$
лежит в $\CC_-$); при $L:=\max_{K_0}|(\Phi^{-1})'|$ интегрирование по отрезку даёт $|\Phi^{-1}(w_1)-\Phi^{-1}(w_2)|\le L|w_1-w_2|$ на $K_0$. Шапочки поколения $m$
внутри $Q_0^1$ --- это $\Phi^{-1}(C(I))$, $I\in\mathcal I_m$ (компонента $\CC_+\setminus S(I)$, не содержащая $T_0'$, есть $C(I)$). Поэтому их диаметры $\le L/(m+1)$;
по шагу 1 то же верно в $Q_0^2,Q_0^3$. Это (c).

\textbf{Шаг 5: выводы.} (a), (b) следуют из $(\mathrm F_m)$ и шага 1 (три шапочки $Q_0^k$ попарно не пересекаются и не пересекаются с $\clos{T_0}^{\,\DD}$, так как
$\CC_+=\clos{T_0'}^{\,\CC_+}\sqcup C(0,1)\sqcup\{\Re w>1\}\sqcup\{\Re w<0\}$). (d): замкнутая дуга $\TT$ от $a_0^1$ до $a_0^2$ равна $\Phi^{-1}([0,1])=\bigcup_{I\in\mathcal I_n}\Phi^{-1}(I)$,
и $\Phi^{-1}(I)\subset\clos{\Phi^{-1}(C(I))}$; остальные две дуги --- поворотом.

(e) Замыкание шапочки $Q$ содержит точки $\TT$, поэтому для $z\in Q$: $1-|z|=\dist(z,\TT)\le\diam Q$. Пусть $z\in\DD$; выберем $n$ с $L/(n+1)<1-|z|$. Тогда $z$
не лежит ни в одной шапочке $n$-го поколения и по (b) лежит в замыкании (в $\DD$) треугольника поколения $\le n$. Треугольники поколений $>n$ лежат в шапочках
поколения $n$ (по (a) и индукции), т.\,е. в кольце $\{|z|>1-L/(n+1)\}$; поэтому компакт $K\subset\{|z|\le r\}$, $r<1$, пересекается лишь с треугольниками
поколений $\le n$ при $L/(n+1)<1-r$, а их конечное число. Каждая сторона треугольника --- это либо дуга, по которой он получен отражением из родителя,
либо свободная сторона, по которой отражением получен ребёнок; в обоих случаях она лежит в некоторой паре $P$. Значит, $\bigcup_nD_n=\DD$.
\end{proof}

\begin{detail}{пара $P=T\sqcup\gamma\sqcup T^*$ --- область}
Пусть $T^*=\sigma_\Gamma(T)$, $\gamma\subset\Gamma$ --- общая сторона, $p\in\gamma$. В модели $\CC_+$ (шаг 3) $\Phi(\gamma)=S(x,y)$, $\Phi(T^*)=\Delta\langle x,c,y\rangle$;
так как замкнутые круги $\clos B_{xc}$, $\clos B_{cy}$ не содержат $\Phi(p)$, найдётся круг $U\ni\Phi(p)$, у которого $U\cap C(x,y)\subset\Phi(T^*)$. Множество
$V=U\cap\sigma(U)$ ($\sigma=\sigma_{S(x,y)}$) --- открытая окрестность $\Phi(p)$, $\sigma(V)=V$; тогда $V\cap C(x,y)\subset\Phi(T^*)$, а
$V\setminus\clos B_{xy}=\sigma(V\cap C(x,y))\subset\sigma(\Phi(T^*))=\Phi(T)$. Итак, $V\subset\Phi(P)$, т.\,е. $P$ открыто; связность $P$ очевидна
($T$, $T^*$ связны и примыкают к $\gamma$). Кроме того, $T$ и $T^*$ лежат по разные стороны от $\Gamma$ --- это и требуется в Теореме 13.1.
\end{detail}

\begin{remarkbox}{Замечание (о слайде 487)}
На слайде утверждается, что $D=\bigcup D_n$ --- односвязная область, <<так как она звёздна относительно $0$>>; звёздность не доказывается и
не очевидна. Лемма 42.3 даёт сразу $D=\DD$, откуда односвязность тривиальна. Отметим также, что вершины треугольников всех поколений
лежат в $\Phi^{-1}(\QQ\cup\{\infty\})$ (в модели $\CC_+$ вершины --- несократимые дроби, соседние вершины --- <<соседи по Фарею>>, $|rq-ps|=1$).
\end{remarkbox}

\subsection{Модулярная функция $\mu$ и её свойства}

\textbf{Построение (слайды 485--487).} Для треугольников всех поколений определим по индукции гомеоморфизмы $\mu_T\colon\clos T\to\clos{\CC_\pm}$
(замыкания в $\Ch$; $\clos T$ включает три вершины на $\TT$): $\mu_{T_0}=\mu_0$, и если $T'=\sigma_\Gamma(T)$ --- ребёнок $T$ через сторону $\gamma\subset\Gamma$, то
\[
\mu_{T'}(z)=\overline{\mu_T(\sigma_\Gamma(z))},\qquad z\in\clos{T'} .
\]
Положим $\mu(z)=\mu_T(z)$, если $z$ лежит в треугольнике $T$ или на его стороне. Обозначим $\Omega=\CC\setminus\{0,1\}$ и
$\Omega_k=(\CC\setminus\RR)\cup I_0^k$, $k=1,2,3$; пару $P=T\sqcup\gamma\sqcup T'$ назовём \emph{парой типа $k$}, если $\mu(\gamma)=I_0^k$.

\begin{thmbox}{Теорема 42.4 (модулярная функция)}
Функция $\mu\colon\DD\to\CC$ корректно определена, и:
\begin{enumerate}[label=\arabic*)]
\item для треугольника $T$ поколения $n$ отображение $\mu|_T$ конформно отображает $T$ на $\CC_+$ ($n$ чётно) или на $\CC_-$ ($n$ нечётно), а $\mu_T$
переводит три вершины $T$ в $0,1,\infty$ и три стороны --- в $I_0^1,I_0^2,I_0^3$ (по одной);
\item для пары $P=T\sqcup\gamma\sqcup T'$ типа $k$: $\mu|_P\in\mathrm{Conf}(P,\Omega_k)$ и $\mu(\sigma_\Gamma(z))=\overline{\mu(z)}$, $z\in P$;
\item $\mu\in H(\DD)$, $\mu'(z)\ne0$ в $\DD$ ($\mu$ локально конформна), $\mu(\DD)=\Omega$; $\mu$ не принимает значений $0,1$ (и $\infty$), а каждое
значение из $\Omega$ принимает в счётном множестве точек;
\item (\emph{накрытие}) каждый круг $B\subset\Omega$ лежит в некоторой $\Omega_k$, и $\mu^{-1}(B)=\bigsqcup_{P}(\mu|_P)^{-1}(B)$, где $P$ пробегает
(попарно непересекающиеся) пары типа $k$; каждое множество $(\mu|_P)^{-1}(B)$ отображается $\mu$ конформно на $B$;
\item $\TT$ --- естественная граница функции $\mu$.
\end{enumerate}
\end{thmbox}
\begin{proof}
1) Индукция по поколению. Для $T_0$ это Предложение 42.2. Пусть утверждение верно для $T$, и $T'=\sigma_\Gamma(T)$. Отображение $\sigma_\Gamma$ ---
гомеоморфизм $\clos{T'}\to\clos T$, переводящий вершины в вершины и стороны в стороны (Лемма 42.1), поэтому $\mu_{T'}$ --- гомеоморфизм $\clos{T'}$ на множество, комплексно сопряжённое к
$\clos{\CC_\pm}$, т.\,е. на $\clos{\CC_\mp}$; вершины $T'$ идут в числа, сопряжённые к $0,1,\infty$, т.\,е. снова в $0,1,\infty$, стороны --- в сопряжённые
к $I_0^k$, т.\,е. в те же интервалы. На общей стороне $\gamma$ имеем $\sigma_\Gamma=\mathrm{id}$ и $\mu_T(\gamma)\subset\RR$, поэтому $\mu_{T'}=\mu_T$ на $\gamma$. Конформность
$\mu_{T'}$ на $T'$ --- из Теоремы 13.1, применённой к $D=T$, $W=\mu(T)=\CC_\pm$, дуге $\gamma\subset\Gamma$ и интервалу $\sigma=\mu_T(\gamma)\subset\widehat\RR$
(условия проверены в Лемме 42.3 и детали о парах): продолжение $\tilde f$ на $D^*=T'$ даётся как раз формулой $\overline{\mu_T(\sigma_\Gamma(z))}$.
Это доказывает и 2), так как $\widetilde W=\CC_\pm\cup I_0^k\cup\CC_\mp=\Omega_k$.

\emph{Корректность.} По Лемме 42.3 треугольники попарно не пересекаются, и каждая точка $\DD$ лежит либо ровно в одном треугольнике, либо на стороне;
сторона $\gamma$ лежит в открытой паре $P=T\sqcup\gamma\sqcup T'$, а другие треугольники с $P$ не пересекаются, поэтому точка $\gamma$ не может лежать на
стороне какого-то третьего треугольника, а на $\gamma$ значения $\mu_T$ и $\mu_{T'}$ совпадают.

3) Каждая точка $\DD$ лежит в некоторой паре (Лемма 42.3\,e), а на паре $\mu$ конформна по 2); значит, $\mu$ голоморфна и локально инъективна,
откуда $\mu'\ne0$ (критерий локальной однолистности, билет 1). Образ: $\mu(\DD)$ состоит из образов треугольников ($\CC_+$ или $\CC_-$) и сторон
($I_0^1,I_0^2,I_0^3$), т.\,е. $\mu(\DD)=\Omega$. Каждое $w\in\CC_+$ имеет ровно один прообраз в каждом треугольнике чётного поколения, $w\in\CC_-$ ---
в каждом нечётного, $w\in I_0^k$ --- на каждой стороне, отображаемой на $I_0^k$; таких треугольников и сторон счётное число.

4) $B\cap\RR$ --- интервал (или пусто), не содержащий $0$ и $1$, поэтому он лежит в одном из $I_0^k$, и $B\subset\Omega_k$. У каждого треугольника ровно
одна сторона отображается на $I_0^k$, поэтому каждый треугольник лежит ровно в одной паре типа $k$, а каждая сторона типа $k$ --- ровно в одной
(своей) паре. Следовательно, пары типа $k$ попарно не пересекаются, а их объединение равно $\mu^{-1}(\Omega_k)$ (точка треугольника --- в паре этого
треугольника; точка $z$ стороны с $\mu(z)\in I_0^k$ --- на стороне типа $k$). Остаётся применить 2).

5) Пусть $\zeta\in\TT$ и $r>0$. Выберем $n$ с $L/(n+1)<r$ и (Лемма 42.3\,d) шапочку $Q=Q_n^k$ с $\zeta\in\clos Q$; тогда $Q\subset B(\zeta,r)$ по Лемме 42.3\,c.
В $Q$ лежат ребёнок и внук треугольника, построенного на $\gamma_n^k$ (поколения $n+1$ и $n+2$), один из них чётного поколения, и в нём есть точка
$z$ с $\mu(z)=w_*:=e^{i\pi/3}$. Итак, найдутся (лежащие в разных, т.\,е. непересекающихся, треугольниках) различные точки $z_j\to\zeta$ с $\mu(z_j)=w_*$.
Если бы $F\in H(B(\zeta,r))$ и $F=\mu$ на $B(\zeta,r)\cap\DD$, то нули $F-w_*$ накапливались бы к $\zeta\in B(\zeta,r)$, т.\,е. $F\equiv w_*$, и по теореме
единственности $\mu\equiv w_*$ в $\DD$ --- противоречие с 3).
\end{proof}

\begin{remarkbox}{Замечание}
Функция $\lambda=\mu\circ\Phi^{-1}\colon\CC_+\to\Omega$ --- это (с точностью до дробно-линейной замены переменной) классическая модулярная
функция $\lambda(\tau)$, инвариантная относительно группы $\Gamma(2)=\{\tau\mapsto\frac{a\tau+b}{c\tau+d}:\ a,b,c,d\in\ZZ,\ ad-bc=1,\ \begin{psmallmatrix}a&b\\c&d\end{psmallmatrix}\equiv
\begin{psmallmatrix}1&0\\0&1\end{psmallmatrix}\ (\mathrm{mod}\ 2)\}$; отсюда название. В модели $\CC_+$ треугольники поколений --- идеальные
треугольники <<разбиения Фарея>>. Для доказательства теоремы Пикара всё это не нужно.
\end{remarkbox}

\subsection{Обратная функция $\mu^{-1}$ как П.А.Ф.}

Отображение $\mu_0^{-1}\colon\CC_+\to T_0$ голоморфно; пусть $G_i=E_i(\mu_0^{-1})=(B(i,R_i),g_i)$ --- его канонический элемент в точке $i$.

\begin{thmbox}{Предложение 42.5 (слайд 488)}
$R_i=1$.
\end{thmbox}
\begin{proof}
Так как $B(i,1)\subset\CC_+$ и $\mu_0^{-1}\in H(\CC_+)$, ряд Тейлора сходится в $B(i,1)$: $R_i\ge1$. Пусть $R_i>1$. Функции $g_i$ и $\mu_0^{-1}$ голоморфны в
выпуклой (значит, связной) области $\CC_+\cap B(i,R_i)$ и совпадают в $B(i,1)$, поэтому совпадают в ней всей; по непрерывности $\mu_0^{-1}$ на $\clos{\CC_+}$
имеем $g_i=\mu_0^{-1}$ и на $\RR\cap B(i,R_i)\ni0$. Значит, $g_i(0)=a:=a_0^1$, $g_i(t)\in\gamma_0^1$ при малых $t>0$ и $g_i(t)\in\gamma_0^3$ при малых $t<0$
(так как $\mu_0(\gamma_0^1)=(0,1)$, $\mu_0(\gamma_0^3)=(-\infty,0)$), а $|g_i(w)|<1$ при $w\in\CC_+$ вблизи $0$ (ибо $\mu_0^{-1}(\CC_+)=T_0\subset\DD$).
Функция $g_i$ непостоянна; запишем $g_i(w)=a\bigl(1+cw^p+O(w^{p+1})\bigr)$, $c\ne0$, $p\ge1$.

\emph{Шаг 1: $p=1$.} При $w=\rho e^{i\theta}$, $\theta\in(0,\pi)$ фиксировано, $\rho\to0+$:
$|g_i(w)|^2=1+2\rho^p\Re(ce^{ip\theta})+O(\rho^{p+1})<1$, откуда $\Re(ce^{ip\theta})\le0$ для всех $\theta\in(0,\pi)$. Но при $p\ge2$ аргумент
$\arg c+p\theta$ пробегает интервал длины $p\pi\ge2\pi$ и принимает значения с $\Re(ce^{ip\theta})>0$. Значит, $p=1$.

\emph{Шаг 2: противоречие.} При $p=1$ для вещественных $t\to0\pm$ имеем $\frac{g_i(t)-a}{|g_i(t)-a|}\to\pm\frac{ac}{|ac|}$ --- два \emph{противоположных}
направления. С другой стороны, $g_i(t)$ стремится к $a$ по дуге $\gamma_0^1$ (при $t>0$) и по дуге $\gamma_0^3$ (при $t<0$); для точек $z$ гладкой дуги,
стремящихся к её концу $a$, направление хорды $\frac{z-a}{|z-a|}$ стремится к касательному направлению дуги в $a$. Обе окружности $\Gamma_0^1,\Gamma_0^3$
ортогональны $\TT$ в точке $a$, т.\,е. их касательная в $a$ --- прямая $\{a+sa\}$, а так как дуги лежат в $\DD$ и $\Re\bigl((z-a)\bar a\bigr)=\Re(z\bar a)-1<0$
для $z\in\DD$, предельное направление в обоих случаях равно $-a$. Противоречие: $\frac{ac}{|ac|}=-a=-\frac{ac}{|ac|}$.
\end{proof}

\begin{remarkbox}{Замечание (пробел на слайде 488)}
На слайде сказано: бесконечно малый угол величины $\pi$ с вершиной $0$ перешёл бы в угол величины $p\pi$, а должен перейти в угол величины $0$.
Это рассуждение исключает только нечётные $p$: при чётном $p$ угол $p\pi$ <<равен $0$ по модулю $2\pi$>>, и стороны угла действительно
могли бы перейти в касающиеся дуги. Шаг 1 выше (оценка модуля) исключает все $p\ge2$.
\end{remarkbox}

Обозначим через $\mu^{-1}$ П.А.Ф., порождённую элементом $G_i$. Ключ к её описанию --- <<поднятие путей>>.

\begin{thmbox}{Лемма 42.6 (поднятие путей)}
Пусть $\Sigma\colon[0,1]\to\Omega$ --- путь и $\mu(z_0)=\Sigma(0)$, $z_0\in\DD$. Существует единственный путь $\tilde\Sigma\colon[0,1]\to\DD$ с $\tilde\Sigma(0)=z_0$ и
$\mu\circ\tilde\Sigma=\Sigma$. Если $\varphi_t$ --- локальная обратная к $\mu$ вблизи $\Sigma(t)$ с $\varphi_t(\Sigma(t))=\tilde\Sigma(t)$, то семейство $F_t=E_{\Sigma(t)}(\varphi_t)$
--- аналитическое продолжение элемента $F_0$ вдоль $\Sigma$.
\end{thmbox}
\begin{proof}
\emph{Существование.} Пусть $d=\min_t\dist(\Sigma(t),\{0,1\})>0$; по равномерной непрерывности выберем $0=t_0<\dots<t_N=1$ так, что
$\Sigma([t_{j-1},t_j])\subset B_j:=B(\Sigma(t_{j-1}),d)\subset\Omega$. Пусть $\tilde\Sigma$ уже построен на $[0,t_{j-1}]$. По Теореме 42.4\,4) точка
$\tilde\Sigma(t_{j-1})\in\mu^{-1}(B_j)$ лежит ровно в одном <<листе>> $U_j=(\mu|_P)^{-1}(B_j)$; пусть $\psi_j=(\mu|_{U_j})^{-1}\colon B_j\to U_j$ (голоморфна).
Положим $\tilde\Sigma=\psi_j\circ\Sigma$ на $[t_{j-1},t_j]$; это непрерывно и согласовано в $t_{j-1}$.

\emph{Единственность.} Для двух поднятий множество $E=\{t:\tilde\Sigma_1(t)=\tilde\Sigma_2(t)\}$ замкнуто и содержит $0$; оно открыто, так как $\mu$ инъективна в
окрестности $N$ точки $\tilde\Sigma_1(t)$, $t\in E$, и при $s$, близких к $t$, обе точки $\tilde\Sigma_{1,2}(s)\in N$ имеют один образ $\Sigma(s)$. Значит, $E=[0,1]$.

\emph{Продолжение.} Фиксируем $t_0$; пусть $B_0=B(\Sigma(t_0),\dist(\Sigma(t_0),\{0,1\}))$, $U_0$ --- лист над $B_0$, содержащий $\tilde\Sigma(t_0)$, и
$\varphi_{t_0}=(\mu|_{U_0})^{-1}$ на $B_0$. При $|t-t_0|<\delta$ имеем $\Sigma(t)\in B_0$ и $\tilde\Sigma(t)\in U_0$ (непрерывность), откуда
$\varphi_{t_0}(\Sigma(t))=\tilde\Sigma(t)$ (инъективность $\mu$ на $U_0$). Две непрерывные локальные обратные $\varphi_t$, $\varphi_{t_0}$ с одинаковым значением
$\tilde\Sigma(t)$ в точке $\Sigma(t)$ совпадают вблизи $\Sigma(t)$ ($\mu$ инъективна вблизи $\tilde\Sigma(t)$), а значит, по теореме единственности, и на
выпуклом пересечении кругов элементов $F_t$, $F_{t_0}$. Это и есть определение продолжения вдоль пути.
\end{proof}

\begin{thmbox}{Предложение 42.7 (упражнение со слайдов 488--489)}
\begin{enumerate}[label=\alph*)]
\item Элемент $G_i$ продолжается вдоль любого пути $\Sigma\subset\Omega$ с началом $i$; элементы П.А.Ф. $\mu^{-1}$ --- это в точности канонические
элементы $E_w(\varphi)$ всевозможных локальных обратных $\varphi$ к $\mu$ (в частности, все значения $\mu^{-1}$ лежат в $\DD$).
\item $\mu^{-1}$ не продолжается вдоль путей, приходящих в $0$ или $1$; таким образом, $\Omega$ --- область определения $\mu^{-1}$, и все точки
$b\in\Omega$ правильные.
\item $0,1,\infty$ --- логарифмические точки ветвления: над проколотой окрестностью каждой из них $\mu^{-1}$ имеет счётное число
аналитических ветвей, и каждая из них имеет там логарифмическое ветвление.
\end{enumerate}
\end{thmbox}
\begin{proof}
a) Первое --- Лемма 42.6 с $z_0=\mu_0^{-1}(i)$ (на $B(i,1)$ функция $g_i=\mu_0^{-1}$ --- локальная обратная). Обратно, для любого $z\in\DD$ соединим $z_0$ с $z$
путём $\zeta\subset\DD$; тогда $\zeta$ --- поднятие пути $\mu\circ\zeta\subset\Omega$, и продолжение $G_i$ вдоль $\mu\circ\zeta$ приходит к элементу локальной
обратной со значением $z$ в центре. Правильность точек $b\in\Omega$: над кругом $B(b,\delta)\subset\Omega$ ветвь распадается на голоморфные ветви
$(\mu|_U)^{-1}$, $U$ --- листы над $B(b,\delta)$.

c) Пусть $0<r<1$, $b=ir$, $c(s)=ire^{2\pi is}$, $s\in[0,1]$ (обход $0$ против часовой стрелки по окружности $|w|=r$, не задевающий $1$), и $z_b\in\mu^{-1}(b)$;
$z_b$ лежит в треугольнике $T$ чётного поколения. Пусть $v$ --- вершина $T$ с $\mu_T(v)=0$; стороны $T$, выходящие из $v$, отображаются на интервалы с концом $0$,
т.\,е. на $I_0^3$ и $I_0^1$. Поднятие $\tilde c$ пути $c$ из $z_b$: при $s\in[0,\frac14)$ путь $c$ лежит в $\CC_+$, и $\tilde c$ остаётся в $T$ (связная
часть $\mu^{-1}(\CC_+)$, т.\,е. объединения непересекающихся треугольников чётных поколений); при $s=\frac14$ имеем $c=-r\in I_0^3$, и $\tilde c(\frac14)$
лежит на стороне $T$ типа $3$; при $s\in(\frac14,\frac34)$ $\tilde c$ лежит в соседнем через эту сторону треугольнике $T_1$; при $s=\frac34$ ($c=r\in I_0^1$) ---
на стороне $T_1$ типа $1$; затем --- в соседнем треугольнике $T_2$. Поскольку $\mu$ на общей стороне и в её концах одинакова у соседей, $v$ --- вершина
$T_1$ и $T_2$ с $\mu=0$, а $T_{j+1}$ получается из $T_j$ отражением относительно \emph{второй} стороны, выходящей из $v$ (не общей с $T_{j-1}$).
$n$ обходов приводят в треугольник $T_{2n}$.

Покажем, что $T_{2n}\ne T$ при $n\ge1$. Пусть $A\in\operatorname{Aut}(\CC_+)$ переводит $\Phi(v)\in\widehat\RR$ в $\infty$ ($A(t)=-1/(t-\Phi(v))$, или $A=\mathrm{id}$),
$\Psi=A\circ\Phi$. Тогда $\Psi$ переводит треугольники с вершиной $v$ в идеальные треугольники $\Delta\langle x,y,\infty\rangle=\{x<\Re w<y\}\setminus\clos B_{xy}$
(вертикальные полуполосы), стороны через $v$ --- в вертикальные прямые, а отражение относительно прямой $\Re w=y$ переводит
$\Delta\langle x,y,\infty\rangle$ в $\Delta\langle y,2y-x,\infty\rangle$ (по \eqref{t42:inv}). Поэтому $\Psi(T_j)$ --- полуполосы одной ширины $d$, каждая следующая
сдвинута на $d$ в одну и ту же сторону; все они различны. Итак, после $n\ge1$ обходов поднятие заканчивается в точке $\ne z_b$, т.\,е. продолжение
элемента $E_b(\varphi)$, $\varphi(b)=z_b$, вдоль $c^n$ даёт элемент с другим значением в центре: ветвление логарифмическое. Для точки $1$ (окружность
$|w-1|=r<1$, интервалы $I_0^1,I_0^2$) и для $\infty$ (окружность $|w|=R>1$, интервалы $I_0^2,I_0^3$) рассуждение то же.
Аналитическая ветвь над $\{0<|w|<r\}$ определяется вершиной $v$ (при любых путях в проколотом круге поднятие пересекает лишь стороны через $v$),
а у разных треугольников <<веера>> вершины $v$ эта вершина общая; вершин с $\mu=0$ бесконечно много (в каждой шапочке есть целый треугольник),
но счётное число --- отсюда счётное число логарифмических ветвлений над $0$.

b) Пусть $G_i$ продолжается вдоль пути $\sigma$, $\sigma([0,1))\subset\Omega$, $\sigma(1)=0$, и $F_1=(B(0,R),g)$. При $t$, близких к $1$, $F_t=E_{\sigma(t)}(g)$;
возьмём такое $t$ с $0<\rho:=|\sigma(t)|<\min(R,1)$, $b'=\sigma(t)$. Путь $\ell$: дуга окружности $|w|=\rho$ от $b'$ до $i\rho$, затем $c$ (с $r=\rho$),
затем та же дуга обратно. Он лежит в $B(0,R)\setminus\{0\}$, где $g$ голоморфна, поэтому продолжение $E_{b'}(g)$ вдоль $\ell$ возвращает $E_{b'}(g)$.
С другой стороны, поднятие $\ell$ из $g(b')$: первая дуга приводит в точку $z'$, обход $c$ --- в точку $z''\ne z'$ (см. доказательство пункта c)), обратная дуга --- в точку $z'''$.
Если $z'''=g(b')$, то обращённое поднятие третьей части --- поднятие первой дуги из $g(b')$ с концом $z''$, и по единственности $z''=z'$ ---
противоречие. Значит, значение в центре изменилось --- противоречие. Для точки $1$ аналогично.
\end{proof}

\subsection{Малая теорема Пикара}

\begin{thmbox}{Теорема 17.1 (малая теорема Пикара для целых функций)}
Пусть $f\in H(\CC)$ и $f\not\equiv\mathrm{const}$. Тогда $f$ принимает (в $\CC$) все значения из $\CC$, кроме, быть может, одного.
\end{thmbox}
\begin{proof}
Пусть, от противного, $f$ непостоянна и не принимает двух значений $a\ne b$. Тогда $f_1=(f-a)/(b-a)\in H(\CC)$ непостоянна и не принимает значений $0$
и $1$, т.\,е. $f_1(\CC)\subset\Omega$. Будем считать $f_1=f$.

\emph{Шаг 1: $f(\CC)\cap B(i,1)\ne\varnothing$.} Иначе $|f(z)-i|\ge1$ на $\CC$, и $1/(f-i)\in H(\CC)$ ограничена единицей; по теореме Лиувилля она
постоянна, а тогда постоянна и $f$. Выберем $z_0$ с $w_0=f(z_0)\in B(i,1)$.

\emph{Шаг 2: исходные элементы.} Пусть $G_{w_0}=E_{w_0}(\mu_0^{-1})$. Так как $R_i=1$ (Предложение 42.5), $w_0$ лежит в круге элемента $G_i$, и $G_{w_0}$ получается
из $G_i$ переразложением ряда Тейлора, т.\,е. принадлежит П.А.Ф. $\mu^{-1}$. По Лемме 42.6 (функция $\mu_0^{-1}$ --- локальная обратная к $\mu$) $G_{w_0}$
продолжается вдоль любого пути в $\Omega$ с началом $w_0$, причём продолжающее семейство состоит из элементов локальных обратных к $\mu$; в частности,
значения всех функций этих элементов \emph{в центрах} лежат в $\DD$.

\emph{Шаг 3: продолжение композиции.} Пусть $\sigma\colon[0,1]\to\CC$ --- любой путь, $\sigma(0)=z_0$. Элемент $E_{z_0}(f)$ продолжается вдоль $\sigma$ (семейство
$F_t=E_{\sigma(t)}(f)$, так как $f$ --- целая). Путь $\Sigma(t)=f(\sigma(t))$ лежит в $\Omega$ и начинается в $w_0$, поэтому $G_{w_0}$ продолжается вдоль $\Sigma$
семейством $G_t=E_{\Sigma(t)}(\varphi_t)$, где $\varphi_t$ --- локальная обратная к $\mu$. По теореме о продолжении сложного элемента (Предложение 15.6, билеты 25--27)
элемент $H_{z_0}=E_{z_0}(\mu_0^{-1}\circ f)$ продолжается вдоль $\sigma$, и продолжающее семейство --- $H_t=E_{\sigma(t)}(\varphi_t\circ f)$.

\emph{Шаг 4: монодромия и Лиувилль.} Итак, $H_{z_0}$ продолжается вдоль любого пути в односвязной области $\CC$. По теореме о монодромии
(Теорема 15.2, билет 26) результат продолжения в точку $b$ не зависит от пути; обозначим его $(B_b,h_b)$ и положим $h(b)=h_b(b)$.
\begin{detail}{$h$ --- целая функция, и $E_b(h)=(B_b,h_b)$}
Если $b'\in B_b$, то, продолжив $H_{z_0}$ сначала в $b$, а затем по отрезку $[b,b']\subset B_b$, получим элемент $E_{b'}(h_b)$ (продолжение элемента вдоль
пути внутри его круга даётся переразложениями той же функции). По независимости от пути $h(b')=h_b(b')$ при $b'\in B_b$, т.\,е. $h=h_b$ в $B_b$, и $h$ голоморфна.
\end{detail}
При этом $h(b)=\varphi_1(f(b))\in\DD$ (значение функции элемента $H_1$ в центре), т.\,е. $|h|<1$ в $\CC$. По теореме Лиувилля $h\equiv C$. Вблизи $z_0$
имеем $h=\mu_0^{-1}\circ f$, откуда $f=\mu_0\circ h\equiv\mu_0(C)$ в окрестности $z_0$, и по теореме единственности $f$ постоянна в $\CC$ --- противоречие.
\end{proof}

\begin{remarkbox}{Замечания к доказательству}
\begin{itemize}
\item Шаг 1 нужен лишь для удобства (чтобы стартовать с ветви $\mu_0^{-1}$); по Предложению 42.7\,a) можно взять любое $z_0$ и любой элемент
$\mu^{-1}$ в точке $f(z_0)$.
\item На слайде 490 сказано, что значения всех функций продолжений лежат в <<$D=\DD$>>. Для доказательства достаточно $D\subset\DD$, т.\,е.
упражнение $D=\DD$ (Лемма 42.3) здесь не используется; однако нужны непересекаемость треугольников и свойство накрытия (Теорема 42.4\,4)),
которые тоже входят в Лемму 42.3.
\item На языке топологии: $\mu\colon\DD\to\Omega$ --- накрытие (даже универсальное), а $h$ --- поднятие $f\colon\CC\to\Omega$ на $\DD$. Для дважды
проколотой плоскости накрывающая --- круг, и поднятие ограничено; для один раз проколотой $\CC\setminus\{0\}$ накрытие $\exp\colon\CC\to\CC\setminus\{0\}$
даёт лишь $f=e^g$ без противоречия --- поэтому одно значение пропустить можно, а два нельзя.
\end{itemize}
\end{remarkbox}

\begin{thmbox}{Теорема 17.2 (малая теорема Пикара для мероморфных функций)}
Пусть $f$ мероморфна в $\CC$ и $f\not\equiv\mathrm{const}$. Тогда $f$ принимает (в $\CC$) все значения из $\Ch$, кроме, быть может, двух.
\end{thmbox}
\begin{proof}
Пусть $f$ не принимает трёх различных значений $a,b,c\in\Ch$. Если одно из них, скажем $c$, равно $\infty$, то у $f$ нет полюсов, $f\in H(\CC)$ непостоянна
и не принимает двух конечных значений $a\ne b$ --- противоречие с Теоремой 17.1. Пусть $a,b,c\in\CC$. Функция $f_1=1/(f-a)$ голоморфна вне полюсов
$f$ (там $f\ne a$, так что знаменатель не обращается в нуль), а в полюсе $z_1$ функции $f$ имеем $f_1(z)\to0$, т.\,е. $z_1$ --- устранимая особенность
(нуль $f_1$). После устранения особенностей $f_1\in H(\CC)$, $f_1$ непостоянна (иначе постоянна $f=a+1/f_1$), и $f_1$ не принимает значений
$\frac1{b-a}\ne\frac1{c-a}$ (если $f_1(z)=\frac1{b-a}$, то $z$ не полюс и $f(z)=b$). Противоречие с Теоремой 17.1.
\end{proof}

\begin{remarkbox}{Примеры (точность теорем)}
$e^z$ пропускает одно значение $0$ (Теорема 17.1 точна); $\operatorname{tg}z=\frac{e^{2iz}-1}{i(e^{2iz}+1)}$ мероморфна и пропускает $\pm i$, а $1/(1+e^z)$
пропускает $0$ и $1$ --- Теорема 17.2 точна. Многочлен степени $\ge1$ принимает все значения из $\CC$.
\end{remarkbox}

\subsection{Связь с доказательством из главы 12}

\begin{remarkbox}{Теорема 12.1 и её доказательство для класса $\mathcal E$ (слайды 268--271)}
В главе 12 малая теорема Пикара доказана лишь для целых функций класса $\mathcal E=\bigcup_n\mathcal E^{(n)}$, где
$\mathcal E^{(n)}=\{f:\ \limsup_{r\to\infty}\log^{(n)}M_f(r)/\log r<\infty\}$, $\log^{(1)}=\log$, $\log^{(n+1)}=\log\log^{(n)}$ ($\mathcal E^{(1)}$ --- многочлены,
$\mathcal E^{(2)}$ --- функции конечного порядка). Схема: индукция по $n$. Если $f\in\mathcal E^{(N+1)}$ и $f\ne0$, то $f=e^g$ (Лемма 12.3, существование
логарифма в односвязной области), $\log M_f(r)=A_g(r):=\max_{|z|=r}\Re g$, а лемма Бореля--Каратеодори при $R=2r$ даёт
$M_g(r)\le2A_g(2r)+3|g(0)|\le3A_g(2r)$ при больших $r$, откуда $g\in\mathcal E^{(N)}$. По предположению индукции $g$ пропускает не более одного
значения, поэтому для любого $w$ хотя бы одно (на самом деле бесконечно много) из чисел $w+2\pi ik$ принимается $g$, и $f$ принимает
$e^w$. Значит, $f$ принимает все значения, кроме $0$; общий случай --- сдвигом $f-a$.

\emph{Сравнение.} Это рассуждение использует оценки роста и <<разворачивает>> только один прокол (через $\exp$); оно не покрывает функций,
растущих быстрее любой итерированной экспоненты (такие целые функции существуют). Доказательство через модулярную функцию вовсе не
использует рост $f$: вся информация о <<пропуске двух значений>> упакована в накрытие $\mu\colon\DD\to\CC\setminus\{0,1\}$, после чего работает
одна теорема Лиувилля.
\end{remarkbox}

\begin{remarkbox}{Коротко: план ответа на экзамене}
\begin{enumerate}
\item Точки $a_0^1,a_0^2=i,a_0^3$ на $\TT$; дуги окружностей $\perp\TT$ через пары точек; треугольник $T_0$ с нулевыми углами.
\item Теорема Римана + Каратеодори: $\mu_0\colon\clos{T_0}\to\clos{\CC_+}$, вершины $\mapsto0,1,\infty$, стороны $\mapsto(0,1),(1,\infty),(-\infty,0)$.
\item Принцип симметрии (сведение к прямой ДЛО): отражения треугольников по сторонам, продолжение $\mu$ на пары
$P\to(\CC\setminus\RR)\cup I_0^k$; поколения $T_n^j$, области $D_n$.
\item Исчерпание: в $\CC_+$ (через $\Phi$) треугольники --- треугольники Фарея, шапочки $C([p/q,r/s])$ поколения $m$ имеют размер $\le\frac1{m+1}$;
отсюда $\bigcup D_n=\DD$ и локальная конечность.
\item Свойства $\mu$: голоморфна, $\mu'\ne0$, $\mu(\DD)=\Omega$, пропускает $0,1,\infty$; накрытие; $\TT$ --- естественная граница.
\item $\mu^{-1}$: $R_i=1$; поднятие путей $\Rightarrow$ продолжение вдоль любых путей в $\Omega$; $0,1,\infty$ --- логарифмические точки ветвления.
\item Теорема 17.1: $f$ пропускает $0,1$ $\Rightarrow$ $\mu^{-1}\circ f$ продолжается вдоль всех путей $\Rightarrow$ (монодромия) целая $h$, $|h|<1$
$\Rightarrow$ (Лиувилль) $h=\mathrm{const}\Rightarrow f=\mathrm{const}$.
\item Теорема 17.2: сведение к 17.1 заменой $1/(f-a)$.
\end{enumerate}
\end{remarkbox}

\end{document}
